% ======================================================================
%  Глава 1. Физическая кристаллография
% ======================================================================
\chapter{Физическая кристаллография}\label{ch:cryst}
% ----------------------------------------------------------------------
%  Введение к главе 1
% ----------------------------------------------------------------------
Кристаллография~— язык, на котором говорит всё материаловедение. Чтобы
описать дислокацию (гл.~2), зону Бриллюэна или анизотропию свойств
(гл.~3), текстуру прокатанного листа или ориентационное соотношение при
мартенситном превращении, нужно уметь задавать направления и плоскости в
кристалле, вычислять углы и расстояния между ними, изображать их на
проекциях и понимать, какие ограничения накладывает на свойства
симметрия. Современная кристаллография включает теорию симметрии,
методы определения структуры, изучение реальной структуры, роста и
физических свойств кристаллов. Задача этой главы~— научить
\emph{кристаллографическому подходу} к анализу эксперимента.

\begin{figure}[ht]
\centering
\begin{tikzpicture}[
  box/.style={draw=#1, fill=#1!8, rounded corners=2pt, align=center,
    font=\small, inner sep=4pt, text width=38mm, minimum height=11mm, line width=0.6pt},
  flow/.style={-{Stealth[length=2mm]}, line width=0.6pt, fmGray}]
  \node[box=fmBlue] (a) at (0,0) {\textbf{1.1} Кристаллическое\\состояние};
  \node[box=fmBlue] (b) at (0,-1.6) {\textbf{1.2} Основы\\кристаллографии};
  \node[box=fmGreen] (c) at (-4.6,-3.6) {\textbf{1.3} Симметрия\\кристаллов};
  \node[box=fmGreen] (d) at (-4.6,-5.3) {\textbf{1.4} Структура\\кристаллов};
  \node[box=fmAmber] (e) at (0,-3.6) {\textbf{1.5} Физические\\свойства кристаллов};
  \node[box=fmRed] (f) at (4.6,-3.6) {\textbf{1.6} Кристаллография\\пластической деформации};
  \node[box=fmRed] (g) at (4.6,-5.3) {\textbf{1.7} Границы зёрен};
  \node[box=fmRed] (h) at (4.6,-7.0) {\textbf{1.8} Мартенситные\\превращения};
  \draw[flow] (a) -- (b);
  \draw[flow] (b.west) -| (c.north); \draw[flow] (c) -- (d);
  \draw[flow] (b) -- (e); \draw[flow] (c.east) -- (e.west);
  \draw[flow] (b.east) -| (f.north); \draw[flow] (f) -- (g); \draw[flow] (g) -- (h);
\end{tikzpicture}
\caption{Карта главы. Аппарат кристаллографии (синий блок) нужен для
описания симметрии и структуры (зелёный), физических свойств (жёлтый) и
реальных процессов в металлах~— деформации, границ, превращений
(красный)}
\label{fig:ch1-map}
\end{figure}

\begin{vnimanie}[Как читать эту главу]
Глава насыщена геометрией и матрицами. Все матричные формулы
дублируются простыми частными случаями для кубической решётки~— на
первом проходе можно ограничиться ими. Выводы, нужные лишь для
вычислений «в общем виде» (косоугольные системы, метрические матрицы),
вынесены в блоки «Подробнее».
\end{vnimanie}


\section{Кристаллическое состояние}\label{sec:cryst-state}
% ----------------------------------------------------------------------
%  1.1. Кристаллическое состояние
% ----------------------------------------------------------------------
\begin{orientir}
Что такое кристалл с точки зрения внешнего вида (законы огранки), с
точки зрения свойств (однородность, анизотропия, симметрия) и с точки
зрения внутреннего строения (решётка и структура). Главный итог:
\vazhno{структура = решётка + базис}.
\end{orientir}

\begin{opredelenie}[Кристалл]
Твёрдое тело с упорядоченной трёхмерно-периодической атомной
структурой. Кристаллическое состояние~— термодинамически равновесное
состояние твёрдого тела: каждой твёрдой фазе данного состава при данных
внешних условиях отвечает одна определённая кристаллическая структура.
\end{opredelenie}

Естественные плоские грани кристалла (огранка)~— внешнее проявление его
внутреннего атомного строения. Именно изучая огранку, кристаллографы
пришли к идее периодичности задолго до того, как удалось «увидеть» атомы.

\subsection{Основные законы кристаллографии}\label{ssec:laws}

Законы кристаллографии получены из наблюдений за огранкой, поэтому их
называют геометрическими (эмпирическими). Внешний облик кристаллов~—
\term{габитус}: размеры, форма и даже число граней~— у разных кристаллов
одного вещества может быть разным. Но их всегда можно расположить так,
что некоторые грани окажутся параллельными; такие грани называют
\term{соответственными}.

\begin{glavnoe}[Закон постоянства углов (Н.~Стенон, 1669)]
Углы между соответственными гранями кристаллов данного вещества
постоянны и характерны для этого вещества.
\end{glavnoe}

Стенон установил закон на кварце и гематите. Форма граней может сильно
меняться в зависимости от условий роста (рис.~\ref{fig:steno}), но углы
остаются прежними~— поэтому по углам, измеренным оптическим
\emph{гониометром}, определяют вещество.

\begin{figure}[ht]
\centering
\begin{tikzpicture}[scale=0.95]
  \def\hexR{1.2}
  \draw[fmBlue, thick, fill=fmBlue!8] (0:\hexR) \foreach \a in {60,120,...,300}{-- (\a:\hexR)} -- cycle;
  \node[font=\small] at (0,-1.55) {\textit{а}};
  \begin{scope}[xshift=4.3cm]
    \coordinate (P1) at (0:1.9); \coordinate (P2) at ($(P1)+(120:0.7)$);
    \coordinate (P3) at ($(P2)+(180:2.4)$); \coordinate (P4) at ($(P3)+(240:1.6)$);
    \coordinate (P5) at ($(P4)+(300:0.7)$); \coordinate (P6) at ($(P5)+(0:2.4)$);
    \draw[fmBlue, thick, fill=fmBlue!8] (P1) -- (P2) -- (P3) -- (P4) -- (P5) -- (P6) -- cycle;
    \foreach \p in {P1,P2,P3,P4,P5,P6}{\fill[fmRed] (\p) circle (1.2pt);}
    \node[font=\scriptsize, fmRed] at ($(P2)+(-0.15,-0.35)$) {$120^\circ$};
    \node[font=\scriptsize, fmRed] at ($(P5)+(0.15,0.35)$) {$120^\circ$};
    \node[font=\small] at (0.3,-1.55) {\textit{б}};
  \end{scope}
\end{tikzpicture}
\caption{Закон постоянства углов: поперечные сечения двух кристаллов
кварца~— «идеального» (\textit{а}) и искажённого при росте (\textit{б}).
Длины граней разные, но углы между соответственными гранями одинаковы
(здесь все по $120^\circ$)}
\label{fig:steno}
\end{figure}

\begin{glavnoe}[Закон рациональных параметров (Р.\,Ж.~Гаюи, 1784)]
Если за оси координат принять три ребра кристалла, то отрезки, которые
любая грань отсекает на осях, относятся как \emph{целые} числа, если их
выразить в некоторых осевых единицах $a$, $b$, $c$ (свои для каждого
вещества).
\end{glavnoe}

Существование осевых единиц, одинаковых для всех граней, означает, что
кристалл построен из одинаковых «кирпичиков»: вместе с представлением об
атомах закон Гаюи с необходимостью ведёт к выводу о \vazhno{трёхмерной
периодичности атомного строения}.

Дальнейшие вехи:
\begin{itemize}
\item Е.\,С.~Фёдоров (1890) и А.~Шёнфлис (1891) теоретически вывели все
\emph{230 пространственных групп} симметрии периодических структур
(п.~1.3.6);
\item первым прямым доказательством существования пространственной
решётки стала \emph{дифракция рентгеновских лучей} на кристаллах
(М.~Лауэ, В.~Фридрих, П.~Книппинг, 1912);
\item в середине XX~в. понятие симметрии обобщили: в группах
\emph{антисимметрии} (А.\,В.~Шубников) и \emph{цветной симметрии}
(Н.\,В.~Белов) кроме трёх пространственных координат есть четвёртая
«координата» иной физической природы~— время, направление спина, знак
заряда. Такие группы (шубниковские и беловские) описывают, например,
магнитные структуры.
\end{itemize}

\subsection{Макроскопические характеристики кристаллов}\label{ssec:macro}

С макроскопической точки зрения кристалл обладает тремя главными
особенностями: однородностью, анизотропией и симметрией. Пусть $F$~—
какое-то свойство: скалярное (плотность, теплоёмкость), векторное
(поляризация) или в общем случае тензорное (упругость).

\paragraph{Однородность}
Свойства во всех участках кристалла одинаковы:
\begin{equation}\label{eq:homog}
F(\vb x)=F(\vb x+\vb x') ,
\end{equation}
т.\,е. свойства не меняются при произвольном переносе начала координат.
Это позволяет рассматривать кристалл как \emph{сплошную среду} и описывать
многие свойства феноменологически, не вспоминая об атомах. Однородны,
впрочем, и жидкости, и газы, и аморфные тела.

\paragraph{Анизотропия}
Отличает кристалл от этих сред зависимость свойств от направления $\vb n$,
вдоль которого их измеряют:
\begin{equation}\label{eq:anis}
F(\vb n_1)\ne F(\vb n_2) .
\end{equation}
Анизотропны также текстуры, жидкие кристаллы, полимеры~— анизотропия
всегда определяется строением вещества.

\begin{primer}[Анизотропия упругости железа]
Модуль Юнга монокристалла железа вдоль ребра куба
$E_{\nfam{100}}\approx131$~ГПа, а вдоль пространственной диагонали
$E_{\nfam{111}}\approx283$~ГПа~— более чем вдвое больше
(рис.~\ref{fig:young-polar}). Для сравнения: вольфрам упруго почти
изотропен ($E\approx410$~ГПа во всех направлениях), хотя тоже имеет
кубическую решётку.
\end{primer}

\begin{figure}[ht]
\centering
\begin{tikzpicture}
\begin{axis}[width=0.62\textwidth, height=0.62\textwidth, axis equal image,
  xmin=-460, xmax=460, ymin=-460, ymax=460, axis lines=none,
  legend style={at={(1.0,0.0)}, anchor=south west, font=\scriptsize}]
  \foreach \r in {100,200,300,400}{\edef\tmp{\noexpand\draw[fmGray!40, thin] (axis cs:0,0) circle[radius=\r];}\tmp}
  \draw[fmGray!60] (axis cs:-440,0) -- (axis cs:440,0) (axis cs:0,-440) -- (axis cs:0,440);
  \draw[fmGray!60, dashed] (axis cs:-380,-268.7) -- (axis cs:380,268.7);
  \draw[fmGray!60, dashed] (axis cs:380,-268.7) -- (axis cs:-380,268.7);
  \addplot[fmRed, thick] table[x=xf, y=yf] {figures/cr-young-polar.dat};
  \addlegendentry{Fe}
  \addplot[fmBlue, thick, dashed] table[x=xw, y=yw] {figures/cr-young-polar.dat};
  \addlegendentry{W}
  \node[font=\scriptsize, anchor=south] at (axis cs:0,440) {$[001]$};
  \node[font=\scriptsize, anchor=west] at (axis cs:440,0) {$[110]$};
  \node[font=\scriptsize, anchor=south west] at (axis cs:380,268.7) {$[111]$};
  \foreach \r in {100,200,300,400}{\edef\tmp{\noexpand\node[font=\noexpand\tiny, fmGray, anchor=north west] at (axis cs:\r,0) {\r};}\tmp}
\end{axis}
\end{tikzpicture}
\caption{Полярная диаграмма модуля Юнга (ГПа) монокристаллов железа и
вольфрама в плоскости $(1\br10)$: длина радиус-вектора равна модулю в
данном направлении (расчёт по упругим постоянным, п.~1.5.4)}
\label{fig:young-polar}
\end{figure}

\paragraph{Симметрия}
Анизотропия не означает, что свойства различны во \emph{всех}
направлениях. Наоборот, для некоторых направлений они закономерно
одинаковы (на рис.~\ref{fig:young-polar} модуль одинаков для всех
направлений типа $\nfam{111}$). Это и есть симметрия:
\begin{equation}\label{eq:symm}
F(\vb x)=F(\vb x')=\ldots=F(\vb x^{(n)}) ,
\end{equation}
т.\,е. функция $F$ инвариантна относительно определённых преобразований
аргумента. На языке симметрии однородность~— инвариантность
относительно любого параллельного переноса, а анизотропия проявляется
\emph{в рамках} симметрии: сами функции, описывающие свойства,
симметричны.

\begin{glavnoe}
По макроскопическим признакам кристаллическое вещество~— это
\vazhno{однородная анизотропная симметричная среда}.
\end{glavnoe}

\subsection{Решётка и структура кристаллов}\label{ssec:lattice-structure}

Конкретное расположение атомов в кристалле называют его
\term{структурой}. Любой кристалл можно разбить на одинаковые ячейки,
содержащие один или несколько определённым образом расположенных атомов;
бесконечное периодическое повторение таких ячеек~— \term{кристаллическое
пространство}.

\begin{opredelenie}[Преобразование симметрии и гомологичные точки]
\term{Преобразование симметрии}~— операция, переводящая точки
кристаллического пространства в новые положения без изменения расстояний
между ними, после которой пространство нельзя отличить от исходного ни
по свойствам, ни по геометрии. Две точки $O$ и $O'$ \term{гомологичны}
(эквивалентны), если существует преобразование симметрии, переводящее $O$
в $O'$.
\end{opredelenie}

Математически кристаллическое пространство строится на трёх
постулатах:
\begin{enumerate}
\item не все точки пространства одинаковы;
\item пространство \emph{однородно}: существует шар постоянного радиуса $R$
(шар однородности), внутри которого, где бы его ни поместить, найдётся
точка, гомологичная любой заданной; рентгеновские данные показывают, что
$R$ не превышает нескольких нанометров;
\item пространство \emph{дискретно}: есть точки, вокруг которых в шаре
некоторого радиуса $r$ (шар дискретности) нет ни одной гомологичной им
точки.
\end{enumerate}
Этих постулатов достаточно, чтобы построить всю геометрическую
кристаллографию.

\paragraph{Узловой ряд и решётка}
Возьмём любую точку $O$ и шар однородности вокруг неё. В шаре найдётся
хотя бы одна точка, гомологичная $O$; ближайшую из них обозначим $1$, а
расстояние до неё~— $a_1$. Раз точка $1$ ничем не отличается от $O$, а у
$O$ справа на расстоянии $a_1$ есть сосед $1$, то и у точки $1$ на том же
расстоянии справа есть гомологичная точка $2$, и так до бесконечности; то
же влево. Получается бесконечный ряд гомологичных точек~— \term{узловой
ряд} с периодом $a_1$, называемым \term{длиной трансляции}. Значит,
\emph{трансляция}~— перенос на $a_1$~— есть преобразование симметрии
бесконечного кристаллического пространства. Повторив рассуждение в двух
других направлениях, получим, что гомологичные точки образуют
\term{пространственную решётку}.

\begin{figure}[ht]
\centering
\begin{tikzpicture}[x=1cm, y=1cm]
  \foreach \i in {0,1,2,3}{
    \node[atom=4.2mm] at (2.6*\i,0) {\scriptsize A};
    \node[atom=3.2mm, fill=atomB] at (2.6*\i+0.85,0) {\scriptsize B};
    \node[atom=2.6mm, fill=atomI] at (2.6*\i+1.7,0) {\tiny C};}
  \node[font=\small, anchor=east] at (-0.5,0) {\textit{а}};
  \foreach \i in {0,1,2,3}{\fill[fmBlue] (2.6*\i,-1.1) circle (2pt);}
  \draw[fmBlue] (-0.3,-1.1) -- (8.4,-1.1);
  \draw[dimline] (0,-1.45) -- node[below, font=\scriptsize]{$a_1$} (2.6,-1.45);
  \node[font=\small, anchor=east] at (-0.5,-1.1) {\textit{б}};
  \foreach \i in {0,1,2,3}{\draw[fmRed] (2.6*\i+0.425,-2.3) node[cross out, draw, inner sep=1.6pt]{};}
  \draw[fmRed] (-0.3,-2.3) -- (8.4,-2.3);
  \node[font=\small, anchor=east] at (-0.5,-2.3) {\textit{в}};
\end{tikzpicture}
\caption{Структура и решётка: \textit{а}~— одномерная структура из атомов
трёх сортов; \textit{б}~— узловой ряд с узлами в центрах атомов A;
\textit{в}~— тот же ряд, сдвинутый так, что узлы оказались посередине
между атомами A и B. Положение узлов произвольно, период один и тот же}
\label{fig:row}
\end{figure}

Разницу между структурой и решёткой хорошо видно на одномерном примере
(рис.~\ref{fig:row}). Поместив узел в центр атома A, найдём остальные узлы
в центрах всех атомов A. Поместив его посередине между A и B, найдём
остальные узлы в таких же серединах. Узлы~— не атомы, а гомологичные
точки, и положение начального узла выбирается произвольно.

\begin{glavnoe}[Структура = решётка + базис]
Структура кристалла~— физическая реальность~— состоит из атомов.
Пространственная решётка~— математическая абстракция, способ описать
\emph{периодичность}: она состоит из узлов. Чтобы из решётки получить
структуру, в каждый узел надо «повесить» одинаковую группу атомов~—
\term{базис} (на рис.~\ref{fig:row} базис~— тройка атомов A, B, C).
\end{glavnoe}

\begin{itogi}
\begin{itemize}
\item Кристалл~— твёрдое тело с трёхмерно-периодической структурой;
кристаллическое состояние равновесно.
\item Законы постоянства углов (Стенон) и рациональных параметров (Гаюи)
привели к идее периодичности; её прямое доказательство~— дифракция
рентгеновских лучей (Лауэ, 1912); всего 230 пространственных групп
(Фёдоров, Шёнфлис).
\item Макроскопически кристалл~— однородная анизотропная симметричная
среда.
\item Гомологичные точки образуют пространственную решётку; структура =
решётка + базис.
\end{itemize}
\end{itogi}

\begin{proverka}
\item Почему кристаллы одного вещества могут иметь разную форму, но
одинаковые углы между гранями?
\item Какой вывод о строении кристаллов следует из закона рациональных
параметров?
\item Чем однородность отличается от изотропии? Приведите пример
однородной, но анизотропной среды.
\item Почему узлы решётки не обязательно совпадают с атомами?
\item Сколько атомов в базисе структуры на рис.~\ref{fig:row}?
\end{proverka}


\section{Основы кристаллографии}\label{sec:cryst-basics}
% ----------------------------------------------------------------------
%  1.2.1. Кристаллографические проекции
% ----------------------------------------------------------------------
\begin{orientir}
Раздел даёт рабочий аппарат кристаллографа: как изобразить на плоскости
ориентацию граней и направлений (проекции, 1.2.1), как описать решётку
(1.2.2), как обозначать узлы, направления и плоскости (индексы Миллера,
1.2.3), что такое обратная решётка (1.2.4), как переходить между
системами координат (1.2.5--1.2.6) и как считать углы и расстояния
(1.2.7). Завершает раздел ячейка Вигнера--Зейтца (1.2.8).
\end{orientir}

\subsection{Кристаллографические проекции}\label{ssec:projections}

Для многих задач важна не форма кристалла, а только \emph{взаимная
ориентация} его граней и рёбер. Поэтому вместо кристалла рассматривают
\term{кристаллографический комплекс}: все параллельные грани заменяют
одной плоскостью, все параллельные рёбра~— одной прямой, и все они
проходят через одну точку~— \emph{центр комплекса}. Так, кубооктаэдр (куб
со срезанными вершинами) превращается в три взаимно перпендикулярные
плоскости граней куба и четыре наклонные плоскости срезов.

\begin{opredelenie}[Зона]
Совокупность плоскостей, пересекающихся по общей прямой, называют
\term{кристаллографической зоной}, а эту прямую~— \term{осью зоны}.
Комплекс состоит из ряда зон.
\end{opredelenie}

Часто удобнее \term{обратный} (полярный) комплекс: каждую плоскость
заменяют нормалью к ней, а каждую прямую~— перпендикулярной ей плоскостью.
Зона в обратном комплексе~— это пучок нормалей, лежащих в одной
плоскости, перпендикулярной оси зоны (рис.~\ref{fig:zone}). Для
кубического кристалла прямой и обратный комплексы совпадают.

\begin{figure}[ht]
\centering
\begin{tikzpicture}[x={(1cm,0cm)}, y={(0.45cm,0.32cm)}, z={(0cm,1cm)}, scale=1.05]
  \draw[fmRed, very thick] (0,0,-1.9) -- (0,0,2.1) node[above, font=\small]{ось зоны};
  \foreach \a/\c in {0/fmBlue, 60/fmGreen, 120/fmAmber}{
    \pgfmathsetmacro{\cx}{1.6*cos(\a)}\pgfmathsetmacro{\cy}{1.6*sin(\a)}
    \fill[\c, opacity=0.18] (\cx,\cy,-1.5) -- (\cx,\cy,1.5) -- (-\cx,-\cy,1.5) -- (-\cx,-\cy,-1.5) -- cycle;
    \draw[\c, thin] (\cx,\cy,-1.5) -- (\cx,\cy,1.5) -- (-\cx,-\cy,1.5) -- (-\cx,-\cy,-1.5) -- cycle;
    \pgfmathsetmacro{\nx}{1.5*cos(\a+90)}\pgfmathsetmacro{\ny}{1.5*sin(\a+90)}
    \draw[\c, thick, -{Stealth[length=2mm]}] (0,0,0) -- (\nx,\ny,0);}
  \draw[fmGray, dashed] (0,0,0) circle[radius=1.5];
  \fill (0,0,0) circle (1.3pt);
  \node[font=\scriptsize, fmGray, anchor=west] at (1.5,-1.6,0) {плоскость нормалей};
\end{tikzpicture}
\caption{Кристаллографическая зона: три плоскости, пересекающиеся по оси
зоны (красная линия), и их нормали (стрелки). Нормали всех плоскостей
зоны лежат в одной плоскости, перпендикулярной оси зоны,~— так зона
выглядит в обратном комплексе}
\label{fig:zone}
\end{figure}

\paragraph{Виды проекций}
Комплекс нужно изобразить на плоскости. Способы различаются тем, какой
комплекс проектируют (прямой или обратный) и на что
(табл.~\ref{tab:projections}).

\begin{table}[ht]
\centering
\caption{Кристаллографические проекции}\label{tab:projections}
\small
\begin{tabularx}{\textwidth}{@{}>{\raggedright\arraybackslash\hspace{0pt}}p{33mm}>{\raggedright\arraybackslash\hspace{0pt}}p{25mm}>{\raggedright\arraybackslash\hspace{0pt}}X>{\raggedright\arraybackslash\hspace{0pt}}X>{\raggedright\arraybackslash\hspace{0pt}}X@{}}
\toprule
Проекция & Что и куда & Плоскость & Направление & Зона \\
\midrule
линейная & прямой $\to$ плоскость & прямая & точка & прямые через точку \\
гномоническая & обратный $\to$ плоскость & точка & прямая & точки на прямой \\
сферическая & прямой $\to$ сфера & большой круг & две точки & большие круги через две точки \\
гномосферическая & обратный $\to$ сфера & две точки & большой круг & точки на большом круге \\
\textbf{стереографическая} & прямой $\to$ плоскость & дуга большого круга & точка & дуги через две точки \\
\textbf{гномостереографическая} & обратный $\to$ плоскость & точка (полюс) & дуга большого круга & точки на одной дуге \\
\bottomrule
\end{tabularx}
\end{table}

Линейной и гномонической проекциям для изображения всего комплекса нужна
бесконечная плоскость, а сферические проекции неплоские. На практике
используют почти исключительно \emph{стереографическую} и
\emph{гномостереографическую} проекции. Гномоническую иногда применяют
при обработке лауэграмм.

\paragraph{Стереографическая проекция}
Центр комплекса помещают в центр $O$ сферы; плоскость проекции проходит
через $O$ (рис.~\ref{fig:stereo}\,\textit{а}). Её сечение сферой~—
\term{основной круг} проекции, концы перпендикулярного ей диаметра~—
\emph{полюсы} $N$ (над плоскостью) и $S$ (под ней). Прямую $L$ комплекса
продолжают до пересечения со сферой в точке $M_1$ и соединяют $M_1$ с
южным полюсом $S$; точка $M_1'$, где этот луч пересекает плоскость
проекции,~— \emph{стереографическая проекция прямой}. Второй выход той же
прямой, $M_2$ в южном полушарии, проектируют из $N$ и отмечают
\emph{крестиком} $M_2'$.

\begin{figure}[ht]
\centering
\begin{subfigure}[b]{0.52\textwidth}\centering
\begin{tikzpicture}[scale=1.0]
  \def\R{2.2}
  \draw[atomLine, thick] (0,0) circle (\R);
  \draw[fmGray, thick] (-2.9,0) -- (2.9,0) node[right, font=\small]{$P$};
  \draw[fmGray, dashed] (0,-\R) -- (0,\R);
  \fill (0,\R) circle (1.4pt) node[above, font=\small]{$N$};
  \fill (0,-\R) circle (1.4pt) node[below, font=\small]{$S$};
  \fill (0,0) circle (1.2pt) node[below left, font=\small]{$O$};
  \coordinate (M1) at (50:\R); \coordinate (M2) at (230:\R);
  \draw[fmBlue, very thick] (M2) -- (M1) node[pos=0.88, right, font=\small]{$L$};
  \fill[fmBlue] (M1) circle (1.6pt) node[above right, font=\small]{$M_1$};
  \fill[fmBlue] (M2) circle (1.6pt) node[below left, font=\small]{$M_2$};
  \draw[fmRed] (0,-\R) -- (M1);
  \draw[fmRed, dashed] (0,\R) -- (M2);
  \fill[fmRed] (0.8007,0) circle (1.8pt) node[below right=-1pt, font=\small]{$M_1'$};
  \node[fmRed, font=\small] at (-0.8007,0) {$\times$};
  \node[fmRed, font=\small, above left=-1pt] at (-0.8007,0) {$M_2'$};
\end{tikzpicture}
\caption{}
\end{subfigure}\hfill
\begin{subfigure}[b]{0.44\textwidth}\centering
\begin{tikzpicture}[scale=1.0]
  \def\R{2.2}
  \draw[atomLine, thick] (0,0) circle (\R);
  \fill (0,0) circle (1.2pt);
  \fill[fmRed] (0.8007,0) circle (1.8pt) node[above right=-1pt, font=\small]{$M_1'$};
  \node[fmRed, font=\small] at (-0.8007,0) {$\times$};
  \node[fmRed, font=\small, above left=-1pt] at (-0.8007,0) {$M_2'$};
  \begin{scope}
    \clip (0,0) circle (\R);
    \draw[fmBlue, thick] (-2.622,0) circle (3.4227);
    \draw[fmBlue, thick, dashed] (2.622,0) circle (3.4227);
  \end{scope}
  \node[fmBlue, font=\small] at (1.05,1.25) {$p$};
  \node[fmBlue, font=\small] at (-1.0,1.3) {$p'$};
\end{tikzpicture}
\caption{}
\end{subfigure}
\caption{Стереографическая проекция: \textit{а}~— построение (сечение
сферы плоскостью, содержащей ось $NS$): выход прямой $L$ в северном
полушарии $M_1$ проектируется из $S$ в точку $M_1'$, выход в южном
полушарии $M_2$~— из $N$ в крестик $M_2'$; \textit{б}~— вид на плоскость
проекции: прямая изображается точкой и крестиком на равных расстояниях от
центра, наклонная плоскость~— дугой большого круга $p$ (её часть в южном
полушарии~— штриховой дугой $p'$)}
\label{fig:stereo}
\end{figure}

\begin{keyeq}\label{eq:stereo-rho}
\rho=r\,\operatorname{tg}\frac\theta2
\end{keyeq}
где $\theta$~— угол между прямой и осью $ON$, $r$~— радиус основного круга,
$\rho$~— расстояние проекции от центра. Прямая, перпендикулярная
плоскости проекции ($\theta=0$), проектируется в центр; прямая, лежащая в
плоскости проекции ($\theta=90^\circ$),~— в две точки основного круга.

Плоскость комплекса проектируется дугой большого круга, опирающейся на
диаметр основного круга (рис.~\ref{fig:stereo}\,\textit{б}): сплошной~— её
часть над плоскостью проекции, штриховой~— под ней. Вертикальная
плоскость изображается диаметром, плоскость проекции~— основным кругом.
Зона~— серия дуг, проходящих через две диаметрально противоположные
точки~— проекции оси зоны.

\begin{glavnoe}[Два свойства, ради которых всё это делается]
\begin{enumerate}
\item Любой круг на сфере (большой или малый) изображается на проекции
окружностью или её дугой.
\item Проекция \emph{конформна}: углы между кривыми на сфере равны углам
между их изображениями. Стереографическое проектирование из $S$
устанавливает взаимно однозначное соответствие между точками сферы (кроме
самой $S$) и плоскости.
\end{enumerate}
\end{glavnoe}

\paragraph{Гномостереографическая проекция}
Это стереографическая проекция обратного комплекса: плоскость заменяют
нормалью $\vb n$ и проектируют точку выхода нормали на сферу. Плоскость
изображается одной точкой~— \term{полюсом}, направление~— дугой большого
круга, а все плоскости одной зоны~— полюсами, лежащими на одной дуге. Эта
проекция удобнее для практических задач, и дальше полюсы плоскостей и
выходы направлений мы будем наносить на один чертёж.

\paragraph{Сетки Вульфа и Болдырева}
Измерять углы на проекции позволяют координатные сетки
(рис.~\ref{fig:nets}). \term{Сетка Вульфа}~— стереографическая проекция
двух семейств кругов на сфере:
\begin{itemize}
\item \emph{меридианы}~— большие круги, проходящие через ось $y$, лежащую в
плоскости проекции, повёрнутые друг относительно друга через $2^\circ$
(на рисунке~— через $10^\circ$). Каждый меридиан~— проекция плоскости, а
все вместе~— проекция зоны с осью $y$;
\item \emph{параллели}~— малые круги, проекции соосных конусов с осью $y$,
полураствор которых меняется через $2^\circ$.
\end{itemize}
Угловая шкала по основному кругу равномерна, по экватору и нулевому
меридиану~— неравномерна, но одинакова. \term{Сетка Болдырева}~—
концентрические окружности (проекции соосных конусов с осью $z$,
перпендикулярной плоскости проекции; радиусы $r\operatorname{tg}(\theta/2)$) и
радиальные прямые (проекции плоскостей зоны с осью $z$).

\begin{figure}[ht]
\centering
\tikzset{wulffA/.style={fmBlue!45, line width=0.25pt},
  wulffB/.style={fmBlue!80, line width=0.55pt},
  wulffC/.style={atomLine, line width=0.9pt}}
\begin{subfigure}[b]{0.48\textwidth}\centering
\begin{tikzpicture}[scale=2.9]
  % создано tools/gen_stereo.py: сетка Вульфа (r = 1)
\draw[wulffA] (0.0000,-1.0000) -- (0.0342,-0.9934) -- (0.0679,-0.9856) -- (0.1011,-0.9768) -- (0.1338,-0.9669) -- (0.1660,-0.9560) -- (0.1976,-0.9441) -- (0.2286,-0.9312) -- (0.2591,-0.9174) -- (0.2888,-0.9026) -- (0.3179,-0.8870) -- (0.3464,-0.8706) -- (0.3741,-0.8533) -- (0.4012,-0.8352) -- (0.4275,-0.8164) -- (0.4531,-0.7968) -- (0.4779,-0.7766) -- (0.5020,-0.7557) -- (0.5252,-0.7341) -- (0.5477,-0.7119) -- (0.5695,-0.6891) -- (0.5904,-0.6658) -- (0.6105,-0.6419) -- (0.6297,-0.6175) -- (0.6482,-0.5927) -- (0.6658,-0.5673) -- (0.6826,-0.5416) -- (0.6986,-0.5154) -- (0.7137,-0.4888) -- (0.7280,-0.4619) -- (0.7414,-0.4346) -- (0.7539,-0.4071) -- (0.7656,-0.3792) -- (0.7765,-0.3510) -- (0.7865,-0.3227) -- (0.7956,-0.2940) -- (0.8039,-0.2652) -- (0.8112,-0.2362) -- (0.8178,-0.2070) -- (0.8234,-0.1777) -- (0.8282,-0.1483) -- (0.8321,-0.1188) -- (0.8352,-0.0891) -- (0.8374,-0.0595) -- (0.8387,-0.0297) -- (0.8391,0.0000) -- (0.8387,0.0297) -- (0.8374,0.0595) -- (0.8352,0.0891) -- (0.8321,0.1188) -- (0.8282,0.1483) -- (0.8234,0.1777) -- (0.8178,0.2070) -- (0.8112,0.2362) -- (0.8039,0.2652) -- (0.7956,0.2940) -- (0.7865,0.3227) -- (0.7765,0.3510) -- (0.7656,0.3792) -- (0.7539,0.4071) -- (0.7414,0.4346) -- (0.7280,0.4619) -- (0.7137,0.4888) -- (0.6986,0.5154) -- (0.6826,0.5416) -- (0.6658,0.5673) -- (0.6482,0.5927) -- (0.6297,0.6175) -- (0.6105,0.6419) -- (0.5904,0.6658) -- (0.5695,0.6891) -- (0.5477,0.7119) -- (0.5252,0.7341) -- (0.5020,0.7557) -- (0.4779,0.7766) -- (0.4531,0.7968) -- (0.4275,0.8164) -- (0.4012,0.8352) -- (0.3741,0.8533) -- (0.3464,0.8706) -- (0.3179,0.8870) -- (0.2888,0.9026) -- (0.2591,0.9174) -- (0.2286,0.9312) -- (0.1976,0.9441) -- (0.1660,0.9560) -- (0.1338,0.9669) -- (0.1011,0.9768) -- (0.0679,0.9856) -- (0.0342,0.9934) -- (0.0000,1.0000);
\draw[wulffA] (0.0000,-1.0000) -- (0.0324,-0.9876) -- (0.0640,-0.9743) -- (0.0948,-0.9602) -- (0.1248,-0.9453) -- (0.1540,-0.9296) -- (0.1824,-0.9132) -- (0.2100,-0.8961) -- (0.2367,-0.8784) -- (0.2626,-0.8601) -- (0.2877,-0.8413) -- (0.3120,-0.8219) -- (0.3355,-0.8020) -- (0.3582,-0.7816) -- (0.3801,-0.7608) -- (0.4012,-0.7396) -- (0.4216,-0.7179) -- (0.4411,-0.6959) -- (0.4599,-0.6736) -- (0.4779,-0.6509) -- (0.4952,-0.6280) -- (0.5117,-0.6047) -- (0.5274,-0.5812) -- (0.5425,-0.5575) -- (0.5568,-0.5335) -- (0.5704,-0.5093) -- (0.5833,-0.4850) -- (0.5955,-0.4604) -- (0.6069,-0.4357) -- (0.6177,-0.4108) -- (0.6278,-0.3857) -- (0.6373,-0.3606) -- (0.6460,-0.3353) -- (0.6541,-0.3099) -- (0.6615,-0.2844) -- (0.6683,-0.2588) -- (0.6743,-0.2332) -- (0.6798,-0.2074) -- (0.6846,-0.1816) -- (0.6887,-0.1558) -- (0.6923,-0.1299) -- (0.6951,-0.1040) -- (0.6973,-0.0780) -- (0.6989,-0.0520) -- (0.6999,-0.0260) -- (0.7002,0.0000) -- (0.6999,0.0260) -- (0.6989,0.0520) -- (0.6973,0.0780) -- (0.6951,0.1040) -- (0.6923,0.1299) -- (0.6887,0.1558) -- (0.6846,0.1816) -- (0.6798,0.2074) -- (0.6743,0.2332) -- (0.6683,0.2588) -- (0.6615,0.2844) -- (0.6541,0.3099) -- (0.6460,0.3353) -- (0.6373,0.3606) -- (0.6278,0.3857) -- (0.6177,0.4108) -- (0.6069,0.4357) -- (0.5955,0.4604) -- (0.5833,0.4850) -- (0.5704,0.5093) -- (0.5568,0.5335) -- (0.5425,0.5575) -- (0.5274,0.5812) -- (0.5117,0.6047) -- (0.4952,0.6280) -- (0.4779,0.6509) -- (0.4599,0.6736) -- (0.4411,0.6959) -- (0.4216,0.7179) -- (0.4012,0.7396) -- (0.3801,0.7608) -- (0.3582,0.7816) -- (0.3355,0.8020) -- (0.3120,0.8219) -- (0.2877,0.8413) -- (0.2626,0.8601) -- (0.2367,0.8784) -- (0.2100,0.8961) -- (0.1824,0.9132) -- (0.1540,0.9296) -- (0.1248,0.9453) -- (0.0948,0.9602) -- (0.0640,0.9743) -- (0.0324,0.9876) -- (0.0000,1.0000);
\draw[wulffB] (0.0000,-1.0000) -- (0.0297,-0.9823) -- (0.0584,-0.9639) -- (0.0860,-0.9451) -- (0.1127,-0.9258) -- (0.1384,-0.9061) -- (0.1631,-0.8860) -- (0.1869,-0.8656) -- (0.2098,-0.8448) -- (0.2318,-0.8238) -- (0.2529,-0.8025) -- (0.2732,-0.7809) -- (0.2927,-0.7592) -- (0.3114,-0.7372) -- (0.3293,-0.7151) -- (0.3464,-0.6928) -- (0.3628,-0.6704) -- (0.3785,-0.6479) -- (0.3934,-0.6253) -- (0.4077,-0.6025) -- (0.4213,-0.5797) -- (0.4342,-0.5568) -- (0.4465,-0.5339) -- (0.4582,-0.5109) -- (0.4692,-0.4879) -- (0.4797,-0.4648) -- (0.4896,-0.4416) -- (0.4988,-0.4185) -- (0.5076,-0.3953) -- (0.5157,-0.3721) -- (0.5234,-0.3489) -- (0.5305,-0.3257) -- (0.5370,-0.3025) -- (0.5431,-0.2792) -- (0.5486,-0.2560) -- (0.5537,-0.2327) -- (0.5582,-0.2094) -- (0.5622,-0.1862) -- (0.5658,-0.1629) -- (0.5689,-0.1396) -- (0.5715,-0.1164) -- (0.5736,-0.0931) -- (0.5752,-0.0698) -- (0.5764,-0.0465) -- (0.5771,-0.0233) -- (0.5774,0.0000) -- (0.5771,0.0233) -- (0.5764,0.0465) -- (0.5752,0.0698) -- (0.5736,0.0931) -- (0.5715,0.1164) -- (0.5689,0.1396) -- (0.5658,0.1629) -- (0.5622,0.1862) -- (0.5582,0.2094) -- (0.5537,0.2327) -- (0.5486,0.2560) -- (0.5431,0.2792) -- (0.5370,0.3025) -- (0.5305,0.3257) -- (0.5234,0.3489) -- (0.5157,0.3721) -- (0.5076,0.3953) -- (0.4988,0.4185) -- (0.4896,0.4416) -- (0.4797,0.4648) -- (0.4692,0.4879) -- (0.4582,0.5109) -- (0.4465,0.5339) -- (0.4342,0.5568) -- (0.4213,0.5797) -- (0.4077,0.6025) -- (0.3934,0.6253) -- (0.3785,0.6479) -- (0.3628,0.6704) -- (0.3464,0.6928) -- (0.3293,0.7151) -- (0.3114,0.7372) -- (0.2927,0.7592) -- (0.2732,0.7809) -- (0.2529,0.8025) -- (0.2318,0.8238) -- (0.2098,0.8448) -- (0.1869,0.8656) -- (0.1631,0.8860) -- (0.1384,0.9061) -- (0.1127,0.9258) -- (0.0860,0.9451) -- (0.0584,0.9639) -- (0.0297,0.9823) -- (0.0000,1.0000);
\draw[wulffA] (0.0000,-1.0000) -- (0.0261,-0.9775) -- (0.0511,-0.9548) -- (0.0750,-0.9319) -- (0.0979,-0.9090) -- (0.1197,-0.8859) -- (0.1405,-0.8628) -- (0.1604,-0.8397) -- (0.1794,-0.8166) -- (0.1975,-0.7935) -- (0.2148,-0.7703) -- (0.2313,-0.7473) -- (0.2470,-0.7242) -- (0.2620,-0.7012) -- (0.2763,-0.6783) -- (0.2899,-0.6554) -- (0.3028,-0.6326) -- (0.3151,-0.6098) -- (0.3268,-0.5872) -- (0.3379,-0.5646) -- (0.3484,-0.5421) -- (0.3584,-0.5196) -- (0.3679,-0.4973) -- (0.3768,-0.4750) -- (0.3853,-0.4528) -- (0.3932,-0.4307) -- (0.4007,-0.4087) -- (0.4077,-0.3867) -- (0.4143,-0.3648) -- (0.4204,-0.3430) -- (0.4262,-0.3212) -- (0.4315,-0.2995) -- (0.4364,-0.2778) -- (0.4409,-0.2563) -- (0.4450,-0.2347) -- (0.4488,-0.2132) -- (0.4521,-0.1918) -- (0.4551,-0.1704) -- (0.4578,-0.1490) -- (0.4601,-0.1277) -- (0.4620,-0.1063) -- (0.4635,-0.0850) -- (0.4647,-0.0638) -- (0.4656,-0.0425) -- (0.4661,-0.0212) -- (0.4663,0.0000) -- (0.4661,0.0212) -- (0.4656,0.0425) -- (0.4647,0.0638) -- (0.4635,0.0850) -- (0.4620,0.1063) -- (0.4601,0.1277) -- (0.4578,0.1490) -- (0.4551,0.1704) -- (0.4521,0.1918) -- (0.4488,0.2132) -- (0.4450,0.2347) -- (0.4409,0.2563) -- (0.4364,0.2778) -- (0.4315,0.2995) -- (0.4262,0.3212) -- (0.4204,0.3430) -- (0.4143,0.3648) -- (0.4077,0.3867) -- (0.4007,0.4087) -- (0.3932,0.4307) -- (0.3853,0.4528) -- (0.3768,0.4750) -- (0.3679,0.4973) -- (0.3584,0.5196) -- (0.3484,0.5421) -- (0.3379,0.5646) -- (0.3268,0.5872) -- (0.3151,0.6098) -- (0.3028,0.6326) -- (0.2899,0.6554) -- (0.2763,0.6783) -- (0.2620,0.7012) -- (0.2470,0.7242) -- (0.2313,0.7473) -- (0.2148,0.7703) -- (0.1975,0.7935) -- (0.1794,0.8166) -- (0.1604,0.8397) -- (0.1405,0.8628) -- (0.1197,0.8859) -- (0.0979,0.9090) -- (0.0750,0.9319) -- (0.0511,0.9548) -- (0.0261,0.9775) -- (0.0000,1.0000);
\draw[wulffA] (0.0000,-1.0000) -- (0.0218,-0.9734) -- (0.0426,-0.9470) -- (0.0622,-0.9208) -- (0.0808,-0.8949) -- (0.0985,-0.8692) -- (0.1153,-0.8438) -- (0.1312,-0.8186) -- (0.1463,-0.7937) -- (0.1606,-0.7690) -- (0.1742,-0.7446) -- (0.1871,-0.7204) -- (0.1993,-0.6965) -- (0.2109,-0.6728) -- (0.2219,-0.6494) -- (0.2324,-0.6262) -- (0.2423,-0.6032) -- (0.2516,-0.5804) -- (0.2605,-0.5578) -- (0.2689,-0.5355) -- (0.2769,-0.5133) -- (0.2844,-0.4913) -- (0.2914,-0.4695) -- (0.2981,-0.4479) -- (0.3044,-0.4264) -- (0.3103,-0.4051) -- (0.3159,-0.3839) -- (0.3211,-0.3629) -- (0.3259,-0.3420) -- (0.3304,-0.3212) -- (0.3347,-0.3006) -- (0.3386,-0.2801) -- (0.3422,-0.2596) -- (0.3455,-0.2393) -- (0.3485,-0.2190) -- (0.3512,-0.1989) -- (0.3537,-0.1788) -- (0.3559,-0.1587) -- (0.3578,-0.1388) -- (0.3594,-0.1189) -- (0.3608,-0.0990) -- (0.3620,-0.0791) -- (0.3628,-0.0593) -- (0.3635,-0.0395) -- (0.3638,-0.0198) -- (0.3640,0.0000) -- (0.3638,0.0198) -- (0.3635,0.0395) -- (0.3628,0.0593) -- (0.3620,0.0791) -- (0.3608,0.0990) -- (0.3594,0.1189) -- (0.3578,0.1388) -- (0.3559,0.1587) -- (0.3537,0.1788) -- (0.3512,0.1989) -- (0.3485,0.2190) -- (0.3455,0.2393) -- (0.3422,0.2596) -- (0.3386,0.2801) -- (0.3347,0.3006) -- (0.3304,0.3212) -- (0.3259,0.3420) -- (0.3211,0.3629) -- (0.3159,0.3839) -- (0.3103,0.4051) -- (0.3044,0.4264) -- (0.2981,0.4479) -- (0.2914,0.4695) -- (0.2844,0.4913) -- (0.2769,0.5133) -- (0.2689,0.5355) -- (0.2605,0.5578) -- (0.2516,0.5804) -- (0.2423,0.6032) -- (0.2324,0.6262) -- (0.2219,0.6494) -- (0.2109,0.6728) -- (0.1993,0.6965) -- (0.1871,0.7204) -- (0.1742,0.7446) -- (0.1606,0.7690) -- (0.1463,0.7937) -- (0.1312,0.8186) -- (0.1153,0.8438) -- (0.0985,0.8692) -- (0.0808,0.8949) -- (0.0622,0.9208) -- (0.0426,0.9470) -- (0.0218,0.9734) -- (0.0000,1.0000);
\draw[wulffB] (0.0000,-1.0000) -- (0.0169,-0.9701) -- (0.0329,-0.9407) -- (0.0479,-0.9120) -- (0.0621,-0.8838) -- (0.0755,-0.8561) -- (0.0881,-0.8289) -- (0.1000,-0.8022) -- (0.1113,-0.7760) -- (0.1219,-0.7503) -- (0.1319,-0.7250) -- (0.1414,-0.7001) -- (0.1504,-0.6756) -- (0.1589,-0.6515) -- (0.1669,-0.6277) -- (0.1745,-0.6043) -- (0.1816,-0.5813) -- (0.1884,-0.5585) -- (0.1948,-0.5361) -- (0.2008,-0.5140) -- (0.2065,-0.4921) -- (0.2118,-0.4705) -- (0.2169,-0.4491) -- (0.2216,-0.4280) -- (0.2261,-0.4071) -- (0.2303,-0.3864) -- (0.2342,-0.3659) -- (0.2379,-0.3456) -- (0.2413,-0.3255) -- (0.2445,-0.3055) -- (0.2474,-0.2857) -- (0.2502,-0.2660) -- (0.2527,-0.2465) -- (0.2550,-0.2271) -- (0.2571,-0.2078) -- (0.2590,-0.1886) -- (0.2608,-0.1695) -- (0.2623,-0.1504) -- (0.2636,-0.1315) -- (0.2648,-0.1126) -- (0.2658,-0.0937) -- (0.2665,-0.0749) -- (0.2672,-0.0562) -- (0.2676,-0.0374) -- (0.2679,-0.0187) -- (0.2679,0.0000) -- (0.2679,0.0187) -- (0.2676,0.0374) -- (0.2672,0.0562) -- (0.2665,0.0749) -- (0.2658,0.0937) -- (0.2648,0.1126) -- (0.2636,0.1315) -- (0.2623,0.1504) -- (0.2608,0.1695) -- (0.2590,0.1886) -- (0.2571,0.2078) -- (0.2550,0.2271) -- (0.2527,0.2465) -- (0.2502,0.2660) -- (0.2474,0.2857) -- (0.2445,0.3055) -- (0.2413,0.3255) -- (0.2379,0.3456) -- (0.2342,0.3659) -- (0.2303,0.3864) -- (0.2261,0.4071) -- (0.2216,0.4280) -- (0.2169,0.4491) -- (0.2118,0.4705) -- (0.2065,0.4921) -- (0.2008,0.5140) -- (0.1948,0.5361) -- (0.1884,0.5585) -- (0.1816,0.5813) -- (0.1745,0.6043) -- (0.1669,0.6277) -- (0.1589,0.6515) -- (0.1504,0.6756) -- (0.1414,0.7001) -- (0.1319,0.7250) -- (0.1219,0.7503) -- (0.1113,0.7760) -- (0.1000,0.8022) -- (0.0881,0.8289) -- (0.0755,0.8561) -- (0.0621,0.8838) -- (0.0479,0.9120) -- (0.0329,0.9407) -- (0.0169,0.9701) -- (0.0000,1.0000);
\draw[wulffA] (0.0000,-1.0000) -- (0.0116,-0.9677) -- (0.0224,-0.9362) -- (0.0326,-0.9056) -- (0.0421,-0.8757) -- (0.0511,-0.8467) -- (0.0595,-0.8183) -- (0.0674,-0.7906) -- (0.0749,-0.7635) -- (0.0819,-0.7370) -- (0.0885,-0.7111) -- (0.0948,-0.6858) -- (0.1006,-0.6609) -- (0.1062,-0.6366) -- (0.1114,-0.6127) -- (0.1163,-0.5892) -- (0.1210,-0.5661) -- (0.1254,-0.5435) -- (0.1295,-0.5212) -- (0.1334,-0.4992) -- (0.1371,-0.4776) -- (0.1405,-0.4563) -- (0.1438,-0.4352) -- (0.1468,-0.4145) -- (0.1497,-0.3940) -- (0.1523,-0.3737) -- (0.1549,-0.3537) -- (0.1572,-0.3339) -- (0.1594,-0.3143) -- (0.1614,-0.2949) -- (0.1633,-0.2757) -- (0.1650,-0.2566) -- (0.1667,-0.2377) -- (0.1681,-0.2189) -- (0.1695,-0.2002) -- (0.1707,-0.1816) -- (0.1718,-0.1632) -- (0.1727,-0.1448) -- (0.1736,-0.1265) -- (0.1743,-0.1083) -- (0.1749,-0.0902) -- (0.1754,-0.0721) -- (0.1758,-0.0540) -- (0.1761,-0.0360) -- (0.1763,-0.0180) -- (0.1763,0.0000) -- (0.1763,0.0180) -- (0.1761,0.0360) -- (0.1758,0.0540) -- (0.1754,0.0721) -- (0.1749,0.0902) -- (0.1743,0.1083) -- (0.1736,0.1265) -- (0.1727,0.1448) -- (0.1718,0.1632) -- (0.1707,0.1816) -- (0.1695,0.2002) -- (0.1681,0.2189) -- (0.1667,0.2377) -- (0.1650,0.2566) -- (0.1633,0.2757) -- (0.1614,0.2949) -- (0.1594,0.3143) -- (0.1572,0.3339) -- (0.1549,0.3537) -- (0.1523,0.3737) -- (0.1497,0.3940) -- (0.1468,0.4145) -- (0.1438,0.4352) -- (0.1405,0.4563) -- (0.1371,0.4776) -- (0.1334,0.4992) -- (0.1295,0.5212) -- (0.1254,0.5435) -- (0.1210,0.5661) -- (0.1163,0.5892) -- (0.1114,0.6127) -- (0.1062,0.6366) -- (0.1006,0.6609) -- (0.0948,0.6858) -- (0.0885,0.7111) -- (0.0819,0.7370) -- (0.0749,0.7635) -- (0.0674,0.7906) -- (0.0595,0.8183) -- (0.0511,0.8467) -- (0.0421,0.8757) -- (0.0326,0.9056) -- (0.0224,0.9362) -- (0.0116,0.9677) -- (0.0000,1.0000);
\draw[wulffA] (0.0000,-1.0000) -- (0.0059,-0.9662) -- (0.0113,-0.9334) -- (0.0165,-0.9017) -- (0.0213,-0.8709) -- (0.0258,-0.8410) -- (0.0300,-0.8119) -- (0.0339,-0.7836) -- (0.0376,-0.7560) -- (0.0411,-0.7292) -- (0.0444,-0.7029) -- (0.0475,-0.6773) -- (0.0504,-0.6523) -- (0.0532,-0.6278) -- (0.0557,-0.6038) -- (0.0582,-0.5803) -- (0.0605,-0.5572) -- (0.0626,-0.5346) -- (0.0646,-0.5124) -- (0.0666,-0.4906) -- (0.0684,-0.4691) -- (0.0700,-0.4480) -- (0.0716,-0.4271) -- (0.0731,-0.4066) -- (0.0745,-0.3864) -- (0.0758,-0.3664) -- (0.0770,-0.3466) -- (0.0782,-0.3271) -- (0.0793,-0.3079) -- (0.0802,-0.2888) -- (0.0812,-0.2699) -- (0.0820,-0.2511) -- (0.0828,-0.2325) -- (0.0835,-0.2141) -- (0.0842,-0.1958) -- (0.0847,-0.1776) -- (0.0853,-0.1596) -- (0.0857,-0.1416) -- (0.0862,-0.1237) -- (0.0865,-0.1059) -- (0.0868,-0.0882) -- (0.0871,-0.0705) -- (0.0872,-0.0528) -- (0.0874,-0.0352) -- (0.0875,-0.0176) -- (0.0875,0.0000) -- (0.0875,0.0176) -- (0.0874,0.0352) -- (0.0872,0.0528) -- (0.0871,0.0705) -- (0.0868,0.0882) -- (0.0865,0.1059) -- (0.0862,0.1237) -- (0.0857,0.1416) -- (0.0853,0.1596) -- (0.0847,0.1776) -- (0.0842,0.1958) -- (0.0835,0.2141) -- (0.0828,0.2325) -- (0.0820,0.2511) -- (0.0812,0.2699) -- (0.0802,0.2888) -- (0.0793,0.3079) -- (0.0782,0.3271) -- (0.0770,0.3466) -- (0.0758,0.3664) -- (0.0745,0.3864) -- (0.0731,0.4066) -- (0.0716,0.4271) -- (0.0700,0.4480) -- (0.0684,0.4691) -- (0.0666,0.4906) -- (0.0646,0.5124) -- (0.0626,0.5346) -- (0.0605,0.5572) -- (0.0582,0.5803) -- (0.0557,0.6038) -- (0.0532,0.6278) -- (0.0504,0.6523) -- (0.0475,0.6773) -- (0.0444,0.7029) -- (0.0411,0.7292) -- (0.0376,0.7560) -- (0.0339,0.7836) -- (0.0300,0.8119) -- (0.0258,0.8410) -- (0.0213,0.8709) -- (0.0165,0.9017) -- (0.0113,0.9334) -- (0.0059,0.9662) -- (0.0000,1.0000);
\draw[wulffB] (0.0000,-1.0000) -- (0.0000,-0.9657) -- (0.0000,-0.9325) -- (0.0000,-0.9004) -- (0.0000,-0.8693) -- (0.0000,-0.8391) -- (0.0000,-0.8098) -- (0.0000,-0.7813) -- (0.0000,-0.7536) -- (0.0000,-0.7265) -- (0.0000,-0.7002) -- (0.0000,-0.6745) -- (0.0000,-0.6494) -- (0.0000,-0.6249) -- (0.0000,-0.6009) -- (0.0000,-0.5774) -- (0.0000,-0.5543) -- (0.0000,-0.5317) -- (0.0000,-0.5095) -- (0.0000,-0.4877) -- (0.0000,-0.4663) -- (0.0000,-0.4452) -- (0.0000,-0.4245) -- (0.0000,-0.4040) -- (0.0000,-0.3839) -- (0.0000,-0.3640) -- (0.0000,-0.3443) -- (0.0000,-0.3249) -- (0.0000,-0.3057) -- (0.0000,-0.2867) -- (0.0000,-0.2679) -- (0.0000,-0.2493) -- (0.0000,-0.2309) -- (0.0000,-0.2126) -- (0.0000,-0.1944) -- (0.0000,-0.1763) -- (0.0000,-0.1584) -- (0.0000,-0.1405) -- (0.0000,-0.1228) -- (0.0000,-0.1051) -- (0.0000,-0.0875) -- (0.0000,-0.0699) -- (0.0000,-0.0524) -- (0.0000,-0.0349) -- (0.0000,-0.0175) -- (0.0000,0.0000) -- (0.0000,0.0175) -- (0.0000,0.0349) -- (0.0000,0.0524) -- (0.0000,0.0699) -- (0.0000,0.0875) -- (0.0000,0.1051) -- (0.0000,0.1228) -- (0.0000,0.1405) -- (0.0000,0.1584) -- (0.0000,0.1763) -- (0.0000,0.1944) -- (0.0000,0.2126) -- (0.0000,0.2309) -- (0.0000,0.2493) -- (0.0000,0.2679) -- (0.0000,0.2867) -- (0.0000,0.3057) -- (0.0000,0.3249) -- (0.0000,0.3443) -- (0.0000,0.3640) -- (0.0000,0.3839) -- (0.0000,0.4040) -- (0.0000,0.4245) -- (0.0000,0.4452) -- (0.0000,0.4663) -- (0.0000,0.4877) -- (0.0000,0.5095) -- (0.0000,0.5317) -- (0.0000,0.5543) -- (0.0000,0.5774) -- (0.0000,0.6009) -- (0.0000,0.6249) -- (0.0000,0.6494) -- (0.0000,0.6745) -- (0.0000,0.7002) -- (0.0000,0.7265) -- (0.0000,0.7536) -- (0.0000,0.7813) -- (0.0000,0.8098) -- (0.0000,0.8391) -- (0.0000,0.8693) -- (0.0000,0.9004) -- (0.0000,0.9325) -- (0.0000,0.9657) -- (0.0000,1.0000);
\draw[wulffA] (-0.0000,-1.0000) -- (-0.0059,-0.9662) -- (-0.0113,-0.9334) -- (-0.0165,-0.9017) -- (-0.0213,-0.8709) -- (-0.0258,-0.8410) -- (-0.0300,-0.8119) -- (-0.0339,-0.7836) -- (-0.0376,-0.7560) -- (-0.0411,-0.7292) -- (-0.0444,-0.7029) -- (-0.0475,-0.6773) -- (-0.0504,-0.6523) -- (-0.0532,-0.6278) -- (-0.0557,-0.6038) -- (-0.0582,-0.5803) -- (-0.0605,-0.5572) -- (-0.0626,-0.5346) -- (-0.0646,-0.5124) -- (-0.0666,-0.4906) -- (-0.0684,-0.4691) -- (-0.0700,-0.4480) -- (-0.0716,-0.4271) -- (-0.0731,-0.4066) -- (-0.0745,-0.3864) -- (-0.0758,-0.3664) -- (-0.0770,-0.3466) -- (-0.0782,-0.3271) -- (-0.0793,-0.3079) -- (-0.0802,-0.2888) -- (-0.0812,-0.2699) -- (-0.0820,-0.2511) -- (-0.0828,-0.2325) -- (-0.0835,-0.2141) -- (-0.0842,-0.1958) -- (-0.0847,-0.1776) -- (-0.0853,-0.1596) -- (-0.0857,-0.1416) -- (-0.0862,-0.1237) -- (-0.0865,-0.1059) -- (-0.0868,-0.0882) -- (-0.0871,-0.0705) -- (-0.0872,-0.0528) -- (-0.0874,-0.0352) -- (-0.0875,-0.0176) -- (-0.0875,0.0000) -- (-0.0875,0.0176) -- (-0.0874,0.0352) -- (-0.0872,0.0528) -- (-0.0871,0.0705) -- (-0.0868,0.0882) -- (-0.0865,0.1059) -- (-0.0862,0.1237) -- (-0.0857,0.1416) -- (-0.0853,0.1596) -- (-0.0847,0.1776) -- (-0.0842,0.1958) -- (-0.0835,0.2141) -- (-0.0828,0.2325) -- (-0.0820,0.2511) -- (-0.0812,0.2699) -- (-0.0802,0.2888) -- (-0.0793,0.3079) -- (-0.0782,0.3271) -- (-0.0770,0.3466) -- (-0.0758,0.3664) -- (-0.0745,0.3864) -- (-0.0731,0.4066) -- (-0.0716,0.4271) -- (-0.0700,0.4480) -- (-0.0684,0.4691) -- (-0.0666,0.4906) -- (-0.0646,0.5124) -- (-0.0626,0.5346) -- (-0.0605,0.5572) -- (-0.0582,0.5803) -- (-0.0557,0.6038) -- (-0.0532,0.6278) -- (-0.0504,0.6523) -- (-0.0475,0.6773) -- (-0.0444,0.7029) -- (-0.0411,0.7292) -- (-0.0376,0.7560) -- (-0.0339,0.7836) -- (-0.0300,0.8119) -- (-0.0258,0.8410) -- (-0.0213,0.8709) -- (-0.0165,0.9017) -- (-0.0113,0.9334) -- (-0.0059,0.9662) -- (-0.0000,1.0000);
\draw[wulffA] (-0.0000,-1.0000) -- (-0.0116,-0.9677) -- (-0.0224,-0.9362) -- (-0.0326,-0.9056) -- (-0.0421,-0.8757) -- (-0.0511,-0.8467) -- (-0.0595,-0.8183) -- (-0.0674,-0.7906) -- (-0.0749,-0.7635) -- (-0.0819,-0.7370) -- (-0.0885,-0.7111) -- (-0.0948,-0.6858) -- (-0.1006,-0.6609) -- (-0.1062,-0.6366) -- (-0.1114,-0.6127) -- (-0.1163,-0.5892) -- (-0.1210,-0.5661) -- (-0.1254,-0.5435) -- (-0.1295,-0.5212) -- (-0.1334,-0.4992) -- (-0.1371,-0.4776) -- (-0.1405,-0.4563) -- (-0.1438,-0.4352) -- (-0.1468,-0.4145) -- (-0.1497,-0.3940) -- (-0.1523,-0.3737) -- (-0.1549,-0.3537) -- (-0.1572,-0.3339) -- (-0.1594,-0.3143) -- (-0.1614,-0.2949) -- (-0.1633,-0.2757) -- (-0.1650,-0.2566) -- (-0.1667,-0.2377) -- (-0.1681,-0.2189) -- (-0.1695,-0.2002) -- (-0.1707,-0.1816) -- (-0.1718,-0.1632) -- (-0.1727,-0.1448) -- (-0.1736,-0.1265) -- (-0.1743,-0.1083) -- (-0.1749,-0.0902) -- (-0.1754,-0.0721) -- (-0.1758,-0.0540) -- (-0.1761,-0.0360) -- (-0.1763,-0.0180) -- (-0.1763,0.0000) -- (-0.1763,0.0180) -- (-0.1761,0.0360) -- (-0.1758,0.0540) -- (-0.1754,0.0721) -- (-0.1749,0.0902) -- (-0.1743,0.1083) -- (-0.1736,0.1265) -- (-0.1727,0.1448) -- (-0.1718,0.1632) -- (-0.1707,0.1816) -- (-0.1695,0.2002) -- (-0.1681,0.2189) -- (-0.1667,0.2377) -- (-0.1650,0.2566) -- (-0.1633,0.2757) -- (-0.1614,0.2949) -- (-0.1594,0.3143) -- (-0.1572,0.3339) -- (-0.1549,0.3537) -- (-0.1523,0.3737) -- (-0.1497,0.3940) -- (-0.1468,0.4145) -- (-0.1438,0.4352) -- (-0.1405,0.4563) -- (-0.1371,0.4776) -- (-0.1334,0.4992) -- (-0.1295,0.5212) -- (-0.1254,0.5435) -- (-0.1210,0.5661) -- (-0.1163,0.5892) -- (-0.1114,0.6127) -- (-0.1062,0.6366) -- (-0.1006,0.6609) -- (-0.0948,0.6858) -- (-0.0885,0.7111) -- (-0.0819,0.7370) -- (-0.0749,0.7635) -- (-0.0674,0.7906) -- (-0.0595,0.8183) -- (-0.0511,0.8467) -- (-0.0421,0.8757) -- (-0.0326,0.9056) -- (-0.0224,0.9362) -- (-0.0116,0.9677) -- (-0.0000,1.0000);
\draw[wulffB] (-0.0000,-1.0000) -- (-0.0169,-0.9701) -- (-0.0329,-0.9407) -- (-0.0479,-0.9120) -- (-0.0621,-0.8838) -- (-0.0755,-0.8561) -- (-0.0881,-0.8289) -- (-0.1000,-0.8022) -- (-0.1113,-0.7760) -- (-0.1219,-0.7503) -- (-0.1319,-0.7250) -- (-0.1414,-0.7001) -- (-0.1504,-0.6756) -- (-0.1589,-0.6515) -- (-0.1669,-0.6277) -- (-0.1745,-0.6043) -- (-0.1816,-0.5813) -- (-0.1884,-0.5585) -- (-0.1948,-0.5361) -- (-0.2008,-0.5140) -- (-0.2065,-0.4921) -- (-0.2118,-0.4705) -- (-0.2169,-0.4491) -- (-0.2216,-0.4280) -- (-0.2261,-0.4071) -- (-0.2303,-0.3864) -- (-0.2342,-0.3659) -- (-0.2379,-0.3456) -- (-0.2413,-0.3255) -- (-0.2445,-0.3055) -- (-0.2474,-0.2857) -- (-0.2502,-0.2660) -- (-0.2527,-0.2465) -- (-0.2550,-0.2271) -- (-0.2571,-0.2078) -- (-0.2590,-0.1886) -- (-0.2608,-0.1695) -- (-0.2623,-0.1504) -- (-0.2636,-0.1315) -- (-0.2648,-0.1126) -- (-0.2658,-0.0937) -- (-0.2665,-0.0749) -- (-0.2672,-0.0562) -- (-0.2676,-0.0374) -- (-0.2679,-0.0187) -- (-0.2679,0.0000) -- (-0.2679,0.0187) -- (-0.2676,0.0374) -- (-0.2672,0.0562) -- (-0.2665,0.0749) -- (-0.2658,0.0937) -- (-0.2648,0.1126) -- (-0.2636,0.1315) -- (-0.2623,0.1504) -- (-0.2608,0.1695) -- (-0.2590,0.1886) -- (-0.2571,0.2078) -- (-0.2550,0.2271) -- (-0.2527,0.2465) -- (-0.2502,0.2660) -- (-0.2474,0.2857) -- (-0.2445,0.3055) -- (-0.2413,0.3255) -- (-0.2379,0.3456) -- (-0.2342,0.3659) -- (-0.2303,0.3864) -- (-0.2261,0.4071) -- (-0.2216,0.4280) -- (-0.2169,0.4491) -- (-0.2118,0.4705) -- (-0.2065,0.4921) -- (-0.2008,0.5140) -- (-0.1948,0.5361) -- (-0.1884,0.5585) -- (-0.1816,0.5813) -- (-0.1745,0.6043) -- (-0.1669,0.6277) -- (-0.1589,0.6515) -- (-0.1504,0.6756) -- (-0.1414,0.7001) -- (-0.1319,0.7250) -- (-0.1219,0.7503) -- (-0.1113,0.7760) -- (-0.1000,0.8022) -- (-0.0881,0.8289) -- (-0.0755,0.8561) -- (-0.0621,0.8838) -- (-0.0479,0.9120) -- (-0.0329,0.9407) -- (-0.0169,0.9701) -- (-0.0000,1.0000);
\draw[wulffA] (-0.0000,-1.0000) -- (-0.0218,-0.9734) -- (-0.0426,-0.9470) -- (-0.0622,-0.9208) -- (-0.0808,-0.8949) -- (-0.0985,-0.8692) -- (-0.1153,-0.8438) -- (-0.1312,-0.8186) -- (-0.1463,-0.7937) -- (-0.1606,-0.7690) -- (-0.1742,-0.7446) -- (-0.1871,-0.7204) -- (-0.1993,-0.6965) -- (-0.2109,-0.6728) -- (-0.2219,-0.6494) -- (-0.2324,-0.6262) -- (-0.2423,-0.6032) -- (-0.2516,-0.5804) -- (-0.2605,-0.5578) -- (-0.2689,-0.5355) -- (-0.2769,-0.5133) -- (-0.2844,-0.4913) -- (-0.2914,-0.4695) -- (-0.2981,-0.4479) -- (-0.3044,-0.4264) -- (-0.3103,-0.4051) -- (-0.3159,-0.3839) -- (-0.3211,-0.3629) -- (-0.3259,-0.3420) -- (-0.3304,-0.3212) -- (-0.3347,-0.3006) -- (-0.3386,-0.2801) -- (-0.3422,-0.2596) -- (-0.3455,-0.2393) -- (-0.3485,-0.2190) -- (-0.3512,-0.1989) -- (-0.3537,-0.1788) -- (-0.3559,-0.1587) -- (-0.3578,-0.1388) -- (-0.3594,-0.1189) -- (-0.3608,-0.0990) -- (-0.3620,-0.0791) -- (-0.3628,-0.0593) -- (-0.3635,-0.0395) -- (-0.3638,-0.0198) -- (-0.3640,0.0000) -- (-0.3638,0.0198) -- (-0.3635,0.0395) -- (-0.3628,0.0593) -- (-0.3620,0.0791) -- (-0.3608,0.0990) -- (-0.3594,0.1189) -- (-0.3578,0.1388) -- (-0.3559,0.1587) -- (-0.3537,0.1788) -- (-0.3512,0.1989) -- (-0.3485,0.2190) -- (-0.3455,0.2393) -- (-0.3422,0.2596) -- (-0.3386,0.2801) -- (-0.3347,0.3006) -- (-0.3304,0.3212) -- (-0.3259,0.3420) -- (-0.3211,0.3629) -- (-0.3159,0.3839) -- (-0.3103,0.4051) -- (-0.3044,0.4264) -- (-0.2981,0.4479) -- (-0.2914,0.4695) -- (-0.2844,0.4913) -- (-0.2769,0.5133) -- (-0.2689,0.5355) -- (-0.2605,0.5578) -- (-0.2516,0.5804) -- (-0.2423,0.6032) -- (-0.2324,0.6262) -- (-0.2219,0.6494) -- (-0.2109,0.6728) -- (-0.1993,0.6965) -- (-0.1871,0.7204) -- (-0.1742,0.7446) -- (-0.1606,0.7690) -- (-0.1463,0.7937) -- (-0.1312,0.8186) -- (-0.1153,0.8438) -- (-0.0985,0.8692) -- (-0.0808,0.8949) -- (-0.0622,0.9208) -- (-0.0426,0.9470) -- (-0.0218,0.9734) -- (-0.0000,1.0000);
\draw[wulffA] (-0.0000,-1.0000) -- (-0.0261,-0.9775) -- (-0.0511,-0.9548) -- (-0.0750,-0.9319) -- (-0.0979,-0.9090) -- (-0.1197,-0.8859) -- (-0.1405,-0.8628) -- (-0.1604,-0.8397) -- (-0.1794,-0.8166) -- (-0.1975,-0.7935) -- (-0.2148,-0.7703) -- (-0.2313,-0.7473) -- (-0.2470,-0.7242) -- (-0.2620,-0.7012) -- (-0.2763,-0.6783) -- (-0.2899,-0.6554) -- (-0.3028,-0.6326) -- (-0.3151,-0.6098) -- (-0.3268,-0.5872) -- (-0.3379,-0.5646) -- (-0.3484,-0.5421) -- (-0.3584,-0.5196) -- (-0.3679,-0.4973) -- (-0.3768,-0.4750) -- (-0.3853,-0.4528) -- (-0.3932,-0.4307) -- (-0.4007,-0.4087) -- (-0.4077,-0.3867) -- (-0.4143,-0.3648) -- (-0.4204,-0.3430) -- (-0.4262,-0.3212) -- (-0.4315,-0.2995) -- (-0.4364,-0.2778) -- (-0.4409,-0.2563) -- (-0.4450,-0.2347) -- (-0.4488,-0.2132) -- (-0.4521,-0.1918) -- (-0.4551,-0.1704) -- (-0.4578,-0.1490) -- (-0.4601,-0.1277) -- (-0.4620,-0.1063) -- (-0.4635,-0.0850) -- (-0.4647,-0.0638) -- (-0.4656,-0.0425) -- (-0.4661,-0.0212) -- (-0.4663,0.0000) -- (-0.4661,0.0212) -- (-0.4656,0.0425) -- (-0.4647,0.0638) -- (-0.4635,0.0850) -- (-0.4620,0.1063) -- (-0.4601,0.1277) -- (-0.4578,0.1490) -- (-0.4551,0.1704) -- (-0.4521,0.1918) -- (-0.4488,0.2132) -- (-0.4450,0.2347) -- (-0.4409,0.2563) -- (-0.4364,0.2778) -- (-0.4315,0.2995) -- (-0.4262,0.3212) -- (-0.4204,0.3430) -- (-0.4143,0.3648) -- (-0.4077,0.3867) -- (-0.4007,0.4087) -- (-0.3932,0.4307) -- (-0.3853,0.4528) -- (-0.3768,0.4750) -- (-0.3679,0.4973) -- (-0.3584,0.5196) -- (-0.3484,0.5421) -- (-0.3379,0.5646) -- (-0.3268,0.5872) -- (-0.3151,0.6098) -- (-0.3028,0.6326) -- (-0.2899,0.6554) -- (-0.2763,0.6783) -- (-0.2620,0.7012) -- (-0.2470,0.7242) -- (-0.2313,0.7473) -- (-0.2148,0.7703) -- (-0.1975,0.7935) -- (-0.1794,0.8166) -- (-0.1604,0.8397) -- (-0.1405,0.8628) -- (-0.1197,0.8859) -- (-0.0979,0.9090) -- (-0.0750,0.9319) -- (-0.0511,0.9548) -- (-0.0261,0.9775) -- (-0.0000,1.0000);
\draw[wulffB] (-0.0000,-1.0000) -- (-0.0297,-0.9823) -- (-0.0584,-0.9639) -- (-0.0860,-0.9451) -- (-0.1127,-0.9258) -- (-0.1384,-0.9061) -- (-0.1631,-0.8860) -- (-0.1869,-0.8656) -- (-0.2098,-0.8448) -- (-0.2318,-0.8238) -- (-0.2529,-0.8025) -- (-0.2732,-0.7809) -- (-0.2927,-0.7592) -- (-0.3114,-0.7372) -- (-0.3293,-0.7151) -- (-0.3464,-0.6928) -- (-0.3628,-0.6704) -- (-0.3785,-0.6479) -- (-0.3934,-0.6253) -- (-0.4077,-0.6025) -- (-0.4213,-0.5797) -- (-0.4342,-0.5568) -- (-0.4465,-0.5339) -- (-0.4582,-0.5109) -- (-0.4692,-0.4879) -- (-0.4797,-0.4648) -- (-0.4896,-0.4416) -- (-0.4988,-0.4185) -- (-0.5076,-0.3953) -- (-0.5157,-0.3721) -- (-0.5234,-0.3489) -- (-0.5305,-0.3257) -- (-0.5370,-0.3025) -- (-0.5431,-0.2792) -- (-0.5486,-0.2560) -- (-0.5537,-0.2327) -- (-0.5582,-0.2094) -- (-0.5622,-0.1862) -- (-0.5658,-0.1629) -- (-0.5689,-0.1396) -- (-0.5715,-0.1164) -- (-0.5736,-0.0931) -- (-0.5752,-0.0698) -- (-0.5764,-0.0465) -- (-0.5771,-0.0233) -- (-0.5774,0.0000) -- (-0.5771,0.0233) -- (-0.5764,0.0465) -- (-0.5752,0.0698) -- (-0.5736,0.0931) -- (-0.5715,0.1164) -- (-0.5689,0.1396) -- (-0.5658,0.1629) -- (-0.5622,0.1862) -- (-0.5582,0.2094) -- (-0.5537,0.2327) -- (-0.5486,0.2560) -- (-0.5431,0.2792) -- (-0.5370,0.3025) -- (-0.5305,0.3257) -- (-0.5234,0.3489) -- (-0.5157,0.3721) -- (-0.5076,0.3953) -- (-0.4988,0.4185) -- (-0.4896,0.4416) -- (-0.4797,0.4648) -- (-0.4692,0.4879) -- (-0.4582,0.5109) -- (-0.4465,0.5339) -- (-0.4342,0.5568) -- (-0.4213,0.5797) -- (-0.4077,0.6025) -- (-0.3934,0.6253) -- (-0.3785,0.6479) -- (-0.3628,0.6704) -- (-0.3464,0.6928) -- (-0.3293,0.7151) -- (-0.3114,0.7372) -- (-0.2927,0.7592) -- (-0.2732,0.7809) -- (-0.2529,0.8025) -- (-0.2318,0.8238) -- (-0.2098,0.8448) -- (-0.1869,0.8656) -- (-0.1631,0.8860) -- (-0.1384,0.9061) -- (-0.1127,0.9258) -- (-0.0860,0.9451) -- (-0.0584,0.9639) -- (-0.0297,0.9823) -- (-0.0000,1.0000);
\draw[wulffA] (-0.0000,-1.0000) -- (-0.0324,-0.9876) -- (-0.0640,-0.9743) -- (-0.0948,-0.9602) -- (-0.1248,-0.9453) -- (-0.1540,-0.9296) -- (-0.1824,-0.9132) -- (-0.2100,-0.8961) -- (-0.2367,-0.8784) -- (-0.2626,-0.8601) -- (-0.2877,-0.8413) -- (-0.3120,-0.8219) -- (-0.3355,-0.8020) -- (-0.3582,-0.7816) -- (-0.3801,-0.7608) -- (-0.4012,-0.7396) -- (-0.4216,-0.7179) -- (-0.4411,-0.6959) -- (-0.4599,-0.6736) -- (-0.4779,-0.6509) -- (-0.4952,-0.6280) -- (-0.5117,-0.6047) -- (-0.5274,-0.5812) -- (-0.5425,-0.5575) -- (-0.5568,-0.5335) -- (-0.5704,-0.5093) -- (-0.5833,-0.4850) -- (-0.5955,-0.4604) -- (-0.6069,-0.4357) -- (-0.6177,-0.4108) -- (-0.6278,-0.3857) -- (-0.6373,-0.3606) -- (-0.6460,-0.3353) -- (-0.6541,-0.3099) -- (-0.6615,-0.2844) -- (-0.6683,-0.2588) -- (-0.6743,-0.2332) -- (-0.6798,-0.2074) -- (-0.6846,-0.1816) -- (-0.6887,-0.1558) -- (-0.6923,-0.1299) -- (-0.6951,-0.1040) -- (-0.6973,-0.0780) -- (-0.6989,-0.0520) -- (-0.6999,-0.0260) -- (-0.7002,0.0000) -- (-0.6999,0.0260) -- (-0.6989,0.0520) -- (-0.6973,0.0780) -- (-0.6951,0.1040) -- (-0.6923,0.1299) -- (-0.6887,0.1558) -- (-0.6846,0.1816) -- (-0.6798,0.2074) -- (-0.6743,0.2332) -- (-0.6683,0.2588) -- (-0.6615,0.2844) -- (-0.6541,0.3099) -- (-0.6460,0.3353) -- (-0.6373,0.3606) -- (-0.6278,0.3857) -- (-0.6177,0.4108) -- (-0.6069,0.4357) -- (-0.5955,0.4604) -- (-0.5833,0.4850) -- (-0.5704,0.5093) -- (-0.5568,0.5335) -- (-0.5425,0.5575) -- (-0.5274,0.5812) -- (-0.5117,0.6047) -- (-0.4952,0.6280) -- (-0.4779,0.6509) -- (-0.4599,0.6736) -- (-0.4411,0.6959) -- (-0.4216,0.7179) -- (-0.4012,0.7396) -- (-0.3801,0.7608) -- (-0.3582,0.7816) -- (-0.3355,0.8020) -- (-0.3120,0.8219) -- (-0.2877,0.8413) -- (-0.2626,0.8601) -- (-0.2367,0.8784) -- (-0.2100,0.8961) -- (-0.1824,0.9132) -- (-0.1540,0.9296) -- (-0.1248,0.9453) -- (-0.0948,0.9602) -- (-0.0640,0.9743) -- (-0.0324,0.9876) -- (-0.0000,1.0000);
\draw[wulffA] (-0.0000,-1.0000) -- (-0.0342,-0.9934) -- (-0.0679,-0.9856) -- (-0.1011,-0.9768) -- (-0.1338,-0.9669) -- (-0.1660,-0.9560) -- (-0.1976,-0.9441) -- (-0.2286,-0.9312) -- (-0.2591,-0.9174) -- (-0.2888,-0.9026) -- (-0.3179,-0.8870) -- (-0.3464,-0.8706) -- (-0.3741,-0.8533) -- (-0.4012,-0.8352) -- (-0.4275,-0.8164) -- (-0.4531,-0.7968) -- (-0.4779,-0.7766) -- (-0.5020,-0.7557) -- (-0.5252,-0.7341) -- (-0.5477,-0.7119) -- (-0.5695,-0.6891) -- (-0.5904,-0.6658) -- (-0.6105,-0.6419) -- (-0.6297,-0.6175) -- (-0.6482,-0.5927) -- (-0.6658,-0.5673) -- (-0.6826,-0.5416) -- (-0.6986,-0.5154) -- (-0.7137,-0.4888) -- (-0.7280,-0.4619) -- (-0.7414,-0.4346) -- (-0.7539,-0.4071) -- (-0.7656,-0.3792) -- (-0.7765,-0.3510) -- (-0.7865,-0.3227) -- (-0.7956,-0.2940) -- (-0.8039,-0.2652) -- (-0.8112,-0.2362) -- (-0.8178,-0.2070) -- (-0.8234,-0.1777) -- (-0.8282,-0.1483) -- (-0.8321,-0.1188) -- (-0.8352,-0.0891) -- (-0.8374,-0.0595) -- (-0.8387,-0.0297) -- (-0.8391,0.0000) -- (-0.8387,0.0297) -- (-0.8374,0.0595) -- (-0.8352,0.0891) -- (-0.8321,0.1188) -- (-0.8282,0.1483) -- (-0.8234,0.1777) -- (-0.8178,0.2070) -- (-0.8112,0.2362) -- (-0.8039,0.2652) -- (-0.7956,0.2940) -- (-0.7865,0.3227) -- (-0.7765,0.3510) -- (-0.7656,0.3792) -- (-0.7539,0.4071) -- (-0.7414,0.4346) -- (-0.7280,0.4619) -- (-0.7137,0.4888) -- (-0.6986,0.5154) -- (-0.6826,0.5416) -- (-0.6658,0.5673) -- (-0.6482,0.5927) -- (-0.6297,0.6175) -- (-0.6105,0.6419) -- (-0.5904,0.6658) -- (-0.5695,0.6891) -- (-0.5477,0.7119) -- (-0.5252,0.7341) -- (-0.5020,0.7557) -- (-0.4779,0.7766) -- (-0.4531,0.7968) -- (-0.4275,0.8164) -- (-0.4012,0.8352) -- (-0.3741,0.8533) -- (-0.3464,0.8706) -- (-0.3179,0.8870) -- (-0.2888,0.9026) -- (-0.2591,0.9174) -- (-0.2286,0.9312) -- (-0.1976,0.9441) -- (-0.1660,0.9560) -- (-0.1338,0.9669) -- (-0.1011,0.9768) -- (-0.0679,0.9856) -- (-0.0342,0.9934) -- (-0.0000,1.0000);
\draw[wulffA] (0.1736,0.9848) -- (0.1725,0.9789) -- (0.1712,0.9730) -- (0.1696,0.9673) -- (0.1679,0.9616) -- (0.1660,0.9560) -- (0.1639,0.9505) -- (0.1617,0.9451) -- (0.1593,0.9398) -- (0.1567,0.9347) -- (0.1540,0.9296) -- (0.1512,0.9247) -- (0.1482,0.9198) -- (0.1450,0.9151) -- (0.1418,0.9106) -- (0.1384,0.9061) -- (0.1349,0.9018) -- (0.1312,0.8976) -- (0.1275,0.8936) -- (0.1236,0.8897) -- (0.1197,0.8859) -- (0.1156,0.8823) -- (0.1115,0.8788) -- (0.1072,0.8755) -- (0.1029,0.8722) -- (0.0985,0.8692) -- (0.0940,0.8663) -- (0.0895,0.8635) -- (0.0849,0.8609) -- (0.0802,0.8584) -- (0.0755,0.8561) -- (0.0707,0.8539) -- (0.0658,0.8519) -- (0.0610,0.8500) -- (0.0560,0.8482) -- (0.0511,0.8467) -- (0.0461,0.8452) -- (0.0410,0.8439) -- (0.0360,0.8428) -- (0.0309,0.8418) -- (0.0258,0.8410) -- (0.0206,0.8403) -- (0.0155,0.8398) -- (0.0103,0.8394) -- (0.0052,0.8392) -- (0.0000,0.8391) -- (-0.0052,0.8392) -- (-0.0103,0.8394) -- (-0.0155,0.8398) -- (-0.0206,0.8403) -- (-0.0258,0.8410) -- (-0.0309,0.8418) -- (-0.0360,0.8428) -- (-0.0410,0.8439) -- (-0.0461,0.8452) -- (-0.0511,0.8467) -- (-0.0560,0.8482) -- (-0.0610,0.8500) -- (-0.0658,0.8519) -- (-0.0707,0.8539) -- (-0.0755,0.8561) -- (-0.0802,0.8584) -- (-0.0849,0.8609) -- (-0.0895,0.8635) -- (-0.0940,0.8663) -- (-0.0985,0.8692) -- (-0.1029,0.8722) -- (-0.1072,0.8755) -- (-0.1115,0.8788) -- (-0.1156,0.8823) -- (-0.1197,0.8859) -- (-0.1236,0.8897) -- (-0.1275,0.8936) -- (-0.1312,0.8976) -- (-0.1349,0.9018) -- (-0.1384,0.9061) -- (-0.1418,0.9106) -- (-0.1450,0.9151) -- (-0.1482,0.9198) -- (-0.1512,0.9247) -- (-0.1540,0.9296) -- (-0.1567,0.9347) -- (-0.1593,0.9398) -- (-0.1617,0.9451) -- (-0.1639,0.9505) -- (-0.1660,0.9560) -- (-0.1679,0.9616) -- (-0.1696,0.9673) -- (-0.1712,0.9730) -- (-0.1725,0.9789) -- (-0.1736,0.9848);
\draw[wulffA] (0.1736,-0.9848) -- (0.1725,-0.9789) -- (0.1712,-0.9730) -- (0.1696,-0.9673) -- (0.1679,-0.9616) -- (0.1660,-0.9560) -- (0.1639,-0.9505) -- (0.1617,-0.9451) -- (0.1593,-0.9398) -- (0.1567,-0.9347) -- (0.1540,-0.9296) -- (0.1512,-0.9247) -- (0.1482,-0.9198) -- (0.1450,-0.9151) -- (0.1418,-0.9106) -- (0.1384,-0.9061) -- (0.1349,-0.9018) -- (0.1312,-0.8976) -- (0.1275,-0.8936) -- (0.1236,-0.8897) -- (0.1197,-0.8859) -- (0.1156,-0.8823) -- (0.1115,-0.8788) -- (0.1072,-0.8755) -- (0.1029,-0.8722) -- (0.0985,-0.8692) -- (0.0940,-0.8663) -- (0.0895,-0.8635) -- (0.0849,-0.8609) -- (0.0802,-0.8584) -- (0.0755,-0.8561) -- (0.0707,-0.8539) -- (0.0658,-0.8519) -- (0.0610,-0.8500) -- (0.0560,-0.8482) -- (0.0511,-0.8467) -- (0.0461,-0.8452) -- (0.0410,-0.8439) -- (0.0360,-0.8428) -- (0.0309,-0.8418) -- (0.0258,-0.8410) -- (0.0206,-0.8403) -- (0.0155,-0.8398) -- (0.0103,-0.8394) -- (0.0052,-0.8392) -- (0.0000,-0.8391) -- (-0.0052,-0.8392) -- (-0.0103,-0.8394) -- (-0.0155,-0.8398) -- (-0.0206,-0.8403) -- (-0.0258,-0.8410) -- (-0.0309,-0.8418) -- (-0.0360,-0.8428) -- (-0.0410,-0.8439) -- (-0.0461,-0.8452) -- (-0.0511,-0.8467) -- (-0.0560,-0.8482) -- (-0.0610,-0.8500) -- (-0.0658,-0.8519) -- (-0.0707,-0.8539) -- (-0.0755,-0.8561) -- (-0.0802,-0.8584) -- (-0.0849,-0.8609) -- (-0.0895,-0.8635) -- (-0.0940,-0.8663) -- (-0.0985,-0.8692) -- (-0.1029,-0.8722) -- (-0.1072,-0.8755) -- (-0.1115,-0.8788) -- (-0.1156,-0.8823) -- (-0.1197,-0.8859) -- (-0.1236,-0.8897) -- (-0.1275,-0.8936) -- (-0.1312,-0.8976) -- (-0.1349,-0.9018) -- (-0.1384,-0.9061) -- (-0.1418,-0.9106) -- (-0.1450,-0.9151) -- (-0.1482,-0.9198) -- (-0.1512,-0.9247) -- (-0.1540,-0.9296) -- (-0.1567,-0.9347) -- (-0.1593,-0.9398) -- (-0.1617,-0.9451) -- (-0.1639,-0.9505) -- (-0.1660,-0.9560) -- (-0.1679,-0.9616) -- (-0.1696,-0.9673) -- (-0.1712,-0.9730) -- (-0.1725,-0.9789) -- (-0.1736,-0.9848);
\draw[wulffA] (0.3420,0.9397) -- (0.3378,0.9286) -- (0.3332,0.9178) -- (0.3284,0.9073) -- (0.3233,0.8970) -- (0.3179,0.8870) -- (0.3123,0.8773) -- (0.3065,0.8679) -- (0.3004,0.8587) -- (0.2942,0.8499) -- (0.2877,0.8413) -- (0.2811,0.8330) -- (0.2743,0.8249) -- (0.2673,0.8172) -- (0.2602,0.8097) -- (0.2529,0.8025) -- (0.2455,0.7955) -- (0.2380,0.7888) -- (0.2304,0.7824) -- (0.2226,0.7762) -- (0.2148,0.7703) -- (0.2068,0.7647) -- (0.1988,0.7593) -- (0.1907,0.7542) -- (0.1825,0.7493) -- (0.1742,0.7446) -- (0.1659,0.7402) -- (0.1575,0.7360) -- (0.1490,0.7321) -- (0.1405,0.7284) -- (0.1319,0.7250) -- (0.1233,0.7217) -- (0.1147,0.7187) -- (0.1060,0.7160) -- (0.0973,0.7134) -- (0.0885,0.7111) -- (0.0797,0.7091) -- (0.0709,0.7072) -- (0.0621,0.7055) -- (0.0533,0.7041) -- (0.0444,0.7029) -- (0.0356,0.7019) -- (0.0267,0.7012) -- (0.0178,0.7006) -- (0.0089,0.7003) -- (0.0000,0.7002) -- (-0.0089,0.7003) -- (-0.0178,0.7006) -- (-0.0267,0.7012) -- (-0.0356,0.7019) -- (-0.0444,0.7029) -- (-0.0533,0.7041) -- (-0.0621,0.7055) -- (-0.0709,0.7072) -- (-0.0797,0.7091) -- (-0.0885,0.7111) -- (-0.0973,0.7134) -- (-0.1060,0.7160) -- (-0.1147,0.7187) -- (-0.1233,0.7217) -- (-0.1319,0.7250) -- (-0.1405,0.7284) -- (-0.1490,0.7321) -- (-0.1575,0.7360) -- (-0.1659,0.7402) -- (-0.1742,0.7446) -- (-0.1825,0.7493) -- (-0.1907,0.7542) -- (-0.1988,0.7593) -- (-0.2068,0.7647) -- (-0.2148,0.7703) -- (-0.2226,0.7762) -- (-0.2304,0.7824) -- (-0.2380,0.7888) -- (-0.2455,0.7955) -- (-0.2529,0.8025) -- (-0.2602,0.8097) -- (-0.2673,0.8172) -- (-0.2743,0.8249) -- (-0.2811,0.8330) -- (-0.2877,0.8413) -- (-0.2942,0.8499) -- (-0.3004,0.8587) -- (-0.3065,0.8679) -- (-0.3123,0.8773) -- (-0.3179,0.8870) -- (-0.3233,0.8970) -- (-0.3284,0.9073) -- (-0.3332,0.9178) -- (-0.3378,0.9286) -- (-0.3420,0.9397);
\draw[wulffA] (0.3420,-0.9397) -- (0.3378,-0.9286) -- (0.3332,-0.9178) -- (0.3284,-0.9073) -- (0.3233,-0.8970) -- (0.3179,-0.8870) -- (0.3123,-0.8773) -- (0.3065,-0.8679) -- (0.3004,-0.8587) -- (0.2942,-0.8499) -- (0.2877,-0.8413) -- (0.2811,-0.8330) -- (0.2743,-0.8249) -- (0.2673,-0.8172) -- (0.2602,-0.8097) -- (0.2529,-0.8025) -- (0.2455,-0.7955) -- (0.2380,-0.7888) -- (0.2304,-0.7824) -- (0.2226,-0.7762) -- (0.2148,-0.7703) -- (0.2068,-0.7647) -- (0.1988,-0.7593) -- (0.1907,-0.7542) -- (0.1825,-0.7493) -- (0.1742,-0.7446) -- (0.1659,-0.7402) -- (0.1575,-0.7360) -- (0.1490,-0.7321) -- (0.1405,-0.7284) -- (0.1319,-0.7250) -- (0.1233,-0.7217) -- (0.1147,-0.7187) -- (0.1060,-0.7160) -- (0.0973,-0.7134) -- (0.0885,-0.7111) -- (0.0797,-0.7091) -- (0.0709,-0.7072) -- (0.0621,-0.7055) -- (0.0533,-0.7041) -- (0.0444,-0.7029) -- (0.0356,-0.7019) -- (0.0267,-0.7012) -- (0.0178,-0.7006) -- (0.0089,-0.7003) -- (0.0000,-0.7002) -- (-0.0089,-0.7003) -- (-0.0178,-0.7006) -- (-0.0267,-0.7012) -- (-0.0356,-0.7019) -- (-0.0444,-0.7029) -- (-0.0533,-0.7041) -- (-0.0621,-0.7055) -- (-0.0709,-0.7072) -- (-0.0797,-0.7091) -- (-0.0885,-0.7111) -- (-0.0973,-0.7134) -- (-0.1060,-0.7160) -- (-0.1147,-0.7187) -- (-0.1233,-0.7217) -- (-0.1319,-0.7250) -- (-0.1405,-0.7284) -- (-0.1490,-0.7321) -- (-0.1575,-0.7360) -- (-0.1659,-0.7402) -- (-0.1742,-0.7446) -- (-0.1825,-0.7493) -- (-0.1907,-0.7542) -- (-0.1988,-0.7593) -- (-0.2068,-0.7647) -- (-0.2148,-0.7703) -- (-0.2226,-0.7762) -- (-0.2304,-0.7824) -- (-0.2380,-0.7888) -- (-0.2455,-0.7955) -- (-0.2529,-0.8025) -- (-0.2602,-0.8097) -- (-0.2673,-0.8172) -- (-0.2743,-0.8249) -- (-0.2811,-0.8330) -- (-0.2877,-0.8413) -- (-0.2942,-0.8499) -- (-0.3004,-0.8587) -- (-0.3065,-0.8679) -- (-0.3123,-0.8773) -- (-0.3179,-0.8870) -- (-0.3233,-0.8970) -- (-0.3284,-0.9073) -- (-0.3332,-0.9178) -- (-0.3378,-0.9286) -- (-0.3420,-0.9397);
\draw[wulffB] (0.5000,0.8660) -- (0.4911,0.8512) -- (0.4820,0.8368) -- (0.4726,0.8230) -- (0.4629,0.8097) -- (0.4531,0.7968) -- (0.4430,0.7845) -- (0.4328,0.7726) -- (0.4224,0.7611) -- (0.4119,0.7501) -- (0.4012,0.7396) -- (0.3905,0.7294) -- (0.3796,0.7197) -- (0.3686,0.7103) -- (0.3575,0.7014) -- (0.3464,0.6928) -- (0.3352,0.6846) -- (0.3239,0.6768) -- (0.3126,0.6693) -- (0.3013,0.6622) -- (0.2899,0.6554) -- (0.2784,0.6489) -- (0.2670,0.6428) -- (0.2555,0.6369) -- (0.2439,0.6314) -- (0.2324,0.6262) -- (0.2208,0.6212) -- (0.2092,0.6166) -- (0.1977,0.6122) -- (0.1861,0.6082) -- (0.1745,0.6043) -- (0.1628,0.6008) -- (0.1512,0.5975) -- (0.1396,0.5945) -- (0.1280,0.5917) -- (0.1163,0.5892) -- (0.1047,0.5869) -- (0.0931,0.5849) -- (0.0814,0.5831) -- (0.0698,0.5816) -- (0.0582,0.5803) -- (0.0465,0.5792) -- (0.0349,0.5784) -- (0.0233,0.5778) -- (0.0116,0.5775) -- (0.0000,0.5774) -- (-0.0116,0.5775) -- (-0.0233,0.5778) -- (-0.0349,0.5784) -- (-0.0465,0.5792) -- (-0.0582,0.5803) -- (-0.0698,0.5816) -- (-0.0814,0.5831) -- (-0.0931,0.5849) -- (-0.1047,0.5869) -- (-0.1163,0.5892) -- (-0.1280,0.5917) -- (-0.1396,0.5945) -- (-0.1512,0.5975) -- (-0.1628,0.6008) -- (-0.1745,0.6043) -- (-0.1861,0.6082) -- (-0.1977,0.6122) -- (-0.2092,0.6166) -- (-0.2208,0.6212) -- (-0.2324,0.6262) -- (-0.2439,0.6314) -- (-0.2555,0.6369) -- (-0.2670,0.6428) -- (-0.2784,0.6489) -- (-0.2899,0.6554) -- (-0.3013,0.6622) -- (-0.3126,0.6693) -- (-0.3239,0.6768) -- (-0.3352,0.6846) -- (-0.3464,0.6928) -- (-0.3575,0.7014) -- (-0.3686,0.7103) -- (-0.3796,0.7197) -- (-0.3905,0.7294) -- (-0.4012,0.7396) -- (-0.4119,0.7501) -- (-0.4224,0.7611) -- (-0.4328,0.7726) -- (-0.4430,0.7845) -- (-0.4531,0.7968) -- (-0.4629,0.8097) -- (-0.4726,0.8230) -- (-0.4820,0.8368) -- (-0.4911,0.8512) -- (-0.5000,0.8660);
\draw[wulffB] (0.5000,-0.8660) -- (0.4911,-0.8512) -- (0.4820,-0.8368) -- (0.4726,-0.8230) -- (0.4629,-0.8097) -- (0.4531,-0.7968) -- (0.4430,-0.7845) -- (0.4328,-0.7726) -- (0.4224,-0.7611) -- (0.4119,-0.7501) -- (0.4012,-0.7396) -- (0.3905,-0.7294) -- (0.3796,-0.7197) -- (0.3686,-0.7103) -- (0.3575,-0.7014) -- (0.3464,-0.6928) -- (0.3352,-0.6846) -- (0.3239,-0.6768) -- (0.3126,-0.6693) -- (0.3013,-0.6622) -- (0.2899,-0.6554) -- (0.2784,-0.6489) -- (0.2670,-0.6428) -- (0.2555,-0.6369) -- (0.2439,-0.6314) -- (0.2324,-0.6262) -- (0.2208,-0.6212) -- (0.2092,-0.6166) -- (0.1977,-0.6122) -- (0.1861,-0.6082) -- (0.1745,-0.6043) -- (0.1628,-0.6008) -- (0.1512,-0.5975) -- (0.1396,-0.5945) -- (0.1280,-0.5917) -- (0.1163,-0.5892) -- (0.1047,-0.5869) -- (0.0931,-0.5849) -- (0.0814,-0.5831) -- (0.0698,-0.5816) -- (0.0582,-0.5803) -- (0.0465,-0.5792) -- (0.0349,-0.5784) -- (0.0233,-0.5778) -- (0.0116,-0.5775) -- (0.0000,-0.5774) -- (-0.0116,-0.5775) -- (-0.0233,-0.5778) -- (-0.0349,-0.5784) -- (-0.0465,-0.5792) -- (-0.0582,-0.5803) -- (-0.0698,-0.5816) -- (-0.0814,-0.5831) -- (-0.0931,-0.5849) -- (-0.1047,-0.5869) -- (-0.1163,-0.5892) -- (-0.1280,-0.5917) -- (-0.1396,-0.5945) -- (-0.1512,-0.5975) -- (-0.1628,-0.6008) -- (-0.1745,-0.6043) -- (-0.1861,-0.6082) -- (-0.1977,-0.6122) -- (-0.2092,-0.6166) -- (-0.2208,-0.6212) -- (-0.2324,-0.6262) -- (-0.2439,-0.6314) -- (-0.2555,-0.6369) -- (-0.2670,-0.6428) -- (-0.2784,-0.6489) -- (-0.2899,-0.6554) -- (-0.3013,-0.6622) -- (-0.3126,-0.6693) -- (-0.3239,-0.6768) -- (-0.3352,-0.6846) -- (-0.3464,-0.6928) -- (-0.3575,-0.7014) -- (-0.3686,-0.7103) -- (-0.3796,-0.7197) -- (-0.3905,-0.7294) -- (-0.4012,-0.7396) -- (-0.4119,-0.7501) -- (-0.4224,-0.7611) -- (-0.4328,-0.7726) -- (-0.4430,-0.7845) -- (-0.4531,-0.7968) -- (-0.4629,-0.8097) -- (-0.4726,-0.8230) -- (-0.4820,-0.8368) -- (-0.4911,-0.8512) -- (-0.5000,-0.8660);
\draw[wulffA] (0.6428,0.7660) -- (0.6283,0.7492) -- (0.6137,0.7332) -- (0.5990,0.7178) -- (0.5843,0.7031) -- (0.5695,0.6891) -- (0.5546,0.6757) -- (0.5398,0.6630) -- (0.5249,0.6507) -- (0.5100,0.6391) -- (0.4952,0.6280) -- (0.4803,0.6174) -- (0.4655,0.6073) -- (0.4507,0.5976) -- (0.4360,0.5885) -- (0.4213,0.5797) -- (0.4066,0.5714) -- (0.3920,0.5635) -- (0.3774,0.5560) -- (0.3629,0.5488) -- (0.3484,0.5421) -- (0.3340,0.5357) -- (0.3197,0.5296) -- (0.3053,0.5238) -- (0.2911,0.5184) -- (0.2769,0.5133) -- (0.2627,0.5085) -- (0.2486,0.5040) -- (0.2345,0.4997) -- (0.2205,0.4958) -- (0.2065,0.4921) -- (0.1925,0.4887) -- (0.1786,0.4855) -- (0.1647,0.4826) -- (0.1509,0.4800) -- (0.1371,0.4776) -- (0.1233,0.4754) -- (0.1095,0.4735) -- (0.0958,0.4718) -- (0.0821,0.4703) -- (0.0684,0.4691) -- (0.0547,0.4681) -- (0.0410,0.4673) -- (0.0273,0.4668) -- (0.0137,0.4664) -- (0.0000,0.4663) -- (-0.0137,0.4664) -- (-0.0273,0.4668) -- (-0.0410,0.4673) -- (-0.0547,0.4681) -- (-0.0684,0.4691) -- (-0.0821,0.4703) -- (-0.0958,0.4718) -- (-0.1095,0.4735) -- (-0.1233,0.4754) -- (-0.1371,0.4776) -- (-0.1509,0.4800) -- (-0.1647,0.4826) -- (-0.1786,0.4855) -- (-0.1925,0.4887) -- (-0.2065,0.4921) -- (-0.2205,0.4958) -- (-0.2345,0.4997) -- (-0.2486,0.5040) -- (-0.2627,0.5085) -- (-0.2769,0.5133) -- (-0.2911,0.5184) -- (-0.3053,0.5238) -- (-0.3197,0.5296) -- (-0.3340,0.5357) -- (-0.3484,0.5421) -- (-0.3629,0.5488) -- (-0.3774,0.5560) -- (-0.3920,0.5635) -- (-0.4066,0.5714) -- (-0.4213,0.5797) -- (-0.4360,0.5885) -- (-0.4507,0.5976) -- (-0.4655,0.6073) -- (-0.4803,0.6174) -- (-0.4952,0.6280) -- (-0.5100,0.6391) -- (-0.5249,0.6507) -- (-0.5398,0.6630) -- (-0.5546,0.6757) -- (-0.5695,0.6891) -- (-0.5843,0.7031) -- (-0.5990,0.7178) -- (-0.6137,0.7332) -- (-0.6283,0.7492) -- (-0.6428,0.7660);
\draw[wulffA] (0.6428,-0.7660) -- (0.6283,-0.7492) -- (0.6137,-0.7332) -- (0.5990,-0.7178) -- (0.5843,-0.7031) -- (0.5695,-0.6891) -- (0.5546,-0.6757) -- (0.5398,-0.6630) -- (0.5249,-0.6507) -- (0.5100,-0.6391) -- (0.4952,-0.6280) -- (0.4803,-0.6174) -- (0.4655,-0.6073) -- (0.4507,-0.5976) -- (0.4360,-0.5885) -- (0.4213,-0.5797) -- (0.4066,-0.5714) -- (0.3920,-0.5635) -- (0.3774,-0.5560) -- (0.3629,-0.5488) -- (0.3484,-0.5421) -- (0.3340,-0.5357) -- (0.3197,-0.5296) -- (0.3053,-0.5238) -- (0.2911,-0.5184) -- (0.2769,-0.5133) -- (0.2627,-0.5085) -- (0.2486,-0.5040) -- (0.2345,-0.4997) -- (0.2205,-0.4958) -- (0.2065,-0.4921) -- (0.1925,-0.4887) -- (0.1786,-0.4855) -- (0.1647,-0.4826) -- (0.1509,-0.4800) -- (0.1371,-0.4776) -- (0.1233,-0.4754) -- (0.1095,-0.4735) -- (0.0958,-0.4718) -- (0.0821,-0.4703) -- (0.0684,-0.4691) -- (0.0547,-0.4681) -- (0.0410,-0.4673) -- (0.0273,-0.4668) -- (0.0137,-0.4664) -- (0.0000,-0.4663) -- (-0.0137,-0.4664) -- (-0.0273,-0.4668) -- (-0.0410,-0.4673) -- (-0.0547,-0.4681) -- (-0.0684,-0.4691) -- (-0.0821,-0.4703) -- (-0.0958,-0.4718) -- (-0.1095,-0.4735) -- (-0.1233,-0.4754) -- (-0.1371,-0.4776) -- (-0.1509,-0.4800) -- (-0.1647,-0.4826) -- (-0.1786,-0.4855) -- (-0.1925,-0.4887) -- (-0.2065,-0.4921) -- (-0.2205,-0.4958) -- (-0.2345,-0.4997) -- (-0.2486,-0.5040) -- (-0.2627,-0.5085) -- (-0.2769,-0.5133) -- (-0.2911,-0.5184) -- (-0.3053,-0.5238) -- (-0.3197,-0.5296) -- (-0.3340,-0.5357) -- (-0.3484,-0.5421) -- (-0.3629,-0.5488) -- (-0.3774,-0.5560) -- (-0.3920,-0.5635) -- (-0.4066,-0.5714) -- (-0.4213,-0.5797) -- (-0.4360,-0.5885) -- (-0.4507,-0.5976) -- (-0.4655,-0.6073) -- (-0.4803,-0.6174) -- (-0.4952,-0.6280) -- (-0.5100,-0.6391) -- (-0.5249,-0.6507) -- (-0.5398,-0.6630) -- (-0.5546,-0.6757) -- (-0.5695,-0.6891) -- (-0.5843,-0.7031) -- (-0.5990,-0.7178) -- (-0.6137,-0.7332) -- (-0.6283,-0.7492) -- (-0.6428,-0.7660);
\draw[wulffA] (0.7660,0.6428) -- (0.7456,0.6261) -- (0.7254,0.6102) -- (0.7054,0.5951) -- (0.6855,0.5809) -- (0.6658,0.5673) -- (0.6464,0.5545) -- (0.6271,0.5423) -- (0.6080,0.5307) -- (0.5891,0.5198) -- (0.5704,0.5093) -- (0.5519,0.4995) -- (0.5336,0.4901) -- (0.5154,0.4812) -- (0.4975,0.4728) -- (0.4797,0.4648) -- (0.4621,0.4572) -- (0.4446,0.4500) -- (0.4273,0.4432) -- (0.4102,0.4368) -- (0.3932,0.4307) -- (0.3764,0.4250) -- (0.3597,0.4195) -- (0.3431,0.4144) -- (0.3266,0.4096) -- (0.3103,0.4051) -- (0.2941,0.4008) -- (0.2780,0.3968) -- (0.2620,0.3931) -- (0.2461,0.3897) -- (0.2303,0.3864) -- (0.2145,0.3834) -- (0.1989,0.3807) -- (0.1833,0.3782) -- (0.1678,0.3758) -- (0.1523,0.3737) -- (0.1369,0.3719) -- (0.1216,0.3702) -- (0.1063,0.3687) -- (0.0910,0.3675) -- (0.0758,0.3664) -- (0.0606,0.3655) -- (0.0454,0.3648) -- (0.0303,0.3644) -- (0.0151,0.3641) -- (0.0000,0.3640) -- (-0.0151,0.3641) -- (-0.0303,0.3644) -- (-0.0454,0.3648) -- (-0.0606,0.3655) -- (-0.0758,0.3664) -- (-0.0910,0.3675) -- (-0.1063,0.3687) -- (-0.1216,0.3702) -- (-0.1369,0.3719) -- (-0.1523,0.3737) -- (-0.1678,0.3758) -- (-0.1833,0.3782) -- (-0.1989,0.3807) -- (-0.2145,0.3834) -- (-0.2303,0.3864) -- (-0.2461,0.3897) -- (-0.2620,0.3931) -- (-0.2780,0.3968) -- (-0.2941,0.4008) -- (-0.3103,0.4051) -- (-0.3266,0.4096) -- (-0.3431,0.4144) -- (-0.3597,0.4195) -- (-0.3764,0.4250) -- (-0.3932,0.4307) -- (-0.4102,0.4368) -- (-0.4273,0.4432) -- (-0.4446,0.4500) -- (-0.4621,0.4572) -- (-0.4797,0.4648) -- (-0.4975,0.4728) -- (-0.5154,0.4812) -- (-0.5336,0.4901) -- (-0.5519,0.4995) -- (-0.5704,0.5093) -- (-0.5891,0.5198) -- (-0.6080,0.5307) -- (-0.6271,0.5423) -- (-0.6464,0.5545) -- (-0.6658,0.5673) -- (-0.6855,0.5809) -- (-0.7054,0.5951) -- (-0.7254,0.6102) -- (-0.7456,0.6261) -- (-0.7660,0.6428);
\draw[wulffA] (0.7660,-0.6428) -- (0.7456,-0.6261) -- (0.7254,-0.6102) -- (0.7054,-0.5951) -- (0.6855,-0.5809) -- (0.6658,-0.5673) -- (0.6464,-0.5545) -- (0.6271,-0.5423) -- (0.6080,-0.5307) -- (0.5891,-0.5198) -- (0.5704,-0.5093) -- (0.5519,-0.4995) -- (0.5336,-0.4901) -- (0.5154,-0.4812) -- (0.4975,-0.4728) -- (0.4797,-0.4648) -- (0.4621,-0.4572) -- (0.4446,-0.4500) -- (0.4273,-0.4432) -- (0.4102,-0.4368) -- (0.3932,-0.4307) -- (0.3764,-0.4250) -- (0.3597,-0.4195) -- (0.3431,-0.4144) -- (0.3266,-0.4096) -- (0.3103,-0.4051) -- (0.2941,-0.4008) -- (0.2780,-0.3968) -- (0.2620,-0.3931) -- (0.2461,-0.3897) -- (0.2303,-0.3864) -- (0.2145,-0.3834) -- (0.1989,-0.3807) -- (0.1833,-0.3782) -- (0.1678,-0.3758) -- (0.1523,-0.3737) -- (0.1369,-0.3719) -- (0.1216,-0.3702) -- (0.1063,-0.3687) -- (0.0910,-0.3675) -- (0.0758,-0.3664) -- (0.0606,-0.3655) -- (0.0454,-0.3648) -- (0.0303,-0.3644) -- (0.0151,-0.3641) -- (0.0000,-0.3640) -- (-0.0151,-0.3641) -- (-0.0303,-0.3644) -- (-0.0454,-0.3648) -- (-0.0606,-0.3655) -- (-0.0758,-0.3664) -- (-0.0910,-0.3675) -- (-0.1063,-0.3687) -- (-0.1216,-0.3702) -- (-0.1369,-0.3719) -- (-0.1523,-0.3737) -- (-0.1678,-0.3758) -- (-0.1833,-0.3782) -- (-0.1989,-0.3807) -- (-0.2145,-0.3834) -- (-0.2303,-0.3864) -- (-0.2461,-0.3897) -- (-0.2620,-0.3931) -- (-0.2780,-0.3968) -- (-0.2941,-0.4008) -- (-0.3103,-0.4051) -- (-0.3266,-0.4096) -- (-0.3431,-0.4144) -- (-0.3597,-0.4195) -- (-0.3764,-0.4250) -- (-0.3932,-0.4307) -- (-0.4102,-0.4368) -- (-0.4273,-0.4432) -- (-0.4446,-0.4500) -- (-0.4621,-0.4572) -- (-0.4797,-0.4648) -- (-0.4975,-0.4728) -- (-0.5154,-0.4812) -- (-0.5336,-0.4901) -- (-0.5519,-0.4995) -- (-0.5704,-0.5093) -- (-0.5891,-0.5198) -- (-0.6080,-0.5307) -- (-0.6271,-0.5423) -- (-0.6464,-0.5545) -- (-0.6658,-0.5673) -- (-0.6855,-0.5809) -- (-0.7054,-0.5951) -- (-0.7254,-0.6102) -- (-0.7456,-0.6261) -- (-0.7660,-0.6428);
\draw[wulffB] (0.8660,0.5000) -- (0.8401,0.4853) -- (0.8147,0.4715) -- (0.7898,0.4585) -- (0.7654,0.4462) -- (0.7414,0.4346) -- (0.7178,0.4237) -- (0.6947,0.4134) -- (0.6721,0.4036) -- (0.6498,0.3944) -- (0.6278,0.3857) -- (0.6063,0.3775) -- (0.5851,0.3698) -- (0.5642,0.3624) -- (0.5436,0.3555) -- (0.5234,0.3489) -- (0.5034,0.3427) -- (0.4837,0.3369) -- (0.4643,0.3313) -- (0.4451,0.3261) -- (0.4262,0.3212) -- (0.4075,0.3166) -- (0.3890,0.3122) -- (0.3707,0.3081) -- (0.3526,0.3042) -- (0.3347,0.3006) -- (0.3169,0.2972) -- (0.2993,0.2940) -- (0.2819,0.2910) -- (0.2646,0.2883) -- (0.2474,0.2857) -- (0.2304,0.2833) -- (0.2135,0.2812) -- (0.1967,0.2791) -- (0.1799,0.2773) -- (0.1633,0.2757) -- (0.1467,0.2742) -- (0.1303,0.2729) -- (0.1138,0.2717) -- (0.0975,0.2707) -- (0.0812,0.2699) -- (0.0649,0.2692) -- (0.0486,0.2686) -- (0.0324,0.2683) -- (0.0162,0.2680) -- (0.0000,0.2679) -- (-0.0162,0.2680) -- (-0.0324,0.2683) -- (-0.0486,0.2686) -- (-0.0649,0.2692) -- (-0.0812,0.2699) -- (-0.0975,0.2707) -- (-0.1138,0.2717) -- (-0.1303,0.2729) -- (-0.1467,0.2742) -- (-0.1633,0.2757) -- (-0.1799,0.2773) -- (-0.1967,0.2791) -- (-0.2135,0.2812) -- (-0.2304,0.2833) -- (-0.2474,0.2857) -- (-0.2646,0.2883) -- (-0.2819,0.2910) -- (-0.2993,0.2940) -- (-0.3169,0.2972) -- (-0.3347,0.3006) -- (-0.3526,0.3042) -- (-0.3707,0.3081) -- (-0.3890,0.3122) -- (-0.4075,0.3166) -- (-0.4262,0.3212) -- (-0.4451,0.3261) -- (-0.4643,0.3313) -- (-0.4837,0.3369) -- (-0.5034,0.3427) -- (-0.5234,0.3489) -- (-0.5436,0.3555) -- (-0.5642,0.3624) -- (-0.5851,0.3698) -- (-0.6063,0.3775) -- (-0.6278,0.3857) -- (-0.6498,0.3944) -- (-0.6721,0.4036) -- (-0.6947,0.4134) -- (-0.7178,0.4237) -- (-0.7414,0.4346) -- (-0.7654,0.4462) -- (-0.7898,0.4585) -- (-0.8147,0.4715) -- (-0.8401,0.4853) -- (-0.8660,0.5000);
\draw[wulffB] (0.8660,-0.5000) -- (0.8401,-0.4853) -- (0.8147,-0.4715) -- (0.7898,-0.4585) -- (0.7654,-0.4462) -- (0.7414,-0.4346) -- (0.7178,-0.4237) -- (0.6947,-0.4134) -- (0.6721,-0.4036) -- (0.6498,-0.3944) -- (0.6278,-0.3857) -- (0.6063,-0.3775) -- (0.5851,-0.3698) -- (0.5642,-0.3624) -- (0.5436,-0.3555) -- (0.5234,-0.3489) -- (0.5034,-0.3427) -- (0.4837,-0.3369) -- (0.4643,-0.3313) -- (0.4451,-0.3261) -- (0.4262,-0.3212) -- (0.4075,-0.3166) -- (0.3890,-0.3122) -- (0.3707,-0.3081) -- (0.3526,-0.3042) -- (0.3347,-0.3006) -- (0.3169,-0.2972) -- (0.2993,-0.2940) -- (0.2819,-0.2910) -- (0.2646,-0.2883) -- (0.2474,-0.2857) -- (0.2304,-0.2833) -- (0.2135,-0.2812) -- (0.1967,-0.2791) -- (0.1799,-0.2773) -- (0.1633,-0.2757) -- (0.1467,-0.2742) -- (0.1303,-0.2729) -- (0.1138,-0.2717) -- (0.0975,-0.2707) -- (0.0812,-0.2699) -- (0.0649,-0.2692) -- (0.0486,-0.2686) -- (0.0324,-0.2683) -- (0.0162,-0.2680) -- (0.0000,-0.2679) -- (-0.0162,-0.2680) -- (-0.0324,-0.2683) -- (-0.0486,-0.2686) -- (-0.0649,-0.2692) -- (-0.0812,-0.2699) -- (-0.0975,-0.2707) -- (-0.1138,-0.2717) -- (-0.1303,-0.2729) -- (-0.1467,-0.2742) -- (-0.1633,-0.2757) -- (-0.1799,-0.2773) -- (-0.1967,-0.2791) -- (-0.2135,-0.2812) -- (-0.2304,-0.2833) -- (-0.2474,-0.2857) -- (-0.2646,-0.2883) -- (-0.2819,-0.2910) -- (-0.2993,-0.2940) -- (-0.3169,-0.2972) -- (-0.3347,-0.3006) -- (-0.3526,-0.3042) -- (-0.3707,-0.3081) -- (-0.3890,-0.3122) -- (-0.4075,-0.3166) -- (-0.4262,-0.3212) -- (-0.4451,-0.3261) -- (-0.4643,-0.3313) -- (-0.4837,-0.3369) -- (-0.5034,-0.3427) -- (-0.5234,-0.3489) -- (-0.5436,-0.3555) -- (-0.5642,-0.3624) -- (-0.5851,-0.3698) -- (-0.6063,-0.3775) -- (-0.6278,-0.3857) -- (-0.6498,-0.3944) -- (-0.6721,-0.4036) -- (-0.6947,-0.4134) -- (-0.7178,-0.4237) -- (-0.7414,-0.4346) -- (-0.7654,-0.4462) -- (-0.7898,-0.4585) -- (-0.8147,-0.4715) -- (-0.8401,-0.4853) -- (-0.8660,-0.5000);
\draw[wulffA] (0.9397,0.3420) -- (0.9093,0.3312) -- (0.8797,0.3210) -- (0.8510,0.3114) -- (0.8229,0.3025) -- (0.7956,0.2940) -- (0.7689,0.2861) -- (0.7429,0.2787) -- (0.7175,0.2717) -- (0.6926,0.2651) -- (0.6683,0.2588) -- (0.6444,0.2530) -- (0.6211,0.2474) -- (0.5982,0.2422) -- (0.5757,0.2373) -- (0.5537,0.2327) -- (0.5320,0.2283) -- (0.5107,0.2242) -- (0.4897,0.2203) -- (0.4691,0.2167) -- (0.4488,0.2132) -- (0.4287,0.2100) -- (0.4090,0.2069) -- (0.3895,0.2041) -- (0.3702,0.2014) -- (0.3512,0.1989) -- (0.3324,0.1965) -- (0.3138,0.1943) -- (0.2954,0.1922) -- (0.2771,0.1903) -- (0.2590,0.1886) -- (0.2411,0.1869) -- (0.2233,0.1854) -- (0.2057,0.1840) -- (0.1881,0.1828) -- (0.1707,0.1816) -- (0.1533,0.1806) -- (0.1361,0.1797) -- (0.1189,0.1789) -- (0.1018,0.1782) -- (0.0847,0.1776) -- (0.0677,0.1772) -- (0.0508,0.1768) -- (0.0338,0.1765) -- (0.0169,0.1764) -- (0.0000,0.1763) -- (-0.0169,0.1764) -- (-0.0338,0.1765) -- (-0.0508,0.1768) -- (-0.0677,0.1772) -- (-0.0847,0.1776) -- (-0.1018,0.1782) -- (-0.1189,0.1789) -- (-0.1361,0.1797) -- (-0.1533,0.1806) -- (-0.1707,0.1816) -- (-0.1881,0.1828) -- (-0.2057,0.1840) -- (-0.2233,0.1854) -- (-0.2411,0.1869) -- (-0.2590,0.1886) -- (-0.2771,0.1903) -- (-0.2954,0.1922) -- (-0.3138,0.1943) -- (-0.3324,0.1965) -- (-0.3512,0.1989) -- (-0.3702,0.2014) -- (-0.3895,0.2041) -- (-0.4090,0.2069) -- (-0.4287,0.2100) -- (-0.4488,0.2132) -- (-0.4691,0.2167) -- (-0.4897,0.2203) -- (-0.5107,0.2242) -- (-0.5320,0.2283) -- (-0.5537,0.2327) -- (-0.5757,0.2373) -- (-0.5982,0.2422) -- (-0.6211,0.2474) -- (-0.6444,0.2530) -- (-0.6683,0.2588) -- (-0.6926,0.2651) -- (-0.7175,0.2717) -- (-0.7429,0.2787) -- (-0.7689,0.2861) -- (-0.7956,0.2940) -- (-0.8229,0.3025) -- (-0.8510,0.3114) -- (-0.8797,0.3210) -- (-0.9093,0.3312) -- (-0.9397,0.3420);
\draw[wulffA] (0.9397,-0.3420) -- (0.9093,-0.3312) -- (0.8797,-0.3210) -- (0.8510,-0.3114) -- (0.8229,-0.3025) -- (0.7956,-0.2940) -- (0.7689,-0.2861) -- (0.7429,-0.2787) -- (0.7175,-0.2717) -- (0.6926,-0.2651) -- (0.6683,-0.2588) -- (0.6444,-0.2530) -- (0.6211,-0.2474) -- (0.5982,-0.2422) -- (0.5757,-0.2373) -- (0.5537,-0.2327) -- (0.5320,-0.2283) -- (0.5107,-0.2242) -- (0.4897,-0.2203) -- (0.4691,-0.2167) -- (0.4488,-0.2132) -- (0.4287,-0.2100) -- (0.4090,-0.2069) -- (0.3895,-0.2041) -- (0.3702,-0.2014) -- (0.3512,-0.1989) -- (0.3324,-0.1965) -- (0.3138,-0.1943) -- (0.2954,-0.1922) -- (0.2771,-0.1903) -- (0.2590,-0.1886) -- (0.2411,-0.1869) -- (0.2233,-0.1854) -- (0.2057,-0.1840) -- (0.1881,-0.1828) -- (0.1707,-0.1816) -- (0.1533,-0.1806) -- (0.1361,-0.1797) -- (0.1189,-0.1789) -- (0.1018,-0.1782) -- (0.0847,-0.1776) -- (0.0677,-0.1772) -- (0.0508,-0.1768) -- (0.0338,-0.1765) -- (0.0169,-0.1764) -- (0.0000,-0.1763) -- (-0.0169,-0.1764) -- (-0.0338,-0.1765) -- (-0.0508,-0.1768) -- (-0.0677,-0.1772) -- (-0.0847,-0.1776) -- (-0.1018,-0.1782) -- (-0.1189,-0.1789) -- (-0.1361,-0.1797) -- (-0.1533,-0.1806) -- (-0.1707,-0.1816) -- (-0.1881,-0.1828) -- (-0.2057,-0.1840) -- (-0.2233,-0.1854) -- (-0.2411,-0.1869) -- (-0.2590,-0.1886) -- (-0.2771,-0.1903) -- (-0.2954,-0.1922) -- (-0.3138,-0.1943) -- (-0.3324,-0.1965) -- (-0.3512,-0.1989) -- (-0.3702,-0.2014) -- (-0.3895,-0.2041) -- (-0.4090,-0.2069) -- (-0.4287,-0.2100) -- (-0.4488,-0.2132) -- (-0.4691,-0.2167) -- (-0.4897,-0.2203) -- (-0.5107,-0.2242) -- (-0.5320,-0.2283) -- (-0.5537,-0.2327) -- (-0.5757,-0.2373) -- (-0.5982,-0.2422) -- (-0.6211,-0.2474) -- (-0.6444,-0.2530) -- (-0.6683,-0.2588) -- (-0.6926,-0.2651) -- (-0.7175,-0.2717) -- (-0.7429,-0.2787) -- (-0.7689,-0.2861) -- (-0.7956,-0.2940) -- (-0.8229,-0.3025) -- (-0.8510,-0.3114) -- (-0.8797,-0.3210) -- (-0.9093,-0.3312) -- (-0.9397,-0.3420);
\draw[wulffA] (0.9848,0.1736) -- (0.9515,0.1679) -- (0.9193,0.1625) -- (0.8880,0.1574) -- (0.8577,0.1527) -- (0.8282,0.1483) -- (0.7996,0.1441) -- (0.7717,0.1402) -- (0.7446,0.1366) -- (0.7181,0.1331) -- (0.6923,0.1299) -- (0.6670,0.1269) -- (0.6424,0.1240) -- (0.6182,0.1213) -- (0.5946,0.1187) -- (0.5715,0.1164) -- (0.5488,0.1141) -- (0.5265,0.1120) -- (0.5046,0.1100) -- (0.4831,0.1081) -- (0.4620,0.1063) -- (0.4412,0.1047) -- (0.4206,0.1031) -- (0.4004,0.1016) -- (0.3805,0.1003) -- (0.3608,0.0990) -- (0.3414,0.0978) -- (0.3222,0.0966) -- (0.3032,0.0956) -- (0.2844,0.0946) -- (0.2658,0.0937) -- (0.2473,0.0929) -- (0.2290,0.0921) -- (0.2109,0.0914) -- (0.1928,0.0908) -- (0.1749,0.0902) -- (0.1571,0.0897) -- (0.1394,0.0892) -- (0.1218,0.0888) -- (0.1043,0.0884) -- (0.0868,0.0882) -- (0.0694,0.0879) -- (0.0520,0.0877) -- (0.0347,0.0876) -- (0.0173,0.0875) -- (0.0000,0.0875) -- (-0.0173,0.0875) -- (-0.0347,0.0876) -- (-0.0520,0.0877) -- (-0.0694,0.0879) -- (-0.0868,0.0882) -- (-0.1043,0.0884) -- (-0.1218,0.0888) -- (-0.1394,0.0892) -- (-0.1571,0.0897) -- (-0.1749,0.0902) -- (-0.1928,0.0908) -- (-0.2109,0.0914) -- (-0.2290,0.0921) -- (-0.2473,0.0929) -- (-0.2658,0.0937) -- (-0.2844,0.0946) -- (-0.3032,0.0956) -- (-0.3222,0.0966) -- (-0.3414,0.0978) -- (-0.3608,0.0990) -- (-0.3805,0.1003) -- (-0.4004,0.1016) -- (-0.4206,0.1031) -- (-0.4412,0.1047) -- (-0.4620,0.1063) -- (-0.4831,0.1081) -- (-0.5046,0.1100) -- (-0.5265,0.1120) -- (-0.5488,0.1141) -- (-0.5715,0.1164) -- (-0.5946,0.1187) -- (-0.6182,0.1213) -- (-0.6424,0.1240) -- (-0.6670,0.1269) -- (-0.6923,0.1299) -- (-0.7181,0.1331) -- (-0.7446,0.1366) -- (-0.7717,0.1402) -- (-0.7996,0.1441) -- (-0.8282,0.1483) -- (-0.8577,0.1527) -- (-0.8880,0.1574) -- (-0.9193,0.1625) -- (-0.9515,0.1679) -- (-0.9848,0.1736);
\draw[wulffA] (0.9848,-0.1736) -- (0.9515,-0.1679) -- (0.9193,-0.1625) -- (0.8880,-0.1574) -- (0.8577,-0.1527) -- (0.8282,-0.1483) -- (0.7996,-0.1441) -- (0.7717,-0.1402) -- (0.7446,-0.1366) -- (0.7181,-0.1331) -- (0.6923,-0.1299) -- (0.6670,-0.1269) -- (0.6424,-0.1240) -- (0.6182,-0.1213) -- (0.5946,-0.1187) -- (0.5715,-0.1164) -- (0.5488,-0.1141) -- (0.5265,-0.1120) -- (0.5046,-0.1100) -- (0.4831,-0.1081) -- (0.4620,-0.1063) -- (0.4412,-0.1047) -- (0.4206,-0.1031) -- (0.4004,-0.1016) -- (0.3805,-0.1003) -- (0.3608,-0.0990) -- (0.3414,-0.0978) -- (0.3222,-0.0966) -- (0.3032,-0.0956) -- (0.2844,-0.0946) -- (0.2658,-0.0937) -- (0.2473,-0.0929) -- (0.2290,-0.0921) -- (0.2109,-0.0914) -- (0.1928,-0.0908) -- (0.1749,-0.0902) -- (0.1571,-0.0897) -- (0.1394,-0.0892) -- (0.1218,-0.0888) -- (0.1043,-0.0884) -- (0.0868,-0.0882) -- (0.0694,-0.0879) -- (0.0520,-0.0877) -- (0.0347,-0.0876) -- (0.0173,-0.0875) -- (0.0000,-0.0875) -- (-0.0173,-0.0875) -- (-0.0347,-0.0876) -- (-0.0520,-0.0877) -- (-0.0694,-0.0879) -- (-0.0868,-0.0882) -- (-0.1043,-0.0884) -- (-0.1218,-0.0888) -- (-0.1394,-0.0892) -- (-0.1571,-0.0897) -- (-0.1749,-0.0902) -- (-0.1928,-0.0908) -- (-0.2109,-0.0914) -- (-0.2290,-0.0921) -- (-0.2473,-0.0929) -- (-0.2658,-0.0937) -- (-0.2844,-0.0946) -- (-0.3032,-0.0956) -- (-0.3222,-0.0966) -- (-0.3414,-0.0978) -- (-0.3608,-0.0990) -- (-0.3805,-0.1003) -- (-0.4004,-0.1016) -- (-0.4206,-0.1031) -- (-0.4412,-0.1047) -- (-0.4620,-0.1063) -- (-0.4831,-0.1081) -- (-0.5046,-0.1100) -- (-0.5265,-0.1120) -- (-0.5488,-0.1141) -- (-0.5715,-0.1164) -- (-0.5946,-0.1187) -- (-0.6182,-0.1213) -- (-0.6424,-0.1240) -- (-0.6670,-0.1269) -- (-0.6923,-0.1299) -- (-0.7181,-0.1331) -- (-0.7446,-0.1366) -- (-0.7717,-0.1402) -- (-0.7996,-0.1441) -- (-0.8282,-0.1483) -- (-0.8577,-0.1527) -- (-0.8880,-0.1574) -- (-0.9193,-0.1625) -- (-0.9515,-0.1679) -- (-0.9848,-0.1736);
\draw[wulffB] (-1,0) -- (1,0) (0,-1) -- (0,1);
\draw[wulffA] (1.0000,0.0000) -- (0.9400,0.0000);
\draw[wulffA] (0.9994,0.0349) -- (0.9644,0.0337);
\draw[wulffA] (0.9976,0.0698) -- (0.9626,0.0673);
\draw[wulffA] (0.9945,0.1045) -- (0.9597,0.1009);
\draw[wulffA] (0.9903,0.1392) -- (0.9556,0.1343);
\draw[wulffA] (0.9848,0.1736) -- (0.9257,0.1632);
\draw[wulffA] (0.9781,0.2079) -- (0.9439,0.2006);
\draw[wulffA] (0.9703,0.2419) -- (0.9363,0.2335);
\draw[wulffA] (0.9613,0.2756) -- (0.9276,0.2660);
\draw[wulffA] (0.9511,0.3090) -- (0.9178,0.2982);
\draw[wulffA] (0.9397,0.3420) -- (0.8833,0.3215);
\draw[wulffA] (0.9272,0.3746) -- (0.8947,0.3615);
\draw[wulffA] (0.9135,0.4067) -- (0.8816,0.3925);
\draw[wulffA] (0.8988,0.4384) -- (0.8673,0.4230);
\draw[wulffA] (0.8829,0.4695) -- (0.8520,0.4530);
\draw[wulffA] (0.8660,0.5000) -- (0.8141,0.4700);
\draw[wulffA] (0.8480,0.5299) -- (0.8184,0.5114);
\draw[wulffA] (0.8290,0.5592) -- (0.8000,0.5396);
\draw[wulffA] (0.8090,0.5878) -- (0.7807,0.5672);
\draw[wulffA] (0.7880,0.6157) -- (0.7604,0.5941);
\draw[wulffA] (0.7660,0.6428) -- (0.7201,0.6042);
\draw[wulffA] (0.7431,0.6691) -- (0.7171,0.6457);
\draw[wulffA] (0.7193,0.6947) -- (0.6942,0.6703);
\draw[wulffA] (0.6947,0.7193) -- (0.6703,0.6942);
\draw[wulffA] (0.6691,0.7431) -- (0.6457,0.7171);
\draw[wulffA] (0.6428,0.7660) -- (0.6042,0.7201);
\draw[wulffA] (0.6157,0.7880) -- (0.5941,0.7604);
\draw[wulffA] (0.5878,0.8090) -- (0.5672,0.7807);
\draw[wulffA] (0.5592,0.8290) -- (0.5396,0.8000);
\draw[wulffA] (0.5299,0.8480) -- (0.5114,0.8184);
\draw[wulffA] (0.5000,0.8660) -- (0.4700,0.8141);
\draw[wulffA] (0.4695,0.8829) -- (0.4530,0.8520);
\draw[wulffA] (0.4384,0.8988) -- (0.4230,0.8673);
\draw[wulffA] (0.4067,0.9135) -- (0.3925,0.8816);
\draw[wulffA] (0.3746,0.9272) -- (0.3615,0.8947);
\draw[wulffA] (0.3420,0.9397) -- (0.3215,0.8833);
\draw[wulffA] (0.3090,0.9511) -- (0.2982,0.9178);
\draw[wulffA] (0.2756,0.9613) -- (0.2660,0.9276);
\draw[wulffA] (0.2419,0.9703) -- (0.2335,0.9363);
\draw[wulffA] (0.2079,0.9781) -- (0.2006,0.9439);
\draw[wulffA] (0.1736,0.9848) -- (0.1632,0.9257);
\draw[wulffA] (0.1392,0.9903) -- (0.1343,0.9556);
\draw[wulffA] (0.1045,0.9945) -- (0.1009,0.9597);
\draw[wulffA] (0.0698,0.9976) -- (0.0673,0.9626);
\draw[wulffA] (0.0349,0.9994) -- (0.0337,0.9644);
\draw[wulffA] (0.0000,1.0000) -- (0.0000,0.9400);
\draw[wulffA] (-0.0349,0.9994) -- (-0.0337,0.9644);
\draw[wulffA] (-0.0698,0.9976) -- (-0.0673,0.9626);
\draw[wulffA] (-0.1045,0.9945) -- (-0.1009,0.9597);
\draw[wulffA] (-0.1392,0.9903) -- (-0.1343,0.9556);
\draw[wulffA] (-0.1736,0.9848) -- (-0.1632,0.9257);
\draw[wulffA] (-0.2079,0.9781) -- (-0.2006,0.9439);
\draw[wulffA] (-0.2419,0.9703) -- (-0.2335,0.9363);
\draw[wulffA] (-0.2756,0.9613) -- (-0.2660,0.9276);
\draw[wulffA] (-0.3090,0.9511) -- (-0.2982,0.9178);
\draw[wulffA] (-0.3420,0.9397) -- (-0.3215,0.8833);
\draw[wulffA] (-0.3746,0.9272) -- (-0.3615,0.8947);
\draw[wulffA] (-0.4067,0.9135) -- (-0.3925,0.8816);
\draw[wulffA] (-0.4384,0.8988) -- (-0.4230,0.8673);
\draw[wulffA] (-0.4695,0.8829) -- (-0.4530,0.8520);
\draw[wulffA] (-0.5000,0.8660) -- (-0.4700,0.8141);
\draw[wulffA] (-0.5299,0.8480) -- (-0.5114,0.8184);
\draw[wulffA] (-0.5592,0.8290) -- (-0.5396,0.8000);
\draw[wulffA] (-0.5878,0.8090) -- (-0.5672,0.7807);
\draw[wulffA] (-0.6157,0.7880) -- (-0.5941,0.7604);
\draw[wulffA] (-0.6428,0.7660) -- (-0.6042,0.7201);
\draw[wulffA] (-0.6691,0.7431) -- (-0.6457,0.7171);
\draw[wulffA] (-0.6947,0.7193) -- (-0.6703,0.6942);
\draw[wulffA] (-0.7193,0.6947) -- (-0.6942,0.6703);
\draw[wulffA] (-0.7431,0.6691) -- (-0.7171,0.6457);
\draw[wulffA] (-0.7660,0.6428) -- (-0.7201,0.6042);
\draw[wulffA] (-0.7880,0.6157) -- (-0.7604,0.5941);
\draw[wulffA] (-0.8090,0.5878) -- (-0.7807,0.5672);
\draw[wulffA] (-0.8290,0.5592) -- (-0.8000,0.5396);
\draw[wulffA] (-0.8480,0.5299) -- (-0.8184,0.5114);
\draw[wulffA] (-0.8660,0.5000) -- (-0.8141,0.4700);
\draw[wulffA] (-0.8829,0.4695) -- (-0.8520,0.4530);
\draw[wulffA] (-0.8988,0.4384) -- (-0.8673,0.4230);
\draw[wulffA] (-0.9135,0.4067) -- (-0.8816,0.3925);
\draw[wulffA] (-0.9272,0.3746) -- (-0.8947,0.3615);
\draw[wulffA] (-0.9397,0.3420) -- (-0.8833,0.3215);
\draw[wulffA] (-0.9511,0.3090) -- (-0.9178,0.2982);
\draw[wulffA] (-0.9613,0.2756) -- (-0.9276,0.2660);
\draw[wulffA] (-0.9703,0.2419) -- (-0.9363,0.2335);
\draw[wulffA] (-0.9781,0.2079) -- (-0.9439,0.2006);
\draw[wulffA] (-0.9848,0.1736) -- (-0.9257,0.1632);
\draw[wulffA] (-0.9903,0.1392) -- (-0.9556,0.1343);
\draw[wulffA] (-0.9945,0.1045) -- (-0.9597,0.1009);
\draw[wulffA] (-0.9976,0.0698) -- (-0.9626,0.0673);
\draw[wulffA] (-0.9994,0.0349) -- (-0.9644,0.0337);
\draw[wulffA] (-1.0000,0.0000) -- (-0.9400,0.0000);
\draw[wulffA] (-0.9994,-0.0349) -- (-0.9644,-0.0337);
\draw[wulffA] (-0.9976,-0.0698) -- (-0.9626,-0.0673);
\draw[wulffA] (-0.9945,-0.1045) -- (-0.9597,-0.1009);
\draw[wulffA] (-0.9903,-0.1392) -- (-0.9556,-0.1343);
\draw[wulffA] (-0.9848,-0.1736) -- (-0.9257,-0.1632);
\draw[wulffA] (-0.9781,-0.2079) -- (-0.9439,-0.2006);
\draw[wulffA] (-0.9703,-0.2419) -- (-0.9363,-0.2335);
\draw[wulffA] (-0.9613,-0.2756) -- (-0.9276,-0.2660);
\draw[wulffA] (-0.9511,-0.3090) -- (-0.9178,-0.2982);
\draw[wulffA] (-0.9397,-0.3420) -- (-0.8833,-0.3215);
\draw[wulffA] (-0.9272,-0.3746) -- (-0.8947,-0.3615);
\draw[wulffA] (-0.9135,-0.4067) -- (-0.8816,-0.3925);
\draw[wulffA] (-0.8988,-0.4384) -- (-0.8673,-0.4230);
\draw[wulffA] (-0.8829,-0.4695) -- (-0.8520,-0.4530);
\draw[wulffA] (-0.8660,-0.5000) -- (-0.8141,-0.4700);
\draw[wulffA] (-0.8480,-0.5299) -- (-0.8184,-0.5114);
\draw[wulffA] (-0.8290,-0.5592) -- (-0.8000,-0.5396);
\draw[wulffA] (-0.8090,-0.5878) -- (-0.7807,-0.5672);
\draw[wulffA] (-0.7880,-0.6157) -- (-0.7604,-0.5941);
\draw[wulffA] (-0.7660,-0.6428) -- (-0.7201,-0.6042);
\draw[wulffA] (-0.7431,-0.6691) -- (-0.7171,-0.6457);
\draw[wulffA] (-0.7193,-0.6947) -- (-0.6942,-0.6703);
\draw[wulffA] (-0.6947,-0.7193) -- (-0.6703,-0.6942);
\draw[wulffA] (-0.6691,-0.7431) -- (-0.6457,-0.7171);
\draw[wulffA] (-0.6428,-0.7660) -- (-0.6042,-0.7201);
\draw[wulffA] (-0.6157,-0.7880) -- (-0.5941,-0.7604);
\draw[wulffA] (-0.5878,-0.8090) -- (-0.5672,-0.7807);
\draw[wulffA] (-0.5592,-0.8290) -- (-0.5396,-0.8000);
\draw[wulffA] (-0.5299,-0.8480) -- (-0.5114,-0.8184);
\draw[wulffA] (-0.5000,-0.8660) -- (-0.4700,-0.8141);
\draw[wulffA] (-0.4695,-0.8829) -- (-0.4530,-0.8520);
\draw[wulffA] (-0.4384,-0.8988) -- (-0.4230,-0.8673);
\draw[wulffA] (-0.4067,-0.9135) -- (-0.3925,-0.8816);
\draw[wulffA] (-0.3746,-0.9272) -- (-0.3615,-0.8947);
\draw[wulffA] (-0.3420,-0.9397) -- (-0.3215,-0.8833);
\draw[wulffA] (-0.3090,-0.9511) -- (-0.2982,-0.9178);
\draw[wulffA] (-0.2756,-0.9613) -- (-0.2660,-0.9276);
\draw[wulffA] (-0.2419,-0.9703) -- (-0.2335,-0.9363);
\draw[wulffA] (-0.2079,-0.9781) -- (-0.2006,-0.9439);
\draw[wulffA] (-0.1736,-0.9848) -- (-0.1632,-0.9257);
\draw[wulffA] (-0.1392,-0.9903) -- (-0.1343,-0.9556);
\draw[wulffA] (-0.1045,-0.9945) -- (-0.1009,-0.9597);
\draw[wulffA] (-0.0698,-0.9976) -- (-0.0673,-0.9626);
\draw[wulffA] (-0.0349,-0.9994) -- (-0.0337,-0.9644);
\draw[wulffA] (-0.0000,-1.0000) -- (-0.0000,-0.9400);
\draw[wulffA] (0.0349,-0.9994) -- (0.0337,-0.9644);
\draw[wulffA] (0.0698,-0.9976) -- (0.0673,-0.9626);
\draw[wulffA] (0.1045,-0.9945) -- (0.1009,-0.9597);
\draw[wulffA] (0.1392,-0.9903) -- (0.1343,-0.9556);
\draw[wulffA] (0.1736,-0.9848) -- (0.1632,-0.9257);
\draw[wulffA] (0.2079,-0.9781) -- (0.2006,-0.9439);
\draw[wulffA] (0.2419,-0.9703) -- (0.2335,-0.9363);
\draw[wulffA] (0.2756,-0.9613) -- (0.2660,-0.9276);
\draw[wulffA] (0.3090,-0.9511) -- (0.2982,-0.9178);
\draw[wulffA] (0.3420,-0.9397) -- (0.3215,-0.8833);
\draw[wulffA] (0.3746,-0.9272) -- (0.3615,-0.8947);
\draw[wulffA] (0.4067,-0.9135) -- (0.3925,-0.8816);
\draw[wulffA] (0.4384,-0.8988) -- (0.4230,-0.8673);
\draw[wulffA] (0.4695,-0.8829) -- (0.4530,-0.8520);
\draw[wulffA] (0.5000,-0.8660) -- (0.4700,-0.8141);
\draw[wulffA] (0.5299,-0.8480) -- (0.5114,-0.8184);
\draw[wulffA] (0.5592,-0.8290) -- (0.5396,-0.8000);
\draw[wulffA] (0.5878,-0.8090) -- (0.5672,-0.7807);
\draw[wulffA] (0.6157,-0.7880) -- (0.5941,-0.7604);
\draw[wulffA] (0.6428,-0.7660) -- (0.6042,-0.7201);
\draw[wulffA] (0.6691,-0.7431) -- (0.6457,-0.7171);
\draw[wulffA] (0.6947,-0.7193) -- (0.6703,-0.6942);
\draw[wulffA] (0.7193,-0.6947) -- (0.6942,-0.6703);
\draw[wulffA] (0.7431,-0.6691) -- (0.7171,-0.6457);
\draw[wulffA] (0.7660,-0.6428) -- (0.7201,-0.6042);
\draw[wulffA] (0.7880,-0.6157) -- (0.7604,-0.5941);
\draw[wulffA] (0.8090,-0.5878) -- (0.7807,-0.5672);
\draw[wulffA] (0.8290,-0.5592) -- (0.8000,-0.5396);
\draw[wulffA] (0.8480,-0.5299) -- (0.8184,-0.5114);
\draw[wulffA] (0.8660,-0.5000) -- (0.8141,-0.4700);
\draw[wulffA] (0.8829,-0.4695) -- (0.8520,-0.4530);
\draw[wulffA] (0.8988,-0.4384) -- (0.8673,-0.4230);
\draw[wulffA] (0.9135,-0.4067) -- (0.8816,-0.3925);
\draw[wulffA] (0.9272,-0.3746) -- (0.8947,-0.3615);
\draw[wulffA] (0.9397,-0.3420) -- (0.8833,-0.3215);
\draw[wulffA] (0.9511,-0.3090) -- (0.9178,-0.2982);
\draw[wulffA] (0.9613,-0.2756) -- (0.9276,-0.2660);
\draw[wulffA] (0.9703,-0.2419) -- (0.9363,-0.2335);
\draw[wulffA] (0.9781,-0.2079) -- (0.9439,-0.2006);
\draw[wulffA] (0.9848,-0.1736) -- (0.9257,-0.1632);
\draw[wulffA] (0.9903,-0.1392) -- (0.9556,-0.1343);
\draw[wulffA] (0.9945,-0.1045) -- (0.9597,-0.1009);
\draw[wulffA] (0.9976,-0.0698) -- (0.9626,-0.0673);
\draw[wulffA] (0.9994,-0.0349) -- (0.9644,-0.0337);
\draw[wulffC] (0,0) circle (1);

  \node[font=\scriptsize, anchor=south] at (0,1.02) {$y$};
\end{tikzpicture}
\caption{}
\end{subfigure}\hfill
\begin{subfigure}[b]{0.48\textwidth}\centering
\begin{tikzpicture}[scale=2.9]
  % создано tools/gen_stereo.py: сетка Болдырева (r = 1)
\draw[wulffA] (0,0) circle (0.0875);
\draw[wulffA] (0,0) circle (0.1763);
\draw[wulffB] (0,0) circle (0.2679);
\draw[wulffA] (0,0) circle (0.3640);
\draw[wulffA] (0,0) circle (0.4663);
\draw[wulffB] (0,0) circle (0.5774);
\draw[wulffA] (0,0) circle (0.7002);
\draw[wulffA] (0,0) circle (0.8391);
\draw[wulffB] (1.0000,0.0000) -- (-1.0000,-0.0000);
\draw[wulffA] (0.9848,0.1736) -- (-0.9848,-0.1736);
\draw[wulffA] (0.9397,0.3420) -- (-0.9397,-0.3420);
\draw[wulffB] (0.8660,0.5000) -- (-0.8660,-0.5000);
\draw[wulffA] (0.7660,0.6428) -- (-0.7660,-0.6428);
\draw[wulffA] (0.6428,0.7660) -- (-0.6428,-0.7660);
\draw[wulffB] (0.5000,0.8660) -- (-0.5000,-0.8660);
\draw[wulffA] (0.3420,0.9397) -- (-0.3420,-0.9397);
\draw[wulffA] (0.1736,0.9848) -- (-0.1736,-0.9848);
\draw[wulffB] (0.0000,1.0000) -- (-0.0000,-1.0000);
\draw[wulffA] (-0.1736,0.9848) -- (0.1736,-0.9848);
\draw[wulffA] (-0.3420,0.9397) -- (0.3420,-0.9397);
\draw[wulffB] (-0.5000,0.8660) -- (0.5000,-0.8660);
\draw[wulffA] (-0.6428,0.7660) -- (0.6428,-0.7660);
\draw[wulffA] (-0.7660,0.6428) -- (0.7660,-0.6428);
\draw[wulffB] (-0.8660,0.5000) -- (0.8660,-0.5000);
\draw[wulffA] (-0.9397,0.3420) -- (0.9397,-0.3420);
\draw[wulffA] (-0.9848,0.1736) -- (0.9848,-0.1736);
\draw[wulffC] (0,0) circle (1);

\end{tikzpicture}
\caption{}
\end{subfigure}
\caption{Координатные сетки (расчёт): \textit{а}~— сетка Вульфа (меридианы
и параллели через $10^\circ$, жирные~— через $30^\circ$; штрихи на основном
круге~— через $2^\circ$); \textit{б}~— сетка Болдырева (окружности равного
угла $\theta$ с осью $z$ и радиусы через $10^\circ$)}
\label{fig:nets}
\end{figure}

\begin{podrobnee}[Как построить меридиан и параллель]
Меридиан, изображающий плоскость, наклонённую под углом $\delta$ к
плоскости проекции (и проходящую через ось $y$),~— дуга окружности
радиуса $R_\text{м}=r/\cos\delta$ с центром на экваторе на расстоянии
$OO_\text{м}=r\operatorname{tg}\delta$ от центра. Параллель~— малый круг угловым
радиусом $\beta$ вокруг оси $y$~— дуга окружности радиуса
$R_\text{п}=r\operatorname{tg}\beta$ с центром на оси $y$ на расстоянии
$OO_\text{п}=r/\cos\beta$ (рис.~\ref{fig:merid}). Обе формулы следуют из
\eqref{eq:stereo-rho} и того, что окружность проходит через известные
точки основного круга\isprmark.
\end{podrobnee}
\isprtext{В оригинале (формулы (1.4)--(1.7)) меридиан характеризуется углом
$\alpha=90^\circ-\delta$, а параллель~— угловым расстоянием от экватора
$90^\circ-\beta$; в этих обозначениях $R_\text{м}=r/\sin\alpha$,
$OO_\text{м}=r\operatorname{ctg}\alpha$, $R_\text{п}=r\operatorname{ctg}\beta'$,
$OO_\text{п}=r/\sin\beta'$. Здесь использованы углы, отсчитываемые от
плоскости проекции и от оси сетки,~— так формулы проще запомнить.}

\begin{figure}[ht]
\centering
\begin{tikzpicture}[scale=1.0]
  \def\R{2.2}
  \draw[atomLine, thick] (0,0) circle (\R);
  \draw[fmGray] (-3.2,0) -- (2.5,0) (0,-2.5) -- (0,3.6);
  \fill (0,0) circle (1.2pt) node[below right, font=\small]{$O$};
  \begin{scope}
    \clip (0,0) circle (\R);
    \draw[fmBlue, very thick] (-2.622,0) circle (3.4227);
    \draw[fmRed, very thick] (0,2.872) circle (1.846);
  \end{scope}
  \fill[fmBlue] (-2.622,0) circle (1.4pt) node[below, font=\small]{$O_\text{м}$};
  \draw[fmBlue, dashed] (-2.622,0) -- (0,2.2) node[midway, above left=-2pt, font=\small]{$R_\text{м}$};
  \fill[fmRed] (0,2.872) circle (1.4pt) node[right, font=\small]{$O_\text{п}$};
  \draw[fmRed, dashed] (0,2.872) -- ({2.2*sin(40)},{2.2*cos(40)}) node[pos=0.45, right, font=\small]{$R_\text{п}$};
  \fill (0,2.2) circle (1.2pt) node[above left, font=\small]{$y$};
  \node[fmBlue, font=\small] at (0.75,-1.3) {меридиан};
  \node[fmRed, font=\small] at (0,0.75) {параллель};
\end{tikzpicture}
\caption{Построение меридиана ($\delta=50^\circ$, синий) и параллели
($\beta=40^\circ$, красный): центры $O_\text{м}$ и $O_\text{п}$ лежат вне основного
круга}
\label{fig:merid}
\end{figure}

На практике проекцию строят на кальке, наложенной на сетку того же
диаметра: кальку поворачивают вокруг общего центра (следя, чтобы центр не
смещался) до нужного положения и считывают углы по сетке.

\paragraph{Основные задачи, решаемые с помощью сеток}
\begin{enumerate}
\item \emph{Угол между двумя направлениями.} Два направления комплекса
всегда лежат в одной плоскости, а плоскость на проекции~— меридиан.
Поэтому кальку поворачивают, пока обе точки не окажутся на одном
меридиане, и отсчитывают угол между ними вдоль меридиана (по
пересекаемым параллелям). Если одно направление изображено точкой, а
другое крестиком, их выводят на меридианы, симметричные относительно
нулевого, и складывают углы, отсчитанные до полюса сетки. Можно и
заменить крестик противоположной ему точкой (отложив равный отрезок по
другую сторону от центра) и найти дополнительный угол $180^\circ-\alpha$.
\item \emph{Координаты проекции по углам.} Если направление образует с
осью $z$ угол $\theta$, а его проекция на плоскость $xy$~— угол $\varphi$
с осью $x$, то
\begin{equation}\label{eq:stereo-xy}
R_x=r\operatorname{tg}\frac\theta2\cos\varphi,\qquad R_y=r\operatorname{tg}\frac\theta2\sin\varphi .
\end{equation}
Через направляющие углы $\alpha$, $\beta$, $\gamma$ с осями $x$, $y$, $z$:
$R_x=r\operatorname{tg}(\gamma/2)\cos\alpha/\sin\gamma$,
$R_y=r\operatorname{tg}(\gamma/2)\cos\beta/\sin\gamma$.
\item \emph{Полюс плоскости.} Плоскость~— меридиан; поворачиваем кальку до
совпадения её дуги с меридианом сетки и отсчитываем по экватору
$90^\circ$: полученная точка~— проекция нормали, т.\,е. полюс плоскости.
\item \emph{Угол между плоскостями}~— это угол между их нормалями:
находим полюсы и измеряем угол между ними, как в задаче~1.
\item \emph{Поворот вокруг оси, лежащей в плоскости проекции.} Совмещаем
ось поворота с полюсом сетки; при повороте каждая точка движется по своей
\emph{параллели} (по конусу вокруг оси) на заданный угол. Точка, «перешедшая»
через основной круг, уходит в южное полушарие и изображается крестиком
(её можно перевести обратно в точку отражением через центр). Плоскость
поворачивают через её полюс: поворачивают полюс, затем строят новую дугу.
\item \emph{Поворот вокруг оси, перпендикулярной плоскости проекции}~—
по сетке Болдырева: все точки смещаются по концентрическим окружностям на
один и тот же угол.
\item \emph{Малый круг}~— геометрическое место направлений, образующих угол
$\alpha$ с заданной осью (проекция кругового конуса). Через проекцию оси
проводят меридиан, откладывают по нему в обе стороны $\alpha$ и через две
полученные точки проводят окружность. Её центр \emph{не совпадает} с
проекцией оси из-за неравномерности шкалы. Поворот вокруг \emph{произвольной}
оси переводит точку по такому малому кругу; угол поворота откладывают не
транспортиром, а в плоскости, перпендикулярной оси.
\end{enumerate}

\begin{primer}[Проекция направления {[111]} кубического кристалла]
Плоскость проекции~— $(001)$. Направление $[111]$ образует с $[001]$ угол
$\theta=\arccos(1/\sqrt3)=54{,}74^\circ$, а его проекция на плоскость $xy$
идёт под $\varphi=45^\circ$ к оси $x$. По~\eqref{eq:stereo-xy}
$\rho=r\operatorname{tg}27{,}37^\circ=0{,}518\,r$, $R_x=R_y=0{,}366\,r$. Угол между
$[111]$ и $[001]$, измеренный по сетке Вульфа (обе точки выведены на
нулевой меридиан), равен тем же $54{,}7^\circ$.
\end{primer}

\paragraph{Описание поворота: углы Эйлера и ось--угол}
Поворот твёрдого тела (кристалла) задают двумя способами.
\begin{itemize}
\item \emph{Углы Эйлера} $\varphi$, $\theta$, $\psi$: поворот на $\varphi$ вокруг
оси $z$ (оси $x,y$ переходят в $x',y'$), затем на $\theta$ вокруг нового
положения оси $x'$ (ось $z$ переходит в $z'$), затем на $\psi$ вокруг $z'$.
По проекции с осями до и после поворота углы находят так: $\theta$~— угол
между $z$ и $z'$; $\varphi$~— угол, на котором линия пересечения плоскостей
$xy$ и $x'y'$ (ось $x'$ промежуточной системы) отстоит от $x$; $\psi$~— угол
от этой линии до конечной оси $x''$.
\item \emph{Ось и угол поворота} (чаще используются в кристаллографии,
например для описания разориентации зёрен, п.~1.7). Если направления $1$
и $2$ переходят в $1'$ и $2'$, ось поворота лежит на пересечении двух
\emph{биссекторных} плоскостей: каждая проходит через биссектрису угла
между направлением до и после поворота и через нормаль к плоскости,
содержащей эти два направления. Угол поворота~— центральный угол в
плоскости, перпендикулярной оси.
\end{itemize}
Матричная запись поворотов~— в п.~\ref{ssec:orth}.

% ----------------------------------------------------------------------
%  1.2.2. Пространственная решётка; 1.2.3. Кристаллографические символы;
%  1.2.4. Обратная решётка
% ----------------------------------------------------------------------
\subsection{Пространственная решётка}\label{ssec:lattice}

Любой узел решётки задаётся вектором
\begin{keyeq}\label{eq:Rm}
\vb R_m=m_1\vb a_1+m_2\vb a_2+m_3\vb a_3 ,
\end{keyeq}
где $\vb a_1$, $\vb a_2$, $\vb a_3$~— \term{векторы трансляций}, $m_1$, $m_2$,
$m_3$~— целые числа. Начало координат можно поместить в любой узел.
Решётку можно получить бесконечным повторением параллелепипеда,
построенного на кратчайших некомпланарных трансляциях,~— \term{элементарной
ячейки}. Если на её гранях и внутри нет узлов, ячейку называют
\term{примитивной}: каждый её вершинный узел принадлежит восьми соседним
ячейкам, и на ячейку приходится ровно один узел.

\begin{figure}[ht]
\centering
\begin{tikzpicture}[x={(1cm,0cm)}, y={(0.55cm,0.38cm)}, z={(0.12cm,1cm)}, scale=1.35]
  \coordinate (O) at (0,0,0);
  \coordinate (A) at (2.2,0,0); \coordinate (B) at (0.35,1.7,0); \coordinate (C) at (0.25,0.2,1.8);
  \draw[atomLine!70, dashed] (O) -- (B) -- ($(A)+(B)$) ($(B)$) -- ($(B)+(C)$);
  \draw[atomLine] (A) -- ($(A)+(B)$) -- ($(A)+(B)+(C)$) -- ($(A)+(C)$) -- cycle
    ($(C)$) -- ($(B)+(C)$) -- ($(A)+(B)+(C)$) (C) -- ($(A)+(C)$) (O) -- (C);
  \draw[bvec, fmBlue] (O) -- (A) node[below, font=\small, black]{$\vb a_1=\vb a$};
  \draw[bvec, fmBlue] (O) -- (B) node[above left=-2pt, font=\small, black]{$\vb a_2=\vb b$};
  \draw[bvec, fmBlue] (O) -- (C) node[left, font=\small, black]{$\vb a_3=\vb c$};
  \node[font=\small, fmRed] at (0.42,0.32,0.08) {$\gamma$};
  \node[font=\small, fmRed] at (0.55,0.0,0.45) {$\beta$};
  \node[font=\small, fmRed] at (0.12,0.45,0.55) {$\alpha$};
\end{tikzpicture}
\caption{Элементарная ячейка и её параметры: длины рёбер $a$, $b$, $c$ и
углы $\alpha=\widehat{\vb b\vb c}$, $\beta=\widehat{\vb a\vb c}$,
$\gamma=\widehat{\vb a\vb b}$}
\label{fig:cell}
\end{figure}

Длины $a$, $b$, $c$ векторов трансляций и углы $\alpha$, $\beta$, $\gamma$ между
ними~— \term{параметры решётки} (рис.~\ref{fig:cell}). Кристаллографическая
система координат строится на этих векторах; в общем случае она
\emph{косоугольная} с разными масштабами по осям. По симметрии
расположения узлов решётки делятся на семь \term{сингоний}
(табл.~\ref{tab:syngonies}).

\begin{table}[ht]
\centering
\caption{Сингонии (кристаллографические системы координат)}\label{tab:syngonies}
\small
\begin{tabular}{@{}lcl@{}}
\toprule
Сингония & Число параметров & Параметры \\
\midrule
Триклинная & 6 & $a\ne b\ne c$;\quad $\alpha\ne\beta\ne\gamma\ne90^\circ$ \\
Моноклинная & 4 & $a\ne b\ne c$;\quad $\alpha=\gamma=90^\circ$, $\beta\ne90^\circ$ \\
Ромбическая & 3 & $a\ne b\ne c$;\quad $\alpha=\beta=\gamma=90^\circ$ \\
Тетрагональная & 2 & $a=b\ne c$;\quad $\alpha=\beta=\gamma=90^\circ$ \\
Ромбоэдрическая (тригональная) & 2 & $a=b=c$;\quad $\alpha=\beta=\gamma\ne90^\circ$ \\
Гексагональная & 2 & $a=b\ne c$;\quad $\alpha=\beta=90^\circ$, $\gamma=120^\circ$ \\
Кубическая & 1 & $a=b=c$;\quad $\alpha=\beta=\gamma=90^\circ$ \\
\bottomrule
\end{tabular}
\end{table}

\paragraph{Метрическая матрица}
Все метрические свойства решётки (длины, углы, объёмы) содержатся в
\term{метрической матрице} $G$, составленной из скалярных произведений
базисных векторов:
\begin{equation}\label{eq:G}
g_{ij}=(\vb a_i\vb a_j),\qquad a_i=\sqrt{g_{ii}},\qquad
\cos\alpha_{ij}=\frac{g_{ij}}{\sqrt{g_{ii}g_{jj}}},\qquad V=\sqrt{\det G} .
\end{equation}
Матрица симметрична ($g_{ij}=g_{ji}$). Для важнейших сингоний
\[
G_\text{куб}=\begin{pmatrix}a^2&0&0\\0&a^2&0\\0&0&a^2\end{pmatrix},\quad
G_\text{ромб}=\begin{pmatrix}a^2&0&0\\0&b^2&0\\0&0&c^2\end{pmatrix},\quad
G_\text{гекс}=\begin{pmatrix}a^2&-a^2/2&0\\-a^2/2&a^2&0\\0&0&c^2\end{pmatrix}
\]
(тетрагональная~— как ромбическая с $b=a$). Пользу матрицы $G$ мы увидим в
п.~\ref{ssec:formulas}: с её помощью любые углы и расстояния в
косоугольной системе вычисляются по одной формуле.

\paragraph{Сложные ячейки}
Ради наглядности (чтобы ячейка отражала симметрию решётки) часто
выбирают не примитивные, а \emph{сложные} ячейки, у которых есть
дополнительные узлы в центре объёма или граней (правила выбора~—
п.~1.3.6). Положение узла тогда задаётся суммой $\vb R_m+\vb R_n$, где
$\vb R_n=n_1\vb a_1+n_2\vb a_2+n_3\vb a_3$ с \emph{дробными} $n_i$~— это
\term{базис} решётки:
\begin{itemize}
\item объёмноцентрированная ($I$): $[\![000]\!]$, $[\![\tfrac12\tfrac12\tfrac12]\!]$;
\item гранецентрированная ($F$): $[\![000]\!]$, $[\![\tfrac12\tfrac120]\!]$,
$[\![0\tfrac12\tfrac12]\!]$, $[\![\tfrac120\tfrac12]\!]$;
\item базоцентрированная ($A$, $B$ или $C$~— по центрируемой грани), например
$A$: $[\![000]\!]$, $[\![0\tfrac12\tfrac12]\!]$.
\end{itemize}

\subsection{Кристаллографические символы}\label{ssec:symbols}

\paragraph{Символ узла}
Узел с вектором $\vb R=m\vb a+n\vb b+p\vb c$ обозначают $[\![mnp]\!]$. Цифры
пишут подряд без запятых (запятые~— только для двузначных индексов) и
читают порознь; минус ставят \emph{над} цифрой: $[\![1\br20]\!]$ читается
«один, минус два, ноль». У узлов сложной ячейки индексы дробные: центр
объёма~— $[\![\tfrac12\tfrac12\tfrac12]\!]$, центр грани~—
$[\![0\tfrac12\tfrac12]\!]$.

\paragraph{Символ направления}
Решётка~— семейство параллельных узловых прямых. Его характеризует
прямая, проходящая через начало координат; её символ $[uvw]$~— индексы
ближайшего к началу узла на этой прямой (целые взаимно простые числа).
Оси координат имеют символы $[100]$, $[010]$, $[001]$ \emph{в любой
сингонии}~— одно из главных удобств кристаллографической символики.
Совокупность направлений, связанных симметрией, обозначают угловыми
скобками: в кубической сингонии $\nfam{100}$~— это $[100]$, $[010]$,
$[001]$, $[\br100]$, $[0\br10]$, $[00\br1]$, т.\,е. направления,
отличающиеся перестановкой и знаками индексов.

\begin{figure}[ht]
\centering
\begin{subfigure}[b]{0.42\textwidth}\centering
\begin{tikzpicture}[cubeproj, scale=2.4]
  \CubeEdges
  \foreach \v/\l/\a in {{(1,0,0)}/{[100]}/north, {(0,1,0)}/{[010]}/south east, {(0,0,1)}/{[001]}/east,
      {(1,1,0)}/{[110]}/west, {(1,0,1)}/{[101]}/south west, {(1,1,1)}/{[111]}/south}{
    \draw[bvec, fmBlue, line width=0.9pt] (0,0,0) -- \v;
    \node[font=\scriptsize, anchor=\a, fill=white, inner sep=1pt, fill opacity=0.8, text opacity=1] at \v {$\l$};}
  \draw[bvec, fmRed, line width=0.9pt] (1,0,0) -- (0,1,0);
  \node[font=\scriptsize, fmRed, anchor=north] at (0.5,0.45,0) {$[\br110]$};
  \fill (0,0,0) circle (0.6pt);
\end{tikzpicture}
\caption{}
\end{subfigure}\hfill
\begin{subfigure}[b]{0.55\textwidth}\centering
\tikzset{zone100/.style={fmBlue, line width=0.6pt},
  zone110/.style={fmBlue!50, line width=0.4pt},
  pole100/.style={rectangle, fill=fmRed, inner sep=0pt, minimum size=4pt},
  pole110/.style={ellipse, fill=fmGreen, inner sep=0pt, minimum width=4.5pt, minimum height=3pt},
  pole111/.style={regular polygon, regular polygon sides=3, fill=fmAmber!90!black, inner sep=0pt, minimum size=6pt},
  poleH/.style={circle, fill=atomLine, inner sep=0pt, minimum size=2pt}}
\begin{tikzpicture}[scale=3.0]
  \draw[atomLine, line width=0.8pt] (0,0) circle (1);
  % создано tools/gen_stereo.py: стандартная проекция (001) куба (r = 1)
\draw[zone100] (0.0000,-1.0000) -- (0.0000,-0.9742) -- (0.0000,-0.9490) -- (0.0000,-0.9244) -- (0.0000,-0.9004) -- (0.0000,-0.8770) -- (0.0000,-0.8541) -- (0.0000,-0.8317) -- (0.0000,-0.8098) -- (0.0000,-0.7883) -- (0.0000,-0.7673) -- (0.0000,-0.7467) -- (0.0000,-0.7265) -- (0.0000,-0.7067) -- (0.0000,-0.6873) -- (0.0000,-0.6682) -- (0.0000,-0.6494) -- (0.0000,-0.6310) -- (0.0000,-0.6128) -- (0.0000,-0.5949) -- (0.0000,-0.5774) -- (0.0000,-0.5600) -- (0.0000,-0.5430) -- (0.0000,-0.5261) -- (0.0000,-0.5095) -- (0.0000,-0.4931) -- (0.0000,-0.4770) -- (0.0000,-0.4610) -- (0.0000,-0.4452) -- (0.0000,-0.4296) -- (0.0000,-0.4142) -- (0.0000,-0.3990) -- (0.0000,-0.3839) -- (0.0000,-0.3689) -- (0.0000,-0.3541) -- (0.0000,-0.3395) -- (0.0000,-0.3249) -- (0.0000,-0.3105) -- (0.0000,-0.2962) -- (0.0000,-0.2820) -- (0.0000,-0.2679) -- (0.0000,-0.2540) -- (0.0000,-0.2401) -- (0.0000,-0.2263) -- (0.0000,-0.2126) -- (0.0000,-0.1989) -- (0.0000,-0.1853) -- (0.0000,-0.1718) -- (0.0000,-0.1584) -- (0.0000,-0.1450) -- (0.0000,-0.1317) -- (0.0000,-0.1184) -- (0.0000,-0.1051) -- (0.0000,-0.0919) -- (0.0000,-0.0787) -- (0.0000,-0.0655) -- (0.0000,-0.0524) -- (0.0000,-0.0393) -- (0.0000,-0.0262) -- (0.0000,-0.0131) -- (0.0000,-0.0000) -- (-0.0000,0.0131) -- (-0.0000,0.0262) -- (-0.0000,0.0393) -- (-0.0000,0.0524) -- (-0.0000,0.0655) -- (-0.0000,0.0787) -- (-0.0000,0.0919) -- (-0.0000,0.1051) -- (-0.0000,0.1184) -- (-0.0000,0.1317) -- (-0.0000,0.1450) -- (-0.0000,0.1584) -- (-0.0000,0.1718) -- (-0.0000,0.1853) -- (-0.0000,0.1989) -- (-0.0000,0.2126) -- (-0.0000,0.2263) -- (-0.0000,0.2401) -- (-0.0000,0.2540) -- (-0.0000,0.2679) -- (-0.0000,0.2820) -- (-0.0000,0.2962) -- (-0.0000,0.3105) -- (-0.0000,0.3249) -- (-0.0000,0.3395) -- (-0.0000,0.3541) -- (-0.0000,0.3689) -- (-0.0000,0.3839) -- (-0.0000,0.3990) -- (-0.0000,0.4142) -- (-0.0000,0.4296) -- (-0.0000,0.4452) -- (-0.0000,0.4610) -- (-0.0000,0.4770) -- (-0.0000,0.4931) -- (-0.0000,0.5095) -- (-0.0000,0.5261) -- (-0.0000,0.5430) -- (-0.0000,0.5600) -- (-0.0000,0.5774) -- (-0.0000,0.5949) -- (-0.0000,0.6128) -- (-0.0000,0.6310) -- (-0.0000,0.6494) -- (-0.0000,0.6682) -- (-0.0000,0.6873) -- (-0.0000,0.7067) -- (-0.0000,0.7265) -- (-0.0000,0.7467) -- (-0.0000,0.7673) -- (-0.0000,0.7883) -- (-0.0000,0.8098) -- (-0.0000,0.8317) -- (-0.0000,0.8541) -- (-0.0000,0.8770) -- (-0.0000,0.9004) -- (-0.0000,0.9244) -- (-0.0000,0.9490) -- (-0.0000,0.9742) -- (-0.0000,1.0000);
\draw[zone100] (1.0000,0.0000) -- (0.9742,0.0000) -- (0.9490,0.0000) -- (0.9244,0.0000) -- (0.9004,0.0000) -- (0.8770,0.0000) -- (0.8541,0.0000) -- (0.8317,0.0000) -- (0.8098,0.0000) -- (0.7883,0.0000) -- (0.7673,0.0000) -- (0.7467,0.0000) -- (0.7265,0.0000) -- (0.7067,0.0000) -- (0.6873,0.0000) -- (0.6682,0.0000) -- (0.6494,0.0000) -- (0.6310,0.0000) -- (0.6128,0.0000) -- (0.5949,0.0000) -- (0.5774,0.0000) -- (0.5600,0.0000) -- (0.5430,0.0000) -- (0.5261,0.0000) -- (0.5095,0.0000) -- (0.4931,0.0000) -- (0.4770,0.0000) -- (0.4610,0.0000) -- (0.4452,0.0000) -- (0.4296,0.0000) -- (0.4142,0.0000) -- (0.3990,0.0000) -- (0.3839,0.0000) -- (0.3689,0.0000) -- (0.3541,0.0000) -- (0.3395,0.0000) -- (0.3249,0.0000) -- (0.3105,0.0000) -- (0.2962,0.0000) -- (0.2820,0.0000) -- (0.2679,0.0000) -- (0.2540,0.0000) -- (0.2401,0.0000) -- (0.2263,0.0000) -- (0.2126,0.0000) -- (0.1989,0.0000) -- (0.1853,0.0000) -- (0.1718,0.0000) -- (0.1584,0.0000) -- (0.1450,0.0000) -- (0.1317,0.0000) -- (0.1184,0.0000) -- (0.1051,0.0000) -- (0.0919,0.0000) -- (0.0787,0.0000) -- (0.0655,0.0000) -- (0.0524,0.0000) -- (0.0393,0.0000) -- (0.0262,0.0000) -- (0.0131,0.0000) -- (0.0000,0.0000) -- (-0.0131,-0.0000) -- (-0.0262,-0.0000) -- (-0.0393,-0.0000) -- (-0.0524,-0.0000) -- (-0.0655,-0.0000) -- (-0.0787,-0.0000) -- (-0.0919,-0.0000) -- (-0.1051,-0.0000) -- (-0.1184,-0.0000) -- (-0.1317,-0.0000) -- (-0.1450,-0.0000) -- (-0.1584,-0.0000) -- (-0.1718,-0.0000) -- (-0.1853,-0.0000) -- (-0.1989,-0.0000) -- (-0.2126,-0.0000) -- (-0.2263,-0.0000) -- (-0.2401,-0.0000) -- (-0.2540,-0.0000) -- (-0.2679,-0.0000) -- (-0.2820,-0.0000) -- (-0.2962,-0.0000) -- (-0.3105,-0.0000) -- (-0.3249,-0.0000) -- (-0.3395,-0.0000) -- (-0.3541,-0.0000) -- (-0.3689,-0.0000) -- (-0.3839,-0.0000) -- (-0.3990,-0.0000) -- (-0.4142,-0.0000) -- (-0.4296,-0.0000) -- (-0.4452,-0.0000) -- (-0.4610,-0.0000) -- (-0.4770,-0.0000) -- (-0.4931,-0.0000) -- (-0.5095,-0.0000) -- (-0.5261,-0.0000) -- (-0.5430,-0.0000) -- (-0.5600,-0.0000) -- (-0.5774,-0.0000) -- (-0.5949,-0.0000) -- (-0.6128,-0.0000) -- (-0.6310,-0.0000) -- (-0.6494,-0.0000) -- (-0.6682,-0.0000) -- (-0.6873,-0.0000) -- (-0.7067,-0.0000) -- (-0.7265,-0.0000) -- (-0.7467,-0.0000) -- (-0.7673,-0.0000) -- (-0.7883,-0.0000) -- (-0.8098,-0.0000) -- (-0.8317,-0.0000) -- (-0.8541,-0.0000) -- (-0.8770,-0.0000) -- (-0.9004,-0.0000) -- (-0.9244,-0.0000) -- (-0.9490,-0.0000) -- (-0.9742,-0.0000) -- (-1.0000,-0.0000);
\draw[zone110] (0.7071,-0.7071) -- (0.6888,-0.6888) -- (0.6710,-0.6710) -- (0.6536,-0.6536) -- (0.6367,-0.6367) -- (0.6201,-0.6201) -- (0.6039,-0.6039) -- (0.5881,-0.5881) -- (0.5726,-0.5726) -- (0.5574,-0.5574) -- (0.5426,-0.5426) -- (0.5280,-0.5280) -- (0.5137,-0.5137) -- (0.4997,-0.4997) -- (0.4860,-0.4860) -- (0.4725,-0.4725) -- (0.4592,-0.4592) -- (0.4462,-0.4462) -- (0.4333,-0.4333) -- (0.4207,-0.4207) -- (0.4082,-0.4082) -- (0.3960,-0.3960) -- (0.3839,-0.3839) -- (0.3720,-0.3720) -- (0.3603,-0.3603) -- (0.3487,-0.3487) -- (0.3373,-0.3373) -- (0.3260,-0.3260) -- (0.3148,-0.3148) -- (0.3038,-0.3038) -- (0.2929,-0.2929) -- (0.2821,-0.2821) -- (0.2714,-0.2714) -- (0.2609,-0.2609) -- (0.2504,-0.2504) -- (0.2400,-0.2400) -- (0.2298,-0.2298) -- (0.2196,-0.2196) -- (0.2095,-0.2095) -- (0.1994,-0.1994) -- (0.1895,-0.1895) -- (0.1796,-0.1796) -- (0.1698,-0.1698) -- (0.1600,-0.1600) -- (0.1503,-0.1503) -- (0.1407,-0.1407) -- (0.1311,-0.1311) -- (0.1215,-0.1215) -- (0.1120,-0.1120) -- (0.1025,-0.1025) -- (0.0931,-0.0931) -- (0.0837,-0.0837) -- (0.0743,-0.0743) -- (0.0650,-0.0650) -- (0.0557,-0.0557) -- (0.0463,-0.0463) -- (0.0371,-0.0371) -- (0.0278,-0.0278) -- (0.0185,-0.0185) -- (0.0093,-0.0093) -- (0.0000,-0.0000) -- (-0.0093,0.0093) -- (-0.0185,0.0185) -- (-0.0278,0.0278) -- (-0.0371,0.0371) -- (-0.0463,0.0463) -- (-0.0557,0.0557) -- (-0.0650,0.0650) -- (-0.0743,0.0743) -- (-0.0837,0.0837) -- (-0.0931,0.0931) -- (-0.1025,0.1025) -- (-0.1120,0.1120) -- (-0.1215,0.1215) -- (-0.1311,0.1311) -- (-0.1407,0.1407) -- (-0.1503,0.1503) -- (-0.1600,0.1600) -- (-0.1698,0.1698) -- (-0.1796,0.1796) -- (-0.1895,0.1895) -- (-0.1994,0.1994) -- (-0.2095,0.2095) -- (-0.2196,0.2196) -- (-0.2298,0.2298) -- (-0.2400,0.2400) -- (-0.2504,0.2504) -- (-0.2609,0.2609) -- (-0.2714,0.2714) -- (-0.2821,0.2821) -- (-0.2929,0.2929) -- (-0.3038,0.3038) -- (-0.3148,0.3148) -- (-0.3260,0.3260) -- (-0.3373,0.3373) -- (-0.3487,0.3487) -- (-0.3603,0.3603) -- (-0.3720,0.3720) -- (-0.3839,0.3839) -- (-0.3960,0.3960) -- (-0.4082,0.4082) -- (-0.4207,0.4207) -- (-0.4333,0.4333) -- (-0.4462,0.4462) -- (-0.4592,0.4592) -- (-0.4725,0.4725) -- (-0.4860,0.4860) -- (-0.4997,0.4997) -- (-0.5137,0.5137) -- (-0.5280,0.5280) -- (-0.5426,0.5426) -- (-0.5574,0.5574) -- (-0.5726,0.5726) -- (-0.5881,0.5881) -- (-0.6039,0.6039) -- (-0.6201,0.6201) -- (-0.6367,0.6367) -- (-0.6536,0.6536) -- (-0.6710,0.6710) -- (-0.6888,0.6888) -- (-0.7071,0.7071);
\draw[zone110] (-0.7071,-0.7071) -- (-0.6888,-0.6888) -- (-0.6710,-0.6710) -- (-0.6536,-0.6536) -- (-0.6367,-0.6367) -- (-0.6201,-0.6201) -- (-0.6039,-0.6039) -- (-0.5881,-0.5881) -- (-0.5726,-0.5726) -- (-0.5574,-0.5574) -- (-0.5426,-0.5426) -- (-0.5280,-0.5280) -- (-0.5137,-0.5137) -- (-0.4997,-0.4997) -- (-0.4860,-0.4860) -- (-0.4725,-0.4725) -- (-0.4592,-0.4592) -- (-0.4462,-0.4462) -- (-0.4333,-0.4333) -- (-0.4207,-0.4207) -- (-0.4082,-0.4082) -- (-0.3960,-0.3960) -- (-0.3839,-0.3839) -- (-0.3720,-0.3720) -- (-0.3603,-0.3603) -- (-0.3487,-0.3487) -- (-0.3373,-0.3373) -- (-0.3260,-0.3260) -- (-0.3148,-0.3148) -- (-0.3038,-0.3038) -- (-0.2929,-0.2929) -- (-0.2821,-0.2821) -- (-0.2714,-0.2714) -- (-0.2609,-0.2609) -- (-0.2504,-0.2504) -- (-0.2400,-0.2400) -- (-0.2298,-0.2298) -- (-0.2196,-0.2196) -- (-0.2095,-0.2095) -- (-0.1994,-0.1994) -- (-0.1895,-0.1895) -- (-0.1796,-0.1796) -- (-0.1698,-0.1698) -- (-0.1600,-0.1600) -- (-0.1503,-0.1503) -- (-0.1407,-0.1407) -- (-0.1311,-0.1311) -- (-0.1215,-0.1215) -- (-0.1120,-0.1120) -- (-0.1025,-0.1025) -- (-0.0931,-0.0931) -- (-0.0837,-0.0837) -- (-0.0743,-0.0743) -- (-0.0650,-0.0650) -- (-0.0557,-0.0557) -- (-0.0463,-0.0463) -- (-0.0371,-0.0371) -- (-0.0278,-0.0278) -- (-0.0185,-0.0185) -- (-0.0093,-0.0093) -- (-0.0000,-0.0000) -- (0.0093,0.0093) -- (0.0185,0.0185) -- (0.0278,0.0278) -- (0.0371,0.0371) -- (0.0463,0.0463) -- (0.0557,0.0557) -- (0.0650,0.0650) -- (0.0743,0.0743) -- (0.0837,0.0837) -- (0.0931,0.0931) -- (0.1025,0.1025) -- (0.1120,0.1120) -- (0.1215,0.1215) -- (0.1311,0.1311) -- (0.1407,0.1407) -- (0.1503,0.1503) -- (0.1600,0.1600) -- (0.1698,0.1698) -- (0.1796,0.1796) -- (0.1895,0.1895) -- (0.1994,0.1994) -- (0.2095,0.2095) -- (0.2196,0.2196) -- (0.2298,0.2298) -- (0.2400,0.2400) -- (0.2504,0.2504) -- (0.2609,0.2609) -- (0.2714,0.2714) -- (0.2821,0.2821) -- (0.2929,0.2929) -- (0.3038,0.3038) -- (0.3148,0.3148) -- (0.3260,0.3260) -- (0.3373,0.3373) -- (0.3487,0.3487) -- (0.3603,0.3603) -- (0.3720,0.3720) -- (0.3839,0.3839) -- (0.3960,0.3960) -- (0.4082,0.4082) -- (0.4207,0.4207) -- (0.4333,0.4333) -- (0.4462,0.4462) -- (0.4592,0.4592) -- (0.4725,0.4725) -- (0.4860,0.4860) -- (0.4997,0.4997) -- (0.5137,0.5137) -- (0.5280,0.5280) -- (0.5426,0.5426) -- (0.5574,0.5574) -- (0.5726,0.5726) -- (0.5881,0.5881) -- (0.6039,0.6039) -- (0.6201,0.6201) -- (0.6367,0.6367) -- (0.6536,0.6536) -- (0.6710,0.6710) -- (0.6888,0.6888) -- (0.7071,0.7071);
\draw[zone110] (0.0000,-1.0000) -- (-0.0182,-0.9815) -- (-0.0357,-0.9630) -- (-0.0526,-0.9445) -- (-0.0688,-0.9261) -- (-0.0845,-0.9077) -- (-0.0996,-0.8893) -- (-0.1142,-0.8710) -- (-0.1282,-0.8528) -- (-0.1417,-0.8346) -- (-0.1547,-0.8165) -- (-0.1672,-0.7985) -- (-0.1793,-0.7805) -- (-0.1910,-0.7626) -- (-0.2022,-0.7448) -- (-0.2130,-0.7271) -- (-0.2234,-0.7095) -- (-0.2334,-0.6919) -- (-0.2430,-0.6745) -- (-0.2523,-0.6571) -- (-0.2612,-0.6398) -- (-0.2698,-0.6226) -- (-0.2780,-0.6055) -- (-0.2860,-0.5884) -- (-0.2936,-0.5715) -- (-0.3009,-0.5546) -- (-0.3080,-0.5378) -- (-0.3147,-0.5211) -- (-0.3212,-0.5045) -- (-0.3274,-0.4879) -- (-0.3333,-0.4714) -- (-0.3390,-0.4550) -- (-0.3445,-0.4386) -- (-0.3497,-0.4224) -- (-0.3546,-0.4061) -- (-0.3594,-0.3900) -- (-0.3639,-0.3739) -- (-0.3682,-0.3579) -- (-0.3723,-0.3419) -- (-0.3761,-0.3260) -- (-0.3798,-0.3101) -- (-0.3833,-0.2943) -- (-0.3865,-0.2785) -- (-0.3896,-0.2628) -- (-0.3925,-0.2471) -- (-0.3951,-0.2315) -- (-0.3976,-0.2159) -- (-0.4000,-0.2003) -- (-0.4021,-0.1848) -- (-0.4040,-0.1693) -- (-0.4058,-0.1538) -- (-0.4074,-0.1383) -- (-0.4089,-0.1229) -- (-0.4101,-0.1075) -- (-0.4112,-0.0921) -- (-0.4121,-0.0767) -- (-0.4129,-0.0614) -- (-0.4135,-0.0460) -- (-0.4139,-0.0307) -- (-0.4141,-0.0153) -- (-0.4142,-0.0000) -- (-0.4141,0.0153) -- (-0.4139,0.0307) -- (-0.4135,0.0460) -- (-0.4129,0.0614) -- (-0.4121,0.0767) -- (-0.4112,0.0921) -- (-0.4101,0.1075) -- (-0.4089,0.1229) -- (-0.4074,0.1383) -- (-0.4058,0.1538) -- (-0.4040,0.1693) -- (-0.4021,0.1848) -- (-0.4000,0.2003) -- (-0.3976,0.2159) -- (-0.3951,0.2315) -- (-0.3925,0.2471) -- (-0.3896,0.2628) -- (-0.3865,0.2785) -- (-0.3833,0.2943) -- (-0.3798,0.3101) -- (-0.3761,0.3260) -- (-0.3723,0.3419) -- (-0.3682,0.3579) -- (-0.3639,0.3739) -- (-0.3594,0.3900) -- (-0.3546,0.4061) -- (-0.3497,0.4224) -- (-0.3445,0.4386) -- (-0.3390,0.4550) -- (-0.3333,0.4714) -- (-0.3274,0.4879) -- (-0.3212,0.5045) -- (-0.3147,0.5211) -- (-0.3080,0.5378) -- (-0.3009,0.5546) -- (-0.2936,0.5715) -- (-0.2860,0.5884) -- (-0.2780,0.6055) -- (-0.2698,0.6226) -- (-0.2612,0.6398) -- (-0.2523,0.6571) -- (-0.2430,0.6745) -- (-0.2334,0.6919) -- (-0.2234,0.7095) -- (-0.2130,0.7271) -- (-0.2022,0.7448) -- (-0.1910,0.7626) -- (-0.1793,0.7805) -- (-0.1672,0.7985) -- (-0.1547,0.8165) -- (-0.1417,0.8346) -- (-0.1282,0.8528) -- (-0.1142,0.8710) -- (-0.0996,0.8893) -- (-0.0845,0.9077) -- (-0.0688,0.9261) -- (-0.0526,0.9445) -- (-0.0357,0.9630) -- (-0.0182,0.9815) -- (-0.0000,1.0000);
\draw[zone110] (0.0000,-1.0000) -- (0.0182,-0.9815) -- (0.0357,-0.9630) -- (0.0526,-0.9445) -- (0.0688,-0.9261) -- (0.0845,-0.9077) -- (0.0996,-0.8893) -- (0.1142,-0.8710) -- (0.1282,-0.8528) -- (0.1417,-0.8346) -- (0.1547,-0.8165) -- (0.1672,-0.7985) -- (0.1793,-0.7805) -- (0.1910,-0.7626) -- (0.2022,-0.7448) -- (0.2130,-0.7271) -- (0.2234,-0.7095) -- (0.2334,-0.6919) -- (0.2430,-0.6745) -- (0.2523,-0.6571) -- (0.2612,-0.6398) -- (0.2698,-0.6226) -- (0.2780,-0.6055) -- (0.2860,-0.5884) -- (0.2936,-0.5715) -- (0.3009,-0.5546) -- (0.3080,-0.5378) -- (0.3147,-0.5211) -- (0.3212,-0.5045) -- (0.3274,-0.4879) -- (0.3333,-0.4714) -- (0.3390,-0.4550) -- (0.3445,-0.4386) -- (0.3497,-0.4224) -- (0.3546,-0.4061) -- (0.3594,-0.3900) -- (0.3639,-0.3739) -- (0.3682,-0.3579) -- (0.3723,-0.3419) -- (0.3761,-0.3260) -- (0.3798,-0.3101) -- (0.3833,-0.2943) -- (0.3865,-0.2785) -- (0.3896,-0.2628) -- (0.3925,-0.2471) -- (0.3951,-0.2315) -- (0.3976,-0.2159) -- (0.4000,-0.2003) -- (0.4021,-0.1848) -- (0.4040,-0.1693) -- (0.4058,-0.1538) -- (0.4074,-0.1383) -- (0.4089,-0.1229) -- (0.4101,-0.1075) -- (0.4112,-0.0921) -- (0.4121,-0.0767) -- (0.4129,-0.0614) -- (0.4135,-0.0460) -- (0.4139,-0.0307) -- (0.4141,-0.0153) -- (0.4142,-0.0000) -- (0.4141,0.0153) -- (0.4139,0.0307) -- (0.4135,0.0460) -- (0.4129,0.0614) -- (0.4121,0.0767) -- (0.4112,0.0921) -- (0.4101,0.1075) -- (0.4089,0.1229) -- (0.4074,0.1383) -- (0.4058,0.1538) -- (0.4040,0.1693) -- (0.4021,0.1848) -- (0.4000,0.2003) -- (0.3976,0.2159) -- (0.3951,0.2315) -- (0.3925,0.2471) -- (0.3896,0.2628) -- (0.3865,0.2785) -- (0.3833,0.2943) -- (0.3798,0.3101) -- (0.3761,0.3260) -- (0.3723,0.3419) -- (0.3682,0.3579) -- (0.3639,0.3739) -- (0.3594,0.3900) -- (0.3546,0.4061) -- (0.3497,0.4224) -- (0.3445,0.4386) -- (0.3390,0.4550) -- (0.3333,0.4714) -- (0.3274,0.4879) -- (0.3212,0.5045) -- (0.3147,0.5211) -- (0.3080,0.5378) -- (0.3009,0.5546) -- (0.2936,0.5715) -- (0.2860,0.5884) -- (0.2780,0.6055) -- (0.2698,0.6226) -- (0.2612,0.6398) -- (0.2523,0.6571) -- (0.2430,0.6745) -- (0.2334,0.6919) -- (0.2234,0.7095) -- (0.2130,0.7271) -- (0.2022,0.7448) -- (0.1910,0.7626) -- (0.1793,0.7805) -- (0.1672,0.7985) -- (0.1547,0.8165) -- (0.1417,0.8346) -- (0.1282,0.8528) -- (0.1142,0.8710) -- (0.0996,0.8893) -- (0.0845,0.9077) -- (0.0688,0.9261) -- (0.0526,0.9445) -- (0.0357,0.9630) -- (0.0182,0.9815) -- (0.0000,1.0000);
\draw[zone110] (1.0000,0.0000) -- (0.9815,-0.0182) -- (0.9630,-0.0357) -- (0.9445,-0.0526) -- (0.9261,-0.0688) -- (0.9077,-0.0845) -- (0.8893,-0.0996) -- (0.8710,-0.1142) -- (0.8528,-0.1282) -- (0.8346,-0.1417) -- (0.8165,-0.1547) -- (0.7985,-0.1672) -- (0.7805,-0.1793) -- (0.7626,-0.1910) -- (0.7448,-0.2022) -- (0.7271,-0.2130) -- (0.7095,-0.2234) -- (0.6919,-0.2334) -- (0.6745,-0.2430) -- (0.6571,-0.2523) -- (0.6398,-0.2612) -- (0.6226,-0.2698) -- (0.6055,-0.2780) -- (0.5884,-0.2860) -- (0.5715,-0.2936) -- (0.5546,-0.3009) -- (0.5378,-0.3080) -- (0.5211,-0.3147) -- (0.5045,-0.3212) -- (0.4879,-0.3274) -- (0.4714,-0.3333) -- (0.4550,-0.3390) -- (0.4386,-0.3445) -- (0.4224,-0.3497) -- (0.4061,-0.3546) -- (0.3900,-0.3594) -- (0.3739,-0.3639) -- (0.3579,-0.3682) -- (0.3419,-0.3723) -- (0.3260,-0.3761) -- (0.3101,-0.3798) -- (0.2943,-0.3833) -- (0.2785,-0.3865) -- (0.2628,-0.3896) -- (0.2471,-0.3925) -- (0.2315,-0.3951) -- (0.2159,-0.3976) -- (0.2003,-0.4000) -- (0.1848,-0.4021) -- (0.1693,-0.4040) -- (0.1538,-0.4058) -- (0.1383,-0.4074) -- (0.1229,-0.4089) -- (0.1075,-0.4101) -- (0.0921,-0.4112) -- (0.0767,-0.4121) -- (0.0614,-0.4129) -- (0.0460,-0.4135) -- (0.0307,-0.4139) -- (0.0153,-0.4141) -- (0.0000,-0.4142) -- (-0.0153,-0.4141) -- (-0.0307,-0.4139) -- (-0.0460,-0.4135) -- (-0.0614,-0.4129) -- (-0.0767,-0.4121) -- (-0.0921,-0.4112) -- (-0.1075,-0.4101) -- (-0.1229,-0.4089) -- (-0.1383,-0.4074) -- (-0.1538,-0.4058) -- (-0.1693,-0.4040) -- (-0.1848,-0.4021) -- (-0.2003,-0.4000) -- (-0.2159,-0.3976) -- (-0.2315,-0.3951) -- (-0.2471,-0.3925) -- (-0.2628,-0.3896) -- (-0.2785,-0.3865) -- (-0.2943,-0.3833) -- (-0.3101,-0.3798) -- (-0.3260,-0.3761) -- (-0.3419,-0.3723) -- (-0.3579,-0.3682) -- (-0.3739,-0.3639) -- (-0.3900,-0.3594) -- (-0.4061,-0.3546) -- (-0.4224,-0.3497) -- (-0.4386,-0.3445) -- (-0.4550,-0.3390) -- (-0.4714,-0.3333) -- (-0.4879,-0.3274) -- (-0.5045,-0.3212) -- (-0.5211,-0.3147) -- (-0.5378,-0.3080) -- (-0.5546,-0.3009) -- (-0.5715,-0.2936) -- (-0.5884,-0.2860) -- (-0.6055,-0.2780) -- (-0.6226,-0.2698) -- (-0.6398,-0.2612) -- (-0.6571,-0.2523) -- (-0.6745,-0.2430) -- (-0.6919,-0.2334) -- (-0.7095,-0.2234) -- (-0.7271,-0.2130) -- (-0.7448,-0.2022) -- (-0.7626,-0.1910) -- (-0.7805,-0.1793) -- (-0.7985,-0.1672) -- (-0.8165,-0.1547) -- (-0.8346,-0.1417) -- (-0.8528,-0.1282) -- (-0.8710,-0.1142) -- (-0.8893,-0.0996) -- (-0.9077,-0.0845) -- (-0.9261,-0.0688) -- (-0.9445,-0.0526) -- (-0.9630,-0.0357) -- (-0.9815,-0.0182) -- (-1.0000,-0.0000);
\draw[zone110] (1.0000,0.0000) -- (0.9815,0.0182) -- (0.9630,0.0357) -- (0.9445,0.0526) -- (0.9261,0.0688) -- (0.9077,0.0845) -- (0.8893,0.0996) -- (0.8710,0.1142) -- (0.8528,0.1282) -- (0.8346,0.1417) -- (0.8165,0.1547) -- (0.7985,0.1672) -- (0.7805,0.1793) -- (0.7626,0.1910) -- (0.7448,0.2022) -- (0.7271,0.2130) -- (0.7095,0.2234) -- (0.6919,0.2334) -- (0.6745,0.2430) -- (0.6571,0.2523) -- (0.6398,0.2612) -- (0.6226,0.2698) -- (0.6055,0.2780) -- (0.5884,0.2860) -- (0.5715,0.2936) -- (0.5546,0.3009) -- (0.5378,0.3080) -- (0.5211,0.3147) -- (0.5045,0.3212) -- (0.4879,0.3274) -- (0.4714,0.3333) -- (0.4550,0.3390) -- (0.4386,0.3445) -- (0.4224,0.3497) -- (0.4061,0.3546) -- (0.3900,0.3594) -- (0.3739,0.3639) -- (0.3579,0.3682) -- (0.3419,0.3723) -- (0.3260,0.3761) -- (0.3101,0.3798) -- (0.2943,0.3833) -- (0.2785,0.3865) -- (0.2628,0.3896) -- (0.2471,0.3925) -- (0.2315,0.3951) -- (0.2159,0.3976) -- (0.2003,0.4000) -- (0.1848,0.4021) -- (0.1693,0.4040) -- (0.1538,0.4058) -- (0.1383,0.4074) -- (0.1229,0.4089) -- (0.1075,0.4101) -- (0.0921,0.4112) -- (0.0767,0.4121) -- (0.0614,0.4129) -- (0.0460,0.4135) -- (0.0307,0.4139) -- (0.0153,0.4141) -- (0.0000,0.4142) -- (-0.0153,0.4141) -- (-0.0307,0.4139) -- (-0.0460,0.4135) -- (-0.0614,0.4129) -- (-0.0767,0.4121) -- (-0.0921,0.4112) -- (-0.1075,0.4101) -- (-0.1229,0.4089) -- (-0.1383,0.4074) -- (-0.1538,0.4058) -- (-0.1693,0.4040) -- (-0.1848,0.4021) -- (-0.2003,0.4000) -- (-0.2159,0.3976) -- (-0.2315,0.3951) -- (-0.2471,0.3925) -- (-0.2628,0.3896) -- (-0.2785,0.3865) -- (-0.2943,0.3833) -- (-0.3101,0.3798) -- (-0.3260,0.3761) -- (-0.3419,0.3723) -- (-0.3579,0.3682) -- (-0.3739,0.3639) -- (-0.3900,0.3594) -- (-0.4061,0.3546) -- (-0.4224,0.3497) -- (-0.4386,0.3445) -- (-0.4550,0.3390) -- (-0.4714,0.3333) -- (-0.4879,0.3274) -- (-0.5045,0.3212) -- (-0.5211,0.3147) -- (-0.5378,0.3080) -- (-0.5546,0.3009) -- (-0.5715,0.2936) -- (-0.5884,0.2860) -- (-0.6055,0.2780) -- (-0.6226,0.2698) -- (-0.6398,0.2612) -- (-0.6571,0.2523) -- (-0.6745,0.2430) -- (-0.6919,0.2334) -- (-0.7095,0.2234) -- (-0.7271,0.2130) -- (-0.7448,0.2022) -- (-0.7626,0.1910) -- (-0.7805,0.1793) -- (-0.7985,0.1672) -- (-0.8165,0.1547) -- (-0.8346,0.1417) -- (-0.8528,0.1282) -- (-0.8710,0.1142) -- (-0.8893,0.0996) -- (-0.9077,0.0845) -- (-0.9261,0.0688) -- (-0.9445,0.0526) -- (-0.9630,0.0357) -- (-0.9815,0.0182) -- (-1.0000,0.0000);
\node[pole100] at (1.0000,0.0000) {};
\node[font=\scriptsize, inner sep=1pt, anchor=180] at (1.0000,0.0000) {$100$};
\node[pole100] at (-1.0000,0.0000) {};
\node[font=\scriptsize, inner sep=1pt, anchor=360] at (-1.0000,0.0000) {$\br{1}00$};
\node[pole100] at (0.0000,0.0000) {};
\node[font=\scriptsize, inner sep=1pt, anchor=90] at (0.0000,0.0000) {$001$};
\node[pole100] at (0.0000,1.0000) {};
\node[font=\scriptsize, inner sep=1pt, anchor=270] at (0.0000,1.0000) {$010$};
\node[pole100] at (0.0000,-1.0000) {};
\node[font=\scriptsize, inner sep=1pt, anchor=90] at (0.0000,-1.0000) {$0\br{1}0$};
\node[pole110] at (0.4142,0.0000) {};
\node[font=\scriptsize, inner sep=1pt, anchor=180] at (0.4142,0.0000) {$101$};
\node[pole110] at (-0.4142,0.0000) {};
\node[font=\scriptsize, inner sep=1pt, anchor=360] at (-0.4142,0.0000) {$\br{1}01$};
\node[pole110] at (0.7071,0.7071) {};
\node[font=\scriptsize, inner sep=1pt, anchor=225] at (0.7071,0.7071) {$110$};
\node[pole110] at (0.7071,-0.7071) {};
\node[font=\scriptsize, inner sep=1pt, anchor=135] at (0.7071,-0.7071) {$1\br{1}0$};
\node[pole110] at (-0.7071,0.7071) {};
\node[font=\scriptsize, inner sep=1pt, anchor=315] at (-0.7071,0.7071) {$\br{1}10$};
\node[pole110] at (-0.7071,-0.7071) {};
\node[font=\scriptsize, inner sep=1pt, anchor=45] at (-0.7071,-0.7071) {$\br{1}\br{1}0$};
\node[pole110] at (0.0000,0.4142) {};
\node[font=\scriptsize, inner sep=1pt, anchor=270] at (0.0000,0.4142) {$011$};
\node[pole110] at (0.0000,-0.4142) {};
\node[font=\scriptsize, inner sep=1pt, anchor=90] at (0.0000,-0.4142) {$0\br{1}1$};
\node[pole111] at (0.3660,0.3660) {};
\node[font=\scriptsize, inner sep=1pt, anchor=225] at (0.3660,0.3660) {$111$};
\node[pole111] at (0.3660,-0.3660) {};
\node[font=\scriptsize, inner sep=1pt, anchor=135] at (0.3660,-0.3660) {$1\br{1}1$};
\node[pole111] at (-0.3660,0.3660) {};
\node[font=\scriptsize, inner sep=1pt, anchor=315] at (-0.3660,0.3660) {$\br{1}11$};
\node[pole111] at (-0.3660,-0.3660) {};
\node[font=\scriptsize, inner sep=1pt, anchor=45] at (-0.3660,-0.3660) {$\br{1}\br{1}1$};
\node[poleH] at (0.0000,0.6180) {};
\node[poleH] at (0.0000,-0.6180) {};
\node[poleH] at (0.4472,0.8944) {};
\node[poleH] at (0.4472,-0.8944) {};
\node[poleH] at (-0.4472,0.8944) {};
\node[poleH] at (-0.4472,-0.8944) {};
\node[poleH] at (0.8944,0.4472) {};
\node[poleH] at (0.8944,-0.4472) {};
\node[poleH] at (-0.8944,0.4472) {};
\node[poleH] at (-0.8944,-0.4472) {};
\node[poleH] at (0.6180,0.0000) {};
\node[poleH] at (-0.6180,0.0000) {};
\node[poleH] at (0.0000,0.2361) {};
\node[poleH] at (0.0000,-0.2361) {};
\node[poleH] at (0.2361,0.0000) {};
\node[poleH] at (-0.2361,0.0000) {};
\node[poleH] at (0.2899,0.5798) {};
\node[poleH] at (0.2899,-0.5798) {};
\node[poleH] at (-0.2899,0.5798) {};
\node[poleH] at (-0.2899,-0.5798) {};
\node[poleH] at (0.5798,0.2899) {};
\node[poleH] at (0.5798,-0.2899) {};
\node[poleH] at (-0.5798,0.2899) {};
\node[poleH] at (-0.5798,-0.2899) {};
\node[poleH] at (0.2247,0.2247) {};
\node[poleH] at (0.2247,-0.2247) {};
\node[poleH] at (-0.2247,0.2247) {};
\node[poleH] at (-0.2247,-0.2247) {};
\node[poleH] at (0.2000,0.4000) {};
\node[poleH] at (0.2000,-0.4000) {};
\node[poleH] at (-0.2000,0.4000) {};
\node[poleH] at (-0.2000,-0.4000) {};
\node[poleH] at (0.4000,0.2000) {};
\node[poleH] at (0.4000,-0.2000) {};
\node[poleH] at (-0.4000,0.2000) {};
\node[poleH] at (-0.4000,-0.2000) {};
\node[poleH] at (0.5000,0.5000) {};
\node[poleH] at (0.5000,-0.5000) {};
\node[poleH] at (-0.5000,0.5000) {};
\node[poleH] at (-0.5000,-0.5000) {};

\end{tikzpicture}
\caption{}
\end{subfigure}
\caption{Важнейшие направления кубического кристалла: \textit{а}~— в
элементарной ячейке (начало координат~— в левой нижней передней
вершине); \textit{б}~— стандартная стереографическая проекция $(001)$
(расчёт): полюсы $\{100\}$ (квадраты), $\{110\}$ (овалы), $\{111\}$
(треугольники), $\{210\}$, $\{211\}$, $\{221\}$ (точки) и следы
плоскостей $\{100\}$ и $\{110\}$. В кубическом кристалле направление
$[uvw]$ перпендикулярно плоскости $(uvw)$, поэтому полюсы плоскостей и
выходы направлений совпадают}
\label{fig:cubic-dirs}
\end{figure}

\paragraph{Примитивные ячейки ГЦК и ОЦК}
В сложных ячейках можно выделить примитивные, построенные на кратчайших
трансляциях (рис.~\ref{fig:prim}). Для ГЦК-решётки это векторы
$\frac a2[110]$, $\frac a2[011]$, $\frac a2[101]$, для ОЦК~—
$\frac a2[11\br1]$, $\frac a2[\br111]$, $\frac a2[1\br11]$. В матричной
записи (строки~— векторы)
\begin{equation}\label{eq:prim}
A_\text{ГЦК}=\frac a2\begin{pmatrix}1&1&0\\0&1&1\\1&0&1\end{pmatrix},\qquad
A_\text{ОЦК}=\frac a2\begin{pmatrix}1&1&\br1\\\br1&1&1\\1&\br1&1\end{pmatrix}.
\end{equation}
Объём примитивной ячейки ГЦК равен $a^3/4$, ОЦК~— $a^3/2$ (в
кубической ячейке 4 и 2 узла соответственно).

\begin{figure}[ht]
\centering
\begin{subfigure}[b]{0.45\textwidth}\centering
\begin{tikzpicture}[cubeproj, scale=2.4]
  \CubeEdges[atomLine!60]
  \coordinate (v1) at (0.5,0.5,0); \coordinate (v2) at (0,0.5,0.5); \coordinate (v3) at (0.5,0,0.5);
  \draw[fmRed!70, line width=0.5pt] (v1) -- (0.5,1,0.5) -- (v2) (v1) -- (1,0.5,0.5) -- (v3)
    (v2) -- (0.5,0.5,1) -- (v3) (0.5,1,0.5) -- (1,1,1) (1,0.5,0.5) -- (1,1,1) (0.5,0.5,1) -- (1,1,1);
  \FCCback[2.2mm] \FCCmid[2.2mm] \FCCfront[2.2mm]
  \draw[bvec, fmRed] (0,0,0) -- (v1); \draw[bvec, fmRed] (0,0,0) -- (v2); \draw[bvec, fmRed] (0,0,0) -- (v3);
\end{tikzpicture}
\caption{ГЦК}
\end{subfigure}\hfill
\begin{subfigure}[b]{0.45\textwidth}\centering
\begin{tikzpicture}[cubeproj, scale=2.4]
  \CubeEdges[atomLine!60]
  \foreach \p in {(0,0,0),(1,0,0),(0,1,0),(1,1,0),(0,0,1),(1,0,1),(0,1,1),(1,1,1)}{\node[atom=2.6mm] at \p {};}
  \node[atomB=2.6mm] at (0.5,0.5,0.5) {};
  \draw[bvec, fmRed] (0.5,0.5,0.5) -- (1,1,0);
  \draw[bvec, fmRed] (0.5,0.5,0.5) -- (0,1,1);
  \draw[bvec, fmRed] (0.5,0.5,0.5) -- (1,0,1);
\end{tikzpicture}
\caption{ОЦК}
\end{subfigure}
\caption{Векторы примитивных трансляций (красные стрелки) в ГЦК- и
ОЦК-решётках. Для ГЦК показана и сама примитивная ячейка~—
ромбоэдр с углом $60^\circ$; для ОЦК векторы проведены из центрального узла
(ромбоэдр с углом $109{,}5^\circ$ выходит за пределы куба)}
\label{fig:prim}
\end{figure}

\paragraph{Символ плоскости (индексы Миллера)}
Решётку можно представить и как семейство параллельных узловых
плоскостей. Пусть плоскость отсекает на осях отрезки $A$, $B$, $C$; её
уравнение $x/A+y/B+z/C=1$ выполняется для узлов $x=ma$, $y=nb$, $z=pc$, т.\,е.
$m(a/A)+n(b/B)+p(c/C)=1$. Отсюда отношения $a/A$, $b/B$, $c/C$ рациональны и
их можно записать через три взаимно простых целых числа:
\begin{keyeq}\label{eq:miller}
h:k:l=\frac aA:\frac bB:\frac cC .
\end{keyeq}

\begin{opredelenie}[Индексы Миллера]
Три взаимно простых целых числа $h$, $k$, $l$, обратно пропорциональные
отрезкам, которые плоскость отсекает на осях (в единицах $a$, $b$, $c$).
Символ плоскости записывают в круглых скобках: $(hkl)$.
\end{opredelenie}

\begin{shagi}
\item Найти отрезки, отсекаемые плоскостью на осях, в единицах периодов
(если плоскость параллельна оси, отрезок~— $\infty$; если проходит через
начало~— сдвинуть начало в соседний узел).
\item Взять обратные величины.
\item Привести к наименьшим целым числам; отрицательный индекс отметить
чертой сверху.
\end{shagi}

\begin{primer}[Индексы плоскости]
Плоскость отсекает отрезки $1$, $2$ и $\infty$. Обратные величины $1$,
$\frac12$, $0$; умножая на 2, получаем $(210)$. Плоскость с отрезками
$\frac12$, $1$, $\infty$ имеет те же индексы $(210)$ и параллельна первой:
все плоскости семейства имеют один символ.
\end{primer}

Подставив $h=q\,a/A$, $k=q\,b/B$, $l=q\,c/C$ в уравнение плоскости, получим
\begin{equation}\label{eq:plane-q}
hm+kn+lp=q ,
\end{equation}
где $q$~— целое число, «номер» плоскости в семействе: $q=0$~— плоскость
через начало координат, $q=1$~— ближайшая к нему; она отсекает отрезки
$a/h$, $b/k$, $c/l$. Расположение важнейших плоскостей кубической решётки
показано на рис.~\ref{fig:planes}.

\begin{figure}[ht]
\centering
\tikzset{plshade/.style={fill=fmBlue!25, draw=fmBlue, line width=0.6pt, fill opacity=0.8}}
\foreach \pl/\lab in {
  {(1,0,0) -- (1,1,0) -- (1,1,1) -- (1,0,1) -- cycle}/{(100)},
  {(1,0,0) -- (0,1,0) -- (0,1,1) -- (1,0,1) -- cycle}/{(110)},
  {(1,0,0) -- (0,1,0) -- (0,0,1) -- cycle}/{(111)},
  {(0.5,0,0) -- (0,1,0) -- (0,1,1) -- (0.5,0,1) -- cycle}/{(210)}}{%
\begin{tikzpicture}[cubeproj, scale=1.6, baseline=(b)]
  \coordinate (b) at (0,0,0);
  \CubeEdges
  \filldraw[plshade] \pl;
  \node[font=\small] at (0.6,0,-0.32) {$\lab$};
\end{tikzpicture}\hspace{4mm}}
\caption{Плоскости $(100)$, $(110)$, $(111)$ и $(210)$ в кубической ячейке
(начало координат~— в левой нижней передней вершине)}
\label{fig:planes}
\end{figure}

\paragraph{Кристаллическая форма и множитель повторяемости}
Совокупность плоскостей, связанных элементами симметрии, обозначают
фигурными скобками $\{hkl\}$ и называют \term{кристаллической формой}. В
кубической сингонии $\{100\}$~— это $(100)$, $(010)$, $(001)$, $(\br100)$,
$(0\br10)$, $(00\br1)$~— шесть граней куба. Число плоскостей формы~—
\term{множитель повторяемости} $p$: для кубической сингонии
$p_{\{100\}}=6$, $p_{\{110\}}=12$, $p_{\{111\}}=8$, $p_{\{123\}}=48$. Все плоскости
одной формы имеют одинаковое межплоскостное расстояние.

\paragraph{Четырёхиндексные символы в гексагональной системе}
Базисная плоскость гексагональной решётки имеет три равноправные оси
$\vb a_1$, $\vb a_2$, $\vb a_3=-(\vb a_1+\vb a_2)$ под углами $120^\circ$
(рис.~\ref{fig:hex4}). Поэтому плоскости обозначают четырьмя индексами
$(hkil)$, где третий индекс относится к оси $\vb a_3$. Он не независим:
\begin{keyeq}\label{eq:hkil}
i=-(h+k) .
\end{keyeq}
(Вывод: если след плоскости отсекает на осях $\vb a_1$, $\vb a_2$, $\vb a_3$
отрезки $A_1$, $A_2$, $-A_3$, то из подобия треугольников
$1/A_1+1/A_2=1/A_3$, а индексы обратны отрезкам.) Выгода~— в том, что
симметрично равнозначные плоскости получаются перестановкой первых трёх
индексов: шесть граней призмы первого рода~— $(10\br10)$, $(01\br10)$,
$(\br1100)$, $(\br1010)$, $(0\br110)$, $(1\br100)$, т.\,е. форма $\{10\br10\}$, тогда
как в трёх индексах они выглядят непохоже: $(100)$, $(010)$, $(\br110)$\ldots{}
Вычисления обычно ведут в трёх индексах, а третий индекс заменяют точкой:
$(hk.l)$.

Направления тоже записывают четырьмя индексами $[uvtw]$ с условием
$u+v+t=0$. Если в трёх индексах $\vb R=U\vb a_1+V\vb a_2+W\vb c$, то, подставив
$\vb a_3=-(\vb a_1+\vb a_2)$ в $\vb R=u\vb a_1+v\vb a_2+t\vb a_3+w\vb c$, находим
$U=2u+v$, $V=u+2v$, $W=w$, откуда (с точностью до общего множителя)
\begin{keyeq}\label{eq:uvtw}
u=2U-V,\qquad v=2V-U,\qquad t=-(u+v),\qquad w=3W .
\end{keyeq}

\begin{primer}[Переход к четырём индексам]
$[100]\to[2\br1\br10]$; $[110]\to[11\br20]$; $[001]\to[0001]$;
$[\br110]\to[\br1100]$. Направления $\nfam{11\br20}$~— это шесть направлений
вдоль осей $\pm\vb a_1$, $\pm\vb a_2$, $\pm\vb a_3$ (плотноупакованные ряды в
ГПУ-металлах).
\end{primer}

\begin{figure}[ht]
\centering
\begin{tikzpicture}[scale=1.25]
  \foreach \i in {-2,...,2}{\foreach \j in {-2,...,2}{
    \pgfmathsetmacro{\px}{\i+0.5*\j}\pgfmathsetmacro{\py}{0.866*\j}
    \pgfmathparse{abs(\i+\j)<=2 ? 1 : 0}\ifnum\pgfmathresult=1
      \fill[atomA, draw=atomLine, line width=0.3pt] (\px,\py) circle (0.07);\fi}}
  \draw[fmGray, thin] (0:1) \foreach \a in {60,120,...,300}{-- (\a:1)} -- cycle;
  \draw[bvec, fmBlue] (0,0) -- (1,0) node[below right=-2pt, font=\small, black]{$\vb a_1$};
  \draw[bvec, fmBlue] (0,0) -- (120:1) node[above left=-2pt, font=\small, black]{$\vb a_2$};
  \draw[bvec, fmBlue] (0,0) -- (240:1) node[below left=-2pt, font=\small, black]{$\vb a_3$};
  \draw[fmRed, very thick] (1,0) -- (60:1);
  \node[fmRed, font=\scriptsize, anchor=west] at (30:0.95) {$(10\br10)$};
  \draw[fmGreen!70!black, very thick] ($(1,0)!-0.25!(120:1)$) -- ($(1,0)!1.25!(120:1)$);
  \node[fmGreen!70!black, font=\scriptsize, anchor=south east] at ($(1,0)!1.25!(120:1)$) {$(11\br20)$};
  \draw[bvec, fmAmber!90!black] (0,0) -- (60:1.6) node[above right=-2pt, font=\scriptsize]{$[11\br20]$};
  \node[font=\scriptsize, fmBlue, anchor=north] at (1.65,-0.05) {$[2\br1\br10]$};
\end{tikzpicture}
\caption{Гексагональная система: оси $\vb a_1$, $\vb a_2$, $\vb a_3$ в
базисной плоскости, следы плоскостей призмы первого рода $(10\br10)$ и
второго рода $(11\br20)$, направления $[2\br1\br10]$ (вдоль $\vb a_1$) и
$[11\br20]$ (вдоль $\vb a_1+\vb a_2$)}
\label{fig:hex4}
\end{figure}

\subsection{Обратная решётка}\label{ssec:reciprocal}

\begin{opredelenie}[Обратная решётка]
Решётка, базисные векторы которой $\vb b_1$, $\vb b_2$, $\vb b_3$ (их обозначают
также $\vb a^*$, $\vb b^*$, $\vb c^*$) связаны с базисом прямой решётки
условием
\begin{equation}\label{eq:recip-def}
(\vb a_i\vb b_j)=\delta_{ij}\qquad(\delta_{ii}=1,\ \delta_{ij}=0\text{ при }i\ne j).
\end{equation}
\end{opredelenie}
Из определения $\vb b_1\perp\vb a_2,\vb a_3$, т.\,е. $\vb b_1\parallel[\vb a_2\times\vb a_3]$;
нормировка $(\vb a_1\vb b_1)=1$ даёт
\begin{keyeq}\label{eq:b1}
\vb b_1=\frac{[\vb a_2\times\vb a_3]}{V},\quad
\vb b_2=\frac{[\vb a_3\times\vb a_1]}{V},\quad
\vb b_3=\frac{[\vb a_1\times\vb a_2]}{V},\qquad V=(\vb a_1[\vb a_2\times\vb a_3]) .
\end{keyeq}
Решётки взаимны: прямая решётка~— обратная к обратной,
$\vb a_1=[\vb b_2\times\vb b_3]/V^*$, и объёмы ячеек связаны условием $VV^*=1$. В
физике твёрдого тела обычно используют определение
$(\vb a_i\vb b_j)=2\pi\delta_{ij}$ (так в гл.~3)~— все векторы обратной решётки
просто умножаются на $2\pi$.

Метрическая матрица обратной решётки $g^{ij}=(\vb b_i\vb b_j)$ равна $G^{-1}$:
\[
G^{-1}_\text{куб}=\frac1{a^2}I,\qquad
G^{-1}_\text{ромб}=\operatorname{diag}\Bigl(\frac1{a^2},\frac1{b^2},\frac1{c^2}\Bigr),\qquad
G^{-1}_\text{гекс}=\begin{pmatrix}\frac4{3a^2}&\frac2{3a^2}&0\\\frac2{3a^2}&\frac4{3a^2}&0\\0&0&\frac1{c^2}\end{pmatrix}.
\]

\begin{podrobnee}[Матричная запись и взаимность ГЦК--ОЦК]
Если векторы прямого базиса записать строками матрицы $A$, а векторы
обратного~— столбцами матрицы $B$, условие~\eqref{eq:recip-def} принимает
вид $AB=I$, т.\,е. $B=A^{-1}$. Для ГЦК-решётки из~\eqref{eq:prim}
\[
A_\text{ГЦК}^{-1}=\frac1a\begin{pmatrix}1&\br1&1\\1&1&\br1\\\br1&1&1\end{pmatrix},
\]
столбцы которой~— векторы $\frac1a(1,1,\br1)$, $\frac1a(\br1,1,1)$,
$\frac1a(1,\br1,1)$: это примитивные векторы \emph{ОЦК}-решётки с
кубической ячейкой $2/a$. Итак, обратная к ГЦК-решётке с ребром $a$~—
ОЦК-решётка с ребром $2/a$ (объём кубической ячейки $8/a^3$), и наоборот.
Метрические матрицы связаны: $\vb b=G^{-1}\vb a$, $\vb a=G\vb b$, $GG^{-1}=I$.
\end{podrobnee}

\paragraph{Вектор обратной решётки и его свойства}
Узел обратной решётки задаётся вектором
$\vb H_{HKL}=H\vb b_1+K\vb b_2+L\vb b_3$, где $H=qh$, $K=qk$, $L=ql$ (общий
множитель $q$ может быть больше единицы).

\begin{keyeq}\label{eq:H-props}
\vb H_{hkl}=h\vb b_1+k\vb b_2+l\vb b_3\ \perp\ (hkl),\qquad
|\vb H_{hkl}|=\frac1{d_{hkl}} .
\end{keyeq}
\begin{itemize}
\item \emph{Перпендикулярность.} Для любого узлового вектора
$\vb R=m\vb a_1+n\vb a_2+p\vb a_3$, лежащего в плоскости $(hkl)$, проходящей через
начало, $(\vb H\vb R)=hm+kn+lp=0$ по~\eqref{eq:plane-q} с $q=0$.
\item \emph{Модуль.} Ближайшая к началу плоскость семейства отсекает на
оси $x$ отрезок $\vb a_1/h$; его проекция на единичную нормаль
$\vb H/|\vb H|$~— межплоскостное расстояние:
$d_{hkl}=\frac{\vb a_1}{h}\cdot\frac{h\vb b_1+k\vb b_2+l\vb b_3}{|\vb H|}=\frac1{|\vb H|}$.
\end{itemize}
Так \vazhno{каждому семейству плоскостей прямой решётки отвечает узел
обратной}: направление на него~— нормаль к плоскостям, расстояние~—
обратное межплоскостное (рис.~\ref{fig:H-vector}).

\begin{figure}[ht]
\centering
\begin{tikzpicture}[scale=1.15]
  \begin{scope}
    \clip (-1.6,-1.0) rectangle (4.6,3.0);
    \foreach \q in {-2,...,5}{
      \draw[fmBlue!60] ($\q*(1.2,0)+(-3*0.8,3*1.0)$) -- ($\q*(1.2,0)+(4*0.8,-4*1.0)$);}
    \foreach \m in {-2,...,5}{\foreach \n in {-2,...,4}{
      \fill[atomA, draw=atomLine, line width=0.3pt] ($\m*(1.2,0)+\n*(0.4,1.0)$) circle (0.06);}}
  \end{scope}
  \draw[bvec, atomLine] (0,0) -- (1.2,0) node[below, font=\small]{$\vb a$};
  \draw[bvec, atomLine] (0,0) -- (0.4,1.0) node[left, font=\small]{$\vb b$};
  \draw[bvec, fmRed, line width=1.1pt] (0,0) -- ($1.6*(0.781,0.625)$) node[above, font=\small]{$\vb H_{11}$};
  \draw[dimline, fmGreen!60!black] ($(0,0)+0.0*(0.781,0.625)+(0.39,-0.49)$) -- ($0.937*(0.781,0.625)+(0.39,-0.49)$);
  \node[font=\small, fmGreen!60!black] at (1.35,-0.15) {$d_{11}$};
  \fill (0,0) circle (1.4pt);
\end{tikzpicture}
\caption{Вектор обратной решётки $\vb H$ перпендикулярен семейству узловых
плоскостей (здесь~— узловых прямых $(11)$ косоугольной двумерной решётки),
а его длина обратна расстоянию $d$ между ними}
\label{fig:H-vector}
\end{figure}

\paragraph{Закон зон (соотношение Вейсса)}
Условие $hu+kv+lw=0$ означает, что направление $[uvw]$ лежит в плоскости
$(hkl)$, т.\,е. плоскость принадлежит \emph{зоне} с осью $[uvw]$. Отсюда
практическое правило: ось зоны двух плоскостей $(h_1k_1l_1)$ и
$(h_2k_2l_2)$~— векторное произведение их индексов,
$[uvw]=[k_1l_2-l_1k_2,\ l_1h_2-h_1l_2,\ h_1k_2-k_1h_2]$, и аналогично
плоскость, содержащая два направления.

\begin{primer}[Ось зоны]
Плоскости $(110)$ и $(1\br10)$: $[uvw]=[1\cdot0-0\cdot(-1),\ 0\cdot1-1\cdot0,\
1\cdot(-1)-1\cdot1]=[00\br2]\parallel[001]$. Проверка: $1\cdot0+1\cdot0+0\cdot1=0$
и $1\cdot0-1\cdot0+0\cdot1=0$~— обе плоскости содержат $[001]$.
\end{primer}

\paragraph{Точки первого и второго рода}
В дифракционных методах (например, в методе Лауэ) полюсы плоскостей
называют точками первого рода, выходы направлений~— точками второго
рода. В кубических кристаллах они совпадают ($[hkl]\perp(hkl)$). В
тетрагональных совпадают лишь $\{hk0\}$ с $\nfam{hk0}$ и $\{001\}$ с
$\nfam{001}$; в гексагональных~— $\{hki0\}$ с $\nfam{hki0}$ и $\{0001\}$ с
$\nfam{0001}$ (в четырёхиндексной записи).

\paragraph{Узлы обратной решётки реального кристалла}
У идеального монокристалла узел обратной решётки~— точка. У реальных
объектов он размыт (рис.~\ref{fig:recip-blur}):
\begin{itemize}
\item в \emph{мозаичном} монокристалле (блоки слегка разориентированы)~—
по поверхности сферы радиуса $|\vb H|$, т.\,е. в направлениях,
перпендикулярных $\vb H$;
\item в идеальном \emph{поликристалле} со случайной ориентацией зёрен узел
превращается в сферу радиуса $|\vb H|$;
\item при \emph{неоднородной упругой деформации} межплоскостные расстояния
различны, и узел размывается вдоль $\vb H$.
\end{itemize}

\begin{figure}[ht]
\centering
\foreach \typ/\lab in {1/а, 2/б, 3/в}{%
\begin{tikzpicture}[scale=1.0]
  \fill (0,0) circle (1.4pt) node[below left=-1pt, font=\small]{$O$};
  \draw[fmGray, dashed] (0,0) circle (1.6);
  \ifnum\typ=1
    \fill[fmBlue!40, rotate=50] (1.6,0) ellipse (0.08 and 0.35);
  \fi
  \ifnum\typ=2
    \draw[fmBlue!60, line width=2.2pt] (0,0) circle (1.6);
  \fi
  \ifnum\typ=3
    \fill[fmBlue!40, rotate=50] (1.6,0) ellipse (0.35 and 0.08);
  \fi
  \draw[bvec, fmRed] (0,0) -- (50:1.55) node[midway, left, font=\small]{$\vb H$};
  \node[font=\small] at (0,-2.0) {\textit{\lab}};
\end{tikzpicture}\hspace{6mm}}
\caption{Размытие узла обратной решётки: \textit{а}~— мозаичный
монокристалл (вдоль сферы радиуса $|\vb H|$); \textit{б}~— поликристалл
(сфера); \textit{в}~— монокристалл с неоднородной упругой деформацией (вдоль
$\vb H$)}
\label{fig:recip-blur}
\end{figure}

\begin{podrobnee}[Обратное пространство как фурье-пространство]
Множество узлов можно задать решётчатой функцией~— суммой
$\delta$-функций в узлах:
$\text{Ш}(\vb r)=\sum_{m,n,p}\delta(\vb r-\vb R_{mnp})$ в прямом пространстве и
$\text{Ш}(\vb h)=\sum_{h,k,l}\delta(\vb h-\vb H_{hkl})$ в обратном. Эти функции
связаны преобразованием Фурье, поэтому обратное пространство называют
фурье-пространством. При дифракции рентгеновских лучей в прямом
пространстве задана электронная плотность $\rho(\vb r)$, а в обратном~—
амплитуда рассеяния $A(\vb h)$: обратное пространство~— это
\emph{дифракционное} пространство.
\end{podrobnee}

% ----------------------------------------------------------------------
%  1.2.5. Матрица ортогонального преобразования; 1.2.6. Преобразование
%  индексов; 1.2.7. Формулы структурной кристаллографии; 1.2.8. Ячейки
%  Вигнера–Зейтца
% ----------------------------------------------------------------------
\subsection{Матрица ортогонального преобразования}\label{ssec:orth}

Повороты (и вообще поворотные элементы симметрии) описывают
\term{ортогональными матрицами}~— они сохраняют длины векторов и углы
между ними. Для ортогональной матрицы $A=\{a_{ij}\}$
\begin{equation}\label{eq:orth}
\sum_{i=1}^3a_{ik}a_{il}=\delta_{kl},\qquad A^{-1}=\tilde A,\qquad\det A=\pm1
\end{equation}
($\tilde A$~— транспонированная матрица). При $\det A=+1$ это
\emph{собственный} поворот, при $\det A=-1$~— \emph{несобственный}
(поворот с отражением или инверсией).

\paragraph{Ось и угол поворота по матрице}
Собственные значения вещественной ортогональной матрицы $3\times3$
(собственного поворота)~— $1$, $e^{-i\alpha}$, $e^{+i\alpha}$. Собственный вектор
$\vb u$ с собственным значением $1$ не меняется при повороте~— это
\emph{ось}, а фаза $\alpha$~— \emph{угол} поворота. Их находят прямо по
элементам матрицы:
\begin{keyeq}\label{eq:axis-angle}
\begin{gathered}
\operatorname{tr}A=a_{11}+a_{22}+a_{33}=1+2\cos\alpha,\\
u_1:u_2:u_3=(a_{32}-a_{23}):(a_{13}-a_{31}):(a_{21}-a_{12}) .
\end{gathered}
\end{keyeq}
След матрицы (по-английски trace, по-немецки Spur, $\operatorname{Sp}A$)~—
инвариант: он не зависит от системы координат.

\begin{primer}[Поворот вокруг {[111]}]
Матрица циклической перестановки осей $x\to y\to z\to x$:
\[
A=\begin{pmatrix}0&0&1\\1&0&0\\0&1&0\end{pmatrix}.
\]
$\operatorname{tr}A=0=1+2\cos\alpha$, откуда $\cos\alpha=-\frac12$, $\alpha=120^\circ$.
Ось: $(a_{32}-a_{23}):(a_{13}-a_{31}):(a_{21}-a_{12})=(1-0):(1-0):(1-0)$, т.\,е.
$[111]$. Это поворотная ось 3-го порядка куба (п.~1.3.1).
\end{primer}

Обратно, матрица поворота на угол $\varphi$ вокруг оси с направляющими
косинусами $l_1$, $l_2$, $l_3$\ispr{В оригинале (формула (1.60)) в элементе
$R_{12}$ стоит $-l_3\cos\varphi$; правильно $-l_3\sin\varphi$.}:
\begin{equation}\label{eq:R-axis}
R=\begin{pmatrix}
l_1^2(1-c)+c & l_1l_2(1-c)-l_3s & l_1l_3(1-c)+l_2s\\
l_1l_2(1-c)+l_3s & l_2^2(1-c)+c & l_2l_3(1-c)-l_1s\\
l_1l_3(1-c)-l_2s & l_2l_3(1-c)+l_1s & l_3^2(1-c)+c
\end{pmatrix},\quad c=\cos\varphi,\ s=\sin\varphi
\end{equation}

\paragraph{Углы Эйлера}
Если поворот задан углами Эйлера $\varphi$, $\theta$, $\psi$ (п.~\ref{ssec:projections}),
каждый из трёх поворотов в своей системе координат выглядит просто
(обозначим $c_x\equiv\cos x$, $s_x\equiv\sin x$):
\[
R_z(\varphi)=\begin{pmatrix}c_\varphi&-s_\varphi&0\\s_\varphi&c_\varphi&0\\0&0&1\end{pmatrix},\quad
R_{x'}(\theta)=\begin{pmatrix}1&0&0\\0&c_\theta&-s_\theta\\0&s_\theta&c_\theta\end{pmatrix},\quad
R_{z'}(\psi)=\begin{pmatrix}c_\psi&-s_\psi&0\\s_\psi&c_\psi&0\\0&0&1\end{pmatrix}
\]
В неподвижной системе координат матрицы перемножают в порядке,
\emph{обратном} порядку поворотов: $R(\varphi,\theta,\psi)=R_z(\varphi)R_{x'}(\theta)R_{z'}(\psi)$,
\begin{equation}\label{eq:euler}
R=\begin{pmatrix}
c_\varphi c_\psi-c_\theta s_\varphi s_\psi & -c_\varphi s_\psi-c_\theta s_\varphi c_\psi & s_\theta s_\varphi\\
s_\varphi c_\psi+c_\theta c_\varphi s_\psi & -s_\varphi s_\psi+c_\theta c_\varphi c_\psi & -s_\theta c_\varphi\\
s_\theta s_\psi & s_\theta c_\psi & c_\theta
\end{pmatrix}
\end{equation}
($0\le\varphi,\psi<2\pi$, $0\le\theta\le\pi$)\ispr{В оригинале (формула (1.61)) в
элементе $R_{11}$ пропущен множитель $\cos\theta$ при
$\sin\varphi\sin\psi$.}.
Углы Эйлера популярны в физике твёрдого тела (и при описании текстур,
п.~1.6.2); кристаллографа же чаще интересуют индексы оси и угол поворота.

\paragraph{Активная и пассивная трактовки}
Вектор $\vb r$ при повороте переходит в $\vb r'=A\vb r$. Эту запись можно
понимать двояко (рис.~\ref{fig:active-passive}): \emph{активно}~— вектор
повёрнут на угол $\alpha$, а система координат та же; \emph{пассивно}~—
вектор на месте, а система координат повёрнута на $\alpha$ в обратную
сторону. Удобное правило для составления матрицы: её \vazhno{столбцы~—
координаты векторов $[100]$, $[010]$, $[001]$ после поворота}.

\begin{figure}[ht]
\centering
\begin{subfigure}[b]{0.42\textwidth}\centering
\begin{tikzpicture}[scale=1.4]
  \draw[-{Stealth[length=2mm]}] (-0.2,0) -- (1.8,0) node[right, font=\small]{$y$};
  \draw[-{Stealth[length=2mm]}] (0,-0.2) -- (0,1.8) node[above, font=\small]{$z$};
  \draw[bvec, fmBlue] (0,0) -- (20:1.5) node[right, font=\small]{$\vb r$};
  \draw[bvec, fmRed] (0,0) -- (65:1.5) node[above right=-2pt, font=\small]{$\vb r'$};
  \draw[-{Stealth[length=1.5mm]}, fmGray] (20:0.7) arc (20:65:0.7);
  \node[font=\small, fmGray] at (42:0.9) {$\alpha$};
\end{tikzpicture}
\caption{}
\end{subfigure}\hfill
\begin{subfigure}[b]{0.42\textwidth}\centering
\begin{tikzpicture}[scale=1.4]
  \draw[-{Stealth[length=2mm]}] (-0.2,0) -- (1.8,0) node[right, font=\small]{$y$};
  \draw[-{Stealth[length=2mm]}] (0,-0.2) -- (0,1.8) node[above, font=\small]{$z$};
  \draw[-{Stealth[length=2mm]}, fmRed] (0,0) -- (-45:1.6) node[right, font=\small]{$y'$};
  \draw[-{Stealth[length=2mm]}, fmRed] (0,0) -- (45:1.6) node[above right=-2pt, font=\small]{$z'$};
  \draw[bvec, fmBlue] (0,0) -- (20:1.5) node[right, font=\small]{$\vb r$};
  \draw[-{Stealth[length=1.5mm]}, fmGray] (0:0.7) arc (0:-45:0.7);
  \node[font=\small, fmGray] at (-22:0.95) {$\alpha$};
\end{tikzpicture}
\caption{}
\end{subfigure}
\caption{Две трактовки поворота вокруг оси $x$: \textit{а}~— активная
(поворачивается вектор); \textit{б}~— пассивная (поворачиваются оси, в
обратную сторону): относительно новых осей вектор $\vb r$ имеет те же
координаты, что $\vb r'$ относительно старых}
\label{fig:active-passive}
\end{figure}

\paragraph{Смена системы координат}
Пусть в системе $xyz$ операция задана матрицей $A$: $\vb r_2=A\vb r_1$, а
переход к системе $x'y'z'$~— матрицей $C$: $\vb r'=C\vb r$. Тогда в новой
системе та же операция записывается матрицей
\begin{equation}\label{eq:similar}
A'=CAC^{-1}
\end{equation}
(подставьте $\vb r_{1,2}=C^{-1}\vb r'_{1,2}$ в $\vb r_2=A\vb r_1$). Такие матрицы
называют \emph{подобными}; все их инварианты, в том числе след, одинаковы:
$\operatorname{tr}A'=\operatorname{tr}A=1+2\cos\alpha$. Угол поворота от выбора осей,
естественно, не зависит.

\subsection{Преобразование индексов при смене системы координат}\label{ssec:index-transform}

Часто нужно перейти от одной элементарной ячейки к другой (например, от
кубической к примитивной, от гексагональной к ортогексагональной). Пусть
новые базисные векторы выражаются через старые матрицей $A$:
\begin{equation}\label{eq:basis-transform}
\vb a'_i=\sum_jA_{ij}\vb a_j\qquad(\text{кратко }\vb a'=A\vb a).
\end{equation}
Как изменятся индексы направлений и плоскостей?
\begin{itemize}
\item Вектор $\vb R=\sum m_i\vb a_i=\sum m'_i\vb a'_i$ не меняется, поэтому
индексы \emph{направлений} преобразуются «обратно» базису
(\emph{контравариантно}):
\begin{keyeq}\label{eq:m-transform}
\vb m'=(\tilde A)^{-1}\vb m .
\end{keyeq}
\item Скалярное произведение $hm_1+km_2+lm_3$ (номер плоскости в
семействе~\eqref{eq:plane-q}) тоже не меняется, откуда индексы
\emph{плоскостей} (компоненты векторов обратной решётки) преобразуются
так же, как базисные векторы (\emph{ковариантно}):
\begin{keyeq}\label{eq:h-transform}
\vb h'=A\vb h .
\end{keyeq}
\end{itemize}

\begin{primer}[Удвоенная ячейка]
Возьмём новые оси $\vb a'_1=\vb a_1+\vb a_2$, $\vb a'_2=-\vb a_1+\vb a_2$, $\vb a'_3=\vb a_3$
(ячейка, повёрнутая на $45^\circ$ и вдвое большая по площади основания):
$A=\begin{pmatrix}1&1&0\\\br1&1&0\\0&0&1\end{pmatrix}$. Плоскость $(110)$:
$\vb h'=A\vb h=(2,0,0)$, т.\,е. в новых осях это $(100)$ (точнее, каждая
вторая плоскость семейства $(200)$). Направление $[110]$:
$(\tilde A)^{-1}=\frac12\begin{pmatrix}1&1&0\\\br1&1&0\\0&0&2\end{pmatrix}$,
$\vb m'=(1,0,0)$~— это $[100]$, как и должно быть: $\vb a_1+\vb a_2=\vb a'_1$.
\end{primer}

\subsection{Основные формулы структурной кристаллографии}\label{ssec:formulas}

С метрической матрицей $G$ (п.~\ref{ssec:lattice}) все расчёты в любой
сингонии сводятся к нескольким формулам. Обозначим $\vb m=(m_1,m_2,m_3)^T$~—
индексы направления, $\vb h=(h,k,l)^T$~— индексы плоскости.

\begin{itemize}
\item \emph{Период идентичности}~— расстояние между ближайшими узлами
вдоль $[m_1m_2m_3]$:
\begin{equation}\label{eq:period}
I_{[m]}=\sqrt{\tilde{\vb m}G\vb m}\qquad\bigl(\text{куб: }I=a\sqrt{m_1^2+m_2^2+m_3^2}\bigr).
\end{equation}
\item \emph{Межплоскостное расстояние}: $1/d^2_{hkl}=|\vb H|^2=\tilde{\vb h}G^{-1}\vb h$.
Это выражение~— \term{квадратичная форма} сингонии:
\begin{keyeq}\label{eq:quadform}
\begin{gathered}
\frac1{d^2}=\frac{h^2+k^2+l^2}{a^2}\ \ \text{(куб)},\qquad
\frac1{d^2}=\frac{h^2}{a^2}+\frac{k^2}{b^2}+\frac{l^2}{c^2}\ \ \text{(ромб)},\\
\frac1{d^2}=\frac43\,\frac{h^2+hk+k^2}{a^2}+\frac{l^2}{c^2}\ \ \text{(гекс)}.
\end{gathered}
\end{keyeq}
Тетрагональная~— как ромбическая с $b=a$.
\item \emph{Угол между направлениями}
\begin{equation}\label{eq:angle-dirs}
\cos\varphi=\frac{\tilde{\vb m}^{(2)}G\vb m^{(1)}}{\sqrt{\tilde{\vb m}^{(1)}G\vb m^{(1)}}\sqrt{\tilde{\vb m}^{(2)}G\vb m^{(2)}}} ,
\end{equation}
\emph{угол между плоскостями}~— то же с $G^{-1}$ и индексами плоскостей;
\emph{угол между направлением и нормалью к плоскости}:
$\cos\varphi=\tilde{\vb h}\vb m\big/\bigl(\sqrt{\tilde{\vb h}G^{-1}\vb h}\sqrt{\tilde{\vb m}G\vb m}\bigr)$.
В кубической сингонии все три формулы сводятся к одной:
\begin{keyeq}\label{eq:cos-cubic}
\cos\varphi=\frac{h_1h_2+k_1k_2+l_1l_2}{\sqrt{h_1^2+k_1^2+l_1^2}\sqrt{h_2^2+k_2^2+l_2^2}} .
\end{keyeq}
\item \emph{Объём ячейки}: $V=(\vb a_1[\vb a_2\times\vb a_3])=\sqrt{\det G}$; для
обратной решётки $V^*=\sqrt{\det G^{-1}}=1/V$.
\item \emph{Переход от прямого базиса к обратному}: один и тот же вектор
имеет в прямом базисе компоненты $\vb m$, в обратном~— $\vb h=G\vb m$
(обратно $\vb m=G^{-1}\vb h$). В кубической сингонии $G=a^2I$, и направление
$[uvw]$ параллельно нормали к плоскости $(uvw)$.
\end{itemize}

\begin{primer}[Расчёты для меди и магния]
\textit{Медь} (ГЦК, $a=0{,}3615$~нм): $d_{111}=a/\sqrt3=0{,}2087$~нм,
$d_{200}=a/2=0{,}1808$~нм; угол между $(111)$ и $(100)$:
$\cos\varphi=1/\sqrt3$, $\varphi=54{,}74^\circ$; между $(111)$ и $(1\br11)$:
$\cos\varphi=1/3$, $\varphi=70{,}53^\circ$. Период вдоль $[110]$:
$a\sqrt2=0{,}511$~нм, но кратчайшее межатомное расстояние $a\sqrt2/2=0{,}256$~нм,
так как ГЦК-решётка центрирована.\\
\textit{Магний} (ГПУ, $a=0{,}321$, $c=0{,}521$~нм): $d_{10\br10}=a\sqrt3/2=0{,}278$~нм,
$d_{0002}=c/2=0{,}261$~нм.
\end{primer}

\subsection{Области Вороного, ячейки Вигнера--Зейтца, зоны Бриллюэна}\label{ssec:wigner}

Элементарную ячейку можно выбрать и иначе. Соединим нулевой узел с
ближайшими узлами и проведём через середины этих отрезков
перпендикулярные им плоскости. Они вырежут многогранник, параллельными
переносами которого заполняется всё пространство. Он содержит ровно
один узел (в центре), т.\,е. является примитивной ячейкой, но в отличие от
параллелепипеда обладает \emph{полной симметрией решётки}. В
кристаллографии это \term{область Вороного}, в физике твёрдого
тела~— \term{ячейка Вигнера--Зейтца}.

\begin{itemize}
\item \emph{ГЦК}: 12 ближайших соседей вдоль $\nfam{110}$; ячейка огранена 12
ромбами $\{110\}$~— \emph{ромбический додекаэдр} (рис.~\ref{fig:wigner}\,\textit{а}).
\item \emph{ОЦК}: 8 ближайших соседей вдоль $\nfam{111}$ и 6 следующих вдоль
$\nfam{100}$ (они лишь на 15\,\% дальше); ячейка огранена 8 правильными
шестиугольниками $\{111\}$ и 6 квадратами $\{100\}$~— 14-гранник,
\emph{усечённый октаэдр} (рис.~\ref{fig:wigner}\,\textit{б}).
\end{itemize}

Ячейка Вигнера--Зейтца, построенная в \emph{обратной} решётке,~—
\term{первая зона Бриллюэна} (гл.~3, п.~3.2.2). Поскольку ГЦК- и
ОЦК-решётки взаимно обратны, зона Бриллюэна ГЦК-кристалла~— усечённый
октаэдр, а ОЦК-кристалла~— ромбический додекаэдр.

\begin{figure}[ht]
\centering
\tikzset{bzfront/.style={atomLine, line width=0.7pt, line join=round},
  bzback/.style={atomLine!55, line width=0.4pt, dashed}}
\begin{subfigure}[b]{0.4\textwidth}\centering
\begin{tikzpicture}[scale=1.6]
  % сгенерировано tools/gen_bands.py — не править вручную
\fill[fmBlue!18] (-0.676,-0.539) -- (-0.469,-0.273) -- (-0.676,0.412) -- (-0.883,0.145) -- cycle;
\fill[fmBlue!14] (0.207,0.267) -- (-0.469,-0.273) -- (0.207,-0.684) -- (0.883,-0.145) -- cycle;
\fill[fmBlue!12] (0.207,-0.684) -- (-0.469,-0.273) -- (-0.676,-0.539) -- (0.000,-0.951) -- cycle;
\fill[fmBlue!20] (-0.676,0.412) -- (-0.469,-0.273) -- (0.207,0.267) -- (0.000,0.951) -- cycle;
\fill[fmBlue!8] (0.676,-0.412) -- (0.883,-0.145) -- (0.207,-0.684) -- (0.000,-0.951) -- cycle;
\fill[fmBlue!16] (0.207,0.267) -- (0.883,-0.145) -- (0.676,0.539) -- (0.000,0.951) -- cycle;
\draw[bzback] (0.469,0.273) -- (-0.207,0.684);
\draw[bzback] (0.469,0.273) -- (-0.207,-0.267);
\draw[bzback] (-0.883,0.145) -- (-0.207,-0.267);
\draw[bzback] (-0.883,0.145) -- (-0.207,0.684);
\draw[bzback] (0.469,0.273) -- (0.676,-0.412);
\draw[bzback] (0.469,0.273) -- (0.676,0.539);
\draw[bzfront] (0.883,-0.145) -- (0.676,0.539);
\draw[bzfront] (0.883,-0.145) -- (0.676,-0.412);
\draw[bzfront] (-0.469,-0.273) -- (-0.676,-0.539);
\draw[bzfront] (-0.469,-0.273) -- (-0.676,0.412);
\draw[bzfront] (-0.883,0.145) -- (-0.676,0.412);
\draw[bzfront] (-0.883,0.145) -- (-0.676,-0.539);
\draw[bzfront] (-0.469,-0.273) -- (0.207,0.267);
\draw[bzfront] (-0.469,-0.273) -- (0.207,-0.684);
\draw[bzfront] (0.883,-0.145) -- (0.207,-0.684);
\draw[bzfront] (0.883,-0.145) -- (0.207,0.267);
\draw[bzfront] (0.000,-0.951) -- (0.676,-0.412);
\draw[bzback] (0.000,-0.951) -- (-0.207,-0.267);
\draw[bzback] (0.000,0.951) -- (-0.207,0.684);
\draw[bzfront] (0.000,0.951) -- (0.676,0.539);
\draw[bzfront] (0.000,-0.951) -- (-0.676,-0.539);
\draw[bzfront] (0.000,-0.951) -- (0.207,-0.684);
\draw[bzfront] (0.000,0.951) -- (0.207,0.267);
\draw[bzfront] (0.000,0.951) -- (-0.676,0.412);
\coordinate (bzG) at (0.000,0.000);
\coordinate (bzH) at (-0.469,-0.273);
\coordinate (bzN) at (0.207,-0.209);
\coordinate (bzP) at (0.207,0.267);

\end{tikzpicture}
\caption{ГЦК}
\end{subfigure}\hspace{8mm}
\begin{subfigure}[b]{0.4\textwidth}\centering
\begin{tikzpicture}[scale=1.6]
  % сгенерировано tools/gen_bands.py — не править вручную
\fill[fmBlue!17] (-0.469,0.203) -- (-0.911,-0.200) -- (-0.469,-0.748) -- (-0.028,-0.345) -- cycle;
\fill[fmBlue!11] (0.883,-0.621) -- (1.118,-0.009) -- (0.883,0.330) -- (0.648,-0.281) -- cycle;
\fill[fmBlue!20] (0.441,0.879) -- (0.235,1.087) -- (-0.441,1.024) -- (-0.235,0.815) -- cycle;
\fill[fmBlue!15] (-0.911,-0.200) -- (-1.118,0.009) -- (-0.883,-0.330) -- (-0.441,-0.879) -- (-0.235,-1.087) -- (-0.469,-0.748) -- cycle;
\fill[fmBlue!21] (-0.235,0.815) -- (-0.441,1.024) -- (-0.883,0.621) -- (-1.118,0.009) -- (-0.911,-0.200) -- (-0.469,0.203) -- cycle;
\fill[fmBlue!10] (-0.469,-0.748) -- (-0.235,-1.087) -- (0.441,-1.024) -- (0.883,-0.621) -- (0.648,-0.281) -- (-0.028,-0.345) -- cycle;
\fill[fmBlue!17] (0.648,-0.281) -- (0.883,0.330) -- (0.441,0.879) -- (-0.235,0.815) -- (-0.469,0.203) -- (-0.028,-0.345) -- cycle;
\draw[bzfront] (-0.911,-0.200) -- (-0.469,0.203);
\draw[bzfront] (-0.911,-0.200) -- (-0.469,-0.748);
\draw[bzfront] (-0.469,-0.748) -- (-0.028,-0.345);
\draw[bzfront] (-0.469,0.203) -- (-0.028,-0.345);
\draw[bzback] (0.028,0.345) -- (0.469,-0.203);
\draw[bzback] (0.028,0.345) -- (0.469,0.748);
\draw[bzback] (0.469,0.748) -- (0.911,0.200);
\draw[bzback] (0.469,-0.203) -- (0.911,0.200);
\draw[bzfront] (1.118,-0.009) -- (0.883,-0.621);
\draw[bzfront] (1.118,-0.009) -- (0.883,0.330);
\draw[bzfront] (0.883,0.330) -- (0.648,-0.281);
\draw[bzfront] (0.883,-0.621) -- (0.648,-0.281);
\draw[bzback] (-0.648,0.281) -- (-0.883,0.621);
\draw[bzback] (-0.648,0.281) -- (-0.883,-0.330);
\draw[bzfront] (-0.883,-0.330) -- (-1.118,0.009);
\draw[bzfront] (-0.883,0.621) -- (-1.118,0.009);
\draw[bzfront] (0.235,1.087) -- (0.441,0.879);
\draw[bzfront] (0.235,1.087) -- (-0.441,1.024);
\draw[bzfront] (-0.441,1.024) -- (-0.235,0.815);
\draw[bzfront] (0.441,0.879) -- (-0.235,0.815);
\draw[bzback] (0.235,-0.815) -- (-0.441,-0.879);
\draw[bzback] (0.235,-0.815) -- (0.441,-1.024);
\draw[bzfront] (0.441,-1.024) -- (-0.235,-1.087);
\draw[bzfront] (-0.441,-0.879) -- (-0.235,-1.087);
\draw[bzback] (0.028,0.345) -- (-0.648,0.281);
\draw[bzback] (0.469,-0.203) -- (0.235,-0.815);
\draw[bzfront] (-0.883,-0.330) -- (-0.441,-0.879);
\draw[bzback] (0.469,0.748) -- (0.235,1.087);
\draw[bzfront] (-0.883,0.621) -- (-0.441,1.024);
\draw[bzback] (0.911,0.200) -- (1.118,-0.009);
\draw[bzfront] (0.441,-1.024) -- (0.883,-0.621);
\draw[bzfront] (0.441,0.879) -- (0.883,0.330);
\draw[bzfront] (-1.118,0.009) -- (-0.911,-0.200);
\draw[bzfront] (-0.235,-1.087) -- (-0.469,-0.748);
\draw[bzfront] (-0.235,0.815) -- (-0.469,0.203);
\draw[bzfront] (0.648,-0.281) -- (-0.028,-0.345);
\coordinate (bzG) at (0.000,0.000);
\coordinate (bzX) at (-0.469,-0.273);
\coordinate (bzL) at (0.207,0.267);
\coordinate (bzW) at (-0.028,-0.345);

\end{tikzpicture}
\caption{ОЦК}
\end{subfigure}
\caption{Ячейки Вигнера--Зейтца (расчёт; ось $z$~— вертикально):
\textit{а}~— ромбический додекаэдр ГЦК-решётки; \textit{б}~— усечённый
октаэдр ОЦК-решётки. Те же многогранники~— зоны Бриллюэна ОЦК- и
ГЦК-решёток соответственно}
\label{fig:wigner}
\end{figure}

\begin{itogi}
\begin{itemize}
\item Стереографическая проекция: направление~— точка ($\rho=r\operatorname{tg}\frac\theta2$),
плоскость~— дуга большого круга; гномостереографическая: плоскость~—
полюс. Углы измеряют по сеткам Вульфа и Болдырева.
\item Решётка: $\vb R=m_1\vb a_1+m_2\vb a_2+m_3\vb a_3$; 7 сингоний; метрическая
матрица $G=\{(\vb a_i\vb a_j)\}$.
\item Символы: узел $[\![mnp]\!]$, направление $[uvw]$, $\nfam{uvw}$; плоскость
$(hkl)$~— обратные отрезки; форма $\{hkl\}$; в гексагональной системе
$(hkil)$, $i=-(h+k)$.
\item Обратная решётка: $(\vb a_i\vb b_j)=\delta_{ij}$; $\vb H_{hkl}\perp(hkl)$,
$|\vb H_{hkl}|=1/d_{hkl}$; закон зон $hu+kv+lw=0$.
\item Поворот: $\operatorname{tr}A=1+2\cos\alpha$; индексы направлений
преобразуются контравариантно, плоскостей~— ковариантно.
\item $1/d^2=\tilde{\vb h}G^{-1}\vb h$; в кубе $d=a/\sqrt{h^2+k^2+l^2}$.
\item Ячейка Вигнера--Зейтца ГЦК~— ромбический додекаэдр, ОЦК~— усечённый
октаэдр.
\end{itemize}
\end{itogi}

\begin{proverka}
\item Чем стереографическая проекция плоскости отличается от
гномостереографической? Почему в кубическом кристалле полюс $(hkl)$
совпадает с выходом $[hkl]$?
\item Найдите расстояние от центра проекции $(001)$ до полюса $(011)$ при
$r=10$~см.
\item Определите индексы плоскости, отсекающей на осях отрезки $2a$,
$3b$, $-c$.
\item Переведите в четырёхиндексную запись направления $[210]$ и $[1\br10]$
гексагонального кристалла.
\item Найдите ось зоны плоскостей $(111)$ и $(1\br10)$.
\item Каков угол поворота, если след матрицы поворота равен $-1$?
\item Найдите межплоскостные расстояния $d_{110}$ и $d_{111}$ для
$\alpha$-железа ($a=0{,}2866$~нм).
\item Почему ячейку Вигнера--Зейтца ОЦК-решётки ограничивают не только
плоскости $\{111\}$, но и $\{100\}$?
\end{proverka}


\section{Симметрия кристаллов}\label{sec:symmetry}
% ----------------------------------------------------------------------
%  1.3. Симметрия кристаллов: 1.3.1–1.3.4 (оси симметрии, элементы теории
%  групп)
% ----------------------------------------------------------------------
\begin{orientir}
Какие повороты и отражения допускает кристалл и почему их так мало
(1.3.1--1.3.3); как их комбинации образуют группы (1.3.4); 32 точечные
группы и 14 решёток Браве (1.3.5); 230 пространственных групп с винтовыми
осями и плоскостями скольжения (1.3.6); предельные группы Кюри, нужные
для описания симметрии физических свойств (1.3.7).
\end{orientir}

Симметрия кристаллического пространства задаётся всеми
преобразованиями, которые сохраняют расстояния между точками и
совмещают пространство с самим собой. Элементы симметрии бывают:
\begin{itemize}
\item \term{открытые}~— содержат трансляцию (переносы, винтовые оси,
плоскости скольжения); они описывают симметрию бесконечного
пространства;
\item \term{закрытые}~— оставляют неподвижной хотя бы одну точку и после
конечного числа повторений возвращают пространство в исходное
положение. Они задаются ортогональными матрицами $R$ с $\det R=\pm1$ и
сводятся к \emph{поворотным осям} (собственные повороты, $\det R=+1$) и
\emph{инверсионным осям} (поворот с последующей инверсией в точке на оси,
$\det R=-1$).
\end{itemize}

\subsection{Поворотные оси симметрии}\label{ssec:rot-axes}

\begin{opredelenie}[Ось симметрии]
Прямая, при повороте вокруг которой на угол $\alpha_n=360^\circ/n$
кристаллическое пространство совмещается с собой. Целое $n$~—
\term{порядок оси}, $\alpha_n$~— её элементарный угол: через $n$ поворотов
пространство возвращается в исходное положение.
\end{opredelenie}

\paragraph{Почему возможны только оси 1, 2, 3, 4 и 6}
При повороте узел решётки $\vb x$ переходит в узел $\vb x'=R\vb x$. В
кристаллографической системе координат координаты узлов~— целые числа,
значит, и элементы матрицы $R$, и её след~— целые: $\operatorname{tr}R=N$. След не
зависит от системы координат (п.~\ref{ssec:orth}); в ортогональной системе
с осью $x$ вдоль оси поворота
\begin{equation}\label{eq:Rx}
R'=\begin{pmatrix}1&0&0\\0&\cos\alpha&-\sin\alpha\\0&\sin\alpha&\cos\alpha\end{pmatrix},\qquad
\operatorname{tr}R'=1+2\cos\alpha=N .
\end{equation}
Так как $|\cos\alpha|\le1$, $N$ может быть лишь $-1$, $0$, $1$, $2$, $3$
(табл.~\ref{tab:axes}).

\begin{table}[ht]
\centering
\caption{Возможные поворотные оси}\label{tab:axes}
\small
\begin{tabular}{@{}ccccc@{}}
\toprule
$\operatorname{tr}R$ & $-1$ & $0$ & $+1$ & $+2$ \ \ ($+3$) \\
\midrule
$\cos\alpha$ & $-1$ & $-1/2$ & $0$ & $1/2$ \ \ ($1$) \\
$\alpha$ & $180^\circ$ & $120^\circ$ & $90^\circ$ & $60^\circ$ \ \ ($0^\circ$) \\
Порядок оси $n$ & 2 & 3 & 4 & 6 \ \ (1) \\
\bottomrule
\end{tabular}
\end{table}

\begin{glavnoe}[Кристаллографическое ограничение]
В кристалле возможны только оси симметрии 1, 2, 3, 4 и 6-го порядков.
Осей 5-го, 7-го и более высоких порядков нет~— их запрещает
трансляционная симметрия (пятиугольниками нельзя без зазоров замостить
плоскость). Оси 5-го порядка появляются лишь в четырёхмерной
кристаллографии и у квазикристаллов (п.~\ref{ssec:space-groups}).
\end{glavnoe}

Действие осей удобно изображать на стереографической проекции: одно
направление общего положения и все направления, в которые его
переводят повороты (рис.~\ref{fig:axes-proj}, верхний ряд). В кубическом
кристалле оси 4-го порядка идут вдоль $\nfam{100}$, 3-го~— вдоль
$\nfam{111}$, 2-го~— вдоль $\nfam{110}$ (рис.~\ref{fig:cube-axes}).

\begin{figure}[ht]
\centering
\begin{tikzpicture}[scale=0.82]
  % создано tools/gen_stereo.py: действие осей симметрии
\draw[atomLine, line width=0.6pt] (0.0,0) circle (1);
\fill[fmBlue] (0.583,0.212) circle (0.075);
\node[font=\small] at (0.0,-1.32) {$1$};
\draw[atomLine, line width=0.6pt] (2.6,0) circle (1);
\fill[fmBlue] (3.183,0.212) circle (0.075);
\fill[fmBlue] (2.017,-0.212) circle (0.075);
\node[font=\small] at (2.6,-1.32) {$2$};
\draw[atomLine, line width=0.6pt] (5.2,0) circle (1);
\fill[fmBlue] (5.783,0.212) circle (0.075);
\fill[fmBlue] (4.725,0.399) circle (0.075);
\fill[fmBlue] (5.092,-0.611) circle (0.075);
\node[font=\small] at (5.2,-1.32) {$3$};
\draw[atomLine, line width=0.6pt] (7.800000000000001,0) circle (1);
\fill[fmBlue] (8.383,0.212) circle (0.075);
\fill[fmBlue] (7.588,0.583) circle (0.075);
\fill[fmBlue] (7.217,-0.212) circle (0.075);
\fill[fmBlue] (8.012,-0.583) circle (0.075);
\node[font=\small] at (7.800000000000001,-1.32) {$4$};
\draw[atomLine, line width=0.6pt] (10.4,0) circle (1);
\fill[fmBlue] (10.983,0.212) circle (0.075);
\fill[fmBlue] (10.508,0.611) circle (0.075);
\fill[fmBlue] (9.925,0.399) circle (0.075);
\fill[fmBlue] (9.817,-0.212) circle (0.075);
\fill[fmBlue] (10.292,-0.611) circle (0.075);
\fill[fmBlue] (10.875,-0.399) circle (0.075);
\node[font=\small] at (10.4,-1.32) {$6$};
\draw[atomLine, line width=0.6pt] (0.0,-2.75) circle (1);
\fill[fmBlue] (0.583,-2.538) circle (0.075);
\node[fmRed, font=\small] at (-0.583,-2.962) {$\times$};
\node[font=\small] at (0.0,-4.07) {$\br{1}$};
\draw[atomLine, line width=0.6pt] (2.6,-2.75) circle (1);
\node[fmRed, font=\small] at (3.183,-2.538) {$\times$};
\fill[fmBlue] (3.183,-2.538) circle (0.075);
\node[font=\small] at (2.6,-4.07) {$\br2=m$};
\draw[atomLine, line width=0.6pt] (5.2,-2.75) circle (1);
\fill[fmBlue] (5.783,-2.538) circle (0.075);
\node[fmRed, font=\small] at (5.308,-2.139) {$\times$};
\fill[fmBlue] (4.725,-2.351) circle (0.075);
\node[fmRed, font=\small] at (4.617,-2.962) {$\times$};
\fill[fmBlue] (5.092,-3.361) circle (0.075);
\node[fmRed, font=\small] at (5.675,-3.149) {$\times$};
\node[font=\small] at (5.2,-4.07) {$\br{3}$};
\draw[atomLine, line width=0.6pt] (7.800000000000001,-2.75) circle (1);
\fill[fmBlue] (8.383,-2.538) circle (0.075);
\node[fmRed, font=\small] at (7.588,-2.167) {$\times$};
\fill[fmBlue] (7.217,-2.962) circle (0.075);
\node[fmRed, font=\small] at (8.012,-3.333) {$\times$};
\node[font=\small] at (7.800000000000001,-4.07) {$\br{4}$};
\draw[atomLine, line width=0.6pt] (10.4,-2.75) circle (1);
\node[fmRed, font=\small] at (10.983,-2.538) {$\times$};
\fill[fmBlue] (10.983,-2.538) circle (0.075);
\node[fmRed, font=\small] at (9.925,-2.351) {$\times$};
\fill[fmBlue] (9.925,-2.351) circle (0.075);
\node[fmRed, font=\small] at (10.292,-3.361) {$\times$};
\fill[fmBlue] (10.292,-3.361) circle (0.075);
\node[font=\small] at (10.4,-4.07) {$\br{6}$};

\end{tikzpicture}
\caption{Действие осей симметрии, перпендикулярных плоскости проекции:
верхний ряд~— поворотные оси 1, 2, 3, 4, 6; нижний~— инверсионные оси
$\br1$, $\br2=m$, $\br3$, $\br4$, $\br6$. Точка~— направление, выходящее в
верхнюю полусферу, крестик~— в нижнюю (расчёт)}
\label{fig:axes-proj}
\end{figure}

\begin{figure}[ht]
\centering
\begin{tikzpicture}[cubeproj, scale=2.3]
  \CubeEdges
  \draw[fmRed, very thick] (0.5,0.5,-0.35) -- (0.5,0.5,1.35);
  \node[draw=fmRed, fill=fmRed!25, minimum size=2.4mm, inner sep=0] at (0.5,0.5,1.35) {};
  \node[font=\scriptsize, fmRed, anchor=west] at (0.55,0.5,1.4) {4};
  \draw[fmBlue, very thick] (-0.25,-0.25,-0.25) -- (1.25,1.25,1.25);
  \node[regular polygon, regular polygon sides=3, draw=fmBlue, fill=fmBlue!25, minimum size=3.2mm, inner sep=0] at (1.25,1.25,1.25) {};
  \node[font=\scriptsize, fmBlue, anchor=west] at (1.32,1.25,1.25) {3};
  \draw[fmGreen!60!black, very thick] (-0.25,0.5,-0.25) -- (1.25,0.5,1.25);
  \node[ellipse, draw=fmGreen!60!black, fill=fmGreen!25, minimum width=3.2mm, minimum height=1.8mm, inner sep=0] at (1.25,0.5,1.25) {};
  \node[font=\scriptsize, fmGreen!60!black, anchor=west] at (1.32,0.5,1.25) {2};
\end{tikzpicture}
\caption{Оси симметрии куба: три оси 4 (вдоль $\nfam{100}$, показана одна),
четыре оси 3 (вдоль $\nfam{111}$) и шесть осей 2 (вдоль $\nfam{110}$)}
\label{fig:cube-axes}
\end{figure}

\paragraph{Матрицы поворотов}
Правило из п.~\ref{ssec:orth}: столбцы матрицы~— координаты векторов $[100]$,
$[010]$, $[001]$ после поворота. Их удобно находить по стандартной
проекции (рис.~\ref{fig:cubic-dirs}\,\textit{б}).
\begin{itemize}
\item Поворот на $90^\circ$ вокруг $[100]$ (против часовой стрелки, если
смотреть с конца оси): $[100]\to[100]$, $[010]\to[001]$, $[001]\to[0\br10]$:
\[
C_4^{[100]}=\begin{pmatrix}1&0&0\\0&0&\br1\\0&1&0\end{pmatrix},\quad
\operatorname{tr}=1\ \Rightarrow\ \alpha=90^\circ .
\]
\item Поворот на $120^\circ$ вокруг $[1\br11]$: $[100]\to[001]$,
$[010]\to[\br100]$, $[001]\to[0\br10]$:
\[
C_3^{[1\br11]}=\begin{pmatrix}0&\br1&0\\0&0&\br1\\1&0&0\end{pmatrix},\quad
\operatorname{tr}=0\ \Rightarrow\ \alpha=120^\circ ;
\]
по~\eqref{eq:axis-angle} ось $(0+1):(0-1):(0+1)=[1\br11]$.
\item Поворот на $180^\circ$ вокруг $[011]$: $[100]\perp[011]$ переходит в
$[\br100]$, а $[010]$ и $[001]$ меняются местами:
\[
C_2^{[011]}=\begin{pmatrix}\br1&0&0\\0&0&1\\0&1&0\end{pmatrix},\quad
\operatorname{tr}=-1\ \Rightarrow\ \alpha=180^\circ .
\]
\end{itemize}

\begin{vnimanie}[Ось поворота на $180^\circ$]
Для $\alpha=180^\circ$ матрица симметрична, и формула~\eqref{eq:axis-angle}
даёт $0:0:0$. В этом случае ось находят как любой ненулевой столбец
матрицы $A+I$: для $C_2^{[011]}$ это $(0,1,1)$.
\end{vnimanie}

\subsection{Инверсионные оси}\label{ssec:inv-axes}

Поворот можно совместить с отражением в точке на оси (инверсией).
Получаются \term{инверсионные оси} тех же порядков: $\br1$, $\br2$, $\br3$, $\br4$,
$\br6$. Инверсия в начале координат (центр инверсии)
\begin{equation}\label{eq:inversion}
\bar1=\begin{pmatrix}\br1&0&0\\0&\br1&0\\0&0&\br1\end{pmatrix},\qquad
R_i=\bar1R=\begin{pmatrix}\br1&0&0\\0&-\cos\alpha&\sin\alpha\\0&-\sin\alpha&-\cos\alpha\end{pmatrix},
\end{equation}
$\det R_i=-1$, $\operatorname{tr}R_i=-1-2\cos\alpha$. Разберём оси по очереди
(рис.~\ref{fig:axes-proj}, нижний ряд).
\begin{itemize}
\item $\br1$ ($\alpha=0$)~— просто \emph{центр симметрии} (инверсии).
\item $\br2$ вдоль $[100]$: $C_2\bar1=\operatorname{diag}(\br1,1,1)$~— зеркальное
отражение в плоскости $yz$. Значит, \vazhno{ось $\br2$ эквивалентна
перпендикулярной ей плоскости симметрии} $m$ (от англ.\ mirror).
\item $\br3=3+\br1$: ось 3 и центр симметрии.
\item $\br6=3/m$: ось 3 и перпендикулярная ей плоскость.
\item $\br4$~— самостоятельный элемент; он содержит ось 2 (дважды
применённый $\br4$~— поворот на $180^\circ$).
\end{itemize}
Независимы, таким образом, $\br1$, $\br2=m$ и $\br4$; оси $\br3$ и $\br6$~— составные.

\subsection{Зеркально-поворотные оси. Полярность}\label{ssec:mirror-axes}

Несобственный поворот можно получить и иначе~— сочетанием поворота с
отражением в перпендикулярной плоскости. Это \term{зеркально-поворотные}
оси $\tilde1$, $\tilde2$, $\tilde3$, $\tilde4$, $\tilde6$ с матрицей
$R_m=mR$, $\det R_m=-1$, $\operatorname{tr}R_m=-1+2\cos\alpha$. Новых элементов они не
дают:
\[
\tilde1=\br2=m,\quad\tilde2=\br1,\quad\tilde3=\br6,\quad\tilde4=\br4,\quad\tilde6=\br3 .
\]
Инверсионные оси приняты в международной символике, зеркально-поворотные~—
в символике Шёнфлиса ($S_n$) и Шубникова.

\paragraph{Полярность}
Вдоль прямой можно двигаться в две стороны. Если оба направления
равноценны (их совмещает какая-либо операция симметрии~— плоскость,
перпендикулярная прямой, центр симметрии на ней, ось 2,
перпендикулярная ей), прямая \emph{неполярна}; иначе~— \term{полярна}.
Инверсионные оси всегда неполярны; поворотная ось полярна, если её не
пересекают элементы $\br1$, $m$, $2$. Полярности направлений отвечает
полярность физических свойств (пироэлектричество, пьезоэлектричество,
п.~\ref{ssec:point-groups}).

\subsection{Элементы теории групп}\label{ssec:groups}

Комбинации элементов симметрии удобно изучать с помощью теории групп.

\begin{opredelenie}[Группа]
Множество элементов $G=\{g_1,g_2,\ldots\}$ с операцией «умножения», для
которой:
\begin{enumerate}
\item произведение любых двух элементов~— элемент группы: $g_ig_j=g_k$;
\item умножение ассоциативно: $g_i(g_jg_l)=(g_ig_j)g_l$;
\item есть единственный единичный элемент $e$: $eg_i=g_ie=g_i$;
\item у каждого элемента есть обратный: $g_ig_i^{-1}=e$.
\end{enumerate}
\end{opredelenie}

«Умножение» не обязательно арифметическое: в группах симметрии это
последовательное выполнение операций ($g_ig_j$~— сначала $g_j$, потом $g_i$).
Коммутативность не требуется; группы, в которых $g_ig_j=g_jg_i$,
называют \term{абелевыми}.

\paragraph{Основные понятия}
\begin{itemize}
\item \term{Порядок группы}~— число её элементов (конечное или
бесконечное).
\item \term{Подгруппа}~— часть элементов, сама образующая группу; содержит
$e$. По \emph{теореме Лагранжа} порядок подгруппы~— делитель порядка
группы; поэтому группа простого порядка не имеет нетривиальных подгрупп.
\item \term{Порядок элемента} $g$~— наименьшее $r$, при котором $g^r=e$;
степени $g,g^2,\ldots,g^r=e$ образуют \emph{период} элемента. Если вся группа
исчерпывается степенями одного элемента, она \term{циклическая}, а этот
элемент~— её \emph{генератор}; иначе выделяют несколько генераторов.
\item Группы \term{изоморфны}, если между их элементами есть взаимно
однозначное соответствие, сохраняющее умножение,~— с абстрактной точки
зрения это одна группа с разными «именами» элементов. При соответствии
без взаимной однозначности (одному элементу $H$ отвечает несколько
элементов $G$) группы \term{гомоморфны}.
\item Элемент $g_k$ \term{сопряжён} элементу $g_i$, если $g_k=g_j^{-1}g_ig_j$ для
какого-то $g_j$ из группы. Взаимно сопряжённые элементы образуют
\term{класс}; единица всегда образует класс сама по себе. Геометрически
сопряжённые операции~— «одинаковые» операции, переводимые друг в друга
симметрией группы.
\end{itemize}

\begin{primer*}[Примеры групп]
\begin{enumerate}
\item Целые числа по сложению: единица~— $0$, обратный к $n$~— $-n$;
бесконечная абелева группа.
\item Рациональные числа, \emph{кроме нуля}, по умножению: единица~— $1$,
обратный~— $1/n$; бесконечная абелева группа\isprmark.
\item Невырожденные матрицы $n\times n$ по умножению; бесконечная
неабелева группа.
\item Повороты тела вокруг всевозможных осей: единица~— отсутствие
поворота, обратный~— поворот в обратную сторону; бесконечная неабелева
группа.
\end{enumerate}
\end{primer*}
\isprtext{В оригинале к группе по умножению отнесено всё множество
рациональных чисел $p/q$, $q\ne0$; но у нуля нет обратного элемента, поэтому
его нужно исключить ($p\ne0$).}

\paragraph{Таблица умножения (таблица Кэли)}
Элементы выписывают по вертикали и горизонтали; на пересечении строки
$g_i$ и столбца $g_j$ ставят произведение $g_ig_j$. Абстрактных групп
4-го порядка всего две:
\[
\begin{array}{c|cccc}
 & e & a & b & c\\\hline
e & e & a & b & c\\
a & a & b & c & e\\
b & b & c & e & a\\
c & c & e & a & b
\end{array}
\qquad\qquad
\begin{array}{c|cccc}
 & e & a & b & c\\\hline
e & e & a & b & c\\
a & a & e & c & b\\
b & b & c & e & a\\
c & c & b & a & e
\end{array}
\]
Слева~— \emph{циклическая} группа: $a^2=b$, $a^3=c$, $a^4=e$ (генератор $a$);
её реализация~— повороты вокруг оси 4. Справа~— \emph{четверная группа Клейна}:
$a^2=b^2=c^2=e$, $ab=c$; генераторами могут служить любые два из $a$, $b$,
$c$; её реализация~— группа $222$ (три перпендикулярные оси 2) или группа
$2/m$. Обе группы абелевы, и каждый элемент образует свой класс.
Геометрически разные группы симметрии (у одной элемент~— отражение, у
другой~— поворот на $180^\circ$) могут иметь одинаковую таблицу умножения,
т.\,е. быть изоморфными.

% ----------------------------------------------------------------------
%  1.3.5. Точечные группы; 1.3.6. Пространственные группы;
%  1.3.7. Предельные группы симметрии
% ----------------------------------------------------------------------
\subsection{Точечные группы симметрии}\label{ssec:point-groups}

\begin{opredelenie}[Точечная группа]
Группа операций симметрии, оставляющих неподвижной хотя бы одну точку.
Ею описывают симметрию тел конечных размеров, а в кристаллографии~—
симметрию внешней формы кристалла и его физических свойств. Точечные
группы кристаллов называют также \term{кристаллическими классами}.
\end{opredelenie}

\paragraph{Группы первого рода}
Из одних поворотных осей получаются 11 групп: циклические $C_1$, $C_2$,
$C_3$, $C_4$, $C_6$ (международные символы 1, 2, 3, 4, 6) и сочетания осей
$222$, $32$, $422$, $622$, $23$, $432$.
\begin{itemize}
\item \emph{Группы поворотов} $N$ ($C_n$): единственный элемент~— ось $N$;
группа циклическая порядка $n=N$. Главная ось полярна.
\item \emph{Группы диэдров} $N2$ ($D_n$): главная ось $N$ и $n$
перпендикулярных ей осей 2; порядок $2n$; главная ось неполярна.
\end{itemize}

Примеры (рядом с порядком группы $n$ указано число классов $k$):
\begin{itemize}
\item $222$ ($D_2$): $E$, $C_2^{[100]}$, $C_2^{[010]}$, $C_2^{[001]}$; $n=4$, $k=4$~—
четверная группа Клейна; генераторы $C_2^{[100]}$ и $C_2^{[010]}$.
\item $422$ ($D_4$): ось 4 и четыре оси 2; $n=8$, $k=5$ (классы: $E$; $C_2$ вокруг
главной оси; два поворота на $\pm90^\circ$; две оси 2 вдоль осей координат;
две диагональные оси 2)\ispr{В оригинале для группы $422$ указано
$k=6$; у группы $D_4$ пять классов сопряжённых элементов.}; генераторы~— $C_4$
и одна из осей 2.
\item $23$ ($T$, группа тетраэдра): четыре оси 3 вдоль $\nfam{111}$ и три оси 2
вдоль $\nfam{100}$; $n=12$, $k=4$ ($E$; три $C_2$; четыре поворота на $120^\circ$;
четыре поворота на $240^\circ$)\ispr{В оригинале для группы $23$ указано
$k=6$; у группы $T$ четыре класса: повороты на $120^\circ$ и на $240^\circ$
вокруг осей $\nfam{111}$ образуют два \emph{разных} класса.}; генераторы
$C_2^{[001]}$ и $C_3^{[111]}$.
\item $432$ ($O$, группа октаэдра): три оси 4, четыре оси 3, шесть осей 2;
$n=24$, $k=5$; генераторы $C_4^{[001]}$ и $C_3^{[111]}$.
\end{itemize}

Сравнение $T$ и $O$ поучительно: равенство углов поворота~— необходимое,
но не достаточное условие сопряжённости. Восемь поворотов на $\pm120^\circ$
вокруг $\nfam{111}$ в группе $O$ образуют один класс, а в группе $T$~— два (в
$T$ нет операции, переводящей поворот на $+120^\circ$ в поворот на
$-120^\circ$).

\paragraph{Все 32 класса}
Добавив инверсионные оси, получаем ещё 21 группу (второго рода)~— всего
\vazhno{32 точечные группы} (табл.~\ref{tab:point-groups}). Самая
симметричная, $m\br3m$ ($O_h$),~— группа куба: 48 элементов, 10 классов,
генераторы $C_4^{[001]}$, $C_3^{[111]}$ и $\br1$ (рис.~\ref{fig:m3m}).

\begin{table}[p]
\centering
\caption{Кристаллические классы (точечные группы) по сингониям}\label{tab:point-groups}
\small\renewcommand{\arraystretch}{1.08}
\begin{tabular}{@{}llll@{}}
\toprule
Сингония & Международный & Шёнфлис & Формула симметрии\\
\midrule
Триклинная & $1$ & $C_1$ & $L_1$\\
 & $\br1$ & $C_i=S_2$ & $C$\\
\midrule
Моноклинная & $2$ & $C_2$ & $L_2$\\
 & $m$ & $C_s=C_{1h}$ & $P$\\
 & $2/m$ & $C_{2h}$ & $L_2PC$\\
\midrule
Ромбическая & $222$ & $D_2=V$ & $3L_2$\\
 & $mm2$ & $C_{2v}$ & $L_22P$\\
 & $mmm$ & $D_{2h}=V_h$ & $3L_23PC$\\
\midrule
Тетрагональная & $4$ & $C_4$ & $L_4$\\
 & $\br4$ & $S_4$ & $\bar L_4$\\
 & $4/m$ & $C_{4h}$ & $L_4PC$\\
 & $422$ & $D_4$ & $L_44L_2$\\
 & $4mm$ & $C_{4v}$ & $L_44P$\\
 & $\br42m$ & $D_{2d}=V_d$ & $\bar L_42L_22P$\\
 & $4/mmm$ & $D_{4h}$ & $L_44L_25PC$\\
\midrule
Тригональная & $3$ & $C_3$ & $L_3$\\
 & $\br3$ & $C_{3i}=S_6$ & $L_3C$\\
 & $32$ & $D_3$ & $L_33L_2$\\
 & $3m$ & $C_{3v}$ & $L_33P$\\
 & $\br3m$ & $D_{3d}$ & $L_33L_23PC$\\
\midrule
Гексагональная & $6$ & $C_6$ & $L_6$\\
 & $\br6$ & $C_{3h}$ & $L_3P$\\
 & $6/m$ & $C_{6h}$ & $L_6PC$\\
 & $622$ & $D_6$ & $L_66L_2$\\
 & $6mm$ & $C_{6v}$ & $L_66P$\\
 & $\br6m2$ & $D_{3h}$ & $L_33L_24P$\\
 & $6/mmm$ & $D_{6h}$ & $L_66L_27PC$\\
\midrule
Кубическая & $23$ & $T$ & $4L_33L_2$\\
 & $m\br3$ & $T_h$ & $4L_33L_23PC$\\
 & $432$ & $O$ & $3L_44L_36L_2$\\
 & $\br43m$ & $T_d$ & $3\bar L_44L_36P$\\
 & $m\br3m$ & $O_h$ & $3L_44L_36L_29PC$\\
\bottomrule
\multicolumn{4}{p{0.8\textwidth}}{\footnotesize В формуле симметрии $L_n$~— ось $n$-го порядка, $\bar L_n$~— инверсионная ось, $P$~— плоскость симметрии, $C$~— центр симметрии; число перед символом~— количество элементов.}
\end{tabular}
\end{table}

\paragraph{Обозначения}
В международной символике плоскость, перпендикулярная оси $N$, пишется
через косую черту ($N/m$), а плоскость, содержащая ось,~— рядом ($Nm$).
У Шёнфлиса первым соответствуют группы $C_{nh}$ (h~— горизонтальная
плоскость), вторым~— $C_{nv}$ (v~— вертикальная), $D_n$~— группы диэдров,
$S_n$~— зеркально-поворотные оси, $T$ и $O$~— группы тетраэдра и
октаэдра. В символике Шубникова точка между элементами означает их
параллельность, двоеточие~— перпендикулярность, косая черта~— косое
расположение.

\begin{figure}[ht]
\centering
\tikzset{mirror/.style={fmBlue, line width=0.6pt},
  ax4/.style={rectangle, draw=fmRed, fill=fmRed!35, inner sep=0pt, minimum size=5pt},
  ax3/.style={regular polygon, regular polygon sides=3, draw=fmAmber!90!black, fill=fmAmber!40, inner sep=0pt, minimum size=7pt},
  ax2/.style={ellipse, draw=fmGreen!60!black, fill=fmGreen!35, inner sep=0pt, minimum width=7pt, minimum height=4pt}}
\begin{tikzpicture}[scale=2.4]
  \draw[mirror, line width=1pt] (0,0) circle (1);
  % создано tools/gen_stereo.py: элементы симметрии m-3m на проекции (001)
\draw[mirror] (0.0000,-1.0000) -- (0.0000,-0.9742) -- (0.0000,-0.9490) -- (0.0000,-0.9244) -- (0.0000,-0.9004) -- (0.0000,-0.8770) -- (0.0000,-0.8541) -- (0.0000,-0.8317) -- (0.0000,-0.8098) -- (0.0000,-0.7883) -- (0.0000,-0.7673) -- (0.0000,-0.7467) -- (0.0000,-0.7265) -- (0.0000,-0.7067) -- (0.0000,-0.6873) -- (0.0000,-0.6682) -- (0.0000,-0.6494) -- (0.0000,-0.6310) -- (0.0000,-0.6128) -- (0.0000,-0.5949) -- (0.0000,-0.5774) -- (0.0000,-0.5600) -- (0.0000,-0.5430) -- (0.0000,-0.5261) -- (0.0000,-0.5095) -- (0.0000,-0.4931) -- (0.0000,-0.4770) -- (0.0000,-0.4610) -- (0.0000,-0.4452) -- (0.0000,-0.4296) -- (0.0000,-0.4142) -- (0.0000,-0.3990) -- (0.0000,-0.3839) -- (0.0000,-0.3689) -- (0.0000,-0.3541) -- (0.0000,-0.3395) -- (0.0000,-0.3249) -- (0.0000,-0.3105) -- (0.0000,-0.2962) -- (0.0000,-0.2820) -- (0.0000,-0.2679) -- (0.0000,-0.2540) -- (0.0000,-0.2401) -- (0.0000,-0.2263) -- (0.0000,-0.2126) -- (0.0000,-0.1989) -- (0.0000,-0.1853) -- (0.0000,-0.1718) -- (0.0000,-0.1584) -- (0.0000,-0.1450) -- (0.0000,-0.1317) -- (0.0000,-0.1184) -- (0.0000,-0.1051) -- (0.0000,-0.0919) -- (0.0000,-0.0787) -- (0.0000,-0.0655) -- (0.0000,-0.0524) -- (0.0000,-0.0393) -- (0.0000,-0.0262) -- (0.0000,-0.0131) -- (0.0000,-0.0000) -- (-0.0000,0.0131) -- (-0.0000,0.0262) -- (-0.0000,0.0393) -- (-0.0000,0.0524) -- (-0.0000,0.0655) -- (-0.0000,0.0787) -- (-0.0000,0.0919) -- (-0.0000,0.1051) -- (-0.0000,0.1184) -- (-0.0000,0.1317) -- (-0.0000,0.1450) -- (-0.0000,0.1584) -- (-0.0000,0.1718) -- (-0.0000,0.1853) -- (-0.0000,0.1989) -- (-0.0000,0.2126) -- (-0.0000,0.2263) -- (-0.0000,0.2401) -- (-0.0000,0.2540) -- (-0.0000,0.2679) -- (-0.0000,0.2820) -- (-0.0000,0.2962) -- (-0.0000,0.3105) -- (-0.0000,0.3249) -- (-0.0000,0.3395) -- (-0.0000,0.3541) -- (-0.0000,0.3689) -- (-0.0000,0.3839) -- (-0.0000,0.3990) -- (-0.0000,0.4142) -- (-0.0000,0.4296) -- (-0.0000,0.4452) -- (-0.0000,0.4610) -- (-0.0000,0.4770) -- (-0.0000,0.4931) -- (-0.0000,0.5095) -- (-0.0000,0.5261) -- (-0.0000,0.5430) -- (-0.0000,0.5600) -- (-0.0000,0.5774) -- (-0.0000,0.5949) -- (-0.0000,0.6128) -- (-0.0000,0.6310) -- (-0.0000,0.6494) -- (-0.0000,0.6682) -- (-0.0000,0.6873) -- (-0.0000,0.7067) -- (-0.0000,0.7265) -- (-0.0000,0.7467) -- (-0.0000,0.7673) -- (-0.0000,0.7883) -- (-0.0000,0.8098) -- (-0.0000,0.8317) -- (-0.0000,0.8541) -- (-0.0000,0.8770) -- (-0.0000,0.9004) -- (-0.0000,0.9244) -- (-0.0000,0.9490) -- (-0.0000,0.9742) -- (-0.0000,1.0000);
\draw[mirror] (1.0000,0.0000) -- (0.9742,0.0000) -- (0.9490,0.0000) -- (0.9244,0.0000) -- (0.9004,0.0000) -- (0.8770,0.0000) -- (0.8541,0.0000) -- (0.8317,0.0000) -- (0.8098,0.0000) -- (0.7883,0.0000) -- (0.7673,0.0000) -- (0.7467,0.0000) -- (0.7265,0.0000) -- (0.7067,0.0000) -- (0.6873,0.0000) -- (0.6682,0.0000) -- (0.6494,0.0000) -- (0.6310,0.0000) -- (0.6128,0.0000) -- (0.5949,0.0000) -- (0.5774,0.0000) -- (0.5600,0.0000) -- (0.5430,0.0000) -- (0.5261,0.0000) -- (0.5095,0.0000) -- (0.4931,0.0000) -- (0.4770,0.0000) -- (0.4610,0.0000) -- (0.4452,0.0000) -- (0.4296,0.0000) -- (0.4142,0.0000) -- (0.3990,0.0000) -- (0.3839,0.0000) -- (0.3689,0.0000) -- (0.3541,0.0000) -- (0.3395,0.0000) -- (0.3249,0.0000) -- (0.3105,0.0000) -- (0.2962,0.0000) -- (0.2820,0.0000) -- (0.2679,0.0000) -- (0.2540,0.0000) -- (0.2401,0.0000) -- (0.2263,0.0000) -- (0.2126,0.0000) -- (0.1989,0.0000) -- (0.1853,0.0000) -- (0.1718,0.0000) -- (0.1584,0.0000) -- (0.1450,0.0000) -- (0.1317,0.0000) -- (0.1184,0.0000) -- (0.1051,0.0000) -- (0.0919,0.0000) -- (0.0787,0.0000) -- (0.0655,0.0000) -- (0.0524,0.0000) -- (0.0393,0.0000) -- (0.0262,0.0000) -- (0.0131,0.0000) -- (0.0000,0.0000) -- (-0.0131,-0.0000) -- (-0.0262,-0.0000) -- (-0.0393,-0.0000) -- (-0.0524,-0.0000) -- (-0.0655,-0.0000) -- (-0.0787,-0.0000) -- (-0.0919,-0.0000) -- (-0.1051,-0.0000) -- (-0.1184,-0.0000) -- (-0.1317,-0.0000) -- (-0.1450,-0.0000) -- (-0.1584,-0.0000) -- (-0.1718,-0.0000) -- (-0.1853,-0.0000) -- (-0.1989,-0.0000) -- (-0.2126,-0.0000) -- (-0.2263,-0.0000) -- (-0.2401,-0.0000) -- (-0.2540,-0.0000) -- (-0.2679,-0.0000) -- (-0.2820,-0.0000) -- (-0.2962,-0.0000) -- (-0.3105,-0.0000) -- (-0.3249,-0.0000) -- (-0.3395,-0.0000) -- (-0.3541,-0.0000) -- (-0.3689,-0.0000) -- (-0.3839,-0.0000) -- (-0.3990,-0.0000) -- (-0.4142,-0.0000) -- (-0.4296,-0.0000) -- (-0.4452,-0.0000) -- (-0.4610,-0.0000) -- (-0.4770,-0.0000) -- (-0.4931,-0.0000) -- (-0.5095,-0.0000) -- (-0.5261,-0.0000) -- (-0.5430,-0.0000) -- (-0.5600,-0.0000) -- (-0.5774,-0.0000) -- (-0.5949,-0.0000) -- (-0.6128,-0.0000) -- (-0.6310,-0.0000) -- (-0.6494,-0.0000) -- (-0.6682,-0.0000) -- (-0.6873,-0.0000) -- (-0.7067,-0.0000) -- (-0.7265,-0.0000) -- (-0.7467,-0.0000) -- (-0.7673,-0.0000) -- (-0.7883,-0.0000) -- (-0.8098,-0.0000) -- (-0.8317,-0.0000) -- (-0.8541,-0.0000) -- (-0.8770,-0.0000) -- (-0.9004,-0.0000) -- (-0.9244,-0.0000) -- (-0.9490,-0.0000) -- (-0.9742,-0.0000) -- (-1.0000,-0.0000);
\draw[mirror] (0.7071,-0.7071) -- (0.6888,-0.6888) -- (0.6710,-0.6710) -- (0.6536,-0.6536) -- (0.6367,-0.6367) -- (0.6201,-0.6201) -- (0.6039,-0.6039) -- (0.5881,-0.5881) -- (0.5726,-0.5726) -- (0.5574,-0.5574) -- (0.5426,-0.5426) -- (0.5280,-0.5280) -- (0.5137,-0.5137) -- (0.4997,-0.4997) -- (0.4860,-0.4860) -- (0.4725,-0.4725) -- (0.4592,-0.4592) -- (0.4462,-0.4462) -- (0.4333,-0.4333) -- (0.4207,-0.4207) -- (0.4082,-0.4082) -- (0.3960,-0.3960) -- (0.3839,-0.3839) -- (0.3720,-0.3720) -- (0.3603,-0.3603) -- (0.3487,-0.3487) -- (0.3373,-0.3373) -- (0.3260,-0.3260) -- (0.3148,-0.3148) -- (0.3038,-0.3038) -- (0.2929,-0.2929) -- (0.2821,-0.2821) -- (0.2714,-0.2714) -- (0.2609,-0.2609) -- (0.2504,-0.2504) -- (0.2400,-0.2400) -- (0.2298,-0.2298) -- (0.2196,-0.2196) -- (0.2095,-0.2095) -- (0.1994,-0.1994) -- (0.1895,-0.1895) -- (0.1796,-0.1796) -- (0.1698,-0.1698) -- (0.1600,-0.1600) -- (0.1503,-0.1503) -- (0.1407,-0.1407) -- (0.1311,-0.1311) -- (0.1215,-0.1215) -- (0.1120,-0.1120) -- (0.1025,-0.1025) -- (0.0931,-0.0931) -- (0.0837,-0.0837) -- (0.0743,-0.0743) -- (0.0650,-0.0650) -- (0.0557,-0.0557) -- (0.0463,-0.0463) -- (0.0371,-0.0371) -- (0.0278,-0.0278) -- (0.0185,-0.0185) -- (0.0093,-0.0093) -- (0.0000,-0.0000) -- (-0.0093,0.0093) -- (-0.0185,0.0185) -- (-0.0278,0.0278) -- (-0.0371,0.0371) -- (-0.0463,0.0463) -- (-0.0557,0.0557) -- (-0.0650,0.0650) -- (-0.0743,0.0743) -- (-0.0837,0.0837) -- (-0.0931,0.0931) -- (-0.1025,0.1025) -- (-0.1120,0.1120) -- (-0.1215,0.1215) -- (-0.1311,0.1311) -- (-0.1407,0.1407) -- (-0.1503,0.1503) -- (-0.1600,0.1600) -- (-0.1698,0.1698) -- (-0.1796,0.1796) -- (-0.1895,0.1895) -- (-0.1994,0.1994) -- (-0.2095,0.2095) -- (-0.2196,0.2196) -- (-0.2298,0.2298) -- (-0.2400,0.2400) -- (-0.2504,0.2504) -- (-0.2609,0.2609) -- (-0.2714,0.2714) -- (-0.2821,0.2821) -- (-0.2929,0.2929) -- (-0.3038,0.3038) -- (-0.3148,0.3148) -- (-0.3260,0.3260) -- (-0.3373,0.3373) -- (-0.3487,0.3487) -- (-0.3603,0.3603) -- (-0.3720,0.3720) -- (-0.3839,0.3839) -- (-0.3960,0.3960) -- (-0.4082,0.4082) -- (-0.4207,0.4207) -- (-0.4333,0.4333) -- (-0.4462,0.4462) -- (-0.4592,0.4592) -- (-0.4725,0.4725) -- (-0.4860,0.4860) -- (-0.4997,0.4997) -- (-0.5137,0.5137) -- (-0.5280,0.5280) -- (-0.5426,0.5426) -- (-0.5574,0.5574) -- (-0.5726,0.5726) -- (-0.5881,0.5881) -- (-0.6039,0.6039) -- (-0.6201,0.6201) -- (-0.6367,0.6367) -- (-0.6536,0.6536) -- (-0.6710,0.6710) -- (-0.6888,0.6888) -- (-0.7071,0.7071);
\draw[mirror] (-0.7071,-0.7071) -- (-0.6888,-0.6888) -- (-0.6710,-0.6710) -- (-0.6536,-0.6536) -- (-0.6367,-0.6367) -- (-0.6201,-0.6201) -- (-0.6039,-0.6039) -- (-0.5881,-0.5881) -- (-0.5726,-0.5726) -- (-0.5574,-0.5574) -- (-0.5426,-0.5426) -- (-0.5280,-0.5280) -- (-0.5137,-0.5137) -- (-0.4997,-0.4997) -- (-0.4860,-0.4860) -- (-0.4725,-0.4725) -- (-0.4592,-0.4592) -- (-0.4462,-0.4462) -- (-0.4333,-0.4333) -- (-0.4207,-0.4207) -- (-0.4082,-0.4082) -- (-0.3960,-0.3960) -- (-0.3839,-0.3839) -- (-0.3720,-0.3720) -- (-0.3603,-0.3603) -- (-0.3487,-0.3487) -- (-0.3373,-0.3373) -- (-0.3260,-0.3260) -- (-0.3148,-0.3148) -- (-0.3038,-0.3038) -- (-0.2929,-0.2929) -- (-0.2821,-0.2821) -- (-0.2714,-0.2714) -- (-0.2609,-0.2609) -- (-0.2504,-0.2504) -- (-0.2400,-0.2400) -- (-0.2298,-0.2298) -- (-0.2196,-0.2196) -- (-0.2095,-0.2095) -- (-0.1994,-0.1994) -- (-0.1895,-0.1895) -- (-0.1796,-0.1796) -- (-0.1698,-0.1698) -- (-0.1600,-0.1600) -- (-0.1503,-0.1503) -- (-0.1407,-0.1407) -- (-0.1311,-0.1311) -- (-0.1215,-0.1215) -- (-0.1120,-0.1120) -- (-0.1025,-0.1025) -- (-0.0931,-0.0931) -- (-0.0837,-0.0837) -- (-0.0743,-0.0743) -- (-0.0650,-0.0650) -- (-0.0557,-0.0557) -- (-0.0463,-0.0463) -- (-0.0371,-0.0371) -- (-0.0278,-0.0278) -- (-0.0185,-0.0185) -- (-0.0093,-0.0093) -- (-0.0000,-0.0000) -- (0.0093,0.0093) -- (0.0185,0.0185) -- (0.0278,0.0278) -- (0.0371,0.0371) -- (0.0463,0.0463) -- (0.0557,0.0557) -- (0.0650,0.0650) -- (0.0743,0.0743) -- (0.0837,0.0837) -- (0.0931,0.0931) -- (0.1025,0.1025) -- (0.1120,0.1120) -- (0.1215,0.1215) -- (0.1311,0.1311) -- (0.1407,0.1407) -- (0.1503,0.1503) -- (0.1600,0.1600) -- (0.1698,0.1698) -- (0.1796,0.1796) -- (0.1895,0.1895) -- (0.1994,0.1994) -- (0.2095,0.2095) -- (0.2196,0.2196) -- (0.2298,0.2298) -- (0.2400,0.2400) -- (0.2504,0.2504) -- (0.2609,0.2609) -- (0.2714,0.2714) -- (0.2821,0.2821) -- (0.2929,0.2929) -- (0.3038,0.3038) -- (0.3148,0.3148) -- (0.3260,0.3260) -- (0.3373,0.3373) -- (0.3487,0.3487) -- (0.3603,0.3603) -- (0.3720,0.3720) -- (0.3839,0.3839) -- (0.3960,0.3960) -- (0.4082,0.4082) -- (0.4207,0.4207) -- (0.4333,0.4333) -- (0.4462,0.4462) -- (0.4592,0.4592) -- (0.4725,0.4725) -- (0.4860,0.4860) -- (0.4997,0.4997) -- (0.5137,0.5137) -- (0.5280,0.5280) -- (0.5426,0.5426) -- (0.5574,0.5574) -- (0.5726,0.5726) -- (0.5881,0.5881) -- (0.6039,0.6039) -- (0.6201,0.6201) -- (0.6367,0.6367) -- (0.6536,0.6536) -- (0.6710,0.6710) -- (0.6888,0.6888) -- (0.7071,0.7071);
\draw[mirror] (0.0000,-1.0000) -- (-0.0182,-0.9815) -- (-0.0357,-0.9630) -- (-0.0526,-0.9445) -- (-0.0688,-0.9261) -- (-0.0845,-0.9077) -- (-0.0996,-0.8893) -- (-0.1142,-0.8710) -- (-0.1282,-0.8528) -- (-0.1417,-0.8346) -- (-0.1547,-0.8165) -- (-0.1672,-0.7985) -- (-0.1793,-0.7805) -- (-0.1910,-0.7626) -- (-0.2022,-0.7448) -- (-0.2130,-0.7271) -- (-0.2234,-0.7095) -- (-0.2334,-0.6919) -- (-0.2430,-0.6745) -- (-0.2523,-0.6571) -- (-0.2612,-0.6398) -- (-0.2698,-0.6226) -- (-0.2780,-0.6055) -- (-0.2860,-0.5884) -- (-0.2936,-0.5715) -- (-0.3009,-0.5546) -- (-0.3080,-0.5378) -- (-0.3147,-0.5211) -- (-0.3212,-0.5045) -- (-0.3274,-0.4879) -- (-0.3333,-0.4714) -- (-0.3390,-0.4550) -- (-0.3445,-0.4386) -- (-0.3497,-0.4224) -- (-0.3546,-0.4061) -- (-0.3594,-0.3900) -- (-0.3639,-0.3739) -- (-0.3682,-0.3579) -- (-0.3723,-0.3419) -- (-0.3761,-0.3260) -- (-0.3798,-0.3101) -- (-0.3833,-0.2943) -- (-0.3865,-0.2785) -- (-0.3896,-0.2628) -- (-0.3925,-0.2471) -- (-0.3951,-0.2315) -- (-0.3976,-0.2159) -- (-0.4000,-0.2003) -- (-0.4021,-0.1848) -- (-0.4040,-0.1693) -- (-0.4058,-0.1538) -- (-0.4074,-0.1383) -- (-0.4089,-0.1229) -- (-0.4101,-0.1075) -- (-0.4112,-0.0921) -- (-0.4121,-0.0767) -- (-0.4129,-0.0614) -- (-0.4135,-0.0460) -- (-0.4139,-0.0307) -- (-0.4141,-0.0153) -- (-0.4142,-0.0000) -- (-0.4141,0.0153) -- (-0.4139,0.0307) -- (-0.4135,0.0460) -- (-0.4129,0.0614) -- (-0.4121,0.0767) -- (-0.4112,0.0921) -- (-0.4101,0.1075) -- (-0.4089,0.1229) -- (-0.4074,0.1383) -- (-0.4058,0.1538) -- (-0.4040,0.1693) -- (-0.4021,0.1848) -- (-0.4000,0.2003) -- (-0.3976,0.2159) -- (-0.3951,0.2315) -- (-0.3925,0.2471) -- (-0.3896,0.2628) -- (-0.3865,0.2785) -- (-0.3833,0.2943) -- (-0.3798,0.3101) -- (-0.3761,0.3260) -- (-0.3723,0.3419) -- (-0.3682,0.3579) -- (-0.3639,0.3739) -- (-0.3594,0.3900) -- (-0.3546,0.4061) -- (-0.3497,0.4224) -- (-0.3445,0.4386) -- (-0.3390,0.4550) -- (-0.3333,0.4714) -- (-0.3274,0.4879) -- (-0.3212,0.5045) -- (-0.3147,0.5211) -- (-0.3080,0.5378) -- (-0.3009,0.5546) -- (-0.2936,0.5715) -- (-0.2860,0.5884) -- (-0.2780,0.6055) -- (-0.2698,0.6226) -- (-0.2612,0.6398) -- (-0.2523,0.6571) -- (-0.2430,0.6745) -- (-0.2334,0.6919) -- (-0.2234,0.7095) -- (-0.2130,0.7271) -- (-0.2022,0.7448) -- (-0.1910,0.7626) -- (-0.1793,0.7805) -- (-0.1672,0.7985) -- (-0.1547,0.8165) -- (-0.1417,0.8346) -- (-0.1282,0.8528) -- (-0.1142,0.8710) -- (-0.0996,0.8893) -- (-0.0845,0.9077) -- (-0.0688,0.9261) -- (-0.0526,0.9445) -- (-0.0357,0.9630) -- (-0.0182,0.9815) -- (-0.0000,1.0000);
\draw[mirror] (0.0000,-1.0000) -- (0.0182,-0.9815) -- (0.0357,-0.9630) -- (0.0526,-0.9445) -- (0.0688,-0.9261) -- (0.0845,-0.9077) -- (0.0996,-0.8893) -- (0.1142,-0.8710) -- (0.1282,-0.8528) -- (0.1417,-0.8346) -- (0.1547,-0.8165) -- (0.1672,-0.7985) -- (0.1793,-0.7805) -- (0.1910,-0.7626) -- (0.2022,-0.7448) -- (0.2130,-0.7271) -- (0.2234,-0.7095) -- (0.2334,-0.6919) -- (0.2430,-0.6745) -- (0.2523,-0.6571) -- (0.2612,-0.6398) -- (0.2698,-0.6226) -- (0.2780,-0.6055) -- (0.2860,-0.5884) -- (0.2936,-0.5715) -- (0.3009,-0.5546) -- (0.3080,-0.5378) -- (0.3147,-0.5211) -- (0.3212,-0.5045) -- (0.3274,-0.4879) -- (0.3333,-0.4714) -- (0.3390,-0.4550) -- (0.3445,-0.4386) -- (0.3497,-0.4224) -- (0.3546,-0.4061) -- (0.3594,-0.3900) -- (0.3639,-0.3739) -- (0.3682,-0.3579) -- (0.3723,-0.3419) -- (0.3761,-0.3260) -- (0.3798,-0.3101) -- (0.3833,-0.2943) -- (0.3865,-0.2785) -- (0.3896,-0.2628) -- (0.3925,-0.2471) -- (0.3951,-0.2315) -- (0.3976,-0.2159) -- (0.4000,-0.2003) -- (0.4021,-0.1848) -- (0.4040,-0.1693) -- (0.4058,-0.1538) -- (0.4074,-0.1383) -- (0.4089,-0.1229) -- (0.4101,-0.1075) -- (0.4112,-0.0921) -- (0.4121,-0.0767) -- (0.4129,-0.0614) -- (0.4135,-0.0460) -- (0.4139,-0.0307) -- (0.4141,-0.0153) -- (0.4142,-0.0000) -- (0.4141,0.0153) -- (0.4139,0.0307) -- (0.4135,0.0460) -- (0.4129,0.0614) -- (0.4121,0.0767) -- (0.4112,0.0921) -- (0.4101,0.1075) -- (0.4089,0.1229) -- (0.4074,0.1383) -- (0.4058,0.1538) -- (0.4040,0.1693) -- (0.4021,0.1848) -- (0.4000,0.2003) -- (0.3976,0.2159) -- (0.3951,0.2315) -- (0.3925,0.2471) -- (0.3896,0.2628) -- (0.3865,0.2785) -- (0.3833,0.2943) -- (0.3798,0.3101) -- (0.3761,0.3260) -- (0.3723,0.3419) -- (0.3682,0.3579) -- (0.3639,0.3739) -- (0.3594,0.3900) -- (0.3546,0.4061) -- (0.3497,0.4224) -- (0.3445,0.4386) -- (0.3390,0.4550) -- (0.3333,0.4714) -- (0.3274,0.4879) -- (0.3212,0.5045) -- (0.3147,0.5211) -- (0.3080,0.5378) -- (0.3009,0.5546) -- (0.2936,0.5715) -- (0.2860,0.5884) -- (0.2780,0.6055) -- (0.2698,0.6226) -- (0.2612,0.6398) -- (0.2523,0.6571) -- (0.2430,0.6745) -- (0.2334,0.6919) -- (0.2234,0.7095) -- (0.2130,0.7271) -- (0.2022,0.7448) -- (0.1910,0.7626) -- (0.1793,0.7805) -- (0.1672,0.7985) -- (0.1547,0.8165) -- (0.1417,0.8346) -- (0.1282,0.8528) -- (0.1142,0.8710) -- (0.0996,0.8893) -- (0.0845,0.9077) -- (0.0688,0.9261) -- (0.0526,0.9445) -- (0.0357,0.9630) -- (0.0182,0.9815) -- (0.0000,1.0000);
\draw[mirror] (1.0000,0.0000) -- (0.9815,-0.0182) -- (0.9630,-0.0357) -- (0.9445,-0.0526) -- (0.9261,-0.0688) -- (0.9077,-0.0845) -- (0.8893,-0.0996) -- (0.8710,-0.1142) -- (0.8528,-0.1282) -- (0.8346,-0.1417) -- (0.8165,-0.1547) -- (0.7985,-0.1672) -- (0.7805,-0.1793) -- (0.7626,-0.1910) -- (0.7448,-0.2022) -- (0.7271,-0.2130) -- (0.7095,-0.2234) -- (0.6919,-0.2334) -- (0.6745,-0.2430) -- (0.6571,-0.2523) -- (0.6398,-0.2612) -- (0.6226,-0.2698) -- (0.6055,-0.2780) -- (0.5884,-0.2860) -- (0.5715,-0.2936) -- (0.5546,-0.3009) -- (0.5378,-0.3080) -- (0.5211,-0.3147) -- (0.5045,-0.3212) -- (0.4879,-0.3274) -- (0.4714,-0.3333) -- (0.4550,-0.3390) -- (0.4386,-0.3445) -- (0.4224,-0.3497) -- (0.4061,-0.3546) -- (0.3900,-0.3594) -- (0.3739,-0.3639) -- (0.3579,-0.3682) -- (0.3419,-0.3723) -- (0.3260,-0.3761) -- (0.3101,-0.3798) -- (0.2943,-0.3833) -- (0.2785,-0.3865) -- (0.2628,-0.3896) -- (0.2471,-0.3925) -- (0.2315,-0.3951) -- (0.2159,-0.3976) -- (0.2003,-0.4000) -- (0.1848,-0.4021) -- (0.1693,-0.4040) -- (0.1538,-0.4058) -- (0.1383,-0.4074) -- (0.1229,-0.4089) -- (0.1075,-0.4101) -- (0.0921,-0.4112) -- (0.0767,-0.4121) -- (0.0614,-0.4129) -- (0.0460,-0.4135) -- (0.0307,-0.4139) -- (0.0153,-0.4141) -- (0.0000,-0.4142) -- (-0.0153,-0.4141) -- (-0.0307,-0.4139) -- (-0.0460,-0.4135) -- (-0.0614,-0.4129) -- (-0.0767,-0.4121) -- (-0.0921,-0.4112) -- (-0.1075,-0.4101) -- (-0.1229,-0.4089) -- (-0.1383,-0.4074) -- (-0.1538,-0.4058) -- (-0.1693,-0.4040) -- (-0.1848,-0.4021) -- (-0.2003,-0.4000) -- (-0.2159,-0.3976) -- (-0.2315,-0.3951) -- (-0.2471,-0.3925) -- (-0.2628,-0.3896) -- (-0.2785,-0.3865) -- (-0.2943,-0.3833) -- (-0.3101,-0.3798) -- (-0.3260,-0.3761) -- (-0.3419,-0.3723) -- (-0.3579,-0.3682) -- (-0.3739,-0.3639) -- (-0.3900,-0.3594) -- (-0.4061,-0.3546) -- (-0.4224,-0.3497) -- (-0.4386,-0.3445) -- (-0.4550,-0.3390) -- (-0.4714,-0.3333) -- (-0.4879,-0.3274) -- (-0.5045,-0.3212) -- (-0.5211,-0.3147) -- (-0.5378,-0.3080) -- (-0.5546,-0.3009) -- (-0.5715,-0.2936) -- (-0.5884,-0.2860) -- (-0.6055,-0.2780) -- (-0.6226,-0.2698) -- (-0.6398,-0.2612) -- (-0.6571,-0.2523) -- (-0.6745,-0.2430) -- (-0.6919,-0.2334) -- (-0.7095,-0.2234) -- (-0.7271,-0.2130) -- (-0.7448,-0.2022) -- (-0.7626,-0.1910) -- (-0.7805,-0.1793) -- (-0.7985,-0.1672) -- (-0.8165,-0.1547) -- (-0.8346,-0.1417) -- (-0.8528,-0.1282) -- (-0.8710,-0.1142) -- (-0.8893,-0.0996) -- (-0.9077,-0.0845) -- (-0.9261,-0.0688) -- (-0.9445,-0.0526) -- (-0.9630,-0.0357) -- (-0.9815,-0.0182) -- (-1.0000,-0.0000);
\draw[mirror] (1.0000,0.0000) -- (0.9815,0.0182) -- (0.9630,0.0357) -- (0.9445,0.0526) -- (0.9261,0.0688) -- (0.9077,0.0845) -- (0.8893,0.0996) -- (0.8710,0.1142) -- (0.8528,0.1282) -- (0.8346,0.1417) -- (0.8165,0.1547) -- (0.7985,0.1672) -- (0.7805,0.1793) -- (0.7626,0.1910) -- (0.7448,0.2022) -- (0.7271,0.2130) -- (0.7095,0.2234) -- (0.6919,0.2334) -- (0.6745,0.2430) -- (0.6571,0.2523) -- (0.6398,0.2612) -- (0.6226,0.2698) -- (0.6055,0.2780) -- (0.5884,0.2860) -- (0.5715,0.2936) -- (0.5546,0.3009) -- (0.5378,0.3080) -- (0.5211,0.3147) -- (0.5045,0.3212) -- (0.4879,0.3274) -- (0.4714,0.3333) -- (0.4550,0.3390) -- (0.4386,0.3445) -- (0.4224,0.3497) -- (0.4061,0.3546) -- (0.3900,0.3594) -- (0.3739,0.3639) -- (0.3579,0.3682) -- (0.3419,0.3723) -- (0.3260,0.3761) -- (0.3101,0.3798) -- (0.2943,0.3833) -- (0.2785,0.3865) -- (0.2628,0.3896) -- (0.2471,0.3925) -- (0.2315,0.3951) -- (0.2159,0.3976) -- (0.2003,0.4000) -- (0.1848,0.4021) -- (0.1693,0.4040) -- (0.1538,0.4058) -- (0.1383,0.4074) -- (0.1229,0.4089) -- (0.1075,0.4101) -- (0.0921,0.4112) -- (0.0767,0.4121) -- (0.0614,0.4129) -- (0.0460,0.4135) -- (0.0307,0.4139) -- (0.0153,0.4141) -- (0.0000,0.4142) -- (-0.0153,0.4141) -- (-0.0307,0.4139) -- (-0.0460,0.4135) -- (-0.0614,0.4129) -- (-0.0767,0.4121) -- (-0.0921,0.4112) -- (-0.1075,0.4101) -- (-0.1229,0.4089) -- (-0.1383,0.4074) -- (-0.1538,0.4058) -- (-0.1693,0.4040) -- (-0.1848,0.4021) -- (-0.2003,0.4000) -- (-0.2159,0.3976) -- (-0.2315,0.3951) -- (-0.2471,0.3925) -- (-0.2628,0.3896) -- (-0.2785,0.3865) -- (-0.2943,0.3833) -- (-0.3101,0.3798) -- (-0.3260,0.3761) -- (-0.3419,0.3723) -- (-0.3579,0.3682) -- (-0.3739,0.3639) -- (-0.3900,0.3594) -- (-0.4061,0.3546) -- (-0.4224,0.3497) -- (-0.4386,0.3445) -- (-0.4550,0.3390) -- (-0.4714,0.3333) -- (-0.4879,0.3274) -- (-0.5045,0.3212) -- (-0.5211,0.3147) -- (-0.5378,0.3080) -- (-0.5546,0.3009) -- (-0.5715,0.2936) -- (-0.5884,0.2860) -- (-0.6055,0.2780) -- (-0.6226,0.2698) -- (-0.6398,0.2612) -- (-0.6571,0.2523) -- (-0.6745,0.2430) -- (-0.6919,0.2334) -- (-0.7095,0.2234) -- (-0.7271,0.2130) -- (-0.7448,0.2022) -- (-0.7626,0.1910) -- (-0.7805,0.1793) -- (-0.7985,0.1672) -- (-0.8165,0.1547) -- (-0.8346,0.1417) -- (-0.8528,0.1282) -- (-0.8710,0.1142) -- (-0.8893,0.0996) -- (-0.9077,0.0845) -- (-0.9261,0.0688) -- (-0.9445,0.0526) -- (-0.9630,0.0357) -- (-0.9815,0.0182) -- (-1.0000,0.0000);
\node[ax4] at (1.0000,0.0000) {};
\node[ax4] at (-1.0000,0.0000) {};
\node[ax4] at (0.0000,0.0000) {};
\node[ax4] at (0.0000,1.0000) {};
\node[ax4] at (0.0000,-1.0000) {};
\node[ax2, rotate=0] at (0.4142,0.0000) {};
\node[ax2, rotate=180] at (-0.4142,0.0000) {};
\node[ax2, rotate=45] at (0.7071,0.7071) {};
\node[ax2, rotate=-45] at (0.7071,-0.7071) {};
\node[ax2, rotate=135] at (-0.7071,0.7071) {};
\node[ax2, rotate=-135] at (-0.7071,-0.7071) {};
\node[ax2, rotate=90] at (0.0000,0.4142) {};
\node[ax2, rotate=-90] at (0.0000,-0.4142) {};
\node[ax3] at (0.3660,0.3660) {};
\node[ax3] at (0.3660,-0.3660) {};
\node[ax3] at (-0.3660,0.3660) {};
\node[ax3] at (-0.3660,-0.3660) {};

\end{tikzpicture}
\caption{Элементы симметрии группы $m\br3m$ ($O_h$) на стереографической
проекции $(001)$ (расчёт): оси 4 (квадраты), 3 (треугольники), 2 (овалы);
линии~— девять плоскостей симметрии (основной круг~— плоскость $(001)$).
Без плоскостей это группа $432$; оставив лишь оси 2 вдоль $\nfam{100}$ и
оси 3, получим группу $23$}
\label{fig:m3m}
\end{figure}

\paragraph{Важные подразделения классов}
\begin{itemize}
\item \emph{Центр симметрии.} В классах с центром симметрии нет полярных
направлений, а значит, нет и полярных свойств~— пироэлектричества,
пьезоэлектричества. Дифракционная картина всегда центросимметрична
(закон Фриделя, ниже), поэтому дифракционными методами различают лишь
11 \term{лауэвских классов}~— центросимметричных:
$\br1$, $2/m$, $mmm$, $4/m$, $4/mmm$, $\br3$, $\br3m$, $6/m$, $6/mmm$, $m\br3$, $m\br3m$.
\item \emph{Энантиоморфизм.} Объекты, симметрия которых содержит только
повороты, могут существовать в двух зеркально равных («правой» и
«левой») формах~— \term{энантиоморфных}; это свойство называют также
\emph{хиральностью} (правый и левый кварц). Энантиоморфных классов 11:
$1$, $2$, $3$, $4$, $6$, $222$, $32$, $422$, $622$, $23$, $432$.
\item \emph{Простые формы.} Многогранник, все грани которого получаются из
одной симметрией класса,~— \term{простая форма}. В классе $m\br3m$ из грани
$(100)$ получается куб, из $(111)$~— октаэдр, из $(110)$~— ромбический
додекаэдр (рис.~\ref{fig:forms}).
\end{itemize}

\begin{figure}[ht]
\centering
\tikzset{bzfront/.style={atomLine, line width=0.7pt, line join=round},
  bzback/.style={atomLine!55, line width=0.4pt, dashed}}
\begin{subfigure}[b]{0.3\textwidth}\centering
\begin{tikzpicture}[cubeproj, scale=1.7]
  \fill[fmBlue!12] (0,0,0) -- (1,0,0) -- (1,0,1) -- (0,0,1) -- cycle;
  \fill[fmBlue!20] (1,0,0) -- (1,1,0) -- (1,1,1) -- (1,0,1) -- cycle;
  \fill[fmBlue!8] (0,0,1) -- (1,0,1) -- (1,1,1) -- (0,1,1) -- cycle;
  \CubeEdges
\end{tikzpicture}
\caption{куб $\{100\}$}
\end{subfigure}\hfill
\begin{subfigure}[b]{0.3\textwidth}\centering
\begin{tikzpicture}[cubeproj, scale=1.0]
  \coordinate (X) at (1,0,0); \coordinate (x) at (-1,0,0);
  \coordinate (Y) at (0,1,0); \coordinate (y) at (0,-1,0);
  \coordinate (Z) at (0,0,1.25); \coordinate (z) at (0,0,-1.25);
  \fill[fmBlue!20] (X) -- (y) -- (Z) -- cycle; \fill[fmBlue!12] (x) -- (y) -- (Z) -- cycle;
  \fill[fmBlue!10] (X) -- (y) -- (z) -- cycle; \fill[fmBlue!6] (x) -- (y) -- (z) -- cycle;
  \draw[atomLine!55, dashed, line width=0.4pt] (Y) -- (X) (Y) -- (x) (Y) -- (Z) (Y) -- (z);
  \draw[atomLine, line width=0.7pt, line join=round] (X) -- (Z) -- (x) -- (z) -- cycle (y) -- (X) (y) -- (Z) (y) -- (x) (y) -- (z);
\end{tikzpicture}
\caption{октаэдр $\{111\}$}
\end{subfigure}\hfill
\begin{subfigure}[b]{0.3\textwidth}\centering
\begin{tikzpicture}[scale=1.25]
  % сгенерировано tools/gen_bands.py — не править вручную
\fill[fmBlue!18] (-0.676,-0.539) -- (-0.469,-0.273) -- (-0.676,0.412) -- (-0.883,0.145) -- cycle;
\fill[fmBlue!14] (0.207,0.267) -- (-0.469,-0.273) -- (0.207,-0.684) -- (0.883,-0.145) -- cycle;
\fill[fmBlue!12] (0.207,-0.684) -- (-0.469,-0.273) -- (-0.676,-0.539) -- (0.000,-0.951) -- cycle;
\fill[fmBlue!20] (-0.676,0.412) -- (-0.469,-0.273) -- (0.207,0.267) -- (0.000,0.951) -- cycle;
\fill[fmBlue!8] (0.676,-0.412) -- (0.883,-0.145) -- (0.207,-0.684) -- (0.000,-0.951) -- cycle;
\fill[fmBlue!16] (0.207,0.267) -- (0.883,-0.145) -- (0.676,0.539) -- (0.000,0.951) -- cycle;
\draw[bzback] (0.469,0.273) -- (-0.207,0.684);
\draw[bzback] (0.469,0.273) -- (-0.207,-0.267);
\draw[bzback] (-0.883,0.145) -- (-0.207,-0.267);
\draw[bzback] (-0.883,0.145) -- (-0.207,0.684);
\draw[bzback] (0.469,0.273) -- (0.676,-0.412);
\draw[bzback] (0.469,0.273) -- (0.676,0.539);
\draw[bzfront] (0.883,-0.145) -- (0.676,0.539);
\draw[bzfront] (0.883,-0.145) -- (0.676,-0.412);
\draw[bzfront] (-0.469,-0.273) -- (-0.676,-0.539);
\draw[bzfront] (-0.469,-0.273) -- (-0.676,0.412);
\draw[bzfront] (-0.883,0.145) -- (-0.676,0.412);
\draw[bzfront] (-0.883,0.145) -- (-0.676,-0.539);
\draw[bzfront] (-0.469,-0.273) -- (0.207,0.267);
\draw[bzfront] (-0.469,-0.273) -- (0.207,-0.684);
\draw[bzfront] (0.883,-0.145) -- (0.207,-0.684);
\draw[bzfront] (0.883,-0.145) -- (0.207,0.267);
\draw[bzfront] (0.000,-0.951) -- (0.676,-0.412);
\draw[bzback] (0.000,-0.951) -- (-0.207,-0.267);
\draw[bzback] (0.000,0.951) -- (-0.207,0.684);
\draw[bzfront] (0.000,0.951) -- (0.676,0.539);
\draw[bzfront] (0.000,-0.951) -- (-0.676,-0.539);
\draw[bzfront] (0.000,-0.951) -- (0.207,-0.684);
\draw[bzfront] (0.000,0.951) -- (0.207,0.267);
\draw[bzfront] (0.000,0.951) -- (-0.676,0.412);
\coordinate (bzG) at (0.000,0.000);
\coordinate (bzH) at (-0.469,-0.273);
\coordinate (bzN) at (0.207,-0.209);
\coordinate (bzP) at (0.207,0.267);

\end{tikzpicture}
\caption{ромбододекаэдр $\{110\}$}
\end{subfigure}
\caption{Простые формы класса $m\br3m$: куб (6 граней), октаэдр (8 граней),
ромбический додекаэдр (12 граней); число граней~— множитель
повторяемости формы}
\label{fig:forms}
\end{figure}

\paragraph{Решётки Браве}
О.~Браве показал, что все многообразие трансляционной симметрии
кристаллов исчерпывается \vazhno{14 типами решёток}, распределёнными по
7 сингониям (рис.~\ref{fig:bravais}, табл.~\ref{tab:bravais}). Ячейку Браве
выбирают по трём условиям, применяемым \emph{последовательно} (первое
важнее второго, второе~— третьего):
\begin{enumerate}
\item симметрия ячейки должна соответствовать симметрии кристалла
(высшей симметрии его сингонии), а рёбра~— быть трансляциями решётки;
\item ячейка должна иметь максимальное число прямых углов, равных углов
и равных рёбер;
\item при этом её объём должен быть минимальным.
\end{enumerate}
По расположению узлов ячейки бывают примитивные ($P$), базоцентрированные
($C$, $A$ или $B$~— узлы в центрах пары параллельных граней),
объёмноцентрированные ($I$~— узел в центре) и гранецентрированные ($F$~—
узлы в центрах всех граней); в тригональной сингонии примитивную
ромбоэдрическую ячейку обозначают $R$.

\begin{table}[ht]
\centering
\caption{Ячейки Браве по сингониям}\label{tab:bravais}
\small
\begin{tabular}{@{}lc@{\qquad}lc@{}}
\toprule
Сингония & Ячейки & Сингония & Ячейки\\
\midrule
Триклинная & $P$ & Тетрагональная & $P$, $I$\\
Моноклинная & $P$, $C$ & Тригональная & $R$\\
Ромбическая & $P$, $C$, $I$, $F$ & Гексагональная & $P$\\
 & & Кубическая & $P$, $I$, $F$\\
\bottomrule
\end{tabular}
\end{table}

\begin{figure}[p]
\centering
\tikzset{bravE/.style={atomLine, line width=0.6pt},
  bravH/.style={atomLine!50, line width=0.45pt, dashed},
  bravN/.style={circle, fill=atomA, draw=atomLine, line width=0.3pt, inner sep=0pt, minimum size=4.2pt},
  bravC/.style={circle, fill=atomB, draw=atomLine, line width=0.3pt, inner sep=0pt, minimum size=4.2pt}}
\begin{tikzpicture}[scale=1.15]
  % создано tools/gen_cryst.py: 14 решёток Браве
\draw[bravE] (0.000,0.000) -- (1.100,0.000);
\draw[bravH] (0.000,0.000) -- (0.720,0.300);
\draw[bravE] (0.000,0.000) -- (0.334,1.160);
\draw[bravE] (0.334,1.160) -- (1.434,1.160);
\draw[bravE] (0.334,1.160) -- (1.054,1.460);
\draw[bravH] (0.720,0.300) -- (1.820,0.300);
\draw[bravH] (0.720,0.300) -- (1.054,1.460);
\draw[bravE] (1.054,1.460) -- (2.154,1.460);
\draw[bravE] (1.100,0.000) -- (1.820,0.300);
\draw[bravE] (1.100,0.000) -- (1.434,1.160);
\draw[bravE] (1.434,1.160) -- (2.154,1.460);
\draw[bravE] (1.820,0.300) -- (2.154,1.460);
\node[bravN] at (0.000,0.000) {};
\node[bravN] at (0.334,1.160) {};
\node[bravN] at (0.720,0.300) {};
\node[bravN] at (1.054,1.460) {};
\node[bravN] at (1.100,0.000) {};
\node[bravN] at (1.434,1.160) {};
\node[bravN] at (1.820,0.300) {};
\node[bravN] at (2.154,1.460) {};
\node[font=\small, anchor=north] at (0.750,-0.250) {триклинная $P$};
\draw[bravE] (3.400,0.000) -- (4.500,0.000);
\draw[bravH] (3.400,0.000) -- (3.820,0.300);
\draw[bravE] (3.400,0.000) -- (3.750,1.050);
\draw[bravE] (3.750,1.050) -- (4.850,1.050);
\draw[bravE] (3.750,1.050) -- (4.170,1.350);
\draw[bravH] (3.820,0.300) -- (4.920,0.300);
\draw[bravH] (3.820,0.300) -- (4.170,1.350);
\draw[bravE] (4.170,1.350) -- (5.270,1.350);
\draw[bravE] (4.500,0.000) -- (4.920,0.300);
\draw[bravE] (4.500,0.000) -- (4.850,1.050);
\draw[bravE] (4.850,1.050) -- (5.270,1.350);
\draw[bravE] (4.920,0.300) -- (5.270,1.350);
\node[bravN] at (3.400,0.000) {};
\node[bravN] at (3.750,1.050) {};
\node[bravN] at (3.820,0.300) {};
\node[bravN] at (4.170,1.350) {};
\node[bravN] at (4.500,0.000) {};
\node[bravN] at (4.850,1.050) {};
\node[bravN] at (4.920,0.300) {};
\node[bravN] at (5.270,1.350) {};
\node[font=\small, anchor=north] at (4.150,-0.250) {моноклинная $P$};
\draw[bravE] (6.800,0.000) -- (7.900,0.000);
\draw[bravH] (6.800,0.000) -- (7.220,0.300);
\draw[bravE] (6.800,0.000) -- (7.150,1.050);
\draw[bravE] (7.150,1.050) -- (8.250,1.050);
\draw[bravE] (7.150,1.050) -- (7.570,1.350);
\draw[bravH] (7.220,0.300) -- (8.320,0.300);
\draw[bravH] (7.220,0.300) -- (7.570,1.350);
\draw[bravE] (7.570,1.350) -- (8.670,1.350);
\draw[bravE] (7.900,0.000) -- (8.320,0.300);
\draw[bravE] (7.900,0.000) -- (8.250,1.050);
\draw[bravE] (8.250,1.050) -- (8.670,1.350);
\draw[bravE] (8.320,0.300) -- (8.670,1.350);
\node[bravN] at (6.800,0.000) {};
\node[bravN] at (7.150,1.050) {};
\node[bravN] at (7.220,0.300) {};
\node[bravN] at (7.570,1.350) {};
\node[bravN] at (7.900,0.000) {};
\node[bravN] at (8.250,1.050) {};
\node[bravN] at (8.320,0.300) {};
\node[bravN] at (8.670,1.350) {};
\node[bravC] at (7.560,0.150) {};
\node[bravC] at (7.910,1.200) {};
\node[font=\small, anchor=north] at (7.550,-0.250) {моноклинная $C$};
\draw[bravE] (10.200,0.000) -- (10.304,0.664);
\draw[bravE] (10.200,0.000) -- (10.858,1.059);
\draw[bravE] (10.200,0.000) -- (9.439,0.861);
\draw[bravE] (9.439,0.861) -- (9.542,1.525);
\draw[bravH] (9.439,0.861) -- (10.096,1.920);
\draw[bravE] (10.858,1.059) -- (10.961,1.723);
\draw[bravH] (10.858,1.059) -- (10.096,1.920);
\draw[bravH] (10.096,1.920) -- (10.200,2.584);
\draw[bravE] (10.304,0.664) -- (10.961,1.723);
\draw[bravE] (10.304,0.664) -- (9.542,1.525);
\draw[bravE] (9.542,1.525) -- (10.200,2.584);
\draw[bravE] (10.961,1.723) -- (10.200,2.584);
\node[bravN] at (10.200,0.000) {};
\node[bravN] at (9.439,0.861) {};
\node[bravN] at (10.858,1.059) {};
\node[bravN] at (10.096,1.920) {};
\node[bravN] at (10.304,0.664) {};
\node[bravN] at (9.542,1.525) {};
\node[bravN] at (10.961,1.723) {};
\node[bravN] at (10.200,2.584) {};
\node[font=\small, anchor=north] at (10.950,-0.250) {тригональная $R$};
\draw[bravE] (0.000,-3.050) -- (1.250,-3.050);
\draw[bravH] (0.000,-3.050) -- (0.378,-2.780);
\draw[bravE] (0.000,-3.050) -- (0.000,-2.050);
\draw[bravE] (0.000,-2.050) -- (1.250,-2.050);
\draw[bravE] (0.000,-2.050) -- (0.378,-1.780);
\draw[bravH] (0.378,-2.780) -- (1.628,-2.780);
\draw[bravH] (0.378,-2.780) -- (0.378,-1.780);
\draw[bravE] (0.378,-1.780) -- (1.628,-1.780);
\draw[bravE] (1.250,-3.050) -- (1.628,-2.780);
\draw[bravE] (1.250,-3.050) -- (1.250,-2.050);
\draw[bravE] (1.250,-2.050) -- (1.628,-1.780);
\draw[bravE] (1.628,-2.780) -- (1.628,-1.780);
\node[bravN] at (0.000,-3.050) {};
\node[bravN] at (0.000,-2.050) {};
\node[bravN] at (0.378,-2.780) {};
\node[bravN] at (0.378,-1.780) {};
\node[bravN] at (1.250,-3.050) {};
\node[bravN] at (1.250,-2.050) {};
\node[bravN] at (1.628,-2.780) {};
\node[bravN] at (1.628,-1.780) {};
\node[font=\small, anchor=north] at (0.750,-3.300) {ромбическая $P$};
\draw[bravE] (3.400,-3.050) -- (4.650,-3.050);
\draw[bravH] (3.400,-3.050) -- (3.778,-2.780);
\draw[bravE] (3.400,-3.050) -- (3.400,-2.050);
\draw[bravE] (3.400,-2.050) -- (4.650,-2.050);
\draw[bravE] (3.400,-2.050) -- (3.778,-1.780);
\draw[bravH] (3.778,-2.780) -- (5.028,-2.780);
\draw[bravH] (3.778,-2.780) -- (3.778,-1.780);
\draw[bravE] (3.778,-1.780) -- (5.028,-1.780);
\draw[bravE] (4.650,-3.050) -- (5.028,-2.780);
\draw[bravE] (4.650,-3.050) -- (4.650,-2.050);
\draw[bravE] (4.650,-2.050) -- (5.028,-1.780);
\draw[bravE] (5.028,-2.780) -- (5.028,-1.780);
\node[bravN] at (3.400,-3.050) {};
\node[bravN] at (3.400,-2.050) {};
\node[bravN] at (3.778,-2.780) {};
\node[bravN] at (3.778,-1.780) {};
\node[bravN] at (4.650,-3.050) {};
\node[bravN] at (4.650,-2.050) {};
\node[bravN] at (5.028,-2.780) {};
\node[bravN] at (5.028,-1.780) {};
\node[bravC] at (4.214,-2.915) {};
\node[bravC] at (4.214,-1.915) {};
\node[font=\small, anchor=north] at (4.150,-3.300) {ромбическая $C$};
\draw[bravE] (6.800,-3.050) -- (8.050,-3.050);
\draw[bravH] (6.800,-3.050) -- (7.178,-2.780);
\draw[bravE] (6.800,-3.050) -- (6.800,-2.050);
\draw[bravE] (6.800,-2.050) -- (8.050,-2.050);
\draw[bravE] (6.800,-2.050) -- (7.178,-1.780);
\draw[bravH] (7.178,-2.780) -- (8.428,-2.780);
\draw[bravH] (7.178,-2.780) -- (7.178,-1.780);
\draw[bravE] (7.178,-1.780) -- (8.428,-1.780);
\draw[bravE] (8.050,-3.050) -- (8.428,-2.780);
\draw[bravE] (8.050,-3.050) -- (8.050,-2.050);
\draw[bravE] (8.050,-2.050) -- (8.428,-1.780);
\draw[bravE] (8.428,-2.780) -- (8.428,-1.780);
\node[bravN] at (6.800,-3.050) {};
\node[bravN] at (6.800,-2.050) {};
\node[bravN] at (7.178,-2.780) {};
\node[bravN] at (7.178,-1.780) {};
\node[bravN] at (8.050,-3.050) {};
\node[bravN] at (8.050,-2.050) {};
\node[bravN] at (8.428,-2.780) {};
\node[bravN] at (8.428,-1.780) {};
\node[bravC] at (7.614,-2.415) {};
\node[font=\small, anchor=north] at (7.550,-3.300) {ромбическая $I$};
\draw[bravE] (10.200,-3.050) -- (11.450,-3.050);
\draw[bravH] (10.200,-3.050) -- (10.578,-2.780);
\draw[bravE] (10.200,-3.050) -- (10.200,-2.050);
\draw[bravE] (10.200,-2.050) -- (11.450,-2.050);
\draw[bravE] (10.200,-2.050) -- (10.578,-1.780);
\draw[bravH] (10.578,-2.780) -- (11.828,-2.780);
\draw[bravH] (10.578,-2.780) -- (10.578,-1.780);
\draw[bravE] (10.578,-1.780) -- (11.828,-1.780);
\draw[bravE] (11.450,-3.050) -- (11.828,-2.780);
\draw[bravE] (11.450,-3.050) -- (11.450,-2.050);
\draw[bravE] (11.450,-2.050) -- (11.828,-1.780);
\draw[bravE] (11.828,-2.780) -- (11.828,-1.780);
\node[bravN] at (10.200,-3.050) {};
\node[bravN] at (10.200,-2.050) {};
\node[bravN] at (10.578,-2.780) {};
\node[bravN] at (10.578,-1.780) {};
\node[bravN] at (11.450,-3.050) {};
\node[bravN] at (11.450,-2.050) {};
\node[bravN] at (11.828,-2.780) {};
\node[bravN] at (11.828,-1.780) {};
\node[bravC] at (11.014,-2.915) {};
\node[bravC] at (11.014,-1.915) {};
\node[bravC] at (10.389,-2.415) {};
\node[bravC] at (11.639,-2.415) {};
\node[bravC] at (10.825,-2.550) {};
\node[bravC] at (11.203,-2.280) {};
\node[font=\small, anchor=north] at (10.950,-3.300) {ромбическая $F$};
\draw[bravE] (0.000,-6.100) -- (0.950,-6.100);
\draw[bravH] (0.000,-6.100) -- (0.399,-5.815);
\draw[bravE] (0.000,-6.100) -- (0.000,-4.750);
\draw[bravE] (0.000,-4.750) -- (0.950,-4.750);
\draw[bravE] (0.000,-4.750) -- (0.399,-4.465);
\draw[bravH] (0.399,-5.815) -- (1.349,-5.815);
\draw[bravH] (0.399,-5.815) -- (0.399,-4.465);
\draw[bravE] (0.399,-4.465) -- (1.349,-4.465);
\draw[bravE] (0.950,-6.100) -- (1.349,-5.815);
\draw[bravE] (0.950,-6.100) -- (0.950,-4.750);
\draw[bravE] (0.950,-4.750) -- (1.349,-4.465);
\draw[bravE] (1.349,-5.815) -- (1.349,-4.465);
\node[bravN] at (0.000,-6.100) {};
\node[bravN] at (0.000,-4.750) {};
\node[bravN] at (0.399,-5.815) {};
\node[bravN] at (0.399,-4.465) {};
\node[bravN] at (0.950,-6.100) {};
\node[bravN] at (0.950,-4.750) {};
\node[bravN] at (1.349,-5.815) {};
\node[bravN] at (1.349,-4.465) {};
\node[font=\small, anchor=north] at (0.750,-6.350) {тетрагональная $P$};
\draw[bravE] (3.400,-6.100) -- (4.350,-6.100);
\draw[bravH] (3.400,-6.100) -- (3.799,-5.815);
\draw[bravE] (3.400,-6.100) -- (3.400,-4.750);
\draw[bravE] (3.400,-4.750) -- (4.350,-4.750);
\draw[bravE] (3.400,-4.750) -- (3.799,-4.465);
\draw[bravH] (3.799,-5.815) -- (4.749,-5.815);
\draw[bravH] (3.799,-5.815) -- (3.799,-4.465);
\draw[bravE] (3.799,-4.465) -- (4.749,-4.465);
\draw[bravE] (4.350,-6.100) -- (4.749,-5.815);
\draw[bravE] (4.350,-6.100) -- (4.350,-4.750);
\draw[bravE] (4.350,-4.750) -- (4.749,-4.465);
\draw[bravE] (4.749,-5.815) -- (4.749,-4.465);
\node[bravN] at (3.400,-6.100) {};
\node[bravN] at (3.400,-4.750) {};
\node[bravN] at (3.799,-5.815) {};
\node[bravN] at (3.799,-4.465) {};
\node[bravN] at (4.350,-6.100) {};
\node[bravN] at (4.350,-4.750) {};
\node[bravN] at (4.749,-5.815) {};
\node[bravN] at (4.749,-4.465) {};
\node[bravC] at (4.074,-5.282) {};
\node[font=\small, anchor=north] at (4.150,-6.350) {тетрагональная $I$};
\draw[bravE] (6.800,-6.100) -- (7.900,-6.100);
\draw[bravH] (6.800,-6.100) -- (6.650,-5.814);
\draw[bravE] (6.800,-6.100) -- (6.800,-4.850);
\draw[bravE] (6.800,-4.850) -- (7.900,-4.850);
\draw[bravE] (6.800,-4.850) -- (6.650,-4.564);
\draw[bravH] (6.650,-5.814) -- (7.750,-5.814);
\draw[bravH] (6.650,-5.814) -- (6.650,-4.564);
\draw[bravE] (6.650,-4.564) -- (7.750,-4.564);
\draw[bravE] (7.900,-6.100) -- (7.750,-5.814);
\draw[bravE] (7.900,-6.100) -- (7.900,-4.850);
\draw[bravE] (7.900,-4.850) -- (7.750,-4.564);
\draw[bravE] (7.750,-5.814) -- (7.750,-4.564);
\node[bravN] at (6.800,-6.100) {};
\node[bravN] at (6.800,-4.850) {};
\node[bravN] at (6.650,-5.814) {};
\node[bravN] at (6.650,-4.564) {};
\node[bravN] at (7.900,-6.100) {};
\node[bravN] at (7.900,-4.850) {};
\node[bravN] at (7.750,-5.814) {};
\node[bravN] at (7.750,-4.564) {};
\node[font=\small, anchor=north] at (7.550,-6.350) {гексагональная $P$};
\draw[bravE] (0.000,-9.150) -- (1.100,-9.150);
\draw[bravH] (0.000,-9.150) -- (0.462,-8.820);
\draw[bravE] (0.000,-9.150) -- (0.000,-8.050);
\draw[bravE] (0.000,-8.050) -- (1.100,-8.050);
\draw[bravE] (0.000,-8.050) -- (0.462,-7.720);
\draw[bravH] (0.462,-8.820) -- (1.562,-8.820);
\draw[bravH] (0.462,-8.820) -- (0.462,-7.720);
\draw[bravE] (0.462,-7.720) -- (1.562,-7.720);
\draw[bravE] (1.100,-9.150) -- (1.562,-8.820);
\draw[bravE] (1.100,-9.150) -- (1.100,-8.050);
\draw[bravE] (1.100,-8.050) -- (1.562,-7.720);
\draw[bravE] (1.562,-8.820) -- (1.562,-7.720);
\node[bravN] at (0.000,-9.150) {};
\node[bravN] at (0.000,-8.050) {};
\node[bravN] at (0.462,-8.820) {};
\node[bravN] at (0.462,-7.720) {};
\node[bravN] at (1.100,-9.150) {};
\node[bravN] at (1.100,-8.050) {};
\node[bravN] at (1.562,-8.820) {};
\node[bravN] at (1.562,-7.720) {};
\node[font=\small, anchor=north] at (0.750,-9.400) {кубическая $P$};
\draw[bravE] (3.400,-9.150) -- (4.500,-9.150);
\draw[bravH] (3.400,-9.150) -- (3.862,-8.820);
\draw[bravE] (3.400,-9.150) -- (3.400,-8.050);
\draw[bravE] (3.400,-8.050) -- (4.500,-8.050);
\draw[bravE] (3.400,-8.050) -- (3.862,-7.720);
\draw[bravH] (3.862,-8.820) -- (4.962,-8.820);
\draw[bravH] (3.862,-8.820) -- (3.862,-7.720);
\draw[bravE] (3.862,-7.720) -- (4.962,-7.720);
\draw[bravE] (4.500,-9.150) -- (4.962,-8.820);
\draw[bravE] (4.500,-9.150) -- (4.500,-8.050);
\draw[bravE] (4.500,-8.050) -- (4.962,-7.720);
\draw[bravE] (4.962,-8.820) -- (4.962,-7.720);
\node[bravN] at (3.400,-9.150) {};
\node[bravN] at (3.400,-8.050) {};
\node[bravN] at (3.862,-8.820) {};
\node[bravN] at (3.862,-7.720) {};
\node[bravN] at (4.500,-9.150) {};
\node[bravN] at (4.500,-8.050) {};
\node[bravN] at (4.962,-8.820) {};
\node[bravN] at (4.962,-7.720) {};
\node[bravC] at (4.181,-8.435) {};
\node[font=\small, anchor=north] at (4.150,-9.400) {кубическая $I$};
\draw[bravE] (6.800,-9.150) -- (7.900,-9.150);
\draw[bravH] (6.800,-9.150) -- (7.262,-8.820);
\draw[bravE] (6.800,-9.150) -- (6.800,-8.050);
\draw[bravE] (6.800,-8.050) -- (7.900,-8.050);
\draw[bravE] (6.800,-8.050) -- (7.262,-7.720);
\draw[bravH] (7.262,-8.820) -- (8.362,-8.820);
\draw[bravH] (7.262,-8.820) -- (7.262,-7.720);
\draw[bravE] (7.262,-7.720) -- (8.362,-7.720);
\draw[bravE] (7.900,-9.150) -- (8.362,-8.820);
\draw[bravE] (7.900,-9.150) -- (7.900,-8.050);
\draw[bravE] (7.900,-8.050) -- (8.362,-7.720);
\draw[bravE] (8.362,-8.820) -- (8.362,-7.720);
\node[bravN] at (6.800,-9.150) {};
\node[bravN] at (6.800,-8.050) {};
\node[bravN] at (7.262,-8.820) {};
\node[bravN] at (7.262,-7.720) {};
\node[bravN] at (7.900,-9.150) {};
\node[bravN] at (7.900,-8.050) {};
\node[bravN] at (8.362,-8.820) {};
\node[bravN] at (8.362,-7.720) {};
\node[bravC] at (7.581,-8.985) {};
\node[bravC] at (7.581,-7.885) {};
\node[bravC] at (7.031,-8.435) {};
\node[bravC] at (8.131,-8.435) {};
\node[bravC] at (7.350,-8.600) {};
\node[bravC] at (7.812,-8.270) {};
\node[font=\small, anchor=north] at (7.550,-9.400) {кубическая $F$};

\end{tikzpicture}
\caption{Четырнадцать решёток Браве (расчёт). Оранжевым выделены
центрирующие узлы. Почему, например, нет тетрагональной $C$-решётки?
Потому что она сводится к примитивной тетрагональной $P$ с меньшей
ячейкой, повёрнутой на $45^\circ$; по той же причине тетрагональная $F$
сводится к $I$}
\label{fig:bravais}
\end{figure}

\subsection{Пространственные группы}\label{ssec:space-groups}

\begin{opredelenie}[Пространственная группа]
Совокупность всех преобразований кристаллического пространства,
переводящих каждую точку и каждое направление в эквивалентные.
Элемент группы~— преобразование
\begin{equation}\label{eq:seitz}
\vb r'=R\vb r+\vb t\equiv\{R\,|\,\vb t\}\,\vb r ,
\end{equation}
где $R$~— поворот (собственный или несобственный), $\vb t$~— перенос.
\end{opredelenie}

При $\det R=+1$ и $\vb t\ne0$ получается винтовое движение, при $\det R=-1$~—
отражение со сдвигом (скользящее отражение). Последовательное
применение двух операций:
\begin{equation}\label{eq:seitz-product}
\{R_2|\vb t_2\}\{R_1|\vb t_1\}\vb r=R_2(R_1\vb r+\vb t_1)+\vb t_2
\ \Rightarrow\ \{R_2|\vb t_2\}\{R_1|\vb t_1\}=\{R_2R_1\,|\,R_2\vb t_1+\vb t_2\}
\end{equation}
\ispr{В оригинале результат записан как «$R_2R_1+R_2\vb t_1+\vb t_2$», что
смешивает матрицу и вектор; произведение~— снова операция вида
$\{R|\vb t\}$ с $R=R_2R_1$ и $\vb t=R_2\vb t_1+\vb t_2$.}; единичный элемент~—
$\{E|\vb 0\}$, обратный~— $\{R|\vb t\}^{-1}=\{R^{-1}|-R^{-1}\vb t\}$. Операции $\{E|\vb t\}$
с векторами решётки образуют бесконечную \emph{трансляционную} подгруппу;
всего таких групп 14~— по числу решёток Браве.

\paragraph{Винтовые оси}
\term{Винтовая ось}~— ось симметрии и параллельный ей перенос,
действующие совместно. Обозначение $n_m$: поворот на $360^\circ/n$ и сдвиг на
$m/n$ периода вдоль оси. Например, $3_1$~— поворот на $120^\circ$ и сдвиг на
$\frac13$ периода, $3_2$~— сдвиг на $\frac23$ (рис.~\ref{fig:screw}). Оси $3_1$ и $3_2$
(так же как $4_1$ и $4_3$, $6_1$ и $6_5$)~— правая и левая разновидности, они
энантиоморфны; ось $2_1$ можно считать и правой, и левой. Действие оси
$2_1$ вдоль $\vb a_2$:
\[
\begin{pmatrix}x'\\y'\\z'\end{pmatrix}=
\begin{pmatrix}\br1&0&0\\0&1&0\\0&0&\br1\end{pmatrix}\begin{pmatrix}x\\y\\z\end{pmatrix}+
\begin{pmatrix}0\\\frac12\\0\end{pmatrix}=\begin{pmatrix}-x\\y+\frac12\\-z\end{pmatrix};
\]
двукратное применение даёт $(x,\ y+1,\ z)$~— трансляцию на период.

\paragraph{Плоскости скользящего отражения}
\term{Плоскость скользящего отражения}~— отражение в плоскости и
параллельный ей перенос, действующие совместно (рис.~\ref{fig:glide}).
Различают плоскости $a$, $b$, $c$ (сдвиг на половину периода вдоль
соответствующей оси), $n$ (сдвиг вдоль диагонали грани:
$(\vb a+\vb b)/2$ и т.\,п.) и «алмазную» $d$ (сдвиг на $(\vb a\pm\vb b)/4$,
$(\vb a\pm\vb b\pm\vb c)/4$ и т.\,п.). Для плоскости, перпендикулярной $\vb a_2$, со
сдвигом вдоль $\vb a_3$: $(x,y,z)\to(x,-y,z+\frac12)$, а двукратное применение
снова даёт трансляцию.

\begin{figure}[ht]
\centering
\begin{subfigure}[b]{0.56\textwidth}\centering
\begin{tikzpicture}[scale=1.0]
  \foreach \col/\lab/\m in {0/3/0, 2.3/3_1/1, 4.6/3_2/2}{
    \draw[fmGray, thick] (\col,-0.2) -- (\col,3.3);
    \foreach \k in {0,1,2,3}{
      \pgfmathsetmacro{\ang}{120*\k}
      \pgfmathsetmacro{\hz}{\m==0 ? 0 : mod(\m*\k,3)}
      \pgfmathsetmacro{\hh}{\m==0 ? 0 : \hz + (\k==3 ? 3 : 0)}
      \pgfmathsetmacro{\px}{\col+0.65*cos(\ang)}
      \pgfmathsetmacro{\py}{\hh+0.22*sin(\ang)}
      \ifnum\k<3
        \fill[fmBlue] (\px,\py) circle (0.08);
        \node[font=\tiny, anchor=west] at (\px+0.07,\py) {\pgfmathparse{int(\k+1)}\pgfmathresult};
      \fi}
    \ifnum\m>0
      \foreach \k in {0,1,2}{
        \pgfmathsetmacro{\ang}{120*\k}
        \pgfmathsetmacro{\py}{mod(\m*\k,3)+3+0.22*sin(\ang)}
        \pgfmathsetmacro{\px}{\col+0.65*cos(\ang)}
        \ifdim\py pt<3.35pt \fill[fmBlue!40] (\px,\py) circle (0.08);\fi}
    \fi
    \draw[fmGray, densely dotted] (\col-0.75,0) -- (\col+0.75,0);
    \draw[fmGray, densely dotted] (\col-0.75,3) -- (\col+0.75,3);
    \node[font=\small] at (\col,-0.55) {$\lab$};}
  \node[font=\scriptsize, anchor=east] at (-0.8,0) {$0$};
  \node[font=\scriptsize, anchor=east] at (-0.8,1) {$\frac13$};
  \node[font=\scriptsize, anchor=east] at (-0.8,2) {$\frac23$};
  \node[font=\scriptsize, anchor=east] at (-0.8,3) {$1$};
  \foreach \yy in {1,2}{\draw[fmGray!50, densely dotted] (-0.75,\yy) -- (5.35,\yy);}
\end{tikzpicture}
\caption{}
\end{subfigure}\hfill
\begin{subfigure}[b]{0.4\textwidth}\centering
\begin{tikzpicture}[scale=0.95]
  \draw[fmRed, thick, dashed] (-0.3,0) -- (5.0,0) node[right, font=\small]{$c$};
  \foreach \i in {0,1,2}{
    \begin{scope}[xshift=\i*1.8cm]
      \draw[fmBlue, thick, fill=fmBlue!20] (0,0.35) -- (0,1.0) -- (0.45,1.0) -- (0.45,0.85) -- (0.15,0.85) -- (0.15,0.35) -- cycle;
    \end{scope}
    \begin{scope}[xshift=\i*1.8cm+0.9cm]
      \draw[fmGreen!60!black, thick, fill=fmGreen!20] (0,-0.35) -- (0,-1.0) -- (0.45,-1.0) -- (0.45,-0.85) -- (0.15,-0.85) -- (0.15,-0.35) -- cycle;
    \end{scope}}
  \draw[dimline] (0,-1.35) -- node[below, font=\scriptsize]{$\vb t=\vb c/2$} (0.9,-1.35);
\end{tikzpicture}
\caption{}
\end{subfigure}
\caption{Открытые элементы симметрии: \textit{а}~— поворотная ось 3 и
винтовые оси $3_1$ и $3_2$ (вид сбоку; числа~— порядок, в котором точки
получаются из первой; бледные~— точки следующего периода); \textit{б}~—
плоскость скользящего отражения $c$ (штрих): фигура отражается в
плоскости и сдвигается на половину периода}
\label{fig:screw}\label{fig:glide}
\end{figure}

\paragraph{230 пространственных групп}
Перебрав все сочетания открытых и закрытых элементов симметрии,
Е.\,С.~Фёдоров и А.~Шёнфлис (независимо, 1890--1891) получили \vazhno{230
пространственных групп}. В 73 \emph{симморфных} группах точечная группа
является подгруппой пространственной; в остальных (несимморфных) этого
нет из-за винтовых осей и плоскостей скольжения.

\paragraph{Правильные системы точек}
Совокупность точек, получаемых из одной всеми операциями
пространственной группы,~— \term{правильная система точек} (аналог
простой формы для точечной группы). Если исходная точка лежит на
каком-либо элементе симметрии (или на равных расстояниях от одинаковых
элементов), система \emph{частная}, иначе~— \emph{общая}.
\term{Кратность}~— число точек системы в элементарной ячейке; у общей
системы она больше, чем у частных. Атомы одного сорта в структуре
занимают одну или несколько правильных систем точек.

\paragraph{Международный символ пространственной группы}
Первая буква~— тип решётки Браве; далее~— порождающие элементы
симметрии, каждый на своём месте (перестановка меняет смысл). В
кубической сингонии на третьем месте всегда стоит 3 (оси 3 вдоль
$\nfam{111}$), на втором~— элементы, связанные с $\nfam{100}$, на
четвёртом~— с $\nfam{110}$. Примеры: медь~— $Fm\br3m$, вольфрам~— $Im\br3m$,
алмаз~— $Fd\br3m$ (вдоль $\nfam{100}$ проходит алмазная плоскость $d$),
NaCl~— $Fm\br3m$, CsCl~— $Pm\br3m$.

\paragraph{Антисимметрия и цветная симметрия}
А.\,В.~Шубников ввёл операцию \emph{антисимметрии}, меняющую «знак»
объекта (как электрон и позитрон, негатив и позитив). Чёрно-белые группы
хорошо описывают магнитные структуры (направление спина~— «знак»): к
32 обычным точечным группам добавляется 58 чёрно-белых~— всего 90
магнитных классов; пространственных шубниковских групп 1651. Обобщение на
многоцветные группы предложил Н.\,В.~Белов.

\paragraph{Двумерная и многомерная кристаллография. Квазикристаллы}
На плоскости возможны те же оси 1, 2, 3, 4, 6; сингоний четыре
(косоугольная, прямоугольная, квадратная, гексагональная), решёток Браве
пять (в прямоугольной есть центрированная), плоских групп симметрии~—
17. В четырёхмерной кристаллографии разрешены оси 5, 8, 10 и 12-го
порядков; решёток Браве там 64, точечных групп 227, пространственных
около 4\,800.

Упорядоченные структуры делят на \emph{периодические} (кристаллы) и
\emph{квазипериодические}, у которых обратная решётка квазипериодична.
Среди последних~— \emph{несоразмерные} кристаллы (кристаллографическая
точечная группа) и \term{квазикристаллы} (некристаллографическая точечная
группа, например с осями 5-го порядка). Двумерный пример~— мозаика Пенроуза;
трёхмерный~— икосаэдрические сплавы Al--Mn (Д.~Шехтман, 1984;
Нобелевская премия 2011~г.).

\paragraph{Как симметрию определяют на опыте}
Атомное строение изучают по дифракции рентгеновских лучей, электронов и
нейтронов: длины волн сравнимы с межатомными расстояниями, и кристалл
служит дифракционной решёткой. Максимумы возникают в направлениях,
удовлетворяющих условию Вульфа--Брэгга
\begin{keyeq}\label{eq:bragg}
2d\sin\theta=n\lambda ,
\end{keyeq}
где $d$~— межплоскостное расстояние, $\theta$~— угол скольжения, $n$~— порядок
отражения, $\lambda$~— длина волны.
\begin{itemize}
\item \emph{Метод Лауэ} (неподвижный монокристалл, «белое» излучение):
определяют ориентацию и элементы симметрии; по совпадению точек первого и
второго рода~— сингонию, по симметрии лауэграммы~— порядок осей.
\item По \emph{закону Фриделя} интенсивности отражений $(hkl)$ и
$(\br h\br k\br l)$ равны: дифракционная картина имеет центр симметрии, даже
если у кристалла его нет. Поэтому классы, различающиеся только центром
симметрии, дают одинаковую картину~— отсюда 11 лауэвских классов.
\item \emph{Метод вращения} (монохроматический пучок, кристалл вращается
вокруг оси зоны): находят период идентичности вдоль оси; сравнивая периоды
вдоль $\nfam{100}$, $\nfam{110}$, $\nfam{111}$, определяют тип решётки Браве.
\item \emph{Метод Дебая--Шеррера}~— для поликристаллов (монохроматическое
излучение).
\end{itemize}
Сначала по геометрии дифракционной картины определяют решётку, затем по
интенсивностям~— расположение атомов в ячейке. По диффузному рассеянию
изучают также упругие волны, модули упругости, фононный спектр.

\subsection{Предельные группы симметрии}\label{ssec:curie}

Для физических полей, тел вращения и многих свойств кристаллов нужны
группы с осями \emph{бесконечного} порядка (поворот на любой угол).
Это \term{предельные группы}, или группы Кюри; П.~Кюри показал, что их
семь (рис.~\ref{fig:curie}, табл.~\ref{tab:curie}). Все 32 кристаллографические
точечные группы~— подгруппы предельных.

\begin{figure}[ht]
\centering
\tikzset{cbody/.style={fill=fmBlue!12, draw=atomLine, line width=0.6pt},
  crot/.style={-{Stealth[length=1.6mm]}, fmRed, line width=0.7pt}}
\begin{tikzpicture}[scale=0.82]
  % а: вращающийся конус
  \begin{scope}[xshift=0cm]
    \draw[cbody] (-0.7,0) -- (0,2.0) -- (0.7,0) arc (0:-180:0.7 and 0.22);
    \draw[atomLine, line width=0.6pt, dashed] (0.7,0) arc (0:180:0.7 and 0.22);
    \draw[crot] (-0.45,1.25) arc (190:520:0.45 and 0.13);
    \node[font=\small] at (0,-0.75) {\textit{а}~$\infty$};
  \end{scope}
  % б: неподвижный конус
  \begin{scope}[xshift=2.2cm]
    \draw[cbody] (-0.7,0) -- (0,2.0) -- (0.7,0) arc (0:-180:0.7 and 0.22);
    \draw[atomLine, line width=0.6pt, dashed] (0.7,0) arc (0:180:0.7 and 0.22);
    \node[font=\small] at (0,-0.75) {\textit{б}~$\infty m$};
  \end{scope}
  % в: вращающийся цилиндр
  \begin{scope}[xshift=4.4cm]
    \draw[cbody] (-0.6,0) -- (-0.6,1.8) arc (180:360:0.6 and 0.18) -- (0.6,0) arc (0:-180:0.6 and 0.18);
    \draw[cbody] (0,1.8) ellipse (0.6 and 0.18);
    \draw[crot] (-0.6,0.9) arc (180:500:0.6 and 0.18);
    \node[font=\small] at (0,-0.75) {\textit{в}~$\infty/m$};
  \end{scope}
  % г: закрученный цилиндр
  \begin{scope}[xshift=6.6cm]
    \draw[cbody] (-0.6,0) -- (-0.6,1.8) arc (180:360:0.6 and 0.18) -- (0.6,0) arc (0:-180:0.6 and 0.18);
    \draw[cbody] (0,1.8) ellipse (0.6 and 0.18);
    \foreach \x in {-0.4,0,0.4}{\draw[fmRed, line width=0.5pt] (\x-0.15,0) .. controls (\x,0.6) and (\x,1.2) .. (\x+0.15,1.8);}
    \node[font=\small] at (0,-0.75) {\textit{г}~$\infty2$};
  \end{scope}
  % д: неподвижный цилиндр
  \begin{scope}[xshift=8.8cm]
    \draw[cbody] (-0.6,0) -- (-0.6,1.8) arc (180:360:0.6 and 0.18) -- (0.6,0) arc (0:-180:0.6 and 0.18);
    \draw[cbody] (0,1.8) ellipse (0.6 and 0.18);
    \node[font=\small] at (0,-0.75) {\textit{д}~$\infty/mmm$};
  \end{scope}
  % е: закрученный шар
  \begin{scope}[xshift=11.0cm, yshift=0.9cm]
    \shade[ball color=fmBlue!25] (0,0) circle (0.85);
    \draw[atomLine, line width=0.6pt] (0,0) circle (0.85);
    \foreach \a in {-50,0,50}{\draw[fmRed, line width=0.5pt, rotate=\a] (-0.6,-0.6) .. controls (-0.2,0.1) and (0.2,-0.1) .. (0.6,0.6);}
    \node[font=\small] at (0,-1.65) {\textit{е}~$\infty\infty$};
  \end{scope}
  % ж: шар
  \begin{scope}[xshift=13.2cm, yshift=0.9cm]
    \shade[ball color=fmBlue!25] (0,0) circle (0.85);
    \draw[atomLine, line width=0.6pt] (0,0) circle (0.85);
    \node[font=\small] at (0,-1.65) {\textit{ж}~$\infty\infty m$};
  \end{scope}
\end{tikzpicture}
\caption{Фигуры, изображающие предельные группы Кюри: \textit{а}~—
вращающийся конус; \textit{б}~— неподвижный конус; \textit{в}~— вращающийся
цилиндр; \textit{г}~— цилиндр с закрученными в разные стороны концами;
\textit{д}~— неподвижный цилиндр; \textit{е}~— шар с закрученными
диаметрами; \textit{ж}~— обычный шар}
\label{fig:curie}
\end{figure}

\begin{table}[ht]
\centering
\caption{Предельные группы симметрии и их физические реализации}\label{tab:curie}
\small
\begin{tabularx}{\textwidth}{@{}lp{30mm}>{\raggedright\arraybackslash}p{28mm}>{\raggedright\arraybackslash\hspace{0pt}}X@{}}
\toprule
Группа & Фигура & Свойства & Пример \\
\midrule
$\infty$ & вращающийся конус & полярна, энантиоморфна & предел групп 1, 2, 3, 4, 6 \\
$\infty m$ & неподвижный конус & полярна & однородное электрическое поле (полярный вектор $\vb E$) \\
$\infty/m$ & вращающийся цилиндр & центр симметрии & магнитное поле (аксиальный вектор $\vb H$) \\
$\infty2$ & закрученный цилиндр & энантиоморфна & вращение плоскости поляризации в анизотропной среде \\
$\infty/mmm$ & неподвижный цилиндр & центр симметрии & одноосное растяжение или сжатие \\
$\infty\infty$ & закрученный шар & энантиоморфна & вращение плоскости поляризации в изотропной среде \\
$\infty\infty m$ & шар & все операции & гидростатическое давление, однородный нагрев (скаляр) \\
\bottomrule
\end{tabularx}
\end{table}

Группа $\infty\infty m$ содержит все точечные операции, оставляющие на месте
шар,~— её называют \emph{ортогональной} группой. Предельные группы особенно
полезны в кристаллофизике: они описывают симметрию внешних воздействий
(п.~\ref{ssec:symmetry-principle}).

\begin{itogi}
\begin{itemize}
\item Трансляционная симметрия допускает только оси 1, 2, 3, 4, 6
($\operatorname{tr}R=1+2\cos\alpha$~— целое).
\item Инверсионные оси: $\br1$~— центр, $\br2=m$~— плоскость, $\br3=3+\br1$,
$\br6=3/m$, $\br4$~— независимый элемент.
\item Группа: замкнутость, ассоциативность, единица, обратный элемент;
подгруппы, классы, генераторы, таблица Кэли.
\item 32 точечные группы (11 лауэвских, 11 энантиоморфных); 14 решёток
Браве; 230 пространственных групп (винтовые оси $n_m$, плоскости
скольжения $a$, $b$, $c$, $n$, $d$).
\item 7 предельных групп Кюри описывают симметрию полей и воздействий:
электрическое поле~— $\infty m$, магнитное~— $\infty/m$, одноосное напряжение~—
$\infty/mmm$.
\end{itemize}
\end{itogi}

\begin{proverka}
\item Почему в кристаллах нет оси 5-го порядка, а в квазикристаллах она
есть?
\item Чему равен след матрицы поворота вокруг оси 6? Оси $\br4$?
\item Покажите, что ось $\br2$ эквивалентна плоскости симметрии, а $\br6$~—
сочетанию оси 3 и перпендикулярной плоскости.
\item Каков порядок группы $m\br3m$? Сколько у неё подгрупп с порядком 5?
\item Может ли пьезоэлектриком быть кристалл класса $m\br3m$? Класса $\br43m$?
\item Почему нет тетрагональной $C$-решётки Браве?
\item Чем винтовая ось $4_1$ отличается от $4_3$?
\item Какой предельной группой описывается симметрия однородного
магнитного поля и почему она не совпадает с симметрией электрического?
\end{proverka}


\section{Структура кристаллов}\label{sec:structure}
% ----------------------------------------------------------------------
%  1.4. Структура кристаллов
% ----------------------------------------------------------------------
\begin{orientir}
Как описывают конкретные кристаллические структуры. Главная идея~—
\vazhno{плотнейшие упаковки шаров и их пустоты}: большинство структур
металлов и многих соединений~— это плотная упаковка одних атомов, в
пустотах которой сидят другие. Разберём структурные типы элементов
(1.4.1--1.4.2), соединений AB, AB$_2$ и сложных (1.4.3--1.4.5),
сверхструктуры (1.4.6), фуллерены и квазикристаллы (1.4.7) и структуру
поверхности (1.4.8).
\end{orientir}

Чтобы описать структуру, указывают пространственную группу, координаты
атомов в ячейке, координационные числа и многогранники.

\begin{opredelenie}[Координация]
\term{Координационное число} (к.\,ч.)~— число ближайших атомов (ионов)
данного сорта, находящихся на одинаковом расстоянии от центрального.
\term{Координационный многогранник}~— многогранник, вершины которого
заняты этими соседями; число его вершин равно к.\,ч.
\end{opredelenie}

Кристаллы, структуры которых одинаковы с точностью до подобия, относят к
одному \term{структурному типу}; его называют по типичному представителю
(«тип меди», «тип NaCl»). В символике Strukturbericht буква указывает
стехиометрию: A~— элементы, B~— соединения AB, C~— AB$_2$, D~—
A$_m$B$_n$, E~— более сложные, L~— сверхструктуры; за буквой следует номер
(A1~— тип меди, B1~— тип NaCl и~т.\,д.).

\subsection{Плотнейшие упаковки}\label{ssec:packings}

Будем строить структуру из одинаковых несжимаемых шаров. В плотном
плоском слое (рис.~\ref{fig:cp-layer}) каждый шар касается шести соседей и
окружён шестью лунками, а каждая лунка~— тремя шарами. Лунок вдвое больше,
чем шаров; они двух сортов~— B (треугольники вершиной вниз) и C (вершиной
вверх). Следующий слой можно положить либо в лунки B, либо в лунки C.

\begin{figure}[ht]
\centering
\begin{subfigure}[b]{0.5\textwidth}\centering
\begin{tikzpicture}[scale=0.62]
  \foreach \i in {-3,...,3}{\foreach \j in {-2,...,2}{
    \pgfmathsetmacro{\px}{\i+0.5*\j}\pgfmathsetmacro{\py}{0.866*\j}
    \fill[atomA, draw=atomLine, line width=0.3pt] (\px,\py) circle (0.5);}}
  \foreach \i in {-2,...,2}{\foreach \j in {-1,0,1}{
    \pgfmathsetmacro{\px}{\i+0.5*\j+0.5}\pgfmathsetmacro{\py}{0.866*\j+0.2887}
    \node[font=\scriptsize\bfseries, fmRed] at (\px,\py) {B};
    \pgfmathsetmacro{\qx}{\i+0.5*\j+1.0}\pgfmathsetmacro{\qy}{0.866*\j+0.5774}
    \node[font=\scriptsize\bfseries, fmBlue] at (\qx,\qy) {C};}}
  \node[font=\scriptsize\bfseries] at (0,0) {A};
\end{tikzpicture}
\caption{}
\end{subfigure}\hfill
\begin{subfigure}[b]{0.46\textwidth}\centering
\begin{tikzpicture}[x=1cm, y=0.55cm]
  \foreach \k/\l in {0/A,1/B,2/A,3/B,4/A}{
    \node[font=\small] at (0,\k) {\l};}
  \foreach \k/\l in {0/A,1/B,2/C,3/A,4/B}{
    \node[font=\small] at (2.6,\k) {\l};}
  \draw[fmGray] (-0.5,-0.6) -- (-0.5,4.6) (2.1,-0.6) -- (2.1,4.6);
  \node[font=\small, align=center] at (0,-1.6) {ГПУ\\$\ldots ABAB\ldots$};
  \node[font=\small, align=center] at (2.6,-1.6) {ГЦК\\$\ldots ABCABC\ldots$};
\end{tikzpicture}
\caption{}
\end{subfigure}
\caption{Плотнейшие упаковки: \textit{а}~— плотный слой шаров A и лунки
двух сортов B и C; \textit{б}~— две последовательности укладки слоёв:
двухслойная (гексагональная) и трёхслойная (кубическая)}
\label{fig:cp-layer}
\end{figure}

Во втором слое пустоты уже двух типов:
\begin{itemize}
\item \term{тетраэдрическая} ($T$)~— над лункой первого слоя лежит шар
второго; пустота окружена четырьмя шарами по вершинам тетраэдра;
\item \term{октаэдрическая} ($O$)~— лунка второго слоя над лункой первого;
пустота окружена шестью шарами по вершинам октаэдра.
\end{itemize}
Октаэдрических пустот столько же, сколько шаров, тетраэдрических~— вдвое
больше. В октаэдрическую пустоту помещается шар радиусом $0{,}414R$, в
тетраэдрическую~— $0{,}225R$ ($R$~— радиус шаров упаковки).

Третий слой можно положить в лунки $T$ (над шарами первого слоя)~—
получится \emph{двухслойная} упаковка $\ldots ABAB\ldots$,~— или в лунки $O$~—
\emph{трёхслойная} $\ldots ABCABC\ldots$ Возможны и другие чередования (кроме
повторения одной буквы), но плотнейших всего две.

\begin{glavnoe}[Две плотнейшие упаковки]
\begin{itemize}
\item \emph{Гексагональная} (ГПУ, $\ldots ABAB\ldots$): к.\,ч.~=~12 (6 в слое, по 3
сверху и снизу); идеальное отношение осей
$c/a=\sqrt{8/3}=1{,}633$; октаэдрические пустоты образуют сквозные каналы
вдоль оси $c$ (пути быстрой диффузии примесей).
\item \emph{Кубическая} ($\ldots ABCABC\ldots$): шары~— в узлах ГЦК-решётки,
плотные слои~— $\{111\}$, плотнейшие ряды~— $\nfam{110}$, к.\,ч.~=~12;
сквозных каналов нет.
\end{itemize}
Коэффициент компактности (доля объёма, занятая шарами) у обеих
$\pi/\sqrt{18}=74{,}05\,\%$; у ОЦК~— $68\,\%$, у простой кубической~— $52\,\%$.
\end{glavnoe}

Идея плотнейших упаковок очень плодотворна: в большинстве структур более
крупные атомы образуют плотную (или почти плотную) упаковку, а структуры
различаются тем, какие и сколько пустот заняты другими атомами.

\subsection{Структурные типы химических элементов}\label{ssec:elem-types}

\paragraph{Тип меди (A1, ГЦК)}
Al, Ni, Pb, $\gamma$-Fe, Au, Ag, Ir, Pt и многие другие. Пространственная
группа $Fm\br3m$; базис: $0,0,0$; $\frac12,\frac12,0$; $\frac12,0,\frac12$;
$0,\frac12,\frac12$ (кратко: $0,0,0$; $\frac12,\frac12,0\,[3]$, где
«[3]» означает три циклические перестановки). Атомов в ячейке $Z=4$, к.\,ч.~=~12,
координационный многогранник~— \emph{кубооктаэдр}. Это кубическая
плотнейшая упаковка слоёв $\{111\}$.

Пустоты ГЦК-структуры (рис.~\ref{fig:fcc-voids}):
\begin{itemize}
\item \emph{октаэдрические}~— в центре куба $\frac12,\frac12,\frac12$ и на серединах
рёбер $\frac12,0,0\,[3]$; на ячейку $1+12/4=4$; радиус
$r=(\sqrt2-1)R=0{,}414R$;
\item \emph{тетраэдрические}~— в центрах восьми «малых кубиков»
$\frac14,\frac14,\frac14$ и т.\,д.; на ячейку 8; $r=(\sqrt{3/2}-1)R=0{,}225R$.
\end{itemize}
Атомы : октапустоты : тетрапустоты $=1:1:2$, размеры
$1:0{,}414:0{,}225$. Многие соединения строятся как плотнейшая упаковка
крупных анионов, в пустотах которой сидят мелкие катионы.

\begin{figure}[ht]
\centering
\begin{subfigure}[b]{0.3\textwidth}\centering
\begin{tikzpicture}[cubeproj, scale=2.2]
  \CubeEdges[atomLine!60]
  \FCCback[2.6mm] \FCCmid[2.6mm] \FCCfront[2.6mm]
  \foreach \p in {(0.5,0.5,0.5),(0.5,0,0),(0.5,1,0),(0.5,0,1),(0.5,1,1),(0,0.5,0),(1,0.5,0),(0,0.5,1),(1,0.5,1),(0,0,0.5),(1,0,0.5),(0,1,0.5),(1,1,0.5)}
    \node[circle, fill=fmRed, inner sep=0pt, minimum size=4.2pt] at \p {};
\end{tikzpicture}
\caption{октаэдрические}
\end{subfigure}\hspace{6mm}
\begin{subfigure}[b]{0.3\textwidth}\centering
\begin{tikzpicture}[cubeproj, scale=2.2]
  \CubeEdges[atomLine!60]
  \FCCback[2.6mm] \FCCmid[2.6mm] \FCCfront[2.6mm]
  \foreach \x in {0.25,0.75}{\foreach \y in {0.25,0.75}{\foreach \z in {0.25,0.75}{
    \node[circle, fill=fmGreen!70!black, inner sep=0pt, minimum size=4.2pt] at (\x,\y,\z) {};}}}
\end{tikzpicture}
\caption{тетраэдрические}
\end{subfigure}\hspace{6mm}
\begin{subfigure}[b]{0.3\textwidth}\centering
\begin{tikzpicture}[cubeproj, scale=2.2]
  \CubeEdges[atomLine!60]
  \fill[fmBlue!20, draw=fmBlue, line width=0.5pt] (1,0,0) -- (0,1,0) -- (0,0,1) -- cycle;
  \fill[fmBlue!10, draw=fmBlue!60, line width=0.4pt, dashed] (1,1,0) -- (0,1,1) -- (1,0,1) -- cycle;
  \FCCback[2.6mm] \FCCmid[2.6mm] \FCCfront[2.6mm]
\end{tikzpicture}
\caption{плотные слои $(111)$}
\end{subfigure}
\caption{ГЦК-структура: \textit{а}~— октаэдрические пустоты (центр куба и
середины рёбер); \textit{б}~— тетраэдрические пустоты; \textit{в}~—
плотноупакованные плоскости $(111)$}
\label{fig:fcc-voids}
\end{figure}

\paragraph{Тип $\alpha$-вольфрама (A2, ОЦК)}
Cr, V, Mo, Nb, Ta, $\alpha$-Fe, щелочные металлы, высокотемпературные
$\beta$-Ti, $\beta$-Zr, $\beta$-Hf. Группа $Im\br3m$; базис $0,0,0$; $\frac12,\frac12,\frac12$;
$Z=2$; к.\,ч.~=~8 (многогранник~— куб). Вторая координационная сфера
(6~атомов на расстоянии $a$) лишь на 15\,\% дальше первой ($a\sqrt3/2$), поэтому
часто говорят о к.\,ч.~=~8+6=14. Пустоты ОЦК-структуры
(рис.~\ref{fig:bcc-voids}):
\begin{itemize}
\item \emph{тетраэдрические} (самые крупные!)~— в точках $\frac12,\frac14,0$ и
эквивалентных, в центрах слегка искажённых тетраэдров; их 12 на ячейку
(по 6 на атом); $r=(\sqrt{5/3}-1)R=0{,}291R$;
\item \emph{октаэдрические}~— в центрах граней и серединах рёбер; 6 на
ячейку ($6\cdot\frac12+12\cdot\frac14$); пустота~— сплюснутый октаэдр,
$r=(2/\sqrt3-1)R=0{,}155R$.
\end{itemize}
Атомы : октапустоты : тетрапустоты $=1:3:6$. Хотя тетрапустоты крупнее,
углерод в $\alpha$-железе сидит в \emph{октаэдрических}: сплюснутый октаэдр
раздвигает лишь два ближайших атома (вдоль оси сплющивания), и это
энергетически выгоднее; отсюда тетрагональность мартенсита
(п.~\ref{ssec:martensite-cryst}).

\begin{figure}[ht]
\centering
\begin{tikzpicture}[cubeproj, scale=2.6]
  \CubeEdges[atomLine!60]
  \foreach \p in {(0,0,0),(1,0,0),(0,1,0),(1,1,0),(0,0,1),(1,0,1),(0,1,1),(1,1,1),(0.5,0.5,0.5)}
    \node[atom=3.6mm] at \p {};
  \foreach \p in {(0.5,0.5,0),(0.5,0.5,1),(0.5,0,0.5),(0.5,1,0.5),(0,0.5,0.5),(1,0.5,0.5),
                  (0.5,0,0),(0.5,1,0),(0,0.5,0),(1,0.5,0),(0.5,0,1),(0.5,1,1),(0,0.5,1),(1,0.5,1),(0,0,0.5),(1,0,0.5),(0,1,0.5),(1,1,0.5)}
    \node[circle, fill=fmRed, inner sep=0pt, minimum size=4pt] at \p {};
  \foreach \p in {(0.5,0.25,0),(0.5,0.75,0),(0.25,0.5,0),(0.75,0.5,0),(0.5,0,0.25),(0.5,0,0.75),(0.25,0,0.5),(0.75,0,0.5),
                  (1,0.5,0.25),(1,0.5,0.75),(1,0.25,0.5),(1,0.75,0.5)}
    \node[regular polygon, regular polygon sides=3, fill=fmGreen!70!black, inner sep=0pt, minimum size=5.5pt] at \p {};
\end{tikzpicture}
\caption{ОЦК-структура: октаэдрические пустоты (красные кружки~— центры
граней и середины рёбер) и тетраэдрические пустоты (зелёные треугольники;
показаны только на трёх видимых гранях)}
\label{fig:bcc-voids}
\end{figure}

\paragraph{Тип магния (A3, ГПУ)}
Be, Mg, Zn, Cd, Re, Os, $\alpha$-Ti, $\alpha$-Zr, $\alpha$-Hf, Co и др. Группа
$P6_3/mmc$, решётка Браве гексагональная примитивная; в ячейке 2 атома:
$0,0,0$ и $\frac23,\frac13,\frac12$ (рис.~\ref{fig:hcp}); к.\,ч.~=~12
(многогранник~— гексагональный, или «двойниковый», кубооктаэдр).
Октаэдрические пустоты: $\frac13,\frac23,\frac14$ и $\frac13,\frac23,\frac34$ (2 на ячейку,
$r=0{,}414R$); тетраэдрические: $0,0,\frac38$; $0,0,\frac58$; $\frac23,\frac13,\frac18$;
$\frac23,\frac13,\frac78$ (4 на ячейку, $r=0{,}225R$)~— соотношения те же, что в ГЦК.
У большинства ГПУ-металлов $c/a=1{,}57$--$1{,}63$, близко к идеальному;
исключения~— цинк (1,86) и кадмий (1,89).

\begin{figure}[ht]
\centering
\begin{tikzpicture}[cubeproj, scale=1.25]
  \def\c{1.633}
  \foreach \k in {0,...,5}{
    \pgfmathsetmacro{\xa}{cos(60*\k)}\pgfmathsetmacro{\ya}{sin(60*\k)}
    \pgfmathsetmacro{\xb}{cos(60*\k+60)}\pgfmathsetmacro{\yb}{sin(60*\k+60)}
    \pgfmathtruncatemacro{\back}{(\k==1 || \k==2) ? 1 : 0}
    \ifnum\back=1
      \draw[atomLine!50, dashed] (\xa,\ya,0) -- (\xb,\yb,0);
    \else
      \draw[atomLine] (\xa,\ya,0) -- (\xb,\yb,0);
    \fi
    \draw[atomLine] (\xa,\ya,\c) -- (\xb,\yb,\c);
    \ifnum\k=2 \draw[atomLine!50, dashed] (\xa,\ya,0) -- (\xa,\ya,\c);
    \else \draw[atomLine] (\xa,\ya,0) -- (\xa,\ya,\c);\fi}
  \foreach \k in {0,...,5}{\node[atom=3.4mm] at ({cos(60*\k)},{sin(60*\k)},0) {};}
  \node[atom=3.4mm] at (0,0,0) {};
  \foreach \k in {0,1,2}{\node[atomB=3.4mm] at ({0.577*cos(120*\k+30)},{0.577*sin(120*\k+30)},{\c/2}) {};}
  \foreach \k in {0,...,5}{\node[atom=3.4mm] at ({cos(60*\k)},{sin(60*\k)},\c) {};}
  \node[atom=3.4mm] at (0,0,\c) {};
  \node[font=\small, anchor=west] at (1.15,0,\c) {A};
  \node[font=\small, anchor=west] at (1.15,0,{\c/2}) {B};
  \node[font=\small, anchor=west] at (1.15,0,0) {A};
\end{tikzpicture}
\caption{ГПУ-структура (тип магния): гексагональная призма из трёх
элементарных ячеек; слои A (голубые атомы) и B (оранжевые) чередуются
вдоль оси $c$}
\label{fig:hcp}
\end{figure}

\paragraph{Тип алмаза (A4)}
Алмаз, Si, Ge, серое олово ($\alpha$-Sn). Группа $Fd\br3m$; решётка Браве ГЦК,
базис из двух атомов $0,0,0$ и $\frac14,\frac14,\frac14$, т.\,е. атомы занимают узлы
ГЦК-решётки и половину тетраэдрических пустот (через одну)
(рис.~\ref{fig:diamond-graphite}\,\textit{а}). Иначе: две ГЦК-решётки,
сдвинутые на $\frac14$ пространственной диагонали. $Z=8$, к.\,ч.~=~4,
многогранник~— тетраэдр\ispr{В оригинале для алмаза указано
к.\,ч.~=~12; при тетраэдрическом координационном многограннике
к.\,ч.~=~4.}. Центр симметрии лежит посередине связи между соседними
атомами. Гексагональная модификация~— \emph{лонсдейлит}: двухслойная
упаковка, в которой атомами заполнена половина тетрапустот.

\paragraph{Тип графита (A9)}
Группа $P6_3/mmc$, гексагональная решётка, 4 атома в ячейке, к.\,ч.~=~3
(многогранник~— правильный треугольник). Структура слоистая
(рис.~\ref{fig:diamond-graphite}\,\textit{б}): внутри слоя прочные ковалентные
связи (расстояние 0,142~нм), между слоями~— слабые ван-дер-ваальсовы
(0,335~нм). Отсюда огромная анизотропия: электропроводность в плоскости
слоёв в $\sim10^4$--$10^5$ раз больше, чем вдоль оси $c$\ispr{В оригинале
сказано наоборот: «электропроводность вдоль оси $c$ в $10^5$ раз больше, чем в
поперечных направлениях». Проводимость графита велика именно в плоскости
слоёв.}, а лёгкое скольжение по $(0001)$ делает графит смазкой. Есть и
ромбоэдрический политип с трёхслойной укладкой ($R\br3m$); ещё одна
форма углерода~— линейный полимер \emph{карбин}.

\begin{figure}[ht]
\centering
\begin{subfigure}[b]{0.4\textwidth}\centering
\begin{tikzpicture}[cubeproj, scale=1.9, anion/.style={atom=3.2mm}, cation/.style={atom=3.2mm}]
  \CubeEdges
  % сгенерировано tools/gen_ionic_cells.py — не править вручную
\draw[bondthick] (0.25,0.25,0.25) -- (0,0,0);
\draw[bondthick] (0.25,0.25,0.25) -- (0,0.5,0.5);
\draw[bondthick] (0.25,0.25,0.25) -- (0.5,0,0.5);
\draw[bondthick] (0.25,0.25,0.25) -- (0.5,0.5,0);
\draw[bondthick] (0.75,0.75,0.25) -- (0.5,0.5,0);
\draw[bondthick] (0.75,0.75,0.25) -- (0.5,1,0.5);
\draw[bondthick] (0.75,0.75,0.25) -- (1,0.5,0.5);
\draw[bondthick] (0.75,0.75,0.25) -- (1,1,0);
\draw[bondthick] (0.75,0.25,0.75) -- (0.5,0,0.5);
\draw[bondthick] (0.75,0.25,0.75) -- (0.5,0.5,1);
\draw[bondthick] (0.75,0.25,0.75) -- (1,0,1);
\draw[bondthick] (0.75,0.25,0.75) -- (1,0.5,0.5);
\draw[bondthick] (0.25,0.75,0.75) -- (0,0.5,0.5);
\draw[bondthick] (0.25,0.75,0.75) -- (0,1,1);
\draw[bondthick] (0.25,0.75,0.75) -- (0.5,0.5,1);
\draw[bondthick] (0.25,0.75,0.75) -- (0.5,1,0.5);
\node[anion] at (0,1,0) {};
\node[anion] at (1,1,0) {};
\node[anion] at (0.5,1,0.5) {};
\node[anion] at (0,1,1) {};
\node[anion] at (1,1,1) {};
\node[cation] at (0.75,0.75,0.25) {};
\node[cation] at (0.25,0.75,0.75) {};
\node[anion] at (0.5,0.5,0) {};
\node[anion] at (0,0.5,0.5) {};
\node[anion] at (1,0.5,0.5) {};
\node[anion] at (0.5,0.5,1) {};
\node[cation] at (0.25,0.25,0.25) {};
\node[cation] at (0.75,0.25,0.75) {};
\node[anion] at (0,0,0) {};
\node[anion] at (1,0,0) {};
\node[anion] at (0.5,0,0.5) {};
\node[anion] at (0,0,1) {};
\node[anion] at (1,0,1) {};

\end{tikzpicture}
\caption{}
\end{subfigure}\hfill
\begin{subfigure}[b]{0.55\textwidth}\centering
\begin{tikzpicture}[scale=0.75]
  \clip (-3.4,-2.6) rectangle (3.4,2.6);
  \foreach \i in {-2,...,2}{\foreach \j in {-2,...,2}{
    \pgfmathsetmacro{\cx}{1.732*\i+0.866*\j}\pgfmathsetmacro{\cy}{1.5*\j}
    \foreach \k in {0,...,5}{
      \draw[fmBlue!40, line width=0.6pt, dashed] ($(\cx+0.866,\cy+0.5)+({cos(60*\k+30)},{sin(60*\k+30)})$) -- ($(\cx+0.866,\cy+0.5)+({cos(60*\k+90)},{sin(60*\k+90)})$);
      \draw[atomLine, line width=0.9pt] ($(\cx,\cy)+({cos(60*\k+30)},{sin(60*\k+30)})$) -- ($(\cx,\cy)+({cos(60*\k+90)},{sin(60*\k+90)})$);}}}
\end{tikzpicture}
\caption{}
\end{subfigure}
\caption{\textit{а}~— структура алмаза (атомы в узлах ГЦК-решётки и в
половине тетраэдрических пустот, линии~— ковалентные связи);
\textit{б}~— графит, вид вдоль оси $c$: сетки из шестиугольников
соседних слоёв A (сплошные) и B (штриховые) сдвинуты друг относительно
друга}
\label{fig:diamond-graphite}
\end{figure}

\paragraph{Тип $\beta$-вольфрама (A15)}
Группа $Pm\br3n$, кубическая примитивная решётка, 8 атомов в ячейке двух
сортов (рис.~\ref{fig:a15}): W$_\text{I}$ в $0,0,0$ и $\frac12,\frac12,\frac12$ (окружены
12 атомами W$_\text{II}$), W$_\text{II}$ парами на гранях ($0,\frac14,\frac12$; $0,\frac34,\frac12$
и~т.\,д.; к.\,ч.~=~14: $2+4+8$), образующими цепочки вдоль трёх осей куба.
Название историческое: за «$\beta$-вольфрам» приняли оксид W$_3$O. Сейчас
тип интересен тем, что его имеют многие низкотемпературные
сверхпроводники: Nb$_3$Sn ($T_\text{с}=18$~К), Nb$_3$Ge (23~К).

\begin{figure}[ht]
\centering
\begin{tikzpicture}[cubeproj, scale=2.4]
  \CubeEdges[atomLine!60]
  \foreach \p in {(0,0,0),(1,0,0),(0,1,0),(1,1,0),(0,0,1),(1,0,1),(0,1,1),(1,1,1),(0.5,0.5,0.5)}
    \node[atom=3.4mm] at \p {};
  \foreach \a/\b in {{(0,0.25,0.5)}/{(0,0.75,0.5)}, {(1,0.25,0.5)}/{(1,0.75,0.5)},
                     {(0.5,0,0.25)}/{(0.5,0,0.75)}, {(0.5,1,0.25)}/{(0.5,1,0.75)},
                     {(0.25,0.5,0)}/{(0.75,0.5,0)}, {(0.25,0.5,1)}/{(0.75,0.5,1)}}{
    \draw[fmAmber!80!black, line width=1.4pt] \a -- \b;
    \node[atomB=3.0mm] at \a {}; \node[atomB=3.0mm] at \b {};}
\end{tikzpicture}
\caption{Структура A15 ($\beta$-W, Nb$_3$Sn): атомы W$_\text{I}$ (Sn) в вершинах и
центре куба, атомы W$_\text{II}$ (Nb)~— парами на гранях; пары образуют три
семейства непересекающихся цепочек вдоль осей куба}
\label{fig:a15}
\end{figure}

\subsection{Структурные типы соединений AB}\label{ssec:ab-types}

Ионные структуры устойчивы, если каждый ион касается только ионов
противоположного знака; поэтому координация определяется отношением
радиусов катиона и аниона $r_\text{к}/r_\text{а}$ (подробно~— гл.~3, п.~3.1.5):
при $0{,}41$--$0{,}73$ выгоден октаэдр (к.\,ч.~=~6), при $0{,}73$--$1{,}0$~— куб
(к.\,ч.~=~8), при $0{,}225$--$0{,}41$~— тетраэдр (к.\,ч.~=~4).

\begin{figure}[ht]
\centering
\begin{subfigure}[b]{0.3\textwidth}\centering
\begin{tikzpicture}[cubeproj, scale=1.8, anion/.style={atom=4.4mm}, cation/.style={atomB=2.8mm}]
  \CubeEdges[draw opacity=0]
  % сгенерировано tools/gen_ionic_cells.py — не править вручную
\draw[bondline] (0,0,0) -- (0,0,0.5);
\draw[bondline] (0,0,0) -- (0,0.5,0);
\draw[bondline] (0,0,0) -- (0.5,0,0);
\draw[bondline] (0,0,0.5) -- (0,0,1);
\draw[bondline] (0,0,0.5) -- (0,0.5,0.5);
\draw[bondline] (0,0,0.5) -- (0.5,0,0.5);
\draw[bondline] (0,0,1) -- (0,0.5,1);
\draw[bondline] (0,0,1) -- (0.5,0,1);
\draw[bondline] (0,0.5,0) -- (0,0.5,0.5);
\draw[bondline] (0,0.5,0) -- (0,1,0);
\draw[bondline] (0,0.5,0) -- (0.5,0.5,0);
\draw[bondline] (0,0.5,0.5) -- (0,0.5,1);
\draw[bondline] (0,0.5,0.5) -- (0,1,0.5);
\draw[bondline] (0,0.5,0.5) -- (0.5,0.5,0.5);
\draw[bondline] (0,0.5,1) -- (0,1,1);
\draw[bondline] (0,0.5,1) -- (0.5,0.5,1);
\draw[bondline] (0,1,0) -- (0,1,0.5);
\draw[bondline] (0,1,0) -- (0.5,1,0);
\draw[bondline] (0,1,0.5) -- (0,1,1);
\draw[bondline] (0,1,0.5) -- (0.5,1,0.5);
\draw[bondline] (0,1,1) -- (0.5,1,1);
\draw[bondline] (0.5,0,0) -- (0.5,0,0.5);
\draw[bondline] (0.5,0,0) -- (0.5,0.5,0);
\draw[bondline] (0.5,0,0) -- (1,0,0);
\draw[bondline] (0.5,0,0.5) -- (0.5,0,1);
\draw[bondline] (0.5,0,0.5) -- (0.5,0.5,0.5);
\draw[bondline] (0.5,0,0.5) -- (1,0,0.5);
\draw[bondline] (0.5,0,1) -- (0.5,0.5,1);
\draw[bondline] (0.5,0,1) -- (1,0,1);
\draw[bondline] (0.5,0.5,0) -- (0.5,0.5,0.5);
\draw[bondline] (0.5,0.5,0) -- (0.5,1,0);
\draw[bondline] (0.5,0.5,0) -- (1,0.5,0);
\draw[bondline] (0.5,0.5,0.5) -- (0.5,0.5,1);
\draw[bondline] (0.5,0.5,0.5) -- (0.5,1,0.5);
\draw[bondline] (0.5,0.5,0.5) -- (1,0.5,0.5);
\draw[bondline] (0.5,0.5,1) -- (0.5,1,1);
\draw[bondline] (0.5,0.5,1) -- (1,0.5,1);
\draw[bondline] (0.5,1,0) -- (0.5,1,0.5);
\draw[bondline] (0.5,1,0) -- (1,1,0);
\draw[bondline] (0.5,1,0.5) -- (0.5,1,1);
\draw[bondline] (0.5,1,0.5) -- (1,1,0.5);
\draw[bondline] (0.5,1,1) -- (1,1,1);
\draw[bondline] (1,0,0) -- (1,0,0.5);
\draw[bondline] (1,0,0) -- (1,0.5,0);
\draw[bondline] (1,0,0.5) -- (1,0,1);
\draw[bondline] (1,0,0.5) -- (1,0.5,0.5);
\draw[bondline] (1,0,1) -- (1,0.5,1);
\draw[bondline] (1,0.5,0) -- (1,0.5,0.5);
\draw[bondline] (1,0.5,0) -- (1,1,0);
\draw[bondline] (1,0.5,0.5) -- (1,0.5,1);
\draw[bondline] (1,0.5,0.5) -- (1,1,0.5);
\draw[bondline] (1,0.5,1) -- (1,1,1);
\draw[bondline] (1,1,0) -- (1,1,0.5);
\draw[bondline] (1,1,0.5) -- (1,1,1);
\node[anion] at (0,1,0) {};
\node[cation] at (0.5,1,0) {};
\node[anion] at (1,1,0) {};
\node[cation] at (0,1,0.5) {};
\node[anion] at (0.5,1,0.5) {};
\node[cation] at (1,1,0.5) {};
\node[anion] at (0,1,1) {};
\node[cation] at (0.5,1,1) {};
\node[anion] at (1,1,1) {};
\node[cation] at (0,0.5,0) {};
\node[anion] at (0.5,0.5,0) {};
\node[cation] at (1,0.5,0) {};
\node[anion] at (0,0.5,0.5) {};
\node[cation] at (0.5,0.5,0.5) {};
\node[anion] at (1,0.5,0.5) {};
\node[cation] at (0,0.5,1) {};
\node[anion] at (0.5,0.5,1) {};
\node[cation] at (1,0.5,1) {};
\node[anion] at (0,0,0) {};
\node[cation] at (0.5,0,0) {};
\node[anion] at (1,0,0) {};
\node[cation] at (0,0,0.5) {};
\node[anion] at (0.5,0,0.5) {};
\node[cation] at (1,0,0.5) {};
\node[anion] at (0,0,1) {};
\node[cation] at (0.5,0,1) {};
\node[anion] at (1,0,1) {};

\end{tikzpicture}
\caption{NaCl (B1)}
\end{subfigure}\hfill
\begin{subfigure}[b]{0.3\textwidth}\centering
\begin{tikzpicture}[cubeproj, scale=1.8, anion/.style={atom=4.4mm}, cation/.style={atomB=4.0mm}]
  \CubeEdges
  % сгенерировано tools/gen_ionic_cells.py — не править вручную
\draw[bondline] (0.5,0.5,0.5) -- (0,0,0);
\draw[bondline] (0.5,0.5,0.5) -- (0,0,1);
\draw[bondline] (0.5,0.5,0.5) -- (0,1,0);
\draw[bondline] (0.5,0.5,0.5) -- (0,1,1);
\draw[bondline] (0.5,0.5,0.5) -- (1,0,0);
\draw[bondline] (0.5,0.5,0.5) -- (1,0,1);
\draw[bondline] (0.5,0.5,0.5) -- (1,1,0);
\draw[bondline] (0.5,0.5,0.5) -- (1,1,1);
\node[anion] at (0,1,0) {};
\node[anion] at (1,1,0) {};
\node[anion] at (0,1,1) {};
\node[anion] at (1,1,1) {};
\node[cation] at (0.5,0.5,0.5) {};
\node[anion] at (0,0,0) {};
\node[anion] at (1,0,0) {};
\node[anion] at (0,0,1) {};
\node[anion] at (1,0,1) {};

\end{tikzpicture}
\caption{CsCl (B2)}
\end{subfigure}\hfill
\begin{subfigure}[b]{0.3\textwidth}\centering
\begin{tikzpicture}[cubeproj, scale=1.8, anion/.style={atom=4.2mm}, cation/.style={atomB=2.8mm}]
  \CubeEdges
  % сгенерировано tools/gen_ionic_cells.py — не править вручную
\draw[bondthick] (0.25,0.25,0.25) -- (0,0,0);
\draw[bondthick] (0.25,0.25,0.25) -- (0,0.5,0.5);
\draw[bondthick] (0.25,0.25,0.25) -- (0.5,0,0.5);
\draw[bondthick] (0.25,0.25,0.25) -- (0.5,0.5,0);
\draw[bondthick] (0.75,0.75,0.25) -- (0.5,0.5,0);
\draw[bondthick] (0.75,0.75,0.25) -- (0.5,1,0.5);
\draw[bondthick] (0.75,0.75,0.25) -- (1,0.5,0.5);
\draw[bondthick] (0.75,0.75,0.25) -- (1,1,0);
\draw[bondthick] (0.75,0.25,0.75) -- (0.5,0,0.5);
\draw[bondthick] (0.75,0.25,0.75) -- (0.5,0.5,1);
\draw[bondthick] (0.75,0.25,0.75) -- (1,0,1);
\draw[bondthick] (0.75,0.25,0.75) -- (1,0.5,0.5);
\draw[bondthick] (0.25,0.75,0.75) -- (0,0.5,0.5);
\draw[bondthick] (0.25,0.75,0.75) -- (0,1,1);
\draw[bondthick] (0.25,0.75,0.75) -- (0.5,0.5,1);
\draw[bondthick] (0.25,0.75,0.75) -- (0.5,1,0.5);
\node[anion] at (0,1,0) {};
\node[anion] at (1,1,0) {};
\node[anion] at (0.5,1,0.5) {};
\node[anion] at (0,1,1) {};
\node[anion] at (1,1,1) {};
\node[cation] at (0.75,0.75,0.25) {};
\node[cation] at (0.25,0.75,0.75) {};
\node[anion] at (0.5,0.5,0) {};
\node[anion] at (0,0.5,0.5) {};
\node[anion] at (1,0.5,0.5) {};
\node[anion] at (0.5,0.5,1) {};
\node[cation] at (0.25,0.25,0.25) {};
\node[cation] at (0.75,0.25,0.75) {};
\node[anion] at (0,0,0) {};
\node[anion] at (1,0,0) {};
\node[anion] at (0.5,0,0.5) {};
\node[anion] at (0,0,1) {};
\node[anion] at (1,0,1) {};

\end{tikzpicture}
\caption{ZnS, сфалерит (B3)}
\end{subfigure}
\caption{Структурные типы соединений AB (голубые~— анионы, оранжевые~—
катионы)}
\label{fig:ab-types}
\end{figure}

\begin{itemize}
\item \term{NaCl (B1)}, $Fm\br3m$: анионы Cl$^-$ в узлах ГЦК-решётки ($0,0,0$;
$\frac12,\frac12,0\,[3]$), катионы Na$^+$ во \emph{всех октаэдрических} пустотах
($\frac12,\frac12,\frac12$; $\frac12,0,0\,[3]$); тетрапустоты свободны. К.\,ч.~=~6 для
обоих ионов (октаэдр), во второй сфере~— 12 одноимённых. Так же построены
оксиды TiO, MnO, FeO, NiO, карбиды и нитриды переходных металлов IV--V
групп (TiC, VN\ldots). Для NaCl $r_\text{Na}/r_\text{Cl}=0{,}98/1{,}81=0{,}54$~—
октаэдрический интервал.
\item \term{CsCl (B2)}, $Pm\br3m$: Cl$^-$ в вершинах куба, Cs$^+$ в центре;
к.\,ч.~=~8 (куб). Решётка Браве~— \emph{простая} кубическая (центр и вершины
заняты разными ионами). $r_\text{Cs}/r_\text{Cl}=1{,}65/1{,}81=0{,}91$~— кубический
интервал, к.\,ч.~=~8\ispr{В оригинале для CsCl сказано, что отношение
радиусов 0,91 «соответствует ОЦК структуре с к.\,ч.~=~6»; правильно~—
к.\,ч.~=~8, а решётка Браве простая кубическая.}. Этот тип имеют многие
интерметаллиды с близкими радиусами атомов (упорядоченная $\beta'$-латунь
CuZn, NiAl, FeAl).
\item \term{Сфалерит ZnS (B3)}, $F\br43m$: S в узлах ГЦК (кубическая плотнейшая
упаковка), Zn в \emph{половине тетраэдрических} пустот ($\frac14,\frac14,\frac14$;
$\frac34,\frac34,\frac14\,[3]$); к.\,ч.~=~4 у обоих (тетраэдры). Центра симметрии нет~—
структура полярна (вдоль $\nfam{111}$), что важно для пьезоэффекта. Тот же
тип~— у GaAs, InSb, CdTe, кубического BN.
\item \term{Вюрцит ZnS (B4)}, $P6_3mc$: то же, но на основе
\emph{гексагональной} плотнейшей упаковки атомов S; Zn в половине
тетрапустот, к.\,ч.~=~4. Так построены ZnO, GaN, AlN, BeO.
\item \term{NiAs (B8)}, $P6_3/mmc$: As~— гексагональная плотнейшая упаковка,
Ni во \emph{всех октаэдрических} пустотах ($0,0,0$; $0,0,\frac12$), к.\,ч.~=~6, $Z=2$.
Характерен для соединений переходных металлов с S, Se, Te, As, Sb (FeS, CoS,
NiSb).
\end{itemize}

\subsection{Структурные типы соединений AB$_2$}\label{ssec:ab2-types}

\term{Флюорит CaF$_2$ (C1)}, $Fm\br3m$ (рис.~\ref{fig:fluorite-perov}\,\textit{а}):
ионы Ca$^{2+}$ в узлах ГЦК-решётки, ионы F$^-$ во \emph{всех тетраэдрических}
пустотах ($\frac14,\frac14,\frac14$ и т.\,д.). Каждый F$^-$ окружён тетраэдром из четырёх
Ca$^{2+}$, каждый Ca$^{2+}$~— кубом из восьми F$^-$: к.\,ч.(F)~=~4, к.\,ч.(Ca)~=~8.
Необычно, что плотнейшую упаковку образуют здесь \emph{катионы}. Так
построены UO$_2$, ThO$_2$, ZrO$_2$ (кубический), PuO$_2$~— основа ядерного
топлива. В \emph{антифлюорите} (Li$_2$O, Na$_2$O, Mg$_2$Si) катионы и анионы
меняются местами: плотную упаковку образуют анионы, катионы заполняют все
тетрапустоты.

\subsection{Структурные типы сложных соединений}\label{ssec:complex-types}

\term{Перовскит CaTiO$_3$ (E2$_1$)}, идеализированно $Pm\br3m$
(рис.~\ref{fig:fluorite-perov}\,\textit{б}): Ca в вершинах куба, Ti в центре,
O в центрах граней. К.\,ч.: Ti по O~— 6 (октаэдр TiO$_6$), Ca по O~— 12, O по
Ti~— 2, O по Ca~— 4. Реальный CaTiO$_3$ слегка искажён (октаэдры повёрнуты)
и имеет ромбическую симметрию. Перовскит~— одна из важнейших структур
материаловедения: сегнетоэлектрики (BaTiO$_3$), а также слоистые структуры,
производные от перовскита,~— высокотемпературные сверхпроводники
YBa$_2$Cu$_3$O$_{7-\delta}$ ($T_\text{с}=93$~К), Bi$_2$Sr$_2$CaCu$_2$O$_{8+x}$ (до 94~К),
Tl$_2$Ba$_2$Ca$_2$Cu$_3$O$_{10+x}$ (125~К), HgBa$_2$Ca$_2$Cu$_3$O$_{8+\delta}$ (133~К).

\begin{figure}[ht]
\centering
\begin{subfigure}[b]{0.4\textwidth}\centering
\begin{tikzpicture}[cubeproj, scale=1.9, anion/.style={atomB=3.4mm}, cation/.style={atom=4.4mm}]
  \CubeEdges
  % сгенерировано tools/gen_ionic_cells.py — не править вручную
\draw[bondline] (0.25,0.25,0.25) -- (0,0,0);
\draw[bondline] (0.25,0.25,0.25) -- (0,0.5,0.5);
\draw[bondline] (0.25,0.25,0.25) -- (0.5,0,0.5);
\draw[bondline] (0.25,0.25,0.25) -- (0.5,0.5,0);
\draw[bondline] (0.25,0.25,0.75) -- (0,0,1);
\draw[bondline] (0.25,0.25,0.75) -- (0,0.5,0.5);
\draw[bondline] (0.25,0.25,0.75) -- (0.5,0,0.5);
\draw[bondline] (0.25,0.25,0.75) -- (0.5,0.5,1);
\draw[bondline] (0.25,0.75,0.25) -- (0,0.5,0.5);
\draw[bondline] (0.25,0.75,0.25) -- (0,1,0);
\draw[bondline] (0.25,0.75,0.25) -- (0.5,0.5,0);
\draw[bondline] (0.25,0.75,0.25) -- (0.5,1,0.5);
\draw[bondline] (0.25,0.75,0.75) -- (0,0.5,0.5);
\draw[bondline] (0.25,0.75,0.75) -- (0,1,1);
\draw[bondline] (0.25,0.75,0.75) -- (0.5,0.5,1);
\draw[bondline] (0.25,0.75,0.75) -- (0.5,1,0.5);
\draw[bondline] (0.75,0.25,0.25) -- (0.5,0,0.5);
\draw[bondline] (0.75,0.25,0.25) -- (0.5,0.5,0);
\draw[bondline] (0.75,0.25,0.25) -- (1,0,0);
\draw[bondline] (0.75,0.25,0.25) -- (1,0.5,0.5);
\draw[bondline] (0.75,0.25,0.75) -- (0.5,0,0.5);
\draw[bondline] (0.75,0.25,0.75) -- (0.5,0.5,1);
\draw[bondline] (0.75,0.25,0.75) -- (1,0,1);
\draw[bondline] (0.75,0.25,0.75) -- (1,0.5,0.5);
\draw[bondline] (0.75,0.75,0.25) -- (0.5,0.5,0);
\draw[bondline] (0.75,0.75,0.25) -- (0.5,1,0.5);
\draw[bondline] (0.75,0.75,0.25) -- (1,0.5,0.5);
\draw[bondline] (0.75,0.75,0.25) -- (1,1,0);
\draw[bondline] (0.75,0.75,0.75) -- (0.5,0.5,1);
\draw[bondline] (0.75,0.75,0.75) -- (0.5,1,0.5);
\draw[bondline] (0.75,0.75,0.75) -- (1,0.5,0.5);
\draw[bondline] (0.75,0.75,0.75) -- (1,1,1);
\node[cation] at (0,1,0) {};
\node[cation] at (1,1,0) {};
\node[cation] at (0.5,1,0.5) {};
\node[cation] at (0,1,1) {};
\node[cation] at (1,1,1) {};
\node[anion] at (0.25,0.75,0.25) {};
\node[anion] at (0.75,0.75,0.25) {};
\node[anion] at (0.25,0.75,0.75) {};
\node[anion] at (0.75,0.75,0.75) {};
\node[cation] at (0.5,0.5,0) {};
\node[cation] at (0,0.5,0.5) {};
\node[cation] at (1,0.5,0.5) {};
\node[cation] at (0.5,0.5,1) {};
\node[anion] at (0.25,0.25,0.25) {};
\node[anion] at (0.75,0.25,0.25) {};
\node[anion] at (0.25,0.25,0.75) {};
\node[anion] at (0.75,0.25,0.75) {};
\node[cation] at (0,0,0) {};
\node[cation] at (1,0,0) {};
\node[cation] at (0.5,0,0.5) {};
\node[cation] at (0,0,1) {};
\node[cation] at (1,0,1) {};

\end{tikzpicture}
\caption{CaF$_2$}
\end{subfigure}\hspace{6mm}
\begin{subfigure}[b]{0.4\textwidth}\centering
\begin{tikzpicture}[cubeproj, scale=2.1]
  \CubeEdges[atomLine!60]
  \draw[fmBlue!50, line width=0.5pt] (0.5,0.5,0) -- (0.5,0.5,1) (0,0.5,0.5) -- (1,0.5,0.5) (0.5,0,0.5) -- (0.5,1,0.5);
  \foreach \p in {(0.5,1,0.5),(0.5,0.5,0),(0,0.5,0.5)} \node[atomI=2.8mm] at \p {};
  \node[atomB=3.2mm] at (0.5,0.5,0.5) {};
  \foreach \p in {(1,0.5,0.5),(0.5,0.5,1),(0.5,0,0.5)} \node[atomI=2.8mm] at \p {};
  \foreach \p in {(0,0,0),(1,0,0),(0,1,0),(1,1,0),(0,0,1),(1,0,1),(0,1,1),(1,1,1)} \node[atom=4.2mm] at \p {};
\end{tikzpicture}
\caption{CaTiO$_3$}
\end{subfigure}
\caption{\textit{а}~— флюорит CaF$_2$ (голубые~— Ca$^{2+}$, оранжевые~— F$^-$);
\textit{б}~— идеальный перовскит CaTiO$_3$: Ca (голубые) в вершинах, Ti
(оранжевый) в центре октаэдра из атомов O (красные) в центрах граней}
\label{fig:fluorite-perov}
\end{figure}

\subsection{Сверхструктуры в твёрдых растворах замещения}\label{ssec:superstructures}

Напомним (гл.~3, п.~3.3.2): твёрдые растворы бывают замещения,
внедрения (атомы H, C, N, O, B в пустотах металлической решётки) и
вычитания (на основе соединений, часть узлов пуста). Растворение атомов
меньшего размера уменьшает период решётки, большего~— увеличивает; в
непрерывных растворах замещения период часто меняется линейно с
составом~— \term{закон Вегарда}:
\begin{equation}\label{eq:vegard}
a=a_1c_1+a_2c_2=a_1+(a_2-a_1)c_2 ,
\end{equation}
где $a_i$~— периоды компонентов, $c_i$~— атомные доли.

При упорядочении на рентгенограммах появляются дополнительные,
\emph{сверхструктурные}, линии~— признак сверхструктуры (сверхрешётки).
Важнейшие сверхструктуры (рис.~\ref{fig:superstructures}):
\begin{itemize}
\item \term{Cu$_3$Au (L1$_2$)}: ниже $\approx390\,^\circ$C из неупорядоченной ГЦК
($Fm\br3m$) образуется структура с Au в вершинах куба и Cu в центрах граней;
группа $Pm\br3m$ (решётка стала примитивной).
\item \term{CuAu\,I (L1$_0$)}: ниже $\approx385\,^\circ$C слои $(001)$ заняты
поочерёдно атомами Cu и Au; решётка тетрагональна ($c/a\approx0{,}93$),
$P4/mmm$.
\item \term{CuPt (L1$_1$)}: упорядочение по плоскостям $(111)$ с
ромбоэдрическим искажением.
\item \term{CuAu\,II}~— сверхструктура с \emph{длинным периодом}: десять
ячеек CuAu\,I вдоль оси $b$, причём посередине периода слои меди и золота
меняются местами~— периодически расположенные антифазные границы
(одномерная длиннопериодическая сверхструктура).
\item На основе ОЦК: \term{$\beta'$-латунь CuZn (B2)} ниже $\approx465\,^\circ$C
(переход $Im\br3m\to Pm\br3m$); на ячейке с удвоенным периодом (8 ОЦК-ячеек,
центры которых делятся на позиции $X$ и $Y$, чередующиеся в шахматном
порядке): \term{FeAl} (Fe в вершинах, Al в $X$ и $Y$~— та же B2),
\term{Fe$_3$Al (D0$_3$)} (Fe в вершинах и $X$, Al в $Y$), \term{Cu$_2$MnAl (L2$_1$,
сплав Гейслера)} (Cu в вершинах, Mn в $X$, Al в $Y$; ферромагнетик из
неферромагнитных компонентов, см. гл.~3, п.~3.6.8).
\end{itemize}

\begin{figure}[ht]
\centering
\begin{subfigure}[b]{0.24\textwidth}\centering
\begin{tikzpicture}[cubeproj, scale=1.55]
  \CubeEdges
  \foreach \p in {(0,1,0),(1,1,0),(0,1,1),(1,1,1)} \node[atomB=3.8mm] at \p {};
  \node[atom=3.2mm] at (0.5,1,0.5) {};
  \foreach \p in {(0,0.5,0.5),(1,0.5,0.5),(0.5,0.5,0),(0.5,0.5,1)} \node[atom=3.2mm] at \p {};
  \foreach \p in {(0,0,0),(1,0,0),(0,0,1),(1,0,1)} \node[atomB=3.8mm] at \p {};
  \node[atom=3.2mm] at (0.5,0,0.5) {};
\end{tikzpicture}
\caption{Cu$_3$Au}
\end{subfigure}\hfill
\begin{subfigure}[b]{0.24\textwidth}\centering
\begin{tikzpicture}[cubeproj, x={(1cm,0cm)}, y={(0.42cm,0.30cm)}, z={(0cm,0.93cm)}, scale=1.55]
  \CubeEdges
  \foreach \p in {(0,1,0),(1,1,0),(0,1,1),(1,1,1)} \node[atomB=3.8mm] at \p {};
  \node[atom=3.2mm] at (0.5,1,0.5) {};
  \foreach \p in {(0,0.5,0.5),(1,0.5,0.5)} \node[atom=3.2mm] at \p {};
  \foreach \p in {(0.5,0.5,0),(0.5,0.5,1)} \node[atomB=3.8mm] at \p {};
  \foreach \p in {(0,0,0),(1,0,0),(0,0,1),(1,0,1)} \node[atomB=3.8mm] at \p {};
  \node[atom=3.2mm] at (0.5,0,0.5) {};
\end{tikzpicture}
\caption{CuAu\,I}
\end{subfigure}\hfill
\begin{subfigure}[b]{0.24\textwidth}\centering
\begin{tikzpicture}[cubeproj, scale=1.55]
  \CubeEdges
  \foreach \p in {(0,1,0),(1,1,0),(0,1,1),(1,1,1)} \node[atom=3.6mm] at \p {};
  \node[atomI=3.6mm] at (0.5,0.5,0.5) {};
  \foreach \p in {(0,0,0),(1,0,0),(0,0,1),(1,0,1)} \node[atom=3.6mm] at \p {};
\end{tikzpicture}
\caption{CuZn (B2)}
\end{subfigure}\hfill
\begin{subfigure}[b]{0.24\textwidth}\centering
\begin{tikzpicture}[cubeproj, scale=0.78]
  \draw[atomLine!50, dashed] (0,0,0)--(0,2,0)--(2,2,0) (0,2,0)--(0,2,2);
  \draw[atomLine] (0,0,0)--(2,0,0)--(2,0,2)--(0,0,2)--cycle (2,0,0)--(2,2,0)--(2,2,2)--(2,0,2) (0,0,2)--(0,2,2)--(2,2,2);
  \foreach \x in {0,1,2}{\foreach \y in {2,1,0}{\foreach \z in {0,1,2}{\node[atom=2.2mm] at (\x,\y,\z) {};}}}
  \foreach \x in {0.5,1.5}{\foreach \y in {1.5,0.5}{\foreach \z in {0.5,1.5}{
    \pgfmathtruncatemacro{\pp}{mod(round(\x+\y+\z-0.5),2)}
    \ifnum\pp=1 \node[atom=2.2mm] at (\x,\y,\z) {};\else \node[atomI=2.6mm] at (\x,\y,\z) {};\fi}}}
\end{tikzpicture}
\caption{Fe$_3$Al (D0$_3$)}
\end{subfigure}
\caption{Сверхструктуры: \textit{а}~— Cu$_3$Au (Au оранжевые~— в вершинах,
Cu~— в центрах граней); \textit{б}~— CuAu\,I (слои Au и Cu вдоль оси $c$);
\textit{в}~— CuZn (атомы разного сорта в вершинах и в центре куба);
\textit{г}~— Fe$_3$Al в удвоенной ОЦК-ячейке (Fe~— голубые, Al~— красные в
позициях $Y$)}
\label{fig:superstructures}
\end{figure}

\subsection{Фуллерены, фуллериты, квазикристаллы}\label{ssec:fullerenes}

\paragraph{Фуллерены}
При лазерном испарении графита в струе гелия образуются устойчивые
молекулы из большого чётного числа атомов углерода, чаще всего C$_{60}$ и
C$_{70}$. Форма молекулы C$_{60}$ напоминает купола архитектора Р.~Фуллера~—
отсюда название \term{фуллерены}. Атомы C$_{60}$ сидят в вершинах
\emph{усечённого икосаэдра} из 20 шестиугольников и 12 правильных
пятиугольников (рис.~\ref{fig:c60}); симметрия молекулы икосаэдрическая,
оси 5-го порядка проходят через центры пятиугольников. Все атомы лежат на
сфере радиусом $R=0{,}357$~нм; валентные электроны распределены почти
равномерно по сферическому слою толщиной около $8a_0$ ($a_0=0{,}0529$~нм~—
боровский радиус), а внутри есть полость радиусом около $2a_0$, почти
свободная от электронов. В неё можно поместить атом~— фуллерены
рассматривают как «контейнеры» для отдельных атомов и молекул.

\begin{figure}[ht]
\centering
\begin{subfigure}[b]{0.42\textwidth}\centering
\tikzset{c60front/.style={atomLine, line width=0.7pt},
  c60back/.style={atomLine!30, line width=0.4pt},
  c60pent/.style={fmAmber!45}}
\begin{tikzpicture}[scale=1.25]
  \fill[fmBlue!8] (0,0) circle (1.25);
  % создано tools/gen_cryst.py: молекула C60 (передние рёбра — сплошные)
\fill[c60pent] (0.009,-1.110) -- (0.451,-1.110) -- (0.757,-0.803) -- (0.506,-0.613) -- (0.043,-0.803) -- cycle;
\fill[c60pent] (-0.793,0.847) -- (-0.570,0.692) -- (-0.107,0.882) -- (-0.045,1.154) -- (-0.468,1.132) -- cycle;
\fill[c60pent] (0.188,0.148) -- (0.578,-0.137) -- (0.903,0.148) -- (0.713,0.610) -- (0.272,0.610) -- cycle;
\fill[c60pent] (-0.770,-0.539) -- (-0.347,-0.517) -- (-0.274,-0.042) -- (-0.653,0.230) -- (-0.960,-0.077) -- cycle;
\draw[c60back] (-0.188,-0.148) -- (0.274,0.042);
\draw[c60back] (-0.578,0.137) -- (-0.188,-0.148);
\draw[c60back] (0.274,0.042) -- (0.347,0.517);
\draw[c60back] (-0.272,-0.610) -- (-0.188,-0.148);
\draw[c60back] (0.274,0.042) -- (0.653,-0.230);
\draw[c60back] (-0.578,0.137) -- (-0.506,0.613);
\draw[c60back] (-0.043,0.803) -- (0.347,0.517);
\draw[c60back] (-0.903,-0.148) -- (-0.578,0.137);
\draw[c60back] (-0.272,-0.610) -- (0.107,-0.882);
\draw[c60back] (-0.506,0.613) -- (-0.043,0.803);
\draw[c60back] (0.570,-0.692) -- (0.653,-0.230);
\draw[c60back] (0.347,0.517) -- (0.770,0.539);
\draw[c60back] (-0.713,-0.610) -- (-0.272,-0.610);
\draw[c60back] (0.653,-0.230) -- (0.960,0.077);
\draw[c60back] (0.107,-0.882) -- (0.570,-0.692);
\draw[c60back] (-0.903,-0.148) -- (-0.713,-0.610);
\draw[c60back] (0.770,0.539) -- (0.960,0.077);
\draw[c60back] (-0.757,0.803) -- (-0.506,0.613);
\draw[c60back] (-0.009,1.110) -- (-0.043,0.803);
\draw[c60back] (0.045,-1.154) -- (0.107,-0.882);
\draw[c60back] (0.570,-0.692) -- (0.793,-0.847);
\draw[c60back] (-1.154,0.042) -- (-0.903,-0.148);
\draw[c60back] (0.804,0.847) -- (0.770,0.539);
\draw[c60back] (-0.776,-0.882) -- (-0.713,-0.610);
\draw[c60back] (0.960,0.077) -- (1.183,-0.077);
\draw[c60back] (-0.757,0.803) -- (-0.451,1.110);
\draw[c60back] (-1.082,0.517) -- (-0.757,0.803);
\draw[c60back] (-0.451,1.110) -- (-0.009,1.110);
\draw[c60back] (-0.009,1.110) -- (0.414,1.132);
\draw[c60back] (-1.154,0.042) -- (-1.082,0.517);
\draw[c60back] (0.414,1.132) -- (0.804,0.847);
\draw[c60back] (-0.397,-1.154) -- (0.045,-1.154);
\draw[c60back] (0.045,-1.154) -- (0.468,-1.132);
\draw[c60back] (0.793,-0.847) -- (1.100,-0.539);
\draw[c60back] (0.468,-1.132) -- (0.793,-0.847);
\draw[c60back] (-0.776,-0.882) -- (-0.397,-1.154);
\draw[c60back] (1.100,-0.539) -- (1.183,-0.077);
\draw[c60back] (-1.217,-0.230) -- (-1.154,0.042);
\draw[c60back] (0.804,0.847) -- (1.028,0.692);
\draw[c60back] (-1.028,-0.692) -- (-0.776,-0.882);
\draw[c60back] (1.217,0.230) -- (1.183,-0.077);
\draw[c60back] (-0.468,1.132) -- (-0.451,1.110);
\draw[c60back] (-1.100,0.539) -- (-1.082,0.517);
\draw[c60back] (0.397,1.154) -- (0.414,1.132);
\draw[c60back] (-1.217,-0.230) -- (-1.028,-0.692);
\draw[c60front] (1.028,0.692) -- (1.217,0.230);
\draw[c60front] (-0.414,-1.132) -- (-0.397,-1.154);
\draw[c60front] (1.082,-0.517) -- (1.100,-0.539);
\draw[c60front] (0.451,-1.110) -- (0.468,-1.132);
\draw[c60front] (-1.183,0.077) -- (-1.217,-0.230);
\draw[c60front] (0.776,0.882) -- (1.028,0.692);
\draw[c60front] (-1.028,-0.692) -- (-0.804,-0.847);
\draw[c60front] (1.154,-0.042) -- (1.217,0.230);
\draw[c60front] (-1.183,0.077) -- (-1.100,0.539);
\draw[c60front] (0.397,1.154) -- (0.776,0.882);
\draw[c60front] (-0.793,0.847) -- (-0.468,1.132);
\draw[c60front] (-1.100,0.539) -- (-0.793,0.847);
\draw[c60front] (-0.468,1.132) -- (-0.045,1.154);
\draw[c60front] (-0.045,1.154) -- (0.397,1.154);
\draw[c60front] (-0.804,-0.847) -- (-0.414,-1.132);
\draw[c60front] (1.082,-0.517) -- (1.154,-0.042);
\draw[c60front] (-0.414,-1.132) -- (0.009,-1.110);
\draw[c60front] (0.009,-1.110) -- (0.451,-1.110);
\draw[c60front] (0.757,-0.803) -- (1.082,-0.517);
\draw[c60front] (0.451,-1.110) -- (0.757,-0.803);
\draw[c60front] (-1.183,0.077) -- (-0.960,-0.077);
\draw[c60front] (0.713,0.610) -- (0.776,0.882);
\draw[c60front] (-0.770,-0.539) -- (-0.804,-0.847);
\draw[c60front] (0.903,0.148) -- (1.154,-0.042);
\draw[c60front] (-0.793,0.847) -- (-0.570,0.692);
\draw[c60front] (-0.107,0.882) -- (-0.045,1.154);
\draw[c60front] (0.043,-0.803) -- (0.009,-1.110);
\draw[c60front] (0.506,-0.613) -- (0.757,-0.803);
\draw[c60front] (-0.960,-0.077) -- (-0.770,-0.539);
\draw[c60front] (0.713,0.610) -- (0.903,0.148);
\draw[c60front] (-0.570,0.692) -- (-0.107,0.882);
\draw[c60front] (-0.960,-0.077) -- (-0.653,0.230);
\draw[c60front] (0.272,0.610) -- (0.713,0.610);
\draw[c60front] (-0.770,-0.539) -- (-0.347,-0.517);
\draw[c60front] (-0.653,0.230) -- (-0.570,0.692);
\draw[c60front] (0.043,-0.803) -- (0.506,-0.613);
\draw[c60front] (-0.107,0.882) -- (0.272,0.610);
\draw[c60front] (0.578,-0.137) -- (0.903,0.148);
\draw[c60front] (-0.347,-0.517) -- (0.043,-0.803);
\draw[c60front] (0.506,-0.613) -- (0.578,-0.137);
\draw[c60front] (-0.653,0.230) -- (-0.274,-0.042);
\draw[c60front] (0.188,0.148) -- (0.272,0.610);
\draw[c60front] (-0.347,-0.517) -- (-0.274,-0.042);
\draw[c60front] (0.188,0.148) -- (0.578,-0.137);
\draw[c60front] (-0.274,-0.042) -- (0.188,0.148);
\fill[atomLine] (1.028,0.692) circle (0.03);
\fill[atomLine] (1.217,0.230) circle (0.03);
\fill[atomLine] (-0.468,1.132) circle (0.03);
\fill[atomLine] (-1.100,0.539) circle (0.03);
\fill[atomLine] (0.397,1.154) circle (0.03);
\fill[atomLine] (-0.414,-1.132) circle (0.03);
\fill[atomLine] (1.082,-0.517) circle (0.03);
\fill[atomLine] (0.451,-1.110) circle (0.03);
\fill[atomLine] (-1.183,0.077) circle (0.03);
\fill[atomLine] (0.776,0.882) circle (0.03);
\fill[atomLine] (-0.804,-0.847) circle (0.03);
\fill[atomLine] (-0.793,0.847) circle (0.03);
\fill[atomLine] (1.154,-0.042) circle (0.03);
\fill[atomLine] (-0.045,1.154) circle (0.03);
\fill[atomLine] (0.009,-1.110) circle (0.03);
\fill[atomLine] (0.757,-0.803) circle (0.03);
\fill[atomLine] (-0.960,-0.077) circle (0.03);
\fill[atomLine] (-0.770,-0.539) circle (0.03);
\fill[atomLine] (0.713,0.610) circle (0.03);
\fill[atomLine] (0.903,0.148) circle (0.03);
\fill[atomLine] (-0.570,0.692) circle (0.03);
\fill[atomLine] (-0.107,0.882) circle (0.03);
\fill[atomLine] (0.043,-0.803) circle (0.03);
\fill[atomLine] (0.506,-0.613) circle (0.03);
\fill[atomLine] (-0.653,0.230) circle (0.03);
\fill[atomLine] (0.272,0.610) circle (0.03);
\fill[atomLine] (-0.347,-0.517) circle (0.03);
\fill[atomLine] (0.578,-0.137) circle (0.03);
\fill[atomLine] (-0.274,-0.042) circle (0.03);
\fill[atomLine] (0.188,0.148) circle (0.03);

\end{tikzpicture}
\caption{}
\end{subfigure}\hfill
\begin{subfigure}[b]{0.52\textwidth}\centering
\tikzset{pthin/.style={fill=fmBlue!45}, pthick/.style={fill=fmAmber!35},
  pedge/.style={atomLine, line width=0.25pt}}
\begin{tikzpicture}[scale=0.92]
  \clip (0,0) circle (3.3);
  % создано tools/gen_cryst.py: мозаика Пенроуза
\fill[pthin] (2.762,0.401) -- (2.470,0.307) -- (2.470,0.496) -- cycle;
\fill[pthick] (2.470,0.307) -- (2.762,0.401) -- (2.470,0.000) -- cycle;
\fill[pthick] (2.179,0.401) -- (1.998,0.649) -- (2.470,0.496) -- cycle;
\fill[pthick] (2.470,0.307) -- (2.179,0.401) -- (2.470,0.000) -- cycle;
\fill[pthin] (2.179,0.401) -- (2.470,0.307) -- (2.470,0.496) -- cycle;
\fill[pthick] (2.942,0.649) -- (2.762,0.401) -- (2.762,0.897) -- cycle;
\fill[pthick] (2.942,0.956) -- (2.942,0.649) -- (3.234,1.051) -- cycle;
\fill[pthin] (2.942,0.649) -- (2.942,0.956) -- (2.762,0.897) -- cycle;
\fill[pthick] (2.290,0.744) -- (1.998,0.649) -- (2.470,0.496) -- cycle;
\fill[pthick] (2.582,0.649) -- (2.290,0.744) -- (2.762,0.897) -- cycle;
\fill[pthin] (2.290,0.744) -- (2.582,0.649) -- (2.470,0.496) -- cycle;
\fill[pthin] (2.762,0.401) -- (2.582,0.649) -- (2.470,0.496) -- cycle;
\fill[pthick] (2.582,0.649) -- (2.762,0.401) -- (2.762,0.897) -- cycle;
\fill[pthick] (2.942,-0.153) -- (3.234,-0.248) -- (2.762,-0.401) -- cycle;
\fill[pthick] (2.650,-0.248) -- (2.942,-0.153) -- (2.470,0.000) -- cycle;
\fill[pthin] (2.942,-0.153) -- (2.650,-0.248) -- (2.762,-0.401) -- cycle;
\fill[pthick] (3.053,-0.803) -- (3.234,-1.051) -- (3.234,-0.555) -- cycle;
\fill[pthick] (3.053,-0.496) -- (3.053,-0.803) -- (2.762,-0.401) -- cycle;
\fill[pthin] (3.053,-0.803) -- (3.053,-0.496) -- (3.234,-0.555) -- cycle;
\fill[pthin] (3.234,-0.248) -- (3.053,-0.496) -- (3.234,-0.555) -- cycle;
\fill[pthick] (3.053,-0.496) -- (3.234,-0.248) -- (2.762,-0.401) -- cycle;
\fill[pthick] (2.942,0.649) -- (2.762,0.401) -- (3.234,0.555) -- cycle;
\fill[pthick] (3.234,0.744) -- (2.942,0.649) -- (3.234,1.051) -- cycle;
\fill[pthin] (2.942,0.649) -- (3.234,0.744) -- (3.234,0.555) -- cycle;
\fill[pthick] (3.234,0.059) -- (3.234,-0.248) -- (2.942,0.153) -- cycle;
\fill[pthick] (3.053,0.307) -- (3.234,0.059) -- (3.234,0.555) -- cycle;
\fill[pthin] (3.234,0.059) -- (3.053,0.307) -- (2.942,0.153) -- cycle;
\fill[pthin] (2.762,0.401) -- (3.053,0.307) -- (2.942,0.153) -- cycle;
\fill[pthick] (3.053,0.307) -- (2.762,0.401) -- (3.234,0.555) -- cycle;
\fill[pthin] (2.762,0.401) -- (2.762,0.095) -- (2.942,0.153) -- cycle;
\fill[pthick] (2.762,0.095) -- (2.762,0.401) -- (2.470,0.000) -- cycle;
\fill[pthick] (2.942,-0.153) -- (3.234,-0.248) -- (2.942,0.153) -- cycle;
\fill[pthick] (2.762,0.095) -- (2.942,-0.153) -- (2.470,0.000) -- cycle;
\fill[pthin] (2.942,-0.153) -- (2.762,0.095) -- (2.942,0.153) -- cycle;
\fill[pthick] (1.235,0.095) -- (1.235,0.401) -- (1.527,0.000) -- cycle;
\fill[pthick] (1.415,-0.153) -- (1.235,0.095) -- (1.235,-0.401) -- cycle;
\fill[pthin] (1.235,0.095) -- (1.415,-0.153) -- (1.527,0.000) -- cycle;
\fill[pthick] (1.818,0.401) -- (1.998,0.649) -- (1.527,0.496) -- cycle;
\fill[pthick] (1.527,0.307) -- (1.818,0.401) -- (1.527,0.000) -- cycle;
\fill[pthin] (1.818,0.401) -- (1.527,0.307) -- (1.527,0.496) -- cycle;
\fill[pthin] (1.235,0.401) -- (1.527,0.307) -- (1.527,0.496) -- cycle;
\fill[pthick] (1.527,0.307) -- (1.235,0.401) -- (1.527,0.000) -- cycle;
\fill[pthick] (0.472,-0.153) -- (0.763,-0.248) -- (0.472,0.153) -- cycle;
\fill[pthick] (0.292,0.095) -- (0.472,-0.153) -- (0.000,0.000) -- cycle;
\fill[pthin] (0.472,-0.153) -- (0.292,0.095) -- (0.472,0.153) -- cycle;
\fill[pthick] (0.944,0.307) -- (1.235,0.401) -- (0.944,0.000) -- cycle;
\fill[pthick] (0.763,0.059) -- (0.944,0.307) -- (0.472,0.153) -- cycle;
\fill[pthin] (0.944,0.307) -- (0.763,0.059) -- (0.944,0.000) -- cycle;
\fill[pthin] (0.763,-0.248) -- (0.763,0.059) -- (0.944,0.000) -- cycle;
\fill[pthick] (0.763,0.059) -- (0.763,-0.248) -- (0.472,0.153) -- cycle;
\fill[pthin] (0.763,-0.248) -- (1.055,-0.153) -- (0.944,0.000) -- cycle;
\fill[pthick] (1.055,-0.153) -- (0.763,-0.248) -- (1.235,-0.401) -- cycle;
\fill[pthick] (1.235,0.095) -- (1.235,0.401) -- (0.944,0.000) -- cycle;
\fill[pthick] (1.055,-0.153) -- (1.235,0.095) -- (1.235,-0.401) -- cycle;
\fill[pthin] (1.235,0.095) -- (1.055,-0.153) -- (0.944,0.000) -- cycle;
\fill[pthick] (2.470,-0.496) -- (2.470,-0.803) -- (2.762,-0.401) -- cycle;
\fill[pthick] (2.650,-0.248) -- (2.470,-0.496) -- (2.470,0.000) -- cycle;
\fill[pthin] (2.470,-0.496) -- (2.650,-0.248) -- (2.762,-0.401) -- cycle;
\fill[pthick] (3.053,-0.803) -- (3.234,-1.051) -- (2.762,-0.897) -- cycle;
\fill[pthick] (2.762,-0.708) -- (3.053,-0.803) -- (2.762,-0.401) -- cycle;
\fill[pthin] (3.053,-0.803) -- (2.762,-0.708) -- (2.762,-0.897) -- cycle;
\fill[pthin] (2.470,-0.803) -- (2.762,-0.708) -- (2.762,-0.897) -- cycle;
\fill[pthick] (2.762,-0.708) -- (2.470,-0.803) -- (2.762,-0.401) -- cycle;
\fill[pthick] (1.707,-0.248) -- (1.998,-0.153) -- (1.707,-0.555) -- cycle;
\fill[pthick] (1.527,-0.496) -- (1.707,-0.248) -- (1.235,-0.401) -- cycle;
\fill[pthin] (1.707,-0.248) -- (1.527,-0.496) -- (1.707,-0.555) -- cycle;
\fill[pthick] (2.179,-0.708) -- (2.470,-0.803) -- (2.179,-0.401) -- cycle;
\fill[pthick] (1.998,-0.460) -- (2.179,-0.708) -- (1.707,-0.555) -- cycle;
\fill[pthin] (2.179,-0.708) -- (1.998,-0.460) -- (2.179,-0.401) -- cycle;
\fill[pthin] (1.998,-0.153) -- (1.998,-0.460) -- (2.179,-0.401) -- cycle;
\fill[pthick] (1.998,-0.460) -- (1.998,-0.153) -- (1.707,-0.555) -- cycle;
\fill[pthin] (1.998,-0.153) -- (2.290,-0.248) -- (2.179,-0.401) -- cycle;
\fill[pthick] (2.290,-0.248) -- (1.998,-0.153) -- (2.470,0.000) -- cycle;
\fill[pthick] (2.470,-0.496) -- (2.470,-0.803) -- (2.179,-0.401) -- cycle;
\fill[pthick] (2.290,-0.248) -- (2.470,-0.496) -- (2.470,0.000) -- cycle;
\fill[pthin] (2.470,-0.496) -- (2.290,-0.248) -- (2.179,-0.401) -- cycle;
\fill[pthin] (1.998,-0.153) -- (2.179,0.095) -- (1.998,0.153) -- cycle;
\fill[pthick] (2.179,0.095) -- (1.998,-0.153) -- (2.470,0.000) -- cycle;
\fill[pthick] (2.179,0.401) -- (1.998,0.649) -- (1.998,0.153) -- cycle;
\fill[pthick] (2.179,0.095) -- (2.179,0.401) -- (2.470,0.000) -- cycle;
\fill[pthin] (2.179,0.401) -- (2.179,0.095) -- (1.998,0.153) -- cycle;
\fill[pthick] (1.707,-0.248) -- (1.998,-0.153) -- (1.527,0.000) -- cycle;
\fill[pthick] (1.415,-0.153) -- (1.707,-0.248) -- (1.235,-0.401) -- cycle;
\fill[pthin] (1.707,-0.248) -- (1.415,-0.153) -- (1.527,0.000) -- cycle;
\fill[pthick] (1.818,0.401) -- (1.998,0.649) -- (1.998,0.153) -- cycle;
\fill[pthick] (1.818,0.095) -- (1.818,0.401) -- (1.527,0.000) -- cycle;
\fill[pthin] (1.818,0.401) -- (1.818,0.095) -- (1.998,0.153) -- cycle;
\fill[pthin] (1.998,-0.153) -- (1.818,0.095) -- (1.998,0.153) -- cycle;
\fill[pthick] (1.818,0.095) -- (1.998,-0.153) -- (1.527,0.000) -- cycle;
\fill[pthin] (2.470,1.299) -- (2.179,1.204) -- (2.290,1.051) -- cycle;
\fill[pthick] (2.179,1.204) -- (2.470,1.299) -- (1.998,1.452) -- cycle;
\fill[pthick] (1.998,0.956) -- (1.998,0.649) -- (2.290,1.051) -- cycle;
\fill[pthick] (2.179,1.204) -- (1.998,0.956) -- (1.998,1.452) -- cycle;
\fill[pthin] (1.998,0.956) -- (2.179,1.204) -- (2.290,1.051) -- cycle;
\fill[pthick] (2.762,1.204) -- (2.470,1.299) -- (2.762,0.897) -- cycle;
\fill[pthick] (2.942,0.956) -- (2.762,1.204) -- (3.234,1.051) -- cycle;
\fill[pthin] (2.762,1.204) -- (2.942,0.956) -- (2.762,0.897) -- cycle;
\fill[pthick] (2.290,0.744) -- (1.998,0.649) -- (2.290,1.051) -- cycle;
\fill[pthick] (2.470,0.992) -- (2.290,0.744) -- (2.762,0.897) -- cycle;
\fill[pthin] (2.290,0.744) -- (2.470,0.992) -- (2.290,1.051) -- cycle;
\fill[pthin] (2.470,1.299) -- (2.470,0.992) -- (2.290,1.051) -- cycle;
\fill[pthick] (2.470,0.992) -- (2.470,1.299) -- (2.762,0.897) -- cycle;
\fill[pthick] (2.290,1.853) -- (2.470,2.101) -- (1.998,1.948) -- cycle;
\fill[pthick] (1.998,1.759) -- (2.290,1.853) -- (1.998,1.452) -- cycle;
\fill[pthin] (2.290,1.853) -- (1.998,1.759) -- (1.998,1.948) -- cycle;
\fill[pthick] (1.998,2.444) -- (1.998,2.751) -- (2.290,2.349) -- cycle;
\fill[pthick] (2.179,2.196) -- (1.998,2.444) -- (1.998,1.948) -- cycle;
\fill[pthin] (1.998,2.444) -- (2.179,2.196) -- (2.290,2.349) -- cycle;
\fill[pthin] (2.470,2.101) -- (2.179,2.196) -- (2.290,2.349) -- cycle;
\fill[pthick] (2.179,2.196) -- (2.470,2.101) -- (1.998,1.948) -- cycle;
\fill[pthick] (2.762,1.204) -- (2.470,1.299) -- (2.942,1.452) -- cycle;
\fill[pthick] (3.053,1.299) -- (2.762,1.204) -- (3.234,1.051) -- cycle;
\fill[pthin] (2.762,1.204) -- (3.053,1.299) -- (2.942,1.452) -- cycle;
\fill[pthick] (2.650,1.853) -- (2.470,2.101) -- (2.470,1.605) -- cycle;
\fill[pthick] (2.650,1.547) -- (2.650,1.853) -- (2.942,1.452) -- cycle;
\fill[pthin] (2.650,1.853) -- (2.650,1.547) -- (2.470,1.605) -- cycle;
\fill[pthin] (2.470,1.299) -- (2.650,1.547) -- (2.470,1.605) -- cycle;
\fill[pthick] (2.650,1.547) -- (2.470,1.299) -- (2.942,1.452) -- cycle;
\fill[pthin] (2.470,1.299) -- (2.290,1.547) -- (2.470,1.605) -- cycle;
\fill[pthick] (2.290,1.547) -- (2.470,1.299) -- (1.998,1.452) -- cycle;
\fill[pthick] (2.290,1.853) -- (2.470,2.101) -- (2.470,1.605) -- cycle;
\fill[pthick] (2.290,1.547) -- (2.290,1.853) -- (1.998,1.452) -- cycle;
\fill[pthin] (2.290,1.853) -- (2.290,1.547) -- (2.470,1.605) -- cycle;
\fill[pthick] (1.055,0.649) -- (1.235,0.401) -- (1.235,0.897) -- cycle;
\fill[pthick] (1.055,0.956) -- (1.055,0.649) -- (0.763,1.051) -- cycle;
\fill[pthin] (1.055,0.649) -- (1.055,0.956) -- (1.235,0.897) -- cycle;
\fill[pthick] (1.707,0.744) -- (1.998,0.649) -- (1.527,0.496) -- cycle;
\fill[pthick] (1.415,0.649) -- (1.707,0.744) -- (1.235,0.897) -- cycle;
\fill[pthin] (1.707,0.744) -- (1.415,0.649) -- (1.527,0.496) -- cycle;
\fill[pthin] (1.235,0.401) -- (1.415,0.649) -- (1.527,0.496) -- cycle;
\fill[pthick] (1.415,0.649) -- (1.235,0.401) -- (1.235,0.897) -- cycle;
\fill[pthick] (0.292,0.401) -- (0.472,0.649) -- (0.472,0.153) -- cycle;
\fill[pthick] (0.292,0.095) -- (0.292,0.401) -- (0.000,0.000) -- cycle;
\fill[pthin] (0.292,0.401) -- (0.292,0.095) -- (0.472,0.153) -- cycle;
\fill[pthick] (0.944,0.307) -- (1.235,0.401) -- (0.763,0.555) -- cycle;
\fill[pthick] (0.652,0.401) -- (0.944,0.307) -- (0.472,0.153) -- cycle;
\fill[pthin] (0.944,0.307) -- (0.652,0.401) -- (0.763,0.555) -- cycle;
\fill[pthin] (0.472,0.649) -- (0.652,0.401) -- (0.763,0.555) -- cycle;
\fill[pthick] (0.652,0.401) -- (0.472,0.649) -- (0.472,0.153) -- cycle;
\fill[pthin] (0.472,0.649) -- (0.763,0.744) -- (0.763,0.555) -- cycle;
\fill[pthick] (0.763,0.744) -- (0.472,0.649) -- (0.763,1.051) -- cycle;
\fill[pthick] (1.055,0.649) -- (1.235,0.401) -- (0.763,0.555) -- cycle;
\fill[pthick] (0.763,0.744) -- (1.055,0.649) -- (0.763,1.051) -- cycle;
\fill[pthin] (1.055,0.649) -- (0.763,0.744) -- (0.763,0.555) -- cycle;
\fill[pthick] (1.707,1.853) -- (1.527,2.101) -- (1.998,1.948) -- cycle;
\fill[pthick] (1.998,1.759) -- (1.707,1.853) -- (1.998,1.452) -- cycle;
\fill[pthin] (1.707,1.853) -- (1.998,1.759) -- (1.998,1.948) -- cycle;
\fill[pthick] (1.998,2.444) -- (1.998,2.751) -- (1.707,2.349) -- cycle;
\fill[pthick] (1.818,2.196) -- (1.998,2.444) -- (1.998,1.948) -- cycle;
\fill[pthin] (1.998,2.444) -- (1.818,2.196) -- (1.707,2.349) -- cycle;
\fill[pthin] (1.527,2.101) -- (1.818,2.196) -- (1.707,2.349) -- cycle;
\fill[pthick] (1.818,2.196) -- (1.527,2.101) -- (1.998,1.948) -- cycle;
\fill[pthick] (1.235,1.204) -- (1.527,1.299) -- (1.055,1.452) -- cycle;
\fill[pthick] (0.944,1.299) -- (1.235,1.204) -- (0.763,1.051) -- cycle;
\fill[pthin] (1.235,1.204) -- (0.944,1.299) -- (1.055,1.452) -- cycle;
\fill[pthick] (1.346,1.853) -- (1.527,2.101) -- (1.527,1.605) -- cycle;
\fill[pthick] (1.346,1.547) -- (1.346,1.853) -- (1.055,1.452) -- cycle;
\fill[pthin] (1.346,1.853) -- (1.346,1.547) -- (1.527,1.605) -- cycle;
\fill[pthin] (1.527,1.299) -- (1.346,1.547) -- (1.527,1.605) -- cycle;
\fill[pthick] (1.346,1.547) -- (1.527,1.299) -- (1.055,1.452) -- cycle;
\fill[pthin] (1.527,1.299) -- (1.707,1.547) -- (1.527,1.605) -- cycle;
\fill[pthick] (1.707,1.547) -- (1.527,1.299) -- (1.998,1.452) -- cycle;
\fill[pthick] (1.707,1.853) -- (1.527,2.101) -- (1.527,1.605) -- cycle;
\fill[pthick] (1.707,1.547) -- (1.707,1.853) -- (1.998,1.452) -- cycle;
\fill[pthin] (1.707,1.853) -- (1.707,1.547) -- (1.527,1.605) -- cycle;
\fill[pthin] (1.527,1.299) -- (1.818,1.204) -- (1.707,1.051) -- cycle;
\fill[pthick] (1.818,1.204) -- (1.527,1.299) -- (1.998,1.452) -- cycle;
\fill[pthick] (1.998,0.956) -- (1.998,0.649) -- (1.707,1.051) -- cycle;
\fill[pthick] (1.818,1.204) -- (1.998,0.956) -- (1.998,1.452) -- cycle;
\fill[pthin] (1.998,0.956) -- (1.818,1.204) -- (1.707,1.051) -- cycle;
\fill[pthick] (1.235,1.204) -- (1.527,1.299) -- (1.235,0.897) -- cycle;
\fill[pthick] (1.055,0.956) -- (1.235,1.204) -- (0.763,1.051) -- cycle;
\fill[pthin] (1.235,1.204) -- (1.055,0.956) -- (1.235,0.897) -- cycle;
\fill[pthick] (1.707,0.744) -- (1.998,0.649) -- (1.707,1.051) -- cycle;
\fill[pthick] (1.527,0.992) -- (1.707,0.744) -- (1.235,0.897) -- cycle;
\fill[pthin] (1.707,0.744) -- (1.527,0.992) -- (1.707,1.051) -- cycle;
\fill[pthin] (1.527,1.299) -- (1.527,0.992) -- (1.707,1.051) -- cycle;
\fill[pthick] (1.527,0.992) -- (1.527,1.299) -- (1.235,0.897) -- cycle;
\fill[pthin] (0.472,2.751) -- (0.472,2.444) -- (0.292,2.503) -- cycle;
\fill[pthick] (0.472,2.444) -- (0.472,2.751) -- (0.763,2.349) -- cycle;
\fill[pthick] (0.292,2.196) -- (0.000,2.101) -- (0.292,2.503) -- cycle;
\fill[pthick] (0.472,2.444) -- (0.292,2.196) -- (0.763,2.349) -- cycle;
\fill[pthin] (0.292,2.196) -- (0.472,2.444) -- (0.292,2.503) -- cycle;
\fill[pthick] (0.292,2.999) -- (0.472,2.751) -- (0.000,2.904) -- cycle;
\fill[pthick] (0.000,3.093) -- (0.292,2.999) -- (0.000,3.400) -- cycle;
\fill[pthin] (0.292,2.999) -- (0.000,3.093) -- (0.000,2.904) -- cycle;
\fill[pthick] (0.000,2.408) -- (0.000,2.101) -- (0.292,2.503) -- cycle;
\fill[pthick] (0.180,2.656) -- (0.000,2.408) -- (0.000,2.904) -- cycle;
\fill[pthin] (0.000,2.408) -- (0.180,2.656) -- (0.292,2.503) -- cycle;
\fill[pthin] (0.472,2.751) -- (0.180,2.656) -- (0.292,2.503) -- cycle;
\fill[pthick] (0.180,2.656) -- (0.472,2.751) -- (0.000,2.904) -- cycle;
\fill[pthick] (1.055,2.751) -- (1.235,2.999) -- (1.235,2.503) -- cycle;
\fill[pthick] (1.055,2.444) -- (1.055,2.751) -- (0.763,2.349) -- cycle;
\fill[pthin] (1.055,2.751) -- (1.055,2.444) -- (1.235,2.503) -- cycle;
\fill[pthick] (1.707,2.656) -- (1.998,2.751) -- (1.527,2.904) -- cycle;
\fill[pthick] (1.415,2.751) -- (1.707,2.656) -- (1.235,2.503) -- cycle;
\fill[pthin] (1.707,2.656) -- (1.415,2.751) -- (1.527,2.904) -- cycle;
\fill[pthin] (1.235,2.999) -- (1.415,2.751) -- (1.527,2.904) -- cycle;
\fill[pthick] (1.415,2.751) -- (1.235,2.999) -- (1.235,2.503) -- cycle;
\fill[pthick] (0.292,2.999) -- (0.472,2.751) -- (0.472,3.247) -- cycle;
\fill[pthick] (0.292,3.305) -- (0.292,2.999) -- (0.000,3.400) -- cycle;
\fill[pthin] (0.292,2.999) -- (0.292,3.305) -- (0.472,3.247) -- cycle;
\fill[pthick] (0.944,3.093) -- (1.235,2.999) -- (0.763,2.845) -- cycle;
\fill[pthick] (0.652,2.999) -- (0.944,3.093) -- (0.472,3.247) -- cycle;
\fill[pthin] (0.944,3.093) -- (0.652,2.999) -- (0.763,2.845) -- cycle;
\fill[pthin] (0.472,2.751) -- (0.652,2.999) -- (0.763,2.845) -- cycle;
\fill[pthick] (0.652,2.999) -- (0.472,2.751) -- (0.472,3.247) -- cycle;
\fill[pthin] (0.472,2.751) -- (0.763,2.656) -- (0.763,2.845) -- cycle;
\fill[pthick] (0.763,2.656) -- (0.472,2.751) -- (0.763,2.349) -- cycle;
\fill[pthick] (1.055,2.751) -- (1.235,2.999) -- (0.763,2.845) -- cycle;
\fill[pthick] (0.763,2.656) -- (1.055,2.751) -- (0.763,2.349) -- cycle;
\fill[pthin] (1.055,2.751) -- (0.763,2.656) -- (0.763,2.845) -- cycle;
\fill[pthick] (0.292,1.204) -- (0.000,1.299) -- (0.472,1.452) -- cycle;
\fill[pthick] (0.583,1.299) -- (0.292,1.204) -- (0.763,1.051) -- cycle;
\fill[pthin] (0.292,1.204) -- (0.583,1.299) -- (0.472,1.452) -- cycle;
\fill[pthick] (0.180,1.853) -- (0.000,2.101) -- (0.000,1.605) -- cycle;
\fill[pthick] (0.180,1.547) -- (0.180,1.853) -- (0.472,1.452) -- cycle;
\fill[pthin] (0.180,1.853) -- (0.180,1.547) -- (0.000,1.605) -- cycle;
\fill[pthin] (0.000,1.299) -- (0.180,1.547) -- (0.000,1.605) -- cycle;
\fill[pthick] (0.180,1.547) -- (0.000,1.299) -- (0.472,1.452) -- cycle;
\fill[pthick] (0.292,0.401) -- (0.472,0.649) -- (0.000,0.496) -- cycle;
\fill[pthick] (0.000,0.307) -- (0.292,0.401) -- (0.000,0.000) -- cycle;
\fill[pthin] (0.292,0.401) -- (0.000,0.307) -- (0.000,0.496) -- cycle;
\fill[pthick] (0.000,0.992) -- (0.000,1.299) -- (0.292,0.897) -- cycle;
\fill[pthick] (0.180,0.744) -- (0.000,0.992) -- (0.000,0.496) -- cycle;
\fill[pthin] (0.000,0.992) -- (0.180,0.744) -- (0.292,0.897) -- cycle;
\fill[pthin] (0.472,0.649) -- (0.180,0.744) -- (0.292,0.897) -- cycle;
\fill[pthick] (0.180,0.744) -- (0.472,0.649) -- (0.000,0.496) -- cycle;
\fill[pthin] (0.472,0.649) -- (0.472,0.956) -- (0.292,0.897) -- cycle;
\fill[pthick] (0.472,0.956) -- (0.472,0.649) -- (0.763,1.051) -- cycle;
\fill[pthick] (0.292,1.204) -- (0.000,1.299) -- (0.292,0.897) -- cycle;
\fill[pthick] (0.472,0.956) -- (0.292,1.204) -- (0.763,1.051) -- cycle;
\fill[pthin] (0.292,1.204) -- (0.472,0.956) -- (0.292,0.897) -- cycle;
\fill[pthick] (1.235,2.196) -- (1.527,2.101) -- (1.235,2.503) -- cycle;
\fill[pthick] (1.055,2.444) -- (1.235,2.196) -- (0.763,2.349) -- cycle;
\fill[pthin] (1.235,2.196) -- (1.055,2.444) -- (1.235,2.503) -- cycle;
\fill[pthick] (1.707,2.656) -- (1.998,2.751) -- (1.707,2.349) -- cycle;
\fill[pthick] (1.527,2.408) -- (1.707,2.656) -- (1.235,2.503) -- cycle;
\fill[pthin] (1.707,2.656) -- (1.527,2.408) -- (1.707,2.349) -- cycle;
\fill[pthin] (1.527,2.101) -- (1.527,2.408) -- (1.707,2.349) -- cycle;
\fill[pthick] (1.527,2.408) -- (1.527,2.101) -- (1.235,2.503) -- cycle;
\fill[pthick] (0.763,1.547) -- (0.763,1.853) -- (1.055,1.452) -- cycle;
\fill[pthick] (0.944,1.299) -- (0.763,1.547) -- (0.763,1.051) -- cycle;
\fill[pthin] (0.763,1.547) -- (0.944,1.299) -- (1.055,1.452) -- cycle;
\fill[pthick] (1.346,1.853) -- (1.527,2.101) -- (1.055,1.948) -- cycle;
\fill[pthick] (1.055,1.759) -- (1.346,1.853) -- (1.055,1.452) -- cycle;
\fill[pthin] (1.346,1.853) -- (1.055,1.759) -- (1.055,1.948) -- cycle;
\fill[pthin] (0.763,1.853) -- (1.055,1.759) -- (1.055,1.948) -- cycle;
\fill[pthick] (1.055,1.759) -- (0.763,1.853) -- (1.055,1.452) -- cycle;
\fill[pthin] (0.763,1.853) -- (0.944,2.101) -- (1.055,1.948) -- cycle;
\fill[pthick] (0.944,2.101) -- (0.763,1.853) -- (0.763,2.349) -- cycle;
\fill[pthick] (1.235,2.196) -- (1.527,2.101) -- (1.055,1.948) -- cycle;
\fill[pthick] (0.944,2.101) -- (1.235,2.196) -- (0.763,2.349) -- cycle;
\fill[pthin] (1.235,2.196) -- (0.944,2.101) -- (1.055,1.948) -- cycle;
\fill[pthin] (0.763,1.853) -- (0.583,2.101) -- (0.472,1.948) -- cycle;
\fill[pthick] (0.583,2.101) -- (0.763,1.853) -- (0.763,2.349) -- cycle;
\fill[pthick] (0.292,2.196) -- (0.000,2.101) -- (0.472,1.948) -- cycle;
\fill[pthick] (0.583,2.101) -- (0.292,2.196) -- (0.763,2.349) -- cycle;
\fill[pthin] (0.292,2.196) -- (0.583,2.101) -- (0.472,1.948) -- cycle;
\fill[pthick] (0.763,1.547) -- (0.763,1.853) -- (0.472,1.452) -- cycle;
\fill[pthick] (0.583,1.299) -- (0.763,1.547) -- (0.763,1.051) -- cycle;
\fill[pthin] (0.763,1.547) -- (0.583,1.299) -- (0.472,1.452) -- cycle;
\fill[pthick] (0.180,1.853) -- (0.000,2.101) -- (0.472,1.948) -- cycle;
\fill[pthick] (0.472,1.759) -- (0.180,1.853) -- (0.472,1.452) -- cycle;
\fill[pthin] (0.180,1.853) -- (0.472,1.759) -- (0.472,1.948) -- cycle;
\fill[pthin] (0.763,1.853) -- (0.472,1.759) -- (0.472,1.948) -- cycle;
\fill[pthick] (0.472,1.759) -- (0.763,1.853) -- (0.472,1.452) -- cycle;
\fill[pthin] (-0.472,2.751) -- (-0.472,2.444) -- (-0.292,2.503) -- cycle;
\fill[pthick] (-0.472,2.444) -- (-0.472,2.751) -- (-0.763,2.349) -- cycle;
\fill[pthick] (-0.292,2.196) -- (0.000,2.101) -- (-0.292,2.503) -- cycle;
\fill[pthick] (-0.472,2.444) -- (-0.292,2.196) -- (-0.763,2.349) -- cycle;
\fill[pthin] (-0.292,2.196) -- (-0.472,2.444) -- (-0.292,2.503) -- cycle;
\fill[pthick] (-0.292,2.999) -- (-0.472,2.751) -- (0.000,2.904) -- cycle;
\fill[pthick] (0.000,3.093) -- (-0.292,2.999) -- (0.000,3.400) -- cycle;
\fill[pthin] (-0.292,2.999) -- (0.000,3.093) -- (0.000,2.904) -- cycle;
\fill[pthick] (0.000,2.408) -- (0.000,2.101) -- (-0.292,2.503) -- cycle;
\fill[pthick] (-0.180,2.656) -- (0.000,2.408) -- (0.000,2.904) -- cycle;
\fill[pthin] (0.000,2.408) -- (-0.180,2.656) -- (-0.292,2.503) -- cycle;
\fill[pthin] (-0.472,2.751) -- (-0.180,2.656) -- (-0.292,2.503) -- cycle;
\fill[pthick] (-0.180,2.656) -- (-0.472,2.751) -- (0.000,2.904) -- cycle;
\fill[pthick] (-1.055,2.751) -- (-1.235,2.999) -- (-1.235,2.503) -- cycle;
\fill[pthick] (-1.055,2.444) -- (-1.055,2.751) -- (-0.763,2.349) -- cycle;
\fill[pthin] (-1.055,2.751) -- (-1.055,2.444) -- (-1.235,2.503) -- cycle;
\fill[pthick] (-1.707,2.656) -- (-1.998,2.751) -- (-1.527,2.904) -- cycle;
\fill[pthick] (-1.415,2.751) -- (-1.707,2.656) -- (-1.235,2.503) -- cycle;
\fill[pthin] (-1.707,2.656) -- (-1.415,2.751) -- (-1.527,2.904) -- cycle;
\fill[pthin] (-1.235,2.999) -- (-1.415,2.751) -- (-1.527,2.904) -- cycle;
\fill[pthick] (-1.415,2.751) -- (-1.235,2.999) -- (-1.235,2.503) -- cycle;
\fill[pthick] (-0.292,2.999) -- (-0.472,2.751) -- (-0.472,3.247) -- cycle;
\fill[pthick] (-0.292,3.305) -- (-0.292,2.999) -- (0.000,3.400) -- cycle;
\fill[pthin] (-0.292,2.999) -- (-0.292,3.305) -- (-0.472,3.247) -- cycle;
\fill[pthick] (-0.944,3.093) -- (-1.235,2.999) -- (-0.763,2.845) -- cycle;
\fill[pthick] (-0.652,2.999) -- (-0.944,3.093) -- (-0.472,3.247) -- cycle;
\fill[pthin] (-0.944,3.093) -- (-0.652,2.999) -- (-0.763,2.845) -- cycle;
\fill[pthin] (-0.472,2.751) -- (-0.652,2.999) -- (-0.763,2.845) -- cycle;
\fill[pthick] (-0.652,2.999) -- (-0.472,2.751) -- (-0.472,3.247) -- cycle;
\fill[pthin] (-0.472,2.751) -- (-0.763,2.656) -- (-0.763,2.845) -- cycle;
\fill[pthick] (-0.763,2.656) -- (-0.472,2.751) -- (-0.763,2.349) -- cycle;
\fill[pthick] (-1.055,2.751) -- (-1.235,2.999) -- (-0.763,2.845) -- cycle;
\fill[pthick] (-0.763,2.656) -- (-1.055,2.751) -- (-0.763,2.349) -- cycle;
\fill[pthin] (-1.055,2.751) -- (-0.763,2.656) -- (-0.763,2.845) -- cycle;
\fill[pthick] (-0.292,1.204) -- (0.000,1.299) -- (-0.472,1.452) -- cycle;
\fill[pthick] (-0.583,1.299) -- (-0.292,1.204) -- (-0.763,1.051) -- cycle;
\fill[pthin] (-0.292,1.204) -- (-0.583,1.299) -- (-0.472,1.452) -- cycle;
\fill[pthick] (-0.180,1.853) -- (0.000,2.101) -- (0.000,1.605) -- cycle;
\fill[pthick] (-0.180,1.547) -- (-0.180,1.853) -- (-0.472,1.452) -- cycle;
\fill[pthin] (-0.180,1.853) -- (-0.180,1.547) -- (0.000,1.605) -- cycle;
\fill[pthin] (0.000,1.299) -- (-0.180,1.547) -- (0.000,1.605) -- cycle;
\fill[pthick] (-0.180,1.547) -- (0.000,1.299) -- (-0.472,1.452) -- cycle;
\fill[pthick] (-0.292,0.401) -- (-0.472,0.649) -- (0.000,0.496) -- cycle;
\fill[pthick] (0.000,0.307) -- (-0.292,0.401) -- (0.000,0.000) -- cycle;
\fill[pthin] (-0.292,0.401) -- (0.000,0.307) -- (0.000,0.496) -- cycle;
\fill[pthick] (0.000,0.992) -- (0.000,1.299) -- (-0.292,0.897) -- cycle;
\fill[pthick] (-0.180,0.744) -- (0.000,0.992) -- (0.000,0.496) -- cycle;
\fill[pthin] (0.000,0.992) -- (-0.180,0.744) -- (-0.292,0.897) -- cycle;
\fill[pthin] (-0.472,0.649) -- (-0.180,0.744) -- (-0.292,0.897) -- cycle;
\fill[pthick] (-0.180,0.744) -- (-0.472,0.649) -- (0.000,0.496) -- cycle;
\fill[pthin] (-0.472,0.649) -- (-0.472,0.956) -- (-0.292,0.897) -- cycle;
\fill[pthick] (-0.472,0.956) -- (-0.472,0.649) -- (-0.763,1.051) -- cycle;
\fill[pthick] (-0.292,1.204) -- (0.000,1.299) -- (-0.292,0.897) -- cycle;
\fill[pthick] (-0.472,0.956) -- (-0.292,1.204) -- (-0.763,1.051) -- cycle;
\fill[pthin] (-0.292,1.204) -- (-0.472,0.956) -- (-0.292,0.897) -- cycle;
\fill[pthick] (-1.235,2.196) -- (-1.527,2.101) -- (-1.235,2.503) -- cycle;
\fill[pthick] (-1.055,2.444) -- (-1.235,2.196) -- (-0.763,2.349) -- cycle;
\fill[pthin] (-1.235,2.196) -- (-1.055,2.444) -- (-1.235,2.503) -- cycle;
\fill[pthick] (-1.707,2.656) -- (-1.998,2.751) -- (-1.707,2.349) -- cycle;
\fill[pthick] (-1.527,2.408) -- (-1.707,2.656) -- (-1.235,2.503) -- cycle;
\fill[pthin] (-1.707,2.656) -- (-1.527,2.408) -- (-1.707,2.349) -- cycle;
\fill[pthin] (-1.527,2.101) -- (-1.527,2.408) -- (-1.707,2.349) -- cycle;
\fill[pthick] (-1.527,2.408) -- (-1.527,2.101) -- (-1.235,2.503) -- cycle;
\fill[pthick] (-0.763,1.547) -- (-0.763,1.853) -- (-1.055,1.452) -- cycle;
\fill[pthick] (-0.944,1.299) -- (-0.763,1.547) -- (-0.763,1.051) -- cycle;
\fill[pthin] (-0.763,1.547) -- (-0.944,1.299) -- (-1.055,1.452) -- cycle;
\fill[pthick] (-1.346,1.853) -- (-1.527,2.101) -- (-1.055,1.948) -- cycle;
\fill[pthick] (-1.055,1.759) -- (-1.346,1.853) -- (-1.055,1.452) -- cycle;
\fill[pthin] (-1.346,1.853) -- (-1.055,1.759) -- (-1.055,1.948) -- cycle;
\fill[pthin] (-0.763,1.853) -- (-1.055,1.759) -- (-1.055,1.948) -- cycle;
\fill[pthick] (-1.055,1.759) -- (-0.763,1.853) -- (-1.055,1.452) -- cycle;
\fill[pthin] (-0.763,1.853) -- (-0.944,2.101) -- (-1.055,1.948) -- cycle;
\fill[pthick] (-0.944,2.101) -- (-0.763,1.853) -- (-0.763,2.349) -- cycle;
\fill[pthick] (-1.235,2.196) -- (-1.527,2.101) -- (-1.055,1.948) -- cycle;
\fill[pthick] (-0.944,2.101) -- (-1.235,2.196) -- (-0.763,2.349) -- cycle;
\fill[pthin] (-1.235,2.196) -- (-0.944,2.101) -- (-1.055,1.948) -- cycle;
\fill[pthin] (-0.763,1.853) -- (-0.583,2.101) -- (-0.472,1.948) -- cycle;
\fill[pthick] (-0.583,2.101) -- (-0.763,1.853) -- (-0.763,2.349) -- cycle;
\fill[pthick] (-0.292,2.196) -- (0.000,2.101) -- (-0.472,1.948) -- cycle;
\fill[pthick] (-0.583,2.101) -- (-0.292,2.196) -- (-0.763,2.349) -- cycle;
\fill[pthin] (-0.292,2.196) -- (-0.583,2.101) -- (-0.472,1.948) -- cycle;
\fill[pthick] (-0.763,1.547) -- (-0.763,1.853) -- (-0.472,1.452) -- cycle;
\fill[pthick] (-0.583,1.299) -- (-0.763,1.547) -- (-0.763,1.051) -- cycle;
\fill[pthin] (-0.763,1.547) -- (-0.583,1.299) -- (-0.472,1.452) -- cycle;
\fill[pthick] (-0.180,1.853) -- (0.000,2.101) -- (-0.472,1.948) -- cycle;
\fill[pthick] (-0.472,1.759) -- (-0.180,1.853) -- (-0.472,1.452) -- cycle;
\fill[pthin] (-0.180,1.853) -- (-0.472,1.759) -- (-0.472,1.948) -- cycle;
\fill[pthin] (-0.763,1.853) -- (-0.472,1.759) -- (-0.472,1.948) -- cycle;
\fill[pthick] (-0.472,1.759) -- (-0.763,1.853) -- (-0.472,1.452) -- cycle;
\fill[pthin] (-2.470,1.299) -- (-2.179,1.204) -- (-2.290,1.051) -- cycle;
\fill[pthick] (-2.179,1.204) -- (-2.470,1.299) -- (-1.998,1.452) -- cycle;
\fill[pthick] (-1.998,0.956) -- (-1.998,0.649) -- (-2.290,1.051) -- cycle;
\fill[pthick] (-2.179,1.204) -- (-1.998,0.956) -- (-1.998,1.452) -- cycle;
\fill[pthin] (-1.998,0.956) -- (-2.179,1.204) -- (-2.290,1.051) -- cycle;
\fill[pthick] (-2.762,1.204) -- (-2.470,1.299) -- (-2.762,0.897) -- cycle;
\fill[pthick] (-2.942,0.956) -- (-2.762,1.204) -- (-3.234,1.051) -- cycle;
\fill[pthin] (-2.762,1.204) -- (-2.942,0.956) -- (-2.762,0.897) -- cycle;
\fill[pthick] (-2.290,0.744) -- (-1.998,0.649) -- (-2.290,1.051) -- cycle;
\fill[pthick] (-2.470,0.992) -- (-2.290,0.744) -- (-2.762,0.897) -- cycle;
\fill[pthin] (-2.290,0.744) -- (-2.470,0.992) -- (-2.290,1.051) -- cycle;
\fill[pthin] (-2.470,1.299) -- (-2.470,0.992) -- (-2.290,1.051) -- cycle;
\fill[pthick] (-2.470,0.992) -- (-2.470,1.299) -- (-2.762,0.897) -- cycle;
\fill[pthick] (-2.290,1.853) -- (-2.470,2.101) -- (-1.998,1.948) -- cycle;
\fill[pthick] (-1.998,1.759) -- (-2.290,1.853) -- (-1.998,1.452) -- cycle;
\fill[pthin] (-2.290,1.853) -- (-1.998,1.759) -- (-1.998,1.948) -- cycle;
\fill[pthick] (-1.998,2.444) -- (-1.998,2.751) -- (-2.290,2.349) -- cycle;
\fill[pthick] (-2.179,2.196) -- (-1.998,2.444) -- (-1.998,1.948) -- cycle;
\fill[pthin] (-1.998,2.444) -- (-2.179,2.196) -- (-2.290,2.349) -- cycle;
\fill[pthin] (-2.470,2.101) -- (-2.179,2.196) -- (-2.290,2.349) -- cycle;
\fill[pthick] (-2.179,2.196) -- (-2.470,2.101) -- (-1.998,1.948) -- cycle;
\fill[pthick] (-2.762,1.204) -- (-2.470,1.299) -- (-2.942,1.452) -- cycle;
\fill[pthick] (-3.053,1.299) -- (-2.762,1.204) -- (-3.234,1.051) -- cycle;
\fill[pthin] (-2.762,1.204) -- (-3.053,1.299) -- (-2.942,1.452) -- cycle;
\fill[pthick] (-2.650,1.853) -- (-2.470,2.101) -- (-2.470,1.605) -- cycle;
\fill[pthick] (-2.650,1.547) -- (-2.650,1.853) -- (-2.942,1.452) -- cycle;
\fill[pthin] (-2.650,1.853) -- (-2.650,1.547) -- (-2.470,1.605) -- cycle;
\fill[pthin] (-2.470,1.299) -- (-2.650,1.547) -- (-2.470,1.605) -- cycle;
\fill[pthick] (-2.650,1.547) -- (-2.470,1.299) -- (-2.942,1.452) -- cycle;
\fill[pthin] (-2.470,1.299) -- (-2.290,1.547) -- (-2.470,1.605) -- cycle;
\fill[pthick] (-2.290,1.547) -- (-2.470,1.299) -- (-1.998,1.452) -- cycle;
\fill[pthick] (-2.290,1.853) -- (-2.470,2.101) -- (-2.470,1.605) -- cycle;
\fill[pthick] (-2.290,1.547) -- (-2.290,1.853) -- (-1.998,1.452) -- cycle;
\fill[pthin] (-2.290,1.853) -- (-2.290,1.547) -- (-2.470,1.605) -- cycle;
\fill[pthick] (-1.055,0.649) -- (-1.235,0.401) -- (-1.235,0.897) -- cycle;
\fill[pthick] (-1.055,0.956) -- (-1.055,0.649) -- (-0.763,1.051) -- cycle;
\fill[pthin] (-1.055,0.649) -- (-1.055,0.956) -- (-1.235,0.897) -- cycle;
\fill[pthick] (-1.707,0.744) -- (-1.998,0.649) -- (-1.527,0.496) -- cycle;
\fill[pthick] (-1.415,0.649) -- (-1.707,0.744) -- (-1.235,0.897) -- cycle;
\fill[pthin] (-1.707,0.744) -- (-1.415,0.649) -- (-1.527,0.496) -- cycle;
\fill[pthin] (-1.235,0.401) -- (-1.415,0.649) -- (-1.527,0.496) -- cycle;
\fill[pthick] (-1.415,0.649) -- (-1.235,0.401) -- (-1.235,0.897) -- cycle;
\fill[pthick] (-0.292,0.401) -- (-0.472,0.649) -- (-0.472,0.153) -- cycle;
\fill[pthick] (-0.292,0.095) -- (-0.292,0.401) -- (0.000,0.000) -- cycle;
\fill[pthin] (-0.292,0.401) -- (-0.292,0.095) -- (-0.472,0.153) -- cycle;
\fill[pthick] (-0.944,0.307) -- (-1.235,0.401) -- (-0.763,0.555) -- cycle;
\fill[pthick] (-0.652,0.401) -- (-0.944,0.307) -- (-0.472,0.153) -- cycle;
\fill[pthin] (-0.944,0.307) -- (-0.652,0.401) -- (-0.763,0.555) -- cycle;
\fill[pthin] (-0.472,0.649) -- (-0.652,0.401) -- (-0.763,0.555) -- cycle;
\fill[pthick] (-0.652,0.401) -- (-0.472,0.649) -- (-0.472,0.153) -- cycle;
\fill[pthin] (-0.472,0.649) -- (-0.763,0.744) -- (-0.763,0.555) -- cycle;
\fill[pthick] (-0.763,0.744) -- (-0.472,0.649) -- (-0.763,1.051) -- cycle;
\fill[pthick] (-1.055,0.649) -- (-1.235,0.401) -- (-0.763,0.555) -- cycle;
\fill[pthick] (-0.763,0.744) -- (-1.055,0.649) -- (-0.763,1.051) -- cycle;
\fill[pthin] (-1.055,0.649) -- (-0.763,0.744) -- (-0.763,0.555) -- cycle;
\fill[pthick] (-1.707,1.853) -- (-1.527,2.101) -- (-1.998,1.948) -- cycle;
\fill[pthick] (-1.998,1.759) -- (-1.707,1.853) -- (-1.998,1.452) -- cycle;
\fill[pthin] (-1.707,1.853) -- (-1.998,1.759) -- (-1.998,1.948) -- cycle;
\fill[pthick] (-1.998,2.444) -- (-1.998,2.751) -- (-1.707,2.349) -- cycle;
\fill[pthick] (-1.818,2.196) -- (-1.998,2.444) -- (-1.998,1.948) -- cycle;
\fill[pthin] (-1.998,2.444) -- (-1.818,2.196) -- (-1.707,2.349) -- cycle;
\fill[pthin] (-1.527,2.101) -- (-1.818,2.196) -- (-1.707,2.349) -- cycle;
\fill[pthick] (-1.818,2.196) -- (-1.527,2.101) -- (-1.998,1.948) -- cycle;
\fill[pthick] (-1.235,1.204) -- (-1.527,1.299) -- (-1.055,1.452) -- cycle;
\fill[pthick] (-0.944,1.299) -- (-1.235,1.204) -- (-0.763,1.051) -- cycle;
\fill[pthin] (-1.235,1.204) -- (-0.944,1.299) -- (-1.055,1.452) -- cycle;
\fill[pthick] (-1.346,1.853) -- (-1.527,2.101) -- (-1.527,1.605) -- cycle;
\fill[pthick] (-1.346,1.547) -- (-1.346,1.853) -- (-1.055,1.452) -- cycle;
\fill[pthin] (-1.346,1.853) -- (-1.346,1.547) -- (-1.527,1.605) -- cycle;
\fill[pthin] (-1.527,1.299) -- (-1.346,1.547) -- (-1.527,1.605) -- cycle;
\fill[pthick] (-1.346,1.547) -- (-1.527,1.299) -- (-1.055,1.452) -- cycle;
\fill[pthin] (-1.527,1.299) -- (-1.707,1.547) -- (-1.527,1.605) -- cycle;
\fill[pthick] (-1.707,1.547) -- (-1.527,1.299) -- (-1.998,1.452) -- cycle;
\fill[pthick] (-1.707,1.853) -- (-1.527,2.101) -- (-1.527,1.605) -- cycle;
\fill[pthick] (-1.707,1.547) -- (-1.707,1.853) -- (-1.998,1.452) -- cycle;
\fill[pthin] (-1.707,1.853) -- (-1.707,1.547) -- (-1.527,1.605) -- cycle;
\fill[pthin] (-1.527,1.299) -- (-1.818,1.204) -- (-1.707,1.051) -- cycle;
\fill[pthick] (-1.818,1.204) -- (-1.527,1.299) -- (-1.998,1.452) -- cycle;
\fill[pthick] (-1.998,0.956) -- (-1.998,0.649) -- (-1.707,1.051) -- cycle;
\fill[pthick] (-1.818,1.204) -- (-1.998,0.956) -- (-1.998,1.452) -- cycle;
\fill[pthin] (-1.998,0.956) -- (-1.818,1.204) -- (-1.707,1.051) -- cycle;
\fill[pthick] (-1.235,1.204) -- (-1.527,1.299) -- (-1.235,0.897) -- cycle;
\fill[pthick] (-1.055,0.956) -- (-1.235,1.204) -- (-0.763,1.051) -- cycle;
\fill[pthin] (-1.235,1.204) -- (-1.055,0.956) -- (-1.235,0.897) -- cycle;
\fill[pthick] (-1.707,0.744) -- (-1.998,0.649) -- (-1.707,1.051) -- cycle;
\fill[pthick] (-1.527,0.992) -- (-1.707,0.744) -- (-1.235,0.897) -- cycle;
\fill[pthin] (-1.707,0.744) -- (-1.527,0.992) -- (-1.707,1.051) -- cycle;
\fill[pthin] (-1.527,1.299) -- (-1.527,0.992) -- (-1.707,1.051) -- cycle;
\fill[pthick] (-1.527,0.992) -- (-1.527,1.299) -- (-1.235,0.897) -- cycle;
\fill[pthin] (-2.762,0.401) -- (-2.470,0.307) -- (-2.470,0.496) -- cycle;
\fill[pthick] (-2.470,0.307) -- (-2.762,0.401) -- (-2.470,0.000) -- cycle;
\fill[pthick] (-2.179,0.401) -- (-1.998,0.649) -- (-2.470,0.496) -- cycle;
\fill[pthick] (-2.470,0.307) -- (-2.179,0.401) -- (-2.470,0.000) -- cycle;
\fill[pthin] (-2.179,0.401) -- (-2.470,0.307) -- (-2.470,0.496) -- cycle;
\fill[pthick] (-2.942,0.649) -- (-2.762,0.401) -- (-2.762,0.897) -- cycle;
\fill[pthick] (-2.942,0.956) -- (-2.942,0.649) -- (-3.234,1.051) -- cycle;
\fill[pthin] (-2.942,0.649) -- (-2.942,0.956) -- (-2.762,0.897) -- cycle;
\fill[pthick] (-2.290,0.744) -- (-1.998,0.649) -- (-2.470,0.496) -- cycle;
\fill[pthick] (-2.582,0.649) -- (-2.290,0.744) -- (-2.762,0.897) -- cycle;
\fill[pthin] (-2.290,0.744) -- (-2.582,0.649) -- (-2.470,0.496) -- cycle;
\fill[pthin] (-2.762,0.401) -- (-2.582,0.649) -- (-2.470,0.496) -- cycle;
\fill[pthick] (-2.582,0.649) -- (-2.762,0.401) -- (-2.762,0.897) -- cycle;
\fill[pthick] (-2.942,-0.153) -- (-3.234,-0.248) -- (-2.762,-0.401) -- cycle;
\fill[pthick] (-2.650,-0.248) -- (-2.942,-0.153) -- (-2.470,0.000) -- cycle;
\fill[pthin] (-2.942,-0.153) -- (-2.650,-0.248) -- (-2.762,-0.401) -- cycle;
\fill[pthick] (-3.053,-0.803) -- (-3.234,-1.051) -- (-3.234,-0.555) -- cycle;
\fill[pthick] (-3.053,-0.496) -- (-3.053,-0.803) -- (-2.762,-0.401) -- cycle;
\fill[pthin] (-3.053,-0.803) -- (-3.053,-0.496) -- (-3.234,-0.555) -- cycle;
\fill[pthin] (-3.234,-0.248) -- (-3.053,-0.496) -- (-3.234,-0.555) -- cycle;
\fill[pthick] (-3.053,-0.496) -- (-3.234,-0.248) -- (-2.762,-0.401) -- cycle;
\fill[pthick] (-2.942,0.649) -- (-2.762,0.401) -- (-3.234,0.555) -- cycle;
\fill[pthick] (-3.234,0.744) -- (-2.942,0.649) -- (-3.234,1.051) -- cycle;
\fill[pthin] (-2.942,0.649) -- (-3.234,0.744) -- (-3.234,0.555) -- cycle;
\fill[pthick] (-3.234,0.059) -- (-3.234,-0.248) -- (-2.942,0.153) -- cycle;
\fill[pthick] (-3.053,0.307) -- (-3.234,0.059) -- (-3.234,0.555) -- cycle;
\fill[pthin] (-3.234,0.059) -- (-3.053,0.307) -- (-2.942,0.153) -- cycle;
\fill[pthin] (-2.762,0.401) -- (-3.053,0.307) -- (-2.942,0.153) -- cycle;
\fill[pthick] (-3.053,0.307) -- (-2.762,0.401) -- (-3.234,0.555) -- cycle;
\fill[pthin] (-2.762,0.401) -- (-2.762,0.095) -- (-2.942,0.153) -- cycle;
\fill[pthick] (-2.762,0.095) -- (-2.762,0.401) -- (-2.470,0.000) -- cycle;
\fill[pthick] (-2.942,-0.153) -- (-3.234,-0.248) -- (-2.942,0.153) -- cycle;
\fill[pthick] (-2.762,0.095) -- (-2.942,-0.153) -- (-2.470,0.000) -- cycle;
\fill[pthin] (-2.942,-0.153) -- (-2.762,0.095) -- (-2.942,0.153) -- cycle;
\fill[pthick] (-1.235,0.095) -- (-1.235,0.401) -- (-1.527,0.000) -- cycle;
\fill[pthick] (-1.415,-0.153) -- (-1.235,0.095) -- (-1.235,-0.401) -- cycle;
\fill[pthin] (-1.235,0.095) -- (-1.415,-0.153) -- (-1.527,0.000) -- cycle;
\fill[pthick] (-1.818,0.401) -- (-1.998,0.649) -- (-1.527,0.496) -- cycle;
\fill[pthick] (-1.527,0.307) -- (-1.818,0.401) -- (-1.527,0.000) -- cycle;
\fill[pthin] (-1.818,0.401) -- (-1.527,0.307) -- (-1.527,0.496) -- cycle;
\fill[pthin] (-1.235,0.401) -- (-1.527,0.307) -- (-1.527,0.496) -- cycle;
\fill[pthick] (-1.527,0.307) -- (-1.235,0.401) -- (-1.527,0.000) -- cycle;
\fill[pthick] (-0.472,-0.153) -- (-0.763,-0.248) -- (-0.472,0.153) -- cycle;
\fill[pthick] (-0.292,0.095) -- (-0.472,-0.153) -- (0.000,0.000) -- cycle;
\fill[pthin] (-0.472,-0.153) -- (-0.292,0.095) -- (-0.472,0.153) -- cycle;
\fill[pthick] (-0.944,0.307) -- (-1.235,0.401) -- (-0.944,0.000) -- cycle;
\fill[pthick] (-0.763,0.059) -- (-0.944,0.307) -- (-0.472,0.153) -- cycle;
\fill[pthin] (-0.944,0.307) -- (-0.763,0.059) -- (-0.944,0.000) -- cycle;
\fill[pthin] (-0.763,-0.248) -- (-0.763,0.059) -- (-0.944,0.000) -- cycle;
\fill[pthick] (-0.763,0.059) -- (-0.763,-0.248) -- (-0.472,0.153) -- cycle;
\fill[pthin] (-0.763,-0.248) -- (-1.055,-0.153) -- (-0.944,0.000) -- cycle;
\fill[pthick] (-1.055,-0.153) -- (-0.763,-0.248) -- (-1.235,-0.401) -- cycle;
\fill[pthick] (-1.235,0.095) -- (-1.235,0.401) -- (-0.944,0.000) -- cycle;
\fill[pthick] (-1.055,-0.153) -- (-1.235,0.095) -- (-1.235,-0.401) -- cycle;
\fill[pthin] (-1.235,0.095) -- (-1.055,-0.153) -- (-0.944,0.000) -- cycle;
\fill[pthick] (-2.470,-0.496) -- (-2.470,-0.803) -- (-2.762,-0.401) -- cycle;
\fill[pthick] (-2.650,-0.248) -- (-2.470,-0.496) -- (-2.470,0.000) -- cycle;
\fill[pthin] (-2.470,-0.496) -- (-2.650,-0.248) -- (-2.762,-0.401) -- cycle;
\fill[pthick] (-3.053,-0.803) -- (-3.234,-1.051) -- (-2.762,-0.897) -- cycle;
\fill[pthick] (-2.762,-0.708) -- (-3.053,-0.803) -- (-2.762,-0.401) -- cycle;
\fill[pthin] (-3.053,-0.803) -- (-2.762,-0.708) -- (-2.762,-0.897) -- cycle;
\fill[pthin] (-2.470,-0.803) -- (-2.762,-0.708) -- (-2.762,-0.897) -- cycle;
\fill[pthick] (-2.762,-0.708) -- (-2.470,-0.803) -- (-2.762,-0.401) -- cycle;
\fill[pthick] (-1.707,-0.248) -- (-1.998,-0.153) -- (-1.707,-0.555) -- cycle;
\fill[pthick] (-1.527,-0.496) -- (-1.707,-0.248) -- (-1.235,-0.401) -- cycle;
\fill[pthin] (-1.707,-0.248) -- (-1.527,-0.496) -- (-1.707,-0.555) -- cycle;
\fill[pthick] (-2.179,-0.708) -- (-2.470,-0.803) -- (-2.179,-0.401) -- cycle;
\fill[pthick] (-1.998,-0.460) -- (-2.179,-0.708) -- (-1.707,-0.555) -- cycle;
\fill[pthin] (-2.179,-0.708) -- (-1.998,-0.460) -- (-2.179,-0.401) -- cycle;
\fill[pthin] (-1.998,-0.153) -- (-1.998,-0.460) -- (-2.179,-0.401) -- cycle;
\fill[pthick] (-1.998,-0.460) -- (-1.998,-0.153) -- (-1.707,-0.555) -- cycle;
\fill[pthin] (-1.998,-0.153) -- (-2.290,-0.248) -- (-2.179,-0.401) -- cycle;
\fill[pthick] (-2.290,-0.248) -- (-1.998,-0.153) -- (-2.470,0.000) -- cycle;
\fill[pthick] (-2.470,-0.496) -- (-2.470,-0.803) -- (-2.179,-0.401) -- cycle;
\fill[pthick] (-2.290,-0.248) -- (-2.470,-0.496) -- (-2.470,0.000) -- cycle;
\fill[pthin] (-2.470,-0.496) -- (-2.290,-0.248) -- (-2.179,-0.401) -- cycle;
\fill[pthin] (-1.998,-0.153) -- (-2.179,0.095) -- (-1.998,0.153) -- cycle;
\fill[pthick] (-2.179,0.095) -- (-1.998,-0.153) -- (-2.470,0.000) -- cycle;
\fill[pthick] (-2.179,0.401) -- (-1.998,0.649) -- (-1.998,0.153) -- cycle;
\fill[pthick] (-2.179,0.095) -- (-2.179,0.401) -- (-2.470,0.000) -- cycle;
\fill[pthin] (-2.179,0.401) -- (-2.179,0.095) -- (-1.998,0.153) -- cycle;
\fill[pthick] (-1.707,-0.248) -- (-1.998,-0.153) -- (-1.527,0.000) -- cycle;
\fill[pthick] (-1.415,-0.153) -- (-1.707,-0.248) -- (-1.235,-0.401) -- cycle;
\fill[pthin] (-1.707,-0.248) -- (-1.415,-0.153) -- (-1.527,0.000) -- cycle;
\fill[pthick] (-1.818,0.401) -- (-1.998,0.649) -- (-1.998,0.153) -- cycle;
\fill[pthick] (-1.818,0.095) -- (-1.818,0.401) -- (-1.527,0.000) -- cycle;
\fill[pthin] (-1.818,0.401) -- (-1.818,0.095) -- (-1.998,0.153) -- cycle;
\fill[pthin] (-1.998,-0.153) -- (-1.818,0.095) -- (-1.998,0.153) -- cycle;
\fill[pthick] (-1.818,0.095) -- (-1.998,-0.153) -- (-1.527,0.000) -- cycle;
\fill[pthin] (-1.998,-1.948) -- (-1.818,-1.700) -- (-1.707,-1.853) -- cycle;
\fill[pthick] (-1.818,-1.700) -- (-1.998,-1.948) -- (-1.998,-1.452) -- cycle;
\fill[pthick] (-1.527,-1.605) -- (-1.235,-1.700) -- (-1.707,-1.853) -- cycle;
\fill[pthick] (-1.818,-1.700) -- (-1.527,-1.605) -- (-1.998,-1.452) -- cycle;
\fill[pthin] (-1.527,-1.605) -- (-1.818,-1.700) -- (-1.707,-1.853) -- cycle;
\fill[pthick] (-1.998,-2.255) -- (-1.998,-1.948) -- (-1.707,-2.349) -- cycle;
\fill[pthick] (-1.818,-2.503) -- (-1.998,-2.255) -- (-1.998,-2.751) -- cycle;
\fill[pthin] (-1.998,-2.255) -- (-1.818,-2.503) -- (-1.707,-2.349) -- cycle;
\fill[pthick] (-1.415,-1.948) -- (-1.235,-1.700) -- (-1.707,-1.853) -- cycle;
\fill[pthick] (-1.707,-2.043) -- (-1.415,-1.948) -- (-1.707,-2.349) -- cycle;
\fill[pthin] (-1.415,-1.948) -- (-1.707,-2.043) -- (-1.707,-1.853) -- cycle;
\fill[pthin] (-1.998,-1.948) -- (-1.707,-2.043) -- (-1.707,-1.853) -- cycle;
\fill[pthick] (-1.707,-2.043) -- (-1.998,-1.948) -- (-1.707,-2.349) -- cycle;
\fill[pthick] (-2.470,-1.605) -- (-2.762,-1.700) -- (-2.470,-1.299) -- cycle;
\fill[pthick] (-2.290,-1.357) -- (-2.470,-1.605) -- (-1.998,-1.452) -- cycle;
\fill[pthin] (-2.470,-1.605) -- (-2.290,-1.357) -- (-2.470,-1.299) -- cycle;
\fill[pthick] (-2.942,-1.145) -- (-3.234,-1.051) -- (-2.942,-1.452) -- cycle;
\fill[pthick] (-2.762,-1.393) -- (-2.942,-1.145) -- (-2.470,-1.299) -- cycle;
\fill[pthin] (-2.942,-1.145) -- (-2.762,-1.393) -- (-2.942,-1.452) -- cycle;
\fill[pthin] (-2.762,-1.700) -- (-2.762,-1.393) -- (-2.942,-1.452) -- cycle;
\fill[pthick] (-2.762,-1.393) -- (-2.762,-1.700) -- (-2.470,-1.299) -- cycle;
\fill[pthick] (-1.998,-2.255) -- (-1.998,-1.948) -- (-2.290,-2.349) -- cycle;
\fill[pthick] (-2.179,-2.503) -- (-1.998,-2.255) -- (-1.998,-2.751) -- cycle;
\fill[pthin] (-1.998,-2.255) -- (-2.179,-2.503) -- (-2.290,-2.349) -- cycle;
\fill[pthick] (-2.582,-1.948) -- (-2.762,-1.700) -- (-2.290,-1.853) -- cycle;
\fill[pthick] (-2.290,-2.043) -- (-2.582,-1.948) -- (-2.290,-2.349) -- cycle;
\fill[pthin] (-2.582,-1.948) -- (-2.290,-2.043) -- (-2.290,-1.853) -- cycle;
\fill[pthin] (-1.998,-1.948) -- (-2.290,-2.043) -- (-2.290,-1.853) -- cycle;
\fill[pthick] (-2.290,-2.043) -- (-1.998,-1.948) -- (-2.290,-2.349) -- cycle;
\fill[pthin] (-1.998,-1.948) -- (-2.179,-1.700) -- (-2.290,-1.853) -- cycle;
\fill[pthick] (-2.179,-1.700) -- (-1.998,-1.948) -- (-1.998,-1.452) -- cycle;
\fill[pthick] (-2.470,-1.605) -- (-2.762,-1.700) -- (-2.290,-1.853) -- cycle;
\fill[pthick] (-2.179,-1.700) -- (-2.470,-1.605) -- (-1.998,-1.452) -- cycle;
\fill[pthin] (-2.470,-1.605) -- (-2.179,-1.700) -- (-2.290,-1.853) -- cycle;
\fill[pthick] (-0.944,-0.803) -- (-0.763,-1.051) -- (-1.235,-0.897) -- cycle;
\fill[pthick] (-1.235,-0.708) -- (-0.944,-0.803) -- (-1.235,-0.401) -- cycle;
\fill[pthin] (-0.944,-0.803) -- (-1.235,-0.708) -- (-1.235,-0.897) -- cycle;
\fill[pthick] (-1.235,-1.393) -- (-1.235,-1.700) -- (-0.944,-1.299) -- cycle;
\fill[pthick] (-1.055,-1.145) -- (-1.235,-1.393) -- (-1.235,-0.897) -- cycle;
\fill[pthin] (-1.235,-1.393) -- (-1.055,-1.145) -- (-0.944,-1.299) -- cycle;
\fill[pthin] (-0.763,-1.051) -- (-1.055,-1.145) -- (-0.944,-1.299) -- cycle;
\fill[pthick] (-1.055,-1.145) -- (-0.763,-1.051) -- (-1.235,-0.897) -- cycle;
\fill[pthick] (-0.472,-0.153) -- (-0.763,-0.248) -- (-0.292,-0.401) -- cycle;
\fill[pthick] (-0.180,-0.248) -- (-0.472,-0.153) -- (0.000,0.000) -- cycle;
\fill[pthin] (-0.472,-0.153) -- (-0.180,-0.248) -- (-0.292,-0.401) -- cycle;
\fill[pthick] (-0.583,-0.803) -- (-0.763,-1.051) -- (-0.763,-0.555) -- cycle;
\fill[pthick] (-0.583,-0.496) -- (-0.583,-0.803) -- (-0.292,-0.401) -- cycle;
\fill[pthin] (-0.583,-0.803) -- (-0.583,-0.496) -- (-0.763,-0.555) -- cycle;
\fill[pthin] (-0.763,-0.248) -- (-0.583,-0.496) -- (-0.763,-0.555) -- cycle;
\fill[pthick] (-0.583,-0.496) -- (-0.763,-0.248) -- (-0.292,-0.401) -- cycle;
\fill[pthin] (-0.763,-0.248) -- (-0.944,-0.496) -- (-0.763,-0.555) -- cycle;
\fill[pthick] (-0.944,-0.496) -- (-0.763,-0.248) -- (-1.235,-0.401) -- cycle;
\fill[pthick] (-0.944,-0.803) -- (-0.763,-1.051) -- (-0.763,-0.555) -- cycle;
\fill[pthick] (-0.944,-0.496) -- (-0.944,-0.803) -- (-1.235,-0.401) -- cycle;
\fill[pthin] (-0.944,-0.803) -- (-0.944,-0.496) -- (-0.763,-0.555) -- cycle;
\fill[pthick] (-2.290,-1.051) -- (-2.470,-0.803) -- (-2.470,-1.299) -- cycle;
\fill[pthick] (-2.290,-1.357) -- (-2.290,-1.051) -- (-1.998,-1.452) -- cycle;
\fill[pthin] (-2.290,-1.051) -- (-2.290,-1.357) -- (-2.470,-1.299) -- cycle;
\fill[pthick] (-2.942,-1.145) -- (-3.234,-1.051) -- (-2.762,-0.897) -- cycle;
\fill[pthick] (-2.650,-1.051) -- (-2.942,-1.145) -- (-2.470,-1.299) -- cycle;
\fill[pthin] (-2.942,-1.145) -- (-2.650,-1.051) -- (-2.762,-0.897) -- cycle;
\fill[pthin] (-2.470,-0.803) -- (-2.650,-1.051) -- (-2.762,-0.897) -- cycle;
\fill[pthick] (-2.650,-1.051) -- (-2.470,-0.803) -- (-2.470,-1.299) -- cycle;
\fill[pthick] (-1.527,-0.803) -- (-1.707,-1.051) -- (-1.707,-0.555) -- cycle;
\fill[pthick] (-1.527,-0.496) -- (-1.527,-0.803) -- (-1.235,-0.401) -- cycle;
\fill[pthin] (-1.527,-0.803) -- (-1.527,-0.496) -- (-1.707,-0.555) -- cycle;
\fill[pthick] (-2.179,-0.708) -- (-2.470,-0.803) -- (-1.998,-0.956) -- cycle;
\fill[pthick] (-1.887,-0.803) -- (-2.179,-0.708) -- (-1.707,-0.555) -- cycle;
\fill[pthin] (-2.179,-0.708) -- (-1.887,-0.803) -- (-1.998,-0.956) -- cycle;
\fill[pthin] (-1.707,-1.051) -- (-1.887,-0.803) -- (-1.998,-0.956) -- cycle;
\fill[pthick] (-1.887,-0.803) -- (-1.707,-1.051) -- (-1.707,-0.555) -- cycle;
\fill[pthin] (-1.707,-1.051) -- (-1.998,-1.145) -- (-1.998,-0.956) -- cycle;
\fill[pthick] (-1.998,-1.145) -- (-1.707,-1.051) -- (-1.998,-1.452) -- cycle;
\fill[pthick] (-2.290,-1.051) -- (-2.470,-0.803) -- (-1.998,-0.956) -- cycle;
\fill[pthick] (-1.998,-1.145) -- (-2.290,-1.051) -- (-1.998,-1.452) -- cycle;
\fill[pthin] (-2.290,-1.051) -- (-1.998,-1.145) -- (-1.998,-0.956) -- cycle;
\fill[pthin] (-1.707,-1.051) -- (-1.707,-1.357) -- (-1.527,-1.299) -- cycle;
\fill[pthick] (-1.707,-1.357) -- (-1.707,-1.051) -- (-1.998,-1.452) -- cycle;
\fill[pthick] (-1.527,-1.605) -- (-1.235,-1.700) -- (-1.527,-1.299) -- cycle;
\fill[pthick] (-1.707,-1.357) -- (-1.527,-1.605) -- (-1.998,-1.452) -- cycle;
\fill[pthin] (-1.527,-1.605) -- (-1.707,-1.357) -- (-1.527,-1.299) -- cycle;
\fill[pthick] (-1.527,-0.803) -- (-1.707,-1.051) -- (-1.235,-0.897) -- cycle;
\fill[pthick] (-1.235,-0.708) -- (-1.527,-0.803) -- (-1.235,-0.401) -- cycle;
\fill[pthin] (-1.527,-0.803) -- (-1.235,-0.708) -- (-1.235,-0.897) -- cycle;
\fill[pthick] (-1.235,-1.393) -- (-1.235,-1.700) -- (-1.527,-1.299) -- cycle;
\fill[pthick] (-1.415,-1.145) -- (-1.235,-1.393) -- (-1.235,-0.897) -- cycle;
\fill[pthin] (-1.235,-1.393) -- (-1.415,-1.145) -- (-1.527,-1.299) -- cycle;
\fill[pthin] (-1.707,-1.051) -- (-1.415,-1.145) -- (-1.527,-1.299) -- cycle;
\fill[pthick] (-1.415,-1.145) -- (-1.707,-1.051) -- (-1.235,-0.897) -- cycle;
\fill[pthin] (-1.235,-2.503) -- (-1.055,-2.255) -- (-1.235,-2.196) -- cycle;
\fill[pthick] (-1.055,-2.255) -- (-1.235,-2.503) -- (-0.763,-2.349) -- cycle;
\fill[pthick] (-1.055,-1.948) -- (-1.235,-1.700) -- (-1.235,-2.196) -- cycle;
\fill[pthick] (-1.055,-2.255) -- (-1.055,-1.948) -- (-0.763,-2.349) -- cycle;
\fill[pthin] (-1.055,-1.948) -- (-1.055,-2.255) -- (-1.235,-2.196) -- cycle;
\fill[pthick] (-1.527,-2.597) -- (-1.235,-2.503) -- (-1.707,-2.349) -- cycle;
\fill[pthick] (-1.818,-2.503) -- (-1.527,-2.597) -- (-1.998,-2.751) -- cycle;
\fill[pthin] (-1.527,-2.597) -- (-1.818,-2.503) -- (-1.707,-2.349) -- cycle;
\fill[pthick] (-1.415,-1.948) -- (-1.235,-1.700) -- (-1.235,-2.196) -- cycle;
\fill[pthick] (-1.415,-2.255) -- (-1.415,-1.948) -- (-1.707,-2.349) -- cycle;
\fill[pthin] (-1.415,-1.948) -- (-1.415,-2.255) -- (-1.235,-2.196) -- cycle;
\fill[pthin] (-1.235,-2.503) -- (-1.415,-2.255) -- (-1.235,-2.196) -- cycle;
\fill[pthick] (-1.415,-2.255) -- (-1.235,-2.503) -- (-1.707,-2.349) -- cycle;
\fill[pthick] (-0.763,-2.845) -- (-0.763,-3.152) -- (-0.472,-2.751) -- cycle;
\fill[pthick] (-0.583,-2.597) -- (-0.763,-2.845) -- (-0.763,-2.349) -- cycle;
\fill[pthin] (-0.763,-2.845) -- (-0.583,-2.597) -- (-0.472,-2.751) -- cycle;
\fill[pthick] (-0.180,-3.152) -- (-0.000,-3.400) -- (-0.472,-3.247) -- cycle;
\fill[pthick] (-0.472,-3.057) -- (-0.180,-3.152) -- (-0.472,-2.751) -- cycle;
\fill[pthin] (-0.180,-3.152) -- (-0.472,-3.057) -- (-0.472,-3.247) -- cycle;
\fill[pthin] (-0.763,-3.152) -- (-0.472,-3.057) -- (-0.472,-3.247) -- cycle;
\fill[pthick] (-0.472,-3.057) -- (-0.763,-3.152) -- (-0.472,-2.751) -- cycle;
\fill[pthick] (-1.527,-2.597) -- (-1.235,-2.503) -- (-1.527,-2.904) -- cycle;
\fill[pthick] (-1.707,-2.845) -- (-1.527,-2.597) -- (-1.998,-2.751) -- cycle;
\fill[pthin] (-1.527,-2.597) -- (-1.707,-2.845) -- (-1.527,-2.904) -- cycle;
\fill[pthick] (-1.055,-3.057) -- (-0.763,-3.152) -- (-1.055,-2.751) -- cycle;
\fill[pthick] (-1.235,-2.809) -- (-1.055,-3.057) -- (-1.527,-2.904) -- cycle;
\fill[pthin] (-1.055,-3.057) -- (-1.235,-2.809) -- (-1.055,-2.751) -- cycle;
\fill[pthin] (-1.235,-2.503) -- (-1.235,-2.809) -- (-1.055,-2.751) -- cycle;
\fill[pthick] (-1.235,-2.809) -- (-1.235,-2.503) -- (-1.527,-2.904) -- cycle;
\fill[pthin] (-1.235,-2.503) -- (-0.944,-2.597) -- (-1.055,-2.751) -- cycle;
\fill[pthick] (-0.944,-2.597) -- (-1.235,-2.503) -- (-0.763,-2.349) -- cycle;
\fill[pthick] (-0.763,-2.845) -- (-0.763,-3.152) -- (-1.055,-2.751) -- cycle;
\fill[pthick] (-0.944,-2.597) -- (-0.763,-2.845) -- (-0.763,-2.349) -- cycle;
\fill[pthin] (-0.763,-2.845) -- (-0.944,-2.597) -- (-1.055,-2.751) -- cycle;
\fill[pthick] (-0.472,-1.145) -- (-0.763,-1.051) -- (-0.472,-1.452) -- cycle;
\fill[pthick] (-0.292,-1.393) -- (-0.472,-1.145) -- (-0.000,-1.299) -- cycle;
\fill[pthin] (-0.472,-1.145) -- (-0.292,-1.393) -- (-0.472,-1.452) -- cycle;
\fill[pthick] (-0.944,-1.605) -- (-1.235,-1.700) -- (-0.944,-1.299) -- cycle;
\fill[pthick] (-0.763,-1.357) -- (-0.944,-1.605) -- (-0.472,-1.452) -- cycle;
\fill[pthin] (-0.944,-1.605) -- (-0.763,-1.357) -- (-0.944,-1.299) -- cycle;
\fill[pthin] (-0.763,-1.051) -- (-0.763,-1.357) -- (-0.944,-1.299) -- cycle;
\fill[pthick] (-0.763,-1.357) -- (-0.763,-1.051) -- (-0.472,-1.452) -- cycle;
\fill[pthick] (-0.000,-0.496) -- (-0.000,-0.803) -- (-0.292,-0.401) -- cycle;
\fill[pthick] (-0.180,-0.248) -- (-0.000,-0.496) -- (0.000,0.000) -- cycle;
\fill[pthin] (-0.000,-0.496) -- (-0.180,-0.248) -- (-0.292,-0.401) -- cycle;
\fill[pthick] (-0.583,-0.803) -- (-0.763,-1.051) -- (-0.292,-0.897) -- cycle;
\fill[pthick] (-0.292,-0.708) -- (-0.583,-0.803) -- (-0.292,-0.401) -- cycle;
\fill[pthin] (-0.583,-0.803) -- (-0.292,-0.708) -- (-0.292,-0.897) -- cycle;
\fill[pthin] (-0.000,-0.803) -- (-0.292,-0.708) -- (-0.292,-0.897) -- cycle;
\fill[pthick] (-0.292,-0.708) -- (-0.000,-0.803) -- (-0.292,-0.401) -- cycle;
\fill[pthin] (-0.000,-0.803) -- (-0.180,-1.051) -- (-0.292,-0.897) -- cycle;
\fill[pthick] (-0.180,-1.051) -- (-0.000,-0.803) -- (-0.000,-1.299) -- cycle;
\fill[pthick] (-0.472,-1.145) -- (-0.763,-1.051) -- (-0.292,-0.897) -- cycle;
\fill[pthick] (-0.180,-1.051) -- (-0.472,-1.145) -- (-0.000,-1.299) -- cycle;
\fill[pthin] (-0.472,-1.145) -- (-0.180,-1.051) -- (-0.292,-0.897) -- cycle;
\fill[pthick] (-0.292,-2.503) -- (-0.000,-2.597) -- (-0.472,-2.751) -- cycle;
\fill[pthick] (-0.583,-2.597) -- (-0.292,-2.503) -- (-0.763,-2.349) -- cycle;
\fill[pthin] (-0.292,-2.503) -- (-0.583,-2.597) -- (-0.472,-2.751) -- cycle;
\fill[pthick] (-0.180,-3.152) -- (-0.000,-3.400) -- (-0.000,-2.904) -- cycle;
\fill[pthick] (-0.180,-2.845) -- (-0.180,-3.152) -- (-0.472,-2.751) -- cycle;
\fill[pthin] (-0.180,-3.152) -- (-0.180,-2.845) -- (-0.000,-2.904) -- cycle;
\fill[pthin] (-0.000,-2.597) -- (-0.180,-2.845) -- (-0.000,-2.904) -- cycle;
\fill[pthick] (-0.180,-2.845) -- (-0.000,-2.597) -- (-0.472,-2.751) -- cycle;
\fill[pthick] (-0.292,-1.700) -- (-0.472,-1.948) -- (-0.000,-1.795) -- cycle;
\fill[pthick] (-0.000,-1.605) -- (-0.292,-1.700) -- (-0.000,-1.299) -- cycle;
\fill[pthin] (-0.292,-1.700) -- (-0.000,-1.605) -- (-0.000,-1.795) -- cycle;
\fill[pthick] (-0.000,-2.291) -- (-0.000,-2.597) -- (-0.292,-2.196) -- cycle;
\fill[pthick] (-0.180,-2.043) -- (-0.000,-2.291) -- (-0.000,-1.795) -- cycle;
\fill[pthin] (-0.000,-2.291) -- (-0.180,-2.043) -- (-0.292,-2.196) -- cycle;
\fill[pthin] (-0.472,-1.948) -- (-0.180,-2.043) -- (-0.292,-2.196) -- cycle;
\fill[pthick] (-0.180,-2.043) -- (-0.472,-1.948) -- (-0.000,-1.795) -- cycle;
\fill[pthin] (-0.472,-1.948) -- (-0.472,-2.255) -- (-0.292,-2.196) -- cycle;
\fill[pthick] (-0.472,-2.255) -- (-0.472,-1.948) -- (-0.763,-2.349) -- cycle;
\fill[pthick] (-0.292,-2.503) -- (-0.000,-2.597) -- (-0.292,-2.196) -- cycle;
\fill[pthick] (-0.472,-2.255) -- (-0.292,-2.503) -- (-0.763,-2.349) -- cycle;
\fill[pthin] (-0.292,-2.503) -- (-0.472,-2.255) -- (-0.292,-2.196) -- cycle;
\fill[pthin] (-0.472,-1.948) -- (-0.763,-2.043) -- (-0.763,-1.853) -- cycle;
\fill[pthick] (-0.763,-2.043) -- (-0.472,-1.948) -- (-0.763,-2.349) -- cycle;
\fill[pthick] (-1.055,-1.948) -- (-1.235,-1.700) -- (-0.763,-1.853) -- cycle;
\fill[pthick] (-0.763,-2.043) -- (-1.055,-1.948) -- (-0.763,-2.349) -- cycle;
\fill[pthin] (-1.055,-1.948) -- (-0.763,-2.043) -- (-0.763,-1.853) -- cycle;
\fill[pthick] (-0.292,-1.700) -- (-0.472,-1.948) -- (-0.472,-1.452) -- cycle;
\fill[pthick] (-0.292,-1.393) -- (-0.292,-1.700) -- (-0.000,-1.299) -- cycle;
\fill[pthin] (-0.292,-1.700) -- (-0.292,-1.393) -- (-0.472,-1.452) -- cycle;
\fill[pthick] (-0.944,-1.605) -- (-1.235,-1.700) -- (-0.763,-1.853) -- cycle;
\fill[pthick] (-0.652,-1.700) -- (-0.944,-1.605) -- (-0.472,-1.452) -- cycle;
\fill[pthin] (-0.944,-1.605) -- (-0.652,-1.700) -- (-0.763,-1.853) -- cycle;
\fill[pthin] (-0.472,-1.948) -- (-0.652,-1.700) -- (-0.763,-1.853) -- cycle;
\fill[pthick] (-0.652,-1.700) -- (-0.472,-1.948) -- (-0.472,-1.452) -- cycle;
\fill[pthin] (1.235,-2.503) -- (1.055,-2.255) -- (1.235,-2.196) -- cycle;
\fill[pthick] (1.055,-2.255) -- (1.235,-2.503) -- (0.763,-2.349) -- cycle;
\fill[pthick] (1.055,-1.948) -- (1.235,-1.700) -- (1.235,-2.196) -- cycle;
\fill[pthick] (1.055,-2.255) -- (1.055,-1.948) -- (0.763,-2.349) -- cycle;
\fill[pthin] (1.055,-1.948) -- (1.055,-2.255) -- (1.235,-2.196) -- cycle;
\fill[pthick] (1.527,-2.597) -- (1.235,-2.503) -- (1.707,-2.349) -- cycle;
\fill[pthick] (1.818,-2.503) -- (1.527,-2.597) -- (1.998,-2.751) -- cycle;
\fill[pthin] (1.527,-2.597) -- (1.818,-2.503) -- (1.707,-2.349) -- cycle;
\fill[pthick] (1.415,-1.948) -- (1.235,-1.700) -- (1.235,-2.196) -- cycle;
\fill[pthick] (1.415,-2.255) -- (1.415,-1.948) -- (1.707,-2.349) -- cycle;
\fill[pthin] (1.415,-1.948) -- (1.415,-2.255) -- (1.235,-2.196) -- cycle;
\fill[pthin] (1.235,-2.503) -- (1.415,-2.255) -- (1.235,-2.196) -- cycle;
\fill[pthick] (1.415,-2.255) -- (1.235,-2.503) -- (1.707,-2.349) -- cycle;
\fill[pthick] (0.763,-2.845) -- (0.763,-3.152) -- (0.472,-2.751) -- cycle;
\fill[pthick] (0.583,-2.597) -- (0.763,-2.845) -- (0.763,-2.349) -- cycle;
\fill[pthin] (0.763,-2.845) -- (0.583,-2.597) -- (0.472,-2.751) -- cycle;
\fill[pthick] (0.180,-3.152) -- (-0.000,-3.400) -- (0.472,-3.247) -- cycle;
\fill[pthick] (0.472,-3.057) -- (0.180,-3.152) -- (0.472,-2.751) -- cycle;
\fill[pthin] (0.180,-3.152) -- (0.472,-3.057) -- (0.472,-3.247) -- cycle;
\fill[pthin] (0.763,-3.152) -- (0.472,-3.057) -- (0.472,-3.247) -- cycle;
\fill[pthick] (0.472,-3.057) -- (0.763,-3.152) -- (0.472,-2.751) -- cycle;
\fill[pthick] (1.527,-2.597) -- (1.235,-2.503) -- (1.527,-2.904) -- cycle;
\fill[pthick] (1.707,-2.845) -- (1.527,-2.597) -- (1.998,-2.751) -- cycle;
\fill[pthin] (1.527,-2.597) -- (1.707,-2.845) -- (1.527,-2.904) -- cycle;
\fill[pthick] (1.055,-3.057) -- (0.763,-3.152) -- (1.055,-2.751) -- cycle;
\fill[pthick] (1.235,-2.809) -- (1.055,-3.057) -- (1.527,-2.904) -- cycle;
\fill[pthin] (1.055,-3.057) -- (1.235,-2.809) -- (1.055,-2.751) -- cycle;
\fill[pthin] (1.235,-2.503) -- (1.235,-2.809) -- (1.055,-2.751) -- cycle;
\fill[pthick] (1.235,-2.809) -- (1.235,-2.503) -- (1.527,-2.904) -- cycle;
\fill[pthin] (1.235,-2.503) -- (0.944,-2.597) -- (1.055,-2.751) -- cycle;
\fill[pthick] (0.944,-2.597) -- (1.235,-2.503) -- (0.763,-2.349) -- cycle;
\fill[pthick] (0.763,-2.845) -- (0.763,-3.152) -- (1.055,-2.751) -- cycle;
\fill[pthick] (0.944,-2.597) -- (0.763,-2.845) -- (0.763,-2.349) -- cycle;
\fill[pthin] (0.763,-2.845) -- (0.944,-2.597) -- (1.055,-2.751) -- cycle;
\fill[pthick] (0.472,-1.145) -- (0.763,-1.051) -- (0.472,-1.452) -- cycle;
\fill[pthick] (0.292,-1.393) -- (0.472,-1.145) -- (-0.000,-1.299) -- cycle;
\fill[pthin] (0.472,-1.145) -- (0.292,-1.393) -- (0.472,-1.452) -- cycle;
\fill[pthick] (0.944,-1.605) -- (1.235,-1.700) -- (0.944,-1.299) -- cycle;
\fill[pthick] (0.763,-1.357) -- (0.944,-1.605) -- (0.472,-1.452) -- cycle;
\fill[pthin] (0.944,-1.605) -- (0.763,-1.357) -- (0.944,-1.299) -- cycle;
\fill[pthin] (0.763,-1.051) -- (0.763,-1.357) -- (0.944,-1.299) -- cycle;
\fill[pthick] (0.763,-1.357) -- (0.763,-1.051) -- (0.472,-1.452) -- cycle;
\fill[pthick] (-0.000,-0.496) -- (-0.000,-0.803) -- (0.292,-0.401) -- cycle;
\fill[pthick] (0.180,-0.248) -- (-0.000,-0.496) -- (0.000,0.000) -- cycle;
\fill[pthin] (-0.000,-0.496) -- (0.180,-0.248) -- (0.292,-0.401) -- cycle;
\fill[pthick] (0.583,-0.803) -- (0.763,-1.051) -- (0.292,-0.897) -- cycle;
\fill[pthick] (0.292,-0.708) -- (0.583,-0.803) -- (0.292,-0.401) -- cycle;
\fill[pthin] (0.583,-0.803) -- (0.292,-0.708) -- (0.292,-0.897) -- cycle;
\fill[pthin] (-0.000,-0.803) -- (0.292,-0.708) -- (0.292,-0.897) -- cycle;
\fill[pthick] (0.292,-0.708) -- (-0.000,-0.803) -- (0.292,-0.401) -- cycle;
\fill[pthin] (-0.000,-0.803) -- (0.180,-1.051) -- (0.292,-0.897) -- cycle;
\fill[pthick] (0.180,-1.051) -- (-0.000,-0.803) -- (-0.000,-1.299) -- cycle;
\fill[pthick] (0.472,-1.145) -- (0.763,-1.051) -- (0.292,-0.897) -- cycle;
\fill[pthick] (0.180,-1.051) -- (0.472,-1.145) -- (-0.000,-1.299) -- cycle;
\fill[pthin] (0.472,-1.145) -- (0.180,-1.051) -- (0.292,-0.897) -- cycle;
\fill[pthick] (0.292,-2.503) -- (-0.000,-2.597) -- (0.472,-2.751) -- cycle;
\fill[pthick] (0.583,-2.597) -- (0.292,-2.503) -- (0.763,-2.349) -- cycle;
\fill[pthin] (0.292,-2.503) -- (0.583,-2.597) -- (0.472,-2.751) -- cycle;
\fill[pthick] (0.180,-3.152) -- (-0.000,-3.400) -- (-0.000,-2.904) -- cycle;
\fill[pthick] (0.180,-2.845) -- (0.180,-3.152) -- (0.472,-2.751) -- cycle;
\fill[pthin] (0.180,-3.152) -- (0.180,-2.845) -- (-0.000,-2.904) -- cycle;
\fill[pthin] (-0.000,-2.597) -- (0.180,-2.845) -- (-0.000,-2.904) -- cycle;
\fill[pthick] (0.180,-2.845) -- (-0.000,-2.597) -- (0.472,-2.751) -- cycle;
\fill[pthick] (0.292,-1.700) -- (0.472,-1.948) -- (-0.000,-1.795) -- cycle;
\fill[pthick] (-0.000,-1.605) -- (0.292,-1.700) -- (-0.000,-1.299) -- cycle;
\fill[pthin] (0.292,-1.700) -- (-0.000,-1.605) -- (-0.000,-1.795) -- cycle;
\fill[pthick] (-0.000,-2.291) -- (-0.000,-2.597) -- (0.292,-2.196) -- cycle;
\fill[pthick] (0.180,-2.043) -- (-0.000,-2.291) -- (-0.000,-1.795) -- cycle;
\fill[pthin] (-0.000,-2.291) -- (0.180,-2.043) -- (0.292,-2.196) -- cycle;
\fill[pthin] (0.472,-1.948) -- (0.180,-2.043) -- (0.292,-2.196) -- cycle;
\fill[pthick] (0.180,-2.043) -- (0.472,-1.948) -- (-0.000,-1.795) -- cycle;
\fill[pthin] (0.472,-1.948) -- (0.472,-2.255) -- (0.292,-2.196) -- cycle;
\fill[pthick] (0.472,-2.255) -- (0.472,-1.948) -- (0.763,-2.349) -- cycle;
\fill[pthick] (0.292,-2.503) -- (-0.000,-2.597) -- (0.292,-2.196) -- cycle;
\fill[pthick] (0.472,-2.255) -- (0.292,-2.503) -- (0.763,-2.349) -- cycle;
\fill[pthin] (0.292,-2.503) -- (0.472,-2.255) -- (0.292,-2.196) -- cycle;
\fill[pthin] (0.472,-1.948) -- (0.763,-2.043) -- (0.763,-1.853) -- cycle;
\fill[pthick] (0.763,-2.043) -- (0.472,-1.948) -- (0.763,-2.349) -- cycle;
\fill[pthick] (1.055,-1.948) -- (1.235,-1.700) -- (0.763,-1.853) -- cycle;
\fill[pthick] (0.763,-2.043) -- (1.055,-1.948) -- (0.763,-2.349) -- cycle;
\fill[pthin] (1.055,-1.948) -- (0.763,-2.043) -- (0.763,-1.853) -- cycle;
\fill[pthick] (0.292,-1.700) -- (0.472,-1.948) -- (0.472,-1.452) -- cycle;
\fill[pthick] (0.292,-1.393) -- (0.292,-1.700) -- (-0.000,-1.299) -- cycle;
\fill[pthin] (0.292,-1.700) -- (0.292,-1.393) -- (0.472,-1.452) -- cycle;
\fill[pthick] (0.944,-1.605) -- (1.235,-1.700) -- (0.763,-1.853) -- cycle;
\fill[pthick] (0.652,-1.700) -- (0.944,-1.605) -- (0.472,-1.452) -- cycle;
\fill[pthin] (0.944,-1.605) -- (0.652,-1.700) -- (0.763,-1.853) -- cycle;
\fill[pthin] (0.472,-1.948) -- (0.652,-1.700) -- (0.763,-1.853) -- cycle;
\fill[pthick] (0.652,-1.700) -- (0.472,-1.948) -- (0.472,-1.452) -- cycle;
\fill[pthin] (1.998,-1.948) -- (1.818,-1.700) -- (1.707,-1.853) -- cycle;
\fill[pthick] (1.818,-1.700) -- (1.998,-1.948) -- (1.998,-1.452) -- cycle;
\fill[pthick] (1.527,-1.605) -- (1.235,-1.700) -- (1.707,-1.853) -- cycle;
\fill[pthick] (1.818,-1.700) -- (1.527,-1.605) -- (1.998,-1.452) -- cycle;
\fill[pthin] (1.527,-1.605) -- (1.818,-1.700) -- (1.707,-1.853) -- cycle;
\fill[pthick] (1.998,-2.255) -- (1.998,-1.948) -- (1.707,-2.349) -- cycle;
\fill[pthick] (1.818,-2.503) -- (1.998,-2.255) -- (1.998,-2.751) -- cycle;
\fill[pthin] (1.998,-2.255) -- (1.818,-2.503) -- (1.707,-2.349) -- cycle;
\fill[pthick] (1.415,-1.948) -- (1.235,-1.700) -- (1.707,-1.853) -- cycle;
\fill[pthick] (1.707,-2.043) -- (1.415,-1.948) -- (1.707,-2.349) -- cycle;
\fill[pthin] (1.415,-1.948) -- (1.707,-2.043) -- (1.707,-1.853) -- cycle;
\fill[pthin] (1.998,-1.948) -- (1.707,-2.043) -- (1.707,-1.853) -- cycle;
\fill[pthick] (1.707,-2.043) -- (1.998,-1.948) -- (1.707,-2.349) -- cycle;
\fill[pthick] (2.470,-1.605) -- (2.762,-1.700) -- (2.470,-1.299) -- cycle;
\fill[pthick] (2.290,-1.357) -- (2.470,-1.605) -- (1.998,-1.452) -- cycle;
\fill[pthin] (2.470,-1.605) -- (2.290,-1.357) -- (2.470,-1.299) -- cycle;
\fill[pthick] (2.942,-1.145) -- (3.234,-1.051) -- (2.942,-1.452) -- cycle;
\fill[pthick] (2.762,-1.393) -- (2.942,-1.145) -- (2.470,-1.299) -- cycle;
\fill[pthin] (2.942,-1.145) -- (2.762,-1.393) -- (2.942,-1.452) -- cycle;
\fill[pthin] (2.762,-1.700) -- (2.762,-1.393) -- (2.942,-1.452) -- cycle;
\fill[pthick] (2.762,-1.393) -- (2.762,-1.700) -- (2.470,-1.299) -- cycle;
\fill[pthick] (1.998,-2.255) -- (1.998,-1.948) -- (2.290,-2.349) -- cycle;
\fill[pthick] (2.179,-2.503) -- (1.998,-2.255) -- (1.998,-2.751) -- cycle;
\fill[pthin] (1.998,-2.255) -- (2.179,-2.503) -- (2.290,-2.349) -- cycle;
\fill[pthick] (2.582,-1.948) -- (2.762,-1.700) -- (2.290,-1.853) -- cycle;
\fill[pthick] (2.290,-2.043) -- (2.582,-1.948) -- (2.290,-2.349) -- cycle;
\fill[pthin] (2.582,-1.948) -- (2.290,-2.043) -- (2.290,-1.853) -- cycle;
\fill[pthin] (1.998,-1.948) -- (2.290,-2.043) -- (2.290,-1.853) -- cycle;
\fill[pthick] (2.290,-2.043) -- (1.998,-1.948) -- (2.290,-2.349) -- cycle;
\fill[pthin] (1.998,-1.948) -- (2.179,-1.700) -- (2.290,-1.853) -- cycle;
\fill[pthick] (2.179,-1.700) -- (1.998,-1.948) -- (1.998,-1.452) -- cycle;
\fill[pthick] (2.470,-1.605) -- (2.762,-1.700) -- (2.290,-1.853) -- cycle;
\fill[pthick] (2.179,-1.700) -- (2.470,-1.605) -- (1.998,-1.452) -- cycle;
\fill[pthin] (2.470,-1.605) -- (2.179,-1.700) -- (2.290,-1.853) -- cycle;
\fill[pthick] (0.944,-0.803) -- (0.763,-1.051) -- (1.235,-0.897) -- cycle;
\fill[pthick] (1.235,-0.708) -- (0.944,-0.803) -- (1.235,-0.401) -- cycle;
\fill[pthin] (0.944,-0.803) -- (1.235,-0.708) -- (1.235,-0.897) -- cycle;
\fill[pthick] (1.235,-1.393) -- (1.235,-1.700) -- (0.944,-1.299) -- cycle;
\fill[pthick] (1.055,-1.145) -- (1.235,-1.393) -- (1.235,-0.897) -- cycle;
\fill[pthin] (1.235,-1.393) -- (1.055,-1.145) -- (0.944,-1.299) -- cycle;
\fill[pthin] (0.763,-1.051) -- (1.055,-1.145) -- (0.944,-1.299) -- cycle;
\fill[pthick] (1.055,-1.145) -- (0.763,-1.051) -- (1.235,-0.897) -- cycle;
\fill[pthick] (0.472,-0.153) -- (0.763,-0.248) -- (0.292,-0.401) -- cycle;
\fill[pthick] (0.180,-0.248) -- (0.472,-0.153) -- (0.000,0.000) -- cycle;
\fill[pthin] (0.472,-0.153) -- (0.180,-0.248) -- (0.292,-0.401) -- cycle;
\fill[pthick] (0.583,-0.803) -- (0.763,-1.051) -- (0.763,-0.555) -- cycle;
\fill[pthick] (0.583,-0.496) -- (0.583,-0.803) -- (0.292,-0.401) -- cycle;
\fill[pthin] (0.583,-0.803) -- (0.583,-0.496) -- (0.763,-0.555) -- cycle;
\fill[pthin] (0.763,-0.248) -- (0.583,-0.496) -- (0.763,-0.555) -- cycle;
\fill[pthick] (0.583,-0.496) -- (0.763,-0.248) -- (0.292,-0.401) -- cycle;
\fill[pthin] (0.763,-0.248) -- (0.944,-0.496) -- (0.763,-0.555) -- cycle;
\fill[pthick] (0.944,-0.496) -- (0.763,-0.248) -- (1.235,-0.401) -- cycle;
\fill[pthick] (0.944,-0.803) -- (0.763,-1.051) -- (0.763,-0.555) -- cycle;
\fill[pthick] (0.944,-0.496) -- (0.944,-0.803) -- (1.235,-0.401) -- cycle;
\fill[pthin] (0.944,-0.803) -- (0.944,-0.496) -- (0.763,-0.555) -- cycle;
\fill[pthick] (2.290,-1.051) -- (2.470,-0.803) -- (2.470,-1.299) -- cycle;
\fill[pthick] (2.290,-1.357) -- (2.290,-1.051) -- (1.998,-1.452) -- cycle;
\fill[pthin] (2.290,-1.051) -- (2.290,-1.357) -- (2.470,-1.299) -- cycle;
\fill[pthick] (2.942,-1.145) -- (3.234,-1.051) -- (2.762,-0.897) -- cycle;
\fill[pthick] (2.650,-1.051) -- (2.942,-1.145) -- (2.470,-1.299) -- cycle;
\fill[pthin] (2.942,-1.145) -- (2.650,-1.051) -- (2.762,-0.897) -- cycle;
\fill[pthin] (2.470,-0.803) -- (2.650,-1.051) -- (2.762,-0.897) -- cycle;
\fill[pthick] (2.650,-1.051) -- (2.470,-0.803) -- (2.470,-1.299) -- cycle;
\fill[pthick] (1.527,-0.803) -- (1.707,-1.051) -- (1.707,-0.555) -- cycle;
\fill[pthick] (1.527,-0.496) -- (1.527,-0.803) -- (1.235,-0.401) -- cycle;
\fill[pthin] (1.527,-0.803) -- (1.527,-0.496) -- (1.707,-0.555) -- cycle;
\fill[pthick] (2.179,-0.708) -- (2.470,-0.803) -- (1.998,-0.956) -- cycle;
\fill[pthick] (1.887,-0.803) -- (2.179,-0.708) -- (1.707,-0.555) -- cycle;
\fill[pthin] (2.179,-0.708) -- (1.887,-0.803) -- (1.998,-0.956) -- cycle;
\fill[pthin] (1.707,-1.051) -- (1.887,-0.803) -- (1.998,-0.956) -- cycle;
\fill[pthick] (1.887,-0.803) -- (1.707,-1.051) -- (1.707,-0.555) -- cycle;
\fill[pthin] (1.707,-1.051) -- (1.998,-1.145) -- (1.998,-0.956) -- cycle;
\fill[pthick] (1.998,-1.145) -- (1.707,-1.051) -- (1.998,-1.452) -- cycle;
\fill[pthick] (2.290,-1.051) -- (2.470,-0.803) -- (1.998,-0.956) -- cycle;
\fill[pthick] (1.998,-1.145) -- (2.290,-1.051) -- (1.998,-1.452) -- cycle;
\fill[pthin] (2.290,-1.051) -- (1.998,-1.145) -- (1.998,-0.956) -- cycle;
\fill[pthin] (1.707,-1.051) -- (1.707,-1.357) -- (1.527,-1.299) -- cycle;
\fill[pthick] (1.707,-1.357) -- (1.707,-1.051) -- (1.998,-1.452) -- cycle;
\fill[pthick] (1.527,-1.605) -- (1.235,-1.700) -- (1.527,-1.299) -- cycle;
\fill[pthick] (1.707,-1.357) -- (1.527,-1.605) -- (1.998,-1.452) -- cycle;
\fill[pthin] (1.527,-1.605) -- (1.707,-1.357) -- (1.527,-1.299) -- cycle;
\fill[pthick] (1.527,-0.803) -- (1.707,-1.051) -- (1.235,-0.897) -- cycle;
\fill[pthick] (1.235,-0.708) -- (1.527,-0.803) -- (1.235,-0.401) -- cycle;
\fill[pthin] (1.527,-0.803) -- (1.235,-0.708) -- (1.235,-0.897) -- cycle;
\fill[pthick] (1.235,-1.393) -- (1.235,-1.700) -- (1.527,-1.299) -- cycle;
\fill[pthick] (1.415,-1.145) -- (1.235,-1.393) -- (1.235,-0.897) -- cycle;
\fill[pthin] (1.235,-1.393) -- (1.415,-1.145) -- (1.527,-1.299) -- cycle;
\fill[pthin] (1.707,-1.051) -- (1.415,-1.145) -- (1.527,-1.299) -- cycle;
\fill[pthick] (1.415,-1.145) -- (1.707,-1.051) -- (1.235,-0.897) -- cycle;
\draw[pedge] (2.470,0.496) -- (2.762,0.401) -- (2.470,0.307);
\draw[pedge] (2.470,0.000) -- (2.470,0.307) -- (2.762,0.401);
\draw[pedge] (2.470,0.496) -- (2.179,0.401) -- (1.998,0.649);
\draw[pedge] (2.470,0.000) -- (2.470,0.307) -- (2.179,0.401);
\draw[pedge] (2.470,0.496) -- (2.179,0.401) -- (2.470,0.307);
\draw[pedge] (2.762,0.897) -- (2.942,0.649) -- (2.762,0.401);
\draw[pedge] (3.234,1.051) -- (2.942,0.956) -- (2.942,0.649);
\draw[pedge] (2.762,0.897) -- (2.942,0.649) -- (2.942,0.956);
\draw[pedge] (2.470,0.496) -- (2.290,0.744) -- (1.998,0.649);
\draw[pedge] (2.762,0.897) -- (2.582,0.649) -- (2.290,0.744);
\draw[pedge] (2.470,0.496) -- (2.290,0.744) -- (2.582,0.649);
\draw[pedge] (2.470,0.496) -- (2.762,0.401) -- (2.582,0.649);
\draw[pedge] (2.762,0.897) -- (2.582,0.649) -- (2.762,0.401);
\draw[pedge] (2.762,-0.401) -- (2.942,-0.153) -- (3.234,-0.248);
\draw[pedge] (2.470,0.000) -- (2.650,-0.248) -- (2.942,-0.153);
\draw[pedge] (2.762,-0.401) -- (2.942,-0.153) -- (2.650,-0.248);
\draw[pedge] (3.234,-0.555) -- (3.053,-0.803) -- (3.234,-1.051);
\draw[pedge] (2.762,-0.401) -- (3.053,-0.496) -- (3.053,-0.803);
\draw[pedge] (3.234,-0.555) -- (3.053,-0.803) -- (3.053,-0.496);
\draw[pedge] (3.234,-0.555) -- (3.234,-0.248) -- (3.053,-0.496);
\draw[pedge] (2.762,-0.401) -- (3.053,-0.496) -- (3.234,-0.248);
\draw[pedge] (3.234,0.555) -- (2.942,0.649) -- (2.762,0.401);
\draw[pedge] (3.234,1.051) -- (3.234,0.744) -- (2.942,0.649);
\draw[pedge] (3.234,0.555) -- (2.942,0.649) -- (3.234,0.744);
\draw[pedge] (2.942,0.153) -- (3.234,0.059) -- (3.234,-0.248);
\draw[pedge] (3.234,0.555) -- (3.053,0.307) -- (3.234,0.059);
\draw[pedge] (2.942,0.153) -- (3.234,0.059) -- (3.053,0.307);
\draw[pedge] (2.942,0.153) -- (2.762,0.401) -- (3.053,0.307);
\draw[pedge] (3.234,0.555) -- (3.053,0.307) -- (2.762,0.401);
\draw[pedge] (2.942,0.153) -- (2.762,0.401) -- (2.762,0.095);
\draw[pedge] (2.470,0.000) -- (2.762,0.095) -- (2.762,0.401);
\draw[pedge] (2.942,0.153) -- (2.942,-0.153) -- (3.234,-0.248);
\draw[pedge] (2.470,0.000) -- (2.762,0.095) -- (2.942,-0.153);
\draw[pedge] (2.942,0.153) -- (2.942,-0.153) -- (2.762,0.095);
\draw[pedge] (1.527,0.000) -- (1.235,0.095) -- (1.235,0.401);
\draw[pedge] (1.235,-0.401) -- (1.415,-0.153) -- (1.235,0.095);
\draw[pedge] (1.527,0.000) -- (1.235,0.095) -- (1.415,-0.153);
\draw[pedge] (1.527,0.496) -- (1.818,0.401) -- (1.998,0.649);
\draw[pedge] (1.527,0.000) -- (1.527,0.307) -- (1.818,0.401);
\draw[pedge] (1.527,0.496) -- (1.818,0.401) -- (1.527,0.307);
\draw[pedge] (1.527,0.496) -- (1.235,0.401) -- (1.527,0.307);
\draw[pedge] (1.527,0.000) -- (1.527,0.307) -- (1.235,0.401);
\draw[pedge] (0.472,0.153) -- (0.472,-0.153) -- (0.763,-0.248);
\draw[pedge] (0.000,0.000) -- (0.292,0.095) -- (0.472,-0.153);
\draw[pedge] (0.472,0.153) -- (0.472,-0.153) -- (0.292,0.095);
\draw[pedge] (0.944,0.000) -- (0.944,0.307) -- (1.235,0.401);
\draw[pedge] (0.472,0.153) -- (0.763,0.059) -- (0.944,0.307);
\draw[pedge] (0.944,0.000) -- (0.944,0.307) -- (0.763,0.059);
\draw[pedge] (0.944,0.000) -- (0.763,-0.248) -- (0.763,0.059);
\draw[pedge] (0.472,0.153) -- (0.763,0.059) -- (0.763,-0.248);
\draw[pedge] (0.944,0.000) -- (0.763,-0.248) -- (1.055,-0.153);
\draw[pedge] (1.235,-0.401) -- (1.055,-0.153) -- (0.763,-0.248);
\draw[pedge] (0.944,0.000) -- (1.235,0.095) -- (1.235,0.401);
\draw[pedge] (1.235,-0.401) -- (1.055,-0.153) -- (1.235,0.095);
\draw[pedge] (0.944,0.000) -- (1.235,0.095) -- (1.055,-0.153);
\draw[pedge] (2.762,-0.401) -- (2.470,-0.496) -- (2.470,-0.803);
\draw[pedge] (2.470,0.000) -- (2.650,-0.248) -- (2.470,-0.496);
\draw[pedge] (2.762,-0.401) -- (2.470,-0.496) -- (2.650,-0.248);
\draw[pedge] (2.762,-0.897) -- (3.053,-0.803) -- (3.234,-1.051);
\draw[pedge] (2.762,-0.401) -- (2.762,-0.708) -- (3.053,-0.803);
\draw[pedge] (2.762,-0.897) -- (3.053,-0.803) -- (2.762,-0.708);
\draw[pedge] (2.762,-0.897) -- (2.470,-0.803) -- (2.762,-0.708);
\draw[pedge] (2.762,-0.401) -- (2.762,-0.708) -- (2.470,-0.803);
\draw[pedge] (1.707,-0.555) -- (1.707,-0.248) -- (1.998,-0.153);
\draw[pedge] (1.235,-0.401) -- (1.527,-0.496) -- (1.707,-0.248);
\draw[pedge] (1.707,-0.555) -- (1.707,-0.248) -- (1.527,-0.496);
\draw[pedge] (2.179,-0.401) -- (2.179,-0.708) -- (2.470,-0.803);
\draw[pedge] (1.707,-0.555) -- (1.998,-0.460) -- (2.179,-0.708);
\draw[pedge] (2.179,-0.401) -- (2.179,-0.708) -- (1.998,-0.460);
\draw[pedge] (2.179,-0.401) -- (1.998,-0.153) -- (1.998,-0.460);
\draw[pedge] (1.707,-0.555) -- (1.998,-0.460) -- (1.998,-0.153);
\draw[pedge] (2.179,-0.401) -- (1.998,-0.153) -- (2.290,-0.248);
\draw[pedge] (2.470,0.000) -- (2.290,-0.248) -- (1.998,-0.153);
\draw[pedge] (2.179,-0.401) -- (2.470,-0.496) -- (2.470,-0.803);
\draw[pedge] (2.470,0.000) -- (2.290,-0.248) -- (2.470,-0.496);
\draw[pedge] (2.179,-0.401) -- (2.470,-0.496) -- (2.290,-0.248);
\draw[pedge] (1.998,0.153) -- (1.998,-0.153) -- (2.179,0.095);
\draw[pedge] (2.470,0.000) -- (2.179,0.095) -- (1.998,-0.153);
\draw[pedge] (1.998,0.153) -- (2.179,0.401) -- (1.998,0.649);
\draw[pedge] (2.470,0.000) -- (2.179,0.095) -- (2.179,0.401);
\draw[pedge] (1.998,0.153) -- (2.179,0.401) -- (2.179,0.095);
\draw[pedge] (1.527,0.000) -- (1.707,-0.248) -- (1.998,-0.153);
\draw[pedge] (1.235,-0.401) -- (1.415,-0.153) -- (1.707,-0.248);
\draw[pedge] (1.527,0.000) -- (1.707,-0.248) -- (1.415,-0.153);
\draw[pedge] (1.998,0.153) -- (1.818,0.401) -- (1.998,0.649);
\draw[pedge] (1.527,0.000) -- (1.818,0.095) -- (1.818,0.401);
\draw[pedge] (1.998,0.153) -- (1.818,0.401) -- (1.818,0.095);
\draw[pedge] (1.998,0.153) -- (1.998,-0.153) -- (1.818,0.095);
\draw[pedge] (1.527,0.000) -- (1.818,0.095) -- (1.998,-0.153);
\draw[pedge] (2.290,1.051) -- (2.470,1.299) -- (2.179,1.204);
\draw[pedge] (1.998,1.452) -- (2.179,1.204) -- (2.470,1.299);
\draw[pedge] (2.290,1.051) -- (1.998,0.956) -- (1.998,0.649);
\draw[pedge] (1.998,1.452) -- (2.179,1.204) -- (1.998,0.956);
\draw[pedge] (2.290,1.051) -- (1.998,0.956) -- (2.179,1.204);
\draw[pedge] (2.762,0.897) -- (2.762,1.204) -- (2.470,1.299);
\draw[pedge] (3.234,1.051) -- (2.942,0.956) -- (2.762,1.204);
\draw[pedge] (2.762,0.897) -- (2.762,1.204) -- (2.942,0.956);
\draw[pedge] (2.290,1.051) -- (2.290,0.744) -- (1.998,0.649);
\draw[pedge] (2.762,0.897) -- (2.470,0.992) -- (2.290,0.744);
\draw[pedge] (2.290,1.051) -- (2.290,0.744) -- (2.470,0.992);
\draw[pedge] (2.290,1.051) -- (2.470,1.299) -- (2.470,0.992);
\draw[pedge] (2.762,0.897) -- (2.470,0.992) -- (2.470,1.299);
\draw[pedge] (1.998,1.948) -- (2.290,1.853) -- (2.470,2.101);
\draw[pedge] (1.998,1.452) -- (1.998,1.759) -- (2.290,1.853);
\draw[pedge] (1.998,1.948) -- (2.290,1.853) -- (1.998,1.759);
\draw[pedge] (2.290,2.349) -- (1.998,2.444) -- (1.998,2.751);
\draw[pedge] (1.998,1.948) -- (2.179,2.196) -- (1.998,2.444);
\draw[pedge] (2.290,2.349) -- (1.998,2.444) -- (2.179,2.196);
\draw[pedge] (2.290,2.349) -- (2.470,2.101) -- (2.179,2.196);
\draw[pedge] (1.998,1.948) -- (2.179,2.196) -- (2.470,2.101);
\draw[pedge] (2.942,1.452) -- (2.762,1.204) -- (2.470,1.299);
\draw[pedge] (3.234,1.051) -- (3.053,1.299) -- (2.762,1.204);
\draw[pedge] (2.942,1.452) -- (2.762,1.204) -- (3.053,1.299);
\draw[pedge] (2.470,1.605) -- (2.650,1.853) -- (2.470,2.101);
\draw[pedge] (2.942,1.452) -- (2.650,1.547) -- (2.650,1.853);
\draw[pedge] (2.470,1.605) -- (2.650,1.853) -- (2.650,1.547);
\draw[pedge] (2.470,1.605) -- (2.470,1.299) -- (2.650,1.547);
\draw[pedge] (2.942,1.452) -- (2.650,1.547) -- (2.470,1.299);
\draw[pedge] (2.470,1.605) -- (2.470,1.299) -- (2.290,1.547);
\draw[pedge] (1.998,1.452) -- (2.290,1.547) -- (2.470,1.299);
\draw[pedge] (2.470,1.605) -- (2.290,1.853) -- (2.470,2.101);
\draw[pedge] (1.998,1.452) -- (2.290,1.547) -- (2.290,1.853);
\draw[pedge] (2.470,1.605) -- (2.290,1.853) -- (2.290,1.547);
\draw[pedge] (1.235,0.897) -- (1.055,0.649) -- (1.235,0.401);
\draw[pedge] (0.763,1.051) -- (1.055,0.956) -- (1.055,0.649);
\draw[pedge] (1.235,0.897) -- (1.055,0.649) -- (1.055,0.956);
\draw[pedge] (1.527,0.496) -- (1.707,0.744) -- (1.998,0.649);
\draw[pedge] (1.235,0.897) -- (1.415,0.649) -- (1.707,0.744);
\draw[pedge] (1.527,0.496) -- (1.707,0.744) -- (1.415,0.649);
\draw[pedge] (1.527,0.496) -- (1.235,0.401) -- (1.415,0.649);
\draw[pedge] (1.235,0.897) -- (1.415,0.649) -- (1.235,0.401);
\draw[pedge] (0.472,0.153) -- (0.292,0.401) -- (0.472,0.649);
\draw[pedge] (0.000,0.000) -- (0.292,0.095) -- (0.292,0.401);
\draw[pedge] (0.472,0.153) -- (0.292,0.401) -- (0.292,0.095);
\draw[pedge] (0.763,0.555) -- (0.944,0.307) -- (1.235,0.401);
\draw[pedge] (0.472,0.153) -- (0.652,0.401) -- (0.944,0.307);
\draw[pedge] (0.763,0.555) -- (0.944,0.307) -- (0.652,0.401);
\draw[pedge] (0.763,0.555) -- (0.472,0.649) -- (0.652,0.401);
\draw[pedge] (0.472,0.153) -- (0.652,0.401) -- (0.472,0.649);
\draw[pedge] (0.763,0.555) -- (0.472,0.649) -- (0.763,0.744);
\draw[pedge] (0.763,1.051) -- (0.763,0.744) -- (0.472,0.649);
\draw[pedge] (0.763,0.555) -- (1.055,0.649) -- (1.235,0.401);
\draw[pedge] (0.763,1.051) -- (0.763,0.744) -- (1.055,0.649);
\draw[pedge] (0.763,0.555) -- (1.055,0.649) -- (0.763,0.744);
\draw[pedge] (1.998,1.948) -- (1.707,1.853) -- (1.527,2.101);
\draw[pedge] (1.998,1.452) -- (1.998,1.759) -- (1.707,1.853);
\draw[pedge] (1.998,1.948) -- (1.707,1.853) -- (1.998,1.759);
\draw[pedge] (1.707,2.349) -- (1.998,2.444) -- (1.998,2.751);
\draw[pedge] (1.998,1.948) -- (1.818,2.196) -- (1.998,2.444);
\draw[pedge] (1.707,2.349) -- (1.998,2.444) -- (1.818,2.196);
\draw[pedge] (1.707,2.349) -- (1.527,2.101) -- (1.818,2.196);
\draw[pedge] (1.998,1.948) -- (1.818,2.196) -- (1.527,2.101);
\draw[pedge] (1.055,1.452) -- (1.235,1.204) -- (1.527,1.299);
\draw[pedge] (0.763,1.051) -- (0.944,1.299) -- (1.235,1.204);
\draw[pedge] (1.055,1.452) -- (1.235,1.204) -- (0.944,1.299);
\draw[pedge] (1.527,1.605) -- (1.346,1.853) -- (1.527,2.101);
\draw[pedge] (1.055,1.452) -- (1.346,1.547) -- (1.346,1.853);
\draw[pedge] (1.527,1.605) -- (1.346,1.853) -- (1.346,1.547);
\draw[pedge] (1.527,1.605) -- (1.527,1.299) -- (1.346,1.547);
\draw[pedge] (1.055,1.452) -- (1.346,1.547) -- (1.527,1.299);
\draw[pedge] (1.527,1.605) -- (1.527,1.299) -- (1.707,1.547);
\draw[pedge] (1.998,1.452) -- (1.707,1.547) -- (1.527,1.299);
\draw[pedge] (1.527,1.605) -- (1.707,1.853) -- (1.527,2.101);
\draw[pedge] (1.998,1.452) -- (1.707,1.547) -- (1.707,1.853);
\draw[pedge] (1.527,1.605) -- (1.707,1.853) -- (1.707,1.547);
\draw[pedge] (1.707,1.051) -- (1.527,1.299) -- (1.818,1.204);
\draw[pedge] (1.998,1.452) -- (1.818,1.204) -- (1.527,1.299);
\draw[pedge] (1.707,1.051) -- (1.998,0.956) -- (1.998,0.649);
\draw[pedge] (1.998,1.452) -- (1.818,1.204) -- (1.998,0.956);
\draw[pedge] (1.707,1.051) -- (1.998,0.956) -- (1.818,1.204);
\draw[pedge] (1.235,0.897) -- (1.235,1.204) -- (1.527,1.299);
\draw[pedge] (0.763,1.051) -- (1.055,0.956) -- (1.235,1.204);
\draw[pedge] (1.235,0.897) -- (1.235,1.204) -- (1.055,0.956);
\draw[pedge] (1.707,1.051) -- (1.707,0.744) -- (1.998,0.649);
\draw[pedge] (1.235,0.897) -- (1.527,0.992) -- (1.707,0.744);
\draw[pedge] (1.707,1.051) -- (1.707,0.744) -- (1.527,0.992);
\draw[pedge] (1.707,1.051) -- (1.527,1.299) -- (1.527,0.992);
\draw[pedge] (1.235,0.897) -- (1.527,0.992) -- (1.527,1.299);
\draw[pedge] (0.292,2.503) -- (0.472,2.751) -- (0.472,2.444);
\draw[pedge] (0.763,2.349) -- (0.472,2.444) -- (0.472,2.751);
\draw[pedge] (0.292,2.503) -- (0.292,2.196) -- (0.000,2.101);
\draw[pedge] (0.763,2.349) -- (0.472,2.444) -- (0.292,2.196);
\draw[pedge] (0.292,2.503) -- (0.292,2.196) -- (0.472,2.444);
\draw[pedge] (0.000,2.904) -- (0.292,2.999) -- (0.472,2.751);
\draw[pedge] (0.000,3.400) -- (0.000,3.093) -- (0.292,2.999);
\draw[pedge] (0.000,2.904) -- (0.292,2.999) -- (0.000,3.093);
\draw[pedge] (0.292,2.503) -- (0.000,2.408) -- (0.000,2.101);
\draw[pedge] (0.000,2.904) -- (0.180,2.656) -- (0.000,2.408);
\draw[pedge] (0.292,2.503) -- (0.000,2.408) -- (0.180,2.656);
\draw[pedge] (0.292,2.503) -- (0.472,2.751) -- (0.180,2.656);
\draw[pedge] (0.000,2.904) -- (0.180,2.656) -- (0.472,2.751);
\draw[pedge] (1.235,2.503) -- (1.055,2.751) -- (1.235,2.999);
\draw[pedge] (0.763,2.349) -- (1.055,2.444) -- (1.055,2.751);
\draw[pedge] (1.235,2.503) -- (1.055,2.751) -- (1.055,2.444);
\draw[pedge] (1.527,2.904) -- (1.707,2.656) -- (1.998,2.751);
\draw[pedge] (1.235,2.503) -- (1.415,2.751) -- (1.707,2.656);
\draw[pedge] (1.527,2.904) -- (1.707,2.656) -- (1.415,2.751);
\draw[pedge] (1.527,2.904) -- (1.235,2.999) -- (1.415,2.751);
\draw[pedge] (1.235,2.503) -- (1.415,2.751) -- (1.235,2.999);
\draw[pedge] (0.472,3.247) -- (0.292,2.999) -- (0.472,2.751);
\draw[pedge] (0.000,3.400) -- (0.292,3.305) -- (0.292,2.999);
\draw[pedge] (0.472,3.247) -- (0.292,2.999) -- (0.292,3.305);
\draw[pedge] (0.763,2.845) -- (0.944,3.093) -- (1.235,2.999);
\draw[pedge] (0.472,3.247) -- (0.652,2.999) -- (0.944,3.093);
\draw[pedge] (0.763,2.845) -- (0.944,3.093) -- (0.652,2.999);
\draw[pedge] (0.763,2.845) -- (0.472,2.751) -- (0.652,2.999);
\draw[pedge] (0.472,3.247) -- (0.652,2.999) -- (0.472,2.751);
\draw[pedge] (0.763,2.845) -- (0.472,2.751) -- (0.763,2.656);
\draw[pedge] (0.763,2.349) -- (0.763,2.656) -- (0.472,2.751);
\draw[pedge] (0.763,2.845) -- (1.055,2.751) -- (1.235,2.999);
\draw[pedge] (0.763,2.349) -- (0.763,2.656) -- (1.055,2.751);
\draw[pedge] (0.763,2.845) -- (1.055,2.751) -- (0.763,2.656);
\draw[pedge] (0.472,1.452) -- (0.292,1.204) -- (0.000,1.299);
\draw[pedge] (0.763,1.051) -- (0.583,1.299) -- (0.292,1.204);
\draw[pedge] (0.472,1.452) -- (0.292,1.204) -- (0.583,1.299);
\draw[pedge] (0.000,1.605) -- (0.180,1.853) -- (0.000,2.101);
\draw[pedge] (0.472,1.452) -- (0.180,1.547) -- (0.180,1.853);
\draw[pedge] (0.000,1.605) -- (0.180,1.853) -- (0.180,1.547);
\draw[pedge] (0.000,1.605) -- (0.000,1.299) -- (0.180,1.547);
\draw[pedge] (0.472,1.452) -- (0.180,1.547) -- (0.000,1.299);
\draw[pedge] (0.000,0.496) -- (0.292,0.401) -- (0.472,0.649);
\draw[pedge] (0.000,0.000) -- (0.000,0.307) -- (0.292,0.401);
\draw[pedge] (0.000,0.496) -- (0.292,0.401) -- (0.000,0.307);
\draw[pedge] (0.292,0.897) -- (0.000,0.992) -- (0.000,1.299);
\draw[pedge] (0.000,0.496) -- (0.180,0.744) -- (0.000,0.992);
\draw[pedge] (0.292,0.897) -- (0.000,0.992) -- (0.180,0.744);
\draw[pedge] (0.292,0.897) -- (0.472,0.649) -- (0.180,0.744);
\draw[pedge] (0.000,0.496) -- (0.180,0.744) -- (0.472,0.649);
\draw[pedge] (0.292,0.897) -- (0.472,0.649) -- (0.472,0.956);
\draw[pedge] (0.763,1.051) -- (0.472,0.956) -- (0.472,0.649);
\draw[pedge] (0.292,0.897) -- (0.292,1.204) -- (0.000,1.299);
\draw[pedge] (0.763,1.051) -- (0.472,0.956) -- (0.292,1.204);
\draw[pedge] (0.292,0.897) -- (0.292,1.204) -- (0.472,0.956);
\draw[pedge] (1.235,2.503) -- (1.235,2.196) -- (1.527,2.101);
\draw[pedge] (0.763,2.349) -- (1.055,2.444) -- (1.235,2.196);
\draw[pedge] (1.235,2.503) -- (1.235,2.196) -- (1.055,2.444);
\draw[pedge] (1.707,2.349) -- (1.707,2.656) -- (1.998,2.751);
\draw[pedge] (1.235,2.503) -- (1.527,2.408) -- (1.707,2.656);
\draw[pedge] (1.707,2.349) -- (1.707,2.656) -- (1.527,2.408);
\draw[pedge] (1.707,2.349) -- (1.527,2.101) -- (1.527,2.408);
\draw[pedge] (1.235,2.503) -- (1.527,2.408) -- (1.527,2.101);
\draw[pedge] (1.055,1.452) -- (0.763,1.547) -- (0.763,1.853);
\draw[pedge] (0.763,1.051) -- (0.944,1.299) -- (0.763,1.547);
\draw[pedge] (1.055,1.452) -- (0.763,1.547) -- (0.944,1.299);
\draw[pedge] (1.055,1.948) -- (1.346,1.853) -- (1.527,2.101);
\draw[pedge] (1.055,1.452) -- (1.055,1.759) -- (1.346,1.853);
\draw[pedge] (1.055,1.948) -- (1.346,1.853) -- (1.055,1.759);
\draw[pedge] (1.055,1.948) -- (0.763,1.853) -- (1.055,1.759);
\draw[pedge] (1.055,1.452) -- (1.055,1.759) -- (0.763,1.853);
\draw[pedge] (1.055,1.948) -- (0.763,1.853) -- (0.944,2.101);
\draw[pedge] (0.763,2.349) -- (0.944,2.101) -- (0.763,1.853);
\draw[pedge] (1.055,1.948) -- (1.235,2.196) -- (1.527,2.101);
\draw[pedge] (0.763,2.349) -- (0.944,2.101) -- (1.235,2.196);
\draw[pedge] (1.055,1.948) -- (1.235,2.196) -- (0.944,2.101);
\draw[pedge] (0.472,1.948) -- (0.763,1.853) -- (0.583,2.101);
\draw[pedge] (0.763,2.349) -- (0.583,2.101) -- (0.763,1.853);
\draw[pedge] (0.472,1.948) -- (0.292,2.196) -- (0.000,2.101);
\draw[pedge] (0.763,2.349) -- (0.583,2.101) -- (0.292,2.196);
\draw[pedge] (0.472,1.948) -- (0.292,2.196) -- (0.583,2.101);
\draw[pedge] (0.472,1.452) -- (0.763,1.547) -- (0.763,1.853);
\draw[pedge] (0.763,1.051) -- (0.583,1.299) -- (0.763,1.547);
\draw[pedge] (0.472,1.452) -- (0.763,1.547) -- (0.583,1.299);
\draw[pedge] (0.472,1.948) -- (0.180,1.853) -- (0.000,2.101);
\draw[pedge] (0.472,1.452) -- (0.472,1.759) -- (0.180,1.853);
\draw[pedge] (0.472,1.948) -- (0.180,1.853) -- (0.472,1.759);
\draw[pedge] (0.472,1.948) -- (0.763,1.853) -- (0.472,1.759);
\draw[pedge] (0.472,1.452) -- (0.472,1.759) -- (0.763,1.853);
\draw[pedge] (-0.292,2.503) -- (-0.472,2.751) -- (-0.472,2.444);
\draw[pedge] (-0.763,2.349) -- (-0.472,2.444) -- (-0.472,2.751);
\draw[pedge] (-0.292,2.503) -- (-0.292,2.196) -- (0.000,2.101);
\draw[pedge] (-0.763,2.349) -- (-0.472,2.444) -- (-0.292,2.196);
\draw[pedge] (-0.292,2.503) -- (-0.292,2.196) -- (-0.472,2.444);
\draw[pedge] (0.000,2.904) -- (-0.292,2.999) -- (-0.472,2.751);
\draw[pedge] (0.000,3.400) -- (0.000,3.093) -- (-0.292,2.999);
\draw[pedge] (0.000,2.904) -- (-0.292,2.999) -- (0.000,3.093);
\draw[pedge] (-0.292,2.503) -- (0.000,2.408) -- (0.000,2.101);
\draw[pedge] (0.000,2.904) -- (-0.180,2.656) -- (0.000,2.408);
\draw[pedge] (-0.292,2.503) -- (0.000,2.408) -- (-0.180,2.656);
\draw[pedge] (-0.292,2.503) -- (-0.472,2.751) -- (-0.180,2.656);
\draw[pedge] (0.000,2.904) -- (-0.180,2.656) -- (-0.472,2.751);
\draw[pedge] (-1.235,2.503) -- (-1.055,2.751) -- (-1.235,2.999);
\draw[pedge] (-0.763,2.349) -- (-1.055,2.444) -- (-1.055,2.751);
\draw[pedge] (-1.235,2.503) -- (-1.055,2.751) -- (-1.055,2.444);
\draw[pedge] (-1.527,2.904) -- (-1.707,2.656) -- (-1.998,2.751);
\draw[pedge] (-1.235,2.503) -- (-1.415,2.751) -- (-1.707,2.656);
\draw[pedge] (-1.527,2.904) -- (-1.707,2.656) -- (-1.415,2.751);
\draw[pedge] (-1.527,2.904) -- (-1.235,2.999) -- (-1.415,2.751);
\draw[pedge] (-1.235,2.503) -- (-1.415,2.751) -- (-1.235,2.999);
\draw[pedge] (-0.472,3.247) -- (-0.292,2.999) -- (-0.472,2.751);
\draw[pedge] (0.000,3.400) -- (-0.292,3.305) -- (-0.292,2.999);
\draw[pedge] (-0.472,3.247) -- (-0.292,2.999) -- (-0.292,3.305);
\draw[pedge] (-0.763,2.845) -- (-0.944,3.093) -- (-1.235,2.999);
\draw[pedge] (-0.472,3.247) -- (-0.652,2.999) -- (-0.944,3.093);
\draw[pedge] (-0.763,2.845) -- (-0.944,3.093) -- (-0.652,2.999);
\draw[pedge] (-0.763,2.845) -- (-0.472,2.751) -- (-0.652,2.999);
\draw[pedge] (-0.472,3.247) -- (-0.652,2.999) -- (-0.472,2.751);
\draw[pedge] (-0.763,2.845) -- (-0.472,2.751) -- (-0.763,2.656);
\draw[pedge] (-0.763,2.349) -- (-0.763,2.656) -- (-0.472,2.751);
\draw[pedge] (-0.763,2.845) -- (-1.055,2.751) -- (-1.235,2.999);
\draw[pedge] (-0.763,2.349) -- (-0.763,2.656) -- (-1.055,2.751);
\draw[pedge] (-0.763,2.845) -- (-1.055,2.751) -- (-0.763,2.656);
\draw[pedge] (-0.472,1.452) -- (-0.292,1.204) -- (0.000,1.299);
\draw[pedge] (-0.763,1.051) -- (-0.583,1.299) -- (-0.292,1.204);
\draw[pedge] (-0.472,1.452) -- (-0.292,1.204) -- (-0.583,1.299);
\draw[pedge] (0.000,1.605) -- (-0.180,1.853) -- (0.000,2.101);
\draw[pedge] (-0.472,1.452) -- (-0.180,1.547) -- (-0.180,1.853);
\draw[pedge] (0.000,1.605) -- (-0.180,1.853) -- (-0.180,1.547);
\draw[pedge] (0.000,1.605) -- (0.000,1.299) -- (-0.180,1.547);
\draw[pedge] (-0.472,1.452) -- (-0.180,1.547) -- (0.000,1.299);
\draw[pedge] (0.000,0.496) -- (-0.292,0.401) -- (-0.472,0.649);
\draw[pedge] (0.000,0.000) -- (0.000,0.307) -- (-0.292,0.401);
\draw[pedge] (0.000,0.496) -- (-0.292,0.401) -- (0.000,0.307);
\draw[pedge] (-0.292,0.897) -- (0.000,0.992) -- (0.000,1.299);
\draw[pedge] (0.000,0.496) -- (-0.180,0.744) -- (0.000,0.992);
\draw[pedge] (-0.292,0.897) -- (0.000,0.992) -- (-0.180,0.744);
\draw[pedge] (-0.292,0.897) -- (-0.472,0.649) -- (-0.180,0.744);
\draw[pedge] (0.000,0.496) -- (-0.180,0.744) -- (-0.472,0.649);
\draw[pedge] (-0.292,0.897) -- (-0.472,0.649) -- (-0.472,0.956);
\draw[pedge] (-0.763,1.051) -- (-0.472,0.956) -- (-0.472,0.649);
\draw[pedge] (-0.292,0.897) -- (-0.292,1.204) -- (0.000,1.299);
\draw[pedge] (-0.763,1.051) -- (-0.472,0.956) -- (-0.292,1.204);
\draw[pedge] (-0.292,0.897) -- (-0.292,1.204) -- (-0.472,0.956);
\draw[pedge] (-1.235,2.503) -- (-1.235,2.196) -- (-1.527,2.101);
\draw[pedge] (-0.763,2.349) -- (-1.055,2.444) -- (-1.235,2.196);
\draw[pedge] (-1.235,2.503) -- (-1.235,2.196) -- (-1.055,2.444);
\draw[pedge] (-1.707,2.349) -- (-1.707,2.656) -- (-1.998,2.751);
\draw[pedge] (-1.235,2.503) -- (-1.527,2.408) -- (-1.707,2.656);
\draw[pedge] (-1.707,2.349) -- (-1.707,2.656) -- (-1.527,2.408);
\draw[pedge] (-1.707,2.349) -- (-1.527,2.101) -- (-1.527,2.408);
\draw[pedge] (-1.235,2.503) -- (-1.527,2.408) -- (-1.527,2.101);
\draw[pedge] (-1.055,1.452) -- (-0.763,1.547) -- (-0.763,1.853);
\draw[pedge] (-0.763,1.051) -- (-0.944,1.299) -- (-0.763,1.547);
\draw[pedge] (-1.055,1.452) -- (-0.763,1.547) -- (-0.944,1.299);
\draw[pedge] (-1.055,1.948) -- (-1.346,1.853) -- (-1.527,2.101);
\draw[pedge] (-1.055,1.452) -- (-1.055,1.759) -- (-1.346,1.853);
\draw[pedge] (-1.055,1.948) -- (-1.346,1.853) -- (-1.055,1.759);
\draw[pedge] (-1.055,1.948) -- (-0.763,1.853) -- (-1.055,1.759);
\draw[pedge] (-1.055,1.452) -- (-1.055,1.759) -- (-0.763,1.853);
\draw[pedge] (-1.055,1.948) -- (-0.763,1.853) -- (-0.944,2.101);
\draw[pedge] (-0.763,2.349) -- (-0.944,2.101) -- (-0.763,1.853);
\draw[pedge] (-1.055,1.948) -- (-1.235,2.196) -- (-1.527,2.101);
\draw[pedge] (-0.763,2.349) -- (-0.944,2.101) -- (-1.235,2.196);
\draw[pedge] (-1.055,1.948) -- (-1.235,2.196) -- (-0.944,2.101);
\draw[pedge] (-0.472,1.948) -- (-0.763,1.853) -- (-0.583,2.101);
\draw[pedge] (-0.763,2.349) -- (-0.583,2.101) -- (-0.763,1.853);
\draw[pedge] (-0.472,1.948) -- (-0.292,2.196) -- (0.000,2.101);
\draw[pedge] (-0.763,2.349) -- (-0.583,2.101) -- (-0.292,2.196);
\draw[pedge] (-0.472,1.948) -- (-0.292,2.196) -- (-0.583,2.101);
\draw[pedge] (-0.472,1.452) -- (-0.763,1.547) -- (-0.763,1.853);
\draw[pedge] (-0.763,1.051) -- (-0.583,1.299) -- (-0.763,1.547);
\draw[pedge] (-0.472,1.452) -- (-0.763,1.547) -- (-0.583,1.299);
\draw[pedge] (-0.472,1.948) -- (-0.180,1.853) -- (0.000,2.101);
\draw[pedge] (-0.472,1.452) -- (-0.472,1.759) -- (-0.180,1.853);
\draw[pedge] (-0.472,1.948) -- (-0.180,1.853) -- (-0.472,1.759);
\draw[pedge] (-0.472,1.948) -- (-0.763,1.853) -- (-0.472,1.759);
\draw[pedge] (-0.472,1.452) -- (-0.472,1.759) -- (-0.763,1.853);
\draw[pedge] (-2.290,1.051) -- (-2.470,1.299) -- (-2.179,1.204);
\draw[pedge] (-1.998,1.452) -- (-2.179,1.204) -- (-2.470,1.299);
\draw[pedge] (-2.290,1.051) -- (-1.998,0.956) -- (-1.998,0.649);
\draw[pedge] (-1.998,1.452) -- (-2.179,1.204) -- (-1.998,0.956);
\draw[pedge] (-2.290,1.051) -- (-1.998,0.956) -- (-2.179,1.204);
\draw[pedge] (-2.762,0.897) -- (-2.762,1.204) -- (-2.470,1.299);
\draw[pedge] (-3.234,1.051) -- (-2.942,0.956) -- (-2.762,1.204);
\draw[pedge] (-2.762,0.897) -- (-2.762,1.204) -- (-2.942,0.956);
\draw[pedge] (-2.290,1.051) -- (-2.290,0.744) -- (-1.998,0.649);
\draw[pedge] (-2.762,0.897) -- (-2.470,0.992) -- (-2.290,0.744);
\draw[pedge] (-2.290,1.051) -- (-2.290,0.744) -- (-2.470,0.992);
\draw[pedge] (-2.290,1.051) -- (-2.470,1.299) -- (-2.470,0.992);
\draw[pedge] (-2.762,0.897) -- (-2.470,0.992) -- (-2.470,1.299);
\draw[pedge] (-1.998,1.948) -- (-2.290,1.853) -- (-2.470,2.101);
\draw[pedge] (-1.998,1.452) -- (-1.998,1.759) -- (-2.290,1.853);
\draw[pedge] (-1.998,1.948) -- (-2.290,1.853) -- (-1.998,1.759);
\draw[pedge] (-2.290,2.349) -- (-1.998,2.444) -- (-1.998,2.751);
\draw[pedge] (-1.998,1.948) -- (-2.179,2.196) -- (-1.998,2.444);
\draw[pedge] (-2.290,2.349) -- (-1.998,2.444) -- (-2.179,2.196);
\draw[pedge] (-2.290,2.349) -- (-2.470,2.101) -- (-2.179,2.196);
\draw[pedge] (-1.998,1.948) -- (-2.179,2.196) -- (-2.470,2.101);
\draw[pedge] (-2.942,1.452) -- (-2.762,1.204) -- (-2.470,1.299);
\draw[pedge] (-3.234,1.051) -- (-3.053,1.299) -- (-2.762,1.204);
\draw[pedge] (-2.942,1.452) -- (-2.762,1.204) -- (-3.053,1.299);
\draw[pedge] (-2.470,1.605) -- (-2.650,1.853) -- (-2.470,2.101);
\draw[pedge] (-2.942,1.452) -- (-2.650,1.547) -- (-2.650,1.853);
\draw[pedge] (-2.470,1.605) -- (-2.650,1.853) -- (-2.650,1.547);
\draw[pedge] (-2.470,1.605) -- (-2.470,1.299) -- (-2.650,1.547);
\draw[pedge] (-2.942,1.452) -- (-2.650,1.547) -- (-2.470,1.299);
\draw[pedge] (-2.470,1.605) -- (-2.470,1.299) -- (-2.290,1.547);
\draw[pedge] (-1.998,1.452) -- (-2.290,1.547) -- (-2.470,1.299);
\draw[pedge] (-2.470,1.605) -- (-2.290,1.853) -- (-2.470,2.101);
\draw[pedge] (-1.998,1.452) -- (-2.290,1.547) -- (-2.290,1.853);
\draw[pedge] (-2.470,1.605) -- (-2.290,1.853) -- (-2.290,1.547);
\draw[pedge] (-1.235,0.897) -- (-1.055,0.649) -- (-1.235,0.401);
\draw[pedge] (-0.763,1.051) -- (-1.055,0.956) -- (-1.055,0.649);
\draw[pedge] (-1.235,0.897) -- (-1.055,0.649) -- (-1.055,0.956);
\draw[pedge] (-1.527,0.496) -- (-1.707,0.744) -- (-1.998,0.649);
\draw[pedge] (-1.235,0.897) -- (-1.415,0.649) -- (-1.707,0.744);
\draw[pedge] (-1.527,0.496) -- (-1.707,0.744) -- (-1.415,0.649);
\draw[pedge] (-1.527,0.496) -- (-1.235,0.401) -- (-1.415,0.649);
\draw[pedge] (-1.235,0.897) -- (-1.415,0.649) -- (-1.235,0.401);
\draw[pedge] (-0.472,0.153) -- (-0.292,0.401) -- (-0.472,0.649);
\draw[pedge] (0.000,0.000) -- (-0.292,0.095) -- (-0.292,0.401);
\draw[pedge] (-0.472,0.153) -- (-0.292,0.401) -- (-0.292,0.095);
\draw[pedge] (-0.763,0.555) -- (-0.944,0.307) -- (-1.235,0.401);
\draw[pedge] (-0.472,0.153) -- (-0.652,0.401) -- (-0.944,0.307);
\draw[pedge] (-0.763,0.555) -- (-0.944,0.307) -- (-0.652,0.401);
\draw[pedge] (-0.763,0.555) -- (-0.472,0.649) -- (-0.652,0.401);
\draw[pedge] (-0.472,0.153) -- (-0.652,0.401) -- (-0.472,0.649);
\draw[pedge] (-0.763,0.555) -- (-0.472,0.649) -- (-0.763,0.744);
\draw[pedge] (-0.763,1.051) -- (-0.763,0.744) -- (-0.472,0.649);
\draw[pedge] (-0.763,0.555) -- (-1.055,0.649) -- (-1.235,0.401);
\draw[pedge] (-0.763,1.051) -- (-0.763,0.744) -- (-1.055,0.649);
\draw[pedge] (-0.763,0.555) -- (-1.055,0.649) -- (-0.763,0.744);
\draw[pedge] (-1.998,1.948) -- (-1.707,1.853) -- (-1.527,2.101);
\draw[pedge] (-1.998,1.452) -- (-1.998,1.759) -- (-1.707,1.853);
\draw[pedge] (-1.998,1.948) -- (-1.707,1.853) -- (-1.998,1.759);
\draw[pedge] (-1.707,2.349) -- (-1.998,2.444) -- (-1.998,2.751);
\draw[pedge] (-1.998,1.948) -- (-1.818,2.196) -- (-1.998,2.444);
\draw[pedge] (-1.707,2.349) -- (-1.998,2.444) -- (-1.818,2.196);
\draw[pedge] (-1.707,2.349) -- (-1.527,2.101) -- (-1.818,2.196);
\draw[pedge] (-1.998,1.948) -- (-1.818,2.196) -- (-1.527,2.101);
\draw[pedge] (-1.055,1.452) -- (-1.235,1.204) -- (-1.527,1.299);
\draw[pedge] (-0.763,1.051) -- (-0.944,1.299) -- (-1.235,1.204);
\draw[pedge] (-1.055,1.452) -- (-1.235,1.204) -- (-0.944,1.299);
\draw[pedge] (-1.527,1.605) -- (-1.346,1.853) -- (-1.527,2.101);
\draw[pedge] (-1.055,1.452) -- (-1.346,1.547) -- (-1.346,1.853);
\draw[pedge] (-1.527,1.605) -- (-1.346,1.853) -- (-1.346,1.547);
\draw[pedge] (-1.527,1.605) -- (-1.527,1.299) -- (-1.346,1.547);
\draw[pedge] (-1.055,1.452) -- (-1.346,1.547) -- (-1.527,1.299);
\draw[pedge] (-1.527,1.605) -- (-1.527,1.299) -- (-1.707,1.547);
\draw[pedge] (-1.998,1.452) -- (-1.707,1.547) -- (-1.527,1.299);
\draw[pedge] (-1.527,1.605) -- (-1.707,1.853) -- (-1.527,2.101);
\draw[pedge] (-1.998,1.452) -- (-1.707,1.547) -- (-1.707,1.853);
\draw[pedge] (-1.527,1.605) -- (-1.707,1.853) -- (-1.707,1.547);
\draw[pedge] (-1.707,1.051) -- (-1.527,1.299) -- (-1.818,1.204);
\draw[pedge] (-1.998,1.452) -- (-1.818,1.204) -- (-1.527,1.299);
\draw[pedge] (-1.707,1.051) -- (-1.998,0.956) -- (-1.998,0.649);
\draw[pedge] (-1.998,1.452) -- (-1.818,1.204) -- (-1.998,0.956);
\draw[pedge] (-1.707,1.051) -- (-1.998,0.956) -- (-1.818,1.204);
\draw[pedge] (-1.235,0.897) -- (-1.235,1.204) -- (-1.527,1.299);
\draw[pedge] (-0.763,1.051) -- (-1.055,0.956) -- (-1.235,1.204);
\draw[pedge] (-1.235,0.897) -- (-1.235,1.204) -- (-1.055,0.956);
\draw[pedge] (-1.707,1.051) -- (-1.707,0.744) -- (-1.998,0.649);
\draw[pedge] (-1.235,0.897) -- (-1.527,0.992) -- (-1.707,0.744);
\draw[pedge] (-1.707,1.051) -- (-1.707,0.744) -- (-1.527,0.992);
\draw[pedge] (-1.707,1.051) -- (-1.527,1.299) -- (-1.527,0.992);
\draw[pedge] (-1.235,0.897) -- (-1.527,0.992) -- (-1.527,1.299);
\draw[pedge] (-2.470,0.496) -- (-2.762,0.401) -- (-2.470,0.307);
\draw[pedge] (-2.470,0.000) -- (-2.470,0.307) -- (-2.762,0.401);
\draw[pedge] (-2.470,0.496) -- (-2.179,0.401) -- (-1.998,0.649);
\draw[pedge] (-2.470,0.000) -- (-2.470,0.307) -- (-2.179,0.401);
\draw[pedge] (-2.470,0.496) -- (-2.179,0.401) -- (-2.470,0.307);
\draw[pedge] (-2.762,0.897) -- (-2.942,0.649) -- (-2.762,0.401);
\draw[pedge] (-3.234,1.051) -- (-2.942,0.956) -- (-2.942,0.649);
\draw[pedge] (-2.762,0.897) -- (-2.942,0.649) -- (-2.942,0.956);
\draw[pedge] (-2.470,0.496) -- (-2.290,0.744) -- (-1.998,0.649);
\draw[pedge] (-2.762,0.897) -- (-2.582,0.649) -- (-2.290,0.744);
\draw[pedge] (-2.470,0.496) -- (-2.290,0.744) -- (-2.582,0.649);
\draw[pedge] (-2.470,0.496) -- (-2.762,0.401) -- (-2.582,0.649);
\draw[pedge] (-2.762,0.897) -- (-2.582,0.649) -- (-2.762,0.401);
\draw[pedge] (-2.762,-0.401) -- (-2.942,-0.153) -- (-3.234,-0.248);
\draw[pedge] (-2.470,0.000) -- (-2.650,-0.248) -- (-2.942,-0.153);
\draw[pedge] (-2.762,-0.401) -- (-2.942,-0.153) -- (-2.650,-0.248);
\draw[pedge] (-3.234,-0.555) -- (-3.053,-0.803) -- (-3.234,-1.051);
\draw[pedge] (-2.762,-0.401) -- (-3.053,-0.496) -- (-3.053,-0.803);
\draw[pedge] (-3.234,-0.555) -- (-3.053,-0.803) -- (-3.053,-0.496);
\draw[pedge] (-3.234,-0.555) -- (-3.234,-0.248) -- (-3.053,-0.496);
\draw[pedge] (-2.762,-0.401) -- (-3.053,-0.496) -- (-3.234,-0.248);
\draw[pedge] (-3.234,0.555) -- (-2.942,0.649) -- (-2.762,0.401);
\draw[pedge] (-3.234,1.051) -- (-3.234,0.744) -- (-2.942,0.649);
\draw[pedge] (-3.234,0.555) -- (-2.942,0.649) -- (-3.234,0.744);
\draw[pedge] (-2.942,0.153) -- (-3.234,0.059) -- (-3.234,-0.248);
\draw[pedge] (-3.234,0.555) -- (-3.053,0.307) -- (-3.234,0.059);
\draw[pedge] (-2.942,0.153) -- (-3.234,0.059) -- (-3.053,0.307);
\draw[pedge] (-2.942,0.153) -- (-2.762,0.401) -- (-3.053,0.307);
\draw[pedge] (-3.234,0.555) -- (-3.053,0.307) -- (-2.762,0.401);
\draw[pedge] (-2.942,0.153) -- (-2.762,0.401) -- (-2.762,0.095);
\draw[pedge] (-2.470,0.000) -- (-2.762,0.095) -- (-2.762,0.401);
\draw[pedge] (-2.942,0.153) -- (-2.942,-0.153) -- (-3.234,-0.248);
\draw[pedge] (-2.470,0.000) -- (-2.762,0.095) -- (-2.942,-0.153);
\draw[pedge] (-2.942,0.153) -- (-2.942,-0.153) -- (-2.762,0.095);
\draw[pedge] (-1.527,0.000) -- (-1.235,0.095) -- (-1.235,0.401);
\draw[pedge] (-1.235,-0.401) -- (-1.415,-0.153) -- (-1.235,0.095);
\draw[pedge] (-1.527,0.000) -- (-1.235,0.095) -- (-1.415,-0.153);
\draw[pedge] (-1.527,0.496) -- (-1.818,0.401) -- (-1.998,0.649);
\draw[pedge] (-1.527,0.000) -- (-1.527,0.307) -- (-1.818,0.401);
\draw[pedge] (-1.527,0.496) -- (-1.818,0.401) -- (-1.527,0.307);
\draw[pedge] (-1.527,0.496) -- (-1.235,0.401) -- (-1.527,0.307);
\draw[pedge] (-1.527,0.000) -- (-1.527,0.307) -- (-1.235,0.401);
\draw[pedge] (-0.472,0.153) -- (-0.472,-0.153) -- (-0.763,-0.248);
\draw[pedge] (0.000,0.000) -- (-0.292,0.095) -- (-0.472,-0.153);
\draw[pedge] (-0.472,0.153) -- (-0.472,-0.153) -- (-0.292,0.095);
\draw[pedge] (-0.944,0.000) -- (-0.944,0.307) -- (-1.235,0.401);
\draw[pedge] (-0.472,0.153) -- (-0.763,0.059) -- (-0.944,0.307);
\draw[pedge] (-0.944,0.000) -- (-0.944,0.307) -- (-0.763,0.059);
\draw[pedge] (-0.944,0.000) -- (-0.763,-0.248) -- (-0.763,0.059);
\draw[pedge] (-0.472,0.153) -- (-0.763,0.059) -- (-0.763,-0.248);
\draw[pedge] (-0.944,0.000) -- (-0.763,-0.248) -- (-1.055,-0.153);
\draw[pedge] (-1.235,-0.401) -- (-1.055,-0.153) -- (-0.763,-0.248);
\draw[pedge] (-0.944,0.000) -- (-1.235,0.095) -- (-1.235,0.401);
\draw[pedge] (-1.235,-0.401) -- (-1.055,-0.153) -- (-1.235,0.095);
\draw[pedge] (-0.944,0.000) -- (-1.235,0.095) -- (-1.055,-0.153);
\draw[pedge] (-2.762,-0.401) -- (-2.470,-0.496) -- (-2.470,-0.803);
\draw[pedge] (-2.470,0.000) -- (-2.650,-0.248) -- (-2.470,-0.496);
\draw[pedge] (-2.762,-0.401) -- (-2.470,-0.496) -- (-2.650,-0.248);
\draw[pedge] (-2.762,-0.897) -- (-3.053,-0.803) -- (-3.234,-1.051);
\draw[pedge] (-2.762,-0.401) -- (-2.762,-0.708) -- (-3.053,-0.803);
\draw[pedge] (-2.762,-0.897) -- (-3.053,-0.803) -- (-2.762,-0.708);
\draw[pedge] (-2.762,-0.897) -- (-2.470,-0.803) -- (-2.762,-0.708);
\draw[pedge] (-2.762,-0.401) -- (-2.762,-0.708) -- (-2.470,-0.803);
\draw[pedge] (-1.707,-0.555) -- (-1.707,-0.248) -- (-1.998,-0.153);
\draw[pedge] (-1.235,-0.401) -- (-1.527,-0.496) -- (-1.707,-0.248);
\draw[pedge] (-1.707,-0.555) -- (-1.707,-0.248) -- (-1.527,-0.496);
\draw[pedge] (-2.179,-0.401) -- (-2.179,-0.708) -- (-2.470,-0.803);
\draw[pedge] (-1.707,-0.555) -- (-1.998,-0.460) -- (-2.179,-0.708);
\draw[pedge] (-2.179,-0.401) -- (-2.179,-0.708) -- (-1.998,-0.460);
\draw[pedge] (-2.179,-0.401) -- (-1.998,-0.153) -- (-1.998,-0.460);
\draw[pedge] (-1.707,-0.555) -- (-1.998,-0.460) -- (-1.998,-0.153);
\draw[pedge] (-2.179,-0.401) -- (-1.998,-0.153) -- (-2.290,-0.248);
\draw[pedge] (-2.470,0.000) -- (-2.290,-0.248) -- (-1.998,-0.153);
\draw[pedge] (-2.179,-0.401) -- (-2.470,-0.496) -- (-2.470,-0.803);
\draw[pedge] (-2.470,0.000) -- (-2.290,-0.248) -- (-2.470,-0.496);
\draw[pedge] (-2.179,-0.401) -- (-2.470,-0.496) -- (-2.290,-0.248);
\draw[pedge] (-1.998,0.153) -- (-1.998,-0.153) -- (-2.179,0.095);
\draw[pedge] (-2.470,0.000) -- (-2.179,0.095) -- (-1.998,-0.153);
\draw[pedge] (-1.998,0.153) -- (-2.179,0.401) -- (-1.998,0.649);
\draw[pedge] (-2.470,0.000) -- (-2.179,0.095) -- (-2.179,0.401);
\draw[pedge] (-1.998,0.153) -- (-2.179,0.401) -- (-2.179,0.095);
\draw[pedge] (-1.527,0.000) -- (-1.707,-0.248) -- (-1.998,-0.153);
\draw[pedge] (-1.235,-0.401) -- (-1.415,-0.153) -- (-1.707,-0.248);
\draw[pedge] (-1.527,0.000) -- (-1.707,-0.248) -- (-1.415,-0.153);
\draw[pedge] (-1.998,0.153) -- (-1.818,0.401) -- (-1.998,0.649);
\draw[pedge] (-1.527,0.000) -- (-1.818,0.095) -- (-1.818,0.401);
\draw[pedge] (-1.998,0.153) -- (-1.818,0.401) -- (-1.818,0.095);
\draw[pedge] (-1.998,0.153) -- (-1.998,-0.153) -- (-1.818,0.095);
\draw[pedge] (-1.527,0.000) -- (-1.818,0.095) -- (-1.998,-0.153);
\draw[pedge] (-1.707,-1.853) -- (-1.998,-1.948) -- (-1.818,-1.700);
\draw[pedge] (-1.998,-1.452) -- (-1.818,-1.700) -- (-1.998,-1.948);
\draw[pedge] (-1.707,-1.853) -- (-1.527,-1.605) -- (-1.235,-1.700);
\draw[pedge] (-1.998,-1.452) -- (-1.818,-1.700) -- (-1.527,-1.605);
\draw[pedge] (-1.707,-1.853) -- (-1.527,-1.605) -- (-1.818,-1.700);
\draw[pedge] (-1.707,-2.349) -- (-1.998,-2.255) -- (-1.998,-1.948);
\draw[pedge] (-1.998,-2.751) -- (-1.818,-2.503) -- (-1.998,-2.255);
\draw[pedge] (-1.707,-2.349) -- (-1.998,-2.255) -- (-1.818,-2.503);
\draw[pedge] (-1.707,-1.853) -- (-1.415,-1.948) -- (-1.235,-1.700);
\draw[pedge] (-1.707,-2.349) -- (-1.707,-2.043) -- (-1.415,-1.948);
\draw[pedge] (-1.707,-1.853) -- (-1.415,-1.948) -- (-1.707,-2.043);
\draw[pedge] (-1.707,-1.853) -- (-1.998,-1.948) -- (-1.707,-2.043);
\draw[pedge] (-1.707,-2.349) -- (-1.707,-2.043) -- (-1.998,-1.948);
\draw[pedge] (-2.470,-1.299) -- (-2.470,-1.605) -- (-2.762,-1.700);
\draw[pedge] (-1.998,-1.452) -- (-2.290,-1.357) -- (-2.470,-1.605);
\draw[pedge] (-2.470,-1.299) -- (-2.470,-1.605) -- (-2.290,-1.357);
\draw[pedge] (-2.942,-1.452) -- (-2.942,-1.145) -- (-3.234,-1.051);
\draw[pedge] (-2.470,-1.299) -- (-2.762,-1.393) -- (-2.942,-1.145);
\draw[pedge] (-2.942,-1.452) -- (-2.942,-1.145) -- (-2.762,-1.393);
\draw[pedge] (-2.942,-1.452) -- (-2.762,-1.700) -- (-2.762,-1.393);
\draw[pedge] (-2.470,-1.299) -- (-2.762,-1.393) -- (-2.762,-1.700);
\draw[pedge] (-2.290,-2.349) -- (-1.998,-2.255) -- (-1.998,-1.948);
\draw[pedge] (-1.998,-2.751) -- (-2.179,-2.503) -- (-1.998,-2.255);
\draw[pedge] (-2.290,-2.349) -- (-1.998,-2.255) -- (-2.179,-2.503);
\draw[pedge] (-2.290,-1.853) -- (-2.582,-1.948) -- (-2.762,-1.700);
\draw[pedge] (-2.290,-2.349) -- (-2.290,-2.043) -- (-2.582,-1.948);
\draw[pedge] (-2.290,-1.853) -- (-2.582,-1.948) -- (-2.290,-2.043);
\draw[pedge] (-2.290,-1.853) -- (-1.998,-1.948) -- (-2.290,-2.043);
\draw[pedge] (-2.290,-2.349) -- (-2.290,-2.043) -- (-1.998,-1.948);
\draw[pedge] (-2.290,-1.853) -- (-1.998,-1.948) -- (-2.179,-1.700);
\draw[pedge] (-1.998,-1.452) -- (-2.179,-1.700) -- (-1.998,-1.948);
\draw[pedge] (-2.290,-1.853) -- (-2.470,-1.605) -- (-2.762,-1.700);
\draw[pedge] (-1.998,-1.452) -- (-2.179,-1.700) -- (-2.470,-1.605);
\draw[pedge] (-2.290,-1.853) -- (-2.470,-1.605) -- (-2.179,-1.700);
\draw[pedge] (-1.235,-0.897) -- (-0.944,-0.803) -- (-0.763,-1.051);
\draw[pedge] (-1.235,-0.401) -- (-1.235,-0.708) -- (-0.944,-0.803);
\draw[pedge] (-1.235,-0.897) -- (-0.944,-0.803) -- (-1.235,-0.708);
\draw[pedge] (-0.944,-1.299) -- (-1.235,-1.393) -- (-1.235,-1.700);
\draw[pedge] (-1.235,-0.897) -- (-1.055,-1.145) -- (-1.235,-1.393);
\draw[pedge] (-0.944,-1.299) -- (-1.235,-1.393) -- (-1.055,-1.145);
\draw[pedge] (-0.944,-1.299) -- (-0.763,-1.051) -- (-1.055,-1.145);
\draw[pedge] (-1.235,-0.897) -- (-1.055,-1.145) -- (-0.763,-1.051);
\draw[pedge] (-0.292,-0.401) -- (-0.472,-0.153) -- (-0.763,-0.248);
\draw[pedge] (0.000,0.000) -- (-0.180,-0.248) -- (-0.472,-0.153);
\draw[pedge] (-0.292,-0.401) -- (-0.472,-0.153) -- (-0.180,-0.248);
\draw[pedge] (-0.763,-0.555) -- (-0.583,-0.803) -- (-0.763,-1.051);
\draw[pedge] (-0.292,-0.401) -- (-0.583,-0.496) -- (-0.583,-0.803);
\draw[pedge] (-0.763,-0.555) -- (-0.583,-0.803) -- (-0.583,-0.496);
\draw[pedge] (-0.763,-0.555) -- (-0.763,-0.248) -- (-0.583,-0.496);
\draw[pedge] (-0.292,-0.401) -- (-0.583,-0.496) -- (-0.763,-0.248);
\draw[pedge] (-0.763,-0.555) -- (-0.763,-0.248) -- (-0.944,-0.496);
\draw[pedge] (-1.235,-0.401) -- (-0.944,-0.496) -- (-0.763,-0.248);
\draw[pedge] (-0.763,-0.555) -- (-0.944,-0.803) -- (-0.763,-1.051);
\draw[pedge] (-1.235,-0.401) -- (-0.944,-0.496) -- (-0.944,-0.803);
\draw[pedge] (-0.763,-0.555) -- (-0.944,-0.803) -- (-0.944,-0.496);
\draw[pedge] (-2.470,-1.299) -- (-2.290,-1.051) -- (-2.470,-0.803);
\draw[pedge] (-1.998,-1.452) -- (-2.290,-1.357) -- (-2.290,-1.051);
\draw[pedge] (-2.470,-1.299) -- (-2.290,-1.051) -- (-2.290,-1.357);
\draw[pedge] (-2.762,-0.897) -- (-2.942,-1.145) -- (-3.234,-1.051);
\draw[pedge] (-2.470,-1.299) -- (-2.650,-1.051) -- (-2.942,-1.145);
\draw[pedge] (-2.762,-0.897) -- (-2.942,-1.145) -- (-2.650,-1.051);
\draw[pedge] (-2.762,-0.897) -- (-2.470,-0.803) -- (-2.650,-1.051);
\draw[pedge] (-2.470,-1.299) -- (-2.650,-1.051) -- (-2.470,-0.803);
\draw[pedge] (-1.707,-0.555) -- (-1.527,-0.803) -- (-1.707,-1.051);
\draw[pedge] (-1.235,-0.401) -- (-1.527,-0.496) -- (-1.527,-0.803);
\draw[pedge] (-1.707,-0.555) -- (-1.527,-0.803) -- (-1.527,-0.496);
\draw[pedge] (-1.998,-0.956) -- (-2.179,-0.708) -- (-2.470,-0.803);
\draw[pedge] (-1.707,-0.555) -- (-1.887,-0.803) -- (-2.179,-0.708);
\draw[pedge] (-1.998,-0.956) -- (-2.179,-0.708) -- (-1.887,-0.803);
\draw[pedge] (-1.998,-0.956) -- (-1.707,-1.051) -- (-1.887,-0.803);
\draw[pedge] (-1.707,-0.555) -- (-1.887,-0.803) -- (-1.707,-1.051);
\draw[pedge] (-1.998,-0.956) -- (-1.707,-1.051) -- (-1.998,-1.145);
\draw[pedge] (-1.998,-1.452) -- (-1.998,-1.145) -- (-1.707,-1.051);
\draw[pedge] (-1.998,-0.956) -- (-2.290,-1.051) -- (-2.470,-0.803);
\draw[pedge] (-1.998,-1.452) -- (-1.998,-1.145) -- (-2.290,-1.051);
\draw[pedge] (-1.998,-0.956) -- (-2.290,-1.051) -- (-1.998,-1.145);
\draw[pedge] (-1.527,-1.299) -- (-1.707,-1.051) -- (-1.707,-1.357);
\draw[pedge] (-1.998,-1.452) -- (-1.707,-1.357) -- (-1.707,-1.051);
\draw[pedge] (-1.527,-1.299) -- (-1.527,-1.605) -- (-1.235,-1.700);
\draw[pedge] (-1.998,-1.452) -- (-1.707,-1.357) -- (-1.527,-1.605);
\draw[pedge] (-1.527,-1.299) -- (-1.527,-1.605) -- (-1.707,-1.357);
\draw[pedge] (-1.235,-0.897) -- (-1.527,-0.803) -- (-1.707,-1.051);
\draw[pedge] (-1.235,-0.401) -- (-1.235,-0.708) -- (-1.527,-0.803);
\draw[pedge] (-1.235,-0.897) -- (-1.527,-0.803) -- (-1.235,-0.708);
\draw[pedge] (-1.527,-1.299) -- (-1.235,-1.393) -- (-1.235,-1.700);
\draw[pedge] (-1.235,-0.897) -- (-1.415,-1.145) -- (-1.235,-1.393);
\draw[pedge] (-1.527,-1.299) -- (-1.235,-1.393) -- (-1.415,-1.145);
\draw[pedge] (-1.527,-1.299) -- (-1.707,-1.051) -- (-1.415,-1.145);
\draw[pedge] (-1.235,-0.897) -- (-1.415,-1.145) -- (-1.707,-1.051);
\draw[pedge] (-1.235,-2.196) -- (-1.235,-2.503) -- (-1.055,-2.255);
\draw[pedge] (-0.763,-2.349) -- (-1.055,-2.255) -- (-1.235,-2.503);
\draw[pedge] (-1.235,-2.196) -- (-1.055,-1.948) -- (-1.235,-1.700);
\draw[pedge] (-0.763,-2.349) -- (-1.055,-2.255) -- (-1.055,-1.948);
\draw[pedge] (-1.235,-2.196) -- (-1.055,-1.948) -- (-1.055,-2.255);
\draw[pedge] (-1.707,-2.349) -- (-1.527,-2.597) -- (-1.235,-2.503);
\draw[pedge] (-1.998,-2.751) -- (-1.818,-2.503) -- (-1.527,-2.597);
\draw[pedge] (-1.707,-2.349) -- (-1.527,-2.597) -- (-1.818,-2.503);
\draw[pedge] (-1.235,-2.196) -- (-1.415,-1.948) -- (-1.235,-1.700);
\draw[pedge] (-1.707,-2.349) -- (-1.415,-2.255) -- (-1.415,-1.948);
\draw[pedge] (-1.235,-2.196) -- (-1.415,-1.948) -- (-1.415,-2.255);
\draw[pedge] (-1.235,-2.196) -- (-1.235,-2.503) -- (-1.415,-2.255);
\draw[pedge] (-1.707,-2.349) -- (-1.415,-2.255) -- (-1.235,-2.503);
\draw[pedge] (-0.472,-2.751) -- (-0.763,-2.845) -- (-0.763,-3.152);
\draw[pedge] (-0.763,-2.349) -- (-0.583,-2.597) -- (-0.763,-2.845);
\draw[pedge] (-0.472,-2.751) -- (-0.763,-2.845) -- (-0.583,-2.597);
\draw[pedge] (-0.472,-3.247) -- (-0.180,-3.152) -- (-0.000,-3.400);
\draw[pedge] (-0.472,-2.751) -- (-0.472,-3.057) -- (-0.180,-3.152);
\draw[pedge] (-0.472,-3.247) -- (-0.180,-3.152) -- (-0.472,-3.057);
\draw[pedge] (-0.472,-3.247) -- (-0.763,-3.152) -- (-0.472,-3.057);
\draw[pedge] (-0.472,-2.751) -- (-0.472,-3.057) -- (-0.763,-3.152);
\draw[pedge] (-1.527,-2.904) -- (-1.527,-2.597) -- (-1.235,-2.503);
\draw[pedge] (-1.998,-2.751) -- (-1.707,-2.845) -- (-1.527,-2.597);
\draw[pedge] (-1.527,-2.904) -- (-1.527,-2.597) -- (-1.707,-2.845);
\draw[pedge] (-1.055,-2.751) -- (-1.055,-3.057) -- (-0.763,-3.152);
\draw[pedge] (-1.527,-2.904) -- (-1.235,-2.809) -- (-1.055,-3.057);
\draw[pedge] (-1.055,-2.751) -- (-1.055,-3.057) -- (-1.235,-2.809);
\draw[pedge] (-1.055,-2.751) -- (-1.235,-2.503) -- (-1.235,-2.809);
\draw[pedge] (-1.527,-2.904) -- (-1.235,-2.809) -- (-1.235,-2.503);
\draw[pedge] (-1.055,-2.751) -- (-1.235,-2.503) -- (-0.944,-2.597);
\draw[pedge] (-0.763,-2.349) -- (-0.944,-2.597) -- (-1.235,-2.503);
\draw[pedge] (-1.055,-2.751) -- (-0.763,-2.845) -- (-0.763,-3.152);
\draw[pedge] (-0.763,-2.349) -- (-0.944,-2.597) -- (-0.763,-2.845);
\draw[pedge] (-1.055,-2.751) -- (-0.763,-2.845) -- (-0.944,-2.597);
\draw[pedge] (-0.472,-1.452) -- (-0.472,-1.145) -- (-0.763,-1.051);
\draw[pedge] (-0.000,-1.299) -- (-0.292,-1.393) -- (-0.472,-1.145);
\draw[pedge] (-0.472,-1.452) -- (-0.472,-1.145) -- (-0.292,-1.393);
\draw[pedge] (-0.944,-1.299) -- (-0.944,-1.605) -- (-1.235,-1.700);
\draw[pedge] (-0.472,-1.452) -- (-0.763,-1.357) -- (-0.944,-1.605);
\draw[pedge] (-0.944,-1.299) -- (-0.944,-1.605) -- (-0.763,-1.357);
\draw[pedge] (-0.944,-1.299) -- (-0.763,-1.051) -- (-0.763,-1.357);
\draw[pedge] (-0.472,-1.452) -- (-0.763,-1.357) -- (-0.763,-1.051);
\draw[pedge] (-0.292,-0.401) -- (-0.000,-0.496) -- (-0.000,-0.803);
\draw[pedge] (0.000,0.000) -- (-0.180,-0.248) -- (-0.000,-0.496);
\draw[pedge] (-0.292,-0.401) -- (-0.000,-0.496) -- (-0.180,-0.248);
\draw[pedge] (-0.292,-0.897) -- (-0.583,-0.803) -- (-0.763,-1.051);
\draw[pedge] (-0.292,-0.401) -- (-0.292,-0.708) -- (-0.583,-0.803);
\draw[pedge] (-0.292,-0.897) -- (-0.583,-0.803) -- (-0.292,-0.708);
\draw[pedge] (-0.292,-0.897) -- (-0.000,-0.803) -- (-0.292,-0.708);
\draw[pedge] (-0.292,-0.401) -- (-0.292,-0.708) -- (-0.000,-0.803);
\draw[pedge] (-0.292,-0.897) -- (-0.000,-0.803) -- (-0.180,-1.051);
\draw[pedge] (-0.000,-1.299) -- (-0.180,-1.051) -- (-0.000,-0.803);
\draw[pedge] (-0.292,-0.897) -- (-0.472,-1.145) -- (-0.763,-1.051);
\draw[pedge] (-0.000,-1.299) -- (-0.180,-1.051) -- (-0.472,-1.145);
\draw[pedge] (-0.292,-0.897) -- (-0.472,-1.145) -- (-0.180,-1.051);
\draw[pedge] (-0.472,-2.751) -- (-0.292,-2.503) -- (-0.000,-2.597);
\draw[pedge] (-0.763,-2.349) -- (-0.583,-2.597) -- (-0.292,-2.503);
\draw[pedge] (-0.472,-2.751) -- (-0.292,-2.503) -- (-0.583,-2.597);
\draw[pedge] (-0.000,-2.904) -- (-0.180,-3.152) -- (-0.000,-3.400);
\draw[pedge] (-0.472,-2.751) -- (-0.180,-2.845) -- (-0.180,-3.152);
\draw[pedge] (-0.000,-2.904) -- (-0.180,-3.152) -- (-0.180,-2.845);
\draw[pedge] (-0.000,-2.904) -- (-0.000,-2.597) -- (-0.180,-2.845);
\draw[pedge] (-0.472,-2.751) -- (-0.180,-2.845) -- (-0.000,-2.597);
\draw[pedge] (-0.000,-1.795) -- (-0.292,-1.700) -- (-0.472,-1.948);
\draw[pedge] (-0.000,-1.299) -- (-0.000,-1.605) -- (-0.292,-1.700);
\draw[pedge] (-0.000,-1.795) -- (-0.292,-1.700) -- (-0.000,-1.605);
\draw[pedge] (-0.292,-2.196) -- (-0.000,-2.291) -- (-0.000,-2.597);
\draw[pedge] (-0.000,-1.795) -- (-0.180,-2.043) -- (-0.000,-2.291);
\draw[pedge] (-0.292,-2.196) -- (-0.000,-2.291) -- (-0.180,-2.043);
\draw[pedge] (-0.292,-2.196) -- (-0.472,-1.948) -- (-0.180,-2.043);
\draw[pedge] (-0.000,-1.795) -- (-0.180,-2.043) -- (-0.472,-1.948);
\draw[pedge] (-0.292,-2.196) -- (-0.472,-1.948) -- (-0.472,-2.255);
\draw[pedge] (-0.763,-2.349) -- (-0.472,-2.255) -- (-0.472,-1.948);
\draw[pedge] (-0.292,-2.196) -- (-0.292,-2.503) -- (-0.000,-2.597);
\draw[pedge] (-0.763,-2.349) -- (-0.472,-2.255) -- (-0.292,-2.503);
\draw[pedge] (-0.292,-2.196) -- (-0.292,-2.503) -- (-0.472,-2.255);
\draw[pedge] (-0.763,-1.853) -- (-0.472,-1.948) -- (-0.763,-2.043);
\draw[pedge] (-0.763,-2.349) -- (-0.763,-2.043) -- (-0.472,-1.948);
\draw[pedge] (-0.763,-1.853) -- (-1.055,-1.948) -- (-1.235,-1.700);
\draw[pedge] (-0.763,-2.349) -- (-0.763,-2.043) -- (-1.055,-1.948);
\draw[pedge] (-0.763,-1.853) -- (-1.055,-1.948) -- (-0.763,-2.043);
\draw[pedge] (-0.472,-1.452) -- (-0.292,-1.700) -- (-0.472,-1.948);
\draw[pedge] (-0.000,-1.299) -- (-0.292,-1.393) -- (-0.292,-1.700);
\draw[pedge] (-0.472,-1.452) -- (-0.292,-1.700) -- (-0.292,-1.393);
\draw[pedge] (-0.763,-1.853) -- (-0.944,-1.605) -- (-1.235,-1.700);
\draw[pedge] (-0.472,-1.452) -- (-0.652,-1.700) -- (-0.944,-1.605);
\draw[pedge] (-0.763,-1.853) -- (-0.944,-1.605) -- (-0.652,-1.700);
\draw[pedge] (-0.763,-1.853) -- (-0.472,-1.948) -- (-0.652,-1.700);
\draw[pedge] (-0.472,-1.452) -- (-0.652,-1.700) -- (-0.472,-1.948);
\draw[pedge] (1.235,-2.196) -- (1.235,-2.503) -- (1.055,-2.255);
\draw[pedge] (0.763,-2.349) -- (1.055,-2.255) -- (1.235,-2.503);
\draw[pedge] (1.235,-2.196) -- (1.055,-1.948) -- (1.235,-1.700);
\draw[pedge] (0.763,-2.349) -- (1.055,-2.255) -- (1.055,-1.948);
\draw[pedge] (1.235,-2.196) -- (1.055,-1.948) -- (1.055,-2.255);
\draw[pedge] (1.707,-2.349) -- (1.527,-2.597) -- (1.235,-2.503);
\draw[pedge] (1.998,-2.751) -- (1.818,-2.503) -- (1.527,-2.597);
\draw[pedge] (1.707,-2.349) -- (1.527,-2.597) -- (1.818,-2.503);
\draw[pedge] (1.235,-2.196) -- (1.415,-1.948) -- (1.235,-1.700);
\draw[pedge] (1.707,-2.349) -- (1.415,-2.255) -- (1.415,-1.948);
\draw[pedge] (1.235,-2.196) -- (1.415,-1.948) -- (1.415,-2.255);
\draw[pedge] (1.235,-2.196) -- (1.235,-2.503) -- (1.415,-2.255);
\draw[pedge] (1.707,-2.349) -- (1.415,-2.255) -- (1.235,-2.503);
\draw[pedge] (0.472,-2.751) -- (0.763,-2.845) -- (0.763,-3.152);
\draw[pedge] (0.763,-2.349) -- (0.583,-2.597) -- (0.763,-2.845);
\draw[pedge] (0.472,-2.751) -- (0.763,-2.845) -- (0.583,-2.597);
\draw[pedge] (0.472,-3.247) -- (0.180,-3.152) -- (-0.000,-3.400);
\draw[pedge] (0.472,-2.751) -- (0.472,-3.057) -- (0.180,-3.152);
\draw[pedge] (0.472,-3.247) -- (0.180,-3.152) -- (0.472,-3.057);
\draw[pedge] (0.472,-3.247) -- (0.763,-3.152) -- (0.472,-3.057);
\draw[pedge] (0.472,-2.751) -- (0.472,-3.057) -- (0.763,-3.152);
\draw[pedge] (1.527,-2.904) -- (1.527,-2.597) -- (1.235,-2.503);
\draw[pedge] (1.998,-2.751) -- (1.707,-2.845) -- (1.527,-2.597);
\draw[pedge] (1.527,-2.904) -- (1.527,-2.597) -- (1.707,-2.845);
\draw[pedge] (1.055,-2.751) -- (1.055,-3.057) -- (0.763,-3.152);
\draw[pedge] (1.527,-2.904) -- (1.235,-2.809) -- (1.055,-3.057);
\draw[pedge] (1.055,-2.751) -- (1.055,-3.057) -- (1.235,-2.809);
\draw[pedge] (1.055,-2.751) -- (1.235,-2.503) -- (1.235,-2.809);
\draw[pedge] (1.527,-2.904) -- (1.235,-2.809) -- (1.235,-2.503);
\draw[pedge] (1.055,-2.751) -- (1.235,-2.503) -- (0.944,-2.597);
\draw[pedge] (0.763,-2.349) -- (0.944,-2.597) -- (1.235,-2.503);
\draw[pedge] (1.055,-2.751) -- (0.763,-2.845) -- (0.763,-3.152);
\draw[pedge] (0.763,-2.349) -- (0.944,-2.597) -- (0.763,-2.845);
\draw[pedge] (1.055,-2.751) -- (0.763,-2.845) -- (0.944,-2.597);
\draw[pedge] (0.472,-1.452) -- (0.472,-1.145) -- (0.763,-1.051);
\draw[pedge] (-0.000,-1.299) -- (0.292,-1.393) -- (0.472,-1.145);
\draw[pedge] (0.472,-1.452) -- (0.472,-1.145) -- (0.292,-1.393);
\draw[pedge] (0.944,-1.299) -- (0.944,-1.605) -- (1.235,-1.700);
\draw[pedge] (0.472,-1.452) -- (0.763,-1.357) -- (0.944,-1.605);
\draw[pedge] (0.944,-1.299) -- (0.944,-1.605) -- (0.763,-1.357);
\draw[pedge] (0.944,-1.299) -- (0.763,-1.051) -- (0.763,-1.357);
\draw[pedge] (0.472,-1.452) -- (0.763,-1.357) -- (0.763,-1.051);
\draw[pedge] (0.292,-0.401) -- (-0.000,-0.496) -- (-0.000,-0.803);
\draw[pedge] (0.000,0.000) -- (0.180,-0.248) -- (-0.000,-0.496);
\draw[pedge] (0.292,-0.401) -- (-0.000,-0.496) -- (0.180,-0.248);
\draw[pedge] (0.292,-0.897) -- (0.583,-0.803) -- (0.763,-1.051);
\draw[pedge] (0.292,-0.401) -- (0.292,-0.708) -- (0.583,-0.803);
\draw[pedge] (0.292,-0.897) -- (0.583,-0.803) -- (0.292,-0.708);
\draw[pedge] (0.292,-0.897) -- (-0.000,-0.803) -- (0.292,-0.708);
\draw[pedge] (0.292,-0.401) -- (0.292,-0.708) -- (-0.000,-0.803);
\draw[pedge] (0.292,-0.897) -- (-0.000,-0.803) -- (0.180,-1.051);
\draw[pedge] (-0.000,-1.299) -- (0.180,-1.051) -- (-0.000,-0.803);
\draw[pedge] (0.292,-0.897) -- (0.472,-1.145) -- (0.763,-1.051);
\draw[pedge] (-0.000,-1.299) -- (0.180,-1.051) -- (0.472,-1.145);
\draw[pedge] (0.292,-0.897) -- (0.472,-1.145) -- (0.180,-1.051);
\draw[pedge] (0.472,-2.751) -- (0.292,-2.503) -- (-0.000,-2.597);
\draw[pedge] (0.763,-2.349) -- (0.583,-2.597) -- (0.292,-2.503);
\draw[pedge] (0.472,-2.751) -- (0.292,-2.503) -- (0.583,-2.597);
\draw[pedge] (-0.000,-2.904) -- (0.180,-3.152) -- (-0.000,-3.400);
\draw[pedge] (0.472,-2.751) -- (0.180,-2.845) -- (0.180,-3.152);
\draw[pedge] (-0.000,-2.904) -- (0.180,-3.152) -- (0.180,-2.845);
\draw[pedge] (-0.000,-2.904) -- (-0.000,-2.597) -- (0.180,-2.845);
\draw[pedge] (0.472,-2.751) -- (0.180,-2.845) -- (-0.000,-2.597);
\draw[pedge] (-0.000,-1.795) -- (0.292,-1.700) -- (0.472,-1.948);
\draw[pedge] (-0.000,-1.299) -- (-0.000,-1.605) -- (0.292,-1.700);
\draw[pedge] (-0.000,-1.795) -- (0.292,-1.700) -- (-0.000,-1.605);
\draw[pedge] (0.292,-2.196) -- (-0.000,-2.291) -- (-0.000,-2.597);
\draw[pedge] (-0.000,-1.795) -- (0.180,-2.043) -- (-0.000,-2.291);
\draw[pedge] (0.292,-2.196) -- (-0.000,-2.291) -- (0.180,-2.043);
\draw[pedge] (0.292,-2.196) -- (0.472,-1.948) -- (0.180,-2.043);
\draw[pedge] (-0.000,-1.795) -- (0.180,-2.043) -- (0.472,-1.948);
\draw[pedge] (0.292,-2.196) -- (0.472,-1.948) -- (0.472,-2.255);
\draw[pedge] (0.763,-2.349) -- (0.472,-2.255) -- (0.472,-1.948);
\draw[pedge] (0.292,-2.196) -- (0.292,-2.503) -- (-0.000,-2.597);
\draw[pedge] (0.763,-2.349) -- (0.472,-2.255) -- (0.292,-2.503);
\draw[pedge] (0.292,-2.196) -- (0.292,-2.503) -- (0.472,-2.255);
\draw[pedge] (0.763,-1.853) -- (0.472,-1.948) -- (0.763,-2.043);
\draw[pedge] (0.763,-2.349) -- (0.763,-2.043) -- (0.472,-1.948);
\draw[pedge] (0.763,-1.853) -- (1.055,-1.948) -- (1.235,-1.700);
\draw[pedge] (0.763,-2.349) -- (0.763,-2.043) -- (1.055,-1.948);
\draw[pedge] (0.763,-1.853) -- (1.055,-1.948) -- (0.763,-2.043);
\draw[pedge] (0.472,-1.452) -- (0.292,-1.700) -- (0.472,-1.948);
\draw[pedge] (-0.000,-1.299) -- (0.292,-1.393) -- (0.292,-1.700);
\draw[pedge] (0.472,-1.452) -- (0.292,-1.700) -- (0.292,-1.393);
\draw[pedge] (0.763,-1.853) -- (0.944,-1.605) -- (1.235,-1.700);
\draw[pedge] (0.472,-1.452) -- (0.652,-1.700) -- (0.944,-1.605);
\draw[pedge] (0.763,-1.853) -- (0.944,-1.605) -- (0.652,-1.700);
\draw[pedge] (0.763,-1.853) -- (0.472,-1.948) -- (0.652,-1.700);
\draw[pedge] (0.472,-1.452) -- (0.652,-1.700) -- (0.472,-1.948);
\draw[pedge] (1.707,-1.853) -- (1.998,-1.948) -- (1.818,-1.700);
\draw[pedge] (1.998,-1.452) -- (1.818,-1.700) -- (1.998,-1.948);
\draw[pedge] (1.707,-1.853) -- (1.527,-1.605) -- (1.235,-1.700);
\draw[pedge] (1.998,-1.452) -- (1.818,-1.700) -- (1.527,-1.605);
\draw[pedge] (1.707,-1.853) -- (1.527,-1.605) -- (1.818,-1.700);
\draw[pedge] (1.707,-2.349) -- (1.998,-2.255) -- (1.998,-1.948);
\draw[pedge] (1.998,-2.751) -- (1.818,-2.503) -- (1.998,-2.255);
\draw[pedge] (1.707,-2.349) -- (1.998,-2.255) -- (1.818,-2.503);
\draw[pedge] (1.707,-1.853) -- (1.415,-1.948) -- (1.235,-1.700);
\draw[pedge] (1.707,-2.349) -- (1.707,-2.043) -- (1.415,-1.948);
\draw[pedge] (1.707,-1.853) -- (1.415,-1.948) -- (1.707,-2.043);
\draw[pedge] (1.707,-1.853) -- (1.998,-1.948) -- (1.707,-2.043);
\draw[pedge] (1.707,-2.349) -- (1.707,-2.043) -- (1.998,-1.948);
\draw[pedge] (2.470,-1.299) -- (2.470,-1.605) -- (2.762,-1.700);
\draw[pedge] (1.998,-1.452) -- (2.290,-1.357) -- (2.470,-1.605);
\draw[pedge] (2.470,-1.299) -- (2.470,-1.605) -- (2.290,-1.357);
\draw[pedge] (2.942,-1.452) -- (2.942,-1.145) -- (3.234,-1.051);
\draw[pedge] (2.470,-1.299) -- (2.762,-1.393) -- (2.942,-1.145);
\draw[pedge] (2.942,-1.452) -- (2.942,-1.145) -- (2.762,-1.393);
\draw[pedge] (2.942,-1.452) -- (2.762,-1.700) -- (2.762,-1.393);
\draw[pedge] (2.470,-1.299) -- (2.762,-1.393) -- (2.762,-1.700);
\draw[pedge] (2.290,-2.349) -- (1.998,-2.255) -- (1.998,-1.948);
\draw[pedge] (1.998,-2.751) -- (2.179,-2.503) -- (1.998,-2.255);
\draw[pedge] (2.290,-2.349) -- (1.998,-2.255) -- (2.179,-2.503);
\draw[pedge] (2.290,-1.853) -- (2.582,-1.948) -- (2.762,-1.700);
\draw[pedge] (2.290,-2.349) -- (2.290,-2.043) -- (2.582,-1.948);
\draw[pedge] (2.290,-1.853) -- (2.582,-1.948) -- (2.290,-2.043);
\draw[pedge] (2.290,-1.853) -- (1.998,-1.948) -- (2.290,-2.043);
\draw[pedge] (2.290,-2.349) -- (2.290,-2.043) -- (1.998,-1.948);
\draw[pedge] (2.290,-1.853) -- (1.998,-1.948) -- (2.179,-1.700);
\draw[pedge] (1.998,-1.452) -- (2.179,-1.700) -- (1.998,-1.948);
\draw[pedge] (2.290,-1.853) -- (2.470,-1.605) -- (2.762,-1.700);
\draw[pedge] (1.998,-1.452) -- (2.179,-1.700) -- (2.470,-1.605);
\draw[pedge] (2.290,-1.853) -- (2.470,-1.605) -- (2.179,-1.700);
\draw[pedge] (1.235,-0.897) -- (0.944,-0.803) -- (0.763,-1.051);
\draw[pedge] (1.235,-0.401) -- (1.235,-0.708) -- (0.944,-0.803);
\draw[pedge] (1.235,-0.897) -- (0.944,-0.803) -- (1.235,-0.708);
\draw[pedge] (0.944,-1.299) -- (1.235,-1.393) -- (1.235,-1.700);
\draw[pedge] (1.235,-0.897) -- (1.055,-1.145) -- (1.235,-1.393);
\draw[pedge] (0.944,-1.299) -- (1.235,-1.393) -- (1.055,-1.145);
\draw[pedge] (0.944,-1.299) -- (0.763,-1.051) -- (1.055,-1.145);
\draw[pedge] (1.235,-0.897) -- (1.055,-1.145) -- (0.763,-1.051);
\draw[pedge] (0.292,-0.401) -- (0.472,-0.153) -- (0.763,-0.248);
\draw[pedge] (0.000,0.000) -- (0.180,-0.248) -- (0.472,-0.153);
\draw[pedge] (0.292,-0.401) -- (0.472,-0.153) -- (0.180,-0.248);
\draw[pedge] (0.763,-0.555) -- (0.583,-0.803) -- (0.763,-1.051);
\draw[pedge] (0.292,-0.401) -- (0.583,-0.496) -- (0.583,-0.803);
\draw[pedge] (0.763,-0.555) -- (0.583,-0.803) -- (0.583,-0.496);
\draw[pedge] (0.763,-0.555) -- (0.763,-0.248) -- (0.583,-0.496);
\draw[pedge] (0.292,-0.401) -- (0.583,-0.496) -- (0.763,-0.248);
\draw[pedge] (0.763,-0.555) -- (0.763,-0.248) -- (0.944,-0.496);
\draw[pedge] (1.235,-0.401) -- (0.944,-0.496) -- (0.763,-0.248);
\draw[pedge] (0.763,-0.555) -- (0.944,-0.803) -- (0.763,-1.051);
\draw[pedge] (1.235,-0.401) -- (0.944,-0.496) -- (0.944,-0.803);
\draw[pedge] (0.763,-0.555) -- (0.944,-0.803) -- (0.944,-0.496);
\draw[pedge] (2.470,-1.299) -- (2.290,-1.051) -- (2.470,-0.803);
\draw[pedge] (1.998,-1.452) -- (2.290,-1.357) -- (2.290,-1.051);
\draw[pedge] (2.470,-1.299) -- (2.290,-1.051) -- (2.290,-1.357);
\draw[pedge] (2.762,-0.897) -- (2.942,-1.145) -- (3.234,-1.051);
\draw[pedge] (2.470,-1.299) -- (2.650,-1.051) -- (2.942,-1.145);
\draw[pedge] (2.762,-0.897) -- (2.942,-1.145) -- (2.650,-1.051);
\draw[pedge] (2.762,-0.897) -- (2.470,-0.803) -- (2.650,-1.051);
\draw[pedge] (2.470,-1.299) -- (2.650,-1.051) -- (2.470,-0.803);
\draw[pedge] (1.707,-0.555) -- (1.527,-0.803) -- (1.707,-1.051);
\draw[pedge] (1.235,-0.401) -- (1.527,-0.496) -- (1.527,-0.803);
\draw[pedge] (1.707,-0.555) -- (1.527,-0.803) -- (1.527,-0.496);
\draw[pedge] (1.998,-0.956) -- (2.179,-0.708) -- (2.470,-0.803);
\draw[pedge] (1.707,-0.555) -- (1.887,-0.803) -- (2.179,-0.708);
\draw[pedge] (1.998,-0.956) -- (2.179,-0.708) -- (1.887,-0.803);
\draw[pedge] (1.998,-0.956) -- (1.707,-1.051) -- (1.887,-0.803);
\draw[pedge] (1.707,-0.555) -- (1.887,-0.803) -- (1.707,-1.051);
\draw[pedge] (1.998,-0.956) -- (1.707,-1.051) -- (1.998,-1.145);
\draw[pedge] (1.998,-1.452) -- (1.998,-1.145) -- (1.707,-1.051);
\draw[pedge] (1.998,-0.956) -- (2.290,-1.051) -- (2.470,-0.803);
\draw[pedge] (1.998,-1.452) -- (1.998,-1.145) -- (2.290,-1.051);
\draw[pedge] (1.998,-0.956) -- (2.290,-1.051) -- (1.998,-1.145);
\draw[pedge] (1.527,-1.299) -- (1.707,-1.051) -- (1.707,-1.357);
\draw[pedge] (1.998,-1.452) -- (1.707,-1.357) -- (1.707,-1.051);
\draw[pedge] (1.527,-1.299) -- (1.527,-1.605) -- (1.235,-1.700);
\draw[pedge] (1.998,-1.452) -- (1.707,-1.357) -- (1.527,-1.605);
\draw[pedge] (1.527,-1.299) -- (1.527,-1.605) -- (1.707,-1.357);
\draw[pedge] (1.235,-0.897) -- (1.527,-0.803) -- (1.707,-1.051);
\draw[pedge] (1.235,-0.401) -- (1.235,-0.708) -- (1.527,-0.803);
\draw[pedge] (1.235,-0.897) -- (1.527,-0.803) -- (1.235,-0.708);
\draw[pedge] (1.527,-1.299) -- (1.235,-1.393) -- (1.235,-1.700);
\draw[pedge] (1.235,-0.897) -- (1.415,-1.145) -- (1.235,-1.393);
\draw[pedge] (1.527,-1.299) -- (1.235,-1.393) -- (1.415,-1.145);
\draw[pedge] (1.527,-1.299) -- (1.707,-1.051) -- (1.415,-1.145);
\draw[pedge] (1.235,-0.897) -- (1.415,-1.145) -- (1.707,-1.051);

\end{tikzpicture}
\caption{}
\end{subfigure}
\caption{\textit{а}~— молекула C$_{60}$: усечённый икосаэдр (пятиугольники
закрашены; расчёт); \textit{б}~— мозаика Пенроуза из толстых (жёлтых,
$72^\circ/108^\circ$) и тонких (синих, $36^\circ/144^\circ$) ромбов (расчёт
подразбиением). Мозаика имеет ось 5-го порядка, но не периодична:
отношение числа толстых и тонких ромбов~— иррациональное число
$\tau=(1+\sqrt5)/2$}
\label{fig:c60}\label{fig:penrose}
\end{figure}

\paragraph{Фуллериты и нанотрубки}
Кристаллы из молекул C$_{60}$, C$_{70}$~— \term{фуллериты}. Фуллерит C$_{60}$ при
комнатной температуре имеет ГЦК-решётку с $a=1{,}417$~нм; радиусы
октаэдрических и тетраэдрических пустот 0,207 и 0,112~нм (в них можно
внедрять атомы щелочных металлов; K$_3$C$_{60}$~— сверхпроводник). Фуллерит
C$_{70}$ имеет ГЦК-модификацию с $a=1{,}501$~нм и ГПУ с $a\approx1{,}06$~нм,
$c\approx1{,}72$~нм\ispr{В оригинале для ГПУ-модификации C$_{70}$ указано
$a=1{,}501$~нм; при $c=1{,}718$~нм это дало бы $c/a=1{,}14$, что
невозможно для плотной упаковки молекул. Расстояние между соседними
молекулами в ГЦК-фазе $1{,}501/\sqrt2\approx1{,}06$~нм, и это же значение~—
период $a$ ГПУ-фазы ($c/a\approx1{,}63$).}. Графитовый слой, свёрнутый в цилиндр,~—
\term{углеродная нанотрубка}: диаметр~— единицы нанометров, длина~— до
десятков микрометров; высокая прочность и теплопроводность, большая
допустимая плотность тока.

\paragraph{Квазикристаллы}
В 1984~г. в быстро закалённом (${\sim}10^6$~К/с) сплаве Al--(14--20)~ат.\,\%~Mn
Д.~Шехтман обнаружил дифракционную картину с осью 5-го порядка,
запрещённой для кристаллов. Сейчас известны сотни \term{квазикристаллов},
есть и с осями 8, 10, 12-го порядков. В них нет трансляционной
симметрии, но порядок дальний, «по закону», а не хаотический, как у
аморфных тел,~— поэтому дифракционные максимумы резкие. Модели строят
\emph{проекционным методом}: берут решётку в шестимерном пространстве и
проецируют её узлы, близкие к иррационально ориентированному трёхмерному
подпространству. Двумерная модель~— \term{мозаика Пенроуза}
(рис.~\ref{fig:penrose}): плоскость без пустот и наложений покрывается
ромбами двух сортов~— тонкими ($36^\circ$ и $144^\circ$) и толстыми ($72^\circ$ и
$108^\circ$); отношение их чисел в бесконечной мозаике равно золотому сечению
$\tau=(\sqrt5+1)/2=1{,}618\ldots$, и, поскольку оно иррационально, мозаика не
является двумерным кристаллом.

Свойства квазикристаллов необычны. Они состоят из металлов, но их
электросопротивление при низких температурах аномально велико,
\emph{уменьшается} с ростом температуры и \emph{растёт} с совершенством
структуры (после отжига дефектов). Модули упругости ниже, чем у
кристаллических фаз близкого состава (ближе к аморфным металлам), но
движение дислокаций сильно затруднено, поэтому квазикристаллы хрупки и
интересны как упрочняющие частицы в сплавах.

\subsection{Структура поверхности}\label{ssec:surface}

Поверхность определяет коррозию, трение, износ, адгезию, хрупкость.
Мысленно получим чистую поверхность, рассекая монокристалл по плоскости
$(hkl)$ и удаляя верхнюю часть. Атомы поверхности лишились части связей,
и в приповерхностном слое идут перестройки:
\begin{itemize}
\item \term{релаксация}~— изменение межплоскостного расстояния $d$ в
нескольких верхних слоях по сравнению с объёмным $d_0$ (без изменения
периодичности в плоскости);
\item \term{реконструкция}~— перестройка поверхностных атомов (часто с
участием адсорбированных атомов) с образованием \emph{новой
периодичности} в плоскости поверхности.
\end{itemize}
Структуру поверхности записывают в виде (обозначения Вуда)
\begin{equation}\label{eq:wood}
M(hkl)\text{--}(n\times m)\text{--}C ,
\end{equation}
где $M$~— элемент подложки, $(hkl)$~— грань, $n$ и $m$~— во сколько раз новый
период больше исходного вдоль $\vb a$ и $\vb b$, $C$~— адсорбированный элемент;
приставка «c» означает центрированную ячейку, «R$\alpha$»~— поворот ячейки
на угол $\alpha$ (рис.~\ref{fig:surface}).

\begin{figure}[ht]
\centering
\foreach \typ/\lab in {1/{$(2\times1)$}, 2/{$(2\times2)$}, 3/{$c(2\times2)=(\sqrt2\times\sqrt2)R45^\circ$}}{%
\begin{tikzpicture}[scale=0.55]
  \foreach \i in {0,...,5}{\foreach \j in {0,...,5}{
    \ifnum\typ=1
      \pgfmathsetmacro{\sh}{mod(\i,2)==0 ? 0.18 : -0.18}
      \fill[atomA, draw=atomLine, line width=0.3pt] (\i+\sh,\j) circle (0.3);
    \else
      \fill[atomA, draw=atomLine, line width=0.3pt] (\i,\j) circle (0.3);
    \fi}}
  \ifnum\typ=2
    \foreach \i in {0,2,4}{\foreach \j in {0,2,4}{\fill[fmRed] (\i+0.5,\j+0.5) circle (0.17);}}
    \draw[fmBlue, thick] (0.5,0.5) rectangle (2.5,2.5);
  \fi
  \ifnum\typ=3
    \foreach \i in {0,...,4}{\foreach \j in {0,...,4}{
      \pgfmathtruncatemacro{\pp}{mod(\i+\j,2)}
      \ifnum\pp=0 \fill[fmRed] (\i+0.5,\j+0.5) circle (0.17);\fi}}
    \draw[fmBlue, thick] (0.5,0.5) rectangle (2.5,2.5);
    \draw[fmGreen!60!black, thick, dashed] (1.5,0.5) -- (2.5,1.5) -- (1.5,2.5) -- (0.5,1.5) -- cycle;
  \fi
  \ifnum\typ=1
    \draw[fmBlue, thick] (-0.5,-0.5) rectangle (1.5,0.5);
  \fi
  \node[font=\small] at (2.5,-1.2) {\lab};
\end{tikzpicture}\hspace{5mm}}
\caption{Реконструкция поверхности (схема): \textit{а}~— $(2\times1)$: атомы
сближаются попарно (димеры), период вдоль одной оси удваивается;
\textit{б}~— $(2\times2)$: адсорбированные атомы (красные) в каждой второй
позиции вдоль обеих осей; \textit{в}~— $c(2\times2)$: адатомы в шахматном
порядке; центрированная ячейка (синяя) эквивалентна ячейке
$(\sqrt2\times\sqrt2)$, повёрнутой на $45^\circ$ (зелёная)}
\label{fig:surface}
\end{figure}

Примеры: Ni(111)--$(2\times2)$--O (кислород на никеле); W(110)--c$(9\times5)$--CO;
чистая грань Si(100) имеет структуру $(2\times1)$ (атомы образуют димеры), а
адсорбция водорода восстанавливает исходную периодичность: Si(100)--$(1\times1)$--H.
Чистые грани Au реконструированы: Au(100)--$(5\times20)$, Au(110)--$(1\times2)$
(«пропущенные ряды»), Au(111)--$(22\times\sqrt3)$. У ОЦК-металлов при низких
температурах реконструирована грань (100) W и Mo~— $c(2\times2)$. Релаксация
обычно сводится к \emph{сжатию} верхнего межплоскостного расстояния, тем
большему, чем «рыхлее» грань: у Cu(110) около $-10\,\%$, Al(110) $-5\ldots-10\,\%$,
W(100) $-5\ldots-10\,\%$, Mo(100) около $-11\,\%$, Fe(111) около $-15\,\%$; у плотных граней
релаксация мала (Ni(111) около $-1\,\%$; Al(111) даже слегка расширяется).
Структуру поверхности изучают дифракцией медленных электронов
(энергия 10--200~эВ: такие электроны проникают лишь на несколько
атомных слоёв).

\begin{itogi}
\begin{itemize}
\item Две плотнейшие упаковки~— ГПУ ($ABAB$) и ГЦК ($ABCABC$), компактность
74\,\%, к.\,ч.~=~12; пустот октаэдрических~— по одной, тетраэдрических~— по
две на атом ($0{,}414R$ и $0{,}225R$).
\item ОЦК: к.\,ч.~=~8(+6), компактность 68\,\%; тетрапустоты ($0{,}291R$) крупнее
октапустот ($0{,}155R$).
\item Многие соединения~— плотная упаковка одних атомов с заполнением
пустот другими: NaCl (все октапустоты ГЦК), сфалерит (половина тетрапустот
ГЦК), вюрцит (половина тетрапустот ГПУ), NiAs (все октапустоты ГПУ),
флюорит (все тетрапустоты ГЦК).
\item Сверхструктуры: L1$_2$ (Cu$_3$Au), L1$_0$ (CuAu), B2 (CuZn), D0$_3$
(Fe$_3$Al), L2$_1$ (Cu$_2$MnAl).
\item Квазикристаллы~— дальний порядок без трансляций (оси 5, 8, 10, 12);
поверхность релаксирует и реконструируется.
\end{itemize}
\end{itogi}

\begin{proverka}
\item Почему коэффициент компактности ГЦК и ГПУ одинаков?
\item Сколько октаэдрических и тетраэдрических пустот приходится на атом в
ГЦК- и ОЦК-решётках?
\item Почему углерод в аустените занимает октаэдрические пустоты, а в
феррите~— сплюснутые октаэдрические, хотя тетраэдрические там крупнее?
\item Какую часть тетраэдрических пустот занимают атомы в структуре алмаза
и почему её к.\,ч. равно 4?
\item Чем отличаются структуры сфалерита и вюрцита?
\item Почему CsCl не имеет ОЦК-решётку, хотя атомы расположены по узлам ОЦК?
\item Сколько атомов Fe и Al в ячейке D0$_3$ на рис.~\ref{fig:superstructures}\,\textit{г}?
\item Что означает запись Si(111)--$(7\times7)$?
\end{proverka}


\section{Физические свойства кристаллов}\label{sec:properties}
% ----------------------------------------------------------------------
%  1.5. Физические свойства кристаллов
% ----------------------------------------------------------------------
\begin{orientir}
Кристаллофизика изучает, как симметрия кристалла ограничивает его
свойства. Два главных инструмента: \vazhno{принципы симметрии Неймана и
Кюри} (1.5.1) и \vazhno{тензорное описание свойств} (1.5.2). Применив их,
узнаем, сколько независимых констант нужно для описания проводимости
(1.5.3) и упругости (1.5.4) кристалла каждой сингонии, а в 1.5.5~—
как это число находится общим методом теории представлений групп.
\end{orientir}

Из-за однородности, анизотропии и симметрии в кристаллах возможны
явления, невозможные в изотропных телах (пьезоэффект, двупреломление,
анизотропия упругости и проводимости). Математический аппарат
кристаллофизики~— тензорное исчисление и теория групп.

\subsection{Принципы симметрии в кристаллофизике}\label{ssec:symmetry-principle}

Для макроскопических свойств кристалл~— однородная среда, поэтому
различия между простыми и винтовыми осями, зеркальными плоскостями и
плоскостями скольжения несущественны. Симметрию макроскопических
свойств определяет не пространственная, а \emph{точечная} группа
кристалла.

\begin{glavnoe}[Принцип Неймана]
Элементы симметрии любого физического свойства кристалла включают
элементы симметрии его точечной группы:
\begin{equation}\label{eq:neumann}
G_\text{свойства}\supseteq G_\text{кристалла} .
\end{equation}
Симметрия свойства \emph{не ниже} симметрии кристалла (может быть и
выше). Точечная группа кристалла~— общая часть (пересечение) групп
симметрии всех его свойств.
\end{glavnoe}

Например, электропроводность кубического кристалла изотропна (симметрия
шара $\infty\infty m$~— намного выше $m\br3m$), а модуль Юнга того же
кристалла анизотропен (рис.~\ref{fig:young-polar}) и имеет ровно
кубическую симметрию.

\begin{glavnoe}[Принцип суперпозиции Кюри]
Кристалл во внешнем поле сохраняет только те элементы симметрии, которые
общи для кристалла без воздействия и для воздействия без кристалла (при
их заданной взаимной ориентации):
\begin{equation}\label{eq:curie}
G_\text{кр+возд}=G_\text{кр}\cap G_\text{возд} .
\end{equation}
\end{glavnoe}

Симметрию воздействий задают предельные группы (п.~\ref{ssec:curie}), а
искать общие элементы удобно по стереографическим проекциям
(рис.~\ref{fig:curie-principle}).
\begin{itemize}
\item \emph{Одноосное растяжение} имеет симметрию $\infty/mmm$
(неподвижный цилиндр). Вдоль оси 4 кристалла $m\br3m$ сохраняются: одна ось
4, четыре поперечные оси 2, четыре продольные и одна поперечная
плоскости, центр~— кристалл становится \emph{тетрагональным} ($4/mmm$).
Вдоль оси 3 остаются ось 3, три поперечные оси 2, три продольные
плоскости и центр~— класс $\br3m$ (\emph{тригональный}).
\item \emph{Электрическое поле}~— полярный вектор, симметрия $\infty m$
(неподвижный конус). Поле вдоль оси 4 кристалла $m\br3m$ даёт класс $4mm$,
вдоль оси 3~— $3m$; центр симметрии исчезает, и кристалл становится
полярным.
\end{itemize}
Обратная задача так не решается: по симметрии результата нельзя судить о
симметрии причин. Обобщённая формулировка Кюри: при наложении явлений
их \emph{диссимметрии} (по Шубникову~— отсутствующие элементы симметрии)
складываются.

\begin{figure}[ht]
\centering
\tikzset{mirror/.style={fmBlue, line width=0.6pt},
  ax4/.style={rectangle, draw=fmRed, fill=fmRed!35, inner sep=0pt, minimum size=5pt},
  ax3/.style={regular polygon, regular polygon sides=3, draw=fmAmber!90!black, fill=fmAmber!40, inner sep=0pt, minimum size=7pt},
  ax2/.style={ellipse, draw=fmGreen!60!black, fill=fmGreen!35, inner sep=0pt, minimum width=7pt, minimum height=4pt}}
\begin{subfigure}[b]{0.32\textwidth}\centering
\begin{tikzpicture}[scale=1.55]
  \draw[mirror, line width=1pt] (0,0) circle (1);
  % создано tools/gen_stereo.py: элементы симметрии m-3m на проекции (001)
\draw[mirror] (0.0000,-1.0000) -- (0.0000,-0.9742) -- (0.0000,-0.9490) -- (0.0000,-0.9244) -- (0.0000,-0.9004) -- (0.0000,-0.8770) -- (0.0000,-0.8541) -- (0.0000,-0.8317) -- (0.0000,-0.8098) -- (0.0000,-0.7883) -- (0.0000,-0.7673) -- (0.0000,-0.7467) -- (0.0000,-0.7265) -- (0.0000,-0.7067) -- (0.0000,-0.6873) -- (0.0000,-0.6682) -- (0.0000,-0.6494) -- (0.0000,-0.6310) -- (0.0000,-0.6128) -- (0.0000,-0.5949) -- (0.0000,-0.5774) -- (0.0000,-0.5600) -- (0.0000,-0.5430) -- (0.0000,-0.5261) -- (0.0000,-0.5095) -- (0.0000,-0.4931) -- (0.0000,-0.4770) -- (0.0000,-0.4610) -- (0.0000,-0.4452) -- (0.0000,-0.4296) -- (0.0000,-0.4142) -- (0.0000,-0.3990) -- (0.0000,-0.3839) -- (0.0000,-0.3689) -- (0.0000,-0.3541) -- (0.0000,-0.3395) -- (0.0000,-0.3249) -- (0.0000,-0.3105) -- (0.0000,-0.2962) -- (0.0000,-0.2820) -- (0.0000,-0.2679) -- (0.0000,-0.2540) -- (0.0000,-0.2401) -- (0.0000,-0.2263) -- (0.0000,-0.2126) -- (0.0000,-0.1989) -- (0.0000,-0.1853) -- (0.0000,-0.1718) -- (0.0000,-0.1584) -- (0.0000,-0.1450) -- (0.0000,-0.1317) -- (0.0000,-0.1184) -- (0.0000,-0.1051) -- (0.0000,-0.0919) -- (0.0000,-0.0787) -- (0.0000,-0.0655) -- (0.0000,-0.0524) -- (0.0000,-0.0393) -- (0.0000,-0.0262) -- (0.0000,-0.0131) -- (0.0000,-0.0000) -- (-0.0000,0.0131) -- (-0.0000,0.0262) -- (-0.0000,0.0393) -- (-0.0000,0.0524) -- (-0.0000,0.0655) -- (-0.0000,0.0787) -- (-0.0000,0.0919) -- (-0.0000,0.1051) -- (-0.0000,0.1184) -- (-0.0000,0.1317) -- (-0.0000,0.1450) -- (-0.0000,0.1584) -- (-0.0000,0.1718) -- (-0.0000,0.1853) -- (-0.0000,0.1989) -- (-0.0000,0.2126) -- (-0.0000,0.2263) -- (-0.0000,0.2401) -- (-0.0000,0.2540) -- (-0.0000,0.2679) -- (-0.0000,0.2820) -- (-0.0000,0.2962) -- (-0.0000,0.3105) -- (-0.0000,0.3249) -- (-0.0000,0.3395) -- (-0.0000,0.3541) -- (-0.0000,0.3689) -- (-0.0000,0.3839) -- (-0.0000,0.3990) -- (-0.0000,0.4142) -- (-0.0000,0.4296) -- (-0.0000,0.4452) -- (-0.0000,0.4610) -- (-0.0000,0.4770) -- (-0.0000,0.4931) -- (-0.0000,0.5095) -- (-0.0000,0.5261) -- (-0.0000,0.5430) -- (-0.0000,0.5600) -- (-0.0000,0.5774) -- (-0.0000,0.5949) -- (-0.0000,0.6128) -- (-0.0000,0.6310) -- (-0.0000,0.6494) -- (-0.0000,0.6682) -- (-0.0000,0.6873) -- (-0.0000,0.7067) -- (-0.0000,0.7265) -- (-0.0000,0.7467) -- (-0.0000,0.7673) -- (-0.0000,0.7883) -- (-0.0000,0.8098) -- (-0.0000,0.8317) -- (-0.0000,0.8541) -- (-0.0000,0.8770) -- (-0.0000,0.9004) -- (-0.0000,0.9244) -- (-0.0000,0.9490) -- (-0.0000,0.9742) -- (-0.0000,1.0000);
\draw[mirror] (1.0000,0.0000) -- (0.9742,0.0000) -- (0.9490,0.0000) -- (0.9244,0.0000) -- (0.9004,0.0000) -- (0.8770,0.0000) -- (0.8541,0.0000) -- (0.8317,0.0000) -- (0.8098,0.0000) -- (0.7883,0.0000) -- (0.7673,0.0000) -- (0.7467,0.0000) -- (0.7265,0.0000) -- (0.7067,0.0000) -- (0.6873,0.0000) -- (0.6682,0.0000) -- (0.6494,0.0000) -- (0.6310,0.0000) -- (0.6128,0.0000) -- (0.5949,0.0000) -- (0.5774,0.0000) -- (0.5600,0.0000) -- (0.5430,0.0000) -- (0.5261,0.0000) -- (0.5095,0.0000) -- (0.4931,0.0000) -- (0.4770,0.0000) -- (0.4610,0.0000) -- (0.4452,0.0000) -- (0.4296,0.0000) -- (0.4142,0.0000) -- (0.3990,0.0000) -- (0.3839,0.0000) -- (0.3689,0.0000) -- (0.3541,0.0000) -- (0.3395,0.0000) -- (0.3249,0.0000) -- (0.3105,0.0000) -- (0.2962,0.0000) -- (0.2820,0.0000) -- (0.2679,0.0000) -- (0.2540,0.0000) -- (0.2401,0.0000) -- (0.2263,0.0000) -- (0.2126,0.0000) -- (0.1989,0.0000) -- (0.1853,0.0000) -- (0.1718,0.0000) -- (0.1584,0.0000) -- (0.1450,0.0000) -- (0.1317,0.0000) -- (0.1184,0.0000) -- (0.1051,0.0000) -- (0.0919,0.0000) -- (0.0787,0.0000) -- (0.0655,0.0000) -- (0.0524,0.0000) -- (0.0393,0.0000) -- (0.0262,0.0000) -- (0.0131,0.0000) -- (0.0000,0.0000) -- (-0.0131,-0.0000) -- (-0.0262,-0.0000) -- (-0.0393,-0.0000) -- (-0.0524,-0.0000) -- (-0.0655,-0.0000) -- (-0.0787,-0.0000) -- (-0.0919,-0.0000) -- (-0.1051,-0.0000) -- (-0.1184,-0.0000) -- (-0.1317,-0.0000) -- (-0.1450,-0.0000) -- (-0.1584,-0.0000) -- (-0.1718,-0.0000) -- (-0.1853,-0.0000) -- (-0.1989,-0.0000) -- (-0.2126,-0.0000) -- (-0.2263,-0.0000) -- (-0.2401,-0.0000) -- (-0.2540,-0.0000) -- (-0.2679,-0.0000) -- (-0.2820,-0.0000) -- (-0.2962,-0.0000) -- (-0.3105,-0.0000) -- (-0.3249,-0.0000) -- (-0.3395,-0.0000) -- (-0.3541,-0.0000) -- (-0.3689,-0.0000) -- (-0.3839,-0.0000) -- (-0.3990,-0.0000) -- (-0.4142,-0.0000) -- (-0.4296,-0.0000) -- (-0.4452,-0.0000) -- (-0.4610,-0.0000) -- (-0.4770,-0.0000) -- (-0.4931,-0.0000) -- (-0.5095,-0.0000) -- (-0.5261,-0.0000) -- (-0.5430,-0.0000) -- (-0.5600,-0.0000) -- (-0.5774,-0.0000) -- (-0.5949,-0.0000) -- (-0.6128,-0.0000) -- (-0.6310,-0.0000) -- (-0.6494,-0.0000) -- (-0.6682,-0.0000) -- (-0.6873,-0.0000) -- (-0.7067,-0.0000) -- (-0.7265,-0.0000) -- (-0.7467,-0.0000) -- (-0.7673,-0.0000) -- (-0.7883,-0.0000) -- (-0.8098,-0.0000) -- (-0.8317,-0.0000) -- (-0.8541,-0.0000) -- (-0.8770,-0.0000) -- (-0.9004,-0.0000) -- (-0.9244,-0.0000) -- (-0.9490,-0.0000) -- (-0.9742,-0.0000) -- (-1.0000,-0.0000);
\draw[mirror] (0.7071,-0.7071) -- (0.6888,-0.6888) -- (0.6710,-0.6710) -- (0.6536,-0.6536) -- (0.6367,-0.6367) -- (0.6201,-0.6201) -- (0.6039,-0.6039) -- (0.5881,-0.5881) -- (0.5726,-0.5726) -- (0.5574,-0.5574) -- (0.5426,-0.5426) -- (0.5280,-0.5280) -- (0.5137,-0.5137) -- (0.4997,-0.4997) -- (0.4860,-0.4860) -- (0.4725,-0.4725) -- (0.4592,-0.4592) -- (0.4462,-0.4462) -- (0.4333,-0.4333) -- (0.4207,-0.4207) -- (0.4082,-0.4082) -- (0.3960,-0.3960) -- (0.3839,-0.3839) -- (0.3720,-0.3720) -- (0.3603,-0.3603) -- (0.3487,-0.3487) -- (0.3373,-0.3373) -- (0.3260,-0.3260) -- (0.3148,-0.3148) -- (0.3038,-0.3038) -- (0.2929,-0.2929) -- (0.2821,-0.2821) -- (0.2714,-0.2714) -- (0.2609,-0.2609) -- (0.2504,-0.2504) -- (0.2400,-0.2400) -- (0.2298,-0.2298) -- (0.2196,-0.2196) -- (0.2095,-0.2095) -- (0.1994,-0.1994) -- (0.1895,-0.1895) -- (0.1796,-0.1796) -- (0.1698,-0.1698) -- (0.1600,-0.1600) -- (0.1503,-0.1503) -- (0.1407,-0.1407) -- (0.1311,-0.1311) -- (0.1215,-0.1215) -- (0.1120,-0.1120) -- (0.1025,-0.1025) -- (0.0931,-0.0931) -- (0.0837,-0.0837) -- (0.0743,-0.0743) -- (0.0650,-0.0650) -- (0.0557,-0.0557) -- (0.0463,-0.0463) -- (0.0371,-0.0371) -- (0.0278,-0.0278) -- (0.0185,-0.0185) -- (0.0093,-0.0093) -- (0.0000,-0.0000) -- (-0.0093,0.0093) -- (-0.0185,0.0185) -- (-0.0278,0.0278) -- (-0.0371,0.0371) -- (-0.0463,0.0463) -- (-0.0557,0.0557) -- (-0.0650,0.0650) -- (-0.0743,0.0743) -- (-0.0837,0.0837) -- (-0.0931,0.0931) -- (-0.1025,0.1025) -- (-0.1120,0.1120) -- (-0.1215,0.1215) -- (-0.1311,0.1311) -- (-0.1407,0.1407) -- (-0.1503,0.1503) -- (-0.1600,0.1600) -- (-0.1698,0.1698) -- (-0.1796,0.1796) -- (-0.1895,0.1895) -- (-0.1994,0.1994) -- (-0.2095,0.2095) -- (-0.2196,0.2196) -- (-0.2298,0.2298) -- (-0.2400,0.2400) -- (-0.2504,0.2504) -- (-0.2609,0.2609) -- (-0.2714,0.2714) -- (-0.2821,0.2821) -- (-0.2929,0.2929) -- (-0.3038,0.3038) -- (-0.3148,0.3148) -- (-0.3260,0.3260) -- (-0.3373,0.3373) -- (-0.3487,0.3487) -- (-0.3603,0.3603) -- (-0.3720,0.3720) -- (-0.3839,0.3839) -- (-0.3960,0.3960) -- (-0.4082,0.4082) -- (-0.4207,0.4207) -- (-0.4333,0.4333) -- (-0.4462,0.4462) -- (-0.4592,0.4592) -- (-0.4725,0.4725) -- (-0.4860,0.4860) -- (-0.4997,0.4997) -- (-0.5137,0.5137) -- (-0.5280,0.5280) -- (-0.5426,0.5426) -- (-0.5574,0.5574) -- (-0.5726,0.5726) -- (-0.5881,0.5881) -- (-0.6039,0.6039) -- (-0.6201,0.6201) -- (-0.6367,0.6367) -- (-0.6536,0.6536) -- (-0.6710,0.6710) -- (-0.6888,0.6888) -- (-0.7071,0.7071);
\draw[mirror] (-0.7071,-0.7071) -- (-0.6888,-0.6888) -- (-0.6710,-0.6710) -- (-0.6536,-0.6536) -- (-0.6367,-0.6367) -- (-0.6201,-0.6201) -- (-0.6039,-0.6039) -- (-0.5881,-0.5881) -- (-0.5726,-0.5726) -- (-0.5574,-0.5574) -- (-0.5426,-0.5426) -- (-0.5280,-0.5280) -- (-0.5137,-0.5137) -- (-0.4997,-0.4997) -- (-0.4860,-0.4860) -- (-0.4725,-0.4725) -- (-0.4592,-0.4592) -- (-0.4462,-0.4462) -- (-0.4333,-0.4333) -- (-0.4207,-0.4207) -- (-0.4082,-0.4082) -- (-0.3960,-0.3960) -- (-0.3839,-0.3839) -- (-0.3720,-0.3720) -- (-0.3603,-0.3603) -- (-0.3487,-0.3487) -- (-0.3373,-0.3373) -- (-0.3260,-0.3260) -- (-0.3148,-0.3148) -- (-0.3038,-0.3038) -- (-0.2929,-0.2929) -- (-0.2821,-0.2821) -- (-0.2714,-0.2714) -- (-0.2609,-0.2609) -- (-0.2504,-0.2504) -- (-0.2400,-0.2400) -- (-0.2298,-0.2298) -- (-0.2196,-0.2196) -- (-0.2095,-0.2095) -- (-0.1994,-0.1994) -- (-0.1895,-0.1895) -- (-0.1796,-0.1796) -- (-0.1698,-0.1698) -- (-0.1600,-0.1600) -- (-0.1503,-0.1503) -- (-0.1407,-0.1407) -- (-0.1311,-0.1311) -- (-0.1215,-0.1215) -- (-0.1120,-0.1120) -- (-0.1025,-0.1025) -- (-0.0931,-0.0931) -- (-0.0837,-0.0837) -- (-0.0743,-0.0743) -- (-0.0650,-0.0650) -- (-0.0557,-0.0557) -- (-0.0463,-0.0463) -- (-0.0371,-0.0371) -- (-0.0278,-0.0278) -- (-0.0185,-0.0185) -- (-0.0093,-0.0093) -- (-0.0000,-0.0000) -- (0.0093,0.0093) -- (0.0185,0.0185) -- (0.0278,0.0278) -- (0.0371,0.0371) -- (0.0463,0.0463) -- (0.0557,0.0557) -- (0.0650,0.0650) -- (0.0743,0.0743) -- (0.0837,0.0837) -- (0.0931,0.0931) -- (0.1025,0.1025) -- (0.1120,0.1120) -- (0.1215,0.1215) -- (0.1311,0.1311) -- (0.1407,0.1407) -- (0.1503,0.1503) -- (0.1600,0.1600) -- (0.1698,0.1698) -- (0.1796,0.1796) -- (0.1895,0.1895) -- (0.1994,0.1994) -- (0.2095,0.2095) -- (0.2196,0.2196) -- (0.2298,0.2298) -- (0.2400,0.2400) -- (0.2504,0.2504) -- (0.2609,0.2609) -- (0.2714,0.2714) -- (0.2821,0.2821) -- (0.2929,0.2929) -- (0.3038,0.3038) -- (0.3148,0.3148) -- (0.3260,0.3260) -- (0.3373,0.3373) -- (0.3487,0.3487) -- (0.3603,0.3603) -- (0.3720,0.3720) -- (0.3839,0.3839) -- (0.3960,0.3960) -- (0.4082,0.4082) -- (0.4207,0.4207) -- (0.4333,0.4333) -- (0.4462,0.4462) -- (0.4592,0.4592) -- (0.4725,0.4725) -- (0.4860,0.4860) -- (0.4997,0.4997) -- (0.5137,0.5137) -- (0.5280,0.5280) -- (0.5426,0.5426) -- (0.5574,0.5574) -- (0.5726,0.5726) -- (0.5881,0.5881) -- (0.6039,0.6039) -- (0.6201,0.6201) -- (0.6367,0.6367) -- (0.6536,0.6536) -- (0.6710,0.6710) -- (0.6888,0.6888) -- (0.7071,0.7071);
\draw[mirror] (0.0000,-1.0000) -- (-0.0182,-0.9815) -- (-0.0357,-0.9630) -- (-0.0526,-0.9445) -- (-0.0688,-0.9261) -- (-0.0845,-0.9077) -- (-0.0996,-0.8893) -- (-0.1142,-0.8710) -- (-0.1282,-0.8528) -- (-0.1417,-0.8346) -- (-0.1547,-0.8165) -- (-0.1672,-0.7985) -- (-0.1793,-0.7805) -- (-0.1910,-0.7626) -- (-0.2022,-0.7448) -- (-0.2130,-0.7271) -- (-0.2234,-0.7095) -- (-0.2334,-0.6919) -- (-0.2430,-0.6745) -- (-0.2523,-0.6571) -- (-0.2612,-0.6398) -- (-0.2698,-0.6226) -- (-0.2780,-0.6055) -- (-0.2860,-0.5884) -- (-0.2936,-0.5715) -- (-0.3009,-0.5546) -- (-0.3080,-0.5378) -- (-0.3147,-0.5211) -- (-0.3212,-0.5045) -- (-0.3274,-0.4879) -- (-0.3333,-0.4714) -- (-0.3390,-0.4550) -- (-0.3445,-0.4386) -- (-0.3497,-0.4224) -- (-0.3546,-0.4061) -- (-0.3594,-0.3900) -- (-0.3639,-0.3739) -- (-0.3682,-0.3579) -- (-0.3723,-0.3419) -- (-0.3761,-0.3260) -- (-0.3798,-0.3101) -- (-0.3833,-0.2943) -- (-0.3865,-0.2785) -- (-0.3896,-0.2628) -- (-0.3925,-0.2471) -- (-0.3951,-0.2315) -- (-0.3976,-0.2159) -- (-0.4000,-0.2003) -- (-0.4021,-0.1848) -- (-0.4040,-0.1693) -- (-0.4058,-0.1538) -- (-0.4074,-0.1383) -- (-0.4089,-0.1229) -- (-0.4101,-0.1075) -- (-0.4112,-0.0921) -- (-0.4121,-0.0767) -- (-0.4129,-0.0614) -- (-0.4135,-0.0460) -- (-0.4139,-0.0307) -- (-0.4141,-0.0153) -- (-0.4142,-0.0000) -- (-0.4141,0.0153) -- (-0.4139,0.0307) -- (-0.4135,0.0460) -- (-0.4129,0.0614) -- (-0.4121,0.0767) -- (-0.4112,0.0921) -- (-0.4101,0.1075) -- (-0.4089,0.1229) -- (-0.4074,0.1383) -- (-0.4058,0.1538) -- (-0.4040,0.1693) -- (-0.4021,0.1848) -- (-0.4000,0.2003) -- (-0.3976,0.2159) -- (-0.3951,0.2315) -- (-0.3925,0.2471) -- (-0.3896,0.2628) -- (-0.3865,0.2785) -- (-0.3833,0.2943) -- (-0.3798,0.3101) -- (-0.3761,0.3260) -- (-0.3723,0.3419) -- (-0.3682,0.3579) -- (-0.3639,0.3739) -- (-0.3594,0.3900) -- (-0.3546,0.4061) -- (-0.3497,0.4224) -- (-0.3445,0.4386) -- (-0.3390,0.4550) -- (-0.3333,0.4714) -- (-0.3274,0.4879) -- (-0.3212,0.5045) -- (-0.3147,0.5211) -- (-0.3080,0.5378) -- (-0.3009,0.5546) -- (-0.2936,0.5715) -- (-0.2860,0.5884) -- (-0.2780,0.6055) -- (-0.2698,0.6226) -- (-0.2612,0.6398) -- (-0.2523,0.6571) -- (-0.2430,0.6745) -- (-0.2334,0.6919) -- (-0.2234,0.7095) -- (-0.2130,0.7271) -- (-0.2022,0.7448) -- (-0.1910,0.7626) -- (-0.1793,0.7805) -- (-0.1672,0.7985) -- (-0.1547,0.8165) -- (-0.1417,0.8346) -- (-0.1282,0.8528) -- (-0.1142,0.8710) -- (-0.0996,0.8893) -- (-0.0845,0.9077) -- (-0.0688,0.9261) -- (-0.0526,0.9445) -- (-0.0357,0.9630) -- (-0.0182,0.9815) -- (-0.0000,1.0000);
\draw[mirror] (0.0000,-1.0000) -- (0.0182,-0.9815) -- (0.0357,-0.9630) -- (0.0526,-0.9445) -- (0.0688,-0.9261) -- (0.0845,-0.9077) -- (0.0996,-0.8893) -- (0.1142,-0.8710) -- (0.1282,-0.8528) -- (0.1417,-0.8346) -- (0.1547,-0.8165) -- (0.1672,-0.7985) -- (0.1793,-0.7805) -- (0.1910,-0.7626) -- (0.2022,-0.7448) -- (0.2130,-0.7271) -- (0.2234,-0.7095) -- (0.2334,-0.6919) -- (0.2430,-0.6745) -- (0.2523,-0.6571) -- (0.2612,-0.6398) -- (0.2698,-0.6226) -- (0.2780,-0.6055) -- (0.2860,-0.5884) -- (0.2936,-0.5715) -- (0.3009,-0.5546) -- (0.3080,-0.5378) -- (0.3147,-0.5211) -- (0.3212,-0.5045) -- (0.3274,-0.4879) -- (0.3333,-0.4714) -- (0.3390,-0.4550) -- (0.3445,-0.4386) -- (0.3497,-0.4224) -- (0.3546,-0.4061) -- (0.3594,-0.3900) -- (0.3639,-0.3739) -- (0.3682,-0.3579) -- (0.3723,-0.3419) -- (0.3761,-0.3260) -- (0.3798,-0.3101) -- (0.3833,-0.2943) -- (0.3865,-0.2785) -- (0.3896,-0.2628) -- (0.3925,-0.2471) -- (0.3951,-0.2315) -- (0.3976,-0.2159) -- (0.4000,-0.2003) -- (0.4021,-0.1848) -- (0.4040,-0.1693) -- (0.4058,-0.1538) -- (0.4074,-0.1383) -- (0.4089,-0.1229) -- (0.4101,-0.1075) -- (0.4112,-0.0921) -- (0.4121,-0.0767) -- (0.4129,-0.0614) -- (0.4135,-0.0460) -- (0.4139,-0.0307) -- (0.4141,-0.0153) -- (0.4142,-0.0000) -- (0.4141,0.0153) -- (0.4139,0.0307) -- (0.4135,0.0460) -- (0.4129,0.0614) -- (0.4121,0.0767) -- (0.4112,0.0921) -- (0.4101,0.1075) -- (0.4089,0.1229) -- (0.4074,0.1383) -- (0.4058,0.1538) -- (0.4040,0.1693) -- (0.4021,0.1848) -- (0.4000,0.2003) -- (0.3976,0.2159) -- (0.3951,0.2315) -- (0.3925,0.2471) -- (0.3896,0.2628) -- (0.3865,0.2785) -- (0.3833,0.2943) -- (0.3798,0.3101) -- (0.3761,0.3260) -- (0.3723,0.3419) -- (0.3682,0.3579) -- (0.3639,0.3739) -- (0.3594,0.3900) -- (0.3546,0.4061) -- (0.3497,0.4224) -- (0.3445,0.4386) -- (0.3390,0.4550) -- (0.3333,0.4714) -- (0.3274,0.4879) -- (0.3212,0.5045) -- (0.3147,0.5211) -- (0.3080,0.5378) -- (0.3009,0.5546) -- (0.2936,0.5715) -- (0.2860,0.5884) -- (0.2780,0.6055) -- (0.2698,0.6226) -- (0.2612,0.6398) -- (0.2523,0.6571) -- (0.2430,0.6745) -- (0.2334,0.6919) -- (0.2234,0.7095) -- (0.2130,0.7271) -- (0.2022,0.7448) -- (0.1910,0.7626) -- (0.1793,0.7805) -- (0.1672,0.7985) -- (0.1547,0.8165) -- (0.1417,0.8346) -- (0.1282,0.8528) -- (0.1142,0.8710) -- (0.0996,0.8893) -- (0.0845,0.9077) -- (0.0688,0.9261) -- (0.0526,0.9445) -- (0.0357,0.9630) -- (0.0182,0.9815) -- (0.0000,1.0000);
\draw[mirror] (1.0000,0.0000) -- (0.9815,-0.0182) -- (0.9630,-0.0357) -- (0.9445,-0.0526) -- (0.9261,-0.0688) -- (0.9077,-0.0845) -- (0.8893,-0.0996) -- (0.8710,-0.1142) -- (0.8528,-0.1282) -- (0.8346,-0.1417) -- (0.8165,-0.1547) -- (0.7985,-0.1672) -- (0.7805,-0.1793) -- (0.7626,-0.1910) -- (0.7448,-0.2022) -- (0.7271,-0.2130) -- (0.7095,-0.2234) -- (0.6919,-0.2334) -- (0.6745,-0.2430) -- (0.6571,-0.2523) -- (0.6398,-0.2612) -- (0.6226,-0.2698) -- (0.6055,-0.2780) -- (0.5884,-0.2860) -- (0.5715,-0.2936) -- (0.5546,-0.3009) -- (0.5378,-0.3080) -- (0.5211,-0.3147) -- (0.5045,-0.3212) -- (0.4879,-0.3274) -- (0.4714,-0.3333) -- (0.4550,-0.3390) -- (0.4386,-0.3445) -- (0.4224,-0.3497) -- (0.4061,-0.3546) -- (0.3900,-0.3594) -- (0.3739,-0.3639) -- (0.3579,-0.3682) -- (0.3419,-0.3723) -- (0.3260,-0.3761) -- (0.3101,-0.3798) -- (0.2943,-0.3833) -- (0.2785,-0.3865) -- (0.2628,-0.3896) -- (0.2471,-0.3925) -- (0.2315,-0.3951) -- (0.2159,-0.3976) -- (0.2003,-0.4000) -- (0.1848,-0.4021) -- (0.1693,-0.4040) -- (0.1538,-0.4058) -- (0.1383,-0.4074) -- (0.1229,-0.4089) -- (0.1075,-0.4101) -- (0.0921,-0.4112) -- (0.0767,-0.4121) -- (0.0614,-0.4129) -- (0.0460,-0.4135) -- (0.0307,-0.4139) -- (0.0153,-0.4141) -- (0.0000,-0.4142) -- (-0.0153,-0.4141) -- (-0.0307,-0.4139) -- (-0.0460,-0.4135) -- (-0.0614,-0.4129) -- (-0.0767,-0.4121) -- (-0.0921,-0.4112) -- (-0.1075,-0.4101) -- (-0.1229,-0.4089) -- (-0.1383,-0.4074) -- (-0.1538,-0.4058) -- (-0.1693,-0.4040) -- (-0.1848,-0.4021) -- (-0.2003,-0.4000) -- (-0.2159,-0.3976) -- (-0.2315,-0.3951) -- (-0.2471,-0.3925) -- (-0.2628,-0.3896) -- (-0.2785,-0.3865) -- (-0.2943,-0.3833) -- (-0.3101,-0.3798) -- (-0.3260,-0.3761) -- (-0.3419,-0.3723) -- (-0.3579,-0.3682) -- (-0.3739,-0.3639) -- (-0.3900,-0.3594) -- (-0.4061,-0.3546) -- (-0.4224,-0.3497) -- (-0.4386,-0.3445) -- (-0.4550,-0.3390) -- (-0.4714,-0.3333) -- (-0.4879,-0.3274) -- (-0.5045,-0.3212) -- (-0.5211,-0.3147) -- (-0.5378,-0.3080) -- (-0.5546,-0.3009) -- (-0.5715,-0.2936) -- (-0.5884,-0.2860) -- (-0.6055,-0.2780) -- (-0.6226,-0.2698) -- (-0.6398,-0.2612) -- (-0.6571,-0.2523) -- (-0.6745,-0.2430) -- (-0.6919,-0.2334) -- (-0.7095,-0.2234) -- (-0.7271,-0.2130) -- (-0.7448,-0.2022) -- (-0.7626,-0.1910) -- (-0.7805,-0.1793) -- (-0.7985,-0.1672) -- (-0.8165,-0.1547) -- (-0.8346,-0.1417) -- (-0.8528,-0.1282) -- (-0.8710,-0.1142) -- (-0.8893,-0.0996) -- (-0.9077,-0.0845) -- (-0.9261,-0.0688) -- (-0.9445,-0.0526) -- (-0.9630,-0.0357) -- (-0.9815,-0.0182) -- (-1.0000,-0.0000);
\draw[mirror] (1.0000,0.0000) -- (0.9815,0.0182) -- (0.9630,0.0357) -- (0.9445,0.0526) -- (0.9261,0.0688) -- (0.9077,0.0845) -- (0.8893,0.0996) -- (0.8710,0.1142) -- (0.8528,0.1282) -- (0.8346,0.1417) -- (0.8165,0.1547) -- (0.7985,0.1672) -- (0.7805,0.1793) -- (0.7626,0.1910) -- (0.7448,0.2022) -- (0.7271,0.2130) -- (0.7095,0.2234) -- (0.6919,0.2334) -- (0.6745,0.2430) -- (0.6571,0.2523) -- (0.6398,0.2612) -- (0.6226,0.2698) -- (0.6055,0.2780) -- (0.5884,0.2860) -- (0.5715,0.2936) -- (0.5546,0.3009) -- (0.5378,0.3080) -- (0.5211,0.3147) -- (0.5045,0.3212) -- (0.4879,0.3274) -- (0.4714,0.3333) -- (0.4550,0.3390) -- (0.4386,0.3445) -- (0.4224,0.3497) -- (0.4061,0.3546) -- (0.3900,0.3594) -- (0.3739,0.3639) -- (0.3579,0.3682) -- (0.3419,0.3723) -- (0.3260,0.3761) -- (0.3101,0.3798) -- (0.2943,0.3833) -- (0.2785,0.3865) -- (0.2628,0.3896) -- (0.2471,0.3925) -- (0.2315,0.3951) -- (0.2159,0.3976) -- (0.2003,0.4000) -- (0.1848,0.4021) -- (0.1693,0.4040) -- (0.1538,0.4058) -- (0.1383,0.4074) -- (0.1229,0.4089) -- (0.1075,0.4101) -- (0.0921,0.4112) -- (0.0767,0.4121) -- (0.0614,0.4129) -- (0.0460,0.4135) -- (0.0307,0.4139) -- (0.0153,0.4141) -- (0.0000,0.4142) -- (-0.0153,0.4141) -- (-0.0307,0.4139) -- (-0.0460,0.4135) -- (-0.0614,0.4129) -- (-0.0767,0.4121) -- (-0.0921,0.4112) -- (-0.1075,0.4101) -- (-0.1229,0.4089) -- (-0.1383,0.4074) -- (-0.1538,0.4058) -- (-0.1693,0.4040) -- (-0.1848,0.4021) -- (-0.2003,0.4000) -- (-0.2159,0.3976) -- (-0.2315,0.3951) -- (-0.2471,0.3925) -- (-0.2628,0.3896) -- (-0.2785,0.3865) -- (-0.2943,0.3833) -- (-0.3101,0.3798) -- (-0.3260,0.3761) -- (-0.3419,0.3723) -- (-0.3579,0.3682) -- (-0.3739,0.3639) -- (-0.3900,0.3594) -- (-0.4061,0.3546) -- (-0.4224,0.3497) -- (-0.4386,0.3445) -- (-0.4550,0.3390) -- (-0.4714,0.3333) -- (-0.4879,0.3274) -- (-0.5045,0.3212) -- (-0.5211,0.3147) -- (-0.5378,0.3080) -- (-0.5546,0.3009) -- (-0.5715,0.2936) -- (-0.5884,0.2860) -- (-0.6055,0.2780) -- (-0.6226,0.2698) -- (-0.6398,0.2612) -- (-0.6571,0.2523) -- (-0.6745,0.2430) -- (-0.6919,0.2334) -- (-0.7095,0.2234) -- (-0.7271,0.2130) -- (-0.7448,0.2022) -- (-0.7626,0.1910) -- (-0.7805,0.1793) -- (-0.7985,0.1672) -- (-0.8165,0.1547) -- (-0.8346,0.1417) -- (-0.8528,0.1282) -- (-0.8710,0.1142) -- (-0.8893,0.0996) -- (-0.9077,0.0845) -- (-0.9261,0.0688) -- (-0.9445,0.0526) -- (-0.9630,0.0357) -- (-0.9815,0.0182) -- (-1.0000,0.0000);
\node[ax4] at (1.0000,0.0000) {};
\node[ax4] at (-1.0000,0.0000) {};
\node[ax4] at (0.0000,0.0000) {};
\node[ax4] at (0.0000,1.0000) {};
\node[ax4] at (0.0000,-1.0000) {};
\node[ax2, rotate=0] at (0.4142,0.0000) {};
\node[ax2, rotate=180] at (-0.4142,0.0000) {};
\node[ax2, rotate=45] at (0.7071,0.7071) {};
\node[ax2, rotate=-45] at (0.7071,-0.7071) {};
\node[ax2, rotate=135] at (-0.7071,0.7071) {};
\node[ax2, rotate=-135] at (-0.7071,-0.7071) {};
\node[ax2, rotate=90] at (0.0000,0.4142) {};
\node[ax2, rotate=-90] at (0.0000,-0.4142) {};
\node[ax3] at (0.3660,0.3660) {};
\node[ax3] at (0.3660,-0.3660) {};
\node[ax3] at (-0.3660,0.3660) {};
\node[ax3] at (-0.3660,-0.3660) {};

\end{tikzpicture}
\caption{$m\br3m$}
\end{subfigure}\hfill
\begin{subfigure}[b]{0.32\textwidth}\centering
\begin{tikzpicture}[scale=1.55]
  \draw[mirror, line width=1pt] (0,0) circle (1);
  \foreach \a in {0,45,90,135}{\draw[mirror] (\a:1) -- (\a+180:1);
    \node[ax2, rotate=\a] at (\a:1) {}; \node[ax2, rotate=\a] at (\a+180:1) {};}
  \node[ax4] at (0,0) {};
\end{tikzpicture}
\caption{$4/mmm$ (растяжение)}
\end{subfigure}\hfill
\begin{subfigure}[b]{0.32\textwidth}\centering
\begin{tikzpicture}[scale=1.55]
  \draw[atomLine, line width=0.6pt] (0,0) circle (1);
  \foreach \a in {0,45,90,135}{\draw[mirror] (\a:1) -- (\a+180:1);}
  \node[ax4] at (0,0) {};
\end{tikzpicture}
\caption{$4mm$ (поле $\vb E$)}
\end{subfigure}
\caption{Принцип суперпозиции Кюри: кубический кристалл $m\br3m$
(\textit{а}) при одноосном растяжении вдоль $[001]$ (перпендикулярно
плоскости рисунка) сохраняет элементы класса $4/mmm$ (\textit{б}), а в
электрическом поле вдоль $[001]$~— класса $4mm$ (\textit{в}): исчезают
горизонтальная плоскость (основной круг теперь тонкий) и оси 2}
\label{fig:curie-principle}
\end{figure}

\subsection{Тензорное описание физических свойств}\label{ssec:tensors}

\paragraph{Кристаллофизическая система координат}
Свойства описывают в правой прямоугольной системе $x_1x_2x_3$. В
кубической, тетрагональной и ромбической сингониях её оси совпадают с
кристаллографическими; в остальных её ориентируют по стандартным
правилам (табл.~\ref{tab:physaxes}).

\begin{table}[ht]
\centering
\caption{Правила кристаллофизической установки}\label{tab:physaxes}
\small
\begin{tabular}{@{}llll@{}}
\toprule
Сингония & Ось $x_1$ & Ось $x_2$ & Ось $x_3$\\
\midrule
Триклинная & \multicolumn{2}{l}{в плоскости, перпендикулярной $[001]$} & $[001]$\\
Моноклинная & в плоскости $(100)$ & $[010]$ & $[001]$\\
Ромбическая, тетрагон., кубическая & $[100]$ & $[010]$ & $[001]$\\
Тригон. и гексагональная & $[2\br1\br10]$ & $[01\br10]$ & $[0001]$\\
\bottomrule
\end{tabular}
\end{table}

\paragraph{Свойство как связь воздействия и отклика}
Физическое свойство проявляется как коэффициент в соотношении
\begin{equation}\label{eq:property}
\text{явление (отклик)}=\text{свойство}\times\text{воздействие} .
\end{equation}
Ранг тензора свойства определяется рангами воздействия и отклика:
\begin{itemize}
\item скаляр $\to$ скаляр: свойство~— \emph{скаляр} (тензор нулевого
ранга: плотность, теплоёмкость, энтропия);
\item скаляр $\to$ вектор: \emph{вектор} (тензор 1-го ранга). Пример~—
\emph{пироэффект}: при однородном нагреве (скаляр) кристалл поляризуется
($\vb P$~— вектор); возможен только в 10 полярных классах;
\item вектор $\to$ вектор: \emph{тензор 2-го ранга}: электро- и
теплопроводность, удельное электро- и тепловое сопротивление,
диэлектрическая и магнитная проницаемость и восприимчивость~— все они
\emph{симметричные} тензоры 2-го ранга;
\item тензор 2-го ранга $\to$ вектор: \emph{3-й ранг} (пьезоэффект; эффект
Холла: $\rho_{ij}(\vb B)=\rho_{ij}+R_{ijk}B_k$);
\item тензор 2-го ранга $\to$ тензор 2-го ранга: \emph{4-й ранг} (упругость,
магнитострикция).
\end{itemize}

При повороте прямоугольной системы координат (матрица $c$, п.~\ref{ssec:orth})
компоненты тензора 2-го ранга преобразуются как произведения координат:
\begin{equation}\label{eq:tensor2}
A'_{kl}=c_{ki}c_{lj}A_{ij}
\end{equation}
(правило Эйнштейна: по дважды повторяющемуся индексу подразумевается
суммирование от 1 до 3). Тензор \emph{симметричен}, если $A_{ij}=A_{ji}$, и
\emph{антисимметричен}, если $A_{ij}=-A_{ji}$.

\paragraph{Материальные и полевые тензоры}
\emph{Материальные} тензоры описывают свойства кристалла
(проводимость, упругость); их симметрия подчинена принципу Неймана.
\emph{Полевые} тензоры описывают воздействия и отклики (поле $\vb E$,
напряжение, деформация); их можно задать любыми, независимо от симметрии
кристалла. Но отклик на произвольное воздействие уже зависит от симметрии
(например, деформация при заданном напряжении определяется упругими
константами).

\subsection{Симметрия кристалла и тензоры второго ранга}\label{ssec:tensor2-symmetry}

Возьмём тензор электропроводности $\sigma$: $\vb j=\sigma\vb E$, т.\,е.
$j_i=\sigma_{ij}E_j$. При одновременной смене знаков $\vb E$ и $\vb j$
компоненты $\sigma_{ij}$ не меняются, т.\,е. \vazhno{любое свойство,
описываемое тензором 2-го ранга, центросимметрично}~— его симметрия
выше симметрии кристалла по крайней мере на центр инверсии.

Как симметрия «обнуляет» компоненты? Если операция $R$ (матрица)~—
элемент симметрии кристалла, то по принципу Неймана тензор не меняется:
\begin{equation}\label{eq:tensor-invariance}
\sigma'=R\,\sigma\,R^{-1}=\sigma .
\end{equation}
\begin{itemize}
\item \emph{Триклинный} кристалл: ограничений нет, у симметричного тензора
6 независимых компонент.
\item \emph{Моноклинный}, ось 2 вдоль $x_2$: $C_{2y}=\operatorname{diag}(\br1,1,\br1)$,
и
\[
C_{2y}\sigma C_{2y}^{-1}=\begin{pmatrix}\sigma_{11}&-\sigma_{12}&\sigma_{13}\\-\sigma_{12}&\sigma_{22}&-\sigma_{23}\\\sigma_{13}&-\sigma_{23}&\sigma_{33}\end{pmatrix}=\sigma
\ \Rightarrow\ \sigma_{12}=\sigma_{23}=0 ;
\]
остаются 4 компоненты: $\sigma_{11}$, $\sigma_{22}$, $\sigma_{33}$, $\sigma_{13}$.
\item \emph{Ромбический}: оси 2 вдоль $x_1$ и $x_3$ обнуляют и $\sigma_{13}$~— 3
компоненты.
\item \emph{Кубический}: ось 3 вдоль $[111]$ циклически переставляет оси
$x_1\to x_2\to x_3$ и переводит $\operatorname{diag}(\sigma_{11},\sigma_{22},\sigma_{33})$ в
$\operatorname{diag}(\sigma_{33},\sigma_{11},\sigma_{22})$; значит,
$\sigma_{11}=\sigma_{22}=\sigma_{33}=\sigma$~— одна компонента: \vazhno{кубический
кристалл изотропен по всем свойствам 2-го ранга} (но не по упругости!).
\end{itemize}

\begin{table}[ht]
\centering
\caption{Число независимых компонент симметричного тензора 2-го ранга}\label{tab:tensor2}
\small
\begin{tabular}{@{}lcl@{}}
\toprule
Сингония & Число & Независимые компоненты\\
\midrule
Триклинная & 6 & $\sigma_{11}$, $\sigma_{22}$, $\sigma_{33}$, $\sigma_{12}$, $\sigma_{13}$, $\sigma_{23}$\\
Моноклинная & 4 & $\sigma_{11}$, $\sigma_{22}$, $\sigma_{33}$, $\sigma_{13}$\\
Ромбическая & 3 & $\sigma_{11}$, $\sigma_{22}$, $\sigma_{33}$\\
Тетрагональная, тригональная, & 2 & $\sigma_{11}=\sigma_{22}$, $\sigma_{33}$\\
\quad гексагональная & & \\
Кубическая & 1 & $\sigma_{11}=\sigma_{22}=\sigma_{33}=\sigma$\\
\bottomrule
\end{tabular}
\end{table}

\begin{primer}[Проводимость в одноосном кристалле]
У гексагонального (тетрагонального, тригонального) кристалла
$\sigma_{11}=\sigma_{22}=\sigma_\perp$, $\sigma_{33}=\sigma_\parallel$. Для направления под
углом $\theta$ к оси $c$ проводимость (проекция $\vb j$ на $\vb E$)
$\sigma(\theta)=\sigma_\perp\sin^2\theta+\sigma_\parallel\cos^2\theta$. Для цинка, например,
удельное сопротивление вдоль и поперёк оси $c$ различается лишь на
несколько процентов, а у графита (п.~\ref{ssec:elem-types})~— в $10^4$ раз и более.
\end{primer}

\subsection{Упругие свойства кристаллов}\label{ssec:elastic-crystals}

\paragraph{Напряжения и деформации}
Однородное напряжённое состояние задаётся симметричным тензором
напряжений $\sigma_{ij}$: $\sigma_{11}$, $\sigma_{22}$, $\sigma_{33}$~— нормальные
напряжения, остальные~— касательные (сдвиговые). Приведя тензор к
главным осям, получим главные напряжения $\sigma_1$, $\sigma_2$, $\sigma_3$.

Деформацию описывают через смещения $\vb u$ точек тела. Если точки $P$ и
$Q$ отстоят на $\Delta\vb x$, то после деформации разность их смещений
\begin{equation}\label{eq:distortion}
\Delta u_i=u_{ij}\Delta x_j,\qquad u_{ij}=\frac{\partial u_i}{\partial x_j}
\end{equation}
($u_{ij}$~— \term{тензор дисторсии}). Удлинение $e_{11}=\partial u_1/\partial x_1$~—
относительное изменение длины, сдвиг $e_{12}=\partial u_1/\partial x_2=\operatorname{tg}\alpha$~—
смещение одной части тела относительно другой. Но не всякая дисторсия~—
деформация: тело можно просто повернуть. Разложим дисторсию на
симметричную и антисимметричную части (рис.~\ref{fig:distortion}):
\begin{keyeq}\label{eq:strain}
u_{ij}=\underbrace{\tfrac12(u_{ij}+u_{ji})}_{\varepsilon_{ij}}+\underbrace{\tfrac12(u_{ij}-u_{ji})}_{\omega_{ij}} .
\end{keyeq}
Симметричный $\varepsilon_{ij}$~— \term{тензор деформации} (его можно привести к
главным осям и главным деформациям $\varepsilon_1$, $\varepsilon_2$, $\varepsilon_3$);
антисимметричный $\omega_{ij}$~— \term{тензор поворота}, равносильный
аксиальному вектору поворота $\vb\omega=(\omega_1,\omega_2,\omega_3)$, где
$\omega_{12}=-\omega_3$, $\omega_{13}=\omega_2$, $\omega_{23}=-\omega_1$.

\begin{figure}[ht]
\centering
\begin{tikzpicture}[scale=1.25]
  \tikzset{orig/.style={fmGray, dashed, line width=0.5pt}, def/.style={fmBlue, thick, fill=fmBlue!12}}
  \begin{scope}
    \draw[def] (0,0) -- (1.5,0) -- (1.95,1.5) -- (0.45,1.5) -- cycle;
    \draw[orig] (0,0) rectangle (1.5,1.5);
    \draw[fmRed, -{Stealth[length=1.6mm]}] (0,1.5) -- (0.42,1.5);
    \node[font=\small] at (0.75,-0.45) {дисторсия $u_{12}$};
  \end{scope}
  \node[font=\Large] at (2.5,0.75) {$=$};
  \begin{scope}[xshift=3.1cm]
    \draw[def] (0,0) -- (1.5,0.225) -- (1.725,1.725) -- (0.225,1.5) -- cycle;
    \draw[orig] (0,0) rectangle (1.5,1.5);
    \node[font=\small] at (0.75,-0.45) {деформация $\varepsilon_{12}$};
  \end{scope}
  \node[font=\Large] at (5.5,0.75) {$+$};
  \begin{scope}[xshift=6.1cm]
    \draw[def, rotate around={-8.5:(0,0)}] (0,0) rectangle (1.5,1.5);
    \draw[orig] (0,0) rectangle (1.5,1.5);
    \node[font=\small] at (0.75,-0.45) {поворот $\omega$};
  \end{scope}
\end{tikzpicture}
\caption{Простой сдвиг (дисторсия $u_{12}$) есть сумма чистой деформации
(квадрат вытягивается вдоль диагонали, $\varepsilon_{12}=\varepsilon_{21}=u_{12}/2$) и
поворота тела как целого на угол $u_{12}/2$}
\label{fig:distortion}
\end{figure}

\paragraph{Обобщённый закон Гука}
Для изотропного тела $\sigma=c\varepsilon$ ($c$~— жёсткость) или $\varepsilon=s\sigma$
($s=1/c$~— податливость). В кристалле
\begin{keyeq}\label{eq:hooke-gen}
\sigma_{ij}=c_{ijkl}\varepsilon_{kl},\qquad \varepsilon_{ij}=s_{ijkl}\sigma_{kl},
\end{keyeq}
где $c_{ijkl}$ (упругие жёсткости) и $s_{ijkl}$ (податливости)~— тензоры
4-го ранга. Из $3^4=81$ компоненты независимы лишь 21: симметрия тензоров
напряжений и деформаций сводит их число к 36, а существование упругой
энергии ($c_{ijkl}=c_{klij}$)~— к 21. Дальше работает симметрия кристалла.

Удобна \emph{матричная} (фойгтовская) запись: пары индексов заменяют одним:
$11\to1$, $22\to2$, $33\to3$, $23,32\to4$, $13,31\to5$, $12,21\to6$. Тогда
$\sigma_i=c_{ij}\varepsilon_j$, $\varepsilon_i=s_{ij}\sigma_j$ ($i,j=1\ldots6$), а матрицы
$c$ и $s$ взаимно обратны\isprmark. Для кубических и гексагональных кристаллов
\[
c_\text{куб}=\begin{pmatrix}c_{11}&c_{12}&c_{12}&0&0&0\\c_{12}&c_{11}&c_{12}&0&0&0\\c_{12}&c_{12}&c_{11}&0&0&0\\0&0&0&c_{44}&0&0\\0&0&0&0&c_{44}&0\\0&0&0&0&0&c_{44}\end{pmatrix},
\]
\[
c_\text{гекс}=\begin{pmatrix}c_{11}&c_{12}&c_{13}&0&0&0\\c_{12}&c_{11}&c_{13}&0&0&0\\c_{13}&c_{13}&c_{33}&0&0&0\\0&0&0&c_{44}&0&0\\0&0&0&0&c_{44}&0\\0&0&0&0&0&\frac{c_{11}-c_{12}}2\end{pmatrix}
\]
(у матриц $s$~— та же структура, только в последнем элементе
гексагональной матрицы стоит $2(s_{11}-s_{12})$). Итак, кубический кристалл
описывается \vazhno{тремя} упругими константами, гексагональный~—
\vazhno{пятью}, изотропное тело~— \vazhno{двумя} ($s_{11}=1/E$,
$s_{12}=-\nu/E$, $s_{44}=1/G$, причём $G=E/2(1+\nu)$, п.~3.8.1).
\isprtext{В матричной записи сдвиговые компоненты деформации берут
удвоенными: $\varepsilon_4=2\varepsilon_{23}$, $\varepsilon_5=2\varepsilon_{13}$, $\varepsilon_6=2\varepsilon_{12}$
(так, чтобы $c_{44}$ был обычным модулем сдвига); в оригинале этот
множитель не оговорён. Соответственно $s_{44}=4s_{2323}$ и т.\,п.}

\paragraph{Анизотропия упругости кубических кристаллов}
Её меру даёт \term{фактор анизотропии}
\begin{keyeq}\label{eq:zener}
A=\frac{2c_{44}}{c_{11}-c_{12}}=\frac{2(s_{11}-s_{12})}{s_{44}} ;
\end{keyeq}
для изотропного тела $A=1$. Модуль Юнга в направлении с направляющими
косинусами $l_1$, $l_2$, $l_3$:
\begin{keyeq}\label{eq:E-cubic}
\frac1{E}=s_{11}-2\Bigl(s_{11}-s_{12}-\frac{s_{44}}2\Bigr)\bigl(l_1^2l_2^2+l_2^2l_3^2+l_3^2l_1^2\bigr),
\end{keyeq}
для гексагональных кристаллов (угол $\theta$ с осью $c$, $l_3=\cos\theta$)
\begin{equation}\label{eq:E-hex}
\frac1E=(1-l_3^2)^2s_{11}+l_3^4s_{33}+l_3^2(1-l_3^2)(2s_{13}+s_{44}) .
\end{equation}
Сумма $l_1^2l_2^2+l_2^2l_3^2+l_3^2l_1^2$ равна 0 для $\nfam{100}$ и $\frac13$ для
$\nfam{111}$. Поэтому при $A>1$ (тогда $s_{11}-s_{12}-s_{44}/2>0$) самое
жёсткое направление~— $\nfam{111}$, самое податливое~— $\nfam{100}$; при
$A<1$~— наоборот (табл.~\ref{tab:cubic-elastic}, рис.~\ref{fig:young-polar}).
Это определяется величиной $A$, а не типом решётки: среди ОЦК-металлов
есть и $\alpha$-Fe, Ta ($A>1$), и Cr, Mo, V, Nb ($A<1$)\ispr{В оригинале
сказано, что у ГЦК-кристаллов модуль Юнга растёт при переходе от $\nfam{100}$ к
$\nfam{111}$, а у ОЦК~— убывает. Таблица того же оригинала опровергает это:
у ОЦК-железа и тантала $E_{111}>E_{001}$. Направление роста определяется
знаком $A-1$.}. Почти изотропны вольфрам ($A\approx0{,}98$) и алюминий
($A\approx1{,}2$)\ispr{В оригинале наиболее изотропным среди ОЦК-металлов
назван молибден ($A=0{,}91$), хотя у вольфрама $A=0{,}98$.}.

\begin{table}[ht]
\centering
\caption{Упругие постоянные кубических монокристаллов (ГПа)}\label{tab:cubic-elastic}
\small
\begin{tabular}{@{}lcccccccc@{}}
\toprule
Металл & $c_{11}$ & $c_{12}$ & $c_{44}$ & $E_{001}$ & $E_{011}$ & $E_{111}$ & $E_\text{ср}$ & $A$\\
\midrule
Al & 108 & 61 & 28 & 63 & 72 & 76 & 70 & 1,23\\
Cu & 168 & 121 & 75 & 67 & 130 & 191 & 127 & 3,21\\
Ag & 124 & 93 & 46 & 43 & 83 & 120 & 81 & 3,04\\
Au & 186 & 157 & 42 & 43 & 81 & 116 & 78 & 2,87\\
Ni & 246 & 147 & 125 & 130 & 223 & 294 & 214 & 2,54\\
Pb & 50 & 42 & 15 & 11 & 24 & 40 & 24 & 4,12\\
$\alpha$-Fe & 231 & 135 & 116 & 132 & 220 & 283 & 211 & 2,41\\
Ta & 267 & 161 & 82 & 142 & 191 & 215 & 183 & 1,58\\
V & 228 & 119 & 43 & 146 & 124 & 117 & 128 & 0,78\\
Nb & 246 & 139 & 29 & 152 & 92 & 82 & 105 & 0,51\\
Cr & 350 & 67 & 101 & 328 & 266 & 250 & 280 & 0,71\\
Mo & 460 & 176 & 110 & 344 & 326 & 316 & 327 & 0,91\\
W & 501 & 198 & 151 & 420 & 414 & 412 & 415 & 0,98\\
UO$_2$ & 395 & 121 & 64 & 338 & 199 & 175 & 230 & 0,47\\
\bottomrule
\end{tabular}
\end{table}

\begin{table}[ht]
\centering
\caption{Упругие постоянные гексагональных монокристаллов (ГПа)}\label{tab:hex-elastic}
\small
\begin{tabular}{@{}lcccccccc@{}}
\toprule
Металл & $c_{11}$ & $c_{12}$ & $c_{13}$ & $c_{33}$ & $c_{44}$ & $c_{66}$ & $E_\text{ср}$ & $c/a$\\
\midrule
Be & 292 & 27 & 14 & 336 & 162 & 133 & 311 & 1,57\\
Mg & 53 & 26 & 21 & 61 & 16 & 17 & 44 & 1,62\\
Ti & 162 & 92 & 69 & 181 & 47 & 35 & 114 & 1,59\\
$\alpha$-Zr & 144 & 73 & 66 & 165 & 32 & 36 & 97 & 1,59\\
Hf & 181 & 77 & 66 & 197 & 56 & 52 & 143 & 1,58\\
Co & 307 & 165 & 105 & 357 & 76 & 71 & 216 & 1,62\\
Re & 613 & 270 & 206 & 683 & 162 & 171 & 462 & 1,62\\
Zn & 165 & 31 & 50 & 62 & 40 & 67 & 102 & 1,86\\
Cd & 114 & 40 & 60 & 51 & 20 & 38 & 63 & 1,89\\
BeO & 470 & 168 & 119 & 494 & 153 & 151 & 394 & 1,62\\
\bottomrule
\multicolumn{9}{l}{\footnotesize $c_{66}=(c_{11}-c_{12})/2$.}
\end{tabular}
\end{table}

\noindent В гексагональной таблице исправлены опечатки
оригинала\ispr{В табл.~1.8 оригинала для кобальта $c_{33}=36$~ГПа (правильно
около 357~ГПа; при 36~ГПа кристалл был бы вдесятеро мягче вдоль оси $c$,
чем поперёк), для кадмия $c/a=1{,}62$ (в тексте того же оригинала верно:
1,89), для BeO $c/a=1{,}58$ (у BeO со структурой вюрцита $c/a\approx1{,}62$).}.
Обратите внимание на цинк и кадмий с аномальным $c/a$: вдоль оси $c$ они
гораздо мягче ($c_{33}\ll c_{11}$)~— связь между базисными слоями слабее, чем
внутри слоя.

\subsection{Матричное представление группы}\label{ssec:representations}

Самый общий способ выяснить, сколько независимых компонент оставляет
симметрия у тензора любого ранга,~— теория представлений групп.

\begin{opredelenie}[Представление группы]
Сопоставление каждому элементу $R$ группы $G$ матрицы $D(R)$ так, что
произведению элементов отвечает произведение матриц: если $AB=C$, то
$D(A)D(B)=D(C)$. Число строк матриц~— \term{размерность} представления.
\end{opredelenie}

\begin{itemize}
\item Представления $D$ и $D'$ \emph{эквивалентны}, если
$D'(R)=SD(R)S^{-1}$ для всех $R$ с одной и той же матрицей $S$.
\item Представление \emph{приводимо}, если подходящей заменой базиса все
его матрицы одновременно приводятся к блочно-диагональному виду; тогда
оно~— прямая сумма представлений меньшей размерности:
$D(R)=D^{(1)}(R)\oplus D^{(2)}(R)$. Основную роль играют неприводимые
неэквивалентные представления.
\item \emph{Характер} представления~— набор следов $\chi(R)=\operatorname{tr}D(R)$.
Элементы одного класса имеют одинаковый след, поэтому характер~—
функция класса.
\end{itemize}

\begin{primer}[Тензор 2-го ранга при повороте вокруг оси 4]
Пусть ось 4 направлена вдоль $x_1$. Из $\sigma'=C_{4x}\sigma C_{4x}^{-1}$ получаем
$\sigma'_{11}=\sigma_{11}$, $\sigma'_{22}=\sigma_{33}$, $\sigma'_{33}=\sigma_{22}$, $\sigma'_{12}=-\sigma_{13}$,
$\sigma'_{13}=\sigma_{12}$, $\sigma'_{23}=-\sigma_{23}$. Записав шесть компонент столбцом
$(\sigma_{11},\sigma_{22},\sigma_{33},\sigma_{12},\sigma_{13},\sigma_{23})$, получаем
шестимерную матрицу $D(C_{4x})$, которая распадается в прямую сумму
одномерных $(1)$ и $(-1)$ и двумерных
$\bigl(\begin{smallmatrix}0&1\\1&0\end{smallmatrix}\bigr)$,
$\bigl(\begin{smallmatrix}0&-1\\1&0\end{smallmatrix}\bigr)$ блоков. Её след
$\chi=1+0+0+0+0-1=0$, что совпадает с формулой~\eqref{eq:chi-sym} ниже при
$\varphi=90^\circ$.
\end{primer}

\paragraph{Теоремы ортогональности}
Для неприводимых неэквивалентных представлений $D^{(i)}$ размерности $l_i$
группы порядка $h$
\begin{equation}\label{eq:orthogonality}
\sum_{R\in G}D^{(i)}(R)_{\mu\nu}\,D^{(j)}(R)^*_{\alpha\beta}=\frac h{l_i}\,\delta_{ij}\delta_{\mu\alpha}\delta_{\nu\beta}
\end{equation}
(суммирование~— по всем \emph{элементам группы})\ispr{В оригинале
сказано, что суммирование в теореме ортогональности ведётся «по всем
неэквивалентным, неприводимым представлениям»; на самом деле сумма берётся
по элементам группы $R$, а индексы $i$, $j$ нумеруют представления.}. Отсюда
следуют: $\sum_il_i^2=h$; ортогональность характеров
$\sum_kN_k\chi^{(i)}(C_k)\chi^{(j)}(C_k)^*=h\delta_{ij}$ (сумма по классам $C_k$, $N_k$~—
число элементов в классе) и то, что число неприводимых представлений
равно числу классов.

\paragraph{Сколько независимых компонент у тензора}
Компонента, не меняющаяся ни при одной операции группы, отвечает
\emph{единичному} (полносимметричному) представлению с $\chi\equiv1$. Число
таких независимых комбинаций~— кратность единичного представления в
представлении тензора:
\begin{keyeq}\label{eq:n-invariants}
n=\frac1h\sum_kN_k\,\chi(C_k) ,
\end{keyeq}
где $\chi(C_k)$~— след матрицы, по которой преобразуются компоненты тензора,
для операций класса $C_k$. Для поворота на угол $\varphi$ ($\det R=+1$) и
несобственного поворота ($\det R=-1$; поворот на $\varphi$ с отражением в
перпендикулярной плоскости)
\begin{equation}\label{eq:chi-sym}
\chi_\text{сим}=4\cos^2\varphi\pm2\cos\varphi\quad\text{(симметричный тензор 2-го ранга)},
\end{equation}
\begin{equation}\label{eq:chi-gen}
\chi=(\pm1+2\cos\varphi)^2\quad\text{(произвольный тензор 2-го ранга)}
\end{equation}
(знак «$+$» для собственных, «$-$» для несобственных операций).

\begin{primer}[Ромбический кристалл $mmm$]
Элементы группы ($h=8$, каждый~— свой класс): $E$ ($\varphi=0$, $\chi=6$); три
оси 2 ($\varphi=180^\circ$, $\chi=4-2=2$); центр инверсии~— несобственный
поворот на $180^\circ$ ($\chi=4+2=6$); три плоскости~— несобственные повороты на
$0^\circ$ ($\chi=4-2=2$). По~\eqref{eq:n-invariants}
$n=\frac18(6+3\cdot2+6+3\cdot2)=3$: у симметричного тензора 2-го ранга
ромбического кристалла (например, электропроводности $\alpha$-урана) три
независимые компоненты~— как и в табл.~\ref{tab:tensor2}.
\end{primer}

\begin{itogi}
\begin{itemize}
\item Принцип Неймана: $G_\text{свойства}\supseteq G_\text{кристалла}$; принцип Кюри:
$G_\text{кр+возд}=G_\text{кр}\cap G_\text{возд}$.
\item Свойства~— тензоры: скаляры, векторы (пироэффект), 2-го ранга
(проводимость, проницаемость), 3-го (пьезоэффект), 4-го (упругость).
\item Свойства 2-го ранга центросимметричны; число независимых компонент:
6, 4, 3, 2, 1 (от триклинной до кубической сингонии); кубический
кристалл изотропен по ним.
\item Упругость: $\sigma_{ij}=c_{ijkl}\varepsilon_{kl}$; 21 константа в общем
случае, 5~— у гексагональных, 3~— у кубических, 2~— у изотропных;
анизотропия кубических кристаллов $A=2c_{44}/(c_{11}-c_{12})$.
\item Число независимых компонент: $n=\frac1h\sum_kN_k\chi(C_k)$.
\end{itemize}
\end{itogi}

\begin{proverka}
\item Какую симметрию приобретёт кубический кристалл при сжатии вдоль
$[110]$?
\item Почему кубический кристалл не может быть пироэлектриком?
\item Почему электропроводность кубического кристалла изотропна, а
модуль Юнга~— нет?
\item Чему равен модуль Юнга меди вдоль $\nfam{110}$, если
$c_{11}=168$, $c_{12}=121$, $c_{44}=75$~ГПа?
\item Найдите по~\eqref{eq:n-invariants} число независимых компонент
симметричного тензора 2-го ранга для класса $4/m$.
\item Почему у цинка $c_{33}$ много меньше $c_{11}$?
\end{proverka}

\section{Кристаллография пластической деформации}\label{sec:plastic}
% ----------------------------------------------------------------------
%  1.6. Кристаллография пластической деформации
% ----------------------------------------------------------------------
\begin{orientir}
Как кристалл меняет форму, не разрушаясь? В~1.6.1 разберём два способа
пластической деформации~— \vazhno{скольжение и двойникование}, закон
Шмида и то, какие системы скольжения работают в ГЦК, ОЦК и ГПУ-металлах.
Скольжение поворачивает решётку, поэтому деформированный поликристалл
приобретает \vazhno{текстуру}; в~1.6.2~— как её описывают (полюсные
фигуры и функция распределения ориентаций).
\end{orientir}

Под действием механического напряжения кристалл деформируется. При малых
напряжениях деформация \term{упругая}: после разгрузки форма полностью
восстанавливается, а деформация и напряжение связаны законом Гука
(п.~\ref{ssec:elastic-crystals}). Выше \term{предела текучести} появляется
\term{остаточная} (пластическая) деформация: после разгрузки кристалл
сохраняет изменённую форму. Практически предел упругости (последнее
напряжение, при котором выполняется закон Гука) неотличим от предела
текучести.

\subsection{Геометрия пластической деформации монокристаллов}\label{ssec:slip-geometry}

Монокристалл деформируется пластически двумя способами
(рис.~\ref{fig:slip-twin}).

\begin{opredelenie}
\term{Скольжение}~— сдвиг одной части кристалла относительно другой по
определённой кристаллографической плоскости (\term{плоскости скольжения})
в определённом направлении (\term{направлении скольжения}) на величину,
\emph{кратную} периоду решётки в этом направлении. Пару «плоскость +
направление» называют \term{системой скольжения} и записывают как
$\plfam{hkl}\nfam{uvw}$.

\term{Механическое двойникование}~— сдвиг, в результате которого часть
кристалла переходит в положение, зеркально-симметричное исходному (или
повёрнутое относительно него на $180^\circ$). Атомы при этом смещаются на
\emph{доли} периода, пропорционально расстоянию от плоскости
двойникования.
\end{opredelenie}

\begin{figure}[ht]
\centering
\begin{tikzpicture}[scale=0.38]
  % (а) исходный кристалл
  \begin{scope}
    \foreach \j in {0,...,7}{\foreach \i in {0,...,6}{
      \fill[atomA!80!black] ({\i+0.35*\j},{0.8*\j}) circle (0.17);}}
    \draw[atomLine, line width=0.6pt] (-0.575,-0.4) -- (6.225,-0.4) -- (9.025,6.0) -- (2.225,6.0) -- cycle;
    \node[font=\small] at (4.2,-1.5) {\textit{а}};
  \end{scope}
  % (б) скольжение: верхняя половина сдвинута на два периода
  \begin{scope}[shift={(11,0)}]
    \foreach \j in {0,...,3}{\foreach \i in {0,...,6}{
      \fill[atomA!80!black] ({\i+0.35*\j},{0.8*\j}) circle (0.17);}}
    \foreach \j in {4,...,7}{\foreach \i in {0,...,6}{
      \fill[atomA!80!black] ({\i+0.35*\j+2},{0.8*\j}) circle (0.17);}}
    \draw[atomLine, line width=0.6pt] (-0.575,-0.4) -- (6.225,-0.4) -- (7.625,2.8) -- (9.625,2.8)
      -- (11.025,6.0) -- (4.225,6.0) -- (2.825,2.8) -- (0.825,2.8) -- cycle;
    \draw[fmRed, dashed, line width=0.7pt] (-0.6,2.8) -- (11.4,2.8);
    \draw[arrB] (6.2,6.7) -- (8.2,6.7);
    \node[font=\small] at (5.2,-1.5) {\textit{б}};
  \end{scope}
  % (в) двойникование: слои 1 и 5 — границы двойниковой прослойки
  \begin{scope}[shift={(25,0)}]
    \fill[fmAmber!20] (-1.6,0.8) rectangle (7.0,4.0);
    \foreach \j in {0,...,7}{\foreach \i in {0,...,6}{
      \pgfmathsetmacro{\sh}{\j<=1 ? 0.35*\j : (\j<=5 ? 0.35-0.35*(\j-1) : 0.35-1.4+0.35*(\j-5))}
      \pgfmathsetmacro{\col}{(\j>=2 && \j<=4) ? 1 : 0}
      \ifdim\col pt>0.5pt
        \fill[atomB!90!black] ({\i+\sh},{0.8*\j}) circle (0.17);
      \else
        \fill[atomA!80!black] ({\i+\sh},{0.8*\j}) circle (0.17);
      \fi}}
    \draw[fmAmber!80!black, dashed, line width=0.7pt] (-1.6,0.8) -- (7.0,0.8) (-1.6,4.0) -- (7.0,4.0);
    \node[font=\small] at (2.7,-1.5) {\textit{в}};
  \end{scope}
\end{tikzpicture}
\caption{Два механизма пластической деформации: \textit{а}~— исходный
кристалл; \textit{б}~— скольжение по плоскости, показанной штриховой линией:
верхняя часть сдвинута на целое число периодов, решётка не изменилась,
на боковых гранях появились ступеньки; \textit{в}~— двойникование: в
прослойке (оранжевые атомы) каждый слой сдвинут относительно
предыдущего на долю периода, и решётка зеркально отражена относительно
границ прослойки (штриховые линии)}
\label{fig:slip-twin}
\end{figure}

Главный механизм пластической деформации при обычных и низких
температурах~— скольжение. Объём при скольжении не меняется, а решётка
после сдвига остаётся той же (сдвиг кратен трансляции), поэтому скольжение
называют также \term{трансляционным}. На атомном уровне скольжение~— это
движение дислокаций (гл.~2, п.~2.2.2).

\subsubsection{Закон Шмида}

Растянем цилиндрический монокристалл напряжением $\sigma$ вдоль оси $F$
(рис.~\ref{fig:schmid}). Пусть нормаль $\vb n$ к плоскости скольжения
образует с осью угол $\chi$, а направление скольжения $\vb m$~— угол
$\lambda$. Сила $\sigma A$ действует на площадку $A/\cos\chi$, её проекция
на направление $\vb m$ равна $\sigma A\cos\lambda$. Поэтому касательное
(«скалывающее») напряжение в системе скольжения
\begin{keyeq}\label{eq:schmid}
\tau=\sigma\cos\chi\cos\lambda .
\end{keyeq}
Множитель $\cos\chi\cos\lambda$ называют \term{фактором Шмида}
(ориентационным множителем).

\begin{glavnoe}[Закон Шмида (закон критического скалывающего напряжения)]
Скольжение начинается, когда касательное напряжение хотя бы в одной из
систем скольжения достигает критического значения $\tau_\text{кр}$~—
постоянной материала, \emph{не зависящей} от ориентации кристалла.
Первой включается система с наибольшим фактором Шмида.
\end{glavnoe}

\begin{figure}[ht]
\centering
\begin{tikzpicture}[scale=1.15]
  \fill[fmBlueBg] (-0.9,-2) -- (-0.9,2) arc[start angle=180, end angle=360, x radius=0.9, y radius=0.25]
    -- (0.9,-2) arc[start angle=0, end angle=-180, x radius=0.9, y radius=0.25];
  \draw[atomLine] (-0.9,-2) -- (-0.9,2) (0.9,-2) -- (0.9,2);
  \draw[atomLine] (0,2) ellipse (0.9 and 0.25);
  \draw[atomLine] (-0.9,-2) arc[start angle=180, end angle=360, x radius=0.9, y radius=0.25];
  \draw[atomLine, dashed] (0.9,-2) arc[start angle=0, end angle=180, x radius=0.9, y radius=0.25];
  \fill[fmRed!25, rotate=30] (0,0) ellipse (1.04 and 0.25);
  \draw[fmRed!80!black, rotate=30] (0,0) ellipse (1.04 and 0.25);
  \node[font=\footnotesize, fmRed!80!black] at (-1.35,-0.75) {$S$};
  \draw[arrB, fmBlue] (0,2.3) -- (0,3.1) node[above, font=\small]{$F$};
  \draw[arrB, fmBlue] (0,-2.3) -- (0,-3.1);
  \draw[fmGray, dashed] (0,-2.0) -- (0,2.0);
  \draw[bvec, fmGreen!70!black] (0,0) -- (-0.75,1.30) node[above left=-2pt, font=\small]{$\vb n$};
  \draw[bvec] (0,0) -- (0.62,0.82) node[right, font=\small]{$\vb m$};
  \draw[fmGreen!70!black] (0,0.75) arc[start angle=90, end angle=120, radius=0.75];
  \node[font=\footnotesize, fmGreen!70!black] at (-0.22,0.95) {$\chi$};
  \draw[fmRed] (0,0.45) arc[start angle=90, end angle=53, radius=0.45];
  \node[font=\footnotesize, fmRed] at (0.17,0.62) {$\lambda$};
  \node[font=\small, fmBlue] at (0.35,-2.75) {$\sigma$};
\end{tikzpicture}
\caption{К закону Шмида: ось растяжения $F$, плоскость скольжения $S$
с нормалью $\vb n$ и направление скольжения $\vb m$}
\label{fig:schmid}
\end{figure}

Фактор Шмида не превышает $\frac12$: максимум достигается, когда $\vb n$,
$\vb m$ и $F$ лежат в одной плоскости и $\chi=\lambda=45^\circ$. Если ось
лежит в плоскости скольжения ($\chi=90^\circ$) или перпендикулярна
направлению скольжения ($\lambda=90^\circ$), скольжение в этой системе
невозможно при любом напряжении~— кристалл в такой ориентации кажется
очень прочным.

\begin{masshtab}[где происходит скольжение]{-9.75/-9.35/{$b$ и $d_{111}$~— доли нм}, -8/-6.6/{ступеньки 10--300 нм}, -6.3/-5/{полосы через 1--10 мкм}, -5/-4/{зерно}, -3.3/-1.7/{образец}}
Цилиндр на рис.~\ref{fig:schmid}~— это \emph{образец} (миллиметры--сантиметры),
а плоскость $S$~— не одна атомная плоскость, а \emph{ориентация}: все
параллельные плоскости $\plfam{111}$ кристалла. В меди
($d_{111}=0{,}21$~нм) на 1~мм вдоль нормали их около $5\cdot10^6$, и
фактор Шмида у всех одинаков. Поэтому он характеризует ориентацию
кристалла (зерна) как целого, а не какую-то одну плоскость.

Сдвиг же идёт не по всем этим плоскостям, а по немногим, где работают
источники дислокаций (рис.~\ref{fig:slip-scales}). Каждая дислокация
сдвигает кристалл на $b=0{,}26$~нм; сотни дислокаций, вышедших по одной
плоскости, дают на поверхности ступеньку высотой десятки--сотни
нанометров~— \term{линию скольжения}. Линии собираются в \term{полосы
скольжения}, отстоящие друг от друга на микрометры; их видно в
оптический микроскоп. В поликристалле роль «цилиндра» играет отдельное
зерно (10--100~мкм) со своим фактором Шмида, а предел текучести
поликристалла $\sigma_\text{т}\approx M\tau_\text{кр}$, где $M\approx3{,}1$
для ГЦК (фактор Тейлора, усреднение по ориентациям зёрен) вместо
$1/0{,}5=2$ у наилучшим образом ориентированного монокристалла.
\end{masshtab}

\begin{figure}[ht]
\centering
\begin{tikzpicture}[font=\footnotesize]
  % ---------- (а) образец: миллиметры
  \begin{scope}
    \fill[fmBlueBg] (-0.7,-1.5) rectangle (0.7,1.5);
    \begin{scope}
      \clip (-0.7,-1.5) rectangle (0.7,1.5);
      \foreach \c in {-2.4,-2.25,...,2.4}{\draw[fmBlue!45, line width=0.3pt] (-0.7,{\c-0.42}) -- (0.7,{\c+0.42});}
    \end{scope}
    \draw[atomLine] (-0.7,-1.5) -- (-0.7,1.5) (0.7,-1.5) -- (0.7,1.5);
    \draw[atomLine, fill=fmBlueBg] (0,1.5) ellipse (0.7 and 0.16);
    \draw[atomLine] (-0.7,-1.5) arc[start angle=180, end angle=360, x radius=0.7, y radius=0.16];
    \draw[arrB, fmBlue] (0,1.75) -- (0,2.2) node[above, inner sep=1pt]{$F$};
    \draw[arrB, fmBlue] (0,-1.75) -- (0,-2.2);
    \draw[fmAmber!80!black, line width=0.8pt] (-0.12,-0.12) rectangle (0.12,0.12);
    \draw[atomLine, line width=0.8pt] (-0.7,-2.45) -- (-0.23,-2.45) node[right, inner sep=2pt]{1~мм};
    \node[font=\small] at (0,-2.85) {\textit{а}\ образец};
  \end{scope}
  % ---------- (б) поверхность: микрометры
  \begin{scope}[shift={(3.0,0)}]
    \fill[fmBlueBg] (-1.4,-1.4) rectangle (1.4,1.4);
    \begin{scope}
      \clip (-1.4,-1.4) rectangle (1.4,1.4);
      \foreach \c in {-2.1,-1.35,...,2.2}{
        \foreach \d in {0,0.06,0.12,0.18}{\draw[fmBlue!80!black, line width=0.35pt] (-1.4,{\c+\d-0.84}) -- (1.4,{\c+\d+0.84});}}
    \end{scope}
    \draw[atomLine] (-1.4,-1.4) rectangle (1.4,1.4);
    \draw[fmAmber!80!black, line width=0.8pt] (0.05,0.25) rectangle (0.45,0.55);
    \draw[atomLine, line width=0.8pt] (-1.4,-1.75) -- (-0.65,-1.75) node[right, inner sep=2pt]{5~мкм};
    \node[font=\small] at (0,-2.85) {\textit{б}\ полосы};
  \end{scope}
  \draw[atomLine!50, densely dotted] (0.12,0.12) -- (1.6,1.4) (0.12,-0.12) -- (1.6,-1.4);
  % ---------- (в) профиль ступенек: сотни нанометров
  \begin{scope}[shift={(6.3,0)}]
    \fill[fmBlueBg] (-1.5,-1.4) -- (-1.5,-0.55) -- (-0.75,-0.55) -- (-0.45,-0.05) -- (0.05,-0.05) -- (0.2,0.2)
      -- (0.65,0.2) -- (1.0,0.8) -- (1.5,0.8) -- (1.5,-1.4) -- cycle;
    \draw[atomLine, line width=0.9pt] (-1.5,-0.55) -- (-0.75,-0.55) -- (-0.45,-0.05) -- (0.05,-0.05) -- (0.2,0.2)
      -- (0.65,0.2) -- (1.0,0.8) -- (1.5,0.8);
    \foreach \x/\y in {-0.75/-0.55, 0.05/-0.05, 0.65/0.2}{
      \draw[fmRed!70, dashed, line width=0.5pt] (\x,\y) -- ++(-0.9,-1.5);}
    \draw[dimline, fmRed] (1.15,0.2) -- node[right, inner sep=1pt, fmRed]{$h=nb$} (1.15,0.8);
    \draw[fmAmber!80!black, line width=0.8pt] (-0.86,-0.75) rectangle (-0.56,-0.45);
    \draw[atomLine, line width=0.8pt] (-1.5,-1.75) -- (-0.9,-1.75) node[right, inner sep=2pt]{100~нм};
    \node[font=\small] at (0,-2.85) {\textit{в}\ ступеньки};
  \end{scope}
  \draw[atomLine!50, densely dotted] (3.45,0.55) -- (4.8,1.4) (3.45,0.25) -- (4.8,-1.4);
  % ---------- (г) атомы: нанометры
  \begin{scope}[shift={(9.75,0)}]
    \begin{scope}
      \clip (-1.45,-1.4) rectangle (1.45,1.4);
      \begin{scope}[scale=0.27]\EdgeDislLattice{6}{5}{5}\end{scope}
    \end{scope}
    \draw[atomLine] (-1.45,-1.4) rectangle (1.45,1.4);
    \draw[fmRed, dashed, line width=0.7pt] (-1.45,0) -- (1.45,0);
    \dislT[fmRed]{(0,0)}{0}
    \draw[atomLine, line width=0.8pt] (-1.45,-1.75) -- (-0.4,-1.75) node[right, inner sep=2pt]{1~нм};
    \node[font=\small] at (0,-2.85) {\textit{г}\ дислокация};
  \end{scope}
  \draw[atomLine!50, densely dotted] (5.74,-0.45) -- (8.3,1.4) (5.74,-0.75) -- (8.3,-1.4);
\end{tikzpicture}
\caption{Масштабы скольжения: \textit{а}~— растягиваемый образец с
линиями скольжения на поверхности (миллиметры); \textit{б}~— полосы
скольжения на поверхности зерна или монокристалла: пучки близких линий,
отстоящие друг от друга на микрометры; \textit{в}~— разрез поверхности:
каждая линия скольжения~— ступенька высотой $h=nb$ (десятки--сотни
нанометров), образованная выходом $n\sim10^2$--$10^3$ дислокаций по одной
плоскости (пунктир); \textit{г}~— атомный масштаб: одна краевая
дислокация в плоскости скольжения сдвигает кристалл на $b\approx0{,}26$~нм.
Каждая следующая панель~— увеличение выделенного участка предыдущей}
\label{fig:slip-scales}
\end{figure}

Если кристалл находится под действием произвольного тензора напряжений
$\sigma$, касательное напряжение в системе с единичной нормалью $\vb h$ и
единичным направлением $\vb m$ вычисляют как
\begin{equation}\label{eq:schmid-tensor}
\tau_{\vb h\vb m}=\tilde h\,\sigma\, m=\sum_{i,j}h_i\sigma_{ij}m_j ,
\end{equation}
где $\tilde h$~— строка направляющих косинусов нормали, $m$~— столбец
направляющих косинусов направления скольжения. Для одноосного растяжения
$\sigma_{ij}=\sigma F_iF_j$ формула~\eqref{eq:schmid-tensor} переходит
в~\eqref{eq:schmid}.

\subsubsection{Тензор пластической дисторсии}

При скольжении на величину $\alpha$ (сдвиг на единицу расстояния от
плоскости скольжения) точка $P$ с радиус-вектором $\vb r$ смещается
вдоль $\vb m$ на величину, пропорциональную её расстоянию $\vb r\cdot\vb h$
от плоскости скольжения, проходящей через начало координат
(рис.~\ref{fig:shear-vector}):
\begin{equation}\label{eq:slip-displacement}
\vb u=\vb r'-\vb r=\alpha\,(\vb r\cdot\vb h)\,\vb m .
\end{equation}

\begin{figure}[ht]
\centering
\begin{tikzpicture}[scale=1.1]
  \fill[fmBlueBg] (-0.3,0) rectangle (4.3,2.2);
  \fill[fmRed!10] (-0.3,0) -- (4.3,0) -- (5.2,2.2) -- (0.6,2.2) -- cycle;
  \draw[atomLine, dashed] (-0.3,0) rectangle (4.3,2.2);
  \draw[fmRed!80!black, line width=0.8pt] (-0.3,0) -- (4.3,0) -- (5.2,2.2) -- (0.6,2.2) -- cycle;
  \draw[fmGray, line width=1pt] (-0.6,0) -- (5.4,0) node[right, font=\footnotesize]{плоскость скольжения};
  \fill (0,0) circle (1.5pt) node[below left, font=\small]{$O$};
  \fill (2.2,1.5) circle (1.5pt) node[above left, font=\small]{$P$};
  \fill (2.81,1.5) circle (1.5pt) node[above right, font=\small]{$P'$};
  \draw[arr] (0,0) -- (2.2,1.5) node[midway, above left=-2pt, font=\small]{$\vb r$};
  \draw[arr, fmRed] (2.2,1.5) -- (2.81,1.5);
  \draw[fmGray, densely dotted] (2.2,1.5) -- (2.2,0);
  \node[font=\footnotesize, anchor=west] at (2.25,0.6) {$\vb r\cdot\vb h$};
  \draw[bvec, fmGreen!70!black] (-0.3,0) -- (-0.3,0.9) node[left, font=\small]{$\vb h$};
  \draw[bvec] (0.6,-0.35) -- (1.6,-0.35) node[right, font=\small]{$\vb m$};
\end{tikzpicture}
\caption{Смещение точки при скольжении: $\overrightarrow{PP'}=\alpha(\vb r\cdot\vb h)\vb m$
(пунктир~— форма кристалла до сдвига)}
\label{fig:shear-vector}
\end{figure}

Компоненты \term{тензора пластической дисторсии}
$u^0_{ij}=\partial u_i/\partial x_j$ (сравните с~\eqref{eq:strain}): так как
$u_i=\alpha\,m_i\sum_kx_kh_k$, то $u^0_{ij}=\alpha\,m_ih_j$, или в матричной форме
\begin{keyeq}\label{eq:plastic-distortion}
u^0=\alpha\, m\,\tilde h=\alpha
\begin{pmatrix}
m_1h_1 & m_1h_2 & m_1h_3\\
m_2h_1 & m_2h_2 & m_2h_3\\
m_3h_1 & m_3h_2 & m_3h_3
\end{pmatrix}\!.
\end{keyeq}
Например, $u^0_{11}=\alpha m_1h_1$, $u^0_{23}=\partial u_2/\partial x_3=\alpha m_2h_3$.

Свойства этого тензора:
\begin{itemize}
\item След $u^0_{11}+u^0_{22}+u^0_{33}=\alpha(\vb m\cdot\vb h)=0$, так как
$\vb m\perp\vb h$: \vazhno{скольжение не меняет объём}.
\item Тензор несимметричен. Его симметричная часть
$\varepsilon_{ij}=\frac\alpha2(m_ih_j+m_jh_i)$~— \emph{чистая пластическая
деформация}, антисимметричная
$\omega_{ij}=\frac\alpha2(m_ih_j-m_jh_i)$~— \emph{поворот решётки}.
Именно этот поворот приводит к образованию текстуры (п.~\ref{ssec:texture}).
\item Отрезок материала вдоль оси $F$ после скольжения переходит, по
\eqref{eq:slip-displacement}, в $F'=F+\alpha(\vb h\cdot F)\,\vb m$~— он
\emph{поворачивается к направлению скольжения} $\vb m$, оставаясь в
плоскости, содержащей $F$ и $\vb m$.
\end{itemize}

\begin{smysl}
Тензор деформаций с нулевым следом имеет $6-1=5$ независимых компонент.
Значит, чтобы зерно поликристалла могло принять \emph{любую} форму,
навязанную соседями (без нарушения сплошности), в нём должны работать
не менее \vazhno{пяти независимых систем скольжения} (критерий Мизеса).
ГЦК- и ОЦК-металлы этому условию удовлетворяют и пластичны в
поликристаллическом состоянии; ГПУ-металлы с одним базисным скольжением~—
нет.
\end{smysl}

\subsubsection{Скольжение в ГЦК-металлах}

Плоскостями и направлениями скольжения, как правило, служат плоскости и
направления \emph{плотнейшей упаковки}: у них наибольшее межплоскостное
расстояние и наименьший вектор трансляции (а значит, и вектор Бюргерса
дислокаций). В ГЦК-металлах (Al, Cu, Ni, Ag, Au, $\gamma$-Fe) система
скольжения~— $\plfam{111}\nfam{110}$: четыре плоскости $\plfam{111}$, в
каждой три направления $\nfam{110}$, итого \vazhno{12 систем}
скольжения\ispr{В оригинале: «8 плоскостей скольжения с шестью
направлениями в каждой плоскости»~— это подсчёт с учётом знаков
нормалей и направлений; различных систем 12.}.

Какая из 12 систем включится первой, зависит от ориентации оси
растяжения. Благодаря кубической симметрии достаточно рассматривать оси,
лежащие в одном \term{стандартном стереографическом треугольнике}~—
здесь, как и в оригинале курса, это треугольник
$[100]$--$[110]$--$[111]$ (рис.~\ref{fig:slip-stereo}): любая другая
ориентация переводится в него операцией симметрии. В табл.~\ref{tab:schmid}
приведены факторы Шмида всех 12 систем для оси $F\parallel[531]$, лежащей
внутри треугольника.

\begin{table}[ht]
\centering
\caption{Факторы Шмида систем $\plfam{111}\nfam{110}$ ГЦК-кристалла при
растяжении вдоль $[531]$ (расчёт по~\eqref{eq:schmid})}
\label{tab:schmid}
\small
\begin{tabular}{@{}lcl@{\hspace{8mm}}lc@{}}
\toprule
Система & $\cos\chi\cos\lambda$ & Название & Система & $\cos\chi\cos\lambda$\\
\midrule
$(11\br1)[101]$ & \textbf{0,490} & первичная   & $(111)[1\br10]$   & 0,210\\
$(111)[10\br1]$ & 0,420 & критическая          & $(11\br1)[1\br10]$ & 0,163\\
$(11\br1)[011]$ & 0,327 &                      & $(1\br11)[10\br1]$ & 0,140\\
$(1\br11)[110]$ & 0,280 & сопряжённая          & $(1\br11)[011]$    & 0,140\\
$(111)[01\br1]$ & 0,210 &                      & $(1\br1\br1)[110]$ & 0,093\\
$(1\br1\br1)[101]$ & 0,070 & поперечная        & $(1\br1\br1)[01\br1]$ & 0,023\\
\bottomrule
\end{tabular}
\end{table}

\begin{figure}[t]
\centering
\tikzset{zone100/.style={fmBlue!70, line width=0.5pt},
  zone110/.style={fmBlue!35, line width=0.4pt},
  pole100/.style={rectangle, fill=atomLine, inner sep=0pt, minimum size=4pt},
  pole110/.style={ellipse, fill=atomLine, inner sep=0pt, minimum width=4.5pt, minimum height=3pt},
  pole111/.style={regular polygon, regular polygon sides=3, fill=atomLine, inner sep=0pt, minimum size=6pt},
  ring/.style={circle, draw=#1, line width=1pt, minimum size=8pt, inner sep=0pt},
  poleH/.style={circle, fill=fmGreen!70!black, inner sep=0pt, minimum size=3pt}}
\begin{tikzpicture}
  % создано tools/gen_plastic.py: системы скольжения ГЦК, проекция вдоль [100]
\begin{scope}
\fill[fmAmber!25] (0.0000,0.0000) -- (0.0242,0.0000) -- (0.0483,0.0000) -- (0.0725,0.0000) -- (0.0967,0.0000) -- (0.1209,0.0000) -- (0.1452,0.0000) -- (0.1694,0.0000) -- (0.1937,0.0000) -- (0.2181,0.0000) -- (0.2425,0.0000) -- (0.2669,0.0000) -- (0.2914,0.0000) -- (0.3160,0.0000) -- (0.3406,0.0000) -- (0.3653,0.0000) -- (0.3900,0.0000) -- (0.4149,0.0000) -- (0.4398,0.0000) -- (0.4648,0.0000) -- (0.4900,0.0000) -- (0.5152,0.0000) -- (0.5405,0.0000) -- (0.5660,0.0000) -- (0.5915,0.0000) -- (0.6172,0.0000) -- (0.6431,0.0000) -- (0.6691,0.0000) -- (0.6952,0.0000) -- (0.7214,0.0000) -- (0.7479,0.0000) -- (0.7745,0.0000) -- (0.8012,0.0000) -- (0.8282,0.0000) -- (0.8553,0.0000) -- (0.8827,0.0000) -- (0.9102,0.0000) -- (0.9379,0.0000) -- (0.9659,0.0000) -- (0.9941,0.0000) -- (0.9941,0.0000) -- (0.9940,0.0222) -- (0.9938,0.0444) -- (0.9935,0.0666) -- (0.9930,0.0888) -- (0.9923,0.1110) -- (0.9915,0.1332) -- (0.9906,0.1554) -- (0.9895,0.1776) -- (0.9882,0.1998) -- (0.9868,0.2221) -- (0.9853,0.2444) -- (0.9836,0.2666) -- (0.9818,0.2889) -- (0.9798,0.3112) -- (0.9777,0.3336) -- (0.9754,0.3559) -- (0.9730,0.3783) -- (0.9704,0.4007) -- (0.9676,0.4231) -- (0.9647,0.4455) -- (0.9617,0.4680) -- (0.9585,0.4905) -- (0.9551,0.5130) -- (0.9516,0.5356) -- (0.9479,0.5582) -- (0.9441,0.5808) -- (0.9400,0.6034) -- (0.9359,0.6261) -- (0.9315,0.6489) -- (0.9270,0.6716) -- (0.9223,0.6944) -- (0.9175,0.7173) -- (0.9124,0.7402) -- (0.9072,0.7631) -- (0.9018,0.7861) -- (0.8963,0.8091) -- (0.8905,0.8322) -- (0.8846,0.8553) -- (0.8785,0.8785) -- (0.8785,0.8785) -- (0.8523,0.8523) -- (0.8264,0.8264) -- (0.8008,0.8008) -- (0.7756,0.7756) -- (0.7506,0.7506) -- (0.7259,0.7259) -- (0.7014,0.7014) -- (0.6772,0.6772) -- (0.6532,0.6532) -- (0.6295,0.6295) -- (0.6059,0.6059) -- (0.5826,0.5826) -- (0.5594,0.5594) -- (0.5365,0.5365) -- (0.5137,0.5137) -- (0.4911,0.4911) -- (0.4687,0.4687) -- (0.4464,0.4464) -- (0.4242,0.4242) -- (0.4022,0.4022) -- (0.3803,0.3803) -- (0.3585,0.3585) -- (0.3369,0.3369) -- (0.3153,0.3153) -- (0.2939,0.2939) -- (0.2725,0.2725) -- (0.2512,0.2512) -- (0.2300,0.2300) -- (0.2089,0.2089) -- (0.1878,0.1878) -- (0.1668,0.1668) -- (0.1459,0.1459) -- (0.1249,0.1249) -- (0.1041,0.1041) -- (0.0832,0.0832) -- (0.0624,0.0624) -- (0.0416,0.0416) -- (0.0208,0.0208) -- (0.0000,0.0000) -- cycle;
\draw[zone100] (-0.0000,2.3768) -- (-0.0000,2.2953) -- (-0.0000,2.2164) -- (-0.0000,2.1400) -- (-0.0000,2.0659) -- (-0.0000,1.9941) -- (-0.0000,1.9243) -- (-0.0000,1.8564) -- (-0.0000,1.7903) -- (-0.0000,1.7260) -- (-0.0000,1.6632) -- (0.0000,1.6019) -- (0.0000,1.5421) -- (0.0000,1.4835) -- (0.0000,1.4263) -- (0.0000,1.3702) -- (0.0000,1.3152) -- (0.0000,1.2612) -- (0.0000,1.2082) -- (0.0000,1.1562) -- (0.0000,1.1050) -- (0.0000,1.0546) -- (0.0000,1.0050) -- (0.0000,0.9561) -- (0.0000,0.9079) -- (0.0000,0.8604) -- (0.0000,0.8134) -- (0.0000,0.7670) -- (0.0000,0.7210) -- (0.0000,0.6756) -- (0.0000,0.6306) -- (0.0000,0.5860) -- (0.0000,0.5418) -- (0.0000,0.4980) -- (0.0000,0.4544) -- (0.0000,0.4112) -- (0.0000,0.3682) -- (0.0000,0.3254) -- (0.0000,0.2829) -- (0.0000,0.2405) -- (0.0000,0.1982) -- (0.0000,0.1561) -- (0.0000,0.1141) -- (0.0000,0.0722) -- (0.0000,0.0302) -- (0.0000,-0.0116) -- (0.0000,-0.0535) -- (0.0000,-0.0955) -- (0.0000,-0.1375) -- (0.0000,-0.1795) -- (0.0000,-0.2217) -- (0.0000,-0.2640) -- (0.0000,-0.3065) -- (0.0000,-0.3492) -- (0.0000,-0.3921) -- (0.0000,-0.4352) -- (0.0000,-0.4352) -- (0.0000,-0.4786) -- (0.0000,-0.5223) -- (0.0000,-0.5664) -- (0.0000,-0.6108) -- (0.0000,-0.6556) -- (0.0000,-0.7008) -- (0.0000,-0.7465) -- (0.0000,-0.7927) -- (0.0000,-0.8394) -- (0.0000,-0.8867) -- (0.0000,-0.9347) -- (0.0000,-0.9832) -- (0.0000,-1.0325) -- (0.0000,-1.0825) -- (0.0000,-1.1333) -- (0.0000,-1.1850) -- (0.0000,-1.2376) -- (0.0000,-1.2911) -- (0.0000,-1.3456) -- (0.0000,-1.4012) -- (0.0000,-1.4580) -- (0.0000,-1.5159) -- (0.0000,-1.5752) -- (0.0000,-1.6358) -- (0.0000,-1.6979) -- (0.0000,-1.7616) -- (0.0000,-1.8269) -- (0.0000,-1.8939) -- (0.0000,-1.9628) -- (0.0000,-2.0338) -- (0.0000,-2.1068) -- (0.0000,-2.1821) -- (0.0000,-2.2599) -- (0.0000,-2.3403);
\draw[zone100] (-2.3570,0.0000) -- (-2.2761,0.0000) -- (-2.1978,0.0000) -- (-2.1220,0.0000) -- (-2.0485,0.0000) -- (-1.9771,0.0000) -- (-1.9078,0.0000) -- (-1.8404,0.0000) -- (-1.7747,0.0000) -- (-1.7108,0.0000) -- (-1.6484,0.0000) -- (-1.5874,0.0000) -- (-1.5279,0.0000) -- (-1.4697,0.0000) -- (-1.4127,0.0000) -- (-1.3569,0.0000) -- (-1.3021,0.0000) -- (-1.2484,0.0000) -- (-1.1956,0.0000) -- (-1.1438,0.0000) -- (-1.0928,0.0000) -- (-1.0426,0.0000) -- (-0.9932,0.0000) -- (-0.9445,0.0000) -- (-0.8965,0.0000) -- (-0.8490,0.0000) -- (-0.8022,0.0000) -- (-0.7559,0.0000) -- (-0.7101,0.0000) -- (-0.6648,0.0000) -- (-0.6199,-0.0000) -- (-0.5754,-0.0000) -- (-0.5313,-0.0000) -- (-0.4875,-0.0000) -- (-0.4440,-0.0000) -- (-0.4008,-0.0000) -- (-0.3579,-0.0000) -- (-0.3152,-0.0000) -- (-0.2727,-0.0000) -- (-0.2303,-0.0000) -- (-0.1881,-0.0000) -- (-0.1460,-0.0000) -- (-0.1040,-0.0000) -- (-0.0621,-0.0000) -- (-0.0202,-0.0000) -- (0.0217,-0.0000) -- (0.0636,-0.0000) -- (0.1055,-0.0000) -- (0.1476,-0.0000) -- (0.1896,-0.0000) -- (0.2319,-0.0000) -- (0.2742,-0.0000) -- (0.3167,-0.0000) -- (0.3595,-0.0000) -- (0.4024,-0.0000) -- (0.4456,-0.0000) -- (0.4891,-0.0000) -- (0.5329,-0.0000) -- (0.5770,-0.0000) -- (0.6215,-0.0000) -- (0.6664,-0.0000) -- (0.7117,-0.0000) -- (0.7576,-0.0000) -- (0.8039,-0.0000) -- (0.8507,-0.0000) -- (0.8982,-0.0000) -- (0.9463,-0.0000) -- (0.9950,-0.0000) -- (1.0445,-0.0000) -- (1.0947,-0.0000) -- (1.1457,-0.0000) -- (1.1975,-0.0000) -- (1.2503,-0.0000) -- (1.3041,-0.0000) -- (1.3589,-0.0000) -- (1.4147,0.0000) -- (1.4718,0.0000) -- (1.5300,0.0000) -- (1.5896,0.0000) -- (1.6506,0.0000) -- (1.7130,0.0000) -- (1.7771,0.0000) -- (1.8428,0.0000) -- (1.9103,0.0000) -- (1.9797,0.0000) -- (2.0511,0.0000) -- (2.1247,0.0000) -- (2.2006,0.0000) -- (2.2790,0.0000) -- (2.3600,0.0000);
\draw[zone100] (-1.6754,1.6754) -- (-1.6179,1.6179) -- (-1.5623,1.5623) -- (-1.5084,1.5084) -- (-1.4562,1.4562) -- (-1.4055,1.4055) -- (-1.3563,1.3563) -- (-1.3084,1.3084) -- (-1.2618,1.2618) -- (-1.2164,1.2164) -- (-1.1721,1.1721) -- (-1.1289,1.1289) -- (-1.0867,1.0867) -- (-1.0454,1.0454) -- (-1.0049,1.0049) -- (-0.9653,0.9653) -- (-0.9265,0.9265) -- (-0.8884,0.8884) -- (-0.8510,0.8510) -- (-0.8143,0.8143) -- (-0.7781,0.7781) -- (-0.7426,0.7426) -- (-0.7075,0.7075) -- (-0.6730,0.6730) -- (-0.6390,0.6390) -- (-0.6054,0.6054) -- (-0.5722,0.5722) -- (-0.5394,0.5394) -- (-0.5070,0.5070) -- (-0.4749,0.4749) -- (-0.4431,0.4431) -- (-0.4116,0.4116) -- (-0.3803,0.3803) -- (-0.3493,0.3493) -- (-0.3186,0.3186) -- (-0.2880,0.2880) -- (-0.2576,0.2576) -- (-0.2274,0.2274) -- (-0.1973,0.1973) -- (-0.1674,0.1674) -- (-0.1375,0.1375) -- (-0.1077,0.1077) -- (-0.0780,0.0780) -- (-0.0484,0.0484) -- (-0.0187,0.0187) -- (0.0109,-0.0109) -- (0.0405,-0.0405) -- (0.0702,-0.0702) -- (0.0999,-0.0999) -- (0.1296,-0.1296) -- (0.1595,-0.1595) -- (0.1894,-0.1894) -- (0.2194,-0.2194) -- (0.2496,-0.2496) -- (0.2800,-0.2800) -- (0.3105,-0.3105) -- (0.3412,-0.3412) -- (0.3721,-0.3721) -- (0.4033,-0.4033) -- (0.4347,-0.4347) -- (0.4664,-0.4664) -- (0.4984,-0.4984) -- (0.5308,-0.5308) -- (0.5635,-0.5635) -- (0.5966,-0.5966) -- (0.6300,-0.6300) -- (0.6640,-0.6640) -- (0.6984,-0.6984) -- (0.7332,-0.7332) -- (0.7687,-0.7687) -- (0.8047,-0.8047) -- (0.8412,-0.8412) -- (0.8785,-0.8785) -- (0.8785,-0.8785) -- (0.9164,-0.9164) -- (0.9550,-0.9550) -- (0.9944,-0.9944) -- (1.0346,-1.0346) -- (1.0756,-1.0756) -- (1.1176,-1.1176) -- (1.1606,-1.1606) -- (1.2046,-1.2046) -- (1.2497,-1.2497) -- (1.2960,-1.2960) -- (1.3435,-1.3435) -- (1.3924,-1.3924) -- (1.4427,-1.4427) -- (1.4945,-1.4945) -- (1.5479,-1.5479) -- (1.6030,-1.6030) -- (1.6600,-1.6600);
\draw[zone100] (1.6949,1.6949) -- (1.6367,1.6367) -- (1.5805,1.5805) -- (1.5260,1.5260) -- (1.4733,1.4733) -- (1.4221,1.4221) -- (1.3724,1.3724) -- (1.3241,1.3241) -- (1.2771,1.2771) -- (1.2313,1.2313) -- (1.1866,1.1866) -- (1.1431,1.1431) -- (1.1005,1.1005) -- (1.0589,1.0589) -- (1.0182,1.0182) -- (0.9783,0.9783) -- (0.9393,0.9393) -- (0.9009,0.9009) -- (0.8633,0.8633) -- (0.8263,0.8263) -- (0.7900,0.7900) -- (0.7543,0.7543) -- (0.7191,0.7191) -- (0.6844,0.6844) -- (0.6502,0.6502) -- (0.6164,0.6164) -- (0.5831,0.5831) -- (0.5502,0.5502) -- (0.5176,0.5176) -- (0.4854,0.4854) -- (0.4535,0.4535) -- (0.4220,0.4220) -- (0.3906,0.3906) -- (0.3596,0.3596) -- (0.3287,0.3287) -- (0.2981,0.2981) -- (0.2677,0.2677) -- (0.2374,0.2374) -- (0.2073,0.2073) -- (0.1773,0.1773) -- (0.1474,0.1474) -- (0.1176,0.1176) -- (0.0878,0.0878) -- (0.0582,0.0582) -- (0.0285,0.0285) -- (-0.0011,-0.0011) -- (-0.0307,-0.0307) -- (-0.0604,-0.0604) -- (-0.0900,-0.0900) -- (-0.1198,-0.1198) -- (-0.1496,-0.1496) -- (-0.1795,-0.1795) -- (-0.2095,-0.2095) -- (-0.2396,-0.2396) -- (-0.2699,-0.2699) -- (-0.2699,-0.2699) -- (-0.3004,-0.3004) -- (-0.3310,-0.3310) -- (-0.3619,-0.3619) -- (-0.3930,-0.3930) -- (-0.4243,-0.4243) -- (-0.4559,-0.4559) -- (-0.4878,-0.4878) -- (-0.5200,-0.5200) -- (-0.5526,-0.5526) -- (-0.5856,-0.5856) -- (-0.6189,-0.6189) -- (-0.6527,-0.6527) -- (-0.6869,-0.6869) -- (-0.7217,-0.7217) -- (-0.7569,-0.7569) -- (-0.7927,-0.7927) -- (-0.8291,-0.8291) -- (-0.8661,-0.8661) -- (-0.9038,-0.9038) -- (-0.9421,-0.9421) -- (-0.9813,-0.9813) -- (-1.0212,-1.0212) -- (-1.0620,-1.0620) -- (-1.1036,-1.1036) -- (-1.1463,-1.1463) -- (-1.1899,-1.1899) -- (-1.2347,-1.2347) -- (-1.2806,-1.2806) -- (-1.3277,-1.3277) -- (-1.3761,-1.3761) -- (-1.4259,-1.4259) -- (-1.4772,-1.4772) -- (-1.5300,-1.5300) -- (-1.5846,-1.5846) -- (-1.6410,-1.6410);
\draw[zone110] (-0.0571,2.3415) -- (-0.1121,2.2824) -- (-0.1645,2.2233) -- (-0.2144,2.1644) -- (-0.2619,2.1058) -- (-0.3071,2.0474) -- (-0.3501,1.9893) -- (-0.3910,1.9314) -- (-0.4299,1.8739) -- (-0.4669,1.8168) -- (-0.5022,1.7599) -- (-0.5357,1.7035) -- (-0.5675,1.6474) -- (-0.5978,1.5916) -- (-0.6265,1.5362) -- (-0.6538,1.4812) -- (-0.6798,1.4266) -- (-0.7043,1.3723) -- (-0.7276,1.3183) -- (-0.7497,1.2647) -- (-0.7706,1.2114) -- (-0.7903,1.1584) -- (-0.8089,1.1057) -- (-0.8265,1.0534) -- (-0.8431,1.0013) -- (-0.8586,0.9495) -- (-0.8732,0.8980) -- (-0.8868,0.8467) -- (-0.8995,0.7957) -- (-0.9114,0.7449) -- (-0.9223,0.6943) -- (-0.9325,0.6439) -- (-0.9418,0.5937) -- (-0.9503,0.5437) -- (-0.9580,0.4938) -- (-0.9649,0.4441) -- (-0.9711,0.3945) -- (-0.9765,0.3450) -- (-0.9812,0.2956) -- (-0.9852,0.2463) -- (-0.9852,0.2463) -- (-0.9884,0.1971) -- (-0.9909,0.1479) -- (-0.9927,0.0988) -- (-0.9937,0.0497) -- (-0.9941,0.0006) -- (-0.9938,-0.0485) -- (-0.9927,-0.0976) -- (-0.9909,-0.1467) -- (-0.9885,-0.1958) -- (-0.9853,-0.2451) -- (-0.9813,-0.2944) -- (-0.9767,-0.3437) -- (-0.9713,-0.3932) -- (-0.9651,-0.4428) -- (-0.9582,-0.4925) -- (-0.9505,-0.5424) -- (-0.9420,-0.5924) -- (-0.9327,-0.6426) -- (-0.9226,-0.6930) -- (-0.9117,-0.7436) -- (-0.8998,-0.7944) -- (-0.8871,-0.8454) -- (-0.8735,-0.8967) -- (-0.8590,-0.9482) -- (-0.8435,-1.0000) -- (-0.8269,-1.0521) -- (-0.8094,-1.1044) -- (-0.7908,-1.1571) -- (-0.7711,-1.2100) -- (-0.7502,-1.2633) -- (-0.7282,-1.3169) -- (-0.7049,-1.3709) -- (-0.6804,-1.4252) -- (-0.6545,-1.4799) -- (-0.6272,-1.5349) -- (-0.5985,-1.5902) -- (-0.5683,-1.6460) -- (-0.5365,-1.7021) -- (-0.5030,-1.7585) -- (-0.4678,-1.8153) -- (-0.4309,-1.8725) -- (-0.3920,-1.9300) -- (-0.3511,-1.9878) -- (-0.3082,-2.0459) -- (-0.2630,-2.1043) -- (-0.2156,-2.1630) -- (-0.1658,-2.2218) -- (-0.1135,-2.2809) -- (-0.0585,-2.3400) -- (-0.0007,-2.3993);
\draw[zone110] (0.0216,-2.3782) -- (0.0784,-2.3190) -- (0.1324,-2.2599) -- (0.1838,-2.2009) -- (0.2327,-2.1421) -- (0.2793,-2.0836) -- (0.3237,-2.0253) -- (0.3659,-1.9672) -- (0.4060,-1.9095) -- (0.4442,-1.8522) -- (0.4805,-1.7951) -- (0.5151,-1.7384) -- (0.5480,-1.6821) -- (0.5792,-1.6261) -- (0.6089,-1.5705) -- (0.6371,-1.5153) -- (0.6639,-1.4604) -- (0.6893,-1.4059) -- (0.7133,-1.3517) -- (0.7362,-1.2979) -- (0.7578,-1.2444) -- (0.7782,-1.1912) -- (0.7975,-1.1383) -- (0.8158,-1.0858) -- (0.8329,-1.0335) -- (0.8491,-0.9816) -- (0.8642,-0.9299) -- (0.8785,-0.8785) -- (0.8917,-0.8273) -- (0.9041,-0.7763) -- (0.9156,-0.7256) -- (0.9263,-0.6751) -- (0.9361,-0.6248) -- (0.9451,-0.5746) -- (0.9533,-0.5247) -- (0.9607,-0.4749) -- (0.9674,-0.4252) -- (0.9733,-0.3756) -- (0.9784,-0.3262) -- (0.9828,-0.2768) -- (0.9865,-0.2276) -- (0.9894,-0.1784) -- (0.9917,-0.1292) -- (0.9932,-0.0801) -- (0.9940,-0.0310) -- (0.9941,0.0181) -- (0.9934,0.0671) -- (0.9921,0.1162) -- (0.9901,0.1654) -- (0.9873,0.2146) -- (0.9838,0.2638) -- (0.9796,0.3131) -- (0.9747,0.3625) -- (0.9690,0.4121) -- (0.9626,0.4617) -- (0.9554,0.5115) -- (0.9474,0.5614) -- (0.9386,0.6115) -- (0.9290,0.6618) -- (0.9185,0.7122) -- (0.9073,0.7629) -- (0.8951,0.8138) -- (0.8821,0.8649) -- (0.8681,0.9163) -- (0.8532,0.9679) -- (0.8373,1.0198) -- (0.8204,1.0719) -- (0.8025,1.1244) -- (0.7834,1.1772) -- (0.7633,1.2303) -- (0.7420,1.2837) -- (0.7195,1.3374) -- (0.6958,1.3915) -- (0.6707,1.4459) -- (0.6443,1.5007) -- (0.6165,1.5559) -- (0.5872,1.6114) -- (0.5564,1.6673) -- (0.5240,1.7235) -- (0.4899,1.7801) -- (0.4540,1.8370) -- (0.4163,1.8943) -- (0.3767,1.9520) -- (0.3350,2.0099) -- (0.2913,2.0681) -- (0.2453,2.1266) -- (0.1970,2.1853) -- (0.1462,2.2443) -- (0.0929,2.3034) -- (0.0369,2.3626);
\draw[zone110] (-2.3626,-0.0369) -- (-2.3034,-0.0929) -- (-2.2443,-0.1462) -- (-2.1853,-0.1970) -- (-2.1266,-0.2453) -- (-2.0681,-0.2913) -- (-2.0099,-0.3350) -- (-1.9520,-0.3767) -- (-1.8943,-0.4163) -- (-1.8370,-0.4540) -- (-1.7801,-0.4899) -- (-1.7235,-0.5240) -- (-1.6673,-0.5564) -- (-1.6114,-0.5872) -- (-1.5559,-0.6165) -- (-1.5007,-0.6443) -- (-1.4459,-0.6707) -- (-1.3915,-0.6958) -- (-1.3374,-0.7195) -- (-1.2837,-0.7420) -- (-1.2303,-0.7633) -- (-1.1772,-0.7834) -- (-1.1244,-0.8025) -- (-1.0719,-0.8204) -- (-1.0198,-0.8373) -- (-0.9679,-0.8532) -- (-0.9163,-0.8681) -- (-0.8649,-0.8821) -- (-0.8138,-0.8951) -- (-0.7629,-0.9073) -- (-0.7122,-0.9185) -- (-0.6618,-0.9290) -- (-0.6115,-0.9386) -- (-0.5614,-0.9474) -- (-0.5115,-0.9554) -- (-0.4617,-0.9626) -- (-0.4121,-0.9690) -- (-0.3625,-0.9747) -- (-0.3131,-0.9796) -- (-0.2638,-0.9838) -- (-0.2146,-0.9873) -- (-0.1654,-0.9901) -- (-0.1162,-0.9921) -- (-0.0671,-0.9934) -- (-0.0181,-0.9941) -- (0.0310,-0.9940) -- (0.0801,-0.9932) -- (0.1292,-0.9917) -- (0.1784,-0.9894) -- (0.2276,-0.9865) -- (0.2768,-0.9828) -- (0.3262,-0.9784) -- (0.3756,-0.9733) -- (0.4252,-0.9674) -- (0.4749,-0.9607) -- (0.5247,-0.9533) -- (0.5746,-0.9451) -- (0.6248,-0.9361) -- (0.6751,-0.9263) -- (0.7256,-0.9156) -- (0.7763,-0.9041) -- (0.8273,-0.8917) -- (0.8785,-0.8785) -- (0.9299,-0.8642) -- (0.9816,-0.8491) -- (1.0335,-0.8329) -- (1.0858,-0.8158) -- (1.1383,-0.7975) -- (1.1912,-0.7782) -- (1.2444,-0.7578) -- (1.2979,-0.7362) -- (1.3517,-0.7133) -- (1.4059,-0.6893) -- (1.4604,-0.6639) -- (1.5153,-0.6371) -- (1.5705,-0.6089) -- (1.6261,-0.5792) -- (1.6821,-0.5480) -- (1.7384,-0.5151) -- (1.7951,-0.4805) -- (1.8522,-0.4442) -- (1.9095,-0.4060) -- (1.9672,-0.3659) -- (2.0253,-0.3237) -- (2.0836,-0.2793) -- (2.1421,-0.2327) -- (2.2009,-0.1838) -- (2.2599,-0.1324) -- (2.3190,-0.0784) -- (2.3782,-0.0216);
\draw[zone110] (2.3625,0.0369) -- (2.3033,0.0929) -- (2.2442,0.1463) -- (2.1853,0.1970) -- (2.1266,0.2453) -- (2.0681,0.2913) -- (2.0098,0.3351) -- (1.9519,0.3767) -- (1.8943,0.4163) -- (1.8370,0.4540) -- (1.7800,0.4899) -- (1.7235,0.5240) -- (1.6672,0.5564) -- (1.6113,0.5872) -- (1.5558,0.6165) -- (1.5007,0.6443) -- (1.4459,0.6707) -- (1.3915,0.6958) -- (1.3374,0.7195) -- (1.2836,0.7420) -- (1.2302,0.7633) -- (1.1771,0.7834) -- (1.1244,0.8025) -- (1.0719,0.8204) -- (1.0197,0.8373) -- (0.9678,0.8532) -- (0.9162,0.8681) -- (0.8649,0.8821) -- (0.8137,0.8951) -- (0.7629,0.9073) -- (0.7122,0.9186) -- (0.6617,0.9290) -- (0.6115,0.9386) -- (0.5614,0.9474) -- (0.5114,0.9554) -- (0.4617,0.9626) -- (0.4120,0.9690) -- (0.3625,0.9747) -- (0.3131,0.9796) -- (0.2638,0.9838) -- (0.2145,0.9873) -- (0.1653,0.9901) -- (0.1162,0.9921) -- (0.0671,0.9934) -- (0.0180,0.9941) -- (-0.0311,0.9940) -- (-0.0802,0.9932) -- (-0.1293,0.9917) -- (-0.1784,0.9894) -- (-0.2276,0.9865) -- (-0.2769,0.9828) -- (-0.3262,0.9784) -- (-0.3757,0.9733) -- (-0.4252,0.9674) -- (-0.4749,0.9607) -- (-0.5247,0.9533) -- (-0.5747,0.9451) -- (-0.6248,0.9361) -- (-0.6751,0.9263) -- (-0.7257,0.9156) -- (-0.7764,0.9041) -- (-0.8273,0.8917) -- (-0.8785,0.8784) -- (-0.9299,0.8642) -- (-0.9816,0.8491) -- (-1.0336,0.8329) -- (-1.0858,0.8157) -- (-1.1384,0.7975) -- (-1.1912,0.7782) -- (-1.2444,0.7578) -- (-1.2979,0.7362) -- (-1.3517,0.7133) -- (-1.4059,0.6892) -- (-1.4604,0.6638) -- (-1.4604,0.6638) -- (-1.5153,0.6371) -- (-1.5706,0.6089) -- (-1.6262,0.5792) -- (-1.6822,0.5479) -- (-1.7385,0.5151) -- (-1.7952,0.4805) -- (-1.8522,0.4442) -- (-1.9096,0.4060) -- (-1.9673,0.3658) -- (-2.0253,0.3236) -- (-2.0836,0.2793) -- (-2.1422,0.2327) -- (-2.2010,0.1838) -- (-2.2599,0.1323) -- (-2.3191,0.0783) -- (-2.3783,0.0215);
\draw[fmAmber!80!black, line width=0.8pt] (0.0000,0.0000) -- (0.0242,0.0000) -- (0.0483,0.0000) -- (0.0725,0.0000) -- (0.0967,0.0000) -- (0.1209,0.0000) -- (0.1452,0.0000) -- (0.1694,0.0000) -- (0.1937,0.0000) -- (0.2181,0.0000) -- (0.2425,0.0000) -- (0.2669,0.0000) -- (0.2914,0.0000) -- (0.3160,0.0000) -- (0.3406,0.0000) -- (0.3653,0.0000) -- (0.3900,0.0000) -- (0.4149,0.0000) -- (0.4398,0.0000) -- (0.4648,0.0000) -- (0.4900,0.0000) -- (0.5152,0.0000) -- (0.5405,0.0000) -- (0.5660,0.0000) -- (0.5915,0.0000) -- (0.6172,0.0000) -- (0.6431,0.0000) -- (0.6691,0.0000) -- (0.6952,0.0000) -- (0.7214,0.0000) -- (0.7479,0.0000) -- (0.7745,0.0000) -- (0.8012,0.0000) -- (0.8282,0.0000) -- (0.8553,0.0000) -- (0.8827,0.0000) -- (0.9102,0.0000) -- (0.9379,0.0000) -- (0.9659,0.0000) -- (0.9941,0.0000) -- (0.9941,0.0000) -- (0.9940,0.0222) -- (0.9938,0.0444) -- (0.9935,0.0666) -- (0.9930,0.0888) -- (0.9923,0.1110) -- (0.9915,0.1332) -- (0.9906,0.1554) -- (0.9895,0.1776) -- (0.9882,0.1998) -- (0.9868,0.2221) -- (0.9853,0.2444) -- (0.9836,0.2666) -- (0.9818,0.2889) -- (0.9798,0.3112) -- (0.9777,0.3336) -- (0.9754,0.3559) -- (0.9730,0.3783) -- (0.9704,0.4007) -- (0.9676,0.4231) -- (0.9647,0.4455) -- (0.9617,0.4680) -- (0.9585,0.4905) -- (0.9551,0.5130) -- (0.9516,0.5356) -- (0.9479,0.5582) -- (0.9441,0.5808) -- (0.9400,0.6034) -- (0.9359,0.6261) -- (0.9315,0.6489) -- (0.9270,0.6716) -- (0.9223,0.6944) -- (0.9175,0.7173) -- (0.9124,0.7402) -- (0.9072,0.7631) -- (0.9018,0.7861) -- (0.8963,0.8091) -- (0.8905,0.8322) -- (0.8846,0.8553) -- (0.8785,0.8785) -- (0.8785,0.8785) -- (0.8523,0.8523) -- (0.8264,0.8264) -- (0.8008,0.8008) -- (0.7756,0.7756) -- (0.7506,0.7506) -- (0.7259,0.7259) -- (0.7014,0.7014) -- (0.6772,0.6772) -- (0.6532,0.6532) -- (0.6295,0.6295) -- (0.6059,0.6059) -- (0.5826,0.5826) -- (0.5594,0.5594) -- (0.5365,0.5365) -- (0.5137,0.5137) -- (0.4911,0.4911) -- (0.4687,0.4687) -- (0.4464,0.4464) -- (0.4242,0.4242) -- (0.4022,0.4022) -- (0.3803,0.3803) -- (0.3585,0.3585) -- (0.3369,0.3369) -- (0.3153,0.3153) -- (0.2939,0.2939) -- (0.2725,0.2725) -- (0.2512,0.2512) -- (0.2300,0.2300) -- (0.2089,0.2089) -- (0.1878,0.1878) -- (0.1668,0.1668) -- (0.1459,0.1459) -- (0.1249,0.1249) -- (0.1041,0.1041) -- (0.0832,0.0832) -- (0.0624,0.0624) -- (0.0416,0.0416) -- (0.0208,0.0208) -- (0.0000,0.0000) -- cycle;
\end{scope}
\draw[atomLine, line width=0.7pt] (0,0) circle (2.4);
\begin{scope}
\clip (3.2500,-2.3500) rectangle (8.0000,2.4000);
\fill[fmAmber!8] (3.5500,-2.0500) -- (3.6366,-1.9634) -- (3.7232,-1.8768) -- (3.8099,-1.7901) -- (3.8967,-1.7033) -- (3.9836,-1.6164) -- (4.0706,-1.5294) -- (4.1577,-1.4423) -- (4.2451,-1.3549) -- (4.3326,-1.2674) -- (4.4204,-1.1796) -- (4.5084,-1.0916) -- (4.5968,-1.0032) -- (4.6855,-0.9145) -- (4.7745,-0.8255) -- (4.8639,-0.7361) -- (4.9537,-0.6463) -- (5.0439,-0.5561) -- (5.1346,-0.4654) -- (5.2258,-0.3742) -- (5.3176,-0.2824) -- (5.4099,-0.1901) -- (5.5028,-0.0972) -- (5.5963,-0.0037) -- (5.6905,0.0905) -- (5.7854,0.1854) -- (5.8810,0.2810) -- (5.9774,0.3774) -- (6.0747,0.4747) -- (6.1727,0.5727) -- (6.2717,0.6717) -- (6.3716,0.7716) -- (6.4725,0.8725) -- (6.5744,0.9744) -- (6.6774,1.0774) -- (6.7815,1.1815) -- (6.8868,1.2868) -- (6.9933,1.3933) -- (7.1011,1.5011) -- (7.2103,1.6103) -- (7.2103,1.6103) -- (7.1137,1.6358) -- (7.0174,1.6605) -- (6.9213,1.6844) -- (6.8253,1.7076) -- (6.7296,1.7301) -- (6.6340,1.7518) -- (6.5387,1.7727) -- (6.4435,1.7930) -- (6.3484,1.8125) -- (6.2536,1.8313) -- (6.1589,1.8494) -- (6.0643,1.8668) -- (5.9699,1.8836) -- (5.8757,1.8996) -- (5.7815,1.9150) -- (5.6876,1.9297) -- (5.5937,1.9437) -- (5.5000,1.9571) -- (5.4064,1.9698) -- (5.3129,1.9818) -- (5.2195,1.9933) -- (5.1262,2.0040) -- (5.0329,2.0142) -- (4.9398,2.0237) -- (4.8468,2.0326) -- (4.7538,2.0408) -- (4.6610,2.0484) -- (4.5681,2.0554) -- (4.4754,2.0618) -- (4.3827,2.0676) -- (4.2900,2.0728) -- (4.1974,2.0773) -- (4.1049,2.0812) -- (4.0123,2.0846) -- (3.9198,2.0873) -- (3.8274,2.0894) -- (3.7349,2.0909) -- (3.6424,2.0918) -- (3.5500,2.0921) -- (3.5500,2.0921) -- (3.5500,1.9747) -- (3.5500,1.8581) -- (3.5500,1.7425) -- (3.5500,1.6278) -- (3.5500,1.5139) -- (3.5500,1.4008) -- (3.5500,1.2885) -- (3.5500,1.1769) -- (3.5500,1.0661) -- (3.5500,0.9560) -- (3.5500,0.8465) -- (3.5500,0.7377) -- (3.5500,0.6295) -- (3.5500,0.5219) -- (3.5500,0.4148) -- (3.5500,0.3082) -- (3.5500,0.2022) -- (3.5500,0.0966) -- (3.5500,-0.0085) -- (3.5500,-0.1132) -- (3.5500,-0.2174) -- (3.5500,-0.3213) -- (3.5500,-0.4248) -- (3.5500,-0.5280) -- (3.5500,-0.6309) -- (3.5500,-0.7335) -- (3.5500,-0.8358) -- (3.5500,-0.9378) -- (3.5500,-1.0397) -- (3.5500,-1.1413) -- (3.5500,-1.2427) -- (3.5500,-1.3440) -- (3.5500,-1.4451) -- (3.5500,-1.5461) -- (3.5500,-1.6470) -- (3.5500,-1.7478) -- (3.5500,-1.8486) -- (3.5500,-1.9493) -- (3.5500,-2.0500) -- cycle;
\fill[fmAmber!25] (3.5500,-2.0500) -- (3.6507,-2.0500) -- (3.7514,-2.0500) -- (3.8522,-2.0500) -- (3.9530,-2.0500) -- (4.0539,-2.0500) -- (4.1549,-2.0500) -- (4.2560,-2.0500) -- (4.3573,-2.0500) -- (4.4587,-2.0500) -- (4.5603,-2.0500) -- (4.6622,-2.0500) -- (4.7642,-2.0500) -- (4.8665,-2.0500) -- (4.9691,-2.0500) -- (5.0720,-2.0500) -- (5.1752,-2.0500) -- (5.2787,-2.0500) -- (5.3826,-2.0500) -- (5.4868,-2.0500) -- (5.5915,-2.0500) -- (5.6966,-2.0500) -- (5.8022,-2.0500) -- (5.9082,-2.0500) -- (6.0148,-2.0500) -- (6.1219,-2.0500) -- (6.2295,-2.0500) -- (6.3377,-2.0500) -- (6.4465,-2.0500) -- (6.5560,-2.0500) -- (6.6661,-2.0500) -- (6.7769,-2.0500) -- (6.8885,-2.0500) -- (7.0008,-2.0500) -- (7.1139,-2.0500) -- (7.2278,-2.0500) -- (7.3425,-2.0500) -- (7.4581,-2.0500) -- (7.5747,-2.0500) -- (7.6921,-2.0500) -- (7.6921,-2.0500) -- (7.6918,-1.9576) -- (7.6909,-1.8651) -- (7.6894,-1.7726) -- (7.6873,-1.6802) -- (7.6846,-1.5877) -- (7.6812,-1.4951) -- (7.6773,-1.4026) -- (7.6728,-1.3100) -- (7.6676,-1.2173) -- (7.6618,-1.1246) -- (7.6554,-1.0319) -- (7.6484,-0.9390) -- (7.6408,-0.8462) -- (7.6326,-0.7532) -- (7.6237,-0.6602) -- (7.6142,-0.5671) -- (7.6040,-0.4738) -- (7.5933,-0.3805) -- (7.5818,-0.2871) -- (7.5698,-0.1936) -- (7.5571,-0.1000) -- (7.5437,-0.0063) -- (7.5297,0.0876) -- (7.5150,0.1815) -- (7.4996,0.2757) -- (7.4836,0.3699) -- (7.4668,0.4643) -- (7.4494,0.5589) -- (7.4313,0.6536) -- (7.4125,0.7484) -- (7.3930,0.8435) -- (7.3727,0.9387) -- (7.3518,1.0340) -- (7.3301,1.1296) -- (7.3076,1.2253) -- (7.2844,1.3213) -- (7.2605,1.4174) -- (7.2358,1.5137) -- (7.2103,1.6103) -- (7.2103,1.6103) -- (7.1011,1.5011) -- (6.9933,1.3933) -- (6.8868,1.2868) -- (6.7815,1.1815) -- (6.6774,1.0774) -- (6.5744,0.9744) -- (6.4725,0.8725) -- (6.3716,0.7716) -- (6.2717,0.6717) -- (6.1727,0.5727) -- (6.0747,0.4747) -- (5.9774,0.3774) -- (5.8810,0.2810) -- (5.7854,0.1854) -- (5.6905,0.0905) -- (5.5963,-0.0037) -- (5.5028,-0.0972) -- (5.4099,-0.1901) -- (5.3176,-0.2824) -- (5.2258,-0.3742) -- (5.1346,-0.4654) -- (5.0439,-0.5561) -- (4.9537,-0.6463) -- (4.8639,-0.7361) -- (4.7745,-0.8255) -- (4.6855,-0.9145) -- (4.5968,-1.0032) -- (4.5084,-1.0916) -- (4.4204,-1.1796) -- (4.3326,-1.2674) -- (4.2451,-1.3549) -- (4.1577,-1.4423) -- (4.0706,-1.5294) -- (3.9836,-1.6164) -- (3.8967,-1.7033) -- (3.8099,-1.7901) -- (3.7232,-1.8768) -- (3.6366,-1.9634) -- (3.5500,-2.0500) -- cycle;
\draw[zone100] (3.5500,7.8534) -- (3.5500,7.5135) -- (3.5500,7.1848) -- (3.5500,6.8666) -- (3.5500,6.5580) -- (3.5500,6.2586) -- (3.5500,5.9678) -- (3.5500,5.6850) -- (3.5500,5.4097) -- (3.5500,5.1415) -- (3.5500,4.8800) -- (3.5500,4.6247) -- (3.5500,4.3753) -- (3.5500,4.1314) -- (3.5500,3.8928) -- (3.5500,3.6590) -- (3.5500,3.4298) -- (3.5500,3.2050) -- (3.5500,2.9843) -- (3.5500,2.7674) -- (3.5500,2.5541) -- (3.5500,2.3443) -- (3.5500,2.1376) -- (3.5500,1.9339) -- (3.5500,1.7331) -- (3.5500,1.5348) -- (3.5500,1.3391) -- (3.5500,1.1456) -- (3.5500,0.9543) -- (3.5500,0.7650) -- (3.5500,0.5776) -- (3.5500,0.3918) -- (3.5500,0.2076) -- (3.5500,0.0249) -- (3.5500,-0.1565) -- (3.5500,-0.3367) -- (3.5500,-0.5159) -- (3.5500,-0.6940) -- (3.5500,-0.8714) -- (3.5500,-1.0480) -- (3.5500,-1.2240) -- (3.5500,-1.3995) -- (3.5500,-1.5746) -- (3.5500,-1.7494) -- (3.5500,-1.9240) -- (3.5500,-2.0985) -- (3.5500,-2.2731) -- (3.5500,-2.4478) -- (3.5500,-2.6227) -- (3.5500,-2.7980) -- (3.5500,-2.9738) -- (3.5500,-3.1501) -- (3.5500,-3.3271) -- (3.5500,-3.5049) -- (3.5500,-3.6836) -- (3.5500,-3.8633) -- (3.5500,-3.8633) -- (3.5500,-4.0442) -- (3.5500,-4.2263) -- (3.5500,-4.4099) -- (3.5500,-4.5949) -- (3.5500,-4.7816) -- (3.5500,-4.9700) -- (3.5500,-5.1604) -- (3.5500,-5.3529) -- (3.5500,-5.5476) -- (3.5500,-5.7448) -- (3.5500,-5.9444) -- (3.5500,-6.1468) -- (3.5500,-6.3521) -- (3.5500,-6.5606) -- (3.5500,-6.7723) -- (3.5500,-6.9875) -- (3.5500,-7.2065) -- (3.5500,-7.4295) -- (3.5500,-7.6567) -- (3.5500,-7.8884) -- (3.5500,-8.1248) -- (3.5500,-8.3664) -- (3.5500,-8.6133) -- (3.5500,-8.8659) -- (3.5500,-9.1246) -- (3.5500,-9.3898) -- (3.5500,-9.6619) -- (3.5500,-9.9413) -- (3.5500,-10.2285) -- (3.5500,-10.5240) -- (3.5500,-10.8284) -- (3.5500,-11.1423) -- (3.5500,-11.4663) -- (3.5500,-11.8011);
\draw[zone100] (-6.2708,-2.0500) -- (-5.9336,-2.0500) -- (-5.6075,-2.0500) -- (-5.2916,-2.0500) -- (-4.9853,-2.0500) -- (-4.6881,-2.0500) -- (-4.3992,-2.0500) -- (-4.1182,-2.0500) -- (-3.8447,-2.0500) -- (-3.5782,-2.0500) -- (-3.3182,-2.0500) -- (-3.0643,-2.0500) -- (-2.8163,-2.0500) -- (-2.5737,-2.0500) -- (-2.3362,-2.0500) -- (-2.1036,-2.0500) -- (-1.8755,-2.0500) -- (-1.6517,-2.0500) -- (-1.4319,-2.0500) -- (-1.2159,-2.0500) -- (-1.0034,-2.0500) -- (-0.7944,-2.0500) -- (-0.5884,-2.0500) -- (-0.3855,-2.0500) -- (-0.1852,-2.0500) -- (0.0124,-2.0500) -- (0.2076,-2.0500) -- (0.4005,-2.0500) -- (0.5913,-2.0500) -- (0.7802,-2.0500) -- (0.9672,-2.0500) -- (1.1526,-2.0500) -- (1.3364,-2.0500) -- (1.5188,-2.0500) -- (1.6999,-2.0500) -- (1.8798,-2.0500) -- (2.0587,-2.0500) -- (2.2367,-2.0500) -- (2.4138,-2.0500) -- (2.5903,-2.0500) -- (2.7662,-2.0500) -- (2.9415,-2.0500) -- (3.1166,-2.0500) -- (3.2913,-2.0500) -- (3.4659,-2.0500) -- (3.6404,-2.0500) -- (3.8150,-2.0500) -- (3.9898,-2.0500) -- (4.1648,-2.0500) -- (4.3402,-2.0500) -- (4.5161,-2.0500) -- (4.6926,-2.0500) -- (4.8697,-2.0500) -- (5.0477,-2.0500) -- (5.2267,-2.0500) -- (5.4067,-2.0500) -- (5.5878,-2.0500) -- (5.7703,-2.0500) -- (5.9541,-2.0500) -- (6.1396,-2.0500) -- (6.3267,-2.0500) -- (6.5156,-2.0500) -- (6.7065,-2.0500) -- (6.8995,-2.0500) -- (7.0947,-2.0500) -- (7.2925,-2.0500) -- (7.4928,-2.0500) -- (7.6958,-2.0500) -- (7.9019,-2.0500) -- (8.1111,-2.0500) -- (8.3236,-2.0500) -- (8.5398,-2.0500) -- (8.7597,-2.0500) -- (8.9837,-2.0500) -- (9.2119,-2.0500) -- (9.4447,-2.0500) -- (9.6824,-2.0500) -- (9.9252,-2.0500) -- (10.1734,-2.0500) -- (10.4275,-2.0500) -- (10.6877,-2.0500) -- (10.9545,-2.0500) -- (11.2283,-2.0500) -- (11.5095,-2.0500) -- (11.7987,-2.0500) -- (12.0963,-2.0500) -- (12.4029,-2.0500) -- (12.7191,-2.0500) -- (13.0457,-2.0500) -- (13.3832,-2.0500);
\draw[zone100] (-3.4309,4.9309) -- (-3.1913,4.6913) -- (-2.9595,4.4595) -- (-2.7351,4.2351) -- (-2.5175,4.0175) -- (-2.3064,3.8064) -- (-2.1013,3.6013) -- (-1.9018,3.4018) -- (-1.7076,3.2076) -- (-1.5184,3.0184) -- (-1.3339,2.8339) -- (-1.1537,2.6537) -- (-0.9777,2.4777) -- (-0.8056,2.3056) -- (-0.6372,2.1372) -- (-0.4722,1.9722) -- (-0.3104,1.8104) -- (-0.1517,1.6517) -- (0.0041,1.4959) -- (0.1572,1.3428) -- (0.3078,1.1922) -- (0.4560,1.0440) -- (0.6019,0.8981) -- (0.7458,0.7542) -- (0.8876,0.6124) -- (1.0276,0.4724) -- (1.1659,0.3341) -- (1.3026,0.1974) -- (1.4377,0.0623) -- (1.5714,-0.0714) -- (1.7039,-0.2039) -- (1.8351,-0.3351) -- (1.9653,-0.4653) -- (2.0944,-0.5944) -- (2.2226,-0.7226) -- (2.3499,-0.8499) -- (2.4765,-0.9765) -- (2.6025,-1.1025) -- (2.7278,-1.2278) -- (2.8527,-1.3527) -- (2.9771,-1.4771) -- (3.1011,-1.6011) -- (3.2249,-1.7249) -- (3.3485,-1.8485) -- (3.4720,-1.9720) -- (3.5954,-2.0954) -- (3.7188,-2.2188) -- (3.8424,-2.3424) -- (3.9661,-2.4661) -- (4.0901,-2.5901) -- (4.2144,-2.7144) -- (4.3391,-2.8391) -- (4.4643,-2.9643) -- (4.5901,-3.0901) -- (4.7165,-3.2165) -- (4.8437,-3.3437) -- (4.9716,-3.4716) -- (5.1005,-3.6005) -- (5.2304,-3.7304) -- (5.3613,-3.8613) -- (5.4934,-3.9934) -- (5.6268,-4.1268) -- (5.7616,-4.2616) -- (5.8978,-4.3978) -- (6.0357,-4.5357) -- (6.1752,-4.6752) -- (6.3165,-4.8165) -- (6.4598,-4.9598) -- (6.6052,-5.1052) -- (6.7528,-5.2528) -- (6.9027,-5.4027) -- (7.0552,-5.5552) -- (7.2103,-5.7103) -- (7.2103,-5.7103) -- (7.3682,-5.8682) -- (7.5291,-6.0291) -- (7.6932,-6.1932) -- (7.8607,-6.3607) -- (8.0319,-6.5319) -- (8.2068,-6.7068) -- (8.3858,-6.8858) -- (8.5692,-7.0692) -- (8.7571,-7.2571) -- (8.9499,-7.4499) -- (9.1480,-7.6480) -- (9.3516,-7.8516) -- (9.5611,-8.0611) -- (9.7769,-8.2769) -- (9.9995,-8.4995) -- (10.2293,-8.7293) -- (10.4667,-8.9667);
\draw[zone100] (10.6119,5.0119) -- (10.3696,4.7696) -- (10.1353,4.5353) -- (9.9085,4.3085) -- (9.6887,4.0887) -- (9.4755,3.8755) -- (9.2684,3.6684) -- (9.0671,3.4671) -- (8.8712,3.2712) -- (8.6804,3.0804) -- (8.4944,2.8944) -- (8.3128,2.7128) -- (8.1355,2.5355) -- (7.9621,2.3621) -- (7.7925,2.1925) -- (7.6264,2.0264) -- (7.4636,1.8636) -- (7.3039,1.7039) -- (7.1471,1.5471) -- (6.9931,1.3931) -- (6.8417,1.2417) -- (6.6927,1.0927) -- (6.5461,0.9461) -- (6.4016,0.8016) -- (6.2591,0.6591) -- (6.1185,0.5185) -- (5.9796,0.3796) -- (5.8425,0.2425) -- (5.7068,0.1068) -- (5.5726,-0.0274) -- (5.4398,-0.1602) -- (5.3081,-0.2919) -- (5.1776,-0.4224) -- (5.0482,-0.5518) -- (4.9197,-0.6803) -- (4.7921,-0.8079) -- (4.6652,-0.9348) -- (4.5391,-1.0609) -- (4.4136,-1.1864) -- (4.2886,-1.3114) -- (4.1640,-1.4360) -- (4.0398,-1.5602) -- (3.9160,-1.6840) -- (3.7923,-1.8077) -- (3.6688,-1.9312) -- (3.5454,-2.0546) -- (3.4220,-2.1780) -- (3.2985,-2.3015) -- (3.1748,-2.4252) -- (3.0509,-2.5491) -- (2.9267,-2.6733) -- (2.8022,-2.7978) -- (2.6771,-2.9229) -- (2.5515,-3.0485) -- (2.4253,-3.1747) -- (2.4253,-3.1747) -- (2.2984,-3.3016) -- (2.1708,-3.4292) -- (2.0422,-3.5578) -- (1.9127,-3.6873) -- (1.7821,-3.8179) -- (1.6504,-3.9496) -- (1.5174,-4.0826) -- (1.3831,-4.2169) -- (1.2474,-4.3526) -- (1.1101,-4.4899) -- (0.9711,-4.6289) -- (0.8304,-4.7696) -- (0.6878,-4.9122) -- (0.5431,-5.0569) -- (0.3963,-5.2037) -- (0.2471,-5.3529) -- (0.0955,-5.5045) -- (-0.0587,-5.6587) -- (-0.2156,-5.8156) -- (-0.3756,-5.9756) -- (-0.5386,-6.1386) -- (-0.7050,-6.3050) -- (-0.8749,-6.4749) -- (-1.0485,-6.6485) -- (-1.2262,-6.8262) -- (-1.4081,-7.0081) -- (-1.5945,-7.1945) -- (-1.7856,-7.3856) -- (-1.9819,-7.5819) -- (-2.1836,-7.7836) -- (-2.3912,-7.9912) -- (-2.6049,-8.2049) -- (-2.8251,-8.4251) -- (-3.0525,-8.6525) -- (-3.2873,-8.8873);
\draw[zone110] (3.3121,7.7063) -- (3.0827,7.4598) -- (2.8644,7.2138) -- (2.6566,6.9685) -- (2.4588,6.7241) -- (2.2706,6.4808) -- (2.0914,6.2386) -- (1.9209,5.9977) -- (1.7587,5.7581) -- (1.6044,5.5199) -- (1.4576,5.2831) -- (1.3180,5.0478) -- (1.1853,4.8141) -- (1.0592,4.5818) -- (0.9394,4.3510) -- (0.8256,4.1218) -- (0.7177,3.8940) -- (0.6153,3.6677) -- (0.5182,3.4429) -- (0.4262,3.2194) -- (0.3392,2.9974) -- (0.2570,2.7767) -- (0.1794,2.5572) -- (0.1062,2.3391) -- (0.0373,2.1221) -- (-0.0275,1.9063) -- (-0.0882,1.6916) -- (-0.1450,1.4780) -- (-0.1980,1.2653) -- (-0.2474,1.0537) -- (-0.2931,0.8429) -- (-0.3353,0.6329) -- (-0.3741,0.4237) -- (-0.4095,0.2152) -- (-0.4417,0.0074) -- (-0.4706,-0.1998) -- (-0.4963,-0.4064) -- (-0.5189,-0.6126) -- (-0.5384,-0.8184) -- (-0.5549,-1.0238) -- (-0.5549,-1.0238) -- (-0.5683,-1.2289) -- (-0.5787,-1.4337) -- (-0.5861,-1.6384) -- (-0.5906,-1.8430) -- (-0.5921,-2.0474) -- (-0.5907,-2.2519) -- (-0.5863,-2.4565) -- (-0.5789,-2.6611) -- (-0.5686,-2.8660) -- (-0.5552,-3.0711) -- (-0.5389,-3.2765) -- (-0.5194,-3.4822) -- (-0.4969,-3.6884) -- (-0.4713,-3.8951) -- (-0.4424,-4.1022) -- (-0.4104,-4.3100) -- (-0.3750,-4.5185) -- (-0.3363,-4.7277) -- (-0.2942,-4.9376) -- (-0.2486,-5.1484) -- (-0.1993,-5.3600) -- (-0.1464,-5.5727) -- (-0.0897,-5.7863) -- (-0.0290,-6.0009) -- (0.0356,-6.2167) -- (0.1044,-6.4336) -- (0.1775,-6.6518) -- (0.2550,-6.8712) -- (0.3371,-7.0918) -- (0.4240,-7.3139) -- (0.5158,-7.5373) -- (0.6128,-7.7621) -- (0.7150,-7.9883) -- (0.8229,-8.2161) -- (0.9365,-8.4453) -- (1.0562,-8.6760) -- (1.1821,-8.9082) -- (1.3147,-9.1420) -- (1.4541,-9.3772) -- (1.6007,-9.6139) -- (1.7548,-9.8521) -- (1.9168,-10.0917) -- (2.0871,-10.3326) -- (2.2660,-10.5747) -- (2.4540,-10.8180) -- (2.6515,-11.0624) -- (2.8591,-11.3077) -- (3.0771,-11.5536) -- (3.3062,-11.8001) -- (3.5469,-12.0469);
\draw[zone110] (3.6400,-11.9592) -- (3.8765,-11.7125) -- (4.1016,-11.4662) -- (4.3158,-11.2204) -- (4.5198,-10.9755) -- (4.7139,-10.7315) -- (4.8986,-10.4885) -- (5.0744,-10.2468) -- (5.2417,-10.0064) -- (5.4009,-9.7673) -- (5.5523,-9.5297) -- (5.6963,-9.2935) -- (5.8332,-9.0588) -- (5.9634,-8.8256) -- (6.0871,-8.5939) -- (6.2045,-8.3637) -- (6.3161,-8.1350) -- (6.4219,-7.9078) -- (6.5223,-7.6820) -- (6.6174,-7.4577) -- (6.7074,-7.2348) -- (6.7926,-7.0133) -- (6.8730,-6.7931) -- (6.9490,-6.5741) -- (7.0205,-6.3564) -- (7.0878,-6.1399) -- (7.1510,-5.9245) -- (7.2103,-5.7103) -- (7.2656,-5.4970) -- (7.3172,-5.2847) -- (7.3652,-5.0734) -- (7.4096,-4.8629) -- (7.4505,-4.6533) -- (7.4880,-4.4443) -- (7.5221,-4.2361) -- (7.5530,-4.0286) -- (7.5807,-3.8216) -- (7.6053,-3.6151) -- (7.6267,-3.4091) -- (7.6450,-3.2035) -- (7.6603,-2.9982) -- (7.6726,-2.7932) -- (7.6819,-2.5884) -- (7.6882,-2.3838) -- (7.6915,-2.1793) -- (7.6919,-1.9748) -- (7.6894,-1.7703) -- (7.6838,-1.5657) -- (7.6753,-1.3610) -- (7.6639,-1.1560) -- (7.6494,-0.9508) -- (7.6318,-0.7453) -- (7.6112,-0.5394) -- (7.5875,-0.3331) -- (7.5607,-0.1262) -- (7.5306,0.0812) -- (7.4973,0.2892) -- (7.4607,0.4980) -- (7.4207,0.7074) -- (7.3773,0.9177) -- (7.3303,1.1288) -- (7.2796,1.3408) -- (7.2253,1.5538) -- (7.1671,1.7678) -- (7.1049,1.9829) -- (7.0387,2.1991) -- (6.9683,2.4164) -- (6.8935,2.6351) -- (6.8143,2.8549) -- (6.7304,3.0761) -- (6.6417,3.2987) -- (6.5479,3.5226) -- (6.4490,3.7480) -- (6.3446,3.9748) -- (6.2346,4.2031) -- (6.1187,4.4329) -- (5.9967,4.6641) -- (5.8683,4.8970) -- (5.7332,5.1313) -- (5.5911,5.3671) -- (5.4416,5.6043) -- (5.2846,5.8430) -- (5.1195,6.0831) -- (4.9459,6.3245) -- (4.7636,6.5671) -- (4.5720,6.8109) -- (4.3707,7.0556) -- (4.1593,7.3011) -- (3.9371,7.5473) -- (3.7036,7.7940);
\draw[zone110] (-6.2940,-2.2036) -- (-6.0473,-2.4371) -- (-5.8011,-2.6593) -- (-5.5556,-2.8707) -- (-5.3109,-3.0720) -- (-5.0671,-3.2636) -- (-4.8245,-3.4459) -- (-4.5831,-3.6195) -- (-4.3430,-3.7846) -- (-4.1043,-3.9416) -- (-3.8671,-4.0911) -- (-3.6313,-4.2332) -- (-3.3970,-4.3683) -- (-3.1641,-4.4967) -- (-2.9329,-4.6187) -- (-2.7031,-4.7346) -- (-2.4748,-4.8446) -- (-2.2480,-4.9490) -- (-2.0226,-5.0479) -- (-1.7987,-5.1417) -- (-1.5761,-5.2304) -- (-1.3549,-5.3143) -- (-1.1351,-5.3935) -- (-0.9164,-5.4683) -- (-0.6991,-5.5387) -- (-0.4829,-5.6049) -- (-0.2678,-5.6671) -- (-0.0538,-5.7253) -- (0.1592,-5.7796) -- (0.3712,-5.8303) -- (0.5823,-5.8773) -- (0.7926,-5.9207) -- (1.0020,-5.9607) -- (1.2108,-5.9973) -- (1.4188,-6.0306) -- (1.6262,-6.0607) -- (1.8331,-6.0875) -- (2.0394,-6.1112) -- (2.2453,-6.1318) -- (2.4508,-6.1494) -- (2.6560,-6.1639) -- (2.8610,-6.1753) -- (3.0657,-6.1838) -- (3.2703,-6.1894) -- (3.4748,-6.1919) -- (3.6793,-6.1915) -- (3.8838,-6.1882) -- (4.0884,-6.1819) -- (4.2932,-6.1726) -- (4.4982,-6.1603) -- (4.7035,-6.1450) -- (4.9091,-6.1267) -- (5.1151,-6.1053) -- (5.3216,-6.0807) -- (5.5286,-6.0530) -- (5.7361,-6.0221) -- (5.9443,-5.9880) -- (6.1533,-5.9505) -- (6.3629,-5.9096) -- (6.5734,-5.8652) -- (6.7847,-5.8172) -- (6.9970,-5.7656) -- (7.2103,-5.7103) -- (7.4245,-5.6510) -- (7.6399,-5.5878) -- (7.8564,-5.5205) -- (8.0741,-5.4490) -- (8.2931,-5.3730) -- (8.5133,-5.2926) -- (8.7348,-5.2074) -- (8.9577,-5.1174) -- (9.1820,-5.0223) -- (9.4078,-4.9219) -- (9.6350,-4.8161) -- (9.8637,-4.7045) -- (10.0939,-4.5871) -- (10.3256,-4.4634) -- (10.5588,-4.3332) -- (10.7935,-4.1963) -- (11.0297,-4.0523) -- (11.2673,-3.9009) -- (11.5064,-3.7417) -- (11.7468,-3.5744) -- (11.9885,-3.3986) -- (12.2315,-3.2139) -- (12.4755,-3.0198) -- (12.7204,-2.8158) -- (12.9662,-2.6016) -- (13.2125,-2.3765) -- (13.4592,-2.1400);
\draw[zone110] (13.3938,-1.8962) -- (13.1471,-1.6627) -- (12.9009,-1.4406) -- (12.6554,-1.2291) -- (12.4107,-1.0278) -- (12.1669,-0.8362) -- (11.9243,-0.6539) -- (11.6829,-0.4804) -- (11.4428,-0.3153) -- (11.2041,-0.1582) -- (10.9669,-0.0088) -- (10.7311,0.1333) -- (10.4968,0.2684) -- (10.2640,0.3968) -- (10.0327,0.5188) -- (9.8029,0.6347) -- (9.5746,0.7447) -- (9.3478,0.8491) -- (9.1224,0.9480) -- (8.8985,1.0417) -- (8.6759,1.1305) -- (8.4547,1.2144) -- (8.2349,1.2936) -- (8.0163,1.3684) -- (7.7989,1.4388) -- (7.5827,1.5050) -- (7.3676,1.5671) -- (7.1536,1.6253) -- (6.9406,1.6797) -- (6.7286,1.7303) -- (6.5175,1.7773) -- (6.3072,1.8207) -- (6.0978,1.8607) -- (5.8891,1.8974) -- (5.6810,1.9307) -- (5.4736,1.9607) -- (5.2668,1.9875) -- (5.0604,2.0112) -- (4.8545,2.0318) -- (4.6490,2.0494) -- (4.4438,2.0639) -- (4.2389,2.0753) -- (4.0341,2.0838) -- (3.8295,2.0894) -- (3.6250,2.0919) -- (3.4206,2.0915) -- (3.2160,2.0882) -- (3.0114,2.0819) -- (2.8066,2.0726) -- (2.6016,2.0603) -- (2.3964,2.0450) -- (2.1907,2.0267) -- (1.9847,2.0052) -- (1.7783,1.9807) -- (1.5713,1.9530) -- (1.3637,1.9221) -- (1.1555,1.8879) -- (0.9466,1.8504) -- (0.7369,1.8095) -- (0.5264,1.7651) -- (0.3151,1.7172) -- (0.1028,1.6656) -- (-0.1104,1.6102) -- (-0.3247,1.5510) -- (-0.5401,1.4878) -- (-0.7566,1.4205) -- (-0.9743,1.3489) -- (-1.1932,1.2730) -- (-1.4135,1.1925) -- (-1.6350,1.1073) -- (-1.8579,1.0173) -- (-2.0822,0.9222) -- (-2.3080,0.8218) -- (-2.5352,0.7160) -- (-2.5352,0.7160) -- (-2.7639,0.6044) -- (-2.9941,0.4870) -- (-3.2258,0.3633) -- (-3.4590,0.2331) -- (-3.6937,0.0962) -- (-3.9299,-0.0479) -- (-4.1675,-0.1993) -- (-4.4066,-0.3584) -- (-4.6470,-0.5257) -- (-4.8888,-0.7015) -- (-5.1317,-0.8863) -- (-5.3757,-1.0804) -- (-5.6206,-1.2843) -- (-5.8664,-1.4986) -- (-6.1127,-1.7237) -- (-6.3594,-1.9602);
\draw[fmAmber!80!black, line width=0.8pt] (3.5500,-2.0500) -- (3.6507,-2.0500) -- (3.7514,-2.0500) -- (3.8522,-2.0500) -- (3.9530,-2.0500) -- (4.0539,-2.0500) -- (4.1549,-2.0500) -- (4.2560,-2.0500) -- (4.3573,-2.0500) -- (4.4587,-2.0500) -- (4.5603,-2.0500) -- (4.6622,-2.0500) -- (4.7642,-2.0500) -- (4.8665,-2.0500) -- (4.9691,-2.0500) -- (5.0720,-2.0500) -- (5.1752,-2.0500) -- (5.2787,-2.0500) -- (5.3826,-2.0500) -- (5.4868,-2.0500) -- (5.5915,-2.0500) -- (5.6966,-2.0500) -- (5.8022,-2.0500) -- (5.9082,-2.0500) -- (6.0148,-2.0500) -- (6.1219,-2.0500) -- (6.2295,-2.0500) -- (6.3377,-2.0500) -- (6.4465,-2.0500) -- (6.5560,-2.0500) -- (6.6661,-2.0500) -- (6.7769,-2.0500) -- (6.8885,-2.0500) -- (7.0008,-2.0500) -- (7.1139,-2.0500) -- (7.2278,-2.0500) -- (7.3425,-2.0500) -- (7.4581,-2.0500) -- (7.5747,-2.0500) -- (7.6921,-2.0500) -- (7.6921,-2.0500) -- (7.6918,-1.9576) -- (7.6909,-1.8651) -- (7.6894,-1.7726) -- (7.6873,-1.6802) -- (7.6846,-1.5877) -- (7.6812,-1.4951) -- (7.6773,-1.4026) -- (7.6728,-1.3100) -- (7.6676,-1.2173) -- (7.6618,-1.1246) -- (7.6554,-1.0319) -- (7.6484,-0.9390) -- (7.6408,-0.8462) -- (7.6326,-0.7532) -- (7.6237,-0.6602) -- (7.6142,-0.5671) -- (7.6040,-0.4738) -- (7.5933,-0.3805) -- (7.5818,-0.2871) -- (7.5698,-0.1936) -- (7.5571,-0.1000) -- (7.5437,-0.0063) -- (7.5297,0.0876) -- (7.5150,0.1815) -- (7.4996,0.2757) -- (7.4836,0.3699) -- (7.4668,0.4643) -- (7.4494,0.5589) -- (7.4313,0.6536) -- (7.4125,0.7484) -- (7.3930,0.8435) -- (7.3727,0.9387) -- (7.3518,1.0340) -- (7.3301,1.1296) -- (7.3076,1.2253) -- (7.2844,1.3213) -- (7.2605,1.4174) -- (7.2358,1.5137) -- (7.2103,1.6103) -- (7.2103,1.6103) -- (7.1011,1.5011) -- (6.9933,1.3933) -- (6.8868,1.2868) -- (6.7815,1.1815) -- (6.6774,1.0774) -- (6.5744,0.9744) -- (6.4725,0.8725) -- (6.3716,0.7716) -- (6.2717,0.6717) -- (6.1727,0.5727) -- (6.0747,0.4747) -- (5.9774,0.3774) -- (5.8810,0.2810) -- (5.7854,0.1854) -- (5.6905,0.0905) -- (5.5963,-0.0037) -- (5.5028,-0.0972) -- (5.4099,-0.1901) -- (5.3176,-0.2824) -- (5.2258,-0.3742) -- (5.1346,-0.4654) -- (5.0439,-0.5561) -- (4.9537,-0.6463) -- (4.8639,-0.7361) -- (4.7745,-0.8255) -- (4.6855,-0.9145) -- (4.5968,-1.0032) -- (4.5084,-1.0916) -- (4.4204,-1.1796) -- (4.3326,-1.2674) -- (4.2451,-1.3549) -- (4.1577,-1.4423) -- (4.0706,-1.5294) -- (3.9836,-1.6164) -- (3.8967,-1.7033) -- (3.8099,-1.7901) -- (3.7232,-1.8768) -- (3.6366,-1.9634) -- (3.5500,-2.0500) -- cycle;
\end{scope}
\draw[atomLine, line width=0.5pt] (3.2500,-2.3500) rectangle (8.0000,2.4000);
\draw[atomLine!70, line width=0.4pt] (-0.0720,-0.0720) rectangle (1.0680,1.0680);
\draw[atomLine!50, densely dotted] (1.0680,1.0680) -- (3.2500,2.4000) (1.0680,-0.0720) -- (3.2500,-2.3500);
\node[pole100] at (0.0000,0.0000) {};
\node[pole100] at (0.0000,2.4000) {};
\node[pole100] at (2.4000,0.0000) {};
\node[pole110] at (0.0000,0.9941) {};
\node[pole110] at (0.0000,-0.9941) {};
\node[pole110] at (0.9941,0.0000) {};
\node[pole110] at (-0.9941,0.0000) {};
\node[pole110] at (1.6971,1.6971) {};
\node[pole110] at (1.6971,-1.6971) {};
\node[pole111] at (0.8785,0.8785) {};
\node[pole111] at (0.8785,-0.8785) {};
\node[pole111] at (-0.8785,0.8785) {};
\node[pole111] at (-0.8785,-0.8785) {};
\node[pole100] (B100) at (3.5500,-2.0500) {};
\node[pole110] (B110) at (7.6921,-2.0500) {};
\node[pole110] (B101) at (3.5500,2.0921) {};
\node[pole111] (B111) at (7.2103,1.6103) {};
\node[poleH] (B211) at (5.7974,0.1974) {};
\draw[fmRed, line width=1pt, -{Stealth[length=2mm]}] (6.2982,-1.1339) -- (6.2971,-1.1322) -- (6.2959,-1.1305) -- (6.2947,-1.1287) -- (6.2935,-1.1270) -- (6.2924,-1.1253) -- (6.2912,-1.1235) -- (6.2900,-1.1218) -- (6.2888,-1.1201) -- (6.2877,-1.1184) -- (6.2865,-1.1166) -- (6.2853,-1.1149) -- (6.2841,-1.1132) -- (6.2830,-1.1114) -- (6.2818,-1.1097) -- (6.2806,-1.1080) -- (6.2794,-1.1063) -- (6.2783,-1.1045) -- (6.2771,-1.1028) -- (6.2759,-1.1011) -- (6.2747,-1.0994) -- (6.2735,-1.0976) -- (6.2724,-1.0959) -- (6.2712,-1.0942) -- (6.2700,-1.0925) -- (6.2688,-1.0907) -- (6.2677,-1.0890) -- (6.2665,-1.0873) -- (6.2653,-1.0856) -- (6.2641,-1.0838) -- (6.2629,-1.0821) -- (6.2618,-1.0804) -- (6.2606,-1.0787) -- (6.2594,-1.0769) -- (6.2582,-1.0752) -- (6.2571,-1.0735) -- (6.2559,-1.0718) -- (6.2547,-1.0700) -- (6.2535,-1.0683) -- (6.2523,-1.0666) -- (6.2512,-1.0649) -- (6.2500,-1.0631) -- (6.2488,-1.0614) -- (6.2476,-1.0597) -- (6.2464,-1.0580) -- (6.2452,-1.0563) -- (6.2441,-1.0545) -- (6.2429,-1.0528) -- (6.2417,-1.0511) -- (6.2405,-1.0494) -- (6.2393,-1.0477) -- (6.2382,-1.0459) -- (6.2370,-1.0442) -- (6.2358,-1.0425) -- (6.2346,-1.0408) -- (6.2334,-1.0391) -- (6.2323,-1.0373) -- (6.2311,-1.0356) -- (6.2299,-1.0339) -- (6.2287,-1.0322) -- (6.2275,-1.0305) -- (6.2263,-1.0287) -- (6.2252,-1.0270) -- (6.2240,-1.0253) -- (6.2228,-1.0236) -- (6.2216,-1.0219) -- (6.2204,-1.0201) -- (6.2192,-1.0184) -- (6.2180,-1.0167) -- (6.2169,-1.0150) -- (6.2157,-1.0133) -- (6.2145,-1.0116) -- (6.2133,-1.0098) -- (6.2121,-1.0081) -- (6.2109,-1.0064) -- (6.2098,-1.0047) -- (6.2086,-1.0030) -- (6.2074,-1.0013) -- (6.2062,-0.9995) -- (6.2050,-0.9978) -- (6.2038,-0.9961) -- (6.2026,-0.9944) -- (6.2014,-0.9927) -- (6.2003,-0.9910) -- (6.1991,-0.9892) -- (6.1979,-0.9875) -- (6.1967,-0.9858) -- (6.1955,-0.9841) -- (6.1943,-0.9824) -- (6.1931,-0.9807) -- (6.1920,-0.9790) -- (6.1908,-0.9772) -- (6.1896,-0.9755) -- (6.1884,-0.9738) -- (6.1872,-0.9721) -- (6.1860,-0.9704) -- (6.1848,-0.9687) -- (6.1836,-0.9670) -- (6.1824,-0.9652) -- (6.1813,-0.9635) -- (6.1801,-0.9618) -- (6.1789,-0.9601) -- (6.1777,-0.9584) -- (6.1765,-0.9567) -- (6.1753,-0.9550) -- (6.1741,-0.9533) -- (6.1729,-0.9515) -- (6.1717,-0.9498) -- (6.1705,-0.9481) -- (6.1694,-0.9464) -- (6.1682,-0.9447) -- (6.1670,-0.9430) -- (6.1658,-0.9413) -- (6.1646,-0.9396) -- (6.1634,-0.9379) -- (6.1622,-0.9362) -- (6.1610,-0.9344) -- (6.1598,-0.9327) -- (6.1586,-0.9310) -- (6.1574,-0.9293) -- (6.1562,-0.9276) -- (6.1551,-0.9259) -- (6.1539,-0.9242) -- (6.1527,-0.9225) -- (6.1515,-0.9208) -- (6.1503,-0.9191) -- (6.1491,-0.9174) -- (6.1479,-0.9156) -- (6.1467,-0.9139) -- (6.1455,-0.9122) -- (6.1443,-0.9105) -- (6.1431,-0.9088) -- (6.1419,-0.9071) -- (6.1407,-0.9054) -- (6.1395,-0.9037) -- (6.1383,-0.9020) -- (6.1371,-0.9003) -- (6.1359,-0.8986) -- (6.1348,-0.8969) -- (6.1336,-0.8952) -- (6.1324,-0.8935) -- (6.1312,-0.8918) -- (6.1300,-0.8900) -- (6.1288,-0.8883) -- (6.1276,-0.8866) -- (6.1264,-0.8849) -- (6.1252,-0.8832) -- (6.1240,-0.8815) -- (6.1228,-0.8798) -- (6.1216,-0.8781) -- (6.1204,-0.8764) -- (6.1192,-0.8747) -- (6.1180,-0.8730) -- (6.1168,-0.8713) -- (6.1156,-0.8696) -- (6.1144,-0.8679) -- (6.1132,-0.8662) -- (6.1120,-0.8645) -- (6.1108,-0.8628) -- (6.1096,-0.8611) -- (6.1084,-0.8594) -- (6.1072,-0.8577) -- (6.1060,-0.8560) -- (6.1048,-0.8543) -- (6.1036,-0.8526) -- (6.1024,-0.8509) -- (6.1012,-0.8492) -- (6.1000,-0.8475) -- (6.0988,-0.8458) -- (6.0976,-0.8441) -- (6.0964,-0.8424) -- (6.0952,-0.8407) -- (6.0940,-0.8390) -- (6.0928,-0.8373) -- (6.0916,-0.8356) -- (6.0904,-0.8339) -- (6.0892,-0.8322) -- (6.0880,-0.8305) -- (6.0868,-0.8288) -- (6.0856,-0.8271) -- (6.0844,-0.8254) -- (6.0832,-0.8237) -- (6.0820,-0.8220) -- (6.0808,-0.8203) -- (6.0796,-0.8186) -- (6.0784,-0.8169) -- (6.0772,-0.8152) -- (6.0760,-0.8135) -- (6.0748,-0.8118) -- (6.0736,-0.8101) -- (6.0724,-0.8084) -- (6.0712,-0.8067) -- (6.0700,-0.8050) -- (6.0688,-0.8033) -- (6.0676,-0.8016) -- (6.0664,-0.7999) -- (6.0652,-0.7982) -- (6.0639,-0.7965) -- (6.0627,-0.7948) -- (6.0615,-0.7931) -- (6.0603,-0.7914) -- (6.0591,-0.7897) -- (6.0579,-0.7880) -- (6.0567,-0.7863) -- (6.0555,-0.7846) -- (6.0543,-0.7829) -- (6.0531,-0.7812) -- (6.0519,-0.7795) -- (6.0507,-0.7778) -- (6.0495,-0.7761) -- (6.0483,-0.7744) -- (6.0471,-0.7728) -- (6.0459,-0.7711) -- (6.0446,-0.7694) -- (6.0434,-0.7677) -- (6.0422,-0.7660) -- (6.0410,-0.7643) -- (6.0398,-0.7626) -- (6.0386,-0.7609) -- (6.0374,-0.7592) -- (6.0362,-0.7575) -- (6.0350,-0.7558) -- (6.0338,-0.7541) -- (6.0326,-0.7524) -- (6.0314,-0.7507) -- (6.0301,-0.7490) -- (6.0289,-0.7474) -- (6.0277,-0.7457) -- (6.0265,-0.7440) -- (6.0253,-0.7423) -- (6.0241,-0.7406) -- (6.0229,-0.7389) -- (6.0217,-0.7372) -- (6.0205,-0.7355) -- (6.0193,-0.7338) -- (6.0180,-0.7321) -- (6.0168,-0.7304) -- (6.0156,-0.7287) -- (6.0144,-0.7271) -- (6.0132,-0.7254) -- (6.0120,-0.7237) -- (6.0108,-0.7220) -- (6.0096,-0.7203) -- (6.0084,-0.7186) -- (6.0071,-0.7169) -- (6.0059,-0.7152) -- (6.0047,-0.7135) -- (6.0035,-0.7118) -- (6.0023,-0.7102) -- (6.0011,-0.7085) -- (5.9999,-0.7068) -- (5.9986,-0.7051) -- (5.9974,-0.7034) -- (5.9962,-0.7017) -- (5.9950,-0.7000) -- (5.9938,-0.6983) -- (5.9926,-0.6967) -- (5.9914,-0.6950) -- (5.9901,-0.6933) -- (5.9889,-0.6916) -- (5.9877,-0.6899) -- (5.9865,-0.6882) -- (5.9853,-0.6865) -- (5.9841,-0.6848) -- (5.9829,-0.6832) -- (5.9816,-0.6815) -- (5.9804,-0.6798) -- (5.9792,-0.6781) -- (5.9780,-0.6764) -- (5.9768,-0.6747) -- (5.9756,-0.6730) -- (5.9743,-0.6714) -- (5.9731,-0.6697) -- (5.9719,-0.6680) -- (5.9707,-0.6663) -- (5.9695,-0.6646) -- (5.9683,-0.6629) -- (5.9670,-0.6612) -- (5.9658,-0.6596) -- (5.9646,-0.6579) -- (5.9634,-0.6562) -- (5.9622,-0.6545) -- (5.9609,-0.6528) -- (5.9597,-0.6511) -- (5.9585,-0.6495) -- (5.9573,-0.6478) -- (5.9561,-0.6461) -- (5.9548,-0.6444) -- (5.9536,-0.6427) -- (5.9524,-0.6410) -- (5.9512,-0.6394) -- (5.9500,-0.6377) -- (5.9487,-0.6360) -- (5.9475,-0.6343) -- (5.9463,-0.6326) -- (5.9451,-0.6310) -- (5.9439,-0.6293) -- (5.9426,-0.6276) -- (5.9414,-0.6259) -- (5.9402,-0.6242) -- (5.9390,-0.6225) -- (5.9378,-0.6209) -- (5.9365,-0.6192) -- (5.9353,-0.6175) -- (5.9341,-0.6158) -- (5.9329,-0.6141) -- (5.9316,-0.6125) -- (5.9304,-0.6108) -- (5.9292,-0.6091) -- (5.9280,-0.6074) -- (5.9268,-0.6057) -- (5.9255,-0.6041) -- (5.9243,-0.6024) -- (5.9231,-0.6007) -- (5.9219,-0.5990) -- (5.9206,-0.5973) -- (5.9194,-0.5957) -- (5.9182,-0.5940) -- (5.9170,-0.5923) -- (5.9157,-0.5906) -- (5.9145,-0.5890) -- (5.9133,-0.5873) -- (5.9121,-0.5856) -- (5.9108,-0.5839) -- (5.9096,-0.5822) -- (5.9084,-0.5806) -- (5.9072,-0.5789) -- (5.9059,-0.5772) -- (5.9047,-0.5755) -- (5.9035,-0.5739) -- (5.9022,-0.5722) -- (5.9010,-0.5705) -- (5.8998,-0.5688) -- (5.8986,-0.5672) -- (5.8973,-0.5655) -- (5.8961,-0.5638) -- (5.8949,-0.5621) -- (5.8937,-0.5604) -- (5.8924,-0.5588) -- (5.8912,-0.5571) -- (5.8900,-0.5554) -- (5.8887,-0.5537) -- (5.8875,-0.5521) -- (5.8863,-0.5504) -- (5.8851,-0.5487) -- (5.8838,-0.5470) -- (5.8826,-0.5454) -- (5.8814,-0.5437) -- (5.8801,-0.5420) -- (5.8789,-0.5404) -- (5.8777,-0.5387) -- (5.8764,-0.5370) -- (5.8752,-0.5353) -- (5.8740,-0.5337) -- (5.8728,-0.5320) -- (5.8715,-0.5303) -- (5.8703,-0.5286) -- (5.8691,-0.5270) -- (5.8678,-0.5253) -- (5.8666,-0.5236) -- (5.8654,-0.5219) -- (5.8641,-0.5203) -- (5.8629,-0.5186) -- (5.8617,-0.5169) -- (5.8604,-0.5153) -- (5.8592,-0.5136) -- (5.8580,-0.5119) -- (5.8567,-0.5102) -- (5.8555,-0.5086) -- (5.8543,-0.5069) -- (5.8530,-0.5052) -- (5.8518,-0.5036) -- (5.8506,-0.5019) -- (5.8493,-0.5002) -- (5.8481,-0.4985) -- (5.8469,-0.4969) -- (5.8456,-0.4952) -- (5.8444,-0.4935) -- (5.8432,-0.4919) -- (5.8419,-0.4902) -- (5.8407,-0.4885) -- (5.8395,-0.4869) -- (5.8382,-0.4852) -- (5.8370,-0.4835) -- (5.8357,-0.4819) -- (5.8345,-0.4802) -- (5.8333,-0.4785) -- (5.8320,-0.4768) -- (5.8308,-0.4752) -- (5.8296,-0.4735) -- (5.8283,-0.4718) -- (5.8271,-0.4702) -- (5.8259,-0.4685) -- (5.8246,-0.4668) -- (5.8234,-0.4652) -- (5.8221,-0.4635) -- (5.8209,-0.4618) -- (5.8197,-0.4602) -- (5.8184,-0.4585) -- (5.8172,-0.4568) -- (5.8159,-0.4552) -- (5.8147,-0.4535) -- (5.8135,-0.4518) -- (5.8122,-0.4502) -- (5.8110,-0.4485) -- (5.8098,-0.4468) -- (5.8085,-0.4452) -- (5.8073,-0.4435) -- (5.8060,-0.4418) -- (5.8048,-0.4402) -- (5.8035,-0.4385) -- (5.8023,-0.4368) -- (5.8011,-0.4352) -- (5.7998,-0.4335) -- (5.7986,-0.4318) -- (5.7973,-0.4302) -- (5.7961,-0.4285) -- (5.7949,-0.4268) -- (5.7936,-0.4252) -- (5.7924,-0.4235) -- (5.7911,-0.4219) -- (5.7899,-0.4202) -- (5.7886,-0.4185) -- (5.7874,-0.4169) -- (5.7862,-0.4152) -- (5.7849,-0.4135) -- (5.7837,-0.4119) -- (5.7824,-0.4102) -- (5.7812,-0.4085) -- (5.7799,-0.4069) -- (5.7787,-0.4052) -- (5.7775,-0.4036) -- (5.7762,-0.4019) -- (5.7750,-0.4002) -- (5.7737,-0.3986) -- (5.7725,-0.3969) -- (5.7712,-0.3952) -- (5.7700,-0.3936) -- (5.7687,-0.3919) -- (5.7675,-0.3903) -- (5.7663,-0.3886) -- (5.7650,-0.3869) -- (5.7638,-0.3853) -- (5.7625,-0.3836) -- (5.7613,-0.3820) -- (5.7600,-0.3803) -- (5.7588,-0.3786) -- (5.7575,-0.3770) -- (5.7563,-0.3753) -- (5.7550,-0.3737) -- (5.7538,-0.3720) -- (5.7525,-0.3703) -- (5.7513,-0.3687) -- (5.7500,-0.3670) -- (5.7488,-0.3654) -- (5.7475,-0.3637) -- (5.7463,-0.3620) -- (5.7450,-0.3604) -- (5.7438,-0.3587) -- (5.7425,-0.3571) -- (5.7413,-0.3554) -- (5.7400,-0.3537) -- (5.7388,-0.3521) -- (5.7375,-0.3504) -- (5.7363,-0.3488) -- (5.7350,-0.3471) -- (5.7338,-0.3454) -- (5.7325,-0.3438) -- (5.7313,-0.3421) -- (5.7300,-0.3405) -- (5.7288,-0.3388) -- (5.7275,-0.3372) -- (5.7263,-0.3355) -- (5.7250,-0.3338) -- (5.7238,-0.3322) -- (5.7225,-0.3305) -- (5.7213,-0.3289) -- (5.7200,-0.3272) -- (5.7188,-0.3256) -- (5.7175,-0.3239) -- (5.7163,-0.3222) -- (5.7150,-0.3206) -- (5.7138,-0.3189) -- (5.7125,-0.3173) -- (5.7113,-0.3156) -- (5.7100,-0.3140) -- (5.7087,-0.3123) -- (5.7075,-0.3107) -- (5.7062,-0.3090) -- (5.7050,-0.3073) -- (5.7037,-0.3057) -- (5.7025,-0.3040) -- (5.7012,-0.3024) -- (5.7000,-0.3007) -- (5.6987,-0.2991) -- (5.6975,-0.2974) -- (5.6962,-0.2958) -- (5.6949,-0.2941) -- (5.6937,-0.2925) -- (5.6924,-0.2908) -- (5.6912,-0.2892) -- (5.6899,-0.2875) -- (5.6887,-0.2858) -- (5.6874,-0.2842) -- (5.6861,-0.2825) -- (5.6849,-0.2809) -- (5.6836,-0.2792) -- (5.6824,-0.2776) -- (5.6811,-0.2759) -- (5.6799,-0.2743) -- (5.6786,-0.2726) -- (5.6773,-0.2710) -- (5.6761,-0.2693) -- (5.6748,-0.2677) -- (5.6736,-0.2660) -- (5.6723,-0.2644) -- (5.6710,-0.2627) -- (5.6698,-0.2611) -- (5.6685,-0.2594) -- (5.6673,-0.2578) -- (5.6660,-0.2561) -- (5.6647,-0.2545) -- (5.6635,-0.2528) -- (5.6622,-0.2512) -- (5.6610,-0.2495) -- (5.6597,-0.2479) -- (5.6584,-0.2462) -- (5.6572,-0.2446) -- (5.6559,-0.2429) -- (5.6546,-0.2413) -- (5.6534,-0.2396) -- (5.6521,-0.2380) -- (5.6509,-0.2363) -- (5.6496,-0.2347) -- (5.6483,-0.2330) -- (5.6471,-0.2314) -- (5.6458,-0.2297) -- (5.6445,-0.2281) -- (5.6433,-0.2264) -- (5.6420,-0.2248) -- (5.6408,-0.2231) -- (5.6395,-0.2215) -- (5.6382,-0.2198) -- (5.6370,-0.2182) -- (5.6357,-0.2165) -- (5.6344,-0.2149) -- (5.6332,-0.2132) -- (5.6319,-0.2116) -- (5.6306,-0.2099) -- (5.6294,-0.2083) -- (5.6281,-0.2066) -- (5.6268,-0.2050) -- (5.6256,-0.2033) -- (5.6243,-0.2017) -- (5.6230,-0.2000) -- (5.6218,-0.1984) -- (5.6205,-0.1967) -- (5.6192,-0.1951) -- (5.6180,-0.1935) -- (5.6167,-0.1918) -- (5.6154,-0.1902) -- (5.6142,-0.1885) -- (5.6129,-0.1869) -- (5.6116,-0.1852) -- (5.6104,-0.1836) -- (5.6091,-0.1819) -- (5.6078,-0.1803) -- (5.6066,-0.1786) -- (5.6053,-0.1770) -- (5.6040,-0.1754) -- (5.6027,-0.1737) -- (5.6015,-0.1721) -- (5.6002,-0.1704) -- (5.5989,-0.1688) -- (5.5977,-0.1671) -- (5.5964,-0.1655) -- (5.5951,-0.1638) -- (5.5939,-0.1622) -- (5.5926,-0.1606) -- (5.5913,-0.1589) -- (5.5900,-0.1573) -- (5.5888,-0.1556) -- (5.5875,-0.1540) -- (5.5862,-0.1523) -- (5.5850,-0.1507) -- (5.5837,-0.1490) -- (5.5824,-0.1474) -- (5.5811,-0.1458) -- (5.5799,-0.1441) -- (5.5786,-0.1425) -- (5.5773,-0.1408) -- (5.5760,-0.1392) -- (5.5748,-0.1375) -- (5.5735,-0.1359) -- (5.5722,-0.1343) -- (5.5709,-0.1326) -- (5.5697,-0.1310) -- (5.5684,-0.1293) -- (5.5671,-0.1277) -- (5.5658,-0.1261) -- (5.5646,-0.1244) -- (5.5633,-0.1228) -- (5.5620,-0.1211) -- (5.5607,-0.1195) -- (5.5595,-0.1179) -- (5.5582,-0.1162) -- (5.5569,-0.1146) -- (5.5556,-0.1129) -- (5.5544,-0.1113) -- (5.5531,-0.1096) -- (5.5518,-0.1080) -- (5.5505,-0.1064) -- (5.5493,-0.1047) -- (5.5480,-0.1031) -- (5.5467,-0.1014) -- (5.5454,-0.0998) -- (5.5441,-0.0982) -- (5.5429,-0.0965) -- (5.5416,-0.0949) -- (5.5403,-0.0933) -- (5.5390,-0.0916) -- (5.5377,-0.0900) -- (5.5365,-0.0883) -- (5.5352,-0.0867) -- (5.5339,-0.0851) -- (5.5326,-0.0834) -- (5.5313,-0.0818) -- (5.5301,-0.0801) -- (5.5288,-0.0785) -- (5.5275,-0.0769) -- (5.5262,-0.0752) -- (5.5249,-0.0736);
\draw[fmRed, line width=0.8pt, dashed] (5.5249,-0.0736) -- (5.5207,-0.0681) -- (5.5164,-0.0627) -- (5.5121,-0.0572) -- (5.5079,-0.0518) -- (5.5036,-0.0463) -- (5.4993,-0.0409) -- (5.4950,-0.0354) -- (5.4907,-0.0300) -- (5.4864,-0.0245) -- (5.4821,-0.0191) -- (5.4778,-0.0136) -- (5.4735,-0.0081) -- (5.4692,-0.0027) -- (5.4649,0.0028) -- (5.4606,0.0082) -- (5.4563,0.0137) -- (5.4519,0.0192) -- (5.4476,0.0246) -- (5.4433,0.0301) -- (5.4389,0.0355) -- (5.4346,0.0410) -- (5.4303,0.0465) -- (5.4259,0.0519) -- (5.4216,0.0574) -- (5.4172,0.0628) -- (5.4129,0.0683) -- (5.4085,0.0738) -- (5.4041,0.0792) -- (5.3998,0.0847);
\draw[fmGreen!70!black, line width=1.2pt, -{Stealth[length=1.8mm]}] (5.5249,-0.0736) -- (5.5336,-0.0650) -- (5.5423,-0.0564) -- (5.5509,-0.0478) -- (5.5596,-0.0391) -- (5.5683,-0.0305) -- (5.5769,-0.0219) -- (5.5856,-0.0132) -- (5.5943,-0.0046) -- (5.6030,0.0041) -- (5.6117,0.0127) -- (5.6204,0.0214) -- (5.6291,0.0300) -- (5.6379,0.0387) -- (5.6466,0.0474) -- (5.6553,0.0561) -- (5.6640,0.0648) -- (5.6728,0.0735) -- (5.6815,0.0822) -- (5.6903,0.0909) -- (5.6991,0.0996) -- (5.7078,0.1083) -- (5.7166,0.1170) -- (5.7254,0.1258) -- (5.7342,0.1345) -- (5.7430,0.1432) -- (5.7518,0.1520) -- (5.7606,0.1608) -- (5.7694,0.1695) -- (5.7782,0.1783);
\coordinate (F) at (6.2982,-1.1339);
\coordinate (FA) at (0.6596,0.2199);
\coordinate (Fcross) at (5.5249,-0.0736);
\coordinate (nPrim) at (0.8785,-0.8785);
\coordinate (dPrim) at (0.0000,0.9941);
\coordinate (nConj) at (-0.8785,0.8785);
\coordinate (dConj) at (0.9941,0.0000);
\coordinate (nCrit) at (0.8785,0.8785);
\coordinate (dCrit) at (0.0000,-0.9941);
\coordinate (nCross) at (-0.8785,-0.8785);
\coordinate (A100) at (0.0000,0.0000);
\coordinate (A010) at (2.4000,0.0000);
\coordinate (A001) at (0.0000,2.4000);
\coordinate (A110) at (0.9941,0.0000);

  % (а) вся проекция
  \fill[fmRed] (FA) circle (0.035);
  \node[font=\scriptsize, anchor=north east, inner sep=1pt] at (A100) {$[100]$};
  \node[font=\scriptsize, anchor=west] at (A010) {$[010]$};
  \node[font=\scriptsize, anchor=south] at (A001) {$[001]$};
  \node[ring=fmRed] at (dPrim) {};
  \node[font=\scriptsize, fmRed, anchor=south east, inner sep=1pt] at (dPrim) {$[101]$};
  \node[ring=fmRed] at (nPrim) {};
  \node[font=\scriptsize, fmRed, anchor=north west, inner sep=2pt] at (nPrim) {$(11\br1)$};
  \node[ring=fmGreen] at (nConj) {};
  \node[font=\scriptsize, fmGreen, anchor=south east, inner sep=2pt] at (nConj) {$(1\br11)$};
  \node[ring=fmPurple] at (nCrit) {};
  \node[font=\scriptsize, fmPurple, anchor=south west, inner sep=2pt] at (nCrit) {$(111)$};
  \node[ring=fmTeal] at (nCross) {};
  \node[font=\scriptsize, fmTeal, anchor=north east, inner sep=2pt] at (nCross) {$(1\br1\br1)$};
  \node[font=\small] at (0,-2.75) {\textit{а}};
  % (б) окрестность стандартного треугольника
  \node[font=\scriptsize, anchor=north east, inner sep=1pt] at (B100) {$[100]$};
  \node[font=\scriptsize, anchor=north, inner sep=2pt] at (B110) {$[110]$};
  \node[font=\scriptsize, anchor=south west, inner sep=1pt] at (B111) {$[111]$};
  \node[font=\scriptsize, anchor=west, inner sep=2pt] at (B101) {$[101]$};
  \node[font=\scriptsize, fmGreen!70!black, anchor=north west, inner sep=1pt] at (B211) {$[211]$};
  \fill[fmRed] (F) circle (0.05) node[below right=-1pt, font=\small, fmRed]{$F$};
  \node[font=\small] at (5.6,-2.75) {\textit{б}};
\end{tikzpicture}
\caption{Системы скольжения ГЦК-кристалла на стереографической проекции
вдоль $[100]$: \textit{а}~— вся проекция; кружками отмечены полюса
плоскостей скольжения первичной $(11\br1)$ (направление $[101]$),
сопряжённой $(1\br11)$ (направление $[110]$), критической $(111)$
(направление $[10\br1]$) и поперечной $(1\br1\br1)$ (направление $[101]$)
систем; квадраты~— $\nfam{100}$, овалы~— $\nfam{110}$, треугольники~—
$\nfam{111}$. \textit{б}~— стандартный треугольник $[100]$--$[110]$--$[111]$
крупно: при одиночном скольжении ось $F\parallel[531]$ поворачивается к
$[101]$ (красная линия), немного «перескакивает» сторону $[100]$--$[111]$
(пунктир), после чего при двойном скольжении движется к $[211]$
(зелёная стрелка)}
\label{fig:slip-stereo}
\end{figure}

Системам скольжения дают названия по их роли в деформации
(рис.~\ref{fig:slip-stereo}):
\begin{description}
\item[Первичная] $(11\br1)[101]$~— с наибольшим фактором Шмида для любой
оси внутри треугольника; включается первой.
\item[Сопряжённая] $(1\br11)[110]$~— зеркальное отражение первичной в
стороне $[100]$--$[111]$. При растяжении ось $F$ поворачивается к
направлению скольжения $[101]$ и достигает этой стороны; там факторы
Шмида первичной и сопряжённой систем равны, и начинается \emph{двойное
скольжение}: ось движется вдоль стороны к $[211]=[101]+[110]$. Обычно
ось немного «перескакивает» границу в соседний треугольник
$[100]$--$[101]$--$[111]$ (пунктир), потому что сопряжённая система
упрочнена дислокациями первичной.
\item[Критическая] $(111)[10\br1]$~— отражение первичной в стороне
$[100]$--$[110]$. При \emph{сжатии} ось движется к нормали плоскости
скольжения $(11\br1)$, попадает в треугольник $[100]$--$[110]$--$[11\br1]$
и включает критическую систему.
\item[Поперечная] $(1\br1\br1)[101]$~— имеет то же направление
скольжения (тот же вектор Бюргерса), что и первичная; в эту плоскость
переходят винтовые дислокации при поперечном скольжении (гл.~2,
п.~2.2.2).
\end{description}

\subsubsection{Скольжение в ОЦК-металлах}

В ОЦК-металлах ($\alpha$-Fe, Mo, W, Nb, Ta, $\beta$-Ti) направление
скольжения~— плотноупакованное $\nfam{111}$, а плоскостью скольжения в
зависимости от металла и температуры служат $\plfam{110}$, $\plfam{112}$
или $\plfam{123}$~— все они содержат направление $\nfam{111}$.
Систем: $\plfam{110}\nfam{111}$~— 12, $\plfam{112}\nfam{111}$~— 12,
$\plfam{123}\nfam{111}$~— 24. Часто наблюдается и
\term{некристаллографическое} («карандашное») скольжение: направление
$\nfam{111}$, а плоскость~— любая, проходящая через него, с наибольшим
касательным напряжением ($\chi=90^\circ-\lambda$). Следы такого скольжения
волнистые.

Как правило, скольжение \emph{неполярно}: сдвиг вдоль $+\vb m$ и $-\vb m$
кристаллографически эквивалентен. Но если плоскость скольжения~— не
плоскость симметрии и в ней нет оси 2, переводящей $\vb m$ в $-\vb m$,
то направление скольжения~— \emph{полярный} вектор в этой плоскости, и
сопротивление сдвигу «туда» и «обратно» различно (несимметричный
потенциальный рельеф решётки). Типичный пример~— системы
$\plfam{112}\nfam{111}$ и $\plfam{123}\nfam{111}$ в ОЦК-металлах: в
одном направлении сдвиг по $\plfam{112}$ идёт «в двойниковом», в обратном~—
«в антидвойниковом» смысле. В системе $\plfam{110}\nfam{111}$ полярности
нет: ось 2 вдоль нормали $\nfam{110}$ переводит $[111]$ в $[\br1\br1\br1]$.

\subsubsection{Скольжение в ГПУ-металлах}

В ГПУ-металлах направление скольжения~— обычно плотноупакованное
$\nfam{11\br20}$ (то же семейство, что $\nfam{2\br1\br10}$), а плоскости
скольжения разнообразны (рис.~\ref{fig:hcp-slip}): \term{базисная}
$(0001)$, \term{призматические} $\plfam{10\br10}$ и \term{пирамидальные}
$\plfam{10\br11}$, $\plfam{10\br12}$, $\plfam{11\br21}$,
$\plfam{11\br22}$\ispr{В оригинале пирамидальные плоскости названы
плоскостями I--IV рода в порядке перечисления. В литературе чаще
различают пирамиды I рода $\plfam{h0\br hl}$ (сюда относятся
$\plfam{10\br11}$ и $\plfam{10\br12}$) и II рода $\plfam{hh\,\overline{2h}\,l}$
($\plfam{11\br21}$, $\plfam{11\br22}$).}.

\begin{figure}[ht]
\centering
\begin{tikzpicture}[cubeproj, scale=0.95]
  % создано tools/gen_plastic.py: плоскости скольжения ГПУ (cubeproj; c/a = 1.633)
\begin{scope}[shift={(0.0,0,0.0)}]
\draw[atomLine!45, dashed] (0.5000,0.8660,0.0000) -- (-0.5000,0.8660,0.0000);
\draw[atomLine!45, dashed] (-0.5000,0.8660,0.0000) -- (-1.0000,0.0000,0.0000);
\draw[atomLine!45, dashed] (-0.5000,0.8660,0.0000) -- (-0.5000,0.8660,1.6330);
\fill[fmRed!35, opacity=0.85] (-1.0000,0.0000,0.8165) -- (-0.5000,-0.8660,0.8165) -- (0.5000,-0.8660,0.8165) -- (1.0000,0.0000,0.8165) -- (0.5000,0.8660,0.8165) -- (-0.5000,0.8660,0.8165) -- cycle;
\draw[fmRed!80!black, line width=0.7pt] (-1.0000,0.0000,0.8165) -- (-0.5000,-0.8660,0.8165) -- (0.5000,-0.8660,0.8165) -- (1.0000,0.0000,0.8165) -- (0.5000,0.8660,0.8165) -- (-0.5000,0.8660,0.8165) -- cycle;
\draw[atomLine, line width=0.5pt] (1.0000,0.0000,0.0000) -- (0.5000,0.8660,0.0000);
\draw[atomLine, line width=0.5pt] (-1.0000,0.0000,0.0000) -- (-0.5000,-0.8660,0.0000);
\draw[atomLine, line width=0.5pt] (-0.5000,-0.8660,0.0000) -- (0.5000,-0.8660,0.0000);
\draw[atomLine, line width=0.5pt] (0.5000,-0.8660,0.0000) -- (1.0000,0.0000,0.0000);
\draw[atomLine, line width=0.5pt] (1.0000,0.0000,1.6330) -- (0.5000,0.8660,1.6330);
\draw[atomLine, line width=0.5pt] (0.5000,0.8660,1.6330) -- (-0.5000,0.8660,1.6330);
\draw[atomLine, line width=0.5pt] (-0.5000,0.8660,1.6330) -- (-1.0000,0.0000,1.6330);
\draw[atomLine, line width=0.5pt] (-1.0000,0.0000,1.6330) -- (-0.5000,-0.8660,1.6330);
\draw[atomLine, line width=0.5pt] (-0.5000,-0.8660,1.6330) -- (0.5000,-0.8660,1.6330);
\draw[atomLine, line width=0.5pt] (0.5000,-0.8660,1.6330) -- (1.0000,0.0000,1.6330);
\draw[atomLine, line width=0.5pt] (1.0000,0.0000,0.0000) -- (1.0000,0.0000,1.6330);
\draw[atomLine, line width=0.5pt] (0.5000,0.8660,0.0000) -- (0.5000,0.8660,1.6330);
\draw[atomLine, line width=0.5pt] (-1.0000,0.0000,0.0000) -- (-1.0000,0.0000,1.6330);
\draw[atomLine, line width=0.5pt] (-0.5000,-0.8660,0.0000) -- (-0.5000,-0.8660,1.6330);
\draw[atomLine, line width=0.5pt] (0.5000,-0.8660,0.0000) -- (0.5000,-0.8660,1.6330);
\node[font=\small, anchor=north] at (0,-1.05,0) {$1$: $(0001)$};
\end{scope}
\begin{scope}[shift={(3.5,0,0.0)}]
\draw[atomLine!45, dashed] (0.5000,0.8660,0.0000) -- (-0.5000,0.8660,0.0000);
\draw[atomLine!45, dashed] (-0.5000,0.8660,0.0000) -- (-1.0000,0.0000,0.0000);
\draw[atomLine!45, dashed] (-0.5000,0.8660,0.0000) -- (-0.5000,0.8660,1.6330);
\fill[fmRed!35, opacity=0.85] (-1.0000,0.0000,1.6330) -- (1.0000,0.0000,1.6330) -- (1.0000,0.0000,0.0000) -- (-1.0000,0.0000,0.0000) -- cycle;
\draw[fmRed!80!black, line width=0.7pt] (-1.0000,0.0000,1.6330) -- (1.0000,0.0000,1.6330) -- (1.0000,0.0000,0.0000) -- (-1.0000,0.0000,0.0000) -- cycle;
\draw[atomLine, line width=0.5pt] (1.0000,0.0000,0.0000) -- (0.5000,0.8660,0.0000);
\draw[atomLine, line width=0.5pt] (-1.0000,0.0000,0.0000) -- (-0.5000,-0.8660,0.0000);
\draw[atomLine, line width=0.5pt] (-0.5000,-0.8660,0.0000) -- (0.5000,-0.8660,0.0000);
\draw[atomLine, line width=0.5pt] (0.5000,-0.8660,0.0000) -- (1.0000,0.0000,0.0000);
\draw[atomLine, line width=0.5pt] (1.0000,0.0000,1.6330) -- (0.5000,0.8660,1.6330);
\draw[atomLine, line width=0.5pt] (0.5000,0.8660,1.6330) -- (-0.5000,0.8660,1.6330);
\draw[atomLine, line width=0.5pt] (-0.5000,0.8660,1.6330) -- (-1.0000,0.0000,1.6330);
\draw[atomLine, line width=0.5pt] (-1.0000,0.0000,1.6330) -- (-0.5000,-0.8660,1.6330);
\draw[atomLine, line width=0.5pt] (-0.5000,-0.8660,1.6330) -- (0.5000,-0.8660,1.6330);
\draw[atomLine, line width=0.5pt] (0.5000,-0.8660,1.6330) -- (1.0000,0.0000,1.6330);
\draw[atomLine, line width=0.5pt] (1.0000,0.0000,0.0000) -- (1.0000,0.0000,1.6330);
\draw[atomLine, line width=0.5pt] (0.5000,0.8660,0.0000) -- (0.5000,0.8660,1.6330);
\draw[atomLine, line width=0.5pt] (-1.0000,0.0000,0.0000) -- (-1.0000,0.0000,1.6330);
\draw[atomLine, line width=0.5pt] (-0.5000,-0.8660,0.0000) -- (-0.5000,-0.8660,1.6330);
\draw[atomLine, line width=0.5pt] (0.5000,-0.8660,0.0000) -- (0.5000,-0.8660,1.6330);
\node[font=\small, anchor=north] at (0,-1.05,0) {$2$: $(01\br10)$};
\end{scope}
\begin{scope}[shift={(7.0,0,0.0)}]
\draw[atomLine!45, dashed] (0.5000,0.8660,0.0000) -- (-0.5000,0.8660,0.0000);
\draw[atomLine!45, dashed] (-0.5000,0.8660,0.0000) -- (-1.0000,0.0000,0.0000);
\draw[atomLine!45, dashed] (-0.5000,0.8660,0.0000) -- (-0.5000,0.8660,1.6330);
\fill[fmRed!35, opacity=0.85] (-1.0000,0.0000,1.6330) -- (-0.5000,-0.8660,0.0000) -- (0.5000,-0.8660,0.0000) -- (1.0000,0.0000,1.6330) -- cycle;
\draw[fmRed!80!black, line width=0.7pt] (-1.0000,0.0000,1.6330) -- (-0.5000,-0.8660,0.0000) -- (0.5000,-0.8660,0.0000) -- (1.0000,0.0000,1.6330) -- cycle;
\draw[atomLine, line width=0.5pt] (1.0000,0.0000,0.0000) -- (0.5000,0.8660,0.0000);
\draw[atomLine, line width=0.5pt] (-1.0000,0.0000,0.0000) -- (-0.5000,-0.8660,0.0000);
\draw[atomLine, line width=0.5pt] (-0.5000,-0.8660,0.0000) -- (0.5000,-0.8660,0.0000);
\draw[atomLine, line width=0.5pt] (0.5000,-0.8660,0.0000) -- (1.0000,0.0000,0.0000);
\draw[atomLine, line width=0.5pt] (1.0000,0.0000,1.6330) -- (0.5000,0.8660,1.6330);
\draw[atomLine, line width=0.5pt] (0.5000,0.8660,1.6330) -- (-0.5000,0.8660,1.6330);
\draw[atomLine, line width=0.5pt] (-0.5000,0.8660,1.6330) -- (-1.0000,0.0000,1.6330);
\draw[atomLine, line width=0.5pt] (-1.0000,0.0000,1.6330) -- (-0.5000,-0.8660,1.6330);
\draw[atomLine, line width=0.5pt] (-0.5000,-0.8660,1.6330) -- (0.5000,-0.8660,1.6330);
\draw[atomLine, line width=0.5pt] (0.5000,-0.8660,1.6330) -- (1.0000,0.0000,1.6330);
\draw[atomLine, line width=0.5pt] (1.0000,0.0000,0.0000) -- (1.0000,0.0000,1.6330);
\draw[atomLine, line width=0.5pt] (0.5000,0.8660,0.0000) -- (0.5000,0.8660,1.6330);
\draw[atomLine, line width=0.5pt] (-1.0000,0.0000,0.0000) -- (-1.0000,0.0000,1.6330);
\draw[atomLine, line width=0.5pt] (-0.5000,-0.8660,0.0000) -- (-0.5000,-0.8660,1.6330);
\draw[atomLine, line width=0.5pt] (0.5000,-0.8660,0.0000) -- (0.5000,-0.8660,1.6330);
\node[font=\small, anchor=north] at (0,-1.05,0) {$3$: $(0\br111)$};
\end{scope}
\begin{scope}[shift={(0.0,0,-3.0)}]
\draw[atomLine!45, dashed] (0.5000,0.8660,0.0000) -- (-0.5000,0.8660,0.0000);
\draw[atomLine!45, dashed] (-0.5000,0.8660,0.0000) -- (-1.0000,0.0000,0.0000);
\draw[atomLine!45, dashed] (-0.5000,0.8660,0.0000) -- (-0.5000,0.8660,1.6330);
\fill[fmRed!35, opacity=0.85] (0.5000,0.8660,1.6330) -- (-0.5000,0.8660,1.6330) -- (-1.0000,0.0000,0.8165) -- (-0.5000,-0.8660,0.0000) -- (0.5000,-0.8660,0.0000) -- (1.0000,0.0000,0.8165) -- cycle;
\draw[fmRed!80!black, line width=0.7pt] (0.5000,0.8660,1.6330) -- (-0.5000,0.8660,1.6330) -- (-1.0000,0.0000,0.8165) -- (-0.5000,-0.8660,0.0000) -- (0.5000,-0.8660,0.0000) -- (1.0000,0.0000,0.8165) -- cycle;
\draw[atomLine, line width=0.5pt] (1.0000,0.0000,0.0000) -- (0.5000,0.8660,0.0000);
\draw[atomLine, line width=0.5pt] (-1.0000,0.0000,0.0000) -- (-0.5000,-0.8660,0.0000);
\draw[atomLine, line width=0.5pt] (-0.5000,-0.8660,0.0000) -- (0.5000,-0.8660,0.0000);
\draw[atomLine, line width=0.5pt] (0.5000,-0.8660,0.0000) -- (1.0000,0.0000,0.0000);
\draw[atomLine, line width=0.5pt] (1.0000,0.0000,1.6330) -- (0.5000,0.8660,1.6330);
\draw[atomLine, line width=0.5pt] (0.5000,0.8660,1.6330) -- (-0.5000,0.8660,1.6330);
\draw[atomLine, line width=0.5pt] (-0.5000,0.8660,1.6330) -- (-1.0000,0.0000,1.6330);
\draw[atomLine, line width=0.5pt] (-1.0000,0.0000,1.6330) -- (-0.5000,-0.8660,1.6330);
\draw[atomLine, line width=0.5pt] (-0.5000,-0.8660,1.6330) -- (0.5000,-0.8660,1.6330);
\draw[atomLine, line width=0.5pt] (0.5000,-0.8660,1.6330) -- (1.0000,0.0000,1.6330);
\draw[atomLine, line width=0.5pt] (1.0000,0.0000,0.0000) -- (1.0000,0.0000,1.6330);
\draw[atomLine, line width=0.5pt] (0.5000,0.8660,0.0000) -- (0.5000,0.8660,1.6330);
\draw[atomLine, line width=0.5pt] (-1.0000,0.0000,0.0000) -- (-1.0000,0.0000,1.6330);
\draw[atomLine, line width=0.5pt] (-0.5000,-0.8660,0.0000) -- (-0.5000,-0.8660,1.6330);
\draw[atomLine, line width=0.5pt] (0.5000,-0.8660,0.0000) -- (0.5000,-0.8660,1.6330);
\node[font=\small, anchor=north] at (0,-1.05,0) {$4$: $(0\br112)$};
\end{scope}
\begin{scope}[shift={(3.5,0,-3.0)}]
\draw[atomLine!45, dashed] (0.5000,0.8660,0.0000) -- (-0.5000,0.8660,0.0000);
\draw[atomLine!45, dashed] (-0.5000,0.8660,0.0000) -- (-1.0000,0.0000,0.0000);
\draw[atomLine!45, dashed] (-0.5000,0.8660,0.0000) -- (-0.5000,0.8660,1.6330);
\fill[fmRed!35, opacity=0.85] (-0.7500,0.4330,1.6330) -- (0.7500,-0.4330,1.6330) -- (1.0000,0.0000,0.0000) -- (-0.5000,0.8660,0.0000) -- cycle;
\draw[fmRed!80!black, line width=0.7pt] (-0.7500,0.4330,1.6330) -- (0.7500,-0.4330,1.6330) -- (1.0000,0.0000,0.0000) -- (-0.5000,0.8660,0.0000) -- cycle;
\draw[atomLine, line width=0.5pt] (1.0000,0.0000,0.0000) -- (0.5000,0.8660,0.0000);
\draw[atomLine, line width=0.5pt] (-1.0000,0.0000,0.0000) -- (-0.5000,-0.8660,0.0000);
\draw[atomLine, line width=0.5pt] (-0.5000,-0.8660,0.0000) -- (0.5000,-0.8660,0.0000);
\draw[atomLine, line width=0.5pt] (0.5000,-0.8660,0.0000) -- (1.0000,0.0000,0.0000);
\draw[atomLine, line width=0.5pt] (1.0000,0.0000,1.6330) -- (0.5000,0.8660,1.6330);
\draw[atomLine, line width=0.5pt] (0.5000,0.8660,1.6330) -- (-0.5000,0.8660,1.6330);
\draw[atomLine, line width=0.5pt] (-0.5000,0.8660,1.6330) -- (-1.0000,0.0000,1.6330);
\draw[atomLine, line width=0.5pt] (-1.0000,0.0000,1.6330) -- (-0.5000,-0.8660,1.6330);
\draw[atomLine, line width=0.5pt] (-0.5000,-0.8660,1.6330) -- (0.5000,-0.8660,1.6330);
\draw[atomLine, line width=0.5pt] (0.5000,-0.8660,1.6330) -- (1.0000,0.0000,1.6330);
\draw[atomLine, line width=0.5pt] (1.0000,0.0000,0.0000) -- (1.0000,0.0000,1.6330);
\draw[atomLine, line width=0.5pt] (0.5000,0.8660,0.0000) -- (0.5000,0.8660,1.6330);
\draw[atomLine, line width=0.5pt] (-1.0000,0.0000,0.0000) -- (-1.0000,0.0000,1.6330);
\draw[atomLine, line width=0.5pt] (-0.5000,-0.8660,0.0000) -- (-0.5000,-0.8660,1.6330);
\draw[atomLine, line width=0.5pt] (0.5000,-0.8660,0.0000) -- (0.5000,-0.8660,1.6330);
\node[font=\small, anchor=north] at (0,-1.05,0) {$5$: $(11\br21)$};
\end{scope}
\begin{scope}[shift={(7.0,0,-3.0)}]
\draw[atomLine!45, dashed] (0.5000,0.8660,0.0000) -- (-0.5000,0.8660,0.0000);
\draw[atomLine!45, dashed] (-0.5000,0.8660,0.0000) -- (-1.0000,0.0000,0.0000);
\draw[atomLine!45, dashed] (-0.5000,0.8660,0.0000) -- (-0.5000,0.8660,1.6330);
\fill[fmRed!35, opacity=0.85] (0.5000,-0.8660,1.6330) -- (1.0000,0.0000,0.0000) -- (-0.5000,0.8660,0.0000) -- (-1.0000,0.0000,1.6330) -- cycle;
\draw[fmRed!80!black, line width=0.7pt] (0.5000,-0.8660,1.6330) -- (1.0000,0.0000,0.0000) -- (-0.5000,0.8660,0.0000) -- (-1.0000,0.0000,1.6330) -- cycle;
\draw[atomLine, line width=0.5pt] (1.0000,0.0000,0.0000) -- (0.5000,0.8660,0.0000);
\draw[atomLine, line width=0.5pt] (-1.0000,0.0000,0.0000) -- (-0.5000,-0.8660,0.0000);
\draw[atomLine, line width=0.5pt] (-0.5000,-0.8660,0.0000) -- (0.5000,-0.8660,0.0000);
\draw[atomLine, line width=0.5pt] (0.5000,-0.8660,0.0000) -- (1.0000,0.0000,0.0000);
\draw[atomLine, line width=0.5pt] (1.0000,0.0000,1.6330) -- (0.5000,0.8660,1.6330);
\draw[atomLine, line width=0.5pt] (0.5000,0.8660,1.6330) -- (-0.5000,0.8660,1.6330);
\draw[atomLine, line width=0.5pt] (-0.5000,0.8660,1.6330) -- (-1.0000,0.0000,1.6330);
\draw[atomLine, line width=0.5pt] (-1.0000,0.0000,1.6330) -- (-0.5000,-0.8660,1.6330);
\draw[atomLine, line width=0.5pt] (-0.5000,-0.8660,1.6330) -- (0.5000,-0.8660,1.6330);
\draw[atomLine, line width=0.5pt] (0.5000,-0.8660,1.6330) -- (1.0000,0.0000,1.6330);
\draw[atomLine, line width=0.5pt] (1.0000,0.0000,0.0000) -- (1.0000,0.0000,1.6330);
\draw[atomLine, line width=0.5pt] (0.5000,0.8660,0.0000) -- (0.5000,0.8660,1.6330);
\draw[atomLine, line width=0.5pt] (-1.0000,0.0000,0.0000) -- (-1.0000,0.0000,1.6330);
\draw[atomLine, line width=0.5pt] (-0.5000,-0.8660,0.0000) -- (-0.5000,-0.8660,1.6330);
\draw[atomLine, line width=0.5pt] (0.5000,-0.8660,0.0000) -- (0.5000,-0.8660,1.6330);
\node[font=\small, anchor=north] at (0,-1.05,0) {$6$: $(11\br22)$};
\end{scope}

\end{tikzpicture}
\caption{Плоскости скольжения в ГПУ-кристалле (гексагональная призма,
$c/a=1{,}633$): 1~— базисная $(0001)$; 2~— призматическая $(01\br10)$
из семейства $\plfam{10\br10}$; 3,\,4~— пирамидальные I рода $(0\br111)$ и
$(0\br112)$ из семейств $\plfam{10\br11}$ и $\plfam{10\br12}$; 5,\,6~—
пирамидальные II рода $(11\br21)$ и $(11\br22)$}
\label{fig:hcp-slip}
\end{figure}

Какая система работает, грубо коррелирует с отношением $c/a$
(табл.~\ref{tab:slip-systems}):
\begin{itemize}
\item при $c/a$ больше идеального $1{,}633$ базисная плоскость~— самая
плотноупакованная и далеко отстоящая от соседних; действует
\emph{базисное} скольжение (Zn, $c/a=1{,}86$; Cd, $1{,}89$); у Mg
($c/a=1{,}62\approx1{,}633$)~— тоже базисное;
\item при $c/a<1{,}633$ ($\alpha$-Ti, $\alpha$-Zr, $c/a\approx1{,}59$)
преобладает \emph{призматическое} скольжение;
\item но бериллий ($c/a=1{,}57$) скользит по базису~— правило $c/a$
лишь эмпирическое: решают детали межатомной связи.
\end{itemize}

\begin{table}[ht]
\centering
\caption{Основные системы скольжения металлов}
\label{tab:slip-systems}
\small
\begin{tabular}{@{}llccl@{}}
\toprule
Решётка & Система & Число систем & Независимых & Примеры\\
\midrule
ГЦК & $\plfam{111}\nfam{1\br10}$ & 12 & 5 & Al, Cu, Ni, Au, $\gamma$-Fe\\
ОЦК & $\plfam{110}\nfam{\br111}$ & 12 & 5 & $\alpha$-Fe, Mo, W, Nb\\
    & $\plfam{112}\nfam{\br1\br11}$ & 12 & 5 & $\alpha$-Fe, W\\
    & $\plfam{123}\nfam{\br1\br11}$ & 24 & 5 & $\alpha$-Fe\\
ГПУ & $(0001)\nfam{11\br20}$ & 3 & 2 & Zn, Cd, Mg, Be\\
    & $\plfam{10\br10}\nfam{11\br20}$ & 3 & 2 & Ti, Zr\\
    & $\plfam{10\br11}\nfam{11\br20}$ & 6 & 4 & Ti, Mg\\
    & $\plfam{11\br22}\nfam{11\br2\br3}$ & 6 & 5 & Zn, Cd, Ti ($\vb c+\vb a$)\\
\bottomrule
\end{tabular}
\end{table}

Сдвиг вдоль $\nfam{11\br20}$ не меняет размеров кристалла вдоль оси $c$,
сколько бы таких систем ни работало. Поэтому для произвольного
формоизменения ГПУ-металлу нужны пирамидальное скольжение с вектором
$\vb c+\vb a$ (направление $\nfam{11\br2\br3}$) или двойникование~— иначе
поликристалл хрупок.

\subsubsection{Механическое двойникование}

Двойникование играет в деформации сравнительно небольшую роль и
проявляется, когда скольжение затруднено. Ему способствуют низкая
температура и ударное нагружение: с понижением температуры и ростом
скорости деформации критическое напряжение скольжения растёт быстрее, чем
двойникования. Легче всего двойникуются гексагональные металлы, труднее~—
ОЦК, труднее всего~— ГЦК.

В отличие от скольжения, при двойниковании атомы смещаются на \emph{доли}
вектора трансляции, а смещение каждого атома пропорционально его
расстоянию от плоскости двойникования (рис.~\ref{fig:slip-twin},\,\textit{в}).
Такой \term{простой сдвиг}~— однородная деформация: мысленно выделенный в
кристалле шар превращается в эллипсоид. Опишем её
(рис.~\ref{fig:twin-shear}). Ось $y$ направим вдоль направления сдвига
$\eta_1$, ось $z$~— по нормали к плоскости сдвига $K_1$; верхняя
половина шара единичного радиуса сдвигается: $(y,z)\to(y+sz,\,z)$.

\begin{figure}[ht]
\centering
\begin{tikzpicture}[scale=1.9]
  \def\s{0.707}
  \fill[fmBlueBg] (-1,0) arc[start angle=180, end angle=360, radius=1] -- cycle;
  \fill[fmAmber!15] plot[domain=0:180, samples=73, smooth] ({cos(\x)+\s*sin(\x)},{sin(\x)}) -- cycle;
  \draw[atomLine, dashed] (1,0) arc[start angle=0, end angle=180, radius=1];
  \draw[fmAmber!80!black, line width=0.9pt] plot[domain=0:180, samples=73, smooth] ({cos(\x)+\s*sin(\x)},{sin(\x)});
  \draw[atomLine, line width=0.6pt] (-1,0) arc[start angle=180, end angle=360, radius=1];
  \draw[fmGray, line width=1.1pt] (-1.6,0) -- (1.9,0) node[right, font=\small]{$K_1$};
  \draw[bvec, fmBlue] (1.05,-0.18) -- (1.65,-0.18) node[below, font=\small]{$\eta_1$};
  % неискажённое направление: до (A) и после (A') сдвига
  \pgfmathsetmacro{\nn}{sqrt(1+\s*\s/4)}
  \coordinate (A) at ({-\s/2/\nn},{1/\nn});
  \coordinate (A2) at ({\s/2/\nn},{1/\nn});
  \draw[fmRed, line width=0.8pt] ($-0.5*(A)$) -- ($1.45*(A)$) node[above, font=\small]{$K_2$};
  \draw[fmRed!70!black, line width=0.8pt, dashed] (0,0) -- ($1.45*(A2)$) node[above, font=\small]{$K_2'$};
  \fill (A) circle (0.03) node[left, font=\small]{$A$};
  \fill (A2) circle (0.03) node[right, font=\small]{$A'$};
  \draw[arr, fmRed] (A) -- (A2);
  \node[font=\footnotesize, fmRed] at ($(A)!0.5!(A2)+(0,0.16)$) {$\eta_2$};
  \fill (0,0) circle (0.025) node[below, font=\small]{$O$};
  \draw[arr, fmAmber!80!black] (0,1) -- (\s,1);
  \node[font=\footnotesize, above] at ({\s/2+0.05},1.02) {$s$};
  \draw[fmGray, densely dotted] (0,0) -- (0,1) (\s,1) -- (0,0);
\end{tikzpicture}
\caption{Однородный сдвиг при двойниковании: верхняя половина единичной
сферы (пунктир) переходит в половину эллипсоида. Неискажёнными остаются
плоскость сдвига $K_1$ и плоскость $K_2$: направление $OA$ ($\eta_2$)
переходит в $OA'$ той же длины; $s$~— кристаллографический сдвиг
(на рисунке $s=0{,}707$, как у ГЦК и ОЦК)}
\label{fig:twin-shear}
\end{figure}

У эллипсоида два круговых \emph{неискажённых} сечения: плоскость сдвига
$K_1$ (она не меняется вовсе) и плоскость $K_2$, перпендикулярная
плоскости рисунка: направление $OA$ ($\eta_2$) переходит в $OA'$ без
изменения длины. Действительно, вектор $(y,z)$ сохраняет длину, если
$(y+sz)^2+z^2=y^2+z^2$, т.\,е. $y=-sz/2$. Отсюда угол $2\phi$ между $K_1$
и $K_2$ связан со сдвигом:
\begin{equation}\label{eq:twin-shear}
s=2\operatorname{ctg}2\phi .
\end{equation}

\begin{opredelenie}
\term{Элементы двойникования}: плоскости $K_1$ (плоскость двойникования)
и $K_2$ (вторая неискажённая плоскость), направления $\eta_1$ (направление
двойникования, лежит в $K_1$) и $\eta_2$ (лежит в $K_2$, перпендикулярно
линии пересечения $K_1$ и $K_2$) и величина $s$~— \term{кристаллографический
сдвиг}.
\end{opredelenie}

\begin{table}[ht]
\centering
\caption{Элементы двойникования металлов}
\label{tab:twinning}
\small
\begin{tabular}{@{}lccccl@{}}
\toprule
Решётка & $K_1$ & $\eta_1$ & $K_2$ & $\eta_2$ & $s$\\
\midrule
ГЦК & $(111)$ & $[11\br2]$ & $(11\br1)$ & $[112]$ & $1/\sqrt2=0{,}707$\\
ОЦК & $(112)$ & $[11\br1]$ & $(11\br2)$ & $[111]$ & $1/\sqrt2=0{,}707$\\
ГПУ & $(10\br12)$ & $[\br1011]$ & $(\br1012)$ & $[10\br11]$ & $\dfrac{|3-(c/a)^2|}{\sqrt3\,c/a}$\\
\bottomrule
\end{tabular}
\end{table}

В ГЦК- и ОЦК-кристаллах обе плоскости $K_1$, $K_2$ принадлежат одному
семейству ($\plfam{111}$ и $\plfam{112}$), направления $\eta_1$, $\eta_2$~—
семействам $\nfam{112}$ и $\nfam{111}$ соответственно, а $s=0{,}707$
(табл.~\ref{tab:twinning}; угол между плоскостями $\plfam{111}$
$2\phi=70{,}5^\circ$ и по~\eqref{eq:twin-shear} $s=2\operatorname{ctg}70{,}5^\circ=0{,}707$).
В ГПУ-металлах ($\alpha$-Zr, Ti, Mg, Zn) главная система~—
$\plfam{10\br12}\nfam{\br1011}$, а сдвиг зависит от $c/a$: у Zr
($c/a=1{,}593$) $s=0{,}167$, у Mg $0{,}129$, у Zn $0{,}138$. При
$c/a<\sqrt3$ такой двойник удлиняет кристалл вдоль оси $c$, при $c/a>\sqrt3$~—
укорачивает.

\begin{vnimanie}
Скольжение не меняет ориентацию решётки внутри сдвинутой части (лишь
поворачивает кристалл как целое), а двойникование \emph{меняет}: двойник~—
это область с другой, симметрично связанной ориентацией. Поэтому
двойники видны на шлифе после травления как полосы иного оттенка, а
следы скольжения~— только как ступеньки на поверхности.
\end{vnimanie}

\subsection{Кристаллографическая текстура}\label{ssec:texture}

При пластической деформации поликристалла каждое зерно (если в нём
работают пять независимых систем деформации) меняет форму вместе с
соседями без образования пустот. При этом решётка зерна \emph{поворачивается}
(антисимметричная часть $\omega_{ij}$ тензора дисторсии), и ориентации
разных зёрен постепенно сближаются.

\begin{opredelenie}
\term{Текстура}~— преимущественная ориентация кристаллитов в
поликристалле.
\end{opredelenie}

Текстуры возникают при деформации (волочении, прокатке), рекристаллизации,
литье, осаждении плёнок и сильно влияют на анизотропию свойств изделия:
упругих модулей, магнитной проницаемости трансформаторной стали,
радиационного роста циркониевых оболочек твэлов.

\subsubsection{Аксиальная и ограниченная текстуры}
В простейшем случае определённые кристаллографические направления
$\nfam{mnp}$ всех кристаллитов параллельны одному внешнему направлению
$F$~— \term{оси текстуры}, а вокруг этой оси кристаллиты повёрнуты
произвольно. Такую текстуру называют \term{аксиальной} (осевой,
волокнистой, неограниченной). Она типична для проволоки: в ОЦК-металлах
ось текстуры волочения~— $\nfam{110}$, в ГЦК~— $\nfam{111}$ и
$\nfam{100}$ (доля $\nfam{100}$ растёт с уменьшением энергии дефекта
упаковки).

В прокатанных листах, лентах и фольге возникает \term{ограниченная}
текстура (текстура прокатки): определённая плоскость $(hkl)$
кристаллитов параллельна плоскости прокатки (поверхности листа), а
направление $[uvw]$~— направлению прокатки. Её записывают как
$\plfam{hkl}\nfam{uvw}$. Например, в ГЦК-металлах распространены
компоненты «медная» $\plfam{112}\nfam{111}$, «латунная»
$\plfam{110}\nfam{112}$, Госса $\plfam{110}\nfam{001}$ и после
рекристаллизации~— кубическая $\plfam{001}\nfam{100}$.

Текстуру описывают тремя способами: прямыми полюсными фигурами (ППФ),
обратными полюсными фигурами (ОПФ) и функцией распределения ориентаций
(ФРО). Все три сравнивают между собой две системы координат; начнём с
них, а затем разберём каждый способ.

\subsubsection{Две системы координат}

\begin{itemize}
\item \term{Система координат образца} привязана к изделию. Для листа это
направление прокатки (НП), поперечное направление (ПН) и нормаль к
листу (НН); для проволоки~— её ось. Она одна на весь лист.
\item \term{Система координат кристалла}~— оси $[100]$, $[010]$, $[001]$
решётки. В каждом зерне она своя.
\end{itemize}
\term{Ориентация} зерна~— поворот, совмещающий одну систему с другой
(матрица поворота или углы Эйлера, п.~\ref{ssec:orth}). Удобно представить
каждое зерно кубиком, рёбра которого идут вдоль $\nfam{100}$
(рис.~\ref{fig:texture-frames}). В металле без текстуры кубики повёрнуты
как попало. При \emph{кубической} текстуре $\plfam{001}\nfam{100}$ почти все
они «стоят» одинаково: грани параллельны поверхности и краям листа,
синие грани $\perp[001]$ смотрят вверх (рис.~\ref{fig:texture-frames},\,\textit{б}).

\begin{figure}[p]
\centering
\begin{tikzpicture}
  % создано tools/gen_plastic.py: системы координат образца и кристалла (текстура)
\fill[atomA!25] (4.600,0.900) -- (10.300,0.900) -- (10.300,1.565) -- (4.600,1.565) -- cycle;
\fill[atomA!40] (10.300,0.900) -- (11.896,2.040) -- (11.896,2.705) -- (10.300,1.565) -- cycle;
\fill[atomA!12] (4.600,1.565) -- (10.300,1.565) -- (11.896,2.705) -- (6.196,2.705) -- cycle;
\draw[atomLine!70, line width=0.5pt] (4.600,1.565) -- (5.650,1.565) -- (6.211,1.869) -- (5.235,2.019) -- cycle;
\draw[atomLine!70, line width=0.5pt] (5.650,1.565) -- (7.052,1.565) -- (7.097,1.885) -- (6.686,1.933) -- (6.211,1.869) -- cycle;
\draw[atomLine!70, line width=0.5pt] (7.052,1.565) -- (7.672,1.565) -- (8.495,1.918) -- (7.969,1.981) -- (7.097,1.885) -- cycle;
\draw[atomLine!70, line width=0.5pt] (7.672,1.565) -- (9.209,1.565) -- (9.497,1.854) -- (8.917,1.964) -- (8.495,1.918) -- cycle;
\draw[atomLine!70, line width=0.5pt] (9.209,1.565) -- (10.300,1.565) -- (11.043,2.095) -- (9.497,1.854) -- cycle;
\draw[atomLine!70, line width=0.5pt] (5.580,2.265) -- (5.235,2.019) -- (6.211,1.869) -- (6.686,1.933) -- (7.225,2.262) -- (7.020,2.313) -- cycle;
\draw[atomLine!70, line width=0.5pt] (7.225,2.262) -- (6.686,1.933) -- (7.097,1.885) -- (7.969,1.981) -- (8.179,2.302) -- (7.973,2.331) -- cycle;
\draw[atomLine!70, line width=0.5pt] (8.179,2.302) -- (7.969,1.981) -- (8.495,1.918) -- (8.917,1.964) -- (9.448,2.286) -- (9.086,2.368) -- cycle;
\draw[atomLine!70, line width=0.5pt] (11.043,2.095) -- (11.081,2.123) -- (10.222,2.374) -- (9.448,2.286) -- (8.917,1.964) -- (9.497,1.854) -- cycle;
\draw[atomLine!70, line width=0.5pt] (7.417,2.705) -- (6.196,2.705) -- (5.580,2.265) -- (7.020,2.313) -- cycle;
\draw[atomLine!70, line width=0.5pt] (8.447,2.705) -- (7.417,2.705) -- (7.020,2.313) -- (7.225,2.262) -- (7.973,2.331) -- cycle;
\draw[atomLine!70, line width=0.5pt] (9.474,2.705) -- (8.447,2.705) -- (7.973,2.331) -- (8.179,2.302) -- (9.086,2.368) -- cycle;
\draw[atomLine!70, line width=0.5pt] (10.883,2.705) -- (9.474,2.705) -- (9.086,2.368) -- (9.448,2.286) -- (10.222,2.374) -- cycle;
\draw[atomLine!70, line width=0.5pt] (11.081,2.123) -- (11.896,2.705) -- (10.883,2.705) -- (10.222,2.374) -- cycle;
\draw[fmAmber!80!black, line width=1.3pt] (8.179,2.302) -- (7.969,1.981) -- (8.495,1.918) -- (8.917,1.964) -- (9.448,2.286) -- (9.086,2.368) -- cycle;
\draw[atomLine, line width=0.6pt] (4.600,0.900) -- (10.300,0.900) -- (11.896,2.040) -- (11.896,2.705) -- (6.196,2.705) -- (4.600,1.565) -- cycle (4.600,1.565) -- (10.300,1.565) -- (11.896,2.705) (10.300,1.565) -- (10.300,0.900);
\filldraw[fmRed!35, draw=fmRed!60!black, line width=0.3pt] (9.913,2.552) -- (10.187,2.729) -- (10.162,3.192) -- (9.887,3.015) -- cycle;
\filldraw[fmGreen!35, draw=fmGreen!60!black, line width=0.3pt] (9.478,2.572) -- (9.452,3.035) -- (9.887,3.015) -- (9.913,2.552) -- cycle;
\filldraw[fmBlue!35, draw=fmBlue!60!black, line width=0.3pt] (9.452,3.035) -- (9.887,3.015) -- (10.162,3.192) -- (9.727,3.212) -- cycle;
\filldraw[fmRed!35, draw=fmRed!60!black, line width=0.3pt] (9.069,3.149) -- (8.607,3.252) -- (8.379,2.927) -- (8.841,2.824) -- cycle;
\filldraw[fmGreen!35, draw=fmGreen!60!black, line width=0.3pt] (9.063,2.788) -- (8.835,2.464) -- (8.841,2.824) -- (9.069,3.149) -- cycle;
\filldraw[fmBlue!35, draw=fmBlue!60!black, line width=0.3pt] (8.835,2.464) -- (8.841,2.824) -- (8.379,2.927) -- (8.373,2.567) -- cycle;
\filldraw[fmRed!35, draw=fmRed!60!black, line width=0.3pt] (7.809,2.513) -- (8.015,2.693) -- (8.020,3.155) -- (7.814,2.975) -- cycle;
\filldraw[fmGreen!35, draw=fmGreen!60!black, line width=0.3pt] (7.337,2.536) -- (7.341,2.998) -- (7.814,2.975) -- (7.809,2.513) -- cycle;
\filldraw[fmBlue!35, draw=fmBlue!60!black, line width=0.3pt] (7.341,2.998) -- (7.814,2.975) -- (8.020,3.155) -- (7.547,3.178) -- cycle;
\filldraw[fmRed!35, draw=fmRed!60!black, line width=0.3pt] (6.905,2.501) -- (7.052,2.701) -- (7.019,3.154) -- (6.872,2.955) -- cycle;
\filldraw[fmGreen!35, draw=fmGreen!60!black, line width=0.3pt] (6.412,2.473) -- (6.379,2.926) -- (6.872,2.955) -- (6.905,2.501) -- cycle;
\filldraw[fmBlue!35, draw=fmBlue!60!black, line width=0.3pt] (6.379,2.926) -- (6.872,2.955) -- (7.019,3.154) -- (6.526,3.126) -- cycle;
\filldraw[fmRed!35, draw=fmRed!60!black, line width=0.3pt] (11.287,2.510) -- (11.432,2.603) -- (11.473,3.090) -- (11.327,2.997) -- cycle;
\filldraw[fmGreen!35, draw=fmGreen!60!black, line width=0.3pt] (10.794,2.520) -- (10.835,3.007) -- (11.327,2.997) -- (11.287,2.510) -- cycle;
\filldraw[fmBlue!35, draw=fmBlue!60!black, line width=0.3pt] (10.835,3.007) -- (11.327,2.997) -- (11.473,3.090) -- (10.980,3.100) -- cycle;
\filldraw[fmRed!35, draw=fmRed!60!black, line width=0.3pt] (8.293,2.244) -- (8.732,2.137) -- (8.947,2.598) -- (8.509,2.705) -- cycle;
\filldraw[fmGreen!35, draw=fmGreen!60!black, line width=0.3pt] (8.131,2.393) -- (8.346,2.854) -- (8.509,2.705) -- (8.293,2.244) -- cycle;
\filldraw[fmBlue!35, draw=fmBlue!60!black, line width=0.3pt] (8.346,2.854) -- (8.509,2.705) -- (8.947,2.598) -- (8.785,2.747) -- cycle;
\filldraw[fmRed!35, draw=fmRed!60!black, line width=0.3pt] (10.015,2.167) -- (10.197,2.286) -- (10.205,2.767) -- (10.023,2.648) -- cycle;
\filldraw[fmGreen!35, draw=fmGreen!60!black, line width=0.3pt] (9.533,2.161) -- (9.541,2.642) -- (10.023,2.648) -- (10.015,2.167) -- cycle;
\filldraw[fmBlue!35, draw=fmBlue!60!black, line width=0.3pt] (9.541,2.642) -- (10.023,2.648) -- (10.205,2.767) -- (9.723,2.762) -- cycle;
\filldraw[fmRed!35, draw=fmRed!60!black, line width=0.3pt] (6.505,2.143) -- (6.695,2.307) -- (6.699,2.775) -- (6.509,2.611) -- cycle;
\filldraw[fmGreen!35, draw=fmGreen!60!black, line width=0.3pt] (6.026,2.153) -- (6.030,2.621) -- (6.509,2.611) -- (6.505,2.143) -- cycle;
\filldraw[fmBlue!35, draw=fmBlue!60!black, line width=0.3pt] (6.030,2.621) -- (6.509,2.611) -- (6.699,2.775) -- (6.220,2.785) -- cycle;
\filldraw[fmRed!35, draw=fmRed!60!black, line width=0.3pt] (7.757,2.141) -- (7.939,2.308) -- (7.894,2.773) -- (7.712,2.606) -- cycle;
\filldraw[fmGreen!35, draw=fmGreen!60!black, line width=0.3pt] (7.277,2.101) -- (7.232,2.567) -- (7.712,2.606) -- (7.757,2.141) -- cycle;
\filldraw[fmBlue!35, draw=fmBlue!60!black, line width=0.3pt] (7.232,2.567) -- (7.712,2.606) -- (7.894,2.773) -- (7.414,2.733) -- cycle;
\filldraw[fmRed!35, draw=fmRed!60!black, line width=0.3pt] (8.017,2.107) -- (7.643,2.290) -- (7.313,1.900) -- (7.687,1.718) -- cycle;
\filldraw[fmGreen!35, draw=fmGreen!60!black, line width=0.3pt] (7.513,2.536) -- (7.183,2.146) -- (7.313,1.900) -- (7.643,2.290) -- cycle;
\filldraw[fmBlue!35, draw=fmBlue!60!black, line width=0.3pt] (7.887,2.353) -- (8.017,2.107) -- (7.643,2.290) -- (7.513,2.536) -- cycle;
\filldraw[fmRed!35, draw=fmRed!60!black, line width=0.3pt] (10.311,1.818) -- (10.480,1.872) -- (10.520,2.365) -- (10.351,2.311) -- cycle;
\filldraw[fmGreen!35, draw=fmGreen!60!black, line width=0.3pt] (9.826,1.819) -- (9.866,2.312) -- (10.351,2.311) -- (10.311,1.818) -- cycle;
\filldraw[fmBlue!35, draw=fmBlue!60!black, line width=0.3pt] (9.866,2.312) -- (10.351,2.311) -- (10.520,2.365) -- (10.035,2.366) -- cycle;
\filldraw[fmRed!35, draw=fmRed!60!black, line width=0.3pt] (5.404,1.748) -- (5.652,1.953) -- (5.606,2.405) -- (5.358,2.200) -- cycle;
\filldraw[fmGreen!35, draw=fmGreen!60!black, line width=0.3pt] (4.955,1.751) -- (4.909,2.203) -- (5.358,2.200) -- (5.404,1.748) -- cycle;
\filldraw[fmBlue!35, draw=fmBlue!60!black, line width=0.3pt] (4.909,2.203) -- (5.358,2.200) -- (5.606,2.405) -- (5.157,2.408) -- cycle;
\filldraw[fmRed!35, draw=fmRed!60!black, line width=0.3pt] (8.710,1.734) -- (8.932,1.920) -- (8.863,2.378) -- (8.641,2.192) -- cycle;
\filldraw[fmGreen!35, draw=fmGreen!60!black, line width=0.3pt] (8.250,1.694) -- (8.182,2.152) -- (8.641,2.192) -- (8.710,1.734) -- cycle;
\filldraw[fmBlue!35, draw=fmBlue!60!black, line width=0.3pt] (8.182,2.152) -- (8.641,2.192) -- (8.863,2.378) -- (8.403,2.337) -- cycle;
\filldraw[fmRed!35, draw=fmRed!60!black, line width=0.3pt] (6.716,1.716) -- (6.893,1.887) -- (6.869,2.352) -- (6.692,2.181) -- cycle;
\filldraw[fmGreen!35, draw=fmGreen!60!black, line width=0.3pt] (6.233,1.697) -- (6.209,2.162) -- (6.692,2.181) -- (6.716,1.716) -- cycle;
\filldraw[fmBlue!35, draw=fmBlue!60!black, line width=0.3pt] (6.209,2.162) -- (6.692,2.181) -- (6.869,2.352) -- (6.386,2.333) -- cycle;
\fill[atomA!40] (0.150,1.250) -- (3.550,1.250) -- (3.550,1.410) -- (0.150,1.410) -- cycle;
\fill[atomA!55] (3.550,1.250) -- (4.474,1.910) -- (4.474,2.070) -- (3.550,1.410) -- cycle;
\fill[atomA!18] (0.150,1.410) -- (3.550,1.410) -- (4.474,2.070) -- (1.074,2.070) -- cycle;
\draw[atomLine, line width=0.6pt] (0.150,1.250) -- (3.550,1.250) -- (4.474,1.910) -- (4.474,2.070) -- (1.074,2.070) -- (0.150,1.410) -- cycle (0.150,1.410) -- (3.550,1.410) -- (4.474,2.070) (3.550,1.410) -- (3.550,1.250);
\draw[fmAmber!80!black, line width=0.8pt] (2.744,1.620) -- (3.194,1.620) -- (3.320,1.710) -- (2.870,1.710) -- cycle;
\draw[atomLine!50, densely dotted] (2.870,1.710) -- (6.196,2.705) (2.744,1.620) -- (4.600,1.565);
\draw[-{Stealth[length=1.8mm]}, atomLine, line width=0.8pt] (-0.452,1.070) -- (0.548,1.070) node[font=\scriptsize, anchor=north] {НП};
\draw[-{Stealth[length=1.8mm]}, atomLine, line width=0.8pt] (-0.452,1.070) -- (0.010,1.400) node[font=\scriptsize, anchor=east] {ПН};
\draw[-{Stealth[length=1.8mm]}, atomLine, line width=0.8pt] (-0.452,1.070) -- (-0.452,2.070) node[font=\scriptsize, anchor=south] {НН};
\coordinate (Alab) at (1.550,0.400);
\coordinate (Bsel) at (8.539,2.154);
\coordinate (Bfront) at (7.450,0.900);
\coordinate (Bright) at (11.098,1.802);
\coordinate (Btl) at (6.196,2.705);
\coordinate (Bbl) at (4.600,0.900);
\filldraw[fmRed!35, draw=fmRed!60!black, line width=0.3pt, opacity=0.75] (0.952,-3.561) -- (2.107,-3.842) -- (2.675,-2.629) -- (1.520,-2.348) -- cycle;
\filldraw[fmGreen!35, draw=fmGreen!60!black, line width=0.3pt, opacity=0.75] (0.525,-3.171) -- (1.093,-1.958) -- (1.520,-2.348) -- (0.952,-3.561) -- cycle;
\filldraw[fmBlue!35, draw=fmBlue!60!black, line width=0.3pt, opacity=0.75] (1.093,-1.958) -- (1.520,-2.348) -- (2.675,-2.629) -- (2.248,-2.239) -- cycle;
\draw[-{Stealth[length=2mm]}, fmRed!80!black, line width=1.1pt] (1.6,-2.9) -- (2.129,-3.384) node[font=\scriptsize, fmRed!80!black, anchor=west] {$[100]$};
\draw[-{Stealth[length=2mm]}, fmGreen!80!black, line width=1.1pt] (1.6,-2.9) -- (0.168,-2.551) node[font=\scriptsize, fmGreen!80!black, anchor=east] {$[0\br10]$};
\draw[-{Stealth[length=2mm]}, fmBlue!80!black, line width=1.1pt] (1.6,-2.9) -- (2.304,-1.396) node[font=\scriptsize, fmBlue!80!black, anchor=west] {$[001]$};
\draw[-{Stealth[length=2mm]}, atomLine, dashed, line width=0.8pt] (1.6,-2.9) -- (3.500,-2.900) node[font=\scriptsize, anchor=west] {НП};
\draw[-{Stealth[length=2mm]}, atomLine, dashed, line width=0.8pt] (1.6,-2.9) -- (2.398,-2.330) node[font=\scriptsize, anchor=south west] {ПН};
\draw[-{Stealth[length=2mm]}, atomLine, dashed, line width=0.8pt] (1.6,-2.9) -- (1.600,-1.000) node[font=\scriptsize, anchor=south] {НН};
\coordinate (Csel) at (1.6,-1.5499999999999998);
\draw[atomLine, line width=0.6pt] (6.3,-2.75) circle (1.45);
\draw[atomLine!40] (4.85,-2.75) -- (7.75,-2.75) (6.3,-4.2) -- (6.3,-1.3);
\filldraw[fmRed] (6.141,-1.348) circle (0.055);
\filldraw[fmGreen] (7.558,-2.612) circle (0.055);
\filldraw[fmBlue] (6.204,-2.781) circle (0.055);
\filldraw[fmRed] (6.364,-1.340) circle (0.055);
\filldraw[fmGreen] (7.664,-2.814) circle (0.055);
\filldraw[fmBlue] (6.255,-2.767) circle (0.055);
\filldraw[fmRed] (7.146,-3.411) circle (0.055);
\filldraw[fmGreen] (5.719,-3.231) circle (0.055);
\filldraw[fmBlue] (6.452,-2.244) circle (0.055);
\filldraw[fmRed] (6.222,-1.445) circle (0.055);
\filldraw[fmGreen] (7.619,-2.684) circle (0.055);
\filldraw[fmBlue] (6.236,-2.829) circle (0.055);
\filldraw[fmRed] (6.223,-4.172) circle (0.055);
\filldraw[fmGreen] (5.095,-2.681) circle (0.055);
\filldraw[fmBlue] (6.433,-2.744) circle (0.055);
\filldraw[fmRed] (6.270,-4.162) circle (0.055);
\filldraw[fmGreen] (7.684,-2.778) circle (0.055);
\filldraw[fmBlue] (6.267,-2.730) circle (0.055);
\filldraw[fmRed] (6.343,-1.405) circle (0.055);
\filldraw[fmGreen] (7.676,-2.799) circle (0.055);
\filldraw[fmBlue] (6.261,-2.803) circle (0.055);
\filldraw[fmRed, draw=fmAmber!80!black, line width=0.9pt] (7.299,-3.625) circle (0.075);
\filldraw[fmGreen, draw=fmAmber!80!black, line width=0.9pt] (5.657,-3.428) circle (0.075);
\filldraw[fmBlue, draw=fmAmber!80!black, line width=0.9pt] (6.462,-2.473) circle (0.075);
\filldraw[fmRed] (6.354,-1.303) circle (0.055);
\filldraw[fmGreen] (4.919,-2.698) circle (0.055);
\filldraw[fmBlue] (6.335,-2.752) circle (0.055);
\filldraw[fmRed] (6.444,-1.348) circle (0.055);
\filldraw[fmGreen] (7.572,-2.885) circle (0.055);
\filldraw[fmBlue] (6.208,-2.761) circle (0.055);
\filldraw[fmRed] (6.316,-4.136) circle (0.055);
\filldraw[fmGreen] (7.638,-2.730) circle (0.055);
\filldraw[fmBlue] (6.241,-2.718) circle (0.055);
\filldraw[fmRed] (5.974,-2.603) circle (0.055);
\filldraw[fmGreen] (6.311,-3.915) circle (0.055);
\filldraw[fmBlue] (7.230,-2.648) circle (0.055);
\filldraw[fmRed] (6.062,-1.333) circle (0.055);
\filldraw[fmGreen] (7.623,-2.528) circle (0.055);
\filldraw[fmBlue] (6.245,-2.766) circle (0.055);
\filldraw[fmRed] (6.156,-4.119) circle (0.055);
\filldraw[fmGreen] (4.998,-2.606) circle (0.055);
\filldraw[fmBlue] (6.377,-2.720) circle (0.055);
\coordinate (PFc) at (6.3,-2.75); \coordinate (PFtop) at (6.3,-1.3); \coordinate (PFright) at (7.75,-2.75);
\filldraw[fmAmber!12, draw=atomLine, line width=0.6pt] (9.050,-3.600) -- (9.139,-3.600) -- (9.229,-3.600) -- (9.318,-3.600) -- (9.408,-3.600) -- (9.498,-3.600) -- (9.587,-3.600) -- (9.677,-3.600) -- (9.768,-3.600) -- (9.858,-3.600) -- (9.949,-3.600) -- (10.040,-3.600) -- (10.132,-3.600) -- (10.224,-3.600) -- (10.316,-3.600) -- (10.409,-3.600) -- (10.503,-3.600) -- (10.597,-3.600) -- (10.691,-3.600) -- (10.787,-3.600) -- (10.882,-3.600) -- (10.979,-3.600) -- (11.077,-3.600) -- (11.175,-3.600) -- (11.274,-3.600) -- (11.374,-3.600) -- (11.475,-3.600) -- (11.577,-3.600) -- (11.680,-3.600) -- (11.784,-3.600) -- (11.784,-3.600) -- (11.783,-3.518) -- (11.782,-3.436) -- (11.781,-3.354) -- (11.778,-3.272) -- (11.775,-3.190) -- (11.771,-3.107) -- (11.766,-3.025) -- (11.761,-2.943) -- (11.754,-2.860) -- (11.748,-2.778) -- (11.740,-2.695) -- (11.731,-2.613) -- (11.722,-2.530) -- (11.712,-2.447) -- (11.702,-2.364) -- (11.690,-2.281) -- (11.678,-2.198) -- (11.665,-2.114) -- (11.651,-2.031) -- (11.636,-1.947) -- (11.621,-1.863) -- (11.604,-1.779) -- (11.587,-1.695) -- (11.569,-1.610) -- (11.550,-1.525) -- (11.531,-1.440) -- (11.510,-1.355) -- (11.488,-1.270) -- (11.466,-1.184) -- (11.466,-1.184) -- (11.369,-1.281) -- (11.274,-1.376) -- (11.180,-1.470) -- (11.088,-1.562) -- (10.997,-1.653) -- (10.908,-1.742) -- (10.819,-1.831) -- (10.732,-1.918) -- (10.646,-2.004) -- (10.560,-2.090) -- (10.476,-2.174) -- (10.392,-2.258) -- (10.309,-2.341) -- (10.227,-2.423) -- (10.146,-2.504) -- (10.065,-2.585) -- (9.985,-2.665) -- (9.905,-2.745) -- (9.826,-2.824) -- (9.747,-2.903) -- (9.669,-2.981) -- (9.590,-3.060) -- (9.513,-3.137) -- (9.435,-3.215) -- (9.358,-3.292) -- (9.281,-3.369) -- (9.204,-3.446) -- (9.127,-3.523) -- (9.050,-3.600) -- cycle;
\filldraw[atomLine] (9.498,-3.510) circle (0.055);
\filldraw[atomLine] (9.249,-3.511) circle (0.055);
\filldraw[atomLine] (11.195,-2.509) circle (0.055);
\filldraw[atomLine] (9.392,-3.291) circle (0.055);
\filldraw[atomLine] (9.653,-3.539) circle (0.055);
\filldraw[atomLine] (9.202,-3.512) circle (0.055);
\filldraw[atomLine] (9.296,-3.430) circle (0.055);
\filldraw[atomLine, draw=fmAmber!80!black, line width=0.9pt] (10.480,-3.297) circle (0.075);
\filldraw[atomLine] (9.209,-3.596) circle (0.055);
\filldraw[atomLine] (9.462,-3.505) circle (0.055);
\filldraw[atomLine] (9.315,-3.450) circle (0.055);
\filldraw[atomLine] (10.492,-2.846) circle (0.055);
\filldraw[atomLine] (9.307,-3.570) circle (0.055);
\filldraw[atomLine] (9.385,-3.427) circle (0.055);
\coordinate (I100) at (9.050,-3.600);
\coordinate (I110) at (11.784,-3.600);
\coordinate (I111) at (11.466,-1.184);

  \draw[-{Stealth[length=2mm]}, fmAmber!80!black, line width=0.9pt] (Bsel) to[bend right=25] (Csel);
  \node[font=\small] at (Alab) {\textit{а}\ лист (см)};
  \node[font=\small] at (8.0,0.25) {\textit{б}\ зёрна (10--100 мкм)};
  \node[font=\small] at (1.6,-4.75) {\textit{в}\ одно зерно};
  \node[font=\small] at (6.3,-4.75) {\textit{г}\ ППФ $(100)$};
  \node[font=\small] at (10.4,-4.75) {\textit{д}\ ОПФ для НН};
  \node[font=\scriptsize, anchor=south] at (PFtop) {НП};
  \node[font=\scriptsize, anchor=west] at (PFright) {ПН};
  \node[font=\scriptsize, fill=white, inner sep=0.5pt] at ($(PFc)+(0.42,-0.3)$) {НН};
  \node[font=\scriptsize, anchor=north east] at (I100) {$[100]$};
  \node[font=\scriptsize, anchor=north west] at (I110) {$[110]$};
  \node[font=\scriptsize, anchor=south] at (I111) {$[111]$};
\end{tikzpicture}
\caption{Две системы координат при описании текстуры. \textit{а}~—
лист с осями образца НП, ПН, НН (сантиметры); рамкой отмечен
увеличиваемый участок. \textit{б}~— участок поверхности: каждое зерно
изображено кубиком с рёбрами вдоль $\nfam{100}$ (грани $\perp[100]$~—
красные, $\perp[010]$~— зелёные, $\perp[001]$~— синие); большинство зёрен
имеют кубическую текстуру $\plfam{001}\nfam{100}$ с разбросом
$\sim10^\circ$, три зерна ориентированы случайно. \textit{в}~— выделенное
зерно (оранжевая граница) крупно: оси кристалла (цветные стрелки) и оси
образца (пунктир) из одной точки. \textit{г}~— прямая полюсная фигура
$(100)$ в осях образца (НН в центре): каждая точка~— направление одного
ребра одного кубика, её цвет~— какое это ребро; у кубической текстуры
точки собираются у НН, НП и ПН. \textit{д}~— обратная полюсная фигура
для НН в стандартном треугольнике осей кристалла: каждая точка~— одно
зерно; у кубической текстуры НН $\parallel\nfam{100}$, и точки собираются у
вершины $[100]$. Выделенное зерно обведено оранжевым}
\label{fig:texture-frames}
\end{figure}

Все способы описания отличаются тем, \emph{какую} систему держат
неподвижной:
\begin{itemize}
\item \vazhno{прямая полюсная фигура}: неподвижны оси образца; для каждого
зерна отмечают, куда смотрят его направления $\nfam{100}$ (или нормали к
другим плоскостям $\plfam{hkl}$)~— рис.~\ref{fig:texture-frames},\,\textit{г};
\item \vazhno{обратная полюсная фигура}: неподвижны оси кристалла; для
каждого зерна отмечают, куда смотрит выбранная ось образца, например
НН,~— рис.~\ref{fig:texture-frames},\,\textit{д};
\item \vazhno{функция распределения ориентаций}: отмечают сам поворот
(три угла Эйлера) каждого зерна~— точку в трёхмерном пространстве
ориентаций.
\end{itemize}

\begin{masshtab}[текстура]{-9.7/-9.2/{оси кристалла задаёт решётка}, -5.2/-3.8/{зерно}, -3.3/-1.3/{образец и пятно рентгена}}
Оси кристалла задаёт решётка (ячейка~— доли нанометра), но ориентация
одинакова во всём зерне размером 10--100~мкм. Оси образца привязаны к
изделию (лист, проволока~— сантиметры и больше). Полюсные фигуры
измеряют рентгеном или нейтронами по пятну размером в миллиметры: в него
попадают $10^4$--$10^6$ зёрен, поэтому полюсная фигура~— статистика
ориентаций множества зёрен, а не картинка одного кристалла.
\end{masshtab}

\subsubsection{Прямые полюсные фигуры}
\begin{opredelenie}
\term{Прямая полюсная фигура} (ППФ) $(hkl)$~— стереографическая проекция
в \emph{системе координат образца} распределения нормалей к плоскостям
$\plfam{hkl}$ всех кристаллитов.
\end{opredelenie}

Пусть ось текстуры $F\parallel[123]$. Нормали к плоскостям $\plfam{100}$
образуют с $F$ углы $\alpha_1$, $\alpha_2$, $\alpha_3$:
\[
\cos\alpha_1=\tfrac1{\sqrt{14}},\quad \cos\alpha_2=\tfrac2{\sqrt{14}},\quad
\cos\alpha_3=\tfrac3{\sqrt{14}}\qquad(\alpha_1=74{,}5^\circ,\ \alpha_2=57{,}7^\circ,\ \alpha_3=36{,}7^\circ).
\]
При повороте кристаллита вокруг $F$ каждая нормаль описывает конус.
Поэтому если ось $F$ находится в центре проекции, ППФ $(100)$~— три
концентрические окружности, а если на основном круге~— три пары
«параллелей» сетки Вульфа (рис.~\ref{fig:pf123}).

\begin{figure}[ht]
\centering
\begin{tikzpicture}[scale=1.55]
  % создано tools/gen_plastic.py: ППФ (100) аксиальной текстуры F || [123] (r = 1)
\begin{scope}
\draw[atomLine, line width=0.7pt] (0,0) circle (1);
\draw[fmBlue, line width=1.1pt] (0,0) circle (0.7604);
\draw[fmRed, line width=1.1pt] (0,0) circle (0.5508);
\draw[fmGreen!70!black, line width=1.1pt] (0,0) circle (0.3317);
\fill[fmBlue] (0.5201,0.5547) circle (0.045);
\fill[fmRed] (-0.4616,0.3004) circle (0.045);
\fill[fmGreen!70!black] (0.1402,-0.3006) circle (0.045);
\fill (0,0) circle (0.03);
\end{scope}
\begin{scope}[shift={(2.7,0)}]
\draw[atomLine, line width=0.7pt] (0,0) circle (1);
\draw[fmBlue, line width=1.1pt] (0.9636,0.2673) -- (0.9317,0.2586) -- (0.9007,0.2504) -- (0.8706,0.2428) -- (0.8414,0.2357) -- (0.8130,0.2290) -- (0.7852,0.2227) -- (0.7582,0.2167) -- (0.7319,0.2112) -- (0.7062,0.2059) -- (0.6811,0.2010) -- (0.6565,0.1964) -- (0.6324,0.1920) -- (0.6089,0.1879) -- (0.5858,0.1840) -- (0.5632,0.1804) -- (0.5410,0.1769) -- (0.5191,0.1737) -- (0.4977,0.1706) -- (0.4766,0.1677) -- (0.4558,0.1650) -- (0.4354,0.1625) -- (0.4152,0.1601) -- (0.3953,0.1578) -- (0.3757,0.1557) -- (0.3564,0.1538) -- (0.3372,0.1519) -- (0.3183,0.1502) -- (0.2995,0.1486) -- (0.2810,0.1471) -- (0.2626,0.1457) -- (0.2444,0.1444) -- (0.2264,0.1432) -- (0.2084,0.1421) -- (0.1906,0.1411) -- (0.1730,0.1403) -- (0.1554,0.1395) -- (0.1379,0.1387) -- (0.1205,0.1381) -- (0.1031,0.1376) -- (0.0859,0.1371) -- (0.0686,0.1368) -- (0.0514,0.1365) -- (0.0343,0.1363) -- (0.0171,0.1361) -- (0.0000,0.1361) -- (-0.0171,0.1361) -- (-0.0343,0.1363) -- (-0.0514,0.1365) -- (-0.0686,0.1368) -- (-0.0859,0.1371) -- (-0.1031,0.1376) -- (-0.1205,0.1381) -- (-0.1379,0.1387) -- (-0.1554,0.1395) -- (-0.1730,0.1403) -- (-0.1906,0.1411) -- (-0.2084,0.1421) -- (-0.2264,0.1432) -- (-0.2444,0.1444) -- (-0.2626,0.1457) -- (-0.2810,0.1471) -- (-0.2995,0.1486) -- (-0.3183,0.1502) -- (-0.3372,0.1519) -- (-0.3564,0.1538) -- (-0.3757,0.1557) -- (-0.3953,0.1578) -- (-0.4152,0.1601) -- (-0.4354,0.1625) -- (-0.4558,0.1650) -- (-0.4766,0.1677) -- (-0.4977,0.1706) -- (-0.5191,0.1737) -- (-0.5410,0.1769) -- (-0.5632,0.1804) -- (-0.5858,0.1840) -- (-0.6089,0.1879) -- (-0.6324,0.1920) -- (-0.6565,0.1964) -- (-0.6811,0.2010) -- (-0.7062,0.2059) -- (-0.7319,0.2112) -- (-0.7582,0.2167) -- (-0.7852,0.2227) -- (-0.8130,0.2290) -- (-0.8414,0.2357) -- (-0.8706,0.2428) -- (-0.9007,0.2504) -- (-0.9317,0.2586) -- (-0.9636,0.2673);
\draw[fmBlue, line width=1.1pt] (0.9636,-0.2673) -- (0.9317,-0.2586) -- (0.9007,-0.2504) -- (0.8706,-0.2428) -- (0.8414,-0.2357) -- (0.8130,-0.2290) -- (0.7852,-0.2227) -- (0.7582,-0.2167) -- (0.7319,-0.2112) -- (0.7062,-0.2059) -- (0.6811,-0.2010) -- (0.6565,-0.1964) -- (0.6324,-0.1920) -- (0.6089,-0.1879) -- (0.5858,-0.1840) -- (0.5632,-0.1804) -- (0.5410,-0.1769) -- (0.5191,-0.1737) -- (0.4977,-0.1706) -- (0.4766,-0.1677) -- (0.4558,-0.1650) -- (0.4354,-0.1625) -- (0.4152,-0.1601) -- (0.3953,-0.1578) -- (0.3757,-0.1557) -- (0.3564,-0.1538) -- (0.3372,-0.1519) -- (0.3183,-0.1502) -- (0.2995,-0.1486) -- (0.2810,-0.1471) -- (0.2626,-0.1457) -- (0.2444,-0.1444) -- (0.2264,-0.1432) -- (0.2084,-0.1421) -- (0.1906,-0.1411) -- (0.1730,-0.1403) -- (0.1554,-0.1395) -- (0.1379,-0.1387) -- (0.1205,-0.1381) -- (0.1031,-0.1376) -- (0.0859,-0.1371) -- (0.0686,-0.1368) -- (0.0514,-0.1365) -- (0.0343,-0.1363) -- (0.0171,-0.1361) -- (0.0000,-0.1361) -- (-0.0171,-0.1361) -- (-0.0343,-0.1363) -- (-0.0514,-0.1365) -- (-0.0686,-0.1368) -- (-0.0859,-0.1371) -- (-0.1031,-0.1376) -- (-0.1205,-0.1381) -- (-0.1379,-0.1387) -- (-0.1554,-0.1395) -- (-0.1730,-0.1403) -- (-0.1906,-0.1411) -- (-0.2084,-0.1421) -- (-0.2264,-0.1432) -- (-0.2444,-0.1444) -- (-0.2626,-0.1457) -- (-0.2810,-0.1471) -- (-0.2995,-0.1486) -- (-0.3183,-0.1502) -- (-0.3372,-0.1519) -- (-0.3564,-0.1538) -- (-0.3757,-0.1557) -- (-0.3953,-0.1578) -- (-0.4152,-0.1601) -- (-0.4354,-0.1625) -- (-0.4558,-0.1650) -- (-0.4766,-0.1677) -- (-0.4977,-0.1706) -- (-0.5191,-0.1737) -- (-0.5410,-0.1769) -- (-0.5632,-0.1804) -- (-0.5858,-0.1840) -- (-0.6089,-0.1879) -- (-0.6324,-0.1920) -- (-0.6565,-0.1964) -- (-0.6811,-0.2010) -- (-0.7062,-0.2059) -- (-0.7319,-0.2112) -- (-0.7582,-0.2167) -- (-0.7852,-0.2227) -- (-0.8130,-0.2290) -- (-0.8414,-0.2357) -- (-0.8706,-0.2428) -- (-0.9007,-0.2504) -- (-0.9317,-0.2586) -- (-0.9636,-0.2673);
\draw[fmRed, line width=1.1pt] (0.8452,0.5345) -- (0.8204,0.5192) -- (0.7962,0.5048) -- (0.7723,0.4911) -- (0.7488,0.4783) -- (0.7258,0.4661) -- (0.7031,0.4546) -- (0.6808,0.4438) -- (0.6589,0.4335) -- (0.6373,0.4238) -- (0.6161,0.4147) -- (0.5952,0.4060) -- (0.5746,0.3978) -- (0.5543,0.3900) -- (0.5342,0.3827) -- (0.5145,0.3757) -- (0.4950,0.3692) -- (0.4758,0.3630) -- (0.4568,0.3571) -- (0.4381,0.3516) -- (0.4195,0.3464) -- (0.4012,0.3414) -- (0.3831,0.3368) -- (0.3651,0.3324) -- (0.3474,0.3283) -- (0.3298,0.3245) -- (0.3123,0.3208) -- (0.2950,0.3175) -- (0.2779,0.3143) -- (0.2609,0.3114) -- (0.2440,0.3086) -- (0.2272,0.3061) -- (0.2106,0.3038) -- (0.1940,0.3016) -- (0.1775,0.2997) -- (0.1611,0.2979) -- (0.1448,0.2963) -- (0.1285,0.2949) -- (0.1123,0.2937) -- (0.0962,0.2926) -- (0.0801,0.2917) -- (0.0640,0.2910) -- (0.0480,0.2904) -- (0.0320,0.2900) -- (0.0160,0.2898) -- (0.0000,0.2897) -- (-0.0160,0.2898) -- (-0.0320,0.2900) -- (-0.0480,0.2904) -- (-0.0640,0.2910) -- (-0.0801,0.2917) -- (-0.0962,0.2926) -- (-0.1123,0.2937) -- (-0.1285,0.2949) -- (-0.1448,0.2963) -- (-0.1611,0.2979) -- (-0.1775,0.2997) -- (-0.1940,0.3016) -- (-0.2106,0.3038) -- (-0.2272,0.3061) -- (-0.2440,0.3086) -- (-0.2609,0.3114) -- (-0.2779,0.3143) -- (-0.2950,0.3175) -- (-0.3123,0.3208) -- (-0.3298,0.3245) -- (-0.3474,0.3283) -- (-0.3651,0.3324) -- (-0.3831,0.3368) -- (-0.4012,0.3414) -- (-0.4195,0.3464) -- (-0.4381,0.3516) -- (-0.4568,0.3571) -- (-0.4758,0.3630) -- (-0.4950,0.3692) -- (-0.5145,0.3757) -- (-0.5342,0.3827) -- (-0.5543,0.3900) -- (-0.5746,0.3978) -- (-0.5952,0.4060) -- (-0.6161,0.4147) -- (-0.6373,0.4238) -- (-0.6589,0.4335) -- (-0.6808,0.4438) -- (-0.7031,0.4546) -- (-0.7258,0.4661) -- (-0.7488,0.4783) -- (-0.7723,0.4911) -- (-0.7962,0.5048) -- (-0.8204,0.5192) -- (-0.8452,0.5345);
\draw[fmRed, line width=1.1pt] (0.8452,-0.5345) -- (0.8204,-0.5192) -- (0.7962,-0.5048) -- (0.7723,-0.4911) -- (0.7488,-0.4783) -- (0.7258,-0.4661) -- (0.7031,-0.4546) -- (0.6808,-0.4438) -- (0.6589,-0.4335) -- (0.6373,-0.4238) -- (0.6161,-0.4147) -- (0.5952,-0.4060) -- (0.5746,-0.3978) -- (0.5543,-0.3900) -- (0.5342,-0.3827) -- (0.5145,-0.3757) -- (0.4950,-0.3692) -- (0.4758,-0.3630) -- (0.4568,-0.3571) -- (0.4381,-0.3516) -- (0.4195,-0.3464) -- (0.4012,-0.3414) -- (0.3831,-0.3368) -- (0.3651,-0.3324) -- (0.3474,-0.3283) -- (0.3298,-0.3245) -- (0.3123,-0.3208) -- (0.2950,-0.3175) -- (0.2779,-0.3143) -- (0.2609,-0.3114) -- (0.2440,-0.3086) -- (0.2272,-0.3061) -- (0.2106,-0.3038) -- (0.1940,-0.3016) -- (0.1775,-0.2997) -- (0.1611,-0.2979) -- (0.1448,-0.2963) -- (0.1285,-0.2949) -- (0.1123,-0.2937) -- (0.0962,-0.2926) -- (0.0801,-0.2917) -- (0.0640,-0.2910) -- (0.0480,-0.2904) -- (0.0320,-0.2900) -- (0.0160,-0.2898) -- (0.0000,-0.2897) -- (-0.0160,-0.2898) -- (-0.0320,-0.2900) -- (-0.0480,-0.2904) -- (-0.0640,-0.2910) -- (-0.0801,-0.2917) -- (-0.0962,-0.2926) -- (-0.1123,-0.2937) -- (-0.1285,-0.2949) -- (-0.1448,-0.2963) -- (-0.1611,-0.2979) -- (-0.1775,-0.2997) -- (-0.1940,-0.3016) -- (-0.2106,-0.3038) -- (-0.2272,-0.3061) -- (-0.2440,-0.3086) -- (-0.2609,-0.3114) -- (-0.2779,-0.3143) -- (-0.2950,-0.3175) -- (-0.3123,-0.3208) -- (-0.3298,-0.3245) -- (-0.3474,-0.3283) -- (-0.3651,-0.3324) -- (-0.3831,-0.3368) -- (-0.4012,-0.3414) -- (-0.4195,-0.3464) -- (-0.4381,-0.3516) -- (-0.4568,-0.3571) -- (-0.4758,-0.3630) -- (-0.4950,-0.3692) -- (-0.5145,-0.3757) -- (-0.5342,-0.3827) -- (-0.5543,-0.3900) -- (-0.5746,-0.3978) -- (-0.5952,-0.4060) -- (-0.6161,-0.4147) -- (-0.6373,-0.4238) -- (-0.6589,-0.4335) -- (-0.6808,-0.4438) -- (-0.7031,-0.4546) -- (-0.7258,-0.4661) -- (-0.7488,-0.4783) -- (-0.7723,-0.4911) -- (-0.7962,-0.5048) -- (-0.8204,-0.5192) -- (-0.8452,-0.5345);
\draw[fmGreen!70!black, line width=1.1pt] (0.5976,0.8018) -- (0.5850,0.7854) -- (0.5723,0.7697) -- (0.5594,0.7546) -- (0.5464,0.7402) -- (0.5332,0.7264) -- (0.5200,0.7132) -- (0.5066,0.7005) -- (0.4932,0.6884) -- (0.4798,0.6768) -- (0.4663,0.6657) -- (0.4527,0.6551) -- (0.4392,0.6450) -- (0.4256,0.6353) -- (0.4121,0.6261) -- (0.3985,0.6173) -- (0.3849,0.6089) -- (0.3713,0.6010) -- (0.3578,0.5934) -- (0.3443,0.5861) -- (0.3307,0.5793) -- (0.3173,0.5728) -- (0.3038,0.5666) -- (0.2903,0.5607) -- (0.2769,0.5552) -- (0.2635,0.5500) -- (0.2501,0.5451) -- (0.2368,0.5405) -- (0.2235,0.5362) -- (0.2102,0.5321) -- (0.1969,0.5283) -- (0.1837,0.5248) -- (0.1704,0.5216) -- (0.1572,0.5186) -- (0.1441,0.5159) -- (0.1309,0.5134) -- (0.1177,0.5112) -- (0.1046,0.5092) -- (0.0915,0.5075) -- (0.0784,0.5060) -- (0.0653,0.5047) -- (0.0523,0.5037) -- (0.0392,0.5029) -- (0.0261,0.5023) -- (0.0131,0.5020) -- (0.0000,0.5019) -- (-0.0131,0.5020) -- (-0.0261,0.5023) -- (-0.0392,0.5029) -- (-0.0523,0.5037) -- (-0.0653,0.5047) -- (-0.0784,0.5060) -- (-0.0915,0.5075) -- (-0.1046,0.5092) -- (-0.1177,0.5112) -- (-0.1309,0.5134) -- (-0.1441,0.5159) -- (-0.1572,0.5186) -- (-0.1704,0.5216) -- (-0.1837,0.5248) -- (-0.1969,0.5283) -- (-0.2102,0.5321) -- (-0.2235,0.5362) -- (-0.2368,0.5405) -- (-0.2501,0.5451) -- (-0.2635,0.5500) -- (-0.2769,0.5552) -- (-0.2903,0.5607) -- (-0.3038,0.5666) -- (-0.3173,0.5728) -- (-0.3307,0.5793) -- (-0.3443,0.5861) -- (-0.3578,0.5934) -- (-0.3713,0.6010) -- (-0.3849,0.6089) -- (-0.3985,0.6173) -- (-0.4121,0.6261) -- (-0.4256,0.6353) -- (-0.4392,0.6450) -- (-0.4527,0.6551) -- (-0.4663,0.6657) -- (-0.4798,0.6768) -- (-0.4932,0.6884) -- (-0.5066,0.7005) -- (-0.5200,0.7132) -- (-0.5332,0.7264) -- (-0.5464,0.7402) -- (-0.5594,0.7546) -- (-0.5723,0.7697) -- (-0.5850,0.7854) -- (-0.5976,0.8018);
\draw[fmGreen!70!black, line width=1.1pt] (0.5976,-0.8018) -- (0.5850,-0.7854) -- (0.5723,-0.7697) -- (0.5594,-0.7546) -- (0.5464,-0.7402) -- (0.5332,-0.7264) -- (0.5200,-0.7132) -- (0.5066,-0.7005) -- (0.4932,-0.6884) -- (0.4798,-0.6768) -- (0.4663,-0.6657) -- (0.4527,-0.6551) -- (0.4392,-0.6450) -- (0.4256,-0.6353) -- (0.4121,-0.6261) -- (0.3985,-0.6173) -- (0.3849,-0.6089) -- (0.3713,-0.6010) -- (0.3578,-0.5934) -- (0.3443,-0.5861) -- (0.3307,-0.5793) -- (0.3173,-0.5728) -- (0.3038,-0.5666) -- (0.2903,-0.5607) -- (0.2769,-0.5552) -- (0.2635,-0.5500) -- (0.2501,-0.5451) -- (0.2368,-0.5405) -- (0.2235,-0.5362) -- (0.2102,-0.5321) -- (0.1969,-0.5283) -- (0.1837,-0.5248) -- (0.1704,-0.5216) -- (0.1572,-0.5186) -- (0.1441,-0.5159) -- (0.1309,-0.5134) -- (0.1177,-0.5112) -- (0.1046,-0.5092) -- (0.0915,-0.5075) -- (0.0784,-0.5060) -- (0.0653,-0.5047) -- (0.0523,-0.5037) -- (0.0392,-0.5029) -- (0.0261,-0.5023) -- (0.0131,-0.5020) -- (0.0000,-0.5019) -- (-0.0131,-0.5020) -- (-0.0261,-0.5023) -- (-0.0392,-0.5029) -- (-0.0523,-0.5037) -- (-0.0653,-0.5047) -- (-0.0784,-0.5060) -- (-0.0915,-0.5075) -- (-0.1046,-0.5092) -- (-0.1177,-0.5112) -- (-0.1309,-0.5134) -- (-0.1441,-0.5159) -- (-0.1572,-0.5186) -- (-0.1704,-0.5216) -- (-0.1837,-0.5248) -- (-0.1969,-0.5283) -- (-0.2102,-0.5321) -- (-0.2235,-0.5362) -- (-0.2368,-0.5405) -- (-0.2501,-0.5451) -- (-0.2635,-0.5500) -- (-0.2769,-0.5552) -- (-0.2903,-0.5607) -- (-0.3038,-0.5666) -- (-0.3173,-0.5728) -- (-0.3307,-0.5793) -- (-0.3443,-0.5861) -- (-0.3578,-0.5934) -- (-0.3713,-0.6010) -- (-0.3849,-0.6089) -- (-0.3985,-0.6173) -- (-0.4121,-0.6261) -- (-0.4256,-0.6353) -- (-0.4392,-0.6450) -- (-0.4527,-0.6551) -- (-0.4663,-0.6657) -- (-0.4798,-0.6768) -- (-0.4932,-0.6884) -- (-0.5066,-0.7005) -- (-0.5200,-0.7132) -- (-0.5332,-0.7264) -- (-0.5464,-0.7402) -- (-0.5594,-0.7546) -- (-0.5723,-0.7697) -- (-0.5850,-0.7854) -- (-0.5976,-0.8018);
\end{scope}

  \node[font=\small, anchor=north west] at (0.05,-0.02) {$F$};
  \fill (2.7,1) circle (0.035);
  \node[font=\small, above] at (2.7,1.02) {$F$};
  \fill (2.7,-1) circle (0.035);
  \node[font=\small] at (0,-1.25) {\textit{а}};
  \node[font=\small] at (2.7,-1.25) {\textit{б}};
  \begin{scope}[shift={(4.25,0.45)}, font=\footnotesize]
    \draw[fmGreen!70!black, line width=1.1pt] (0,0.4) -- (0.3,0.4) node[right, black]{$\alpha_3=36{,}7^\circ$};
    \draw[fmRed, line width=1.1pt] (0,0) -- (0.3,0) node[right, black]{$\alpha_2=57{,}7^\circ$};
    \draw[fmBlue, line width=1.1pt] (0,-0.4) -- (0.3,-0.4) node[right, black]{$\alpha_1=74{,}5^\circ$};
  \end{scope}
\end{tikzpicture}
\caption{Прямая полюсная фигура $(100)$ аксиальной текстуры
$F\parallel[123]$: \textit{а}~— ось $F$ в центре проекции (точки~—
полюса $\plfam{100}$ одного кристаллита); \textit{б}~— ось $F$ на основном
круге проекции}
\label{fig:pf123}
\end{figure}

\subsubsection{Обратные полюсные фигуры}
\begin{opredelenie}
\term{Обратная полюсная фигура} (ОПФ)~— стереографическая проекция в
\emph{кристаллографической системе координат} распределения выбранного
внешнего направления (оси образца).
\end{opredelenie}

Благодаря симметрии кристалла ОПФ кубического материала достаточно
строить в одном стандартном треугольнике $[100]$--$[110]$--$[111]$.
Для аксиальной текстуры $F\parallel\nfam{123}$ ОПФ направления $F$~—
\emph{одна точка} $[321]$ (тот из 24 эквивалентных векторов
$\nfam{123}$, который попадает в треугольник). Для текстуры
$F\parallel\nfam{110}$ ОПФ направления $G\perp F$~— множество направлений,
перпендикулярных $[110]$ (и эквивалентных им); в стандартном треугольнике
оно занимает две его стороны: $[110]$--$[111]$ и $[111]$--$[100]$
(рис.~\ref{fig:ipf}).

\begin{figure}[ht]
\centering
\begin{tikzpicture}[scale=9]
  % создано tools/gen_plastic.py: стандартный треугольник [100]-[110]-[111] (r = 1)
\fill[fmAmber!15] (0.0000,0.0000) -- (0.0101,0.0000) -- (0.0201,0.0000) -- (0.0302,0.0000) -- (0.0403,0.0000) -- (0.0504,0.0000) -- (0.0605,0.0000) -- (0.0706,0.0000) -- (0.0807,0.0000) -- (0.0909,0.0000) -- (0.1010,0.0000) -- (0.1112,0.0000) -- (0.1214,0.0000) -- (0.1317,0.0000) -- (0.1419,0.0000) -- (0.1522,0.0000) -- (0.1625,0.0000) -- (0.1729,0.0000) -- (0.1833,0.0000) -- (0.1937,0.0000) -- (0.2042,0.0000) -- (0.2147,0.0000) -- (0.2252,0.0000) -- (0.2358,0.0000) -- (0.2465,0.0000) -- (0.2572,0.0000) -- (0.2679,0.0000) -- (0.2788,0.0000) -- (0.2897,0.0000) -- (0.3006,0.0000) -- (0.3116,0.0000) -- (0.3227,0.0000) -- (0.3338,0.0000) -- (0.3451,0.0000) -- (0.3564,0.0000) -- (0.3678,0.0000) -- (0.3792,0.0000) -- (0.3908,0.0000) -- (0.4025,0.0000) -- (0.4142,0.0000) -- (0.4142,0.0000) -- (0.4142,0.0092) -- (0.4141,0.0185) -- (0.4139,0.0277) -- (0.4137,0.0370) -- (0.4135,0.0462) -- (0.4131,0.0555) -- (0.4127,0.0647) -- (0.4123,0.0740) -- (0.4118,0.0833) -- (0.4112,0.0925) -- (0.4105,0.1018) -- (0.4098,0.1111) -- (0.4091,0.1204) -- (0.4083,0.1297) -- (0.4074,0.1390) -- (0.4064,0.1483) -- (0.4054,0.1576) -- (0.4043,0.1669) -- (0.4032,0.1763) -- (0.4020,0.1856) -- (0.4007,0.1950) -- (0.3994,0.2044) -- (0.3980,0.2138) -- (0.3965,0.2232) -- (0.3950,0.2326) -- (0.3934,0.2420) -- (0.3917,0.2514) -- (0.3899,0.2609) -- (0.3881,0.2704) -- (0.3862,0.2798) -- (0.3843,0.2893) -- (0.3823,0.2989) -- (0.3802,0.3084) -- (0.3780,0.3180) -- (0.3758,0.3275) -- (0.3734,0.3371) -- (0.3710,0.3467) -- (0.3686,0.3564) -- (0.3660,0.3660) -- (0.3660,0.3660) -- (0.3551,0.3551) -- (0.3443,0.3443) -- (0.3337,0.3337) -- (0.3232,0.3232) -- (0.3127,0.3127) -- (0.3024,0.3024) -- (0.2923,0.2923) -- (0.2822,0.2822) -- (0.2722,0.2722) -- (0.2623,0.2623) -- (0.2525,0.2525) -- (0.2427,0.2427) -- (0.2331,0.2331) -- (0.2235,0.2235) -- (0.2141,0.2141) -- (0.2046,0.2046) -- (0.1953,0.1953) -- (0.1860,0.1860) -- (0.1768,0.1768) -- (0.1676,0.1676) -- (0.1585,0.1585) -- (0.1494,0.1494) -- (0.1404,0.1404) -- (0.1314,0.1314) -- (0.1224,0.1224) -- (0.1135,0.1135) -- (0.1047,0.1047) -- (0.0958,0.0958) -- (0.0870,0.0870) -- (0.0783,0.0783) -- (0.0695,0.0695) -- (0.0608,0.0608) -- (0.0521,0.0521) -- (0.0434,0.0434) -- (0.0347,0.0347) -- (0.0260,0.0260) -- (0.0173,0.0173) -- (0.0087,0.0087) -- (0.0000,0.0000) -- cycle;
\draw[atomLine, line width=0.6pt] (0.0000,0.0000) -- (0.0101,0.0000) -- (0.0201,0.0000) -- (0.0302,0.0000) -- (0.0403,0.0000) -- (0.0504,0.0000) -- (0.0605,0.0000) -- (0.0706,0.0000) -- (0.0807,0.0000) -- (0.0909,0.0000) -- (0.1010,0.0000) -- (0.1112,0.0000) -- (0.1214,0.0000) -- (0.1317,0.0000) -- (0.1419,0.0000) -- (0.1522,0.0000) -- (0.1625,0.0000) -- (0.1729,0.0000) -- (0.1833,0.0000) -- (0.1937,0.0000) -- (0.2042,0.0000) -- (0.2147,0.0000) -- (0.2252,0.0000) -- (0.2358,0.0000) -- (0.2465,0.0000) -- (0.2572,0.0000) -- (0.2679,0.0000) -- (0.2788,0.0000) -- (0.2897,0.0000) -- (0.3006,0.0000) -- (0.3116,0.0000) -- (0.3227,0.0000) -- (0.3338,0.0000) -- (0.3451,0.0000) -- (0.3564,0.0000) -- (0.3678,0.0000) -- (0.3792,0.0000) -- (0.3908,0.0000) -- (0.4025,0.0000) -- (0.4142,0.0000) -- (0.4142,0.0000) -- (0.4142,0.0092) -- (0.4141,0.0185) -- (0.4139,0.0277) -- (0.4137,0.0370) -- (0.4135,0.0462) -- (0.4131,0.0555) -- (0.4127,0.0647) -- (0.4123,0.0740) -- (0.4118,0.0833) -- (0.4112,0.0925) -- (0.4105,0.1018) -- (0.4098,0.1111) -- (0.4091,0.1204) -- (0.4083,0.1297) -- (0.4074,0.1390) -- (0.4064,0.1483) -- (0.4054,0.1576) -- (0.4043,0.1669) -- (0.4032,0.1763) -- (0.4020,0.1856) -- (0.4007,0.1950) -- (0.3994,0.2044) -- (0.3980,0.2138) -- (0.3965,0.2232) -- (0.3950,0.2326) -- (0.3934,0.2420) -- (0.3917,0.2514) -- (0.3899,0.2609) -- (0.3881,0.2704) -- (0.3862,0.2798) -- (0.3843,0.2893) -- (0.3823,0.2989) -- (0.3802,0.3084) -- (0.3780,0.3180) -- (0.3758,0.3275) -- (0.3734,0.3371) -- (0.3710,0.3467) -- (0.3686,0.3564) -- (0.3660,0.3660) -- (0.3660,0.3660) -- (0.3551,0.3551) -- (0.3443,0.3443) -- (0.3337,0.3337) -- (0.3232,0.3232) -- (0.3127,0.3127) -- (0.3024,0.3024) -- (0.2923,0.2923) -- (0.2822,0.2822) -- (0.2722,0.2722) -- (0.2623,0.2623) -- (0.2525,0.2525) -- (0.2427,0.2427) -- (0.2331,0.2331) -- (0.2235,0.2235) -- (0.2141,0.2141) -- (0.2046,0.2046) -- (0.1953,0.1953) -- (0.1860,0.1860) -- (0.1768,0.1768) -- (0.1676,0.1676) -- (0.1585,0.1585) -- (0.1494,0.1494) -- (0.1404,0.1404) -- (0.1314,0.1314) -- (0.1224,0.1224) -- (0.1135,0.1135) -- (0.1047,0.1047) -- (0.0958,0.0958) -- (0.0870,0.0870) -- (0.0783,0.0783) -- (0.0695,0.0695) -- (0.0608,0.0608) -- (0.0521,0.0521) -- (0.0434,0.0434) -- (0.0347,0.0347) -- (0.0260,0.0260) -- (0.0173,0.0173) -- (0.0087,0.0087) -- (0.0000,0.0000) -- cycle;
\draw[fmBlue, line width=2.2pt, line cap=round] (0.4142,0.0000) -- (0.4142,0.0092) -- (0.4141,0.0185) -- (0.4139,0.0277) -- (0.4137,0.0370) -- (0.4135,0.0462) -- (0.4131,0.0555) -- (0.4127,0.0647) -- (0.4123,0.0740) -- (0.4118,0.0833) -- (0.4112,0.0925) -- (0.4105,0.1018) -- (0.4098,0.1111) -- (0.4091,0.1204) -- (0.4083,0.1297) -- (0.4074,0.1390) -- (0.4064,0.1483) -- (0.4054,0.1576) -- (0.4043,0.1669) -- (0.4032,0.1763) -- (0.4020,0.1856) -- (0.4007,0.1950) -- (0.3994,0.2044) -- (0.3980,0.2138) -- (0.3965,0.2232) -- (0.3950,0.2326) -- (0.3934,0.2420) -- (0.3917,0.2514) -- (0.3899,0.2609) -- (0.3881,0.2704) -- (0.3862,0.2798) -- (0.3843,0.2893) -- (0.3823,0.2989) -- (0.3802,0.3084) -- (0.3780,0.3180) -- (0.3758,0.3275) -- (0.3734,0.3371) -- (0.3710,0.3467) -- (0.3686,0.3564) -- (0.3660,0.3660) -- (0.3551,0.3551) -- (0.3443,0.3443) -- (0.3337,0.3337) -- (0.3232,0.3232) -- (0.3127,0.3127) -- (0.3024,0.3024) -- (0.2923,0.2923) -- (0.2822,0.2822) -- (0.2722,0.2722) -- (0.2623,0.2623) -- (0.2525,0.2525) -- (0.2427,0.2427) -- (0.2331,0.2331) -- (0.2235,0.2235) -- (0.2141,0.2141) -- (0.2046,0.2046) -- (0.1953,0.1953) -- (0.1860,0.1860) -- (0.1768,0.1768) -- (0.1676,0.1676) -- (0.1585,0.1585) -- (0.1494,0.1494) -- (0.1404,0.1404) -- (0.1314,0.1314) -- (0.1224,0.1224) -- (0.1135,0.1135) -- (0.1047,0.1047) -- (0.0958,0.0958) -- (0.0870,0.0870) -- (0.0783,0.0783) -- (0.0695,0.0695) -- (0.0608,0.0608) -- (0.0521,0.0521) -- (0.0434,0.0434) -- (0.0347,0.0347) -- (0.0260,0.0260) -- (0.0173,0.0173) -- (0.0087,0.0087) -- (0.0000,0.0000);
\coordinate (T100) at (0.0000,0.0000);
\coordinate (T110) at (0.4142,0.0000);
\coordinate (T111) at (0.3660,0.3660);
\coordinate (T321) at (0.2967,0.1483);
\coordinate (T210) at (0.2361,0.0000);
\coordinate (T211) at (0.2247,0.2247);
\coordinate (T221) at (0.4000,0.2000);
\coordinate (T310) at (0.1623,0.0000);
\coordinate (T311) at (0.1583,0.1583);

  \foreach \p/\a/\t in {T100/north east/{[100]}, T110/north west/{[110]}, T111/south/{[111]}}{
    \fill[atomLine] (\p) circle (0.006);
    \node[font=\small, anchor=\a] at (\p) {$\t$};}
  \fill[fmRed] (T321) circle (0.009) node[below=1pt, font=\small, fmRed]{$[321]$};
  \foreach \p in {T211, T210}{\fill[atomLine!60] (\p) circle (0.004);}
  \node[font=\scriptsize, atomLine!80, anchor=north] at (T210) {$[210]$};
  \node[font=\scriptsize, atomLine!80, anchor=south east] at (T211) {$[211]$};
  \node[font=\footnotesize, fmBlue, anchor=west, align=left] at (0.43,0.22)
    {$G\perp F$ для\\текстуры $\nfam{110}$};
\end{tikzpicture}
\caption{Обратные полюсные фигуры в стандартном треугольнике: ось
текстуры $F\parallel\nfam{123}$ изображается точкой $[321]$ (красная);
направления $G\perp F$ при текстуре $F\parallel\nfam{110}$~— сторонами
$[110]$--$[111]$--$[100]$ (синие)}
\label{fig:ipf}
\end{figure}

\subsubsection{Функция распределения ориентаций}
Полюсные фигуры~— двумерные «проекции» трёхмерного распределения
ориентаций, и по одной ППФ или ОПФ ориентацию кристаллитов однозначно не
восстановить. Полное описание даёт ФРО.

\begin{opredelenie}
\term{Функция распределения ориентаций} $f(g)=f(\varphi,\theta,\psi)$~—
плотность вероятности найти кристаллит с ориентацией $g$, заданной
углами Эйлера~\eqref{eq:euler}: объёмная доля кристаллитов с ориентациями
между $g$ и $g+\mathrm dg$ равна $\Delta V/V=f(g)\,\mathrm dg$, где
$\mathrm dg=\sin\theta\,\mathrm d\varphi\,\mathrm d\theta\,\mathrm d\psi/8\pi^2$.
\end{opredelenie}

Углы Эйлера меняются в пределах $0\le\varphi<2\pi$, $0\le\theta\le\pi$,
$0\le\psi<2\pi$, т.\,е. ФРО задана в \term{ориентационном пространстве}~—
«параллелепипеде» $2\pi\times\pi\times2\pi$. Его геометрический объём
$4\pi^3\approx12\pi^2$, но объём с учётом веса $\sin\theta$ равен
\[
\int_0^{2\pi}\!\!\mathrm d\varphi\int_0^\pi\!\sin\theta\,\mathrm d\theta\int_0^{2\pi}\!\!\mathrm d\psi=8\pi^2 .
\]
Вес $\sin\theta$ появляется потому, что ФРО задана на непрерывной группе
вращений, а равные интервалы углов Эйлера соответствуют неравным
«объёмам» ориентаций: при $\theta\to0$ углы $\varphi$ и $\psi$ описывают
почти один и тот же поворот.

\begin{figure}[ht]
\centering
\begin{tikzpicture}[cubeproj, scale=1.0]
  \def\Lx{4} \def\Ly{4} \def\Lz{2}
  \draw[atomLine!50, dashed] (0,\Ly,0) -- (0,0,0) (0,\Ly,0) -- (\Lx,\Ly,0) (0,\Ly,0) -- (0,\Ly,\Lz);
  \draw[atomLine] (0,0,0) -- (\Lx,0,0) -- (\Lx,\Ly,0) -- (\Lx,\Ly,\Lz) -- (0,\Ly,\Lz) -- (0,0,\Lz) -- cycle
    (0,0,\Lz) -- (\Lx,0,\Lz) -- (\Lx,\Ly,\Lz) (\Lx,0,0) -- (\Lx,0,\Lz);
  \draw[arr] (0,0,0) -- (4.6,0,0) node[right, font=\small]{$\varphi$};
  \draw[arr] (0,0,0) -- (0,0,2.6) node[above, font=\small]{$\theta$};
  \draw[arr, atomLine!60] (0,\Ly,0) -- (0,4.7,0) node[above right=-2pt, font=\small, black]{$\psi$};
  \node[font=\scriptsize, below] at (\Lx,0,0) {$2\pi$};
  \node[font=\scriptsize, left] at (0,0,\Lz) {$\pi$};
  \node[font=\scriptsize, right] at (\Lx,\Ly,0) {$2\pi$};
  % стержень: θ = 36,7°, ψ = 26,6°
  \def\th{0.408} \def\ps{0.295}
  \draw[fmGray, densely dotted] (0,\ps,0) -- (0,\ps,\th) (0,0,\th) -- (0,\ps,\th);
  \draw[fmRed, line width=2.2pt] (0,\ps,\th) -- (\Lx,\ps,\th);
  \node[font=\footnotesize, fmRed, anchor=south] at (2.4,\ps,\th) {$\theta,\psi=\text{const}$, $\varphi$~— любое};
\end{tikzpicture}
\caption{ФРО аксиальной текстуры $F\parallel[123]$ в пространстве углов
Эйлера~— стержень, параллельный оси $\varphi$ ($\cos\theta=3/\sqrt{14}$,
$\operatorname{tg}\psi=1/2$; ещё 23 симметрично эквивалентных стержня не показаны)}
\label{fig:odf}
\end{figure}

Для аксиальной текстуры удобно совместить ось $F$ с осью $Z$ образца.
Первые два поворота Эйлера ($\varphi$ вокруг $Z$, затем $\theta$) задают,
где окажется ось $[001]$ кристалла, а последний~— поворот кристалла
вокруг неё. По третьей строке матрицы~\eqref{eq:euler} ось $Z$ образца
имеет в осях кристалла направляющие косинусы
$(\sin\theta\sin\psi,\ \sin\theta\cos\psi,\ \cos\theta)$. Если
$Z\parallel[mnp]$, то
\begin{equation}\label{eq:fiber-euler}
\cos\theta=\frac{p}{\sqrt{m^2+n^2+p^2}},\qquad \operatorname{tg}\psi=\frac mn ,
\end{equation}
а угол $\varphi$ (поворот вокруг оси текстуры) любой. Значит, \vazhno{ФРО
аксиальной текстуры~— стержень}: $\varphi$ меняется, $\theta$ и $\psi$
постоянны (рис.~\ref{fig:odf}). Поскольку у кубического кристалла
24 поворотных элемента симметрии, систему координат кристалла можно
выбрать 24 способами, и в полном ориентационном пространстве таких
стержней 24 (симметрично эквивалентных).

\begin{itogi}
\begin{itemize}
\item Пластическая деформация идёт скольжением (сдвиг на целое число
периодов, решётка не меняется) или двойникованием (сдвиг на доли периода,
решётка зеркально отражается).
\item Закон Шмида: скольжение начинается при
$\tau=\sigma\cos\chi\cos\lambda=\tau_\text{кр}$; фактор Шмида $\le\frac12$.
\item Тензор пластической дисторсии $u^0_{ij}=\alpha m_ih_j$: след нулевой
(объём сохраняется), антисимметричная часть~— поворот решётки. Для
произвольного формоизменения нужны 5 независимых систем.
\item ГЦК: 12 систем $\plfam{111}\nfam{110}$; ОЦК: направление $\nfam{111}$,
плоскости $\plfam{110}$, $\plfam{112}$, $\plfam{123}$; ГПУ: направление
$\nfam{11\br20}$, плоскость зависит от $c/a$.
\item Элементы двойникования $K_1,K_2,\eta_1,\eta_2,s$; в ГЦК и ОЦК
$s=1/\sqrt2$.
\item Текстура: аксиальная или ограниченная; описание~— ППФ (нормали в
осях образца), ОПФ (оси образца в осях кристалла), ФРО в пространстве
углов Эйлера объёмом $8\pi^2$.
\end{itemize}
\end{itogi}

\begin{proverka}
\item Почему при скольжении решётка не меняется, а при двойниковании
меняется?
\item При какой ориентации оси растяжения фактор Шмида равен
$\frac12$? Может ли он быть больше?
\item Найдите фактор Шмида системы $(\br111)[101]$ ГЦК-кристалла при
растяжении вдоль $[001]$. Сколько систем имеют такой же фактор?
\item Почему ГПУ-металлы с одним базисным скольжением хрупки в
поликристаллическом состоянии?
\item Вычислите угол $2\phi$ между плоскостями $K_1$ и $K_2$ двойникования
ОЦК-кристалла и проверьте формулу~\eqref{eq:twin-shear}.
\item Какой будет ППФ $(100)$ для аксиальной текстуры $F\parallel[001]$?
\end{proverka}

\section{Кристаллография границ зёрен}\label{sec:boundaries}
% ----------------------------------------------------------------------
%  1.7. Кристаллография границ зёрен
% ----------------------------------------------------------------------
\begin{orientir}
Граница зерна~— место встречи двух одинаковых решёток с разной
ориентацией. Сначала выясним, \vazhno{сколько параметров} задают границу и
чем границы наклона отличаются от границ кручения. Затем разберём
\vazhno{малоугловые} границы~— стенки и сетки дислокаций (1.7.1) и
\vazhno{высокоугловые}, строение которых описывают решётки совпадающих
узлов (1.7.2).
\end{orientir}

Во всяком поликристалле есть внутренние границы: между зёрнами одной фазы
(\term{границы зёрен}) и между разными фазами (\term{межфазные границы}).
Граница~— двумерный дефект: в двух измерениях её размеры макроскопические,
а в третьем (толщина)~— атомные. Поэтому граница~— естественный
\emph{наноструктурный} элемент материала.

\begin{masshtab}[граница зерна]{-9.5/-9/{толщина 0,5--1 нм}, -8/-7.4/{дислокации в субгранице}, -5.5/-3.7/{размер зерна}}
Граница тонкая~— несколько межатомных расстояний (0,5--1~нм), но
протяжённая~— во весь размер зерна (1--100~мкм). В малоугловой
субгранице дислокации отстоят на $D=b/\theta$: при $\theta=1^\circ$ это
около 15~нм. Суммарная площадь границ огромна: в 1~см$^3$ металла с
зерном 50~мкм она около $2/d\approx400$~см$^2$. По этой «внутренней
поверхности» быстрее идёт диффузия и на ней скапливаются примеси.
\end{masshtab}

\paragraph{Сколько параметров задают границу}
Рассмотрим простейший случай~— плоскую границу в бикристалле. С
макроскопической точки зрения нужны \vazhno{5 параметров}:
\begin{itemize}
\item 3~— для взаимного разворота кристаллов: три угла Эйлера
$(\varphi,\theta,\psi)$ или угол $\theta$ и ось разворота $\vb u\parallel[mnp]$
(у единичного вектора оси независимы лишь две компоненты, так как
$l_1^2+l_2^2+l_3^2=1$);
\item 2~— для нормали $\vb n$ к плоскости границы.
\end{itemize}
На атомном уровне добавляется вектор $\vb t=t_1\vb a_1+t_2\vb a_2+t_3\vb a_3$
относительного сдвига решёток (3 параметра в долях трансляций), т.\,е.
всего \vazhno{8 параметров}.

\begin{opredelenie}
Если ось разворота лежит в плоскости границы ($\vb u\perp\vb n$), граница
называется \term{границей наклона}; если ось перпендикулярна границе
($\vb u\parallel\vb n$)~— \term{границей кручения} (рис.~\ref{fig:tilt-twist}).
Границы с углом разориентации $\theta\lesssim10$--$15^\circ$ называют
\term{малоугловыми} (это границы субзёрен), с бо́льшим~—
\term{высокоугловыми}.
\end{opredelenie}

\begin{figure}[ht]
\centering
\begin{tikzpicture}[cubeproj, scale=1.1]
  % (а) граница наклона: граница — плоскость x = 0, ось u — вдоль y
  \begin{scope}
    \foreach \x/\c in {-2/fmBlueBg, 0/fmRedBg}{
      \fill[\c] (\x,0,0) -- (\x+2,0,0) -- (\x+2,0,2) -- (\x,0,2) -- cycle;
      \fill[\c!70!white] (\x,0,2) -- (\x+2,0,2) -- (\x+2,1.6,2) -- (\x,1.6,2) -- cycle;}
    \fill[fmRedBg!60!white] (2,0,0) -- (2,1.6,0) -- (2,1.6,2) -- (2,0,2) -- cycle;
    \begin{scope}
      \clip (-2,0,0) -- (0,0,0) -- (0,0,2) -- (-2,0,2) -- cycle;
      \foreach \k in {0,...,16}{\draw[fmBlue!70, line width=0.5pt] ({-3+0.3*\k},0,-0.5) -- ++({3*cos(78)},0,{3*sin(78)});}
    \end{scope}
    \begin{scope}
      \clip (0,0,0) -- (2,0,0) -- (2,0,2) -- (0,0,2) -- cycle;
      \foreach \k in {0,...,16}{\draw[fmRed!70, line width=0.5pt] ({-0.8+0.3*\k},0,-0.5) -- ++({3*cos(102)},0,{3*sin(102)});}
    \end{scope}
    \draw[atomLine] (-2,0,0) -- (2,0,0) -- (2,0,2) -- (-2,0,2) -- cycle
      (-2,0,2) -- (-2,1.6,2) -- (2,1.6,2) -- (2,0,2) (2,0,0) -- (2,1.6,0) -- (2,1.6,2);
    \draw[fmGray, line width=1.3pt] (0,0,0) -- (0,0,2) -- (0,1.6,2);
    \draw[bvec, fmGreen!70!black] (0,-0.2,2.3) -- (0,1.9,2.3) node[above right=-2pt, font=\small]{$\vb u$};
    \draw[bvec, fmBlue] (0,0,2.6) -- (0.9,0,2.6) node[right, font=\small]{$\vb n$};
    \node[font=\small] at (0,0,-0.5) {\textit{а}};
  \end{scope}
  % (б) граница кручения: граница — плоскость z = 1, ось u — вдоль z (раздвинуто)
  \begin{scope}[shift={(4.6,0,0)}]
    \foreach \z/\c/\a in {0/fmBlueBg/-10, 1.6/fmRedBg/10}{
      \fill[\c] (0,0,\z) -- (2,0,\z) -- (2,0,\z+1) -- (0,0,\z+1) -- cycle;
      \fill[\c!60!white] (2,0,\z) -- (2,2,\z) -- (2,2,\z+1) -- (2,0,\z+1) -- cycle;
      \fill[\c!80!white] (0,0,\z+1) -- (2,0,\z+1) -- (2,2,\z+1) -- (0,2,\z+1) -- cycle;
      \begin{scope}
        \clip (0,0,\z+1) -- (2,0,\z+1) -- (2,2,\z+1) -- (0,2,\z+1) -- cycle;
        \foreach \k in {-6,...,6}{
          \draw[atomLine!60, line width=0.4pt] ({1+0.3*\k*cos(\a)-2*sin(\a)},{1+0.3*\k*sin(\a)+2*cos(\a)},\z+1)
            -- ({1+0.3*\k*cos(\a)+2*sin(\a)},{1+0.3*\k*sin(\a)-2*cos(\a)},\z+1);}
      \end{scope}
      \draw[atomLine] (0,0,\z) -- (2,0,\z) -- (2,0,\z+1) -- (0,0,\z+1) -- cycle
        (2,0,\z) -- (2,2,\z) -- (2,2,\z+1) -- (0,2,\z+1) -- (0,0,\z+1) (2,0,\z+1) -- (2,2,\z+1);}
    \draw[bvec, fmGreen!70!black] (1,1,2.6) -- (1,1,3.6) node[above, font=\small]{$\vb u\parallel\vb n$};
    \draw[fmGreen!70!black, dashed] (1,1,1.0) -- (1,1,2.6);
    \node[font=\footnotesize, anchor=west] at (2.15,0,1.3) {граница};
    \draw[fmGray, densely dotted] (2,0,1.0) -- (2,0,1.6);
    \node[font=\small] at (1,0,-0.5) {\textit{б}};
  \end{scope}
\end{tikzpicture}
\caption{Граница наклона (\textit{а}): ось разворота $\vb u$ лежит в плоскости
границы, атомные плоскости двух зёрен наклонены в разные стороны.
Граница кручения (\textit{б}, зёрна раздвинуты): ось разворота
перпендикулярна границе, и решётки двух зёрен повёрнуты относительно друг
друга в плоскости границы}
\label{fig:tilt-twist}
\end{figure}

\subsection{Малоугловые границы}\label{ssec:low-angle}

Малоугловые границы~— \emph{дислокационные} структуры (подробнее~— гл.~2,
п.~2.2.5). \term{Симметричная граница наклона}~— стенка параллельных
краевых дислокаций одного знака (рис.~\ref{fig:lab-wall}). Если расстояние
между дислокациями $D$, вектор Бюргерса $b$ и угол разориентации $\theta$,
то $\sin(\theta/2)=b/2D$, а при малых углах
\begin{keyeq}\label{eq:wall-spacing}
D=\frac b\theta .
\end{keyeq}
Чем больше угол, тем ближе дислокации друг к другу. При $\theta\gtrsim10$--$15^\circ$
расстояние $D$ становится порядка нескольких $b$, ядра дислокаций
сливаются, и дислокационная модель теряет смысл~— граница становится
высокоугловой.

\begin{figure}[ht]
\centering
\begin{tikzpicture}[scale=0.42]
  % Сгенерировано tools/gen_tilt_boundary.py — не редактировать вручную
\begin{scope}
\clip (-6.30,-6.30) rectangle (6.30,6.30);
\draw[atomLine!55, line width=0.45pt] (-7.059,-0.102)--(-7.178,0.897) (-7.059,-0.102)--(-6.940,-1.008) (-7.059,-0.102)--(-6.060,-0.008) (-7.178,0.897)--(-7.297,1.896) (-7.178,0.897)--(-6.180,0.992) (-7.297,1.896)--(-6.299,1.990) (-6.535,3.893)--(-6.653,4.890) (-6.535,3.893)--(-6.419,2.989) (-6.535,3.893)--(-5.539,3.987) (-6.653,4.890)--(-6.772,5.887) (-6.653,4.890)--(-5.658,4.984) (-6.772,5.887)--(-6.890,6.884) (-6.772,5.887)--(-5.778,5.980) (-6.890,6.884)--(-5.897,6.976) (-7.222,-7.020)--(-6.216,-6.922) (-6.581,-4.011)--(-6.701,-3.010) (-6.581,-4.011)--(-6.457,-4.916) (-6.581,-4.011)--(-5.578,-3.913) (-6.701,-3.010)--(-6.820,-2.008) (-6.701,-3.010)--(-5.699,-2.911) (-6.820,-2.008)--(-6.940,-1.008) (-6.820,-2.008)--(-5.819,-1.910) (-6.940,-1.008)--(-5.940,-0.910) (-6.060,-0.008)--(-6.180,0.992) (-6.060,-0.008)--(-5.940,-0.910) (-6.060,-0.008)--(-5.060,0.090) (-6.180,0.992)--(-6.299,1.990) (-6.180,0.992)--(-5.181,1.090) (-6.299,1.990)--(-6.419,2.989) (-6.299,1.990)--(-5.301,2.089) (-6.419,2.989)--(-5.422,3.087) (-5.539,3.987)--(-5.658,4.984) (-5.539,3.987)--(-5.422,3.087) (-5.539,3.987)--(-4.543,4.084) (-5.658,4.984)--(-5.778,5.980) (-5.658,4.984)--(-4.663,5.081) (-5.778,5.980)--(-5.897,6.976) (-5.778,5.980)--(-4.784,6.078) (-5.897,6.976)--(-4.904,7.073) (-6.216,-6.922)--(-6.337,-5.919) (-6.216,-6.922)--(-5.211,-6.820) (-6.337,-5.919)--(-6.457,-4.916) (-6.337,-5.919)--(-5.331,-5.816) (-6.457,-4.916)--(-5.453,-4.814) (-5.578,-3.913)--(-5.699,-2.911) (-5.578,-3.913)--(-5.453,-4.814) (-5.578,-3.913)--(-4.575,-3.811) (-5.699,-2.911)--(-5.819,-1.910) (-5.699,-2.911)--(-4.697,-2.808) (-5.819,-1.910)--(-5.940,-0.910) (-5.819,-1.910)--(-4.817,-1.807) (-5.940,-0.910)--(-4.939,-0.807) (-5.060,0.090)--(-5.181,1.090) (-5.060,0.090)--(-4.939,-0.807) (-5.060,0.090)--(-4.061,0.193) (-5.181,1.090)--(-5.301,2.089) (-5.181,1.090)--(-4.183,1.193) (-5.301,2.089)--(-5.422,3.087) (-5.301,2.089)--(-4.303,2.192) (-5.422,3.087)--(-4.425,3.189) (-4.543,4.084)--(-4.663,5.081) (-4.543,4.084)--(-4.425,3.189) (-4.543,4.084)--(-3.547,4.186) (-4.663,5.081)--(-4.784,6.078) (-4.663,5.081)--(-3.669,5.184) (-4.784,6.078)--(-4.904,7.073) (-4.784,6.078)--(-3.789,6.180) (-4.904,7.073)--(-3.910,7.174) (-5.211,-6.820)--(-5.331,-5.816) (-5.211,-6.820)--(-4.208,-6.713) (-5.331,-5.816)--(-5.453,-4.814) (-5.331,-5.816)--(-4.325,-5.709) (-5.453,-4.814)--(-4.447,-4.708) (-4.575,-3.811)--(-4.697,-2.808) (-4.575,-3.811)--(-4.447,-4.708) (-4.575,-3.811)--(-3.574,-3.706) (-4.697,-2.808)--(-4.817,-1.807) (-4.697,-2.808)--(-3.697,-2.700) (-4.817,-1.807)--(-4.939,-0.807) (-4.817,-1.807)--(-3.814,-1.699) (-4.939,-0.807)--(-3.937,-0.701) (-4.061,0.193)--(-4.183,1.193) (-4.061,0.193)--(-3.937,-0.701) (-4.061,0.193)--(-3.063,0.299) (-4.183,1.193)--(-4.303,2.192) (-4.183,1.193)--(-3.186,1.301) (-4.303,2.192)--(-4.425,3.189) (-4.303,2.192)--(-3.303,2.300) (-4.425,3.189)--(-3.426,3.294) (-3.547,4.186)--(-3.669,5.184) (-3.547,4.186)--(-3.426,3.294) (-3.547,4.186)--(-2.553,4.292) (-3.669,5.184)--(-3.789,6.180) (-3.669,5.184)--(-2.675,5.291) (-3.789,6.180)--(-3.910,7.174) (-3.789,6.180)--(-2.792,6.287) (-3.910,7.174)--(-2.915,7.279) (-4.208,-6.713)--(-4.325,-5.709) (-4.208,-6.713)--(-3.208,-6.598) (-4.325,-5.709)--(-4.447,-4.708) (-4.325,-5.709)--(-3.315,-5.594) (-4.447,-4.708)--(-3.438,-4.601) (-3.574,-3.706)--(-3.697,-2.700) (-3.574,-3.706)--(-3.438,-4.601) (-3.574,-3.706)--(-2.578,-3.598) (-3.697,-2.700)--(-3.814,-1.699) (-3.697,-2.700)--(-2.701,-2.585) (-3.814,-1.699)--(-3.937,-0.701) (-3.814,-1.699)--(-2.807,-1.583) (-3.937,-0.701)--(-2.930,-0.594) (-3.063,0.299)--(-3.186,1.301) (-3.063,0.299)--(-2.930,-0.594) (-3.063,0.299)--(-2.070,0.406) (-3.186,1.301)--(-3.303,2.300) (-3.186,1.301)--(-2.193,1.417) (-3.303,2.300)--(-3.426,3.294) (-3.303,2.300)--(-2.299,2.415) (-3.426,3.294)--(-2.422,3.402) (-2.553,4.292)--(-2.675,5.291) (-2.553,4.292)--(-2.422,3.402) (-2.553,4.292)--(-1.562,4.399) (-2.675,5.291)--(-2.792,6.287) (-2.675,5.291)--(-1.685,5.406) (-2.792,6.287)--(-2.915,7.279) (-2.792,6.287)--(-1.792,6.402) (-3.208,-6.598)--(-3.315,-5.594) (-3.208,-6.598)--(-2.221,-6.473) (-3.315,-5.594)--(-3.438,-4.601) (-3.315,-5.594)--(-2.293,-5.469) (-3.438,-4.601)--(-2.411,-4.495) (-2.578,-3.598)--(-2.701,-2.585) (-2.578,-3.598)--(-2.411,-4.495) (-2.578,-3.598)--(-1.598,-3.492) (-2.701,-2.585)--(-2.807,-1.583) (-2.701,-2.585)--(-1.716,-2.460) (-2.807,-1.583)--(-2.930,-0.594) (-2.807,-1.583)--(-1.789,-1.458) (-2.930,-0.594)--(-1.906,-0.487) (-2.070,0.406)--(-2.193,1.417) (-2.070,0.406)--(-1.906,-0.487) (-2.070,0.406)--(-1.094,0.513) (-2.193,1.417)--(-2.299,2.415) (-2.193,1.417)--(-1.211,1.542) (-2.299,2.415)--(-2.422,3.402) (-2.299,2.415)--(-1.284,2.540) (-2.422,3.402)--(-1.402,3.508) (-1.562,4.399)--(-1.685,5.406) (-1.562,4.399)--(-1.402,3.508) (-1.562,4.399)--(-0.589,4.505) (-1.685,5.406)--(-1.792,6.402) (-1.685,5.406)--(-0.707,5.531) (-1.792,6.402)--(-0.779,6.527) (-2.221,-6.473)--(-2.293,-5.469) (-2.221,-6.473)--(-1.244,-6.356) (-2.293,-5.469)--(-2.411,-4.495) (-2.293,-5.469)--(-1.261,-5.352) (-2.411,-4.495)--(-1.317,-4.365) (-1.598,-3.492)--(-1.716,-2.460) (-1.598,-3.492)--(-1.317,-4.365) (-1.598,-3.492)--(-0.686,-3.362) (-1.716,-2.460)--(-1.789,-1.458) (-1.716,-2.460)--(-0.742,-2.343) (-1.789,-1.458)--(-1.906,-0.487) (-1.789,-1.458)--(-0.759,-1.341) (-1.906,-0.487)--(-0.816,-0.357) (-1.094,0.513)--(-1.211,1.542) (-1.094,0.513)--(-0.816,-0.357) (-1.094,0.513)--(-0.184,0.643) (-1.211,1.542)--(-1.284,2.540) (-1.211,1.542)--(-0.241,1.659) (-1.284,2.540)--(-1.402,3.508) (-1.284,2.540)--(-0.258,2.657) (-1.402,3.508)--(-0.314,3.638) (-0.589,4.505)--(-0.707,5.531) (-0.589,4.505)--(-0.314,3.638) (-0.589,4.505)--(0.317,4.635) (-0.707,5.531)--(-0.779,6.527) (-0.707,5.531)--(0.261,5.648) (-0.779,6.527)--(0.244,6.644) (-1.244,-6.356)--(-1.261,-5.352) (-1.244,-6.356)--(-0.256,-6.356) (-1.261,-5.352)--(-1.317,-4.365) (-1.261,-5.352)--(-0.239,-5.352) (-1.317,-4.365)--(-0.686,-3.362) (-1.317,-4.365)--(-0.183,-4.365) (-0.686,-3.362)--(-0.742,-2.343) (-0.686,-3.362)--(-0.183,-4.365) (-0.686,-3.362)--(0.186,-3.362) (-0.742,-2.343)--(-0.759,-1.341) (-0.742,-2.343)--(0.242,-2.343) (-0.759,-1.341)--(-0.816,-0.357) (-0.759,-1.341)--(0.259,-1.341) (-0.816,-0.357)--(-0.184,0.643) (-0.816,-0.357)--(0.316,-0.357) (-0.184,0.643)--(-0.241,1.659) (-0.184,0.643)--(0.316,-0.357) (-0.184,0.643)--(0.684,0.643) (-0.241,1.659)--(-0.258,2.657) (-0.241,1.659)--(0.741,1.659) (-0.258,2.657)--(-0.314,3.638) (-0.258,2.657)--(0.758,2.657) (-0.314,3.638)--(0.317,4.635) (-0.314,3.638)--(0.814,3.638) (0.317,4.635)--(0.261,5.648) (0.317,4.635)--(0.814,3.638) (0.317,4.635)--(1.183,4.635) (0.261,5.648)--(0.244,6.644) (0.261,5.648)--(1.239,5.648) (0.244,6.644)--(1.256,6.644) (-0.256,-6.356)--(-0.239,-5.352) (-0.256,-6.356)--(0.721,-6.473) (-0.239,-5.352)--(-0.183,-4.365) (-0.239,-5.352)--(0.793,-5.469) (-0.183,-4.365)--(0.186,-3.362) (-0.183,-4.365)--(0.911,-4.495) (0.186,-3.362)--(0.242,-2.343) (0.186,-3.362)--(1.098,-3.492) (0.242,-2.343)--(0.259,-1.341) (0.242,-2.343)--(1.216,-2.460) (0.259,-1.341)--(0.316,-0.357) (0.259,-1.341)--(1.289,-1.458) (0.316,-0.357)--(0.684,0.643) (0.316,-0.357)--(1.406,-0.487) (0.684,0.643)--(0.741,1.659) (0.684,0.643)--(1.594,0.513) (0.741,1.659)--(0.758,2.657) (0.741,1.659)--(1.711,1.542) (0.758,2.657)--(0.814,3.638) (0.758,2.657)--(1.784,2.540) (0.814,3.638)--(1.183,4.635) (0.814,3.638)--(1.902,3.508) (1.183,4.635)--(1.239,5.648) (1.183,4.635)--(2.089,4.505) (1.239,5.648)--(1.256,6.644) (1.239,5.648)--(2.207,5.531) (1.256,6.644)--(2.279,6.527) (0.721,-6.473)--(0.793,-5.469) (0.721,-6.473)--(1.708,-6.598) (0.793,-5.469)--(0.911,-4.495) (0.793,-5.469)--(1.815,-5.594) (0.911,-4.495)--(1.098,-3.492) (0.911,-4.495)--(1.938,-4.601) (1.098,-3.492)--(1.216,-2.460) (1.098,-3.492)--(2.078,-3.598) (1.216,-2.460)--(1.289,-1.458) (1.216,-2.460)--(2.201,-2.585) (1.289,-1.458)--(1.406,-0.487) (1.289,-1.458)--(2.307,-1.583) (1.406,-0.487)--(1.594,0.513) (1.406,-0.487)--(2.430,-0.594) (1.594,0.513)--(1.711,1.542) (1.594,0.513)--(2.570,0.406) (1.711,1.542)--(1.784,2.540) (1.711,1.542)--(2.693,1.417) (1.784,2.540)--(1.902,3.508) (1.784,2.540)--(2.799,2.415) (1.902,3.508)--(2.089,4.505) (1.902,3.508)--(2.922,3.402) (2.089,4.505)--(2.207,5.531) (2.089,4.505)--(3.062,4.399) (2.207,5.531)--(2.279,6.527) (2.207,5.531)--(3.185,5.406) (2.279,6.527)--(3.292,6.402) (1.708,-6.598)--(1.815,-5.594) (1.708,-6.598)--(2.708,-6.713) (1.815,-5.594)--(1.938,-4.601) (1.815,-5.594)--(2.825,-5.709) (1.938,-4.601)--(2.078,-3.598) (1.938,-4.601)--(2.947,-4.708) (2.078,-3.598)--(2.201,-2.585) (2.078,-3.598)--(3.074,-3.706) (2.201,-2.585)--(2.307,-1.583) (2.201,-2.585)--(3.197,-2.700) (2.307,-1.583)--(2.430,-0.594) (2.307,-1.583)--(3.314,-1.699) (2.430,-0.594)--(2.570,0.406) (2.430,-0.594)--(3.437,-0.701) (2.570,0.406)--(2.693,1.417) (2.570,0.406)--(3.563,0.299) (2.693,1.417)--(2.799,2.415) (2.693,1.417)--(3.686,1.301) (2.799,2.415)--(2.922,3.402) (2.799,2.415)--(3.803,2.300) (2.922,3.402)--(3.062,4.399) (2.922,3.402)--(3.926,3.294) (3.062,4.399)--(3.185,5.406) (3.062,4.399)--(4.053,4.292) (3.185,5.406)--(3.292,6.402) (3.185,5.406)--(4.175,5.291) (3.292,6.402)--(4.292,6.287) (2.708,-6.713)--(2.825,-5.709) (2.708,-6.713)--(3.711,-6.820) (2.825,-5.709)--(2.947,-4.708) (2.825,-5.709)--(3.831,-5.816) (2.947,-4.708)--(3.074,-3.706) (2.947,-4.708)--(3.953,-4.814) (3.074,-3.706)--(3.197,-2.700) (3.074,-3.706)--(4.075,-3.811) (3.197,-2.700)--(3.314,-1.699) (3.197,-2.700)--(4.197,-2.808) (3.314,-1.699)--(3.437,-0.701) (3.314,-1.699)--(4.317,-1.807) (3.437,-0.701)--(3.563,0.299) (3.437,-0.701)--(4.439,-0.807) (3.563,0.299)--(3.686,1.301) (3.563,0.299)--(4.561,0.193) (3.686,1.301)--(3.803,2.300) (3.686,1.301)--(4.683,1.193) (3.803,2.300)--(3.926,3.294) (3.803,2.300)--(4.803,2.192) (3.926,3.294)--(4.053,4.292) (3.926,3.294)--(4.925,3.189) (4.053,4.292)--(4.175,5.291) (4.053,4.292)--(5.047,4.186) (4.175,5.291)--(4.292,6.287) (4.175,5.291)--(5.169,5.184) (4.292,6.287)--(4.415,7.279) (4.292,6.287)--(5.289,6.180) (4.415,7.279)--(5.410,7.174) (3.711,-6.820)--(3.831,-5.816) (3.711,-6.820)--(4.716,-6.922) (3.831,-5.816)--(3.953,-4.814) (3.831,-5.816)--(4.837,-5.919) (3.953,-4.814)--(4.075,-3.811) (3.953,-4.814)--(4.957,-4.916) (4.075,-3.811)--(4.197,-2.808) (4.075,-3.811)--(5.078,-3.913) (4.197,-2.808)--(4.317,-1.807) (4.197,-2.808)--(5.199,-2.911) (4.317,-1.807)--(4.439,-0.807) (4.317,-1.807)--(5.319,-1.910) (4.439,-0.807)--(4.561,0.193) (4.439,-0.807)--(5.440,-0.910) (4.561,0.193)--(4.683,1.193) (4.561,0.193)--(5.560,0.090) (4.683,1.193)--(4.803,2.192) (4.683,1.193)--(5.681,1.090) (4.803,2.192)--(4.925,3.189) (4.803,2.192)--(5.801,2.089) (4.925,3.189)--(5.047,4.186) (4.925,3.189)--(5.922,3.087) (5.047,4.186)--(5.169,5.184) (5.047,4.186)--(6.043,4.084) (5.169,5.184)--(5.289,6.180) (5.169,5.184)--(6.163,5.081) (5.289,6.180)--(5.410,7.174) (5.289,6.180)--(6.284,6.078) (5.410,7.174)--(6.404,7.073) (4.716,-6.922)--(4.837,-5.919) (4.716,-6.922)--(5.722,-7.020) (4.837,-5.919)--(4.957,-4.916) (4.837,-5.919)--(5.842,-6.016) (4.957,-4.916)--(5.078,-3.913) (4.957,-4.916)--(5.961,-5.013) (5.078,-3.913)--(5.199,-2.911) (5.078,-3.913)--(6.081,-4.011) (5.199,-2.911)--(5.319,-1.910) (5.199,-2.911)--(6.201,-3.010) (5.319,-1.910)--(5.440,-0.910) (5.319,-1.910)--(6.320,-2.008) (5.440,-0.910)--(5.560,0.090) (5.440,-0.910)--(6.440,-1.008) (5.560,0.090)--(5.681,1.090) (5.560,0.090)--(6.560,-0.008) (5.681,1.090)--(5.801,2.089) (5.681,1.090)--(6.680,0.992) (5.801,2.089)--(5.922,3.087) (5.801,2.089)--(6.799,1.990) (5.922,3.087)--(6.043,4.084) (5.922,3.087)--(6.919,2.989) (6.043,4.084)--(6.163,5.081) (6.043,4.084)--(7.039,3.987) (6.163,5.081)--(6.284,6.078) (6.163,5.081)--(7.158,4.984) (6.284,6.078)--(6.404,7.073) (6.284,6.078)--(7.278,5.980) (5.722,-7.020)--(5.842,-6.016) (5.722,-7.020)--(6.728,-7.113) (5.842,-6.016)--(5.961,-5.013) (5.842,-6.016)--(6.847,-6.110) (5.961,-5.013)--(6.081,-4.011) (5.961,-5.013)--(6.965,-5.107) (6.081,-4.011)--(6.201,-3.010) (6.081,-4.011)--(7.084,-4.105) (6.201,-3.010)--(6.320,-2.008) (6.201,-3.010)--(7.203,-3.104) (6.320,-2.008)--(6.440,-1.008) (6.440,-1.008)--(6.560,-0.008) (6.560,-0.008)--(6.680,0.992) (6.680,0.992)--(6.799,1.990) (6.799,1.990)--(6.919,2.989) (6.919,2.989)--(7.039,3.987) (7.039,3.987)--(7.158,4.984) (7.158,4.984)--(7.278,5.980) (6.728,-7.113)--(6.847,-6.110) (6.847,-6.110)--(6.965,-5.107) (6.965,-5.107)--(7.084,-4.105) (7.084,-4.105)--(7.203,-3.104);
\fill[atomA, draw=atomLine, line width=0.35pt] (-7.059,-0.102) circle (0.20);
\fill[atomA, draw=atomLine, line width=0.35pt] (-7.178,0.897) circle (0.20);
\fill[atomA, draw=atomLine, line width=0.35pt] (-7.297,1.896) circle (0.20);
\fill[atomA, draw=atomLine, line width=0.35pt] (-6.535,3.893) circle (0.20);
\fill[atomA, draw=atomLine, line width=0.35pt] (-6.653,4.890) circle (0.20);
\fill[atomA, draw=atomLine, line width=0.35pt] (-6.772,5.887) circle (0.20);
\fill[atomA, draw=atomLine, line width=0.35pt] (-6.890,6.884) circle (0.20);
\fill[atomA, draw=atomLine, line width=0.35pt] (-7.222,-7.020) circle (0.20);
\fill[atomA, draw=atomLine, line width=0.35pt] (-6.581,-4.011) circle (0.20);
\fill[atomA, draw=atomLine, line width=0.35pt] (-6.701,-3.010) circle (0.20);
\fill[atomA, draw=atomLine, line width=0.35pt] (-6.820,-2.008) circle (0.20);
\fill[atomA, draw=atomLine, line width=0.35pt] (-6.940,-1.008) circle (0.20);
\fill[atomA, draw=atomLine, line width=0.35pt] (-6.060,-0.008) circle (0.20);
\fill[atomA, draw=atomLine, line width=0.35pt] (-6.180,0.992) circle (0.20);
\fill[atomA, draw=atomLine, line width=0.35pt] (-6.299,1.990) circle (0.20);
\fill[atomA, draw=atomLine, line width=0.35pt] (-6.419,2.989) circle (0.20);
\fill[atomA, draw=atomLine, line width=0.35pt] (-5.539,3.987) circle (0.20);
\fill[atomA, draw=atomLine, line width=0.35pt] (-5.658,4.984) circle (0.20);
\fill[atomA, draw=atomLine, line width=0.35pt] (-5.778,5.980) circle (0.20);
\fill[atomA, draw=atomLine, line width=0.35pt] (-5.897,6.976) circle (0.20);
\fill[atomA, draw=atomLine, line width=0.35pt] (-6.216,-6.922) circle (0.20);
\fill[atomA, draw=atomLine, line width=0.35pt] (-6.337,-5.919) circle (0.20);
\fill[atomA, draw=atomLine, line width=0.35pt] (-6.457,-4.916) circle (0.20);
\fill[atomA, draw=atomLine, line width=0.35pt] (-5.578,-3.913) circle (0.20);
\fill[atomA, draw=atomLine, line width=0.35pt] (-5.699,-2.911) circle (0.20);
\fill[atomA, draw=atomLine, line width=0.35pt] (-5.819,-1.910) circle (0.20);
\fill[atomA, draw=atomLine, line width=0.35pt] (-5.940,-0.910) circle (0.20);
\fill[atomA, draw=atomLine, line width=0.35pt] (-5.060,0.090) circle (0.20);
\fill[atomA, draw=atomLine, line width=0.35pt] (-5.181,1.090) circle (0.20);
\fill[atomA, draw=atomLine, line width=0.35pt] (-5.301,2.089) circle (0.20);
\fill[atomA, draw=atomLine, line width=0.35pt] (-5.422,3.087) circle (0.20);
\fill[atomA, draw=atomLine, line width=0.35pt] (-4.543,4.084) circle (0.20);
\fill[atomA, draw=atomLine, line width=0.35pt] (-4.663,5.081) circle (0.20);
\fill[atomA, draw=atomLine, line width=0.35pt] (-4.784,6.078) circle (0.20);
\fill[atomA, draw=atomLine, line width=0.35pt] (-4.904,7.073) circle (0.20);
\fill[atomA, draw=atomLine, line width=0.35pt] (-5.211,-6.820) circle (0.20);
\fill[atomA, draw=atomLine, line width=0.35pt] (-5.331,-5.816) circle (0.20);
\fill[atomA, draw=atomLine, line width=0.35pt] (-5.453,-4.814) circle (0.20);
\fill[atomA, draw=atomLine, line width=0.35pt] (-4.575,-3.811) circle (0.20);
\fill[atomA, draw=atomLine, line width=0.35pt] (-4.697,-2.808) circle (0.20);
\fill[atomA, draw=atomLine, line width=0.35pt] (-4.817,-1.807) circle (0.20);
\fill[atomA, draw=atomLine, line width=0.35pt] (-4.939,-0.807) circle (0.20);
\fill[atomA, draw=atomLine, line width=0.35pt] (-4.061,0.193) circle (0.20);
\fill[atomA, draw=atomLine, line width=0.35pt] (-4.183,1.193) circle (0.20);
\fill[atomA, draw=atomLine, line width=0.35pt] (-4.303,2.192) circle (0.20);
\fill[atomA, draw=atomLine, line width=0.35pt] (-4.425,3.189) circle (0.20);
\fill[atomA, draw=atomLine, line width=0.35pt] (-3.547,4.186) circle (0.20);
\fill[atomA, draw=atomLine, line width=0.35pt] (-3.669,5.184) circle (0.20);
\fill[atomA, draw=atomLine, line width=0.35pt] (-3.789,6.180) circle (0.20);
\fill[atomA, draw=atomLine, line width=0.35pt] (-3.910,7.174) circle (0.20);
\fill[atomA, draw=atomLine, line width=0.35pt] (-4.208,-6.713) circle (0.20);
\fill[atomA, draw=atomLine, line width=0.35pt] (-4.325,-5.709) circle (0.20);
\fill[atomA, draw=atomLine, line width=0.35pt] (-4.447,-4.708) circle (0.20);
\fill[atomA, draw=atomLine, line width=0.35pt] (-3.574,-3.706) circle (0.20);
\fill[atomA, draw=atomLine, line width=0.35pt] (-3.697,-2.700) circle (0.20);
\fill[atomA, draw=atomLine, line width=0.35pt] (-3.814,-1.699) circle (0.20);
\fill[atomA, draw=atomLine, line width=0.35pt] (-3.937,-0.701) circle (0.20);
\fill[atomA, draw=atomLine, line width=0.35pt] (-3.063,0.299) circle (0.20);
\fill[atomA, draw=atomLine, line width=0.35pt] (-3.186,1.301) circle (0.20);
\fill[atomA, draw=atomLine, line width=0.35pt] (-3.303,2.300) circle (0.20);
\fill[atomA, draw=atomLine, line width=0.35pt] (-3.426,3.294) circle (0.20);
\fill[atomA, draw=atomLine, line width=0.35pt] (-2.553,4.292) circle (0.20);
\fill[atomA, draw=atomLine, line width=0.35pt] (-2.675,5.291) circle (0.20);
\fill[atomA, draw=atomLine, line width=0.35pt] (-2.792,6.287) circle (0.20);
\fill[atomA, draw=atomLine, line width=0.35pt] (-2.915,7.279) circle (0.20);
\fill[atomA, draw=atomLine, line width=0.35pt] (-3.208,-6.598) circle (0.20);
\fill[atomA, draw=atomLine, line width=0.35pt] (-3.315,-5.594) circle (0.20);
\fill[atomA, draw=atomLine, line width=0.35pt] (-3.438,-4.601) circle (0.20);
\fill[atomA, draw=atomLine, line width=0.35pt] (-2.578,-3.598) circle (0.20);
\fill[atomA, draw=atomLine, line width=0.35pt] (-2.701,-2.585) circle (0.20);
\fill[atomA, draw=atomLine, line width=0.35pt] (-2.807,-1.583) circle (0.20);
\fill[atomA, draw=atomLine, line width=0.35pt] (-2.930,-0.594) circle (0.20);
\fill[atomA, draw=atomLine, line width=0.35pt] (-2.070,0.406) circle (0.20);
\fill[atomA, draw=atomLine, line width=0.35pt] (-2.193,1.417) circle (0.20);
\fill[atomA, draw=atomLine, line width=0.35pt] (-2.299,2.415) circle (0.20);
\fill[atomA, draw=atomLine, line width=0.35pt] (-2.422,3.402) circle (0.20);
\fill[atomA, draw=atomLine, line width=0.35pt] (-1.562,4.399) circle (0.20);
\fill[atomA, draw=atomLine, line width=0.35pt] (-1.685,5.406) circle (0.20);
\fill[atomA, draw=atomLine, line width=0.35pt] (-1.792,6.402) circle (0.20);
\fill[atomA, draw=atomLine, line width=0.35pt] (-2.221,-6.473) circle (0.20);
\fill[atomA, draw=atomLine, line width=0.35pt] (-2.293,-5.469) circle (0.20);
\fill[atomA, draw=atomLine, line width=0.35pt] (-2.411,-4.495) circle (0.20);
\fill[atomA, draw=atomLine, line width=0.35pt] (-1.598,-3.492) circle (0.20);
\fill[atomA, draw=atomLine, line width=0.35pt] (-1.716,-2.460) circle (0.20);
\fill[atomA, draw=atomLine, line width=0.35pt] (-1.789,-1.458) circle (0.20);
\fill[atomA, draw=atomLine, line width=0.35pt] (-1.906,-0.487) circle (0.20);
\fill[atomA, draw=atomLine, line width=0.35pt] (-1.094,0.513) circle (0.20);
\fill[atomA, draw=atomLine, line width=0.35pt] (-1.211,1.542) circle (0.20);
\fill[atomA, draw=atomLine, line width=0.35pt] (-1.284,2.540) circle (0.20);
\fill[atomA, draw=atomLine, line width=0.35pt] (-1.402,3.508) circle (0.20);
\fill[atomA, draw=atomLine, line width=0.35pt] (-0.589,4.505) circle (0.20);
\fill[atomA, draw=atomLine, line width=0.35pt] (-0.707,5.531) circle (0.20);
\fill[atomA, draw=atomLine, line width=0.35pt] (-0.779,6.527) circle (0.20);
\fill[atomA, draw=atomLine, line width=0.35pt] (-1.244,-6.356) circle (0.20);
\fill[atomA, draw=atomLine, line width=0.35pt] (-1.261,-5.352) circle (0.20);
\fill[atomA, draw=atomLine, line width=0.35pt] (-1.317,-4.365) circle (0.20);
\fill[atomA, draw=atomLine, line width=0.35pt] (-0.686,-3.362) circle (0.20);
\fill[atomA, draw=atomLine, line width=0.35pt] (-0.742,-2.343) circle (0.20);
\fill[atomA, draw=atomLine, line width=0.35pt] (-0.759,-1.341) circle (0.20);
\fill[atomA, draw=atomLine, line width=0.35pt] (-0.816,-0.357) circle (0.20);
\fill[atomA, draw=atomLine, line width=0.35pt] (-0.184,0.643) circle (0.20);
\fill[atomA, draw=atomLine, line width=0.35pt] (-0.241,1.659) circle (0.20);
\fill[atomA, draw=atomLine, line width=0.35pt] (-0.258,2.657) circle (0.20);
\fill[atomA, draw=atomLine, line width=0.35pt] (-0.314,3.638) circle (0.20);
\fill[atomA, draw=atomLine, line width=0.35pt] (0.317,4.635) circle (0.20);
\fill[atomA, draw=atomLine, line width=0.35pt] (0.261,5.648) circle (0.20);
\fill[atomA, draw=atomLine, line width=0.35pt] (0.244,6.644) circle (0.20);
\fill[atomA, draw=atomLine, line width=0.35pt] (-0.256,-6.356) circle (0.20);
\fill[atomA, draw=atomLine, line width=0.35pt] (-0.239,-5.352) circle (0.20);
\fill[atomA, draw=atomLine, line width=0.35pt] (-0.183,-4.365) circle (0.20);
\fill[atomA, draw=atomLine, line width=0.35pt] (0.186,-3.362) circle (0.20);
\fill[atomA, draw=atomLine, line width=0.35pt] (0.242,-2.343) circle (0.20);
\fill[atomA, draw=atomLine, line width=0.35pt] (0.259,-1.341) circle (0.20);
\fill[atomA, draw=atomLine, line width=0.35pt] (0.316,-0.357) circle (0.20);
\fill[atomA, draw=atomLine, line width=0.35pt] (0.684,0.643) circle (0.20);
\fill[atomA, draw=atomLine, line width=0.35pt] (0.741,1.659) circle (0.20);
\fill[atomA, draw=atomLine, line width=0.35pt] (0.758,2.657) circle (0.20);
\fill[atomA, draw=atomLine, line width=0.35pt] (0.814,3.638) circle (0.20);
\fill[atomA, draw=atomLine, line width=0.35pt] (1.183,4.635) circle (0.20);
\fill[atomA, draw=atomLine, line width=0.35pt] (1.239,5.648) circle (0.20);
\fill[atomA, draw=atomLine, line width=0.35pt] (1.256,6.644) circle (0.20);
\fill[atomA, draw=atomLine, line width=0.35pt] (0.721,-6.473) circle (0.20);
\fill[atomA, draw=atomLine, line width=0.35pt] (0.793,-5.469) circle (0.20);
\fill[atomA, draw=atomLine, line width=0.35pt] (0.911,-4.495) circle (0.20);
\fill[atomA, draw=atomLine, line width=0.35pt] (1.098,-3.492) circle (0.20);
\fill[atomA, draw=atomLine, line width=0.35pt] (1.216,-2.460) circle (0.20);
\fill[atomA, draw=atomLine, line width=0.35pt] (1.289,-1.458) circle (0.20);
\fill[atomA, draw=atomLine, line width=0.35pt] (1.406,-0.487) circle (0.20);
\fill[atomA, draw=atomLine, line width=0.35pt] (1.594,0.513) circle (0.20);
\fill[atomA, draw=atomLine, line width=0.35pt] (1.711,1.542) circle (0.20);
\fill[atomA, draw=atomLine, line width=0.35pt] (1.784,2.540) circle (0.20);
\fill[atomA, draw=atomLine, line width=0.35pt] (1.902,3.508) circle (0.20);
\fill[atomA, draw=atomLine, line width=0.35pt] (2.089,4.505) circle (0.20);
\fill[atomA, draw=atomLine, line width=0.35pt] (2.207,5.531) circle (0.20);
\fill[atomA, draw=atomLine, line width=0.35pt] (2.279,6.527) circle (0.20);
\fill[atomA, draw=atomLine, line width=0.35pt] (1.708,-6.598) circle (0.20);
\fill[atomA, draw=atomLine, line width=0.35pt] (1.815,-5.594) circle (0.20);
\fill[atomA, draw=atomLine, line width=0.35pt] (1.938,-4.601) circle (0.20);
\fill[atomA, draw=atomLine, line width=0.35pt] (2.078,-3.598) circle (0.20);
\fill[atomA, draw=atomLine, line width=0.35pt] (2.201,-2.585) circle (0.20);
\fill[atomA, draw=atomLine, line width=0.35pt] (2.307,-1.583) circle (0.20);
\fill[atomA, draw=atomLine, line width=0.35pt] (2.430,-0.594) circle (0.20);
\fill[atomA, draw=atomLine, line width=0.35pt] (2.570,0.406) circle (0.20);
\fill[atomA, draw=atomLine, line width=0.35pt] (2.693,1.417) circle (0.20);
\fill[atomA, draw=atomLine, line width=0.35pt] (2.799,2.415) circle (0.20);
\fill[atomA, draw=atomLine, line width=0.35pt] (2.922,3.402) circle (0.20);
\fill[atomA, draw=atomLine, line width=0.35pt] (3.062,4.399) circle (0.20);
\fill[atomA, draw=atomLine, line width=0.35pt] (3.185,5.406) circle (0.20);
\fill[atomA, draw=atomLine, line width=0.35pt] (3.292,6.402) circle (0.20);
\fill[atomA, draw=atomLine, line width=0.35pt] (2.708,-6.713) circle (0.20);
\fill[atomA, draw=atomLine, line width=0.35pt] (2.825,-5.709) circle (0.20);
\fill[atomA, draw=atomLine, line width=0.35pt] (2.947,-4.708) circle (0.20);
\fill[atomA, draw=atomLine, line width=0.35pt] (3.074,-3.706) circle (0.20);
\fill[atomA, draw=atomLine, line width=0.35pt] (3.197,-2.700) circle (0.20);
\fill[atomA, draw=atomLine, line width=0.35pt] (3.314,-1.699) circle (0.20);
\fill[atomA, draw=atomLine, line width=0.35pt] (3.437,-0.701) circle (0.20);
\fill[atomA, draw=atomLine, line width=0.35pt] (3.563,0.299) circle (0.20);
\fill[atomA, draw=atomLine, line width=0.35pt] (3.686,1.301) circle (0.20);
\fill[atomA, draw=atomLine, line width=0.35pt] (3.803,2.300) circle (0.20);
\fill[atomA, draw=atomLine, line width=0.35pt] (3.926,3.294) circle (0.20);
\fill[atomA, draw=atomLine, line width=0.35pt] (4.053,4.292) circle (0.20);
\fill[atomA, draw=atomLine, line width=0.35pt] (4.175,5.291) circle (0.20);
\fill[atomA, draw=atomLine, line width=0.35pt] (4.292,6.287) circle (0.20);
\fill[atomA, draw=atomLine, line width=0.35pt] (4.415,7.279) circle (0.20);
\fill[atomA, draw=atomLine, line width=0.35pt] (3.711,-6.820) circle (0.20);
\fill[atomA, draw=atomLine, line width=0.35pt] (3.831,-5.816) circle (0.20);
\fill[atomA, draw=atomLine, line width=0.35pt] (3.953,-4.814) circle (0.20);
\fill[atomA, draw=atomLine, line width=0.35pt] (4.075,-3.811) circle (0.20);
\fill[atomA, draw=atomLine, line width=0.35pt] (4.197,-2.808) circle (0.20);
\fill[atomA, draw=atomLine, line width=0.35pt] (4.317,-1.807) circle (0.20);
\fill[atomA, draw=atomLine, line width=0.35pt] (4.439,-0.807) circle (0.20);
\fill[atomA, draw=atomLine, line width=0.35pt] (4.561,0.193) circle (0.20);
\fill[atomA, draw=atomLine, line width=0.35pt] (4.683,1.193) circle (0.20);
\fill[atomA, draw=atomLine, line width=0.35pt] (4.803,2.192) circle (0.20);
\fill[atomA, draw=atomLine, line width=0.35pt] (4.925,3.189) circle (0.20);
\fill[atomA, draw=atomLine, line width=0.35pt] (5.047,4.186) circle (0.20);
\fill[atomA, draw=atomLine, line width=0.35pt] (5.169,5.184) circle (0.20);
\fill[atomA, draw=atomLine, line width=0.35pt] (5.289,6.180) circle (0.20);
\fill[atomA, draw=atomLine, line width=0.35pt] (5.410,7.174) circle (0.20);
\fill[atomA, draw=atomLine, line width=0.35pt] (4.716,-6.922) circle (0.20);
\fill[atomA, draw=atomLine, line width=0.35pt] (4.837,-5.919) circle (0.20);
\fill[atomA, draw=atomLine, line width=0.35pt] (4.957,-4.916) circle (0.20);
\fill[atomA, draw=atomLine, line width=0.35pt] (5.078,-3.913) circle (0.20);
\fill[atomA, draw=atomLine, line width=0.35pt] (5.199,-2.911) circle (0.20);
\fill[atomA, draw=atomLine, line width=0.35pt] (5.319,-1.910) circle (0.20);
\fill[atomA, draw=atomLine, line width=0.35pt] (5.440,-0.910) circle (0.20);
\fill[atomA, draw=atomLine, line width=0.35pt] (5.560,0.090) circle (0.20);
\fill[atomA, draw=atomLine, line width=0.35pt] (5.681,1.090) circle (0.20);
\fill[atomA, draw=atomLine, line width=0.35pt] (5.801,2.089) circle (0.20);
\fill[atomA, draw=atomLine, line width=0.35pt] (5.922,3.087) circle (0.20);
\fill[atomA, draw=atomLine, line width=0.35pt] (6.043,4.084) circle (0.20);
\fill[atomA, draw=atomLine, line width=0.35pt] (6.163,5.081) circle (0.20);
\fill[atomA, draw=atomLine, line width=0.35pt] (6.284,6.078) circle (0.20);
\fill[atomA, draw=atomLine, line width=0.35pt] (6.404,7.073) circle (0.20);
\fill[atomA, draw=atomLine, line width=0.35pt] (5.722,-7.020) circle (0.20);
\fill[atomA, draw=atomLine, line width=0.35pt] (5.842,-6.016) circle (0.20);
\fill[atomA, draw=atomLine, line width=0.35pt] (5.961,-5.013) circle (0.20);
\fill[atomA, draw=atomLine, line width=0.35pt] (6.081,-4.011) circle (0.20);
\fill[atomA, draw=atomLine, line width=0.35pt] (6.201,-3.010) circle (0.20);
\fill[atomA, draw=atomLine, line width=0.35pt] (6.320,-2.008) circle (0.20);
\fill[atomA, draw=atomLine, line width=0.35pt] (6.440,-1.008) circle (0.20);
\fill[atomA, draw=atomLine, line width=0.35pt] (6.560,-0.008) circle (0.20);
\fill[atomA, draw=atomLine, line width=0.35pt] (6.680,0.992) circle (0.20);
\fill[atomA, draw=atomLine, line width=0.35pt] (6.799,1.990) circle (0.20);
\fill[atomA, draw=atomLine, line width=0.35pt] (6.919,2.989) circle (0.20);
\fill[atomA, draw=atomLine, line width=0.35pt] (7.039,3.987) circle (0.20);
\fill[atomA, draw=atomLine, line width=0.35pt] (7.158,4.984) circle (0.20);
\fill[atomA, draw=atomLine, line width=0.35pt] (7.278,5.980) circle (0.20);
\fill[atomA, draw=atomLine, line width=0.35pt] (6.728,-7.113) circle (0.20);
\fill[atomA, draw=atomLine, line width=0.35pt] (6.847,-6.110) circle (0.20);
\fill[atomA, draw=atomLine, line width=0.35pt] (6.965,-5.107) circle (0.20);
\fill[atomA, draw=atomLine, line width=0.35pt] (7.084,-4.105) circle (0.20);
\fill[atomA, draw=atomLine, line width=0.35pt] (7.203,-3.104) circle (0.20);
\end{scope}
\dislT[fmRed]{(0.000,-3.055)}{0}
\dislT[fmRed]{(0.000,0.945)}{0}
\dislT[fmRed]{(0.000,4.945)}{0}
\coordinate (TB0) at (0,0.995);
\coordinate (TB1) at (0,4.995);

  \draw[dimline, fmRed] ($(TB0)+(1.3,0)$) -- node[right, slbl, fmRed]{$D$} ($(TB1)+(1.3,0)$);
\end{tikzpicture}
\caption{Симметричная граница наклона в простой кубической решётке~—
стенка краевых дислокаций через $D=4b$; угол разориентации
$\theta=b/D\approx14^\circ$ (положения атомов рассчитаны наложением упругих
полей дислокаций)}
\label{fig:lab-wall}
\end{figure}

\begin{primer}[Расстояние между дислокациями в субгранице]
В алюминии ($a=0{,}405$~нм, $b=a/\sqrt2=0{,}286$~нм) субграница наклона
с $\theta=2^\circ=0{,}035$~рад состоит из дислокаций, отстоящих на
$D=b/\theta=0{,}286/0{,}035\approx8{,}2$~нм, т.\,е. примерно через $29b$.
\end{primer}

В \emph{несимметричной} границе наклона (плоскость границы не делит угол
между решётками пополам) присутствуют дислокации с двумя разными
векторами Бюргерса. Малоугловая \term{граница кручения}~— сетка из двух
семейств параллельных \emph{винтовых} дислокаций; формула~\eqref{eq:wall-spacing}
справедлива и для неё, если $D$~— расстояние между дислокациями одного
семейства.

В общем случае дислокационное строение малоугловой границы задаёт
\term{формула Франка}:
\begin{equation}\label{eq:frank}
\vb d=\sum_i\vb b_i=2\sin\frac\theta2\,(\vb r\times\vb u),
\end{equation}
где $\vb r$~— произвольный вектор в плоскости границы, $\vb u$~— единичный
вектор оси разворота, $\vb d$~— суммарный вектор Бюргерса дислокаций,
которые пересекает вектор $\vb r$. Для симметричной границы наклона
($\vb r\perp\vb u$, оба в плоскости границы) $|\vb d|=2|\vb r|\sin(\theta/2)\approx|\vb r|\theta$:
вектор длины $|\vb r|$ пересекает $|\vb r|/D$ дислокаций, откуда снова
$D=b/\theta$.

\paragraph{Энергия малоугловой границы}
Энергия единицы площади границы растёт с углом разориентации
(\term{формула Рида--Шокли}):
\begin{keyeq}\label{eq:read-shockley}
E_\text{гр}=E_0\,\theta\,(A-\ln\theta),
\end{keyeq}
где $E_0\propto Gb$ ($G$~— модуль сдвига) и $A$~— константа, зависящая от
энергии ядра дислокации. Множитель $\theta$~— это плотность дислокаций в
стенке ($1/D=\theta/b$), а $-\ln\theta$ отражает взаимное экранирование
полей соседних дислокаций: чем они ближе, тем меньше энергия каждой.
Формула~\eqref{eq:read-shockley} имеет максимум при $\theta_m=e^{A-1}$, и её
удобно записать в приведённом виде
$E/E_m=(\theta/\theta_m)\,[1-\ln(\theta/\theta_m)]$ (рис.~\ref{fig:gb-energy},\,\textit{а}).
Практически она применима до $\theta\approx10$--$15^\circ$; выше энергия
границы почти не зависит от угла~— кроме особых разориентаций.

\subsection{Высокоугловые границы}\label{ssec:high-angle}

\paragraph{От аморфной прослойки к решёткам совпадения}
Первоначально границу зерна представляли \emph{аморфной прослойкой},
обеспечивающей сцепление зёрен. Эта модель противоречит опыту:
\begin{enumerate}
\item свойства аморфной границы не зависели бы от взаимной ориентации
зёрен, а на самом деле энергия границы, скорость диффузии по ней,
скорость миграции и зернограничное проскальзывание от разориентации
зависят; значит, граница имеет определённое строение;
\item чтобы зёрна неправильной формы могли проскальзывать, прослойка
должна быть толщиной порядка сотни атомных диаметров, а прямые
наблюдения (ионный проектор, электронная микроскопия высокого разрешения)
дают толщину всего в несколько межатомных расстояний.
\end{enumerate}
Затем предлагались модель \emph{переходной решётки} (атомы занимают
промежуточные положения между узлами решёток соседних зёрен) и
\emph{островковая модель Мотта}: граница состоит из участков «хорошего» и
«плохого» сопряжения решёток. В исходном виде модель Мотта не используют,
но сама идея чередования участков хорошего и плохого сопряжения лежит в
основе большинства современных моделей.

\paragraph{Решётка совпадающих узлов}
Мысленно продолжим решётки обоих зёрен на всё пространство. При
определённых, строго фиксированных осях и углах разворота часть узлов
одной решётки совпадает с узлами другой, и совпадающие узлы сами образуют
решётку (рис.~\ref{fig:csl5}).

\begin{opredelenie}
\term{Решётка совпадающих узлов} (РСУ, англ. CSL~— coincidence site lattice)~—
общая часть (пересечение) решёток двух зёрен, $L_\text{РСУ}=L_1\cap L_2$.
Величина $\Sigma$~— \term{обратная плотность совпадающих узлов}: число узлов
решётки, приходящееся на один совпадающий узел (отношение объёмов
элементарных ячеек РСУ и решётки).
\end{opredelenie}

Так, при повороте вокруг $[001]$ на угол
$\theta=2\operatorname{arctg}(1/3)=36{,}87^\circ$\ispr{В оригинале
$\theta=2\operatorname{arctg}(1/2)=36{,}9^\circ$. Но $2\operatorname{arctg}(1/2)=53{,}13^\circ$;
этот угол даёт ту же разориентацию $\Sigma5$ (так как $53{,}13^\circ=90^\circ-36{,}87^\circ$,
а поворот на $90^\circ$ вокруг $[001]$~— операция симметрии куба), но
значению $36{,}9^\circ$ соответствует $2\operatorname{arctg}(1/3)$.}
совпадает каждый пятый узел: $\Sigma=5$, а ячейка РСУ~— квадрат со
стороной $a\sqrt5$ (в трёх измерениях~— тетрагональная ячейка).
$\Sigma=1$ означает полное совпадение решёток (границы нет); $\Sigma=3$~—
двойник в ГЦК-решётке (поворот на $60^\circ$ вокруг $\nfam{111}$). Для
кубических решёток $\Sigma$ всегда \emph{нечётно}\ispr{В оригинале:
«число $\Sigma$ всегда простое». Это неверно: бывают $\Sigma=9$, 15, 21,
25, 27\ldots\ (табл.~\ref{tab:csl}). Верно другое: в кубических решётках
$\Sigma$ всегда нечётно.}.

\begin{figure}[ht]
\centering
\begin{tikzpicture}[scale=0.85]
  % создано tools/gen_boundaries.py: решётка совпадающих узлов Σ5 [001]
\begin{scope}
\clip (-3.25,-3.25) rectangle (3.25,3.25);
\draw[fmGreen!70!black, line width=0.5pt] (3,-9) -- ++(1,2) -- ++(-2,1);
\draw[fmGreen!70!black, line width=0.5pt] (1,-8) -- ++(1,2) -- ++(-2,1);
\draw[fmGreen!70!black, line width=0.5pt] (-1,-7) -- ++(1,2) -- ++(-2,1);
\draw[fmGreen!70!black, line width=0.5pt] (-3,-6) -- ++(1,2) -- ++(-2,1);
\draw[fmGreen!70!black, line width=0.5pt] (-5,-5) -- ++(1,2) -- ++(-2,1);
\draw[fmGreen!70!black, line width=0.5pt] (-7,-4) -- ++(1,2) -- ++(-2,1);
\draw[fmGreen!70!black, line width=0.5pt] (-9,-3) -- ++(1,2) -- ++(-2,1);
\draw[fmGreen!70!black, line width=0.5pt] (4,-7) -- ++(1,2) -- ++(-2,1);
\draw[fmGreen!70!black, line width=0.5pt] (2,-6) -- ++(1,2) -- ++(-2,1);
\draw[fmGreen!70!black, line width=0.5pt] (0,-5) -- ++(1,2) -- ++(-2,1);
\draw[fmGreen!70!black, line width=0.5pt] (-2,-4) -- ++(1,2) -- ++(-2,1);
\draw[fmGreen!70!black, line width=0.5pt] (-4,-3) -- ++(1,2) -- ++(-2,1);
\draw[fmGreen!70!black, line width=0.5pt] (-6,-2) -- ++(1,2) -- ++(-2,1);
\draw[fmGreen!70!black, line width=0.5pt] (-8,-1) -- ++(1,2) -- ++(-2,1);
\draw[fmGreen!70!black, line width=0.5pt] (5,-5) -- ++(1,2) -- ++(-2,1);
\draw[fmGreen!70!black, line width=0.5pt] (3,-4) -- ++(1,2) -- ++(-2,1);
\draw[fmGreen!70!black, line width=0.5pt] (1,-3) -- ++(1,2) -- ++(-2,1);
\draw[fmGreen!70!black, line width=0.5pt] (-1,-2) -- ++(1,2) -- ++(-2,1);
\draw[fmGreen!70!black, line width=0.5pt] (-3,-1) -- ++(1,2) -- ++(-2,1);
\draw[fmGreen!70!black, line width=0.5pt] (-5,0) -- ++(1,2) -- ++(-2,1);
\draw[fmGreen!70!black, line width=0.5pt] (-7,1) -- ++(1,2) -- ++(-2,1);
\draw[fmGreen!70!black, line width=0.5pt] (6,-3) -- ++(1,2) -- ++(-2,1);
\draw[fmGreen!70!black, line width=0.5pt] (4,-2) -- ++(1,2) -- ++(-2,1);
\draw[fmGreen!70!black, line width=0.5pt] (2,-1) -- ++(1,2) -- ++(-2,1);
\draw[fmGreen!70!black, line width=0.5pt] (0,0) -- ++(1,2) -- ++(-2,1);
\draw[fmGreen!70!black, line width=0.5pt] (-2,1) -- ++(1,2) -- ++(-2,1);
\draw[fmGreen!70!black, line width=0.5pt] (-4,2) -- ++(1,2) -- ++(-2,1);
\draw[fmGreen!70!black, line width=0.5pt] (-6,3) -- ++(1,2) -- ++(-2,1);
\draw[fmGreen!70!black, line width=0.5pt] (7,-1) -- ++(1,2) -- ++(-2,1);
\draw[fmGreen!70!black, line width=0.5pt] (5,0) -- ++(1,2) -- ++(-2,1);
\draw[fmGreen!70!black, line width=0.5pt] (3,1) -- ++(1,2) -- ++(-2,1);
\draw[fmGreen!70!black, line width=0.5pt] (1,2) -- ++(1,2) -- ++(-2,1);
\draw[fmGreen!70!black, line width=0.5pt] (-1,3) -- ++(1,2) -- ++(-2,1);
\draw[fmGreen!70!black, line width=0.5pt] (-3,4) -- ++(1,2) -- ++(-2,1);
\draw[fmGreen!70!black, line width=0.5pt] (-5,5) -- ++(1,2) -- ++(-2,1);
\draw[fmGreen!70!black, line width=0.5pt] (8,1) -- ++(1,2) -- ++(-2,1);
\draw[fmGreen!70!black, line width=0.5pt] (6,2) -- ++(1,2) -- ++(-2,1);
\draw[fmGreen!70!black, line width=0.5pt] (4,3) -- ++(1,2) -- ++(-2,1);
\draw[fmGreen!70!black, line width=0.5pt] (2,4) -- ++(1,2) -- ++(-2,1);
\draw[fmGreen!70!black, line width=0.5pt] (0,5) -- ++(1,2) -- ++(-2,1);
\draw[fmGreen!70!black, line width=0.5pt] (-2,6) -- ++(1,2) -- ++(-2,1);
\draw[fmGreen!70!black, line width=0.5pt] (-4,7) -- ++(1,2) -- ++(-2,1);
\draw[fmGreen!70!black, line width=0.5pt] (9,3) -- ++(1,2) -- ++(-2,1);
\draw[fmGreen!70!black, line width=0.5pt] (7,4) -- ++(1,2) -- ++(-2,1);
\draw[fmGreen!70!black, line width=0.5pt] (5,5) -- ++(1,2) -- ++(-2,1);
\draw[fmGreen!70!black, line width=0.5pt] (3,6) -- ++(1,2) -- ++(-2,1);
\draw[fmGreen!70!black, line width=0.5pt] (1,7) -- ++(1,2) -- ++(-2,1);
\draw[fmGreen!70!black, line width=0.5pt] (-1,8) -- ++(1,2) -- ++(-2,1);
\draw[fmGreen!70!black, line width=0.5pt] (-3,9) -- ++(1,2) -- ++(-2,1);
\fill[fmGreen!15] (0,0) -- (1,2) -- (-1,3) -- (-2,1) -- cycle;
\draw[fmGreen!70!black, line width=1.1pt] (0,0) -- (1,2) -- (-1,3) -- (-2,1) -- cycle;
\fill[fmGreen!60!black] (0.000,0.000) circle (0.045);
\fill[fmGreen!60!black] (-0.400,0.200) circle (0.045);
\fill[fmGreen!60!black] (-0.800,0.400) circle (0.045);
\fill[fmGreen!60!black] (-1.200,0.600) circle (0.045);
\fill[fmGreen!60!black] (-1.600,0.800) circle (0.045);
\fill[fmGreen!60!black] (-2.000,1.000) circle (0.045);
\fill[fmGreen!60!black] (0.200,0.400) circle (0.045);
\fill[fmGreen!60!black] (-0.200,0.600) circle (0.045);
\fill[fmGreen!60!black] (-0.600,0.800) circle (0.045);
\fill[fmGreen!60!black] (-1.000,1.000) circle (0.045);
\fill[fmGreen!60!black] (-1.400,1.200) circle (0.045);
\fill[fmGreen!60!black] (-1.800,1.400) circle (0.045);
\fill[fmGreen!60!black] (0.400,0.800) circle (0.045);
\fill[fmGreen!60!black] (0.000,1.000) circle (0.045);
\fill[fmGreen!60!black] (-0.400,1.200) circle (0.045);
\fill[fmGreen!60!black] (-0.800,1.400) circle (0.045);
\fill[fmGreen!60!black] (-1.200,1.600) circle (0.045);
\fill[fmGreen!60!black] (-1.600,1.800) circle (0.045);
\fill[fmGreen!60!black] (0.600,1.200) circle (0.045);
\fill[fmGreen!60!black] (0.200,1.400) circle (0.045);
\fill[fmGreen!60!black] (-0.200,1.600) circle (0.045);
\fill[fmGreen!60!black] (-0.600,1.800) circle (0.045);
\fill[fmGreen!60!black] (-1.000,2.000) circle (0.045);
\fill[fmGreen!60!black] (-1.400,2.200) circle (0.045);
\fill[fmGreen!60!black] (0.800,1.600) circle (0.045);
\fill[fmGreen!60!black] (0.400,1.800) circle (0.045);
\fill[fmGreen!60!black] (0.000,2.000) circle (0.045);
\fill[fmGreen!60!black] (-0.400,2.200) circle (0.045);
\fill[fmGreen!60!black] (-0.800,2.400) circle (0.045);
\fill[fmGreen!60!black] (-1.200,2.600) circle (0.045);
\fill[fmGreen!60!black] (1.000,2.000) circle (0.045);
\fill[fmGreen!60!black] (0.600,2.200) circle (0.045);
\fill[fmGreen!60!black] (0.200,2.400) circle (0.045);
\fill[fmGreen!60!black] (-0.200,2.600) circle (0.045);
\fill[fmGreen!60!black] (-0.600,2.800) circle (0.045);
\fill[fmGreen!60!black] (-1.000,3.000) circle (0.045);
\draw[fmRed, line width=0.7pt] (-2.6000,-3.2000) circle (0.15);
\draw[fmRed, line width=0.7pt] (-3.2000,-2.4000) circle (0.15);
\fill[fmBlue] (-3,-3) circle (0.09);
\fill[fmBlue] (-3,-2) circle (0.09);
\draw[fmRed, line width=0.7pt] (-1.2000,-3.4000) circle (0.15);
\fill[fmBlue] (-3,-1) circle (0.09);
\draw[fmRed, line width=0.7pt] (-1.8000,-2.6000) circle (0.15);
\fill[fmBlue] (-3,0) circle (0.09);
\draw[fmRed, line width=0.7pt] (-2.4000,-1.8000) circle (0.15);
\fill[fmBlue] (-3,1) circle (0.09);
\draw[fmRed, line width=0.7pt] (-3.0000,-1.0000) circle (0.15);
\draw[black, line width=1.1pt] (-3.0000,-1.0000) circle (0.25);
\fill[fmBlue] (-3,2) circle (0.09);
\fill[fmBlue] (-3,3) circle (0.09);
\fill[fmBlue] (-2,-3) circle (0.09);
\fill[fmBlue] (-2,-2) circle (0.09);
\draw[fmRed, line width=0.7pt] (-0.4000,-2.8000) circle (0.15);
\fill[fmBlue] (-2,-1) circle (0.09);
\draw[fmRed, line width=0.7pt] (-1.0000,-2.0000) circle (0.15);
\draw[black, line width=1.1pt] (-1.0000,-2.0000) circle (0.25);
\fill[fmBlue] (-2,0) circle (0.09);
\draw[fmRed, line width=0.7pt] (-1.6000,-1.2000) circle (0.15);
\fill[fmBlue] (-2,1) circle (0.09);
\draw[fmRed, line width=0.7pt] (-2.2000,-0.4000) circle (0.15);
\fill[fmBlue] (-2,2) circle (0.09);
\draw[fmRed, line width=0.7pt] (-2.8000,0.4000) circle (0.15);
\fill[fmBlue] (-2,3) circle (0.09);
\draw[fmRed, line width=0.7pt] (-3.4000,1.2000) circle (0.15);
\fill[fmBlue] (-1,-3) circle (0.09);
\draw[fmRed, line width=0.7pt] (1.0000,-3.0000) circle (0.15);
\draw[black, line width=1.1pt] (1.0000,-3.0000) circle (0.25);
\fill[fmBlue] (-1,-2) circle (0.09);
\draw[fmRed, line width=0.7pt] (0.4000,-2.2000) circle (0.15);
\fill[fmBlue] (-1,-1) circle (0.09);
\draw[fmRed, line width=0.7pt] (-0.2000,-1.4000) circle (0.15);
\fill[fmBlue] (-1,0) circle (0.09);
\draw[fmRed, line width=0.7pt] (-0.8000,-0.6000) circle (0.15);
\fill[fmBlue] (-1,1) circle (0.09);
\draw[fmRed, line width=0.7pt] (-1.4000,0.2000) circle (0.15);
\fill[fmBlue] (-1,2) circle (0.09);
\draw[fmRed, line width=0.7pt] (-2.0000,1.0000) circle (0.15);
\draw[black, line width=1.1pt] (-2.0000,1.0000) circle (0.25);
\fill[fmBlue] (-1,3) circle (0.09);
\draw[fmRed, line width=0.7pt] (-2.6000,1.8000) circle (0.15);
\draw[fmRed, line width=0.7pt] (-3.2000,2.6000) circle (0.15);
\draw[fmRed, line width=0.7pt] (2.4000,-3.2000) circle (0.15);
\fill[fmBlue] (0,-3) circle (0.09);
\draw[fmRed, line width=0.7pt] (1.8000,-2.4000) circle (0.15);
\fill[fmBlue] (0,-2) circle (0.09);
\draw[fmRed, line width=0.7pt] (1.2000,-1.6000) circle (0.15);
\fill[fmBlue] (0,-1) circle (0.09);
\draw[fmRed, line width=0.7pt] (0.6000,-0.8000) circle (0.15);
\fill[fmBlue] (0,0) circle (0.09);
\draw[fmRed, line width=0.7pt] (0.0000,0.0000) circle (0.15);
\draw[black, line width=1.1pt] (0.0000,0.0000) circle (0.25);
\fill[fmBlue] (0,1) circle (0.09);
\draw[fmRed, line width=0.7pt] (-0.6000,0.8000) circle (0.15);
\fill[fmBlue] (0,2) circle (0.09);
\draw[fmRed, line width=0.7pt] (-1.2000,1.6000) circle (0.15);
\fill[fmBlue] (0,3) circle (0.09);
\draw[fmRed, line width=0.7pt] (-1.8000,2.4000) circle (0.15);
\draw[fmRed, line width=0.7pt] (-2.4000,3.2000) circle (0.15);
\draw[fmRed, line width=0.7pt] (3.2000,-2.6000) circle (0.15);
\fill[fmBlue] (1,-3) circle (0.09);
\draw[fmRed, line width=0.7pt] (2.6000,-1.8000) circle (0.15);
\fill[fmBlue] (1,-2) circle (0.09);
\draw[fmRed, line width=0.7pt] (2.0000,-1.0000) circle (0.15);
\draw[black, line width=1.1pt] (2.0000,-1.0000) circle (0.25);
\fill[fmBlue] (1,-1) circle (0.09);
\draw[fmRed, line width=0.7pt] (1.4000,-0.2000) circle (0.15);
\fill[fmBlue] (1,0) circle (0.09);
\draw[fmRed, line width=0.7pt] (0.8000,0.6000) circle (0.15);
\fill[fmBlue] (1,1) circle (0.09);
\draw[fmRed, line width=0.7pt] (0.2000,1.4000) circle (0.15);
\fill[fmBlue] (1,2) circle (0.09);
\draw[fmRed, line width=0.7pt] (-0.4000,2.2000) circle (0.15);
\fill[fmBlue] (1,3) circle (0.09);
\draw[fmRed, line width=0.7pt] (-1.0000,3.0000) circle (0.15);
\draw[black, line width=1.1pt] (-1.0000,3.0000) circle (0.25);
\fill[fmBlue] (2,-3) circle (0.09);
\draw[fmRed, line width=0.7pt] (3.4000,-1.2000) circle (0.15);
\fill[fmBlue] (2,-2) circle (0.09);
\draw[fmRed, line width=0.7pt] (2.8000,-0.4000) circle (0.15);
\fill[fmBlue] (2,-1) circle (0.09);
\draw[fmRed, line width=0.7pt] (2.2000,0.4000) circle (0.15);
\fill[fmBlue] (2,0) circle (0.09);
\draw[fmRed, line width=0.7pt] (1.6000,1.2000) circle (0.15);
\fill[fmBlue] (2,1) circle (0.09);
\draw[fmRed, line width=0.7pt] (1.0000,2.0000) circle (0.15);
\draw[black, line width=1.1pt] (1.0000,2.0000) circle (0.25);
\fill[fmBlue] (2,2) circle (0.09);
\draw[fmRed, line width=0.7pt] (0.4000,2.8000) circle (0.15);
\fill[fmBlue] (2,3) circle (0.09);
\fill[fmBlue] (3,-3) circle (0.09);
\fill[fmBlue] (3,-2) circle (0.09);
\fill[fmBlue] (3,-1) circle (0.09);
\draw[fmRed, line width=0.7pt] (3.0000,1.0000) circle (0.15);
\draw[black, line width=1.1pt] (3.0000,1.0000) circle (0.25);
\fill[fmBlue] (3,0) circle (0.09);
\draw[fmRed, line width=0.7pt] (2.4000,1.8000) circle (0.15);
\fill[fmBlue] (3,1) circle (0.09);
\draw[fmRed, line width=0.7pt] (1.8000,2.6000) circle (0.15);
\fill[fmBlue] (3,2) circle (0.09);
\draw[fmRed, line width=0.7pt] (1.2000,3.4000) circle (0.15);
\fill[fmBlue] (3,3) circle (0.09);
\draw[fmRed, line width=0.7pt] (3.2000,2.4000) circle (0.15);
\draw[fmRed, line width=0.7pt] (2.6000,3.2000) circle (0.15);
\end{scope}
\draw[atomLine] (-3.25,-3.25) rectangle (3.25,3.25);
\begin{scope}[shift={(7.6,0)}]
\clip (-3.25,-3.25) rectangle (3.25,3.25);
\fill[fmBlue!7] (-3.25,-3.25) rectangle (0,3.25);
\fill[fmRed!7] (0,-3.25) rectangle (3.25,3.25);
\node[atom=2.3mm] at (-3.4785,2.2136) {};
\node[atom=2.3mm] at (-3.1623,3.1623) {};
\node[atom=2.3mm] at (-3.4785,-0.9487) {};
\node[atom=2.3mm] at (-3.1623,0.0000) {};
\node[atom=2.3mm] at (-2.8460,0.9487) {};
\node[atom=2.3mm] at (-2.5298,1.8974) {};
\node[atom=2.3mm] at (-2.2136,2.8460) {};
\node[atom=2.3mm] at (-3.1623,-3.1623) {};
\node[atom=2.3mm] at (-2.8460,-2.2136) {};
\node[atom=2.3mm] at (-2.5298,-1.2649) {};
\node[atom=2.3mm] at (-2.2136,-0.3162) {};
\node[atom=2.3mm] at (-1.8974,0.6325) {};
\node[atom=2.3mm] at (-1.5811,1.5811) {};
\node[atom=2.3mm] at (-1.2649,2.5298) {};
\node[atom=2.3mm] at (-0.9487,3.4785) {};
\node[atom=2.3mm] at (-2.2136,-3.4785) {};
\node[atom=2.3mm] at (-1.8974,-2.5298) {};
\node[atom=2.3mm, fill=atomI] at (0.0000,-3.1623) {};
\draw[black, line width=1.1pt] (0.0000,-3.1623) circle (0.25);
\node[atom=2.3mm] at (-1.5811,-1.5811) {};
\node[atom=2.3mm] at (-1.2649,-0.6325) {};
\node[atom=2.3mm] at (-0.9487,0.3162) {};
\node[atom=2.3mm] at (-0.6325,1.2649) {};
\node[atom=2.3mm] at (-0.3162,2.2136) {};
\node[atom=2.3mm, fill=atomI] at (0.0000,3.1623) {};
\draw[black, line width=1.1pt] (0.0000,3.1623) circle (0.25);
\node[atom=2.3mm] at (-0.9487,-2.8460) {};
\node[atomB=2.3mm] at (0.9487,-2.8460) {};
\node[atom=2.3mm] at (-0.6325,-1.8974) {};
\node[atomB=2.3mm] at (0.6325,-1.8974) {};
\node[atom=2.3mm] at (-0.3162,-0.9487) {};
\node[atomB=2.3mm] at (0.3162,-0.9487) {};
\node[atom=2.3mm, fill=atomI] at (0.0000,0.0000) {};
\draw[black, line width=1.1pt] (0.0000,0.0000) circle (0.25);
\node[atomB=2.3mm] at (2.2136,-3.4785) {};
\node[atomB=2.3mm] at (1.8974,-2.5298) {};
\node[atomB=2.3mm] at (1.5811,-1.5811) {};
\node[atomB=2.3mm] at (1.2649,-0.6325) {};
\node[atomB=2.3mm] at (0.9487,0.3162) {};
\node[atomB=2.3mm] at (0.6325,1.2649) {};
\node[atomB=2.3mm] at (0.3162,2.2136) {};
\node[atomB=2.3mm] at (3.1623,-3.1623) {};
\node[atomB=2.3mm] at (2.8460,-2.2136) {};
\node[atomB=2.3mm] at (2.5298,-1.2649) {};
\node[atomB=2.3mm] at (2.2136,-0.3162) {};
\node[atomB=2.3mm] at (1.8974,0.6325) {};
\node[atomB=2.3mm] at (1.5811,1.5811) {};
\node[atomB=2.3mm] at (1.2649,2.5298) {};
\node[atomB=2.3mm] at (0.9487,3.4785) {};
\node[atomB=2.3mm] at (3.4785,-0.9487) {};
\node[atomB=2.3mm] at (3.1623,0.0000) {};
\node[atomB=2.3mm] at (2.8460,0.9487) {};
\node[atomB=2.3mm] at (2.5298,1.8974) {};
\node[atomB=2.3mm] at (2.2136,2.8460) {};
\node[atomB=2.3mm] at (3.4785,2.2136) {};
\node[atomB=2.3mm] at (3.1623,3.1623) {};
\end{scope}
\draw[atomLine] (7.6-3.25,-3.25) rectangle (7.6+3.25,3.25);
\draw[fmGray, dashed, line width=0.6pt] (7.6,-3.25) -- (7.6,3.25);

  \node[font=\small] at (0,-3.85) {\textit{а}};
  \node[font=\small] at (7.6,-3.85) {\textit{б}};
\end{tikzpicture}
\caption{Решётка совпадающих узлов $\Sigma5$ (поворот на $36{,}87^\circ$
вокруг $[001]$, показана плоскость $(001)$). \textit{а}~— наложение двух
решёток: синие точки и красные кружки; совпадающие узлы обведены, ячейка
РСУ выделена зелёным; зелёные точки внутри неё~— узлы полной решётки
наложения (ПРН), шаг которой в $\sqrt5$ раз меньше периода. \textit{б}~—
симметричная граница наклона $\Sigma5\,(310)$: слева решётка первого
зерна, справа второго; на границе (штриховая линия) каждый узел~—
совпадающий (красные атомы)}
\label{fig:csl5}
\end{figure}

\begin{table}[ht]
\centering
\caption{Разориентации кубических решёток с малыми $\Sigma$ (ось и
наименьший угол поворота) и допустимые отклонения по критерию
Брэндона~\eqref{eq:brandon}}
\label{tab:csl}
\small
\begin{tabular}{@{}rclc@{\hspace{10mm}}rclc@{}}
\toprule
$\Sigma$ & Ось & Угол & $\Delta\theta_\text{max}$ & $\Sigma$ & Ось & Угол & $\Delta\theta_\text{max}$\\
\midrule
3   & $\nfam{111}$ & $60{,}00^\circ$ & $8{,}7^\circ$ & 15  & $\nfam{210}$ & $48{,}19^\circ$ & $3{,}9^\circ$\\
5   & $\nfam{100}$ & $36{,}87^\circ$ & $6{,}7^\circ$ & 17a & $\nfam{100}$ & $28{,}07^\circ$ & $3{,}6^\circ$\\
7   & $\nfam{111}$ & $38{,}21^\circ$ & $5{,}7^\circ$ & 19a & $\nfam{110}$ & $26{,}53^\circ$ & $3{,}4^\circ$\\
9   & $\nfam{110}$ & $38{,}94^\circ$ & $5{,}0^\circ$ & 21a & $\nfam{111}$ & $21{,}79^\circ$ & $3{,}3^\circ$\\
11  & $\nfam{110}$ & $50{,}48^\circ$ & $4{,}5^\circ$ & 25a & $\nfam{100}$ & $16{,}26^\circ$ & $3{,}0^\circ$\\
13a & $\nfam{100}$ & $22{,}62^\circ$ & $4{,}2^\circ$ & 27a & $\nfam{110}$ & $31{,}59^\circ$ & $2{,}9^\circ$\\
13b & $\nfam{111}$ & $27{,}80^\circ$ & $4{,}2^\circ$ & 29a & $\nfam{100}$ & $43{,}60^\circ$ & $2{,}8^\circ$\\
\bottomrule
\end{tabular}
\end{table}

Для поворота вокруг оси $[hkl]$ угол, дающий РСУ, удовлетворяет условию
$\operatorname{tg}(\theta/2)=(n/m)\sqrt{N}$, где $N=h^2+k^2+l^2$, а $m$ и
$n$~— взаимно простые целые; тогда $\Sigma=m^2+Nn^2$, делённое на 2 столько
раз, сколько возможно (соотношения Кронберга--Вильсона в форме Ранганатана).
Например, для $\Sigma5$: ось $[100]$ ($N=1$), $m=3$, $n=1$:
$\Sigma=(9+1)/2=5$; для $\Sigma3$: ось $[111]$, $m=3$, $n=1$:
$(9+3)/4=3$ (табл.~\ref{tab:csl}).

\paragraph{Специальные границы}
Если граница совпадает с плотноупакованной плоскостью РСУ, в ней много
совпадающих узлов, решётки сопрягаются хорошо, и энергия границы мала. Такие
границы называют \term{специальными} (особыми); на зависимости энергии от
угла разворота им соответствуют резкие провалы (рис.~\ref{fig:gb-energy},\,\textit{б}).
Практически специальными считают границы с $1\le\Sigma\lesssim25$--$29$.
Самая «специальная»~— когерентная двойниковая граница $\Sigma3$ по $\plfam{111}$
в ГЦК-металлах: её энергия в десятки раз меньше, чем у обычной
высокоугловой границы (у меди $\approx0{,}02$ и $\approx0{,}6$~Дж/м$^2$).

\begin{figure}[ht]
\centering
\begin{tikzpicture}
\begin{axis}[name=rs, width=0.44\textwidth, height=0.36\textwidth,
  xlabel={$\theta/\theta_m$}, ylabel={$E/E_m$}, xmin=0, xmax=2.2, ymin=0, ymax=1.1,
  tick label style={font=\footnotesize}, label style={font=\small}]
  \addplot[fmBlue, line width=1.1pt, domain=0.001:1, samples=120] {x*(1-ln(x))};
  \addplot[fmBlue, line width=1.1pt, dashed, domain=1:2.2, samples=2] {1};
  \node[font=\footnotesize, anchor=north west, align=left] at (axis cs:0.35,0.55) {Рид--Шокли};
  \node[font=\footnotesize, anchor=north] at (axis cs:1.6,0.95) {высокоугловые};
\end{axis}
\begin{axis}[at={(rs.east)}, anchor=west, xshift=12mm, width=0.5\textwidth, height=0.36\textwidth,
  xlabel={угол разворота вокруг $\nfam{110}$, град}, ylabel={$E/E_\text{пл}$},
  xmin=0, xmax=180, ymin=0, ymax=1.15, xtick={0,30,...,180},
  tick label style={font=\footnotesize}, label style={font=\small}]
  \addplot[fmRed, line width=1.1pt] table[x=theta, y=E] {figures/gb-energy.dat};
  \foreach \t/\l/\y in {38.94/9/0.78, 50.48/11/0.53, 109.47/3/0.53, 129.52/11/0.73, 141.06/9/0.78}{
    \edef\tmp{\noexpand\node[font=\noexpand\scriptsize, anchor=north] at (axis cs:\t,\y) {$\noexpand\Sigma\l$};}\tmp}
  \node[font=\scriptsize, anchor=west] at (axis cs:73,0.07) {$\Sigma3$};
\end{axis}
\end{tikzpicture}
\caption{Энергия границы зерна: \textit{а}~— малоугловые границы,
формула Рида--Шокли~\eqref{eq:read-shockley} в приведённых координатах;
\textit{б}~— схема зависимости энергии симметричных границ наклона вокруг
$\nfam{110}$ в ГЦК-металле от угла разворота: провалы при специальных
разориентациях; глубже всех~— когерентный двойник $\Sigma3$ $(111)$ при
$70{,}5^\circ$}
\label{fig:gb-energy}
\end{figure}

Если граница идёт под небольшим углом к плотноупакованной плоскости РСУ,
она становится \emph{ступенчатой} (фасетированной): большей частью своей
площади она располагается по плоскостям с максимальной плотностью
совпадающих узлов.

Специальные свойства сохраняются и при небольших отклонениях
разориентации от специальной. Отклонение «принимают на себя»
\term{зернограничные дислокации}: они приспосабливают (аккомодируют)
структуру границы к фактической разориентации~— так же, как решёточные
дислокации образуют малоугловую границу в совершенном кристалле.
Предельное отклонение, при котором это ещё возможно, оценивают по
\term{критерию Брэндона}:
\begin{keyeq}\label{eq:brandon}
\Delta\theta_\text{max}=\frac{\theta_0}{\sqrt\Sigma},\qquad\theta_0\approx15^\circ
\end{keyeq}
(табл.~\ref{tab:csl})\ispr{В оригинале (формула (1.152)) критерий
записан как $\Delta\theta=(10\div15)/(\pi\sqrt\Sigma)$~рад, что даёт для
$\Sigma3$ около $160^\circ$. Критерий Брэндона: $\Delta\theta=\theta_0/\sqrt\Sigma$,
где $\theta_0=10$--$15^\circ$ (граница малоугловых границ);
в радианах $\theta_0=(10\div15)\pi/180$.}. Чем больше $\Sigma$, тем реже
совпадающие узлы и тем меньше допустимое отклонение.

\paragraph{Пространство разориентаций}
Разворот двух зёрен задаётся осью $\vb l$ и углом $\theta$ (или матрицей
поворота $A_0(\vb l,\theta)$, формула~\eqref{eq:R-axis}). Все возможные
развороты можно изобразить точками шара радиуса $\pi$ (направление из центра~—
ось, расстояние~— угол). С учётом веса, как и для углов Эйлера, «объём»
этого пространства равен $8\pi^2$.

Но из-за симметрии кристаллов один и тот же разворот описывается
несколькими способами. Если $R_i$~— поворотные элементы симметрии
кристалла, то эквивалентны описания
\begin{equation}\label{eq:equiv-misor}
A_i=A_0R_i^{-1}
\end{equation}
(поворот кристалла как целого операцией симметрии ничего не меняет). Для
кубических кристаллов (24 поворотных элемента группы $432$) существует
\vazhno{24 эквивалентных описания} разориентации с разными осями и углами
$\theta_i$. Описание с наименьшим углом называют \term{разориентировкой}
(англ. disorientation); обычно именно его и приводят.

Поэтому угол разориентации не может быть сколь угодно большим.
\term{Предельный угол} вокруг оси $[mnp]$~— наибольший угол, при котором
описание ещё остаётся наименьшим; при $\theta>\theta_\text{пр}$ найдётся
эквивалентное описание с меньшим углом. Для кубических кристаллов
предельные углы равны $45^\circ$ вокруг $\nfam{100}$, $60{,}7^\circ$ вокруг
$\nfam{110}$ и $60^\circ$ вокруг $\nfam{111}$, а наибольший возможный угол
разориентации~— $62{,}8^\circ$ вокруг оси $\nfam{1\,1\,\sqrt2-1}$, близкой
к $\nfam{221}$. Если зёрна ориентированы случайно, углы разориентации
распределены по закону Маккензи (рис.~\ref{fig:mackenzie}): чаще всего
встречаются углы около $45^\circ$, средний угол $40{,}7^\circ$, а
малоугловых границ мало.

\begin{figure}[ht]
\centering
\begin{tikzpicture}
\begin{axis}[width=0.62\textwidth, height=0.34\textwidth,
  xlabel={угол разориентации $\theta$, град}, ylabel={доля, 1/град},
  xmin=0, xmax=65, ymin=0, ymax=0.048, xtick={0,10,...,60},
  yticklabel style={/pgf/number format/fixed, /pgf/number format/precision=2},
  scaled y ticks=false,
  tick label style={font=\footnotesize}, label style={font=\small}]
  \addplot[const plot mark mid, fill=fmBlue!25, draw=fmBlue!70] table[x=theta, y=p] {figures/gb-mackenzie.dat} \closedcycle;
  \draw[fmRed, dashed] (axis cs:62.8,0) -- (axis cs:62.8,0.044);
  \node[font=\scriptsize, fmRed, anchor=south east] at (axis cs:62.8,0.040) {$62{,}8^\circ$};
\end{axis}
\end{tikzpicture}
\caption{Распределение Маккензи: доля границ с данным углом разориентации
в поликристалле со случайно ориентированными кубическими зёрнами (расчёт
методом Монте-Карло; учтены 24 эквивалентных описания~\eqref{eq:equiv-misor})}
\label{fig:mackenzie}
\end{figure}

\paragraph{Полная решётка наложения и O-решётка}
Если РСУ~— пересечение решёток двух зёрен, то \term{полная решётка наложения}
(ПРН, англ. DSC~— displacement shift complete lattice)~— наименьшая
решётка, содержащая узлы \emph{обеих} решёток\ispr{В оригинале ПРН
определена как объединение множеств узлов $L_1\cup L_2$. Объединение двух
решёток, вообще говоря, не является решёткой; ПРН~— наименее плотная
решётка, содержащая это объединение (порождённая им).}. Её векторы~—
наименьшие смещения одной решётки относительно другой, сохраняющие картину
совпадений (рис.~\ref{fig:csl5},\,\textit{а}). Объём ячейки РСУ в $\Sigma$ раз
больше, а объём ячейки ПРН в $\Sigma$ раз меньше объёма ячейки решётки
$V$:
\[
V_\text{РСУ}=\Sigma V,\qquad V_\text{ПРН}=V/\Sigma,\qquad
V_\text{РСУ}V_\text{ПРН}=V^2 .
\]
Векторы ПРН~— это возможные векторы Бюргерса зернограничных дислокаций;
они \emph{короче} векторов решётки (для $\Sigma5$ в плоскости $(001)$~—
в $\sqrt5$ раз). Наиболее плотноупакованное направление ПРН задаёт вектор
Бюргерса этих дислокаций.

Для расчёта дислокационного строения \emph{любой} границы (а не только
специальной) Боллман ввёл \term{O-решётку}~— множество точек, в которых
совпадают не только узлы, но и любые физически эквивалентные точки
элементарных ячеек двух решёток. O-решётка существует при любом развороте;
РСУ~— её частный случай. Для малоугловых границ шаг O-решётки
$b/(2\sin\frac\theta2)$ совпадает с расстоянием между дислокациями~—
формула Франка получается как частный случай.

\begin{itogi}
\begin{itemize}
\item Граница задаётся 5 макроскопическими параметрами (3~— разворот, 2~—
нормаль) и 3 микроскопическими (сдвиг $\vb t$).
\item Граница наклона: ось разворота в плоскости границы; кручения:
ось перпендикулярна границе.
\item Малоугловые границы~— стенки и сетки дислокаций: $D=b/\theta$;
формула Франка $\vb d=2\sin\frac\theta2(\vb r\times\vb u)$; энергия
$E=E_0\theta(A-\ln\theta)$.
\item Высокоугловые границы описывают РСУ: $\Sigma$~— число узлов на один
совпадающий (в кубических решётках нечётно); специальные границы
($\Sigma\lesssim25$--$29$) имеют пониженную энергию; критерий Брэндона
$\Delta\theta=15^\circ/\sqrt\Sigma$.
\item У кубических кристаллов 24 эквивалентных описания разворота;
наибольший угол разориентации $62{,}8^\circ$.
\item ПРН (DSC)~— решётка векторов Бюргерса зернограничных дислокаций,
$V_\text{РСУ}V_\text{ПРН}=V^2$; O-решётка Боллмана обобщает РСУ на любые
развороты.
\end{itemize}
\end{itogi}

\enlargethispage{2\baselineskip}
\begin{proverka}
\item Сколько параметров задают плоскую границу зерна с макроскопической
и с атомной точки зрения?
\item Чем граница наклона отличается от границы кручения? Из каких
дислокаций они состоят, если малоугловые?
\item Субграница наклона в меди ($b=0{,}256$~нм) содержит дислокации через
50~нм. Найдите угол разориентации.
\item Почему поворот на $53{,}13^\circ$ вокруг $[001]$ даёт ту же РСУ $\Sigma5$,
что и поворот на $36{,}87^\circ$?
\item Почему двойниковая граница $\Sigma3$ в ГЦК-металлах имеет очень
малую энергию?
\item Может ли угол разориентации двух кубических зёрен равняться
$70^\circ$?
\end{proverka}

\section{Кристаллография мартенситных превращений}\label{sec:martensite}
% ----------------------------------------------------------------------
%  1.8. Кристаллография мартенситных превращений
% ----------------------------------------------------------------------
\begin{orientir}
Мартенситное превращение~— перестройка решётки \vazhno{сдвигом}, без
диффузии. В~1.8.1 по рельефу на шлифе выясним, что макроскопически это
\vazhno{деформация с инвариантной плоскостью} и почему, кроме деформации
решётки, нужна ещё деформация, не меняющая решётку (скольжение или
двойникование). В~1.8.2~— как ГЦК-аустенит переходит в
мартенсит (деформация Бейна) и какие \vazhno{ориентационные соотношения}
наблюдаются на опыте.
\end{orientir}

Название «мартенсит» (по имени немецкого металловеда А.~Мартенса)
первоначально относилось к твёрдой составляющей закалённой стали. Сейчас
мартенситные (их называют также \emph{сдвиговыми}) превращения известны в
сплавах, чистых металлах, керамиках, минералах, неорганических
соединениях, затвердевших газах и полимерах.

\begin{opredelenie}
\term{Мартенситное превращение}~— бездиффузионное фазовое превращение, при
котором решётка перестраивается \emph{кооперативным} смещением атомов на
расстояния меньше межатомного; каждый атом сохраняет своих соседей, а
состав фаз одинаков.
\end{opredelenie}

Превращение идёт с очень большой скоростью: кристалл мартенсита в стали
растёт со скоростью порядка $10^5$~см/с ($\sim1$~км/с), сравнимой со
скоростью звука. Как правило, мартенсит образуется \emph{атермически}:
количество мартенсита определяется температурой охлаждения ниже точки
$M_\text{н}$ и почти не растёт при выдержке (хотя известно и изотермическое
мартенситное превращение).

\begin{masshtab}[мартенсит]{-10/-9.6/{смещения атомов}, -8.3/-7.3/{внутренние двойники}, -7/-4/{рейки и пластины}}
Каждый атом смещается меньше чем на межатомное расстояние (на десятые
доли нанометра и меньше), но согласованно с миллиардами соседей.
Кристаллы мартенсита~— рейки толщиной 0,1--0,5~мкм в низкоуглеродистой
стали и пластины размером вплоть до зерна (10--100~мкм) в
высокоуглеродистой; внутренние двойники в пластинах~— 5--50~нм.
Пластина длиной 100~мкм при скорости роста $\sim1$~км/с вырастает за
$\sim10^{-7}$~с.
\end{masshtab}

\subsection{Морфология мартенситных превращений}\label{ssec:martensite-morph}

Важнейшие сведения о геометрии превращения дал металлографический
анализ \emph{рельефа}, возникающего на предварительно полированной
поверхности шлифа там, где образовался кристалл мартенсита
(рис.~\ref{fig:mart-relief}). Установлено:
\begin{itemize}
\item поверхность превращённой области наклоняется к плоскости шлифа,
но остаётся \emph{плоской};
\item прямые риски, нанесённые до превращения, остаются \emph{непрерывными}
и \emph{прямыми}, но испытывают излом там, где выходит на поверхность
граница кристалла мартенсита~— \term{плоскость габитуса}.
\end{itemize}

\begin{figure}[ht]
\centering
\begin{tikzpicture}[cubeproj, scale=1.35]
  \def\dz{0.55}
  % исходная фаза (x от 0 до 2)
  \fill[fmBlueBg] (0,0,0) -- (2,0,0) -- (2,0,1.5) -- (0,0,1.5) -- cycle;
  \fill[fmBlueBg!70!white] (0,0,1.5) -- (2,0,1.5) -- (2,2.5,1.5) -- (0,2.5,1.5) -- cycle;
  % мартенсит (x от 2 до 4): верхняя грань наклонена
  \fill[fmAmberBg] (2,0,0) -- (4,0,0) -- (4,0,1.5+\dz) -- (2,0,1.5) -- cycle;
  \fill[fmAmber!25] (2,0,1.5) -- (4,0,1.5+\dz) -- (4,2.5,1.5+\dz) -- (2,2.5,1.5) -- cycle;
  \fill[fmAmber!15] (4,0,0) -- (4,2.5,0) -- (4,2.5,1.5+\dz) -- (4,0,1.5+\dz) -- cycle;
  \draw[atomLine] (0,0,0) -- (4,0,0) -- (4,0,1.5+\dz) -- (2,0,1.5) -- (0,0,1.5) -- cycle
    (0,0,1.5) -- (0,2.5,1.5) -- (2,2.5,1.5) -- (4,2.5,1.5+\dz) -- (4,0,1.5+\dz)
    (4,0,0) -- (4,2.5,0) -- (4,2.5,1.5+\dz);
  \draw[atomLine!50, dashed] (0,2.5,1.5) -- (0,2.5,0) -- (4,2.5,0) (0,2.5,0) -- (0,0,0);
  % плоскость габитуса ABC
  \fill[fmRed!20, opacity=0.6] (2,0,0) -- (2,0,1.5) -- (2,2.5,1.5) -- (2,2.5,0) -- cycle;
  \draw[fmRed!80!black, line width=0.9pt] (2,0,0) -- (2,0,1.5) -- (2,2.5,1.5);
  \draw[fmRed!80!black, dashed] (2,2.5,1.5) -- (2,2.5,0) -- (2,0,0);
  \node[font=\small, fmRed!80!black, anchor=north] at (2,0,0) {$A$};
  \node[font=\small, fmRed!80!black, anchor=east] at (2,0,1.5) {$B$};
  \node[font=\small, fmRed!80!black, anchor=south] at (2,2.5,1.5) {$C$};
  % риски
  \foreach \y in {0.6,1.25,1.9}{
    \draw[fmBlue, line width=1pt] (0,\y,1.5) -- (2,\y,1.5) -- (4,\y,1.5+\dz);}
  \node[font=\small] at (1,0,0.6) {1};
  \node[font=\small] at (3,0,0.6) {2};
  \node[font=\small, fmBlue] at (0.6,1.9,1.75) {3};
\end{tikzpicture}
\caption{Рельеф на поверхности шлифа при образовании кристалла
мартенсита: 1~— исходная фаза; 2~— мартенсит; 3~— прямые риски,
нанесённые на поверхность до превращения; $ABC$~— плоскость габитуса.
Риски остаются прямыми и непрерывными, но изламываются на следе
плоскости габитуса}
\label{fig:mart-relief}
\end{figure}

Это означает, что рельеф создан \emph{однородной деформацией с
инвариантной плоскостью}: существует плоскость (габитус), которая не
искажается и не поворачивается, а все смещения направлены в одну
сторону и пропорциональны расстоянию от неё. Если $\vb p$~— единичная
нормаль к инвариантной плоскости, а $\vb d$~— вектор смещения точек,
отстоящих от неё на единицу длины, то
\begin{keyeq}\label{eq:ips}
\vb r'=\vb r+\vb d\,(\vb p\cdot\vb r),\qquad P=I+\vb d\,\tilde{\vb p}
\end{keyeq}
(ср. со смещением при скольжении~\eqref{eq:slip-displacement}). Два
предельных случая (рис.~\ref{fig:ips}): $\vb d\parallel\vb p$~—
одноосное растяжение или сжатие (дилатация), $\vb d\perp\vb p$~— простой
сдвиг. Любую деформацию с инвариантной плоскостью можно представить как
их комбинацию. У мартенсита в сталях сдвиговая составляющая $\approx0{,}2$, а
нормальная~— несколько процентов (удельный объём растёт на 3--4\,\%).

\begin{figure}[ht]
\centering
\begin{tikzpicture}[scale=1.4]
  \foreach \X/\h/\s/\l in {0/1.35/0/а, 2.6/1.0/0.5/б, 5.2/1.3/0.4/в}{
    \begin{scope}[shift={(\X,0)}]
      \fill[fmAmberBg] (0,0) -- (1.4,0) -- (1.4+\s,\h) -- (\s,\h) -- cycle;
      \draw[atomLine, dashed] (0,0) rectangle (1.4,1);
      \draw[fmAmber!80!black, line width=0.8pt] (0,0) -- (1.4,0) -- (1.4+\s,\h) -- (\s,\h) -- cycle;
      \draw[fmRed, line width=2pt] (-0.25,0) -- (1.65,0);
      \draw[arr, fmBlue] (0.7,1) -- ({0.7+\s},\h);
      \node[font=\small] at (0.7,-0.35) {\textit{\l}};
    \end{scope}}
  \node[font=\footnotesize, fmBlue, anchor=west] at (0.75,1.18) {$\vb d$};
  \node[font=\footnotesize, fmRed, anchor=west] at (7.05,0) {инвариантная плоскость};
\end{tikzpicture}
\caption{Деформации с инвариантной плоскостью (красная): \textit{а}~—
дилатация ($\vb d\parallel\vb p$); \textit{б}~— сдвиг ($\vb d\perp\vb p$);
\textit{в}~— сдвиг с дилатацией. Пунктир~— исходная форма}
\label{fig:ips}
\end{figure}

Если перестройку решётки представлять как деформацию с инвариантной
плоскостью, то граница между когерентными областями новой и исходной фаз
и есть инвариантная плоскость~— плоскость габитуса. Чтобы сохранить
контакт с непревращённой частью кристалла, деформация должна включать не
только чистую деформацию (растяжение, сжатие, сдвиг), но и \emph{поворот}.

\paragraph{Деформация решётки и деформация с инвариантной решёткой}
Опыт показал, что макроскопическая деформация превращённого объёма
(найденная по рельефу) \emph{отличается} от деформации, переводящей
решётку аустенита в решётку мартенсита. Разницу обеспечивает
дополнительная деформация, которая не меняет решётку,~— \term{деформация
с инвариантной решёткой}: скольжение или двойникование
(рис.~\ref{fig:lis}). Деформация решётки задаёт \emph{структуру} новой фазы,
а деформация с инвариантной решёткой~— её \emph{форму}.

\begin{figure}[ht]
\centering
\begin{tikzpicture}[scale=0.95]
  \def\w{0.42} \def\h{0.38} \def\g{0.105}
  % (а) исходная решётка
  \begin{scope}
    \foreach \j in {0,...,7}{\foreach \i in {0,...,3}{
      \draw[fmBlue, fill=fmBlueBg] ({\i*\w},{\j*\h}) rectangle ++(\w,\h);}}
    \node[font=\small] at (0.84,-0.4) {\textit{а}};
  \end{scope}
  % (б) деформация решётки: каждая ячейка скошена, форма изменилась
  \begin{scope}[shift={(2.6,0)}]
    \draw[atomLine!50, dashed] (0,0) rectangle ({4*\w},{8*\h});
    \foreach \j in {0,...,7}{\foreach \i in {0,...,3}{
      \draw[fmAmber!80!black, fill=fmAmberBg] ({\i*\w+\j*\g},{\j*\h}) -- ++(\w,0) -- ++(\g,\h) -- ++(-\w,0) -- cycle;}}
    \node[font=\small] at (1.2,-0.4) {\textit{б}};
  \end{scope}
  % (в) + скольжение: каждые 4 слоя сдвинуты назад на период
  \begin{scope}[shift={(5.6,0)}]
    \draw[atomLine!50, dashed] (0,0) rectangle ({4*\w},{8*\h});
    \foreach \j in {0,...,7}{
      \pgfmathsetmacro{\off}{\j*\g-\w*floor(\j/4)}
      \foreach \i in {0,...,3}{
        \draw[fmAmber!80!black, fill=fmAmberBg] ({\i*\w+\off},{\j*\h}) -- ++(\w,0) -- ++(\g,\h) -- ++(-\w,0) -- cycle;}}
    \draw[fmRed, line width=1pt] (-0.15,{4*\h}) -- ({4*\w+0.6},{4*\h});
    \node[font=\small] at (1.0,-0.4) {\textit{в}};
  \end{scope}
  % (г) + двойникование: чередование слоёв с противоположным скосом
  \begin{scope}[shift={(8.6,0)}]
    \draw[atomLine!50, dashed] (0,0) rectangle ({4*\w},{8*\h});
    \foreach \j in {0,...,7}{
      \pgfmathtruncatemacro{\m}{mod(\j,4)}
      \ifnum\m<2
        \pgfmathsetmacro{\off}{\m*\g}
        \foreach \i in {0,...,3}{
          \draw[fmAmber!80!black, fill=fmAmberBg] ({\i*\w+\off},{\j*\h}) -- ++(\w,0) -- ++(\g,\h) -- ++(-\w,0) -- cycle;}
      \else
        \pgfmathsetmacro{\off}{(4-\m)*\g}
        \foreach \i in {0,...,3}{
          \draw[fmRed!70!black, fill=fmRedBg] ({\i*\w+\off},{\j*\h}) -- ++(\w,0) -- ++(-\g,\h) -- ++(-\w,0) -- cycle;}
      \fi}
    \node[font=\small] at (0.95,-0.4) {\textit{г}};
  \end{scope}
\end{tikzpicture}
\caption{Деформация решётки и деформация с инвариантной решёткой:
\textit{а}~— исходная решётка; \textit{б}~— деформация решётки меняет и
структуру, и форму (пунктир~— исходная форма); \textit{в}~— та же
структура, но скольжение (красная линия~— плоскость скольжения)
возвращает макроскопическую форму, оставляя ступеньки; \textit{г}~— то же
достигается двойникованием: слои двух двойниковых ориентаций (оранжевые и
розовые) чередуются, и средняя форма не меняется}
\label{fig:lis}
\end{figure}

Так объясняется \emph{субструктура} мартенсита: в низкоуглеродистых сталях
это пакеты реек с высокой плотностью дислокаций (деформация с
инвариантной решёткой~— скольжением), в высокоуглеродистых и в сплавах
Fe--Ni~— линзовидные пластины с тонкими внутренними двойниками
(двойникованием).

\subsection{Кристаллография мартенситного превращения в сталях}\label{ssec:martensite-cryst}

При закалке стали ГЦК-решётка аустенита (твёрдого раствора углерода в
$\gamma$-Fe) превращается в \emph{объёмноцентрированную тетрагональную}
(ОЦТ) решётку мартенсита. Как это возможно при минимальных смещениях
атомов, первым показал Бейн.

\paragraph{Соответствие решёток по Бейну}
В двух соседних ГЦК-ячейках можно выделить ОЦТ-ячейку
(рис.~\ref{fig:bain},\,\textit{а}): её основание~— квадрат со стороной
$a_\gamma/\sqrt2$, построенный на направлениях $[110]_\gamma$ и
$[\br110]_\gamma$, высота~— $a_\gamma$ вдоль $[001]_\gamma$, а в центре~—
атом из центра общей грани. Отношение $c/a$ этой ячейки $\sqrt2=1{,}41$.
Чтобы получить ОЦК (или слабо тетрагональную) ячейку мартенсита, её
нужно сжать вдоль $c$ и растянуть в основании~— это и есть
\term{деформация Бейна}:
\begin{keyeq}\label{eq:bain}
B=\begin{pmatrix}\eta_1&0&0\\0&\eta_1&0\\0&0&\eta_3\end{pmatrix}
\approx\begin{pmatrix}1{,}12&0&0\\0&1{,}12&0\\0&0&0{,}83\end{pmatrix},\qquad
\eta_1=\frac{\sqrt2\,a_\alpha}{a_\gamma},\quad \eta_3=\frac{c_\alpha}{a_\gamma}
\end{keyeq}
(в осях $[100]_\alpha$, $[010]_\alpha$, $[001]_\alpha$): растяжение примерно
на 12\,\% вдоль $[100]_\alpha$ и $[010]_\alpha$ и сжатие примерно на 17\,\%
вдоль $[001]_\alpha$. Объём при этом увеличивается в
$\eta_1^2\eta_3\approx1{,}04$ раза.

\begin{figure}[ht]
\centering
\begin{tikzpicture}[cubeproj, scale=2.0]
  % (а) две ГЦК-ячейки и выделенная ОЦТ-ячейка
  \draw[atomLine!45, dashed] (0,1,0) -- (2,1,0) (0,1,0) -- (0,0,0) (0,1,0) -- (0,1,1) (1,1,0) -- (1,1,1);
  \draw[atomLine] (0,0,0) -- (2,0,0) -- (2,1,0) -- (2,1,1) -- (0,1,1) -- (0,0,1) -- cycle
    (0,0,1) -- (2,0,1) -- (2,1,1) (2,0,0) -- (2,0,1) (1,0,0) -- (1,0,1) -- (1,1,1);
  \foreach \p in {(1,0,0),(1.5,0.5,0),(1,1,0),(0.5,0.5,0)}{\draw[fmRed!80!black, line width=1.1pt] \p -- ++(0,0,1);}
  \draw[fmRed!80!black, line width=1.1pt] (1,0,0) -- (1.5,0.5,0) -- (1,1,0) -- (0.5,0.5,0) -- cycle
    (1,0,1) -- (1.5,0.5,1) -- (1,1,1) -- (0.5,0.5,1) -- cycle;
  \foreach \x in {0,1,2}{\foreach \y in {0,1}{\foreach \z in {0,1}{
    \node[atom=2.6mm] at (\x,\y,\z) {};}}}
  \foreach \x in {0.5,1.5}{\foreach \z in {0,1}{\node[atom=2.6mm] at (\x,0.5,\z) {};}
    \foreach \y in {0,1}{\node[atom=2.6mm] at (\x,\y,0.5) {};}}
  \foreach \x in {0,1,2}{\node[atom=2.6mm] at (\x,0.5,0.5) {};}
  \foreach \p in {(1,0,0),(1.5,0.5,0),(1,1,0),(0.5,0.5,0),(1,0,1),(1.5,0.5,1),(1,1,1),(0.5,0.5,1)}{
    \node[atomB=2.6mm] at \p {};}
  \node[atomI=2.6mm] at (1,0.5,0.5) {};
  \node[font=\small] at (1,0,-0.35) {\textit{а}};
  % (б) деформация Бейна
  \begin{scope}[shift={(2.9,0,0)}]
    \def\a{0.707}
    \draw[fmRed!80!black, line width=0.8pt] (0,0,0) -- (\a,0,0) -- (\a,\a,0) (0,0,0) -- (0,0,1) -- (\a,0,1) -- (\a,\a,1) -- (0,\a,1) -- (0,0,1) (\a,0,0) -- (\a,0,1) (\a,\a,0) -- (\a,\a,1);
    \draw[fmRed!50, dashed] (0,0,0) -- (0,\a,0) -- (\a,\a,0) (0,\a,0) -- (0,\a,1);
    \foreach \p in {(0,0,0),(\a,0,0),(\a,\a,0),(0,0,1),(\a,0,1),(\a,\a,1),(0,\a,1),(0,\a,0)}{\node[atomB=2.4mm] at \p {};}
    \node[atomI=2.4mm] at ({\a/2},{\a/2},0.5) {};
    \node[font=\scriptsize, anchor=north] at ({\a/2},0,0) {$a_\gamma/\sqrt2$};
    \node[font=\scriptsize, anchor=east] at (0,0,0.5) {$a_\gamma$};
    \draw[arrB] (1.05,0.35,0.5) -- (1.55,0.35,0.5);
    \node[font=\scriptsize, above] at (1.3,0.35,0.55) {$B$};
    \begin{scope}[shift={(1.85,0,0.08)}]
      \def\b{0.79} \def\c{0.83}
      \draw[atomLine, line width=0.8pt] (0,0,0) -- (\b,0,0) -- (\b,\b,0) (0,0,0) -- (0,0,\c) -- (\b,0,\c) -- (\b,\b,\c) -- (0,\b,\c) -- (0,0,\c) (\b,0,0) -- (\b,0,\c) (\b,\b,0) -- (\b,\b,\c);
      \draw[atomLine!50, dashed] (0,0,0) -- (0,\b,0) -- (\b,\b,0) (0,\b,0) -- (0,\b,\c);
      \foreach \p in {(0,0,0),(\b,0,0),(\b,\b,0),(0,0,\c),(\b,0,\c),(\b,\b,\c),(0,\b,\c),(0,\b,0)}{\node[atomB=2.4mm] at \p {};}
      \node[atomI=2.4mm] at ({\b/2},{\b/2},{\c/2}) {};
      \node[font=\scriptsize, anchor=north] at ({\b/2},0,0) {$a_\alpha$};
      \node[font=\scriptsize, anchor=west] at (\b,0,{\c/2}) {$c_\alpha$};
      \draw[arr, fmBlue] ({\b/2},{\b/2},{\c+0.45}) -- ({\b/2},{\b/2},{\c+0.12});
      \draw[arr, fmBlue] (-0.45,0,{\c/2}) -- (-0.12,0,{\c/2});
    \end{scope}
    \node[font=\small] at (1.5,0,-0.35) {\textit{б}};
  \end{scope}
\end{tikzpicture}
\caption{Соответствие решёток по Бейну: \textit{а}~— две ГЦК-ячейки
аустенита и выделенная в них ОЦТ-ячейка ($c/a=\sqrt2$; оранжевые атомы~—
её вершины, красный~— центр); \textit{б}~— деформация Бейна: сжатие вдоль
$c$ примерно на 17\,\% и растяжение в основании на 12\,\% переводят её в
ячейку мартенсита}
\label{fig:bain}
\end{figure}

Соответствие Бейна сразу объясняет \emph{тетрагональность} мартенсита.
Атомы углерода в аустените занимают октаэдрические пустоты~— середины
рёбер и центр ГЦК-ячейки (п.~\ref{ssec:elem-types}). При соответствии
Бейна все эти пустоты попадают в \emph{одну} из трёх подрешёток
октаэдрических пустот ОЦТ-решётки~— ту, у которой сплюснутый октаэдр
вытянут вдоль оси $c$. Внедрённые атомы углерода раздвигают решётку вдоль
$c$, и мартенсит оказывается тетрагональным:
$c/a\approx1+0{,}046\,C$, где $C$~— содержание углерода в массовых
процентах.

Деформации Бейна соответствуют \term{ориентационные соотношения Бейна}:
\[
(001)_\gamma\parallel(001)_\alpha,\qquad [100]_\gamma\parallel[110]_\alpha .
\]
В чистом виде они наблюдаются лишь при превращении в очень тонких плёнках.

\paragraph{Деформация Бейна не имеет инвариантной плоскости}
Применим деформацию Бейна к сфере единичного радиуса: она превратится в
эллипсоид с полуосями $\eta_1,\eta_1,\eta_3$ (рис.~\ref{fig:bain-cone}).
Неискажёнными (сохранившими длину) остаются лишь направления, по которым
эллипсоид пересекает сферу,~— они образуют \emph{конус} $L'$, в который
перешёл конус $L$ исходных направлений. Найдём его: точка эллипсоида
$(x,z)$ лежит на сфере, если $x^2+z^2=1$ и
$x^2/\eta_1^2+z^2/\eta_3^2=1$. При $\eta_1=1{,}12$, $\eta_3=0{,}83$
угол конуса $L'$ с осью $c$ равен $56{,}2^\circ$, а исходного конуса $L$~—
$47{,}9^\circ$: неискажённые направления \emph{поворачиваются}.

\begin{figure}[ht]
\centering
\begin{tikzpicture}[scale=1.9]
  \fill[fmBlueBg] (0,0) circle (1);
  \fill[fmAmber!25, opacity=0.8] (0,0) ellipse (1.12 and 0.83);
  \draw[atomLine, line width=0.7pt] (0,0) circle (1);
  \draw[fmAmber!80!black, line width=0.9pt] (0,0) ellipse (1.12 and 0.83);
  \draw[arr, atomLine!70] (-1.35,0) -- (1.45,0) node[right, font=\small, black]{$[100]_\alpha$};
  \draw[arr, atomLine!70] (0,-1.15) -- (0,1.3) node[above, font=\small, black]{$[001]_\alpha$};
  \foreach \sx in {1,-1}{\foreach \sz in {1,-1}{
    \draw[fmBlue, dashed, line width=0.8pt] (0,0) -- ({\sx*0.7417*1.2},{\sz*0.6707*1.2});
    \draw[fmRed, line width=1pt] (0,0) -- ({\sx*0.8307*1.2},{\sz*0.5567*1.2});
    \fill[fmBlue] ({\sx*0.7417},{\sz*0.6707}) circle (0.025);
    \fill[fmRed] ({\sx*0.8307},{\sz*0.5567}) circle (0.025);}}
  \draw[arr, fmGray] (0.7417,0.6707) to[bend left=20] (0.8307,0.5567);
  \node[font=\small, fmBlue, anchor=south] at ({0.7417*1.2},{0.6707*1.2}) {$L$};
  \node[font=\small, fmRed, anchor=west] at ({0.8307*1.2},{0.5567*1.2}) {$L'$};
  \node[font=\footnotesize, anchor=north west] at (1.0,-0.05) {$\eta_1$};
  \node[font=\footnotesize, anchor=south west] at (0.02,0.83) {$\eta_3$};
  \node[font=\footnotesize, fmBlue!70!black, anchor=east] at (-0.72,0.75) {сфера};
  \node[font=\footnotesize, fmAmber!80!black, anchor=east] at (-1.1,-0.35) {эллипсоид};
\end{tikzpicture}
\caption{Деформация Бейна в сечении, содержащем ось $c$: сфера
единичного радиуса переходит в эллипсоид. Неискажённые направления
лежат на конусе $L'$ (пересечение сферы и эллипсоида); до деформации они
лежали на конусе $L$. Плоскости, которая осталась бы и неискажённой, и
неповёрнутой, нет}
\label{fig:bain-cone}
\end{figure}

Пусть плоскость $P'$ после деформации проходит через два неискажённых
направления $a'$ и $b'$ конуса $L'$; до деформации она проходила через
направления $a$ и $b$ (плоскость $P$). Но углы между неискажёнными
направлениями не сохраняются, поэтому \emph{ни одна} плоскость не
остаётся неискажённой. Причина в главных удлинениях: инвариантная
плоскость возможна только если одно из них равно 1, а два других~— по
разные стороны от 1. У деформации Бейна два удлинения больше 1, одно
меньше, и ни одно не равно 1.

\paragraph{Как получается инвариантная плоскость}
Добавление к деформации Бейна обычного скольжения не помогает: скольжение
сопровождается поворотом кристалла, и инвариантная плоскость по-прежнему
не возникает. Феноменологическая теория мартенситного превращения
(Векслер, Либерман, Рид; Боулс, Маккензи) строит наблюдаемую
макроскопическую деформацию $P_1$ из трёх частей:
\begin{keyeq}\label{eq:wlr}
P_1=R\,P\,B ,
\end{keyeq}
где $B$~— деформация Бейна (задаёт решётку мартенсита), $P$~—
деформация с инвариантной решёткой (двойникование или скольжение по
подходящей системе), $R$~— поворот. Величину сдвига $P$ подбирают так,
чтобы у произведения $PB$ одно главное удлинение стало равно 1; тогда
появляется неискажённая плоскость, а поворот $R$ возвращает её в исходное
положение~— получается инвариантная плоскость, т.\,е. плоскость
габитуса. Геометрически: одно неискажённое направление выбирают в
плоскости двойникования $K_1$ (оно сохраняется при двойниковании и
переходит в неискажённое при деформации Бейна), второе~— в плоскости
$K_2$; плоскость, проходящая через них, после двойникования и деформации
Бейна поворачивается, и поворот $R$ совмещает её с исходной. Теория
предсказывает плоскость габитуса, ориентационное соотношение и
макроскопический сдвиг; все три хорошо согласуются с опытом.

\paragraph{Наблюдаемые ориентационные соотношения}
Опыт даёт ориентационные соотношения, близкие к соотношению Бейна, но
отличающиеся от него поворотом $R$ на несколько градусов
(табл.~\ref{tab:or}).

\begin{table}[ht]
\centering
\caption{Ориентационные соотношения при мартенситных превращениях}
\label{tab:or}
\small
\begin{tabular}{@{}lllcl@{}}
\toprule
Соотношение & Плоскости & Направления & Вариантов & Где наблюдается\\
\midrule
Бейна & $(001)_\gamma\parallel(001)_\alpha$ & $[100]_\gamma\parallel[110]_\alpha$ & 3 & тонкие плёнки\\
Курдюмова--Закса & $(111)_\gamma\parallel(011)_\alpha$ & $[10\br1]_\gamma\parallel[11\br1]_\alpha$ & 24 & стали, 0,5--1,4\,\% C\\
Нисиямы & $(111)_\gamma\parallel(011)_\alpha$ & $[\br1\br12]_\gamma\parallel[0\br11]_\alpha$ & 12 & сплавы Fe--Ni\\
Бюргерса & $\plfam{110}_\beta\parallel(0001)_\alpha$ & $\nfam{111}_\beta\parallel\nfam{2\br1\br10}_\alpha$ & 12 & Zr, Ti (ОЦК$\to$ГПУ)\\
\bottomrule
\end{tabular}
\end{table}

В массивных образцах сталей с 0,5--1,4\,\% (масс.) углерода выполняются
соотношения \term{Курдюмова--Закса}: плотноупакованная плоскость $\plfam{111}$
аустенита параллельна плотноупакованной плоскости $\plfam{110}$ мартенсита,
а плотноупакованные направления $\nfam{110}_\gamma$ и $\nfam{111}_\alpha$
в этих плоскостях параллельны. Плоскость габитуса лежит в области
$\plfam{225}_\gamma$--$\plfam{259}_\gamma$. В сплавах Fe--Ni выполняются
соотношения \term{Нисиямы} (Нисиямы--Вассермана)\ispr{В оригинале
направления в соотношении Нисиямы записаны как
$[112]_\gamma\parallel[11\br1]_\alpha$, т.\,е. $\nfam{112}_\gamma\parallel\nfam{111}_\alpha$.
Это неверно: в соотношении Нисиямы--Вассермана
$\nfam{112}_\gamma\parallel\nfam{110}_\alpha$ (эквивалентно
$\nfam{110}_\gamma\parallel\nfam{100}_\alpha$). Оно отличается от
соотношения Курдюмова--Закса поворотом на $5{,}26^\circ$ вокруг нормали к
$(111)_\gamma$.}: плоскости те же, но в плоскости $(111)_\gamma$ решётка
повёрнута на $5{,}26^\circ$, так что параллельны $\nfam{112}_\gamma$ и
$\nfam{110}_\alpha$; габитус около $\plfam{3\,10\,15}_\gamma$.

\begin{smysl}
Ориентационное соотношение~— «отпечаток» механизма превращения:
плотноупакованные плоскости и направления двух фаз совмещены, потому что
решётка перестраивается сдвигом с минимальными смещениями атомов. Число
вариантов (24 для соотношения Курдюмова--Закса) показывает, сколько
по-разному ориентированных кристаллов мартенсита может возникнуть в одном
зерне аустенита.
\end{smysl}

Мартенситный характер имеют и некоторые полиморфные превращения. Так, при
переходе ОЦК $\beta$-Zr в ГПУ $\alpha$-Zr (и в титане) выполняются
\term{соотношения Бюргерса}: $\plfam{110}_\beta\parallel(0001)_\alpha$,
$\nfam{111}_\beta\parallel\nfam{2\br1\br10}_\alpha$. Более подробный анализ
требует рассмотрения дислокационных механизмов, кинетики и термодинамики
мартенситных превращений.

\begin{itogi}
\begin{itemize}
\item Мартенситное превращение~— бездиффузионная кооперативная
перестройка решётки сдвигом; скорость роста $\sim1$~км/с.
\item Макроскопически~— деформация с инвариантной плоскостью (габитусом):
$\vb r'=\vb r+\vb d(\vb p\cdot\vb r)$; на поверхности возникает рельеф.
\item Полная деформация $P_1=RPB$: деформация решётки $B$ (Бейн:
$\eta_1\approx1{,}12$, $\eta_3\approx0{,}83$), деформация с инвариантной
решёткой $P$ (скольжение или двойникование) и поворот $R$.
\item Деформация Бейна сама по себе не имеет инвариантной плоскости
(лишь конус неискажённых направлений).
\item Соответствие Бейна объясняет тетрагональность мартенсита:
$c/a\approx1+0{,}046\,C$.
\item Ориентационные соотношения: Курдюмова--Закса
$\plfam{111}_\gamma\parallel\plfam{110}_\alpha$,
$\nfam{110}_\gamma\parallel\nfam{111}_\alpha$; Нисиямы; Бюргерса
($\beta\to\alpha$ в Zr, Ti).
\end{itemize}
\end{itogi}

\begin{proverka}
\item Какие наблюдения рельефа показывают, что мартенситное превращение~—
деформация с инвариантной плоскостью?
\item Чем деформация решётки отличается от деформации с инвариантной
решёткой? Приведите примеры последней.
\item Найдите по~\eqref{eq:bain} изменение объёма при превращении аустенита
($a_\gamma=0{,}3584$~нм) в мартенсит ($a=0{,}2850$, $c=0{,}2958$~нм).
\item Почему деформация Бейна не имеет инвариантной плоскости?
\item Проверьте, что направление $[10\br1]_\gamma$ лежит в плоскости
$(111)_\gamma$, а $[11\br1]_\alpha$~— в плоскости $(011)_\alpha$.
\end{proverka}

% ----------------------------------------------------------------------
%  Итоги главы 1: сводка, ответы, задачи, литература
% ----------------------------------------------------------------------
\clearpage
\section*{Итоги главы}
\addcontentsline{toc}{section}{Итоги главы}
\markright{Итоги главы 1}

Глава прошла путь от внешней формы кристалла к его внутренней геометрии
и от неё~— к свойствам и превращениям. Ключевая мысль: \vazhno{кристалл~—
это решётка плюс базис}, а почти всё остальное (индексы, симметрия,
число независимых констант, системы скольжения, ориентационные
соотношения) следует из геометрии решётки и её симметрии.
Таблица~\ref{tab:ch1-summary} сводит вместе инструменты главы.

\medskip\noindent\begin{minipage}{\textwidth}
\centering
\captionof{table}{Инструменты физической кристаллографии и их применение}
\label{tab:ch1-summary}
\footnotesize
\begin{tabularx}{\textwidth}{@{}>{\raggedright\arraybackslash\hspace{0pt}}p{30mm}>{\raggedright\arraybackslash\hspace{0pt}}X>{\raggedright\arraybackslash\hspace{0pt}}p{40mm}@{}}
\toprule
Инструмент & Что даёт & Где используется \\
\midrule
Стереографическая проекция, сетка Вульфа & углы между направлениями и плоскостями на плоском рисунке & ориентировка монокристаллов, полюсные фигуры, системы скольжения \\
Индексы Миллера (Миллера--Браве), обратная решётка & описание направлений и плоскостей; $d_{hkl}=1/|\vb H_{hkl}|$, закон зон & дифракция, зоны Бриллюэна, все расчёты углов \\
Матрицы поворота & смена систем координат, преобразование индексов & оси симметрии, разориентация зёрен, текстуры \\
Точечные и пространственные группы & 32 класса, 14 решёток Браве, 230 групп & классификация кристаллов, правила отбора свойств \\
Принципы Неймана и Кюри & симметрия свойств и кристалла во внешнем поле & пьезо- и пироэффект, анизотропия \\
Тензоры свойств & число независимых констант по сингониям & проводимость, упругость ($c_{ij}$), тепловое расширение \\
Фактор Шмида, тензор дисторсии & какая система скольжения работает и как поворачивается решётка & пластичность монокристаллов, текстуры \\
РСУ, $\Sigma$, ПРН & строение и энергия границ зёрен & зернограничная инженерия, диффузия по границам \\
Деформация Бейна, $P_1=RPB$ & решётка и форма мартенсита & закалка сталей, сплавы с памятью формы \\
\bottomrule
\end{tabularx}
\end{minipage}

\begin{itogi}[Главные формулы главы]
\begin{itemize}
\item Стереографическая проекция: $\rho=r\operatorname{tg}(\theta/2)$
\eqref{eq:stereo-rho}.
\item Закон зон: $hu+kv+lw=0$; межплоскостное расстояние $d_{hkl}=1/|\vb H_{hkl}|$,
для кубической решётки $d=a/\sqrt{h^2+k^2+l^2}$.
\item Угол между направлениями в кубическом кристалле
$\cos\varphi=\dfrac{u_1u_2+v_1v_2+w_1w_2}{\sqrt{u_1^2+v_1^2+w_1^2}\sqrt{u_2^2+v_2^2+w_2^2}}$
\eqref{eq:cos-cubic}.
\item Четырёхиндексные обозначения: $i=-(h+k)$; $[UVW]\to[uvtw]$ по
\eqref{eq:uvtw}.
\item Угол поворота по следу матрицы: $\operatorname{tr}R=\pm(1+2\cos\varphi)$.
\item Принцип Неймана $G_\text{свойства}\supseteq G_\text{кристалла}$;
принцип Кюри $G_\text{кр+возд}=G_\text{кр}\cap G_\text{возд}$.
\item Закон Гука $\sigma_{ij}=c_{ijkl}\varepsilon_{kl}$; фактор анизотропии
$A=2c_{44}/(c_{11}-c_{12})$ \eqref{eq:zener}; модуль Юнга кубического кристалла
\eqref{eq:E-cubic}.
\item Закон Шмида $\tau=\sigma\cos\chi\cos\lambda$ \eqref{eq:schmid};
тензор пластической дисторсии $u^0_{ij}=\alpha m_ih_j$ \eqref{eq:plastic-distortion}.
\item Двойникование: $s=2\operatorname{ctg}2\phi$ \eqref{eq:twin-shear}.
\item Малоугловая граница $D=b/\theta$; формула Франка \eqref{eq:frank};
энергия $E=E_0\theta(A-\ln\theta)$; критерий Брэндона
$\Delta\theta=15^\circ/\sqrt\Sigma$.
\item Деформация с инвариантной плоскостью $P=I+\vb d\tilde{\vb p}$;
деформация Бейна $\eta_1=\sqrt2a_\alpha/a_\gamma$, $\eta_3=c_\alpha/a_\gamma$.
\end{itemize}
\end{itogi}

%=======================================================================
\section*{Ответы на вопросы для самопроверки}
\addcontentsline{toc}{section}{Ответы на вопросы для самопроверки}
\markright{Ответы на вопросы главы 1}

{\small
\paragraph{1.1 (кристаллическое состояние)}
\begin{enumerate}[label=\arabic*.]
\item Форма (габитус) зависит от условий роста: скорости роста разных
граней различны, и одни грани развиваются сильнее других. Но все грани
параллельны одним и тем же плоскостям решётки, поэтому углы между
соответственными гранями одинаковы (закон постоянства углов).
\item Отношения параметров граней~— отношения небольших целых чисел;
значит, грани проходят через узлы трёхмерной периодической решётки:
кристалл имеет решётчатое строение.
\item Однородность~— одинаковость свойств в параллельно расположенных
точках (симметрия относительно трансляций); изотропия~— одинаковость
свойств по всем направлениям. Кристалл однороден, но анизотропен
(например, модуль Юнга меди вдоль $\nfam{111}$ почти втрое больше, чем
вдоль $\nfam{100}$).
\item Решётка~— геометрическое множество эквивалентных точек; её начало
можно поместить в любую точку структуры. Атомы базиса располагаются
относительно узлов, и в общем случае ни один из них не обязан попадать
в узел.
\item Три атома (A, B, C).
\end{enumerate}

\paragraph{1.2 (основы кристаллографии)}
\begin{enumerate}[label=\arabic*.]
\item Стереографическая проекция плоскости~— её след (дуга большого
круга), гномостереографическая~— полюс (проекция нормали). В
кубическом кристалле направление $[hkl]$ перпендикулярно плоскости
$(hkl)$, поэтому полюс $(hkl)$ совпадает с выходом $[hkl]$.
\item Угол между $[001]$ и нормалью к $(011)$ равен $45^\circ$:
$\rho=r\operatorname{tg}22{,}5^\circ=4{,}14$~см.
\item Обратные отрезки $\frac12:\frac13:(-1)$, умноженные на 6: $(32\br6)$.
\item $[210]\to[10\br10]$, $[1\br10]\to[1\br100]$.
\item Векторное произведение нормалей: $[11\br2]$.
\item $1+2\cos\varphi=-1$, $\varphi=180^\circ$ (ось 2).
\item $d_{110}=a/\sqrt2=0{,}2027$~нм, $d_{111}=a/\sqrt3=0{,}1655$~нм. (В
ОЦК-решётке посередине между плоскостями $(111)$ лежат такие же атомные
плоскости, поэтому реальное расстояние между атомными слоями $\parallel(111)$~—
$d_{222}=0{,}083$~нм; отражение $111$ погасает.)
\item Вторые соседи в ОЦК-решётке (на расстоянии $a$) лишь на 15\,\%
дальше первых ($a\sqrt3/2$), и плоскости $\plfam{100}$, делящие пополам
отрезки до них, срезают вершины октаэдра из плоскостей $\plfam{111}$:
получается усечённый октаэдр (8 шестиугольников и 6 квадратов).
\end{enumerate}

\paragraph{1.3 (симметрия кристаллов)}
\begin{enumerate}[label=\arabic*.]
\item Поворот на $72^\circ$ несовместим с трансляционной периодичностью:
пятиугольниками нельзя без зазоров замостить плоскость. У квазикристаллов
трансляционной периодичности нет (порядок квазипериодический), поэтому
запрет на них не распространяется.
\item Ось 6: $1+2\cos60^\circ=2$; ось $\br4$: $-(1+2\cos90^\circ)=-1$.
\item $\br2=iC_2$: при оси 2 вдоль $z$ матрица $\operatorname{diag}(-1,-1,1)$,
умноженная на $-I$, даёт $\operatorname{diag}(1,1,-1)$~— отражение в
плоскости $\perp z$. $\br6=iC_6$: $(iC_6)^2=C_3$, $(iC_6)^3=iC_2=m_\perp$,
т.\,е. $\br6=3/m$.
\item 48. Подгрупп порядка 5 нет: по теореме Лагранжа порядок подгруппы
делит порядок группы, а 5 не делит 48.
\item Класс $m\br3m$ центросимметричен~— пьезоэффекта нет. Класс $\br43m$
нецентросимметричен и пьезоэлектричен (сфалерит, GaAs).
\item Тетрагональная $C$-решётка~— та же примитивная тетрагональная
решётка $P$ с ячейкой в два раза меньшего объёма, повёрнутой на $45^\circ$
вокруг оси $c$ (рёбра основания $a/\sqrt2$).
\item За поворот на $90^\circ$ ось $4_1$ сдвигает на $c/4$, ось $4_3$~— на
$3c/4$ (эквивалентно $-c/4$): это правый и левый винты, зеркальные
двойники друг друга.
\item Магнитное поле~— аксиальный вектор: $\infty/m$ (вращающийся цилиндр:
центр симметрии и плоскость $\perp\vb H$, но нет продольных плоскостей).
Электрическое поле~— полярный вектор: $\infty mm$ (конус).
\end{enumerate}

\paragraph{1.4 (структура кристаллов)}
\begin{enumerate}[label=\arabic*.]
\item Обе~— плотнейшие упаковки одинаковых плотных слоёв, отличающиеся
лишь порядком чередования слоёв (ABC и AB): плотность заполнения
одинакова, 0,74.
\item ГЦК: одна октаэдрическая и две тетраэдрические пустоты на атом.
ОЦК: три октаэдрические (сплюснутые) и шесть тетраэдрических.
\item В аустените октапустоты крупные ($r=0{,}414R$). В феррите атом
углерода в сплюснутом октаэдре раздвигает лишь два ближайших атома вдоль
оси сплющивания; в тетрапустоте пришлось бы раздвигать четыре атома, что
энергетически невыгодно.
\item Половину (через одну); у каждого атома четыре ближайших соседа в
вершинах тетраэдра~— направленные $sp^3$-связи.
\item Обе построены из тетраэдров; в сфалерите слои упакованы по закону
ABC (кубическая решётка, $F\br43m$), в вюрците~— AB (гексагональная,
$P6_3mc$).
\item Атомы в вершинах (Cs) и в центре (Cl) разные, поэтому трансляции
$\frac12\frac12\frac12$ нет: решётка простая кубическая с базисом из двух
атомов.
\item 12 атомов Fe и 4 атома Al (16 атомов, Fe$_3$Al).
\item Поверхностная реконструкция грани $(111)$ кремния: двумерная ячейка
поверхности в $7\times7$ раз больше ячейки идеальной (нереконструированной)
грани.
\end{enumerate}

\paragraph{1.5 (физические свойства)}
\begin{enumerate}[label=\arabic*.]
\item По принципу Кюри $m\br3m\cap\infty/mmm$ (ось цилиндра $\parallel[110]$)
$=mmm$: остаются оси 2 вдоль $[110]$, $[\br110]$, $[001]$ и
перпендикулярные им плоскости (ромбическая симметрия).
\item Пироэлектрик должен иметь единственное полярное направление, а в
кубических группах четыре оси 3 вдоль $\nfam{111}$ не допускают его.
\item Проводимость~— тензор 2-го ранга, а у кубического кристалла все
такие свойства изотропны; модуль Юнга определяется тензором 4-го ранга,
у которого три независимые компоненты, и анизотропия сохраняется.
\item $E_{\nfam{110}}\approx130$~ГПа (для сравнения $E_{\nfam{100}}=67$,
$E_{\nfam{111}}=190$~ГПа).
\item Два ($\sigma_{11}=\sigma_{22}$ и $\sigma_{33}$): в~\eqref{eq:n-invariants}
вклады $E$, $C_2$, $i$, $\sigma_h$ равны 6, 2, 6, 2, а у $C_4$, $C_4^3$, $S_4$,
$S_4^3$~— нули; $n=16/8=2$, как у всех тетрагональных классов.
\item При $c/a=1{,}86$ базисные слои цинка далеко отстоят друг от друга и
связаны слабо; поэтому кристалл легко сжимается вдоль оси $c$.
\end{enumerate}

\paragraph{1.6 (пластическая деформация)}
\begin{enumerate}[label=\arabic*.]
\item Скольжение сдвигает на вектор трансляции, и узлы попадают в узлы той
же решётки. При двойниковании атомы смещаются на доли трансляции, и
решётка переходит в зеркально-симметричную ориентацию.
\item Когда $\vb n$, $\vb m$ и $F$ лежат в одной плоскости и
$\chi=\lambda=45^\circ$. Больше $\frac12$ быть не может: $\lambda\ge90^\circ-\chi$,
поэтому $\cos\chi\cos\lambda\le\cos\chi\sin\chi\le\frac12$.
\item $\cos\chi=1/\sqrt3$, $\cos\lambda=1/\sqrt2$, фактор $0{,}408$. Такой
же фактор у 8 систем (в каждой из четырёх плоскостей $\plfam{111}$~— два
направления, не перпендикулярные $[001]$).
\item Базисное скольжение даёт лишь две независимые системы (меньше
пяти) и не может изменить размер кристалла вдоль оси $c$: зёрна не могут
согласованно деформироваться, на границах возникают трещины.
\item $\cos2\phi=\frac{1\cdot1+1\cdot1+2\cdot(-2)}{6}=-\frac13$, т.\,е. угол
между плоскостями $70{,}53^\circ$; $s=2\operatorname{ctg}70{,}53^\circ=0{,}707$.
\item Точка в центре проекции (полюс $(001)$) и основной круг (полюса
$(100)$ и $(010)$ при вращении вокруг $F$).
\end{enumerate}

\paragraph{1.7 (границы зёрен)}
\begin{enumerate}[label=\arabic*.]
\item Пять (три~— разворот, два~— нормаль к границе); на атомном уровне
ещё три (вектор $\vb t$), всего восемь.
\item У границы наклона ось разворота лежит в плоскости границы, у границы
кручения перпендикулярна ей. Малоугловая граница наклона~— стенка
краевых дислокаций, кручения~— сетка винтовых.
\item $\theta=b/D=0{,}256/50=5{,}1\cdot10^{-3}$~рад $\approx0{,}29^\circ$.
\item Поворот на $90^\circ$ вокруг $[001]$~— операция симметрии куба, а
$53{,}13^\circ=90^\circ-36{,}87^\circ$: это одна и та же разориентация,
описанная по-разному (и с переставленными зёрнами); совпадающие узлы те же.
\item Граница $\Sigma3$ по $\plfam{111}$~— зеркальная плоскость: у атомов
сохраняются все ближайшие соседи, меняется лишь порядок чередования
плотных слоёв (\ldots ABCACB\ldots); энергия порядка половины энергии
дефекта упаковки.
\item Нет: наибольший угол разориентации кубических кристаллов
$62{,}8^\circ$, и у любого поворота на $70^\circ$ есть эквивалентное описание
с меньшим углом (например, $70{,}5^\circ$ вокруг $\nfam{110}$~— это
$60^\circ$ вокруг $\nfam{111}$, двойник $\Sigma3$).
\end{enumerate}

\paragraph{1.8 (мартенситные превращения)}
\begin{enumerate}[label=\arabic*.]
\item Поверхность превращённой области остаётся плоской (деформация
однородна), а риски~— прямыми и непрерывными, изламываясь только на
следе габитуса (плоскость габитуса не искажается и не поворачивается).
\item Деформация решётки меняет кристаллическую структуру (ГЦК$\to$ОЦТ);
деформация с инвариантной решёткой меняет только форму, не меняя
решётки: скольжение, двойникование.
\item $\eta_1=\sqrt2\cdot0{,}2850/0{,}3584=1{,}125$, $\eta_3=0{,}2958/0{,}3584=0{,}825$;
$\Delta V/V=\eta_1^2\eta_3-1\approx+4{,}4\,\%$.
\item У неё два главных удлинения больше 1, одно меньше и ни одно не
равно 1; неискажённые направления образуют конус, но не плоскость.
\item $1\cdot1+0\cdot1+(-1)\cdot1=0$ и $0\cdot1+1\cdot1+1\cdot(-1)=0$: оба
условия зональности выполнены.
\end{enumerate}
}

%=======================================================================
\section*{Задачи}
\addcontentsline{toc}{section}{Задачи}
\markright{Задачи к главе 1}

Задачи взяты из раздела «Контрольные вопросы, задачи и упражнения»
исходного учебника (номера в скобках) и снабжены ответами.

\begin{enumerate}[label={\sffamily\bfseries\color{fmHead}\arabic*.}]
\item (13) Найдите индексы направлений, проходящих через узлы
$\uzel{000}$ и $\uzel{1\,1\,\br{\tfrac12}}$; $\uzel{\tfrac13\,1\,1}$ и
$\uzel{0\,\tfrac34\,1}$; $\uzel{\br{\tfrac23}\,0\,1}$ и
$\uzel{1\,0\,\tfrac12}$.\\
\emph{Ответ:} $[22\br1]$; $[430]$; $[10\,0\,\br3]$.
\item (15) Найдите индексы плоскостей, отсекающих на осях отрезки
$1,\,-\frac12,\,\frac13$; $-\frac13,\,1,\,\frac12$; $\frac12,\,\frac14,\,\frac12$.\\
\emph{Ответ:} $(1\br23)$; $(\br312)$; $(121)$.
\item (17) Входят ли плоскости $(12\br3)$, $(111)$, $(01\br1)$ в зону с осью
$[111]$?\\
\emph{Ответ:} $(12\br3)$ и $(01\br1)$~— да ($hu+kv+lw=0$), $(111)$~— нет.
\item (19) Найдите ось зоны, содержащей плоскости $(01\br11)$ и $(000\br1)$,
в трёх- и четырёхиндексной записи.\\
\emph{Ответ:} $[100]=[2\br1\br10]$.
\item (20) Найдите в кубическом кристалле углы между $[100]$ и $[110]$;
$[110]$ и $[11\br1]$; $[123]$ и $(123)$; $(100)$ и $(1\br10)$; $(100)$ и
$(1\br1\br1)$.\\
\emph{Ответ:} $45^\circ$; $35{,}26^\circ$; $0^\circ$; $45^\circ$; $54{,}74^\circ$.
\item (26) Найдите угол между базисными векторами примитивной ячейки
ГЦК-решётки.\\
\emph{Ответ:} $60^\circ$ (векторы $\frac a2[110]$, $\frac a2[011]$,
$\frac a2[101]$).
\item (28) Какие индексы будут у $(1\br10)$ и $[1\br10]$ кубического кристалла в
базисе примитивной ячейки ГЦК-решётки~\eqref{eq:prim}?\\
\emph{Ответ:} $h'=Ah$ даёт $(0\,\br{\tfrac12}\,\tfrac12)\to(0\br11)$; вектор
$a[1\br10]=2(\vb p_3-\vb p_2)$, т.\,е. $[0\br22]$, или направление $[0\br11]$.
\item (29--32) Укажите число и индексы граней ячейки Вигнера--Зейтца и
первой зоны Бриллюэна ГЦК- и ОЦК-решёток.\\
\emph{Ответ:} ячейка ВЗ ГЦК~— ромбододекаэдр, 12 граней $\plfam{110}$;
зона Бриллюэна ГЦК~— усечённый октаэдр, 8 граней $\plfam{111}$ и 6 граней
$\plfam{200}$; ячейка ВЗ ОЦК~— усечённый октаэдр, 8 граней $\plfam{111}$ и
6 граней $\plfam{100}$; зона Бриллюэна ОЦК~— ромбододекаэдр, 12 граней
$\plfam{110}$.
\item (34--36) Запишите матрицы поворотов: $C_4$ вокруг $[010]$ на
$90^\circ$ по часовой стрелке; $C_3$ вокруг $[\br111]$ на $120^\circ$
против часовой стрелки; $C_2$ вокруг $[110]$.\\
\emph{Ответ:}
$\begin{pmatrix}0&0&\br1\\0&1&0\\1&0&0\end{pmatrix}$,
$\begin{pmatrix}0&\br1&0\\0&0&1\\\br1&0&0\end{pmatrix}$,
$\begin{pmatrix}0&1&0\\1&0&0\\0&0&\br1\end{pmatrix}$.
\item (37) Найдите $C=C_2[01\br1]\,C_4[001]$ ($C_4$~— на $90^\circ$ против
часовой стрелки) и определите тип и индексы оси.\\
\emph{Ответ:} $C=\begin{pmatrix}0&1&0\\0&0&\br1\\\br1&0&0\end{pmatrix}$,
$\det C=1$, $\operatorname{tr}C=0$: поворотная ось 3 вдоль $[11\br1]$.
\item (38--41) Найдите результат последовательного выполнения
операций: $\sigma_{(1\br10)}\sigma_{(101)}$; $\sigma_{(011)}\,i$;
$C_3[\br111]\,i$; $C_2[101]\,\sigma_{(01\br1)}$.\\
\emph{Ответ:} поворотная ось 3 вдоль $[11\br1]$ (две плоскости под
углом $60^\circ$ дают поворот на $120^\circ$ вокруг линии их пересечения);
поворотная ось 2 вдоль $[011]$; инверсионная ось $\br3$ вдоль $[\br111]$;
инверсионная ось $\br3$ вдоль $[\br111]$ (определитель $-1$, угол
$120^\circ$).
\item (46) Почему в кубической сингонии нет базоцентрированной решётки
Браве?\\
\emph{Ответ:} центрирование одной пары граней несовместимо с осями 3
вдоль $\nfam{111}$, которые переводят грани $(001)$ в $(100)$ и $(010)$;
такая решётка имела бы лишь тетрагональную симметрию.
\item (49, 50, 55) Найдите координационные числа и радиусы первых пяти
координационных сфер в ГЦК, ОЦК-структурах и в алмазе.\\
\emph{Ответ:} ГЦК: $a\sqrt{n/2}$ ($n=1\ldots5$), числа 12, 6, 24, 12, 24.
ОЦК: $a\sqrt3/2$, $a$, $a\sqrt2$, $a\sqrt{11}/2$, $a\sqrt3$; числа 8, 6, 12,
24, 8. Алмаз: $\frac a4\sqrt3$, $\frac a4\sqrt8$, $\frac a4\sqrt{11}$, $a$,
$\frac a4\sqrt{19}$; числа 4, 12, 12, 6, 12.
\item (53) Докажите, что в идеальной ГПУ-структуре $c/a=\sqrt{8/3}=1{,}633$.\\
\emph{Ответ:} атом слоя B лежит над центром правильного треугольника из
атомов слоя A со стороной $a$ на расстоянии $a$ от каждого: высота
правильного тетраэдра $h=a\sqrt{2/3}$, а $c=2h$.
\item (57, 58) Найдите координационные числа и расстояния в структурах
UO$_2$ (флюорит) и CaTiO$_3$ (перовскит).\\
\emph{Ответ:} UO$_2$: U по O~— 8 на $a\sqrt3/4$, O по U~— 4 на
$a\sqrt3/4$. Перовскит (Ti в центре куба): Ca--Ti~— 8 на $a\sqrt3/2$,
Ti--O~— 6 на $a/2$, O--Ca~— 4 на $a/\sqrt2$, Ti--Ca~— 8 на $a\sqrt3/2$,
O--Ti~— 2 на $a/2$, Ca--O~— 12 на $a/\sqrt2$.
\item (59--61) Какие элементы симметрии сохранятся: у класса $\br43m$ при
растяжении вдоль $[010]$; у $m\br3m$ в электрическом поле вдоль $[1\br11]$; у
$6/mmm$ в магнитном поле вдоль $[01\br10]$?\\
\emph{Ответ:} $\br42m$ (ось $\br4$ вдоль $[010]$, оси 2 вдоль $[100]$ и
$[001]$, плоскости $(101)$ и $(10\br1)$); $3m$ (ось 3 вдоль $[1\br11]$ и
три плоскости $\plfam{110}$, проходящие через неё); $2/m$ (ось 2 вдоль
поля, плоскость $\perp$ полю, центр инверсии).
\item (62, 63) Найдите модули Юнга вдоль $\nfam{100}$, $\nfam{110}$,
$\nfam{111}$ монокристаллов Cu ($c_{11}=168$, $c_{12}=121$, $c_{44}=76$~ГПа)
и Nb ($c_{11}=246$, $c_{12}=134$, $c_{44}=29$~ГПа).\\
\emph{Ответ:} Cu: 67, 131, 192~ГПа ($A=3{,}2$, максимум вдоль $\nfam{111}$);
Nb: 152, 93, 82~ГПа ($A=0{,}52$, максимум вдоль $\nfam{100}$).
\item (64) То же для Be вдоль $[2\br1\br10]$, $[1\br100]$, $[\br1011]$ ($c_{11}=292$,
$c_{12}=27$, $c_{13}=14$, $c_{33}=336$, $c_{44}=162$~ГПа, $c/a=1{,}57$).\\
\emph{Ответ:} в базисной плоскости гексагональный кристалл упруго
изотропен: $E=289$~ГПа вдоль обоих направлений; вдоль $[\br1011]$
(угол $47{,}8^\circ$ с осью $c$) $E\approx321$~ГПа; вдоль $[0001]$~— 335~ГПа.
\item (65) ГЦК-кристалл растянут напряжением $\sigma$ вдоль $[\br132]$. Найдите
касательные напряжения в первичной, сопряжённой, критической и
поперечной системах.\\
\emph{Ответ:} ось лежит в треугольнике $[010]$--$[011]$--$[\br111]$;
первичная $(111)[\br110]$: $0{,}467\sigma$; сопряжённая $(\br11\br1)[011]$:
$0{,}292\sigma$; критическая $(\br111)[110]$: $0{,}350\sigma$; поперечная
$(11\br1)[\br110]$: 0 (ось лежит в этой плоскости).
\item (67) Первичная система ГЦК-кристалла $(111)[10\br1]$. Какова
поперечная система?\\
\emph{Ответ:} $(1\br11)[10\br1]$~— вторая плоскость $\plfam{111}$, содержащая
то же направление скольжения.
\item (70--73) Сколько независимых компонент у тензора пластической
дисторсии? Сколько систем скольжения $\plfam{112}\nfam{111}$ в ОЦК-металлах,
базисных и призматических систем в ГПУ-металлах?\\
\emph{Ответ:} 8 (след равен нулю: 5~— деформация, 3~— поворот);
12; 3; 3.
\item (78) Текстура прокатки $\plfam{101}\nfam{112}$. Каково поперечное
направление?\\
\emph{Ответ:} $\text{ПН}=\text{НН}\times\text{НП}$, например
$(101)\times[\br121]\parallel[\br1\br11]$, т.\,е. $\nfam{111}$.
\item (80) Найдите углы Эйлера $\theta$, $\psi$ стержня ФРО для аксиальной
текстуры $\nfam{mnp}$.\\
\emph{Ответ:} $\cos\theta=p/\sqrt{m^2+n^2+p^2}$, $\operatorname{tg}\psi=m/n$
\eqref{eq:fiber-euler}.
\item (81, 82) Найдите направление линии краевой дислокации: в
ГЦК-кристалле с $\vb b\parallel[110]$ в плоскости $\plfam{111}$; в ОЦК с
$\vb b\parallel[111]$ в плоскости $\plfam{112}$.\\
\emph{Ответ:} $\vb l\parallel\vb b\times\vb n$: в плоскости $(1\br11)$~—
$[1\br1\br2]$, в $(\br111)$~— $[1\br12]$ (тип $\nfam{112}$); в ОЦК в
плоскости $(11\br2)$~— $[\br110]$ (тип $\nfam{110}$).
\item (83) На какие частичные дислокации Шокли расщепляется полная
дислокация $\frac a2[1\br10]$ в плоскости $(111)$?\\
\emph{Ответ:} $\frac a2[1\br10]\to\frac a6[2\br1\br1]+\frac a6[1\br21]$.
\item (87, 88) Найдите расстояние между дислокациями в симметричной
границе наклона в Cu при $\theta=1{,}5^\circ$ ($a=0{,}362$~нм) и угол
разворота в W при расстоянии 2,5~мкм ($a=0{,}316$~нм).\\
\emph{Ответ:} Cu: $b=a/\sqrt2=0{,}256$~нм, $D=b/\theta=9{,}8$~нм;
W: $b=a\sqrt3/2=0{,}274$~нм, $\theta=1{,}1\cdot10^{-4}$~рад $\approx0{,}006^\circ$.
\item (89) Найдите плоскость малоугловой границы кручения из винтовых
дислокаций с $\vb b_1=\frac a2[101]$ и $\vb b_2=\frac a2[110]$.\\
\emph{Ответ:} обе дислокации лежат в границе: $\vb n\parallel\vb b_1\times\vb b_2$,
плоскость $(\br111)$.
\item (92) Найдите матрицу деформации Бейна при превращении аустенита
($a=0{,}3584$~нм) в мартенсит ($a=0{,}2850$, $c=0{,}2958$~нм).\\
\emph{Ответ:} $B=\operatorname{diag}(1{,}125;\ 1{,}125;\ 0{,}825)$.
\item (93) Найдите углы между направлениями $[101]$ и $[11\br1]$, $[011]$ и
$[\br101]$ аустенита ($a=0{,}3610$~нм) после превращения в мартенсит
($a=0{,}2845$, $c=0{,}3020$~нм; ось $c$ вдоль $[001]_\gamma$).\\
\emph{Ответ:} $\eta_1=1{,}115$, $\eta_3=0{,}837$; углы меняются с
$90^\circ$ на $77{,}4^\circ$ и с $60^\circ$ на $68{,}9^\circ$.
\end{enumerate}

\paragraph{Вопросы для повторения (по исходному учебнику)}
{\small
\begin{enumerate}[label=\arabic*.]
\item Что такое кристалл в современном понимании?
\item Что такое габитус кристалла, соответственные плоскости кристалла?
\item Сформулируйте закон постоянства углов и закон рациональных
параметров.
\item Какие свойства кристалла являются макроскопическими, какие~—
микроскопическими?
\item Что такое гомологические точки кристалла? Сформулируйте
математические постулаты кристаллического пространства.
\item Что такое шар однородности и шар дискретности?
\item В чём отличие решётки и структуры кристалла?
\item На стереографической проекции показана зона плоскостей. Как будет
выглядеть та же зона на гномостереографической проекции?
\item Что такое кристаллографические индексы направления и плоскости?
\item Покажите в элементарной ячейке плоскости $(321)$, $(\br120)$; в
гексагональной ячейке~— плоскости $(1\br102)$, $(11\br20)$ и направления
$[1\br102]$, $[11\br20]$.
\item Совпадает ли точка первого рода $(011)$ с точкой второго рода $[011]$
в тетрагональном кристалле? Докажите, что в гексагональном кристалле
$\plfam{hki0}$ и $\nfam{hki0}$ совпадают.
\item Почему в кристаллах нет поворотных осей пятого порядка?
\item Перечислите групповые постулаты. Что такое генератор группы? Найдите
обратный элемент для $a^t$ циклической группы порядка $n$.
\item Какой физический смысл имеет вхождение элементов симметрии в один
класс сопряжённых элементов?
\item Покажите плоскость скольжения $n$-типа в ОЦК-решётке, $c$-типа в
ГЦК-решётке, $d$-типа и винтовую ось $4_1$ в алмазе.
\item Какая решётка Браве у ГПУ-структуры $\alpha$-Zr? Обладает ли
сфалерит центром симметрии?
\item Изобразите ППФ $(211)$ для текстуры $F\parallel\nfam{110}$; ОПФ направления
$F$ и направления $G\perp F$ для текстуры $\nfam{112}$; ФРО для текстуры
$\nfam{112}$.
\item Найдите выражения для касательного напряжения в системе
$(1\br1\br1)[10\br1]$ кубического кристалла при двухосном нагружении ($\sigma_1$
вдоль $[100]$, $\sigma_2$ вдоль $[001]$).
\end{enumerate}
}

%=======================================================================
\section*{Литература}
\addcontentsline{toc}{section}{Литература к главе 1}
\markright{Литература к главе 1}

{\small
\paragraph{Использованная в исходном учебнике}
\begin{enumerate}[label=\arabic*.]
\item Вайнштейн Б.\,К. Современная кристаллография. Т.~1. Симметрия
кристаллов. Методы структурной кристаллографии.~— М.: Наука, 1979.~— 384~с.
\item Вайнштейн Б.\,К., Фридкин В.\,М., Инденбом В.\,Л. Современная
кристаллография. Т.~2. Структура кристаллов.~— М.: Наука, 1979.~— 312~с.
\item Шувалов Л.\,А., Урусовская А.\,А., Желудев И.\,С. и др. Современная
кристаллография. Т.~4. Физические свойства кристаллов.~— М.: Наука,
1981.~— 496~с.
\item Кристиан Дж. Теория превращений в металлах и сплавах. Ч.~1.~— М.:
Мир, 1978.~— 806~с.
\item Васильев Д.\,М. Физическая кристаллография.~— М.: Металлургия,
1972.~— 280~с.; 2-е изд.~— 1981.~— 256~с.
\item Русаков А.\,А. Рентгенография металлов.~— М.: Атомиздат, 1977.~— 480~с.
\item Русаков А.\,А., Яльцев В.\,Н. Основы рентгенографии металлов. Ч.~I.~—
М.: МИФИ, 1994.~— 172~с.
\item Чупрунов Е.\,В., Хохлов А.\,Ф., Фаддеев М.\,А. Основы
кристаллографии.~— М.: Физматлит, 2004.~— 500~с.
\item Уманский Я.\,С., Скаков Ю.\,А., Иванов А.\,Н., Расторгуев Л.\,Н.
Кристаллография, рентгенография и электронная микроскопия.~— М.:
Металлургия, 1982.~— 632~с.
\item Шаскольская М.\,П. Кристаллография.~— М.: Высшая школа, 1976.~— 391~с.
\item Штремель М.\,А. Прочность сплавов. Ч.~I. Дефекты решётки.~— М.:
Металлургия, 1982.~— 280~с.; Ч.~II. Деформация.~— М.: МИСиС, 1997.~— 527~с.
\item Новиков И.\,И., Розин К.\,М. Кристаллография и дефекты
кристаллической решётки.~— М.: Металлургия, 1990.~— 336~с.
\item Орлов А.\,Н., Перевезенцев В.\,Н., Рыбин В.\,В. Границы зёрен в
металлах.~— М.: Металлургия, 1980.~— 156~с.
\item Келли А., Гровс Г. Кристаллография и дефекты в кристаллах.~— М.:
Мир, 1974.~— 496~с.
\item Физическое металловедение / под ред. Р.\,У.~Кана. Вып.~1--3.~— М.:
Мир, 1967--1968.
\item Физическое металловедение / под ред. Р.\,У.~Кана, П.~Хаазена. 3-е
изд. В 3~т.~— М.: Металлургия, 1987.
\item Пенкаля Т. Очерки кристаллохимии.~— Л.: Химия, 1974.~— 496~с.
\item Баррет Ч.\,С., Масальский Т.\,Б. Структура металлов. В 2~ч.~— М.:
Металлургия, 1984.~— 352~с.
\item Вустер У. Применение тензоров и теории групп для описания
физических свойств кристаллов.~— М.: Мир, 1977.~— 384~с.
\end{enumerate}

\paragraph{Дополнительно}
\begin{enumerate}[label=\arabic*., resume]
\item Най Дж. Физические свойства кристаллов и их описание при помощи
тензоров и матриц.~— М.: Мир, 1967.
\item Сиротин Ю.\,И., Шаскольская М.\,П. Основы кристаллофизики.~— М.:
Наука, 1979.
\item Вассерман Г., Гревен И. Текстуры металлических материалов.~— М.:
Металлургия, 1969.
\item Кайбышев О.\,А., Валиев Р.\,З. Границы зёрен и свойства металлов.~—
М.: Металлургия, 1987.
\item Курдюмов Г.\,В., Утевский Л.\,М., Энтин Р.\,И. Превращения в железе и
стали.~— М.: Наука, 1977.
\end{enumerate}
}

